\chapter{The Lean Cone}
\label{app:lean-cone}


Throughout this book, claims that lean on a mechanized result cite
the result by its Lean name, set in typewriter face. This appendix
collects those citations into one place: for every cited declaration,
it reproduces the declaration's statement from the \path{lean/}
trees of \path{maitria-coalition/qtsl}, with its file and line. We call this set the \emph{Lean cone} of the book: the set
of mechanized results the book's prose binds to, closed under
nothing --- what you see is what the text actually cites, re-derived
from the sources on every build, so the listing cannot drift from
the tree it describes.

A word on how to read the entries. Statements are reproduced at
signature level: a theorem up to its proof, a definition up to its
body, a structure up to its fields (an ellipsis marks the elision),
and long source lines re-wrapped to the measure of the page.
Each entry is shown in the namespace where it lives, and surrounding
context (\texttt{open} and \texttt{variable} declarations,
notation) lives in the cited file, so an entry here reads as a
faithful excerpt rather than a self-contained module. To verify an
entry, clone the repository and run \texttt{lake build} inside the
relevant Lake package under \path{lean/} (for example
\path{lean/qtsl} or \path{lean/su-functor}); the toolchain is pinned
by each package's \path{lean-toolchain} file, and a build that
succeeds has kernel-checked every statement listed here, since none
of the trees admits \texttt{sorry}.

Two boundaries of the cone, so that its edges are as legible as its
contents. First, the Rule Catalogue displays a Lean or HOL4
statement beside each rule already, extracted from the mechanized
rule sources and hash-bound at emission; this appendix lists the
catalogue's citations that resolve into the \path{lean/} trees of
this repository and defers to the catalogue's own apparatus for its
rule-source and HOL4 layers. Second, module citations (a dotted
name denoting a whole file or directory of results) are listed per
chapter without statements; the module's file is the reference.

The listing follows the book's chapter order, and each declaration
appears once, under the first chapter that cites it; later chapters
that cite the same result say so in a closing line rather than
repeating the statement.


\section{The Semantic Universe}
Citations from Chapter~\ref{ch:qtsl-semantic-universe}.

\par\medskip\noindent{\small \texttt{Qtsl.DiamondSite.}\texttt{\bfseries ChronologyKit}}\hfill{\footnotesize \path{lean/qtsl-aux/QtslAux/DiamondSite.lean}, line 79}\par\nopagebreak
\begin{leansrc}
structure ChronologyKit (X : Type*) [PartialOrder X] where …
\end{leansrc}
\par\medskip\noindent{\small \texttt{Qtsl.DiamondSite.}\texttt{\bfseries minkowskiKit}}\hfill{\footnotesize \path{lean/qtsl-aux/QtslAux/DiamondTopology.lean}, line 283}\par\nopagebreak
\begin{leansrc}
def minkowskiKit : ChronologyKit (Mink E) where …
\end{leansrc}
\par\medskip\noindent{\small \texttt{Qtsl.DiamondSite.}\texttt{\bfseries diamondTopology}}\hfill{\footnotesize \path{lean/qtsl-aux/QtslAux/DiamondTopology.lean}, line 114}\par\nopagebreak
\begin{leansrc}
def diamondTopology (K : ChronologyKit X) :
    GrothendieckTopology (Diamond X) where …
\end{leansrc}
\par\medskip\noindent{\small \texttt{Qtsl.DiamondSite.}\texttt{\bfseries mem\_of\_not\_thick}}\hfill{\footnotesize \path{lean/qtsl-aux/QtslAux/DiamondTopology.lean}, line 142}\par\nopagebreak
\begin{leansrc}
theorem mem_of_not_thick {D : Diamond X} (h : ¬ Thick K D)
    (S : Sieve D) : S ∈ diamondTopology K D
\end{leansrc}
\par\medskip\noindent{\small \texttt{Qtsl.DiamondSite.}\texttt{\bfseries mem\_diamondTopology\_degenerate\_iff}}\hfill{\footnotesize \path{lean/qtsl-aux/QtslAux/DiamondTopology.lean}, line 158}\par\nopagebreak
\begin{leansrc}
theorem mem_diamondTopology_degenerate_iff
    (hK : ∀ a b : X, K.chron a b ↔ a ≤ b) {D : Diamond X} {S : Sieve D} :
    S ∈ diamondTopology K D ↔ S = ⊤
\end{leansrc}
\par\medskip\noindent{\small \texttt{Qtsl.DiamondSite.}\texttt{\bfseries int\_lt\_no\_chronologyKit}}\hfill{\footnotesize \path{lean/qtsl-aux/QtslAux/DiamondTopology.lean}, line 383}\par\nopagebreak
\begin{leansrc}
theorem int_lt_no_chronologyKit :
    ¬ ∃ K : ChronologyKit ℤ, ∀ a b : ℤ, K.chron a b ↔ a < b
\end{leansrc}
\par\medskip\noindent{\small \texttt{SuFunctor.Site.TTT.}\texttt{\bfseries anchorQuot}}\hfill{\footnotesize \path{lean/su-functor/SuFunctor/Site/TemporalTypeTheory.lean}, line 156}\par\nopagebreak
\begin{leansrc}
def anchorQuot : Window ⥤ IntervalBase where …
\end{leansrc}
\par\medskip\noindent{\small \texttt{SuFunctor.Site.TTT.}\texttt{\bfseries anchorTrajIso}}\hfill{\footnotesize \path{lean/su-functor/SuFunctor/Site/TemporalTypeTheory.lean}, line 323}\par\nopagebreak
\begin{leansrc}
def anchorTrajIso : anchorQuot.op ⋙ ssTraj X ≅ windowTraj X
\end{leansrc}
\par\medskip\noindent{\small \texttt{SuFunctor.Site.}\texttt{\bfseries ttObject}}\hfill{\footnotesize \path{lean/su-functor/SuFunctor/Site/BehaviorSpec.lean}, line 246}\par\nopagebreak
\begin{leansrc}
def ttObject (U : ι → Set T) {A : WireSpan El} (B : WBehavior U (states A)) :
    SiteSpecTotal U El
\end{leansrc}
\par\medskip\noindent{\small \texttt{SuFunctor.Site.}\texttt{\bfseries satAt\_specOfBehavior\_iff}}\hfill{\footnotesize \path{lean/su-functor/SuFunctor/Site/BehaviorSpec.lean}, line 208}\par\nopagebreak
\begin{leansrc}
theorem satAt_specOfBehavior_iff (U : ι → Set T) (B : WBehavior U X) (j : ι)
    (ν : Joint (Traj (U j) X)) :
    SatAt ν (specOfBehavior U B j) ↔ (↑ν.wt.support : Set (Traj (U j) X)) ⊆ B.cls j
\end{leansrc}
\par\medskip\noindent{\small \texttt{SuFunctor.Site.}\texttt{\bfseries wSpecMap\_specOfBehavior}}\hfill{\footnotesize \path{lean/su-functor/SuFunctor/Site/BehaviorSpec.lean}, line 256}\par\nopagebreak
\begin{leansrc}
theorem wSpecMap_specOfBehavior (U : ι → Set T) {A B : WireSpan El}
    (φ : A ⟶ B) (Bh : WBehavior U (states A)) :
    wSpecMap U φ (specOfBehaviorW U Bh) =
      specOfBehaviorW U (Bh.pull (res φ))
\end{leansrc}
\par\medskip\noindent{\small \texttt{SuFunctor.Site.TTT.}\texttt{\bfseries mem\_interiorTopology\_iff\_derived}}\hfill{\footnotesize \path{lean/su-functor/SuFunctor/Site/CausalConcordance.lean}, line 110}\par\nopagebreak
\begin{leansrc}
theorem mem_interiorTopology_iff_derived {W : Window} (S : Sieve W) :
    S ∈ interiorTopology W ↔
      S.functorPushforward windowToDiamond ∈
        diamondTopology (linearOrderKit ℝ) W.toDiamond
\end{leansrc}
\par\medskip\noindent{\small \texttt{SuFunctor.Site.TTT.}\texttt{\bfseries wavy\_mem\_iff\_derived}}\hfill{\footnotesize \path{lean/su-functor/SuFunctor/Site/CausalConcordance.lean}, line 133}\par\nopagebreak
\begin{leansrc}
theorem wavy_mem_iff_derived {W : Window} (S : Sieve W) :
    S.functorPushforward anchorQuot ∈ wavyTopology ⟨W.len⟩ ↔
      S.functorPushforward windowToDiamond ∈
        diamondTopology (linearOrderKit ℝ) W.toDiamond
\end{leansrc}
\par\medskip\noindent{\small \texttt{Qtsl.QuantumStates.}\texttt{\bfseries recordShadowQ\_separates\_bell\_mixed}}\hfill{\footnotesize \path{lean/qtsl-aux/QtslAux/QuantumStates.lean}, line 939}\par\nopagebreak
\begin{leansrc}
theorem recordShadowQ_separates_bell_mixed :
    recordShadowQ zzRecord zzRecord_mem_top
        (stateClassMk qubitNet Dstar ⟨bell, bell_isState⟩) ≠
      recordShadowQ zzRecord zzRecord_mem_top
        (stateClassMk qubitNet Dstar ⟨mixed, mixed_isState⟩)
\end{leansrc}
\par\medskip\noindent{\small \texttt{Qtsl.DiamondSite.}\texttt{\bfseries minkowski\_diamond\_spacelike\_pair}}\hfill{\footnotesize \path{lean/qtsl-aux/QtslAux/DiamondTopology.lean}, line 245}\par\nopagebreak
\begin{leansrc}
theorem minkowski_diamond_spacelike_pair {x₀ : E} (hx : x₀ ≠ 0)
    (hx1 : ‖x₀‖ ≤ 1) :
    ∃ p q : Mink E,
      ((⟨0, 0⟩ : Mink E) ≤ p ∧ p ≤ ⟨2, 0⟩) ∧
      ((⟨0, 0⟩ : Mink E) ≤ q ∧ q ≤ ⟨2, 0⟩) ∧
      ¬ p ≤ q ∧ ¬ q ≤ p
\end{leansrc}
\par\medskip\noindent{\small \texttt{SuFunctor.Universe.FellCloseds.}\texttt{\bfseries compactSpace}}\hfill{\footnotesize \path{lean/su-functor/SuFunctor/Universe/FellTopology.lean}, line 268}\par\nopagebreak
\begin{leansrc}
instance compactSpace : CompactSpace (FellCloseds α)
\end{leansrc}
\par\medskip\noindent{\small \texttt{SuFunctor.Universe.}\texttt{\bfseries fell\_eq\_vietoris}}\hfill{\footnotesize \path{lean/su-functor/SuFunctor/Universe/FellTopology.lean}, line 181}\par\nopagebreak
\begin{leansrc}
theorem fell_eq_vietoris [CompactSpace α] [T2Space α] :
    SuFunctor.Universe.fell α = TopologicalSpace.vietoris α
\end{leansrc}
\par\medskip\noindent{\small \texttt{SuFunctor.Universe.fell.}\texttt{\bfseries continuous\_restrictOpen}}\hfill{\footnotesize \path{lean/su-functor/SuFunctor/Universe/FellRestrict.lean}, line 81}\par\nopagebreak
\begin{leansrc}
theorem continuous_restrictOpen (hU : IsOpen U) :
    Continuous (restrictOpen U)
\end{leansrc}
\par\medskip\noindent{\small \texttt{SuFunctor.Universe.fell.}\texttt{\bfseries continuous\_image\_of\_isProperMap}}\hfill{\footnotesize \path{lean/su-functor/SuFunctor/Universe/FellRestrict.lean}, line 130}\par\nopagebreak
\begin{leansrc}
theorem continuous_image_of_isProperMap (hf : IsProperMap f) :
    Continuous (Set.image f : Set α → Set β)
\end{leansrc}
\par\medskip\noindent{\small \texttt{SuFunctor.Universe.FellCloseds.}\texttt{\bfseries projLimHomeomorph}}\hfill{\footnotesize \path{lean/su-functor/SuFunctor/Universe/FellDescent.lean}, line 315}\par\nopagebreak
\begin{leansrc}
def projLimHomeomorph (U : ι → Set α) (hUo : ∀ i, IsOpen (U i))
    (hdir : Directed (· ⊆ ·) U) (hcov : ⋃ i, U i = univ) :
    FellCloseds α ≃ₜ ProjLim U where …
\end{leansrc}
\par\medskip\noindent{\small \texttt{SuFunctor.Universe.FellCloseds.IsLocalLocus.}\texttt{\bfseries isClosed}}\hfill{\footnotesize \path{lean/su-functor/SuFunctor/Universe/FellRecordLocus.lean}, line 99}\par\nopagebreak
\begin{leansrc}
theorem IsLocalLocus.isClosed (hloc : IsLocalLocus U R₀ R)
    (hUo : ∀ i, IsOpen (U i)) (hRcl : ∀ i, IsClosed (R i)) :
    IsClosed R₀
\end{leansrc}
\par\medskip\noindent{\small \texttt{SuFunctor.Universe.FellCloseds.IsLocalLocus.}\texttt{\bfseries isCompact}}\hfill{\footnotesize \path{lean/su-functor/SuFunctor/Universe/FellRecordLocus.lean}, line 113}\par\nopagebreak
\begin{leansrc}
theorem IsLocalLocus.isCompact (hloc : IsLocalLocus U R₀ R)
    (hUo : ∀ i, IsOpen (U i)) (hRcl : ∀ i, IsClosed (R i)) :
    IsCompact R₀
\end{leansrc}
\par\medskip\noindent{\small \texttt{SuFunctor.Universe.FellCloseds.IsLocalLocusEq.}\texttt{\bfseries glue\_mem}}\hfill{\footnotesize \path{lean/su-functor/SuFunctor/Adequacy/SolutionSetMesh.lean}, line 125}\par\nopagebreak
\begin{leansrc}
theorem IsLocalLocusEq.glue_mem {U : ι → Set α} {R₀ : Set (FellCloseds α)}
    {R : ∀ i, Set (FellCloseds ↥(U i))} (hloc : IsLocalLocusEq U R₀ R)
    (hUo : ∀ i, IsOpen (U i)) (hcov : ⋃ i, U i = univ)
    {p : Equalizer U} (hp : ∀ i, p.1 i ∈ R i) :
    glueEqualizer U hUo hcov p ∈ R₀
\end{leansrc}
\par\medskip\noindent{\small \texttt{SuFunctor.Universe.FellCloseds.}\texttt{\bfseries speedLocus}}\hfill{\footnotesize \path{lean/su-functor/SuFunctor/Universe/FellRecordLocus.lean}, line 234}\par\nopagebreak
\begin{leansrc}
def speedLocus : Set (FellCloseds β)
\end{leansrc}
\par\medskip\noindent{\small \texttt{SuFunctor.Universe.FellCloseds.}\texttt{\bfseries isClosed\_speedLocus}}\hfill{\footnotesize \path{lean/su-functor/SuFunctor/Universe/FellRecordLocus.lean}, line 259}\par\nopagebreak
\begin{leansrc}
theorem isClosed_speedLocus (hτ : Continuous τ) (hσ : Continuous σ)
    (c : ℝ≥0) : IsClosed (speedLocus τ σ c)
\end{leansrc}
\par\medskip\noindent{\small \texttt{SuFunctor.Universe.FellCloseds.}\texttt{\bfseries isSpeedGraph\_iff\_forall\_clip}}\hfill{\footnotesize \path{lean/su-functor/SuFunctor/Universe/FellRecordLocus.lean}, line 308}\par\nopagebreak
\begin{leansrc}
theorem isSpeedGraph_iff_forall_clip (hdir : Directed (· ⊆ ·) U)
    (hcov : ⋃ i, U i = univ) (F : FellCloseds β) :
    IsSpeedGraph τ σ c (F : Set β)
      ↔ ∀ i, FellCloseds.clip (U i) F
          ∈ speedLocus (τ ∘ (↑)) (σ ∘ (↑)) c
\end{leansrc}
\par\medskip\noindent{\small \texttt{SuFunctor.Universe.FellCloseds.}\texttt{\bfseries speedLocusHomeomorph}}\hfill{\footnotesize \path{lean/su-functor/SuFunctor/Universe/FellRecordLocus.lean}, line 346}\par\nopagebreak
\begin{leansrc}
noncomputable def speedLocusHomeomorph (hUo : ∀ i, IsOpen (U i))
    (hdir : Directed (· ⊆ ·) U) (hcov : ⋃ i, U i = univ) :
    ↥(speedLocus τ σ c)
      ≃ₜ ↥(FellCloseds.levelLocus U
          fun _ => speedLocus (τ ∘ (↑)) (σ ∘ (↑)) c)
\end{leansrc}
\par\medskip\noindent{\small \texttt{SuFunctor.Universe.FellCloseds.}\texttt{\bfseries bot\_mem\_speedLocus}}\hfill{\footnotesize \path{lean/su-functor/SuFunctor/Universe/FellRecordLocus.lean}, line 242}\par\nopagebreak
\begin{leansrc}
theorem bot_mem_speedLocus : (⊥ : FellCloseds β) ∈ speedLocus τ σ c
\end{leansrc}
\par\medskip\noindent{\small \texttt{SuFunctor.Universe.FellCloseds.}\texttt{\bfseries isLocallyClosed\_speedLocus}}\hfill{\footnotesize \path{lean/su-functor/SuFunctor/Universe/FellRecordLocus.lean}, line 294}\par\nopagebreak
\begin{leansrc}
theorem isLocallyClosed_speedLocus (hτ : Continuous τ)
    (hσ : Continuous σ) (c : ℝ≥0) :
    IsLocallyClosed (speedLocus τ σ c)
\end{leansrc}
\par\medskip\noindent{\small \texttt{SuFunctor.Universe.FellCloseds.}\texttt{\bfseries isSpeedGraph\_iff\_exists\_lipschitzOnWith}}\hfill{\footnotesize \path{lean/su-functor/SuFunctor/Universe/FellRecordLocus.lean}, line 440}\par\nopagebreak
\begin{leansrc}
theorem isSpeedGraph_iff_exists_lipschitzOnWith [CompleteSpace E]
    [Nonempty E] {F : FellCloseds (T × E)} :
    IsSpeedGraph Prod.fst Prod.snd c (F : Set (T × E))
      ↔ ∃ D : Set T, IsClosed D ∧ ∃ f : T → E, LipschitzOnWith c f D ∧
          (F : Set (T × E)) = (fun t => (t, f t)) '' D
\end{leansrc}
\par\medskip\noindent{\small \texttt{SuFunctor.Universe.FellCloseds.}\texttt{\bfseries topology\_onePoint\_eq\_vietoris}}\hfill{\footnotesize \path{lean/su-functor/SuFunctor/Universe/FellOnePoint.lean}, line 319}\par\nopagebreak
\begin{leansrc}
theorem topology_onePoint_eq_vietoris [WeaklyLocallyCompactSpace α] [T2Space α] :
    (FellCloseds.topology : TopologicalSpace (FellCloseds (OnePoint α)))
      = .induced (↑) (TopologicalSpace.vietoris (OnePoint α))
\end{leansrc}
\par\medskip\noindent{\small \texttt{SuFunctor.Universe.FellCloseds.}\texttt{\bfseries nablaGen}}\hfill{\footnotesize \path{lean/su-functor/SuFunctor/Universe/FellCapacity.lean}, line 75}\par\nopagebreak
\begin{leansrc}
def nablaGen (B : ι → Set Y) (hits misses : Finset ι) : Set (FellCloseds Y)
\end{leansrc}
\par\medskip\noindent{\small \texttt{SuFunctor.Universe.FellCloseds.}\texttt{\bfseries isTopologicalBasis\_nablaGen}}\hfill{\footnotesize \path{lean/su-functor/SuFunctor/Universe/FellCapacity.lean}, line 136}\par\nopagebreak
\begin{leansrc}
theorem isTopologicalBasis_nablaGen [LocallyCompactSpace Y] [T2Space Y]
    {B : ι → Set Y} (hB : IsTopologicalBasis (Set.range B))
    (hBcl : ∀ i, IsCompact (closure (B i))) :
    IsTopologicalBasis
      (Set.range fun p : Finset ι × Finset ι => nablaGen B p.1 p.2)
\end{leansrc}
\par\medskip\noindent{\small \texttt{SuFunctor.Universe.FellCloseds.}\texttt{\bfseries measure\_ext\_of\_nablaGen}}\hfill{\footnotesize \path{lean/su-functor/SuFunctor/Universe/FellCapacity.lean}, line 278}\par\nopagebreak
\begin{leansrc}
theorem measure_ext_of_nablaGen [LocallyCompactSpace Y] [T2Space Y] [Countable ι]
    {B : ι → Set Y} (hB : IsTopologicalBasis (Set.range B))
    (hBcl : ∀ i, IsCompact (closure (B i)))
    [MeasurableSpace (FellCloseds Y)] [BorelSpace (FellCloseds Y)]
    {μ ν : Measure (FellCloseds Y)} [IsFiniteMeasure μ]
    (h : ∀ hits misses : Finset ι,
      μ (nablaGen B hits misses) = ν (nablaGen B hits misses)) :
    μ = ν
\end{leansrc}
\par\medskip\noindent{\small \texttt{SuFunctor.Universe.FellCloseds.}\texttt{\bfseries exists\_nablaGen\_measure\_ne}}\hfill{\footnotesize \path{lean/su-functor/SuFunctor/Universe/FellCapacity.lean}, line 302}\par\nopagebreak
\begin{leansrc}
theorem exists_nablaGen_measure_ne [LocallyCompactSpace Y] [T2Space Y] [Countable ι]
    {B : ι → Set Y} (hB : IsTopologicalBasis (Set.range B))
    (hBcl : ∀ i, IsCompact (closure (B i)))
    [MeasurableSpace (FellCloseds Y)] [BorelSpace (FellCloseds Y)]
    {μ ν : Measure (FellCloseds Y)} [IsFiniteMeasure μ]
    (hμν : μ ≠ ν) :
    ∃ hits misses : Finset ι,
      μ (nablaGen B hits misses) ≠ ν (nablaGen B hits misses)
\end{leansrc}
\par\medskip\noindent{\small \texttt{SuFunctor.Universe.FellCloseds.}\texttt{\bfseries scottContinuous\_ext\_of\_nablaGen}}\hfill{\footnotesize \path{lean/su-functor/SuFunctor/Universe/FellCapacity.lean}, line 326}\par\nopagebreak
\begin{leansrc}
theorem scottContinuous_ext_of_nablaGen [LocallyCompactSpace Y] [T2Space Y]
    {B : ι → Set Y} (hB : IsTopologicalBasis (Set.range B))
    (hBcl : ∀ i, IsCompact (closure (B i)))
    {β : Type*} [PartialOrder β] {f g : Opens (FellCloseds Y) → β}
    (hf : ScottContinuous f) (hg : ScottContinuous g)
    (h : ∀ U : Opens (FellCloseds Y),
      (∃ s : Finset (Finset ι × Finset ι),
        (U : Set (FellCloseds Y)) = ⋃ p ∈ s, nablaGen B p.1 p.2) →
      f U = g U) :
    f = g
\end{leansrc}
\par\medskip\noindent{\small \texttt{SuFunctor.Universe.FellCloseds.}\texttt{\bfseries not\_isOpen\_disjoint\_univ}}\hfill{\footnotesize \path{lean/su-functor/SuFunctor/Universe/FellCapacity.lean}, line 375}\par\nopagebreak
\begin{leansrc}
theorem not_isOpen_disjoint_univ [T1Space Y] (hY : ¬CompactSpace Y) :
    ¬IsOpen {F : FellCloseds Y | Disjoint (F : Set Y) univ}
\end{leansrc}
\par\medskip\noindent{\small \texttt{SuFunctor.Universe.FellCloseds.}\texttt{\bfseries exists\_scottContinuous\_eq\_on\_lower\_ne}}\hfill{\footnotesize \path{lean/su-functor/SuFunctor/Universe/FellCapacity.lean}, line 401}\par\nopagebreak
\begin{leansrc}
theorem exists_scottContinuous_eq_on_lower_ne [T1Space Y] [Nonempty Y] :
    ∃ f g : Opens (FellCloseds Y) → ℝ≥0∞,
      ScottContinuous f ∧ ScottContinuous g ∧
      (∀ U : Opens (FellCloseds Y),
        IsOpen[TopologicalSpace.generateFrom
          ((fun V => {F : FellCloseds Y | ((F : Set Y) ∩ V).Nonempty})
            '' {V : Set Y | IsOpen V})]
          (U : Set (FellCloseds Y)) → f U = g U) ∧
      f ≠ g
\end{leansrc}
\par\medskip\noindent{\small \texttt{Pointfree.NablaData.Valuation.}\texttt{\bfseries ext\_of\_nabla}}\hfill{\footnotesize \path{lean/pointfree/Pointfree/Frame/HyperspacePresentation.lean}, line 259}\par\nopagebreak
\begin{leansrc}
theorem ext_of_nabla
    (hμS : μ.IsScottContinuous) (hνS : ν.IsScottContinuous)
    (hfin : ∀ p, μ (nabla geom p) ≠ ∞)
    (h : ∀ p, μ (nabla geom p) = ν (nabla geom p)) : μ = ν
\end{leansrc}
\par\medskip\noindent{\small \texttt{Pointfree.NablaData.Valuation.}\texttt{\bfseries exists\_nabla\_ne}}\hfill{\footnotesize \path{lean/pointfree/Pointfree/Frame/HyperspacePresentation.lean}, line 271}\par\nopagebreak
\begin{leansrc}
theorem exists_nabla_ne
    (hμS : μ.IsScottContinuous) (hνS : ν.IsScottContinuous)
    (hfin : ∀ p, μ (nabla geom p) ≠ ∞)
    (hne : μ ≠ ν) : ∃ p, μ (nabla geom p) ≠ ν (nabla geom p)
\end{leansrc}
\par\medskip\noindent{\small \texttt{Pointfree.Valuation.}\texttt{\bfseries ext\_of\_generators}}\hfill{\footnotesize \path{lean/pointfree/Pointfree/Valuation/Determination.lean}, line 131}\par\nopagebreak
\begin{leansrc}
theorem ext_of_generators
    (hμS : μ.IsScottContinuous) (hνS : ν.IsScottContinuous)
    (hGinf : ∀ i j, ∃ k, G i ⊓ G j = G k)
    (hGgen : ∀ x : L, ∃ S : Set κ, x = ⨆ i ∈ S, G i)
    (hfin : ∀ i, μ (G i) ≠ ∞)
    (h : ∀ i, μ (G i) = ν (G i)) : μ = ν
\end{leansrc}
\par\medskip\noindent{\small \texttt{Pointfree.Valuation.}\texttt{\bfseries ext\_of\_isTopologicalBasis}}\hfill{\footnotesize \path{lean/pointfree/Pointfree/Valuation/SpatialDetermination.lean}, line 60}\par\nopagebreak
\begin{leansrc}
theorem ext_of_isTopologicalBasis
    (hB : IsTopologicalBasis (Set.range fun i => (N i : Set X)))
    (hNinf : ∀ i j, ∃ k, N i ⊓ N j = N k)
    {μ ν : Valuation (Opens X)}
    (hμS : μ.IsScottContinuous) (hνS : ν.IsScottContinuous)
    (hfin : ∀ i, μ (N i) ≠ ∞)
    (h : ∀ i, μ (N i) = ν (N i)) : μ = ν
\end{leansrc}
\par\medskip\noindent{\small \texttt{Pointfree.SubProb.}\texttt{\bfseries promoteHomeomorph}}\hfill{\footnotesize \path{lean/pointfree/Pointfree/Valuation/Topology.lean}, line 230}\par\nopagebreak
\begin{leansrc}
noncomputable def promoteHomeomorph : SubProb L ≃ₜ Prob (L × Prop)
\end{leansrc}
\par\medskip\noindent{\small \texttt{Pointfree.SubProb.}\texttt{\bfseries induced\_promote\_eq}}\hfill{\footnotesize \path{lean/pointfree/Pointfree/Valuation/Topology.lean}, line 171}\par\nopagebreak
\begin{leansrc}
theorem induced_promote_eq :
    TopologicalSpace.induced (promote : SubProb L → Prob (L × Prop))
        (Prob.instTopologicalSpace (L × Prop))
      = SubProb.instTopologicalSpace L
\end{leansrc}
\par\medskip\noindent{\small \texttt{SuFunctor.Universe.General.}\texttt{\bfseries SubProbabilityMeasure}}\hfill{\footnotesize \path{lean/su-functor/SuFunctor/Universe/General/Carrier.lean}, line 51}\par\nopagebreak
\begin{leansrc}
def SubProbabilityMeasure : Type _
\end{leansrc}
\par\medskip\noindent{\small \texttt{Pointfree.SubProbabilityMeasure.}\texttt{\bfseries continuous\_toSubProb}}\hfill{\footnotesize \path{lean/pointfree/Pointfree/Valuation/BridgeContinuity.lean}, line 180}\par\nopagebreak
\begin{leansrc}
theorem SubProbabilityMeasure.continuous_toSubProb :
    Continuous (SubProbabilityMeasure.toSubProb (X := X))
\end{leansrc}
\par\medskip\noindent{\small \texttt{Pointfree.}\texttt{\bfseries daggerChartMeasure}}\hfill{\footnotesize \path{lean/pointfree/Pointfree/Measure/DaggerChart.lean}, line 390}\par\nopagebreak
\begin{leansrc}
noncomputable def daggerChartMeasure :
    SubProbabilityMeasure X ≃ₜ ProbabilityMeasure (X ⊕ Unit) where …
\end{leansrc}
\par\medskip\noindent{\small \texttt{Pointfree.}\texttt{\bfseries compactSpace\_of\_daggerChartMeasure}}\hfill{\footnotesize \path{lean/pointfree/Pointfree/Measure/DaggerChart.lean}, line 415}\par\nopagebreak
\begin{leansrc}
theorem compactSpace_of_daggerChartMeasure [T2Space X] [BorelSpace X]
    [CompactSpace X] : CompactSpace (SubProbabilityMeasure X)
\end{leansrc}
\par\medskip\noindent{\small \texttt{Pointfree.}\texttt{\bfseries metrizableSpace\_subProbabilityMeasure}}\hfill{\footnotesize \path{lean/pointfree/Pointfree/Measure/DaggerChart.lean}, line 425}\par\nopagebreak
\begin{leansrc}
theorem metrizableSpace_subProbabilityMeasure [MetrizableSpace X]
    [SeparableSpace X] [BorelSpace X] :
    MetrizableSpace (SubProbabilityMeasure X)
\end{leansrc}
\par\medskip\noindent{\small \texttt{SuFunctor.Universe.}\texttt{\bfseries belCarrier}}\hfill{\footnotesize \path{lean/su-functor/SuFunctor/Universe/Carrier.lean}, line 63}\par\nopagebreak
\begin{leansrc}
def belCarrier : Set (V → ℝ)
\end{leansrc}
\par\medskip\noindent{\small \texttt{SuFunctor.Universe.}\texttt{\bfseries BelState}}\hfill{\footnotesize \path{lean/su-functor/SuFunctor/Universe/Carrier.lean}, line 119}\par\nopagebreak
\begin{leansrc}
abbrev BelState : Type _
\end{leansrc}
\par\medskip\noindent{\small \texttt{SuFunctor.Universe.}\texttt{\bfseries convex\_belCarrier}}\hfill{\footnotesize \path{lean/su-functor/SuFunctor/Universe/Carrier.lean}, line 72}\par\nopagebreak
\begin{leansrc}
theorem convex_belCarrier : Convex ℝ (belCarrier V)
\end{leansrc}
\par\medskip\noindent{\small \texttt{SuFunctor.Universe.}\texttt{\bfseries isCompact\_belCarrier}}\hfill{\footnotesize \path{lean/su-functor/SuFunctor/Universe/Carrier.lean}, line 107}\par\nopagebreak
\begin{leansrc}
theorem isCompact_belCarrier : IsCompact (belCarrier V)
\end{leansrc}
\par\medskip\noindent{\small \texttt{SuFunctor.Universe.}\texttt{\bfseries probCarrier}}\hfill{\footnotesize \path{lean/su-functor/SuFunctor/Universe/ProbCarrier.lean}, line 45}\par\nopagebreak
\begin{leansrc}
def probCarrier : Set (V → ℝ)
\end{leansrc}
\par\medskip\noindent{\small \texttt{SuFunctor.Universe.}\texttt{\bfseries ProbState}}\hfill{\footnotesize \path{lean/su-functor/SuFunctor/Universe/ProbCarrier.lean}, line 85}\par\nopagebreak
\begin{leansrc}
abbrev ProbState : Type _
\end{leansrc}
\par\medskip\noindent{\small \texttt{SuFunctor.Universe.General.}\texttt{\bfseries sat\_sup\_iff}}\hfill{\footnotesize \path{lean/su-functor/SuFunctor/Universe/General/Faces.lean}, line 72}\par\nopagebreak
\begin{leansrc}
theorem sat_sup_iff {σ φ ψ : StatePenalty X} :
    Sat σ (φ ⊔ ψ) ↔ Sat σ φ ∧ Sat σ ψ
\end{leansrc}
\par\medskip\noindent{\small \texttt{SuFunctor.Universe.General.}\texttt{\bfseries entails\_inf\_iff}}\hfill{\footnotesize \path{lean/su-functor/SuFunctor/Universe/General/Faces.lean}, line 101}\par\nopagebreak
\begin{leansrc}
theorem entails_inf_iff {φ ψ ρ : StatePenalty X} :
    ρ ≤ φ ⊓ ψ ↔ ρ ≤ φ ∧ ρ ≤ ψ
\end{leansrc}
\par\medskip\noindent{\small \texttt{SuFunctor.Universe.General.}\texttt{\bfseries exists\_sat\_inf\_not\_sat}}\hfill{\footnotesize \path{lean/su-functor/SuFunctor/Universe/General/Faces.lean}, line 112}\par\nopagebreak
\begin{leansrc}
theorem exists_sat_inf_not_sat [Nonempty X] :
    ∃ σ φ ψ : StatePenalty X, Sat σ (φ ⊓ ψ) ∧ ¬ Sat σ φ ∧ ¬ Sat σ ψ
\end{leansrc}
\par\medskip\noindent{\small \texttt{SuFunctor.Universe.}\texttt{\bfseries iSup\_boxRows\_eq\_self}}\hfill{\footnotesize \path{lean/su-functor/SuFunctor/Universe/Generation.lean}, line 146}\par\nopagebreak
\begin{leansrc}
theorem iSup_boxRows_eq_self {f : ProbPenalty V}
    (hf : LowerSemicontinuous f) :
    ⨆ (q : ℚ) (l : V → ℚ) (u : V → ℚ) (_ : 0 < q)
        (_ : ∀ ν, boxRow q l u ν ≤ f ν), boxRow q l u = f
\end{leansrc}
\par\medskip\noindent{\small \texttt{SuFunctor.Universe.}\texttt{\bfseries iSup\_tents\_eq\_self}}\hfill{\footnotesize \path{lean/su-functor/SuFunctor/Universe/Tents.lean}, line 76}\par\nopagebreak
\begin{leansrc}
theorem iSup_tents_eq_self {σ : Penalty V} (hσ : LowerSemicontinuous σ) :
    ⨆ (μ₀ : BelState V) (L : ℝ) (M : ℝ)
        (_ : ∀ ν, tent μ₀ L M ν ≤ σ ν), tent μ₀ L M = σ
\end{leansrc}
\par\medskip\noindent{\small \texttt{SuFunctor.Universe.}\texttt{\bfseries ProbPenalty}}\hfill{\footnotesize \path{lean/su-functor/SuFunctor/Universe/ProbCarrier.lean}, line 157}\par\nopagebreak
\begin{leansrc}
abbrev ProbPenalty : Type _
\end{leansrc}
\par\medskip\noindent{\small \texttt{SuFunctor.Universe.General.}\texttt{\bfseries sat\_sup\_iff}}\hfill{\footnotesize \path{lean/su-functor/SuFunctor/Universe/General/Faces.lean}, line 72}\par\nopagebreak
\begin{leansrc}
theorem sat_sup_iff {σ φ ψ : StatePenalty X} :
    Sat σ (φ ⊔ ψ) ↔ Sat σ φ ∧ Sat σ ψ
\end{leansrc}
\par\medskip\noindent{\small \texttt{SuFunctor.Universe.General.}\texttt{\bfseries entails\_inf\_iff}}\hfill{\footnotesize \path{lean/su-functor/SuFunctor/Universe/General/Faces.lean}, line 101}\par\nopagebreak
\begin{leansrc}
theorem entails_inf_iff {φ ψ ρ : StatePenalty X} :
    ρ ≤ φ ⊓ ψ ↔ ρ ≤ φ ∧ ρ ≤ ψ
\end{leansrc}
\par\medskip\noindent{\small \texttt{SuFunctor.Universe.General.}\texttt{\bfseries exists\_sat\_inf\_not\_sat}}\hfill{\footnotesize \path{lean/su-functor/SuFunctor/Universe/General/Faces.lean}, line 112}\par\nopagebreak
\begin{leansrc}
theorem exists_sat_inf_not_sat [Nonempty X] :
    ∃ σ φ ψ : StatePenalty X, Sat σ (φ ⊓ ψ) ∧ ¬ Sat σ φ ∧ ¬ Sat σ ψ
\end{leansrc}
\par\medskip\noindent{\small \texttt{SuFunctor.Universe.General.}\texttt{\bfseries lowerSemicontinuous\_wall\_iff}}\hfill{\footnotesize \path{lean/su-functor/SuFunctor/Universe/General/Faces.lean}, line 222}\par\nopagebreak
\begin{leansrc}
theorem lowerSemicontinuous_wall_iff {C : Set (SubProbabilityMeasure X)} :
    LowerSemicontinuous (wall C) ↔ IsClosed C
\end{leansrc}
\par\medskip\noindent{\small \texttt{SuFunctor.Universe.General.}\texttt{\bfseries iSup\_gambleRows\_eq\_self}}\hfill{\footnotesize \path{lean/su-functor/SuFunctor/Universe/General/GenerationGeneral.lean}, line 248}\par\nopagebreak
\begin{leansrc}
theorem iSup_gambleRows_eq_self {ι : Y → FiniteMeasure X}
    (hι : Topology.IsInducing ι) {f : Y → ℝ≥0∞} (hf : LowerSemicontinuous f) :
    ⨆ (q : ℚ) (F : Finset ((X →ᵇ ℝ≥0) × ℚ × ℚ)) (_ : 0 < q)
        (_ : ∀ ν, gambleRow ι (ENNReal.ofReal q) F ν ≤ f ν),
      gambleRow ι (ENNReal.ofReal q) F = f
\end{leansrc}
\par\medskip\noindent{\small \texttt{SuFunctor.Universe.General.}\texttt{\bfseries isInducing\_testProfile}}\hfill{\footnotesize \path{lean/su-functor/SuFunctor/Universe/General/GenerationGeneral.lean}, line 176}\par\nopagebreak
\begin{leansrc}
theorem isInducing_testProfile {ι : Y → FiniteMeasure X}
    (hι : Topology.IsInducing ι) : Topology.IsInducing (testProfile ι)
\end{leansrc}
\par\medskip\noindent{\small \texttt{SuFunctor.Universe.General.}\texttt{\bfseries exists\_countable\_gambleRows\_eq\_self}}\hfill{\footnotesize \path{lean/su-functor/SuFunctor/Universe/General/GenerationGeneral.lean}, line 305}\par\nopagebreak
\begin{leansrc}
theorem exists_countable_gambleRows_eq_self [SecondCountableTopology Y]
    {ι : Y → FiniteMeasure X} (hι : Topology.IsInducing ι) {f : Y → ℝ≥0∞}
    (hf : LowerSemicontinuous f) :
    ∃ R : Set (ℚ × Finset ((X →ᵇ ℝ≥0) × ℚ × ℚ)), R.Countable
      ∧ (∀ p ∈ R, ∀ ν, gambleRow ι (ENNReal.ofReal p.1) p.2 ν ≤ f ν)
      ∧ ⨆ p ∈ R, gambleRow ι (ENNReal.ofReal p.1) p.2 = f
\end{leansrc}
\par\medskip\noindent{\small \texttt{SuFunctor.Universe.General.}\texttt{\bfseries secondCountableTopology\_probabilityMeasure}}\hfill{\footnotesize \path{lean/su-functor/SuFunctor/Universe/General/GenerationGeneral.lean}, line 433}\par\nopagebreak
\begin{leansrc}
theorem secondCountableTopology_probabilityMeasure
    [CompactSpace X] [TopologicalSpace.MetrizableSpace X] [BorelSpace X] :
    SecondCountableTopology (ProbabilityMeasure X)
\end{leansrc}
\par\medskip\noindent{\small \texttt{SuFunctor.Plain.Joint.}\texttt{\bfseries toMeas\_surjective}}\hfill{\footnotesize \path{lean/su-functor/SuFunctor/Plain/SeamDensity.lean}, line 333}\par\nopagebreak
\begin{leansrc}
theorem Joint.toMeas_surjective :
    Function.Surjective (Joint.toMeas : Joint X → MeasJoint X)
\end{leansrc}
\par\medskip\noindent{\small \texttt{Qtsl.Imprecise.}\texttt{\bfseries sum\_dominates\_sup}}\hfill{\footnotesize \path{lean/qtsl-aux/QtslAux/Imprecise/MarginalCompat.lean}, line 103}\par\nopagebreak
\begin{leansrc}
theorem sum_dominates_sup (m : ∀ i, FinPMF X → FinPMF (Xi i))
    (f : ∀ i, FinPMF (Xi i) → ENNRat) (ν : FinPMF X) :
    supPenalty m f ν ≤ compositePenalty m f ν
\end{leansrc}
\par\medskip\noindent{\small \texttt{Qtsl.Imprecise.}\texttt{\bfseries sublevel\_marginal\_of\_dominates}}\hfill{\footnotesize \path{lean/qtsl-aux/QtslAux/Imprecise/MarginalCompat.lean}, line 136}\par\nopagebreak
\begin{leansrc}
theorem sublevel_marginal_of_dominates (m : ∀ i, FinPMF X → FinPMF (Xi i))
    (f : ∀ i, FinPMF (Xi i) → ENNRat) {F : FinPMF X → ENNRat} {ν : FinPMF X}
    (hdom : supPenalty m f ν ≤ F ν) {t : ENNRat}
    (h : F ν ≤ t) (i : ι) : f i (m i ν) ≤ t
\end{leansrc}
\par\medskip\noindent{\small \texttt{SuFunctor.Universe.General.}\texttt{\bfseries compileEquiv}}\hfill{\footnotesize \path{lean/su-functor/SuFunctor/Universe/RadialUniverse.lean}, line 99}\par\nopagebreak
\begin{leansrc}
noncomputable def compileEquiv [Nonempty X] :
    {f : Penalty X // LowerSemicontinuous f} ≃ RadialUniverse X where …
\end{leansrc}
\par\medskip\noindent{\small \texttt{SuFunctor.Adequacy.}\texttt{\bfseries rwall\_ne\_top}}\hfill{\footnotesize \path{lean/su-functor/SuFunctor/Universe/RadialUniverse.lean}, line 349}\par\nopagebreak
\begin{leansrc}
theorem rwall_ne_top (c : SubPMF V → Prop) :
    rwall c ≠ (⊤ : Plain.Pen V)
\end{leansrc}
\par\medskip\noindent{\small \texttt{SuFunctor.Adequacy.}\texttt{\bfseries contradictory\_composite\_denotes\_bot}}\hfill{\footnotesize \path{lean/su-functor/SuFunctor/Universe/RadialUniverse.lean}, line 589}\par\nopagebreak
\begin{leansrc}
theorem contradictory_composite_denotes_bot :
    rwall xyEqC ⊔ rwall xFalseC ⊔ rwall yTrueC
      = rayPoint (SubPMF.zero (Bool × Bool))
\end{leansrc}
\par\medskip\noindent{\small \texttt{SuFunctor.Universe.General.}\texttt{\bfseries compile\_ne\_singleton}}\hfill{\footnotesize \path{lean/su-functor/SuFunctor/Universe/RadialUniverse.lean}, line 172}\par\nopagebreak
\begin{leansrc}
theorem compile_ne_singleton (f : Penalty X) {ν : SubProbabilityMeasure X}
    (hν : ν ≠ 0) : compile f ≠ ({ν} : Set (SubProbabilityMeasure X))
\end{leansrc}
\par\medskip\noindent{\small \texttt{SuFunctor.StateSpace.}\texttt{\bfseries sortedCone}}\hfill{\footnotesize \path{lean/su-functor/SuFunctor/StateSpace/FamCompletion.lean}, line 165}\par\nopagebreak
\begin{leansrc}
def sortedCone (n : ℕ) : OmegaSA
\end{leansrc}
\par\medskip\noindent{\small \texttt{SuFunctor.StateSpace.}\texttt{\bfseries mem\_sortedCone\_iff\_monotone}}\hfill{\footnotesize \path{lean/su-functor/SuFunctor/StateSpace/FamCompletion.lean}, line 132}\par\nopagebreak
\begin{leansrc}
theorem mem_sortedCone_iff_monotone {n : ℕ} (x : Fin n → ℝ) :
    x ∈ BasicSet (sortedConds n) ↔ Monotone x
\end{leansrc}
\par\medskip\noindent{\small \texttt{SuFunctor.BagFormer.}\texttt{\bfseries bagObj}}\hfill{\footnotesize \path{lean/su-functor/SuFunctor/BagFormer.lean}, line 99}\par\nopagebreak
\begin{leansrc}
def bagObj (A : WireSpan El) : WireSpan El
\end{leansrc}
\par\medskip\noindent{\small \texttt{SuFunctor.StateSpace.}\texttt{\bfseries bagIsColimit}}\hfill{\footnotesize \path{lean/su-functor/SuFunctor/StateSpace/FamCompletion.lean}, line 188}\par\nopagebreak
\begin{leansrc}
def bagIsColimit :
    IsColimit (famComponentsCofan bagObj)
\end{leansrc}
\par\medskip\noindent{\small \texttt{Pointfree.Prob.}\texttt{\bfseries lowerSemicontinuous\_eval}}\hfill{\footnotesize \path{lean/pointfree/Pointfree/Valuation/Topology.lean}, line 79}\par\nopagebreak
\begin{leansrc}
theorem lowerSemicontinuous_eval (a : M) :
    LowerSemicontinuous fun μ : Prob M => μ a
\end{leansrc}
\par\medskip\noindent{\small \texttt{Pointfree.SubProb.}\texttt{\bfseries upperSemicontinuous\_mass}}\hfill{\footnotesize \path{lean/pointfree/Pointfree/Valuation/Topology.lean}, line 122}\par\nopagebreak
\begin{leansrc}
theorem upperSemicontinuous_mass :
    UpperSemicontinuous fun μ : SubProb L => μ ⊤
\end{leansrc}
\par\medskip\noindent{\small \texttt{Pointfree.SubProb.}\texttt{\bfseries isClosed\_liveness\_wall}}\hfill{\footnotesize \path{lean/pointfree/Pointfree/Valuation/Topology.lean}, line 132}\par\nopagebreak
\begin{leansrc}
theorem isClosed_liveness_wall (p : ℝ≥0∞) :
    IsClosed {μ : SubProb L | p ≤ μ ⊤}
\end{leansrc}
\par\medskip\noindent{\small \texttt{SuFunctor.}\texttt{\bfseries satCredal\_lower\_rayC\_geC}}\hfill{\footnotesize \path{lean/su-functor/SuFunctor/Adequacy/Envelope.lean}, line 390}\par\nopagebreak
\begin{leansrc}
theorem satCredal_lower_rayC_geC {μ : SubPMF V} (hμ : μ ≠ SubPMF.zero V)
    (A : Set V) (q : ℚ) :
    SatCredal .lower (rayC μ) (geC A q) ↔ q ≤ μ.prob A
\end{leansrc}
\par\medskip\noindent{\small \texttt{SuFunctor.}\texttt{\bfseries maxSlice\_rayC\_empty\_iff}}\hfill{\footnotesize \path{lean/su-functor/SuFunctor/Adequacy/Envelope.lean}, line 216}\par\nopagebreak
\begin{leansrc}
theorem maxSlice_rayC_empty_iff {μ : SubPMF V} :
    (∀ ν, ¬ maxSlice (rayC μ) ν) ↔ μ = SubPMF.zero V
\end{leansrc}
\par\medskip\noindent{\small \texttt{SuFunctor.}\texttt{\bfseries not\_satCredal\_upper\_rayC\_geC}}\hfill{\footnotesize \path{lean/su-functor/SuFunctor/Adequacy/Envelope.lean}, line 381}\par\nopagebreak
\begin{leansrc}
theorem not_satCredal_upper_rayC_geC (μ : SubPMF V) {A : Set V} {q : ℚ}
    (hq : 0 < q) : ¬ SatCredal .upper (rayC μ) (geC A q)
\end{leansrc}
\par\medskip\noindent{\small \texttt{SuFunctor.Universe.SpecDnf.}\texttt{\bfseries top\_notMem\_modelClass}}\hfill{\footnotesize \path{lean/su-functor/SuFunctor/Universe/SpecDnf.lean}, line 216}\par\nopagebreak
\begin{leansrc}
theorem top_notMem_modelClass (D : SpecDnf K) :
    (⊤ : K → ℝ≥0∞) ∉ D.modelClass
\end{leansrc}
\par\medskip\noindent{\small \texttt{SuFunctor.Universe.General.}\texttt{\bfseries isClosed\_compile\_iff}}\hfill{\footnotesize \path{lean/su-functor/SuFunctor/Universe/General/Closedness.lean}, line 308}\par\nopagebreak
\begin{leansrc}
theorem isClosed_compile_iff [Nonempty X] {f : Penalty X} :
    IsClosed (compile f) ↔ LowerSemicontinuous f
\end{leansrc}
\par\medskip\noindent{\small \texttt{SuFunctor.Universe.General.}\texttt{\bfseries exists\_lsc\_compile\_eq\_iff}}\hfill{\footnotesize \path{lean/su-functor/SuFunctor/Universe/General/Closedness.lean}, line 294}\par\nopagebreak
\begin{leansrc}
theorem exists_lsc_compile_eq_iff [Nonempty X]
    {R : Set (SubProbabilityMeasure X)} :
    (∃ f : Penalty X, LowerSemicontinuous f ∧ compile f = R)
      ↔ R.Nonempty ∧ IsClosed R ∧ IsRadial R
\end{leansrc}
\par\medskip\noindent{\small \texttt{SuFunctor.Universe.General.}\texttt{\bfseries compile\_rayPenalty}}\hfill{\footnotesize \path{lean/su-functor/SuFunctor/Universe/General/Closedness.lean}, line 260}\par\nopagebreak
\begin{leansrc}
theorem compile_rayPenalty [Nonempty X] {R : Set (SubProbabilityMeasure X)}
    (hcl : IsClosed R) (hrad : IsRadial R) (hne : R.Nonempty) :
    compile (rayPenalty R) = R
\end{leansrc}
\par\medskip\noindent{\small \texttt{SuFunctor.Universe.General.}\texttt{\bfseries compile\_injective}}\hfill{\footnotesize \path{lean/su-functor/SuFunctor/Universe/General/Dictionary.lean}, line 238}\par\nopagebreak
\begin{leansrc}
theorem compile_injective [Nonempty X] :
    Function.Injective (compile (X := X))
\end{leansrc}
\par\medskip\noindent{\small \texttt{SuFunctor.Universe.General.}\texttt{\bfseries rayPenalty\_compile}}\hfill{\footnotesize \path{lean/su-functor/SuFunctor/Universe/General/Dictionary.lean}, line 231}\par\nopagebreak
\begin{leansrc}
theorem rayPenalty_compile [Nonempty X] (f : Penalty X) :
    rayPenalty (compile f) = f
\end{leansrc}
\par\medskip\noindent{\small \texttt{SuFunctor.Universe.General.}\texttt{\bfseries upperSemicontinuous\_rayGauge}}\hfill{\footnotesize \path{lean/su-functor/SuFunctor/Universe/General/Closedness.lean}, line 233}\par\nopagebreak
\begin{leansrc}
theorem upperSemicontinuous_rayGauge {R : Set (SubProbabilityMeasure X)}
    (hcl : IsClosed R) (hrad : IsRadial R) (hne : R.Nonempty) :
    UpperSemicontinuous (rayGauge R)
\end{leansrc}
\par\medskip\noindent{\small \texttt{SuFunctor.Universe.General.}\texttt{\bfseries le\_gradeN\_of\_forall\_lt}}\hfill{\footnotesize \path{lean/su-functor/SuFunctor/Universe/General/Carrier.lean}, line 326}\par\nopagebreak
\begin{leansrc}
theorem le_gradeN_of_forall_lt {a : ℝ≥0∞} {c : ℝ≥0} (hc : c ≤ 1)
    (h : ∀ y, y < a → c ≤ gradeN y) : c ≤ gradeN a
\end{leansrc}
\par\medskip\noindent{\small \texttt{SuFunctor.Universe.}\texttt{\bfseries compile}}\hfill{\footnotesize \path{lean/su-functor/SuFunctor/Universe/VietorisDictionary.lean}, line 168}\par\nopagebreak
\begin{leansrc}
def compile (f : ProbPenalty V) : Set (BelState V)
\end{leansrc}
\par\medskip\noindent{\small \texttt{SuFunctor.Universe.EpistemicStates.}\texttt{\bfseries condTarget\_mix}}\hfill{\footnotesize \path{lean/su-functor/SuFunctor/Universe/EpistemicStates.lean}, line 935}\par\nopagebreak
\begin{leansrc}
theorem condTarget_mix (k : B → ProbState A) {μ ν : ProbState (A × B)}
    {s t : ℝ} (hs : 0 ≤ s) (ht : 0 ≤ t) (hst : s + t = 1) :
    condTarget k (ProbState.mix s t hs ht hst μ ν)
      = ProbState.mix s t hs ht hst (condTarget k μ)
          (condTarget k ν)
\end{leansrc}
\par\medskip\noindent{\small \texttt{SuFunctor.Universe.EpistemicStates.}\texttt{\bfseries klq\_mix\_mix\_le}}\hfill{\footnotesize \path{lean/su-functor/SuFunctor/Universe/EpistemicStates.lean}, line 429}\par\nopagebreak
\begin{leansrc}
theorem klq_mix_mix_le {μ ν ρ σ : ProbState V} {s t : ℝ}
    (hs : 0 < s) (ht : 0 < t) (hst : s + t = 1) :
    klq (ProbState.mix s t hs.le ht.le hst μ ν)
        (ProbState.mix s t hs.le ht.le hst ρ σ)
      ≤ ENNReal.ofReal s * klq μ ρ + ENNReal.ofReal t * klq ν σ
\end{leansrc}
\par\medskip\noindent{\small \texttt{SuFunctor.Universe.EpistemicStates.}\texttt{\bfseries lowerSemicontinuous\_condPenalty}}\hfill{\footnotesize \path{lean/su-functor/SuFunctor/Universe/EpistemicStates.lean}, line 1056}\par\nopagebreak
\begin{leansrc}
theorem lowerSemicontinuous_condPenalty {β : ℝ≥0∞} (hβ : β ≠ ⊤)
    (k : B → ProbState A) :
    LowerSemicontinuous (condPenalty β k)
\end{leansrc}
\par\medskip\noindent{\small \texttt{SuFunctor.Universe.EpistemicStates.}\texttt{\bfseries klq\_condTarget\_eq\_zero\_iff}}\hfill{\footnotesize \path{lean/su-functor/SuFunctor/Universe/EpistemicStates.lean}, line 974}\par\nopagebreak
\begin{leansrc}
theorem klq_condTarget_eq_zero_iff (k : B → ProbState A)
    (μ : ProbState (A × B)) :
    klq μ (condTarget k μ) = 0
      ↔ ∀ b, 0 < (sndMarg μ : B → ℝ) b
          → ∀ a, (μ : A × B → ℝ) (a, b) / (sndMarg μ : B → ℝ) b
              = (k b : A → ℝ) a
\end{leansrc}
\par\medskip\noindent{\small \texttt{SuFunctor.Universe.EpistemicStates.}\texttt{\bfseries mutualInfoQ\_eq\_zero\_iff}}\hfill{\footnotesize \path{lean/su-functor/SuFunctor/Universe/EpistemicStates.lean}, line 844}\par\nopagebreak
\begin{leansrc}
theorem mutualInfoQ_eq_zero_iff (μ : ProbState (A × B)) :
    mutualInfoQ μ = 0 ↔ IsProductLaw μ
\end{leansrc}
\par\medskip\noindent{\small \texttt{SuFunctor.Universe.EpistemicStates.}\texttt{\bfseries continuous\_mutualInfoQ}}\hfill{\footnotesize \path{lean/su-functor/SuFunctor/Universe/EpistemicStates.lean}, line 882}\par\nopagebreak
\begin{leansrc}
theorem continuous_mutualInfoQ :
    Continuous (mutualInfoQ (A := A) (B := B))
\end{leansrc}
\par\medskip\noindent{\small \texttt{SuFunctor.Universe.EpistemicStates.Witness.}\texttt{\bfseries not\_penaltyConvex\_indepPenalty}}\hfill{\footnotesize \path{lean/su-functor/SuFunctor/Universe/EpistemicStates.lean}, line 1262}\par\nopagebreak
\begin{leansrc}
theorem not_penaltyConvex_indepPenalty {β : ℝ≥0∞} (hβ : β ≠ 0) :
    ¬ PenaltyConvex (indepPenalty (A := Fin 2) (B := Fin 2) β)
\end{leansrc}
\par\medskip\noindent{\small \texttt{SuFunctor.Universe.EpistemicStates.}\texttt{\bfseries indepBody\_credalBody\_not\_convex}}\hfill{\footnotesize \path{lean/su-functor/SuFunctor/Universe/EpistemicStates.lean}, line 1435}\par\nopagebreak
\begin{leansrc}
theorem indepBody_credalBody_not_convex :
    IsCredalBody (indepBody (Fin 2) (Fin 2))
      ∧ ¬ IsConvexCredalBody (indepBody (Fin 2) (Fin 2))
\end{leansrc}
\par\medskip\noindent{\small \texttt{SuFunctor.Universe.EpistemicStates.}\texttt{\bfseries compile\_marginalWallSum}}\hfill{\footnotesize \path{lean/su-functor/SuFunctor/Universe/EpistemicStates.lean}, line 1145}\par\nopagebreak
\begin{leansrc}
theorem compile_marginalWallSum {MX : Set (ProbState A)}
    {MY : Set (ProbState B)} (hX : MX.Nonempty) (hY : MY.Nonempty) :
    compile (fun μ : ProbState (A × B) =>
        probWall MX (fstMarg μ) + probWall MY (sndMarg μ))
      = {x : BelState (A × B) | ∃ μ ∈ frechetClass MX MY, ∃ γ : ℝ,
          0 ≤ γ ∧ γ ≤ 1 ∧ (x : A × B → ℝ) = γ • (μ : A × B → ℝ)}
\end{leansrc}
\par\medskip\noindent{\small \texttt{SuFunctor.Universe.EpistemicStates.}\texttt{\bfseries marginalSum\_eq\_zero\_iff}}\hfill{\footnotesize \path{lean/su-functor/SuFunctor/Universe/EpistemicStates.lean}, line 1078}\par\nopagebreak
\begin{leansrc}
theorem marginalSum_eq_zero_iff (g : ProbState A → ℝ≥0∞)
    (h : ProbState B → ℝ≥0∞) (μ : ProbState (A × B)) :
    g (fstMarg μ) + h (sndMarg μ) = 0
      ↔ g (fstMarg μ) = 0 ∧ h (sndMarg μ) = 0
\end{leansrc}
\par\medskip\noindent{\small \texttt{SuFunctor.Universe.EpistemicStates.Witness.}\texttt{\bfseries frechet\_does\_not\_imply\_independence}}\hfill{\footnotesize \path{lean/su-functor/SuFunctor/Universe/EpistemicStates.lean}, line 1321}\par\nopagebreak
\begin{leansrc}
theorem frechet_does_not_imply_independence :
    (probWall {uniform2} (fstMarg chi)
        + probWall {uniform2} (sndMarg chi) = 0)
      ∧ ¬ IsProductLaw chi ∧ mutualInfoQ chi ≠ 0
\end{leansrc}
\par\medskip\noindent{\small \texttt{SuFunctor.Universe.CredalPoset.}\texttt{\bfseries house\_includes\_all}}\hfill{\footnotesize \path{lean/su-functor/SuFunctor/Universe/CredalPoset.lean}, line 47}\par\nopagebreak
\begin{leansrc}
theorem house_includes_all {V : Type*} [Fintype V] [Nonempty V]
    (S : Set (BelState V)) :
    (IsConvexCredalBody S → IsCredalBody S) ∧
    (IsRichardsonBody S → IsCredalBody S) ∧
    (IsUltra S → IsCredalBody S) ∧
    (IsMVS S → IsCredalBody S) ∧
    (IsHL S → IsCredalBody S) ∧
    (IsCozman S → IsCredalBody S)
\end{leansrc}
\par\medskip\noindent{\small \texttt{SuFunctor.Universe.General.CredalPoset.}\texttt{\bfseries scaleSegment\_eq\_compile\_dirWall}}\hfill{\footnotesize \path{lean/su-functor/SuFunctor/Universe/General/CredalPoset.lean}, line 446}\par\nopagebreak
\begin{leansrc}
theorem scaleSegment_eq_compile_dirWall {K : Set (ProbabilityMeasure X)}
    (hne : K.Nonempty) : scaleSegment K = compile (dirWall K)
\end{leansrc}
\par\medskip\noindent{\small \texttt{SuFunctor.Universe.CredalPoset.}\texttt{\bfseries exists\_penaltyConvex\_not\_closed}}\hfill{\footnotesize \path{lean/su-functor/SuFunctor/Universe/CredalPoset/Adjudication.lean}, line 172}\par\nopagebreak
\begin{leansrc}
theorem exists_penaltyConvex_not_closed :
    ∃ f : ProbPenalty (Fin 2),
      PenaltyConvex f ∧ ¬ IsClosed (compile f)
\end{leansrc}
\par\medskip\noindent{\small \texttt{SuFunctor.Universe.CredalPoset.}\texttt{\bfseries exists\_penaltyConvex\_compile\_not\_richardson}}\hfill{\footnotesize \path{lean/su-functor/SuFunctor/Universe/CredalPoset/Adjudication.lean}, line 201}\par\nopagebreak
\begin{leansrc}
theorem exists_penaltyConvex_compile_not_richardson :
    ∃ f : ProbPenalty (Fin 2),
      PenaltyConvex f ∧ ¬ IsRichardsonBody (compile f)
\end{leansrc}
\par\medskip\noindent{\small \texttt{SuFunctor.Universe.CredalPoset.IsRichardsonBody.}\texttt{\bfseries isConvexCredalBody}}\hfill{\footnotesize \path{lean/su-functor/SuFunctor/Universe/CredalPoset/Inclusions.lean}, line 176}\par\nopagebreak
\begin{leansrc}
theorem IsRichardsonBody.isConvexCredalBody [Nonempty V] {S : Set (BelState V)}
    (h : IsRichardsonBody S) : IsConvexCredalBody S
\end{leansrc}
\par\medskip\noindent{\small \texttt{SuFunctor.Universe.General.CredalPoset.}\texttt{\bfseries GradeConvex}}\hfill{\footnotesize \path{lean/su-functor/SuFunctor/Universe/General/CredalPoset.lean}, line 165}\par\nopagebreak
\begin{leansrc}
def GradeConvex (f : Penalty X) : Prop
\end{leansrc}
\par\medskip\noindent{\small \texttt{SuFunctor.Universe.General.CredalPoset.}\texttt{\bfseries convexBody\_compile\_of\_gradeConvex}}\hfill{\footnotesize \path{lean/su-functor/SuFunctor/Universe/General/CredalPoset.lean}, line 252}\par\nopagebreak
\begin{leansrc}
theorem convexBody_compile_of_gradeConvex [Nonempty X] {f : Penalty X}
    (hgc : GradeConvex f) : ConvexBody (compile f)
\end{leansrc}
\par\medskip\noindent{\small \texttt{SuFunctor.Universe.CredalPoset.IsUltra.}\texttt{\bfseries isConvexCredalBody}}\hfill{\footnotesize \path{lean/su-functor/SuFunctor/Universe/CredalPoset/Defs.lean}, line 196}\par\nopagebreak
\begin{leansrc}
theorem IsUltra.isConvexCredalBody {S : Set (BelState V)}
    (h : IsUltra S) : IsConvexCredalBody S
\end{leansrc}
\par\medskip\noindent{\small \texttt{SuFunctor.Universe.CredalPoset.IsPrevisionBody.}\texttt{\bfseries isUltra}}\hfill{\footnotesize \path{lean/su-functor/SuFunctor/Universe/CredalPoset/Previsions.lean}, line 50}\par\nopagebreak
\begin{leansrc}
theorem IsPrevisionBody.isUltra {S : Set (BelState V)}
    (hP : IsPrevisionBody S) : IsUltra S
\end{leansrc}
\par\medskip\noindent{\small \texttt{SuFunctor.Universe.CredalPoset.}\texttt{\bfseries isPrevisionBody\_iff\_isUltra}}\hfill{\footnotesize \path{lean/su-functor/SuFunctor/Universe/CredalPoset/Previsions.lean}, line 202}\par\nopagebreak
\begin{leansrc}
theorem isPrevisionBody_iff_isUltra {S : Set (BelState V)} :
    IsPrevisionBody S ↔ IsUltra S
\end{leansrc}
\par\medskip\noindent{\small \texttt{SuFunctor.Universe.General.CredalPoset.}\texttt{\bfseries isClosed\_ptwiseDownHull}}\hfill{\footnotesize \path{lean/su-functor/SuFunctor/Universe/General/CredalPoset.lean}, line 588}\par\nopagebreak
\begin{leansrc}
theorem isClosed_ptwiseDownHull
    (horder : IsClosed {p : FiniteMeasure X × FiniteMeasure X |
      (p.1 : Measure X) ≤ (p.2 : Measure X)})
    (F : Finset (SubProbabilityMeasure X)) :
    IsClosed (ptwiseDownHull F)
\end{leansrc}
\par\medskip\noindent{\small \texttt{SuFunctor.Universe.CredalPoset.IsMVS.}\texttt{\bfseries isUltra}}\hfill{\footnotesize \path{lean/su-functor/SuFunctor/Universe/CredalPoset/Inclusions.lean}, line 270}\par\nopagebreak
\begin{leansrc}
theorem IsMVS.isUltra {S : Set (BelState V)} (h : IsMVS S) : IsUltra S
\end{leansrc}
\par\medskip\noindent{\small \texttt{SuFunctor.}\texttt{\bfseries maxSlice\_rayC}}\hfill{\footnotesize \path{lean/su-functor/SuFunctor/Adequacy/Envelope.lean}, line 176}\par\nopagebreak
\begin{leansrc}
theorem maxSlice_rayC {μ : SubPMF V} (hμ : μ ≠ SubPMF.zero V)
    (ν : SubPMF V) : maxSlice (rayC μ) ν ↔ singletonC μ ν
\end{leansrc}
\par\medskip\noindent{\small \texttt{SuFunctor.}\texttt{\bfseries satCredal\_lower\_rayC}}\hfill{\footnotesize \path{lean/su-functor/SuFunctor/Adequacy/Envelope.lean}, line 327}\par\nopagebreak
\begin{leansrc}
theorem satCredal_lower_rayC {μ : SubPMF V} (hμ : μ ≠ SubPMF.zero V)
    (c : SubPMF V → Prop) : SatCredal .lower (rayC μ) c ↔ c μ
\end{leansrc}
\par\medskip\noindent{\small \texttt{SuFunctor.}\texttt{\bfseries DownClosed}}\hfill{\footnotesize \path{lean/su-functor/SuFunctor/Adequacy/Envelope.lean}, line 305}\par\nopagebreak
\begin{leansrc}
def DownClosed (c : SubPMF V → Prop) : Prop
\end{leansrc}
\par\medskip\noindent{\small \texttt{SuFunctor.}\texttt{\bfseries rayC\_eq\_singletonC\_iff}}\hfill{\footnotesize \path{lean/su-functor/SuFunctor/Adequacy/Envelope.lean}, line 251}\par\nopagebreak
\begin{leansrc}
theorem rayC_eq_singletonC_iff {μ : SubPMF V} :
    (∀ ν, rayC μ ν ↔ singletonC μ ν) ↔ μ = SubPMF.zero V
\end{leansrc}
\par\medskip\noindent{\small \texttt{SuFunctor.}\texttt{\bfseries maxSlice\_unionRays\_subset}}\hfill{\footnotesize \path{lean/su-functor/SuFunctor/Adequacy/EnvelopeUnion.lean}, line 138}\par\nopagebreak
\begin{leansrc}
theorem maxSlice_unionRays_subset {μs : ι → SubPMF V} {ν : SubPMF V}
    (h : maxSlice (unionRays μs) ν) : unionSingletons μs ν
\end{leansrc}
\par\medskip\noindent{\small \texttt{SuFunctor.}\texttt{\bfseries maxSlice\_unionRays}}\hfill{\footnotesize \path{lean/su-functor/SuFunctor/Adequacy/EnvelopeUnion.lean}, line 183}\par\nopagebreak
\begin{leansrc}
theorem maxSlice_unionRays {μs : ι → SubPMF V} {m : ℚ}
    (hm : ∀ i, (μs i).mass = m) (hm0 : m ≠ 0) (ν : SubPMF V) :
    maxSlice (unionRays μs) ν ↔ unionSingletons μs ν
\end{leansrc}
\par\medskip\noindent{\small \texttt{Qtsl.Imprecise.}\texttt{\bfseries natExtG\_le\_lowerEnvG}}\hfill{\footnotesize \path{lean/qtsl-aux/QtslAux/Imprecise/NaturalExtensionGeneral.lean}, line 372}\par\nopagebreak
\begin{leansrc}
theorem natExtG_le_lowerEnvG {M' : Set (ProbabilityMeasure X)}
    (hM' : M'.Nonempty) (hdom : ∀ i, c i ≤ lowerEnvG M' (h i))
    (g : BddGamble X) : natExtG h c g ≤ lowerEnvG M' g
\end{leansrc}
\par\medskip\noindent{\small \texttt{Qtsl.Imprecise.}\texttt{\bfseries avoidsSureLossG\_of\_dominating}}\hfill{\footnotesize \path{lean/qtsl-aux/QtslAux/Imprecise/NaturalExtensionGeneral.lean}, line 381}\par\nopagebreak
\begin{leansrc}
theorem avoidsSureLossG_of_dominating {M' : Set (ProbabilityMeasure X)}
    (hM' : M'.Nonempty) (hdom : ∀ i, c i ≤ lowerEnvG M' (h i)) :
    AvoidsSureLossG h c
\end{leansrc}
\par\medskip\noindent{\small \texttt{Pointfree.}\texttt{\bfseries Prevision}}\hfill{\footnotesize \path{lean/pointfree/Pointfree/Prevision/Basic.lean}, line 286}\par\nopagebreak
\begin{leansrc}
structure Prevision (L : Type*) [Order.Frame L] where …
\end{leansrc}
\par\medskip\noindent{\small \texttt{SuFunctor.Universe.}\texttt{\bfseries singletons\_not\_definable}}\hfill{\footnotesize \path{lean/su-functor/SuFunctor/Universe/Points.lean}, line 165}\par\nopagebreak
\begin{leansrc}
theorem singletons_not_definable {σ : ProbPenalty V}
    (hσ : LowerSemicontinuous σ) (hproper : ∃ μ, σ μ ≠ ⊤)
    {n : ℕ} (u : Fin n → ProbPenalty V) (𝒱 : Fin n → Set (ProbPenalty V))
    (hfin : ∀ i, (𝒱 i).Finite) :
    ¬ ∀ τ : ProbPenalty V, LowerSemicontinuous τ →
        ((∃ i, CoupledSat τ (u i) (𝒱 i)) ↔ τ = σ)
\end{leansrc}
\par\medskip\noindent{\small \texttt{SuFunctor.Plain.}\texttt{\bfseries idealPoint\_le\_iff}}\hfill{\footnotesize \path{lean/su-functor/SuFunctor/Plain/Battery.lean}, line 117}\par\nopagebreak
\begin{leansrc}
theorem idealPoint_le_iff (σ : Pen (states A)) (𝒮 : SpecClass A) :
    ({σ} : SpecClass A) ⊆ 𝒮 ↔ σ ∈ 𝒮
\end{leansrc}
\par\medskip\noindent{\small \texttt{SuFunctor.Universe.}\texttt{\bfseries crescentClass}}\hfill{\footnotesize \path{lean/su-functor/SuFunctor/Universe/VietorisStratum.lean}, line 315}\par\nopagebreak
\begin{leansrc}
def crescentClass (u : ProbPenalty V) (𝒱 : Set (ProbPenalty V)) :
    Set (NonemptyCompacts (BelState V))
\end{leansrc}
\par\medskip\noindent{\small \texttt{SuFunctor.Universe.General.}\texttt{\bfseries isLocallyClosed\_crescentClass}}\hfill{\footnotesize \path{lean/su-functor/SuFunctor/Universe/General/CrescentStratum.lean}, line 246}\par\nopagebreak
\begin{leansrc}
theorem isLocallyClosed_crescentClass [Nonempty X] {u : Penalty X}
    {𝒱 : Set (Penalty X)} (hu : LowerSemicontinuous u)
    (h𝒱 : 𝒱.Finite) (hlsc : ∀ v ∈ 𝒱, LowerSemicontinuous v) :
    IsLocallyClosed (crescentClass u 𝒱)
\end{leansrc}
\par\medskip\noindent{\small \texttt{SuFunctor.Universe.SpecDnf.}\texttt{\bfseries compileNC\_mem\_crescentClass\_iff}}\hfill{\footnotesize \path{lean/su-functor/SuFunctor/Universe/SpecDnf.lean}, line 512}\par\nopagebreak
\begin{leansrc}
theorem SpecDnf.compileNC_mem_crescentClass_iff [Nonempty V]
    {σ : ProbPenalty V} (hσ : LowerSemicontinuous σ)
    (D : SpecDnf (ProbState V)) :
    compileNC σ hσ ∈ D.crescentClass ↔ σ ∈ D.modelClass
\end{leansrc}
\par\medskip\noindent{\small \texttt{SuFunctor.Universe.General.}\texttt{\bfseries compile\_subset\_compile\_iff}}\hfill{\footnotesize \path{lean/su-functor/SuFunctor/Universe/General/Dictionary.lean}, line 122}\par\nopagebreak
\begin{leansrc}
theorem compile_subset_compile_iff {f h : Penalty X} :
    compile f ⊆ compile h ↔ h ≤ f
\end{leansrc}
\par\medskip\noindent{\small \texttt{SuFunctor.Universe.}\texttt{\bfseries mem\_modelClass\_reachAvoid\_iff}}\hfill{\footnotesize \path{lean/su-functor/SuFunctor/Universe/ReachAvoid.lean}, line 94}\par\nopagebreak
\begin{leansrc}
theorem mem_modelClass_reachAvoid_iff {u ℓ σ : K → ℝ≥0∞} :
    σ ∈ (SpecAtom.reachAvoid u ℓ).modelClass ↔ u ≤ σ ∧ Dip σ ℓ
\end{leansrc}
\par\medskip\noindent{\small \texttt{SuFunctor.Universe.}\texttt{\bfseries coupledSat\_crispRA\_iff}}\hfill{\footnotesize \path{lean/su-functor/SuFunctor/Universe/ReachAvoid.lean}, line 494}\par\nopagebreak
\begin{leansrc}
theorem coupledSat_crispRA_iff {σ : K → ℝ≥0∞} {S T : Set K} :
    σ ∈ (SpecAtom.reachAvoid (wallOn S) (wallOn Tᶜ)).modelClass
      ↔ (∀ x, x ∉ S → σ x = ⊤) ∧ ∃ x ∈ T, σ x < ⊤
\end{leansrc}
\par\medskip\noindent{\small \texttt{SuFunctor.Universe.}\texttt{\bfseries exists\_genClosure\_reachAvoid\_not\_crispGenerated}}\hfill{\footnotesize \path{lean/su-functor/SuFunctor/Universe/ReachAvoid.lean}, line 563}\par\nopagebreak
\begin{leansrc}
theorem exists_genClosure_reachAvoid_not_crispGenerated :
    ∃ C : Set (Unit → ℝ≥0∞),
      GenClosure (reachAvoidClasses Unit) C
        ∧ ¬ GenClosure (crispReachAvoidClasses Unit) C
\end{leansrc}
\par\medskip\noindent{\small \texttt{SuFunctor.Universe.}\texttt{\bfseries genClosure\_reachAvoid\_iff\_exists\_specDnf}}\hfill{\footnotesize \path{lean/su-functor/SuFunctor/Universe/ReachAvoid.lean}, line 259}\par\nopagebreak
\begin{leansrc}
theorem genClosure_reachAvoid_iff_exists_specDnf {C : Set (K → ℝ≥0∞)} :
    GenClosure (reachAvoidClasses K) C
      ↔ ∃ D : SpecDnf K, C = D.modelClass
\end{leansrc}
\par\medskip\noindent{\small \texttt{SuFunctor.Universe.}\texttt{\bfseries genClosure\_reachAvoidCres\_iff\_exists\_specDnf}}\hfill{\footnotesize \path{lean/su-functor/SuFunctor/Universe/ReachAvoid.lean}, line 370}\par\nopagebreak
\begin{leansrc}
theorem genClosure_reachAvoidCres_iff_exists_specDnf
    {C : Set (NonemptyCompacts (BelState V))} :
    GenClosure (reachAvoidCrescents V) C
      ↔ ∃ D : SpecDnf (ProbState V), C = D.crescentClass
\end{leansrc}
\par\medskip\noindent{\small \texttt{SuFunctor.Universe.}\texttt{\bfseries isConstructibleUnion\_of\_genClosure\_reachAvoidCrescentsLsc}}\hfill{\footnotesize \path{lean/su-functor/SuFunctor/Universe/ReachAvoid.lean}, line 456}\par\nopagebreak
\begin{leansrc}
theorem isConstructibleUnion_of_genClosure_reachAvoidCrescentsLsc
    {C : Set (NonemptyCompacts (BelState V))}
    (h : GenClosure (reachAvoidCrescentsLsc V) C) :
    IsConstructibleUnion C
\end{leansrc}
\par\medskip\noindent{\small \texttt{SuFunctor.Universe.}\texttt{\bfseries modelClass\_eq\_iInter\_reachAvoid}}\hfill{\footnotesize \path{lean/su-functor/SuFunctor/Universe/ReachAvoid.lean}, line 122}\par\nopagebreak
\begin{leansrc}
theorem modelClass_eq_iInter_reachAvoid (a : SpecAtom K) :
    a.modelClass
      = ⋂ ℓ ∈ a.marks, (SpecAtom.reachAvoid a.floor ℓ).modelClass
\end{leansrc}
\par\medskip\noindent{\small \texttt{SuFunctor.Universe.}\texttt{\bfseries crescentClass\_eq\_iInter\_reachAvoid}}\hfill{\footnotesize \path{lean/su-functor/SuFunctor/Universe/ReachAvoid.lean}, line 296}\par\nopagebreak
\begin{leansrc}
theorem crescentClass_eq_iInter_reachAvoid (a : SpecAtom (ProbState V)) :
    crescentClass a.floor a.marks
      = ⋂ ℓ ∈ a.marks, crescentClass a.floor {ℓ, ⊤}
\end{leansrc}
\par\medskip\noindent{\small \texttt{SuFunctor.Universe.}\texttt{\bfseries exists\_modelClass\_ne\_iUnion\_reachAvoid}}\hfill{\footnotesize \path{lean/su-functor/SuFunctor/Universe/ReachAvoid.lean}, line 141}\par\nopagebreak
\begin{leansrc}
theorem exists_modelClass_ne_iUnion_reachAvoid :
    ∃ a : SpecAtom Unit,
      a.modelClass
        ≠ ⋃ ℓ ∈ a.marks, (SpecAtom.reachAvoid a.floor ℓ).modelClass
\end{leansrc}
\par\medskip\noindent{\small \texttt{SuFunctor.Universe.}\texttt{\bfseries dip\_sup\_iff}}\hfill{\footnotesize \path{lean/su-functor/SuFunctor/Universe/CoupledJudgment.lean}, line 226}\par\nopagebreak
\begin{leansrc}
theorem dip_sup_iff {σ v₁ v₂ : K → ℝ≥0∞} :
    Dip σ (v₁ ⊔ v₂) ↔ Dip σ v₁ ∨ Dip σ v₂
\end{leansrc}
\par\medskip\noindent{\small \texttt{SuFunctor.Universe.}\texttt{\bfseries IsConstructibleUnion}}\hfill{\footnotesize \path{lean/su-functor/SuFunctor/Universe/LocallyClosedLayer.lean}, line 406}\par\nopagebreak
\begin{leansrc}
def IsConstructibleUnion (s : Set X) : Prop
\end{leansrc}
\par\medskip\noindent{\small \texttt{SuFunctor.Universe.General.}\texttt{\bfseries crescentClass\_inter\_crescentClass}}\hfill{\footnotesize \path{lean/su-functor/SuFunctor/Universe/General/CrescentStratum.lean}, line 286}\par\nopagebreak
\begin{leansrc}
theorem crescentClass_inter_crescentClass {u₁ u₂ : Penalty X}
    {𝒱₁ 𝒱₂ : Set (Penalty X)} :
    crescentClass u₁ 𝒱₁ ∩ crescentClass u₂ 𝒱₂
      = crescentClass (u₁ ⊔ u₂) (𝒱₁ ∪ 𝒱₂)
\end{leansrc}
\par\medskip\noindent{\small \texttt{SuFunctor.Universe.}\texttt{\bfseries inter\_dnfClass\_append}}\hfill{\footnotesize \path{lean/su-functor/SuFunctor/Universe/DisjunctionElim.lean}, line 204}\par\nopagebreak
\begin{leansrc}
theorem inter_dnfClass_append {n m : ℕ}
    (γ : Set (NonemptyCompacts (BelState V)))
    (u : Fin n → ProbPenalty V) (𝒱 : Fin n → Set (ProbPenalty V))
    (u' : Fin m → ProbPenalty V) (𝒱' : Fin m → Set (ProbPenalty V)) :
    γ ∩ dnfClass (Fin.append u u') (Fin.append 𝒱 𝒱')
      = (γ ∩ dnfClass u 𝒱) ∪ (γ ∩ dnfClass u' 𝒱')
\end{leansrc}
\par\medskip\noindent{\small \texttt{SuFunctor.Universe.General.}\texttt{\bfseries compile\_pullPen}}\hfill{\footnotesize \path{lean/su-functor/SuFunctor/Universe/General/TransportGeneral.lean}, line 176}\par\nopagebreak
\begin{leansrc}
theorem compile_pullPen [Nonempty X] (hθ : Continuous θ) (f : Penalty Y) :
    compile (pullPen hθ f) = push θ ⁻¹' compile f
\end{leansrc}
\par\medskip\noindent{\small \texttt{SuFunctor.Universe.General.}\texttt{\bfseries transportClass\_union}}\hfill{\footnotesize \path{lean/su-functor/SuFunctor/Universe/General/TransportGeneral.lean}, line 278}\par\nopagebreak
\begin{leansrc}
theorem transportClass_union (θ : X → Y)
    (C₁ C₂ : Set (FellCloseds (SubProbabilityMeasure Y))) :
    transportClass θ (C₁ ∪ C₂)
      = transportClass θ C₁ ∪ transportClass θ C₂
\end{leansrc}
\par\medskip\noindent{\small \texttt{SuFunctor.Universe.}\texttt{\bfseries ctxOrE\_valid\_iff}}\hfill{\footnotesize \path{lean/su-functor/SuFunctor/Universe/DisjunctionElim.lean}, line 99}\par\nopagebreak
\begin{leansrc}
theorem ctxOrE_valid_iff :
    (∀ γ ψ : Set α, γ ∩ A ⊆ ψ → γ ∩ B ⊆ ψ → γ ∩ J ⊆ ψ) ↔ J ⊆ A ∪ B
\end{leansrc}
\par\medskip\noindent{\small \texttt{SuFunctor.Universe.}\texttt{\bfseries orI\_ctxOrE\_valid\_iff}}\hfill{\footnotesize \path{lean/su-functor/SuFunctor/Universe/DisjunctionElim.lean}, line 107}\par\nopagebreak
\begin{leansrc}
theorem orI_ctxOrE_valid_iff :
    (A ⊆ J ∧ B ⊆ J ∧
        ∀ γ ψ : Set α, γ ∩ A ⊆ ψ → γ ∩ B ⊆ ψ → γ ∩ J ⊆ ψ)
      ↔ J = A ∪ B
\end{leansrc}
\par\medskip\noindent{\small \texttt{SuFunctor.Universe.}\texttt{\bfseries exists\_ctxOrE\_countermodel\_of\_ssubset}}\hfill{\footnotesize \path{lean/su-functor/SuFunctor/Universe/DisjunctionElim.lean}, line 123}\par\nopagebreak
\begin{leansrc}
theorem exists_ctxOrE_countermodel_of_ssubset (h : A ∪ B ⊂ J) :
    ∃ γ ψ : Set α, γ ∩ A ⊆ ψ ∧ γ ∩ B ⊆ ψ ∧ ¬ γ ∩ J ⊆ ψ
\end{leansrc}
\par\medskip\noindent{\small \texttt{SuFunctor.Universe.}\texttt{\bfseries exists\_union\_crescentClass\_not\_isLocallyClosed}}\hfill{\footnotesize \path{lean/su-functor/SuFunctor/Universe/DisjunctionElim.lean}, line 403}\par\nopagebreak
\begin{leansrc}
theorem exists_union_crescentClass_not_isLocallyClosed :
    ∃ u v : ProbPenalty Bool,
      LowerSemicontinuous u ∧ LowerSemicontinuous v ∧
        ¬ IsLocallyClosed
          (crescentClass u ∅
            ∪ crescentClass (0 : ProbPenalty Bool) {v})
\end{leansrc}
\par\medskip\noindent{\small \texttt{SuFunctor.Universe.}\texttt{\bfseries dnfClass\_append}}\hfill{\footnotesize \path{lean/su-functor/SuFunctor/Universe/DisjunctionElim.lean}, line 146}\par\nopagebreak
\begin{leansrc}
theorem dnfClass_append {n m : ℕ} (u : Fin n → ProbPenalty V)
    (𝒱 : Fin n → Set (ProbPenalty V)) (u' : Fin m → ProbPenalty V)
    (𝒱' : Fin m → Set (ProbPenalty V)) :
    dnfClass (Fin.append u u') (Fin.append 𝒱 𝒱')
      = dnfClass u 𝒱 ∪ dnfClass u' 𝒱'
\end{leansrc}
\par\medskip\noindent{\small \texttt{SuFunctor.Universe.}\texttt{\bfseries isConstructibleUnion\_dnfClass\_append}}\hfill{\footnotesize \path{lean/su-functor/SuFunctor/Universe/DisjunctionElim.lean}, line 226}\par\nopagebreak
\begin{leansrc}
theorem isConstructibleUnion_dnfClass_append {n m : ℕ}
    {u : Fin n → ProbPenalty V} {𝒱 : Fin n → Set (ProbPenalty V)}
    {u' : Fin m → ProbPenalty V} {𝒱' : Fin m → Set (ProbPenalty V)}
    (hu : ∀ i, LowerSemicontinuous (u i)) (h𝒱 : ∀ i, (𝒱 i).Finite)
    (hlsc : ∀ i, ∀ v ∈ 𝒱 i, LowerSemicontinuous v)
    (hu' : ∀ i, LowerSemicontinuous (u' i)) (h𝒱' : ∀ i, (𝒱' i).Finite)
    (hlsc' : ∀ i, ∀ v ∈ 𝒱' i, LowerSemicontinuous v) :
    IsConstructibleUnion
      (dnfClass (Fin.append u u') (Fin.append 𝒱 𝒱'))
\end{leansrc}
\par\medskip\noindent{\small \texttt{SuFunctor.Universe.}\texttt{\bfseries ctxOrE\_dnfClass}}\hfill{\footnotesize \path{lean/su-functor/SuFunctor/Universe/DisjunctionElim.lean}, line 216}\par\nopagebreak
\begin{leansrc}
theorem ctxOrE_dnfClass {n m : ℕ} {u : Fin n → ProbPenalty V}
    {𝒱 : Fin n → Set (ProbPenalty V)} {u' : Fin m → ProbPenalty V}
    {𝒱' : Fin m → Set (ProbPenalty V)}
    {γ ψ : Set (NonemptyCompacts (BelState V))}
    (h₁ : γ ∩ dnfClass u 𝒱 ⊆ ψ) (h₂ : γ ∩ dnfClass u' 𝒱' ⊆ ψ) :
    γ ∩ dnfClass (Fin.append u u') (Fin.append 𝒱 𝒱') ⊆ ψ
\end{leansrc}
\par\medskip\noindent{\small \texttt{SuFunctor.Universe.}\texttt{\bfseries isClosed\_pairHull}}\hfill{\footnotesize \path{lean/su-functor/SuFunctor/Universe/HullClosure.lean}, line 270}\par\nopagebreak
\begin{leansrc}
theorem isClosed_pairHull {C D : Set (ProbabilityMeasure X)}
    (hC : IsClosed C) (hD : IsClosed D) : IsClosed (pairHull C D)
\end{leansrc}
\par\medskip\noindent{\small \texttt{SuFunctor.Universe.}\texttt{\bfseries isClosed\_binaryHull}}\hfill{\footnotesize \path{lean/su-functor/SuFunctor/Universe/HullClosure.lean}, line 395}\par\nopagebreak
\begin{leansrc}
theorem isClosed_binaryHull {C D : Set (ProbabilityMeasure X)}
    (hC : IsClosed C) (hD : IsClosed D) : IsClosed (binaryHull C D)
\end{leansrc}
\par\medskip\noindent{\small \texttt{SuFunctor.Universe.}\texttt{\bfseries sublevel\_qCoproduct\_eq}}\hfill{\footnotesize \path{lean/su-functor/SuFunctor/Universe/HullClosure.lean}, line 686}\par\nopagebreak
\begin{leansrc}
theorem sublevel_qCoproduct_eq {σ τ' : ProbabilityMeasure X → ℝ≥0∞}
    (hσ : LowerSemicontinuous σ) (hτ' : LowerSemicontinuous τ') (c : ℝ≥0∞) :
    {μ | qCoproduct σ τ' μ ≤ c}
      = binaryHull {ν | σ ν ≤ c} {ν | τ' ν ≤ c}
\end{leansrc}
\par\medskip\noindent{\small \texttt{SuFunctor.Universe.}\texttt{\bfseries qCoproduct\_le\_min}}\hfill{\footnotesize \path{lean/su-functor/SuFunctor/Universe/HullClosure.lean}, line 663}\par\nopagebreak
\begin{leansrc}
theorem qCoproduct_le_min (σ τ' : ProbabilityMeasure X → ℝ≥0∞)
    (μ : ProbabilityMeasure X) : qCoproduct σ τ' μ ≤ min (σ μ) (τ' μ)
\end{leansrc}
\par\medskip\noindent{\small \texttt{SuFunctor.Universe.}\texttt{\bfseries le\_qCoproduct\_of\_forall\_mixConvex}}\hfill{\footnotesize \path{lean/su-functor/SuFunctor/Universe/HullClosure.lean}, line 670}\par\nopagebreak
\begin{leansrc}
theorem le_qCoproduct_of_forall_mixConvex {σ τ' h : ProbabilityMeasure X → ℝ≥0∞}
    (hqc : ∀ c, MixConvex {ν | h ν ≤ c}) (h1 : ∀ ν, h ν ≤ σ ν)
    (h2 : ∀ ν, h ν ≤ τ' ν) (μ : ProbabilityMeasure X) :
    h μ ≤ qCoproduct σ τ' μ
\end{leansrc}
\par\medskip\noindent{\small \texttt{SuFunctor.Universe.}\texttt{\bfseries boxClass\_union\_ssubset}}\hfill{\footnotesize \path{lean/su-functor/SuFunctor/Universe/DisjunctionElim.lean}, line 275}\par\nopagebreak
\begin{leansrc}
theorem boxClass_union_ssubset {x y : α} (hxy : x ≠ y) :
    boxClass ({x} : Set α) ∪ boxClass {y} ⊂ boxClass ({x, y} : Set α)
\end{leansrc}
\par\medskip\noindent{\small \texttt{SuFunctor.Universe.}\texttt{\bfseries exists\_ctxOrE\_countermodel\_body\_reading\_carrier}}\hfill{\footnotesize \path{lean/su-functor/SuFunctor/Universe/DisjunctionElim.lean}, line 311}\par\nopagebreak
\begin{leansrc}
theorem exists_ctxOrE_countermodel_body_reading_carrier :
    ∃ γ ψ : Set (NonemptyCompacts (BelState Bool)),
      γ ∩ boxClass {BelState.dirac true} ⊆ ψ
        ∧ γ ∩ boxClass {BelState.dirac false} ⊆ ψ
        ∧ ¬ γ ∩ boxClass
            ({BelState.dirac true, BelState.dirac false}
              : Set (BelState Bool)) ⊆ ψ
\end{leansrc}
\par\medskip\noindent{\small \texttt{SuFunctor.Universe.}\texttt{\bfseries mem\_dnfClass\_append\_iff}}\hfill{\footnotesize \path{lean/su-functor/SuFunctor/Universe/DisjunctionElim.lean}, line 175}\par\nopagebreak
\begin{leansrc}
theorem mem_dnfClass_append_iff {n m : ℕ} {u : Fin n → ProbPenalty V}
    {𝒱 : Fin n → Set (ProbPenalty V)} {u' : Fin m → ProbPenalty V}
    {𝒱' : Fin m → Set (ProbPenalty V)} {K : NonemptyCompacts (BelState V)} :
    K ∈ dnfClass (Fin.append u u') (Fin.append 𝒱 𝒱')
      ↔ K ∈ dnfClass u 𝒱 ∨ K ∈ dnfClass u' 𝒱'
\end{leansrc}
\par\medskip\noindent{\small \texttt{SuFunctor.Universe.}\texttt{\bfseries coupled\_union\_refinement\_iff}}\hfill{\footnotesize \path{lean/su-functor/SuFunctor/Universe/LocallyClosedLayer.lean}, line 138}\par\nopagebreak
\begin{leansrc}
theorem coupled_union_refinement_iff {P : (K → ℝ≥0∞) → Prop} {u : K → ℝ≥0∞}
    {𝒱 : Set (K → ℝ≥0∞)} {u' : Fin m → K → ℝ≥0∞}
    {𝒱' : Fin m → Set (K → ℝ≥0∞)}
    (hPc : ∀ c, SelValid 𝒱' c → P (selSup u c)) :
    (∀ σ, P σ → CoupledSat σ u 𝒱 → ∃ j, CoupledSat σ (u' j) (𝒱' j))
      ↔ ∀ c, SelValid 𝒱' c → Blocked u 𝒱 u' c
\end{leansrc}
\par\medskip\noindent{\small \texttt{SuFunctor.Universe.}\texttt{\bfseries exists\_naive\_cover\_incomplete}}\hfill{\footnotesize \path{lean/su-functor/SuFunctor/Universe/LocallyClosedLayer.lean}, line 235}\par\nopagebreak
\begin{leansrc}
theorem exists_naive_cover_incomplete :
    ∃ u v : ProbPenalty (Fin 2),
      LowerSemicontinuous u ∧ LowerSemicontinuous v
        ∧ (∀ σ, CoupledSat σ u {u ⊔ v} ↔ CoupledSat σ u {v})
        ∧ ¬ u ⊔ v ≤ v
\end{leansrc}
\par\medskip\noindent{\small \texttt{SuFunctor.Universe.}\texttt{\bfseries exists\_essential\_floor\_raise}}\hfill{\footnotesize \path{lean/su-functor/SuFunctor/Universe/LocallyClosedLayer.lean}, line 294}\par\nopagebreak
\begin{leansrc}
theorem exists_essential_floor_raise :
    ∃ u₁ w : ProbPenalty (Fin 2),
      LowerSemicontinuous u₁ ∧ LowerSemicontinuous w
        ∧ (∀ σ : ProbPenalty (Fin 2), CoupledSat σ 0 {(⊤ : ProbPenalty (Fin 2))} →
            CoupledSat σ u₁ {(⊤ : ProbPenalty (Fin 2))} ∨ CoupledSat σ 0 {w})
        ∧ ¬ (∀ σ : ProbPenalty (Fin 2), LowerSemicontinuous σ →
              CoupledSat σ 0 {(⊤ : ProbPenalty (Fin 2))} →
              CoupledSat σ u₁ {(⊤ : ProbPenalty (Fin 2))})
        ∧ ¬ (∀ σ : ProbPenalty (Fin 2), LowerSemicontinuous σ →
              CoupledSat σ 0 {(⊤ : ProbPenalty (Fin 2))} →
              CoupledSat σ 0 {w})
        ∧ ¬ u₁ ≤ (0 : ProbPenalty (Fin 2))
\end{leansrc}
\par\medskip\noindent{\small \texttt{SuFunctor.Universe.}\texttt{\bfseries le\_infProj\_iff}}\hfill{\footnotesize \path{lean/su-functor/SuFunctor/Universe/Marginalization.lean}, line 79}\par\nopagebreak
\begin{leansrc}
theorem le_infProj_iff {p : Z → W} {σ : Z → ℝ≥0∞} {φ : W → ℝ≥0∞} :
    φ ≤ infProj p σ ↔ φ ∘ p ≤ σ
\end{leansrc}
\par\medskip\noindent{\small \texttt{SuFunctor.Universe.}\texttt{\bfseries le\_hide\_iff}}\hfill{\footnotesize \path{lean/su-functor/SuFunctor/Universe/Marginalization.lean}, line 195}\par\nopagebreak
\begin{leansrc}
theorem le_hide_iff {p : Z → W} {σ : Z → ℝ≥0∞} {φ : W → ℝ≥0∞}
    (hφ : LowerSemicontinuous φ) :
    φ ≤ hide p σ ↔ φ ∘ p ≤ σ
\end{leansrc}
\par\medskip\noindent{\small \texttt{SuFunctor.Universe.}\texttt{\bfseries hide\_sup\_le}}\hfill{\footnotesize \path{lean/su-functor/SuFunctor/Universe/HideFrobenius.lean}, line 156}\par\nopagebreak
\begin{leansrc}
theorem hide_sup_le (p : Z → W) (σ : Z → ℝ≥0∞) {g : W → ℝ≥0∞}
    (hg : LowerSemicontinuous g) :
    hide p σ ⊔ g ≤ hide p (σ ⊔ g ∘ p)
\end{leansrc}
\par\medskip\noindent{\small \texttt{SuFunctor.Universe.}\texttt{\bfseries hide\_add\_le}}\hfill{\footnotesize \path{lean/su-functor/SuFunctor/Universe/HideFrobenius.lean}, line 168}\par\nopagebreak
\begin{leansrc}
theorem hide_add_le (p : Z → W) (σ : Z → ℝ≥0∞) {g : W → ℝ≥0∞}
    (hg : LowerSemicontinuous g) :
    hide p σ + g ≤ hide p (σ + g ∘ p)
\end{leansrc}
\par\medskip\noindent{\small \texttt{SuFunctor.Universe.}\texttt{\bfseries exists\_lowerSemicontinuous\_hide\_add\_ne}}\hfill{\footnotesize \path{lean/su-functor/SuFunctor/Universe/HideFrobenius.lean}, line 309}\par\nopagebreak
\begin{leansrc}
theorem exists_lowerSemicontinuous_hide_add_ne :
    ∃ (σ : ProbabilityMeasure ℕ → ℝ≥0∞) (g : ℝ≥0∞ → ℝ≥0∞),
      LowerSemicontinuous σ ∧ LowerSemicontinuous g ∧
      hide liveMass (σ + g ∘ liveMass) ≠ hide liveMass σ + g
\end{leansrc}
\par\medskip\noindent{\small \texttt{SuFunctor.Universe.}\texttt{\bfseries lscEnv\_charFn}}\hfill{\footnotesize \path{lean/su-functor/SuFunctor/Universe/HideFrobenius.lean}, line 111}\par\nopagebreak
\begin{leansrc}
theorem lscEnv_charFn (C : Set Z) : lscEnv (charFn C) = charFn (closure C)
\end{leansrc}
\par\medskip\noindent{\small \texttt{SuFunctor.Universe.}\texttt{\bfseries exists\_judgment\_separating\_hide\_add}}\hfill{\footnotesize \path{lean/su-functor/SuFunctor/Universe/HideFrobenius.lean}, line 347}\par\nopagebreak
\begin{leansrc}
theorem exists_judgment_separating_hide_add :
    ∃ (σ : ProbabilityMeasure ℕ → ℝ≥0∞) (g φ : ℝ≥0∞ → ℝ≥0∞),
      LowerSemicontinuous σ ∧ LowerSemicontinuous g ∧
      LowerSemicontinuous φ ∧
      φ ≤ hide liveMass (σ + g ∘ liveMass) ∧
      ¬ φ ≤ hide liveMass σ + g
\end{leansrc}
\par\medskip\noindent{\small \texttt{SuFunctor.Universe.}\texttt{\bfseries hide\_sup\_comp}}\hfill{\footnotesize \path{lean/su-functor/SuFunctor/Universe/HideFrobenius.lean}, line 189}\par\nopagebreak
\begin{leansrc}
theorem hide_sup_comp [T2Space W] {p : Z → W} (hp : Continuous p)
    {σ : Z → ℝ≥0∞} (hσ : LowerSemicontinuous σ)
    (hcpt : ∀ c : ℝ≥0∞, c ≠ ⊤ → IsCompact {z | σ z ≤ c})
    {g : W → ℝ≥0∞} (hg : LowerSemicontinuous g) :
    hide p (σ ⊔ g ∘ p) = hide p σ ⊔ g
\end{leansrc}
\par\medskip\noindent{\small \texttt{SuFunctor.Universe.}\texttt{\bfseries hide\_add\_comp}}\hfill{\footnotesize \path{lean/su-functor/SuFunctor/Universe/HideFrobenius.lean}, line 209}\par\nopagebreak
\begin{leansrc}
theorem hide_add_comp [T2Space W] {p : Z → W} (hp : Continuous p)
    {σ : Z → ℝ≥0∞} (hσ : LowerSemicontinuous σ)
    (hcpt : ∀ c : ℝ≥0∞, c ≠ ⊤ → IsCompact {z | σ z ≤ c})
    {g : W → ℝ≥0∞} (hg : LowerSemicontinuous g) :
    hide p (σ + g ∘ p) = hide p σ + g
\end{leansrc}
\par\medskip\noindent{\small \texttt{Qtsl.Imprecise.}\texttt{\bfseries hurwicz\_ge\_iff\_exists\_chord}}\hfill{\footnotesize \path{lean/qtsl-aux/QtslAux/Imprecise/HurwiczCrescents.lean}, line 67}\par\nopagebreak
\begin{leansrc}
theorem hurwicz_ge_iff_exists_chord (α : ℚ) (hα0 : 0 ≤ α) (hα1 : α ≤ 1)
    (M : ι → FinPMF X) (f : X → ℚ) (lo : ℚ) :
    lo ≤ hurwicz α M f
      ↔ ∃ a b : ℚ, lo ≤ α * a + (1 - α) * b
          ∧ a ≤ upperE M f ∧ b ≤ lowerE M f
\end{leansrc}
\par\medskip\noindent{\small \texttt{Qtsl.Imprecise.}\texttt{\bfseries hurwiczSuperlevel\_eq\_iUnion\_crescents}}\hfill{\footnotesize \path{lean/qtsl-aux/QtslAux/Imprecise/HurwiczCrescents.lean}, line 81}\par\nopagebreak
\begin{leansrc}
theorem hurwiczSuperlevel_eq_iUnion_crescents (α : ℚ) (hα0 : 0 ≤ α)
    (hα1 : α ≤ 1) (M : ι → FinPMF X) (lo : ℚ) :
    {f : X → ℚ | lo ≤ hurwicz α M f}
      = ⋃ p : {q : ℚ × ℚ // lo ≤ α * q.1 + (1 - α) * q.2},
          ({f : X → ℚ | p.val.1 ≤ upperE M f}
            ∩ {f : X → ℚ | p.val.2 ≤ lowerE M f})
\end{leansrc}
\par\medskip\noindent{\small \texttt{Qtsl.Imprecise.}\texttt{\bfseries finite\_chords\_strict}}\hfill{\footnotesize \path{lean/qtsl-aux/QtslAux/Imprecise/HurwiczCrescents.lean}, line 112}\par\nopagebreak
\begin{leansrc}
theorem finite_chords_strict {α lo : ℚ} (hα0 : 0 < α) (hα1 : α < 1)
    (S : Finset (ℚ × ℚ)) (hS : ∀ p ∈ S, lo ≤ α * p.1 + (1 - α) * p.2) :
    ∃ (M : Bool → FinPMF Bool) (f : Bool → ℚ),
      lo ≤ hurwicz α M f ∧
        ∀ p ∈ S, ¬(p.2 ≤ lowerE M f ∧ p.1 ≤ upperE M f)
\end{leansrc}
\par\medskip\noindent{\small \texttt{Qtsl.Imprecise.}\texttt{\bfseries le\_hurwicz\_of\_chord}}\hfill{\footnotesize \path{lean/qtsl-aux/QtslAux/Imprecise/HurwiczCrescents.lean}, line 54}\par\nopagebreak
\begin{leansrc}
theorem le_hurwicz_of_chord {α lo a b : ℚ} (hα0 : 0 ≤ α) (hα1 : α ≤ 1)
    {M : ι → FinPMF X} {f : X → ℚ}
    (hchord : lo ≤ α * a + (1 - α) * b)
    (ha : a ≤ upperE M f) (hb : b ≤ lowerE M f) :
    lo ≤ hurwicz α M f
\end{leansrc}
\par\medskip\noindent{\small \texttt{Qtsl.Imprecise.}\texttt{\bfseries penUpperE}}\hfill{\footnotesize \path{lean/qtsl-aux/QtslAux/Imprecise/PenaltyEnvelope.lean}, line 100}\par\nopagebreak
\begin{leansrc}
def penUpperE (M : ι → FinPMF X) (c : ι → ℚ) (f : X → ℚ) : ℚ
\end{leansrc}
\par\medskip\noindent{\small \texttt{Qtsl.Imprecise.}\texttt{\bfseries penLowerE}}\hfill{\footnotesize \path{lean/qtsl-aux/QtslAux/Imprecise/PenaltyEnvelope.lean}, line 105}\par\nopagebreak
\begin{leansrc}
def penLowerE (M : ι → FinPMF X) (c : ι → ℚ) (f : X → ℚ) : ℚ
\end{leansrc}
\par\medskip\noindent{\small \texttt{Qtsl.Imprecise.}\texttt{\bfseries penUpperE\_zero\_cost}}\hfill{\footnotesize \path{lean/qtsl-aux/QtslAux/Imprecise/PenaltyEnvelope.lean}, line 138}\par\nopagebreak
\begin{leansrc}
theorem penUpperE_zero_cost (M : ι → FinPMF X) (f : X → ℚ) :
    penUpperE M 0 f = upperE M f
\end{leansrc}
\par\medskip\noindent{\small \texttt{Qtsl.Imprecise.}\texttt{\bfseries penUpperE\_eq\_sup\_levelProfile}}\hfill{\footnotesize \path{lean/qtsl-aux/QtslAux/Imprecise/PenaltyEnvelope.lean}, line 332}\par\nopagebreak
\begin{leansrc}
theorem penUpperE_eq_sup_levelProfile (M : ι → FinPMF X) (c : ι → ℚ)
    (f : X → ℚ) :
    penUpperE M c f
      = univ.sup' univ_nonempty (fun i₀ => levelProfile M c f i₀ - c i₀)
\end{leansrc}
\par\medskip\noindent{\small \texttt{Qtsl.Imprecise.}\texttt{\bfseries penUpperE\_compFam}}\hfill{\footnotesize \path{lean/qtsl-aux/QtslAux/Imprecise/PenaltyKernelCompose.lean}, line 71}\par\nopagebreak
\begin{leansrc}
theorem penUpperE_compFam (M : ι → FinPMF X) (c : ι → ℚ)
    (K : (x : X) → κ x → FinPMF Y) (d : (x : X) → κ x → ℚ)
    (f : X × Y → ℚ) :
    penUpperE (compFam M K) (penCompCost M c d) f
      = penUpperE M c fun x => penUpperE (K x) (d x) fun y => f (x, y)
\end{leansrc}
\par\medskip\noindent{\small \texttt{Qtsl.Imprecise.}\texttt{\bfseries penUpperE\_mix\_le}}\hfill{\footnotesize \path{lean/qtsl-aux/QtslAux/Imprecise/PenaltyEnvelope.lean}, line 219}\par\nopagebreak
\begin{leansrc}
theorem penUpperE_mix_le (M : ι → FinPMF X) (c : ι → ℚ) (f g : X → ℚ)
    {t : ℚ} (h₀ : 0 ≤ t) (h₁ : t ≤ 1) :
    penUpperE M c (fun x => t * f x + (1 - t) * g x)
      ≤ t * penUpperE M c f + (1 - t) * penUpperE M c g
\end{leansrc}
\par\medskip\noindent{\small \texttt{Qtsl.Imprecise.Mutants.}\texttt{\bfseries penUpperE\_not\_posHomogeneous}}\hfill{\footnotesize \path{lean/qtsl-aux/QtslAux/Imprecise/PenaltyEnvelope.lean}, line 362}\par\nopagebreak
\begin{leansrc}
theorem penUpperE_not_posHomogeneous :
    penUpperE (fun _ : Unit => FinPMF.dirac (default : Unit)) (fun _ => 1)
        (fun x => (2 : ℚ) * (fun _ : Unit => (0 : ℚ)) x)
      ≠ 2 * penUpperE (fun _ : Unit => FinPMF.dirac (default : Unit))
          (fun _ => 1) (fun _ : Unit => (0 : ℚ))
\end{leansrc}
\par\medskip\noindent{\small \texttt{Qtsl.Imprecise.}\texttt{\bfseries penUpperE\_dirac\_family}}\hfill{\footnotesize \path{lean/qtsl-aux/QtslAux/Imprecise/PenaltyEnvelope.lean}, line 292}\par\nopagebreak
\begin{leansrc}
theorem penUpperE_dirac_family (w : M → ℚ) (g : M → ℚ) :
    penUpperE (fun m => FinPMF.dirac m) w g
      = univ.sup' univ_nonempty (fun m => g m - w m)
\end{leansrc}
\par\medskip\noindent{\small \texttt{SuFunctor.Plain.}\texttt{\bfseries MeasJoint}}\hfill{\footnotesize \path{lean/su-functor/SuFunctor/Plain/MeasJoint.lean}, line 110}\par\nopagebreak
\begin{leansrc}
def MeasJoint (X : Type u) [MeasurableSpace X] : Type u
\end{leansrc}
\par\medskip\noindent{\small \texttt{SuFunctor.Plain.Joint.}\texttt{\bfseries toMeas\_push}}\hfill{\footnotesize \path{lean/su-functor/SuFunctor/Plain/GeneralSeam.lean}, line 131}\par\nopagebreak
\begin{leansrc}
theorem toMeas_push (h : X → Y) (hm : Measurable h) (μ : Joint X) :
    toMeas (μ.push h) = (toMeas μ).push h
\end{leansrc}
\par\medskip\noindent{\small \texttt{SuFunctor.Site.}\texttt{\bfseries penExtM}}\hfill{\footnotesize \path{lean/su-functor/SuFunctor/Site/GeneralWindow.lean}, line 86}\par\nopagebreak
\begin{leansrc}
abbrev penExtM {V W : Set T} (hVW : V ⊆ W) {X : Type u} [MeasurableSpace X] :
    MeasPen (Traj V X) → MeasPen (Traj W X)
\end{leansrc}
\par\medskip\noindent{\small \texttt{SuFunctor.Site.}\texttt{\bfseries penExtM\_comp\_toMeas}}\hfill{\footnotesize \path{lean/su-functor/SuFunctor/Site/GeneralWindow.lean}, line 265}\par\nopagebreak
\begin{leansrc}
theorem penExtM_comp_toMeas {V W : Set T} (hVW : V ⊆ W) {X : Type u}
    [MeasurableSpace X] (F : MeasPen (Traj V X)) :
    (penExtM hVW F) ∘ Joint.toMeas = penExt hVW (F ∘ Joint.toMeas)
\end{leansrc}
\par\medskip\noindent{\small \texttt{SuFunctor.Site.IsPersistentM.}\texttt{\bfseries seamRestrict}}\hfill{\footnotesize \path{lean/su-functor/SuFunctor/Site/GeneralWindow.lean}, line 273}\par\nopagebreak
\begin{leansrc}
theorem IsPersistentM.seamRestrict {U : ι → Set T} {X : Type u}
    [MeasurableSpace X] {Φ : ∀ i, Set (MeasPen (Traj (U i) X))}
    (hΦ : IsPersistentM U Φ) :
    IsPersistent U (fun i => (fun F => F ∘ Joint.toMeas) '' Φ i)
\end{leansrc}
\par\medskip\noindent{\small \texttt{SuFunctor.Plain.Joint.}\texttt{\bfseries toMeas\_injective}}\hfill{\footnotesize \path{lean/su-functor/SuFunctor/Plain/GeneralSeam.lean}, line 119}\par\nopagebreak
\begin{leansrc}
theorem toMeas_injective [MeasurableSingletonClass X] :
    Function.Injective (toMeas : Joint X → MeasJoint X)
\end{leansrc}
\par\medskip\noindent{\small \texttt{SuFunctor.Site.}\texttt{\bfseries penExt}}\hfill{\footnotesize \path{lean/su-functor/SuFunctor/Site/WindowPenalty.lean}, line 76}\par\nopagebreak
\begin{leansrc}
abbrev penExt {V W : Set T} (hVW : V ⊆ W) {X : Type u} :
    Pen (Traj V X) → Pen (Traj W X)
\end{leansrc}
\par\medskip\noindent{\small \texttt{SuFunctor.Plain.Joint.}\texttt{\bfseries dense\_ratSeam}}\hfill{\footnotesize \path{lean/su-functor/SuFunctor/Plain/SeamDensity.lean}, line 299}\par\nopagebreak
\begin{leansrc}
theorem Joint.dense_ratSeam
    (hd : DenseRange (Joint.toMeas : Joint X → MeasJoint X)) :
    Dense (Joint.toMeas '' {ν : Joint X | ν.IsRatWeighted})
\end{leansrc}
\par\medskip\noindent{\small \texttt{SuFunctor.Plain.Joint.}\texttt{\bfseries denseRange\_toMeas\_nat}}\hfill{\footnotesize \path{lean/su-functor/SuFunctor/Plain/SeamDensity.lean}, line 443}\par\nopagebreak
\begin{leansrc}
theorem Joint.denseRange_toMeas_nat :
    DenseRange (Joint.toMeas : Joint ℕ → MeasJoint ℕ)
\end{leansrc}
\par\medskip\noindent{\small \texttt{SuFunctor.Plain.MeasPen.}\texttt{\bfseries le\_of\_seam\_le}}\hfill{\footnotesize \path{lean/su-functor/SuFunctor/Plain/SeamDensity.lean}, line 149}\par\nopagebreak
\begin{leansrc}
theorem MeasPen.le_of_seam_le
    (hd : DenseRange (Joint.toMeas : Joint X → MeasJoint X))
    {F G : MeasPen X} (hF : LowerSemicontinuous F)
    (hG : UpperSemicontinuous G)
    (h : ∀ ν : Joint X, F ν.toMeas ≤ G ν.toMeas) : F ≤ G
\end{leansrc}
\par\medskip\noindent{\small \texttt{SuFunctor.Plain.MeasPen.}\texttt{\bfseries iSup\_seam\_eq}}\hfill{\footnotesize \path{lean/su-functor/SuFunctor/Plain/SeamDensity.lean}, line 176}\par\nopagebreak
\begin{leansrc}
theorem MeasPen.iSup_seam_eq
    (hd : DenseRange (Joint.toMeas : Joint X → MeasJoint X))
    {F : MeasPen X} (hF : LowerSemicontinuous F) :
    ⨆ ν : Joint X, F ν.toMeas = ⨆ μ : MeasJoint X, F μ
\end{leansrc}
\par\medskip\noindent{\small \texttt{SuFunctor.Plain.MeasPen.}\texttt{\bfseries eq\_of\_seam\_eq}}\hfill{\footnotesize \path{lean/su-functor/SuFunctor/Plain/SeamDensity.lean}, line 166}\par\nopagebreak
\begin{leansrc}
theorem MeasPen.eq_of_seam_eq
    (hd : DenseRange (Joint.toMeas : Joint X → MeasJoint X))
    {F G : MeasPen X} (hF : Continuous F) (hG : Continuous G)
    (h : ∀ ν : Joint X, F ν.toMeas = G ν.toMeas) : F = G
\end{leansrc}
\par\medskip\noindent{\small \texttt{SuFunctor.Plain.}\texttt{\bfseries dense\_seam\_inf\_gap}}\hfill{\footnotesize \path{lean/su-functor/SuFunctor/Plain/SeamDensity.lean}, line 635}\par\nopagebreak
\begin{leansrc}
theorem dense_seam_inf_gap :
    DenseRange (Joint.toMeas : Joint ℕ → MeasJoint ℕ) ∧
      ∃ F : MeasPen ℕ, LowerSemicontinuous F ∧
        (∀ ν : Joint ℕ, 1 ≤ F ν.toMeas) ∧
        (⨅ ν : Joint ℕ, F ν.toMeas) ≠ ⨅ μ : MeasJoint ℕ, F μ
\end{leansrc}
\par\medskip\noindent{\small \texttt{SuFunctor.Site.}\texttt{\bfseries IsMatching}}\hfill{\footnotesize \path{lean/su-functor/SuFunctor/Site/WindowPenalty.lean}, line 241}\par\nopagebreak
\begin{leansrc}
def IsMatching (U : ι → Set T) {X : Type u}
    (Φ : ∀ i, Set (Pen (Traj (U i) X))) : Prop
\end{leansrc}
\par\medskip\noindent{\small \texttt{SuFunctor.Site.}\texttt{\bfseries IsPersistent}}\hfill{\footnotesize \path{lean/su-functor/SuFunctor/Site/WindowPenalty.lean}, line 234}\par\nopagebreak
\begin{leansrc}
def IsPersistent (U : ι → Set T) {X : Type u}
    (Φ : ∀ i, Set (Pen (Traj (U i) X))) : Prop
\end{leansrc}
\par\medskip\noindent{\small \texttt{SuFunctor.Site.IsPersistent.}\texttt{\bfseries image2\_boxplusW}}\hfill{\footnotesize \path{lean/su-functor/SuFunctor/Site/WindowPenalty.lean}, line 269}\par\nopagebreak
\begin{leansrc}
theorem IsPersistent.image2_boxplusW {U : ι → Set T}
    {A B C : WireSpan El} {M : Pat A B} {N : Pat B C}
    {Φ : ∀ i, Set (Pen (Traj (U i) (Pat.Tot M)))}
    {Ψ : ∀ i, Set (Pen (Traj (U i) (Pat.Tot N)))}
    (hΦ : IsPersistent U Φ) (hΨ : IsPersistent U Ψ) :
    IsPersistent U (fun i => Set.image2 boxplusW (Φ i) (Ψ i))
\end{leansrc}
\par\medskip\noindent{\small \texttt{SuFunctor.Site.}\texttt{\bfseries classDescent}}\hfill{\footnotesize \path{lean/su-functor/SuFunctor/Site/ClassDescent.lean}, line 243}\par\nopagebreak
\begin{leansrc}
def classDescent (hUo : ∀ i, IsOpen (U i)) (hdir : Directed (· ⊆ ·) U)
    (hcov : ⋃ i, U i = univ) :
    Set (FellCloseds α) ≃o Set (ProjLim U) where …
\end{leansrc}
\par\medskip\noindent{\small \texttt{SuFunctor.Site.}\texttt{\bfseries isConstructibleUnion\_classDescent\_iff}}\hfill{\footnotesize \path{lean/su-functor/SuFunctor/Site/ClassDescent.lean}, line 268}\par\nopagebreak
\begin{leansrc}
theorem isConstructibleUnion_classDescent_iff (hUo : ∀ i, IsOpen (U i))
    (hdir : Directed (· ⊆ ·) U) (hcov : ⋃ i, U i = univ)
    {Φ : Set (FellCloseds α)} :
    IsConstructibleUnion (classDescent hUo hdir hcov Φ)
      ↔ IsConstructibleUnion Φ
\end{leansrc}
\par\medskip\noindent{\small \texttt{SuFunctor.Site.IsLocallyClosed.}\texttt{\bfseries classExt}}\hfill{\footnotesize \path{lean/su-functor/SuFunctor/Site/ClassDescent.lean}, line 167}\par\nopagebreak
\begin{leansrc}
theorem IsLocallyClosed.classExt (hV : IsOpen V) (hVW : V ⊆ W)
    {Φ : Set (FellCloseds ↥V)} (hΦ : IsLocallyClosed Φ) :
    IsLocallyClosed (classExt hVW Φ)
\end{leansrc}
\par\medskip\noindent{\small \texttt{SuFunctor.Site.IsConstructibleUnion.}\texttt{\bfseries classExt}}\hfill{\footnotesize \path{lean/su-functor/SuFunctor/Site/ClassDescent.lean}, line 174}\par\nopagebreak
\begin{leansrc}
theorem IsConstructibleUnion.classExt (hV : IsOpen V) (hVW : V ⊆ W)
    {Φ : Set (FellCloseds ↥V)} (hΦ : IsConstructibleUnion Φ) :
    IsConstructibleUnion (Site.classExt hVW Φ)
\end{leansrc}
\par\medskip\noindent{\small \texttt{SuFunctor.Site.}\texttt{\bfseries bijective\_toProjLim}}\hfill{\footnotesize \path{lean/su-functor/SuFunctor/Site/ClassDescent.lean}, line 281}\par\nopagebreak
\begin{leansrc}
theorem bijective_toProjLim (hUo : ∀ i, IsOpen (U i))
    (hdir : Directed (· ⊆ ·) U) (hcov : ⋃ i, U i = univ) :
    Function.Bijective (toProjLim U)
\end{leansrc}
\par\medskip\noindent{\small \texttt{SuFunctor.Site.}\texttt{\bfseries clip\_mem\_classExt\_iff}}\hfill{\footnotesize \path{lean/su-functor/SuFunctor/Site/ClassDescent.lean}, line 159}\par\nopagebreak
\begin{leansrc}
theorem clip_mem_classExt_iff (hVW : V ⊆ W) {Φ : Set (FellCloseds ↥V)}
    (F : FellCloseds α) :
    clip W F ∈ classExt hVW Φ ↔ clip V F ∈ Φ
\end{leansrc}
\par\medskip\noindent{\small \texttt{SuFunctor.}\texttt{\bfseries SingleSiteSemanticUniverse}}\hfill{\footnotesize \path{lean/su-functor/SuFunctor/SingleSiteSemanticUniverse.lean}, line 489}\par\nopagebreak
\begin{leansrc}
abbrev SingleSiteSemanticUniverse (El : Type u)
\end{leansrc}
\par\medskip\noindent{\small \texttt{SuFunctor.QtslSem.}\texttt{\bfseries SpecJudgLayer}}\hfill{\footnotesize \path{lean/su-functor/SuFunctor/QtslSem/SpecJudg.lean}, line 108}\par\nopagebreak
\begin{leansrc}
abbrev SpecJudgLayer (El : Type u)
\end{leansrc}
\par\medskip\noindent{\small \texttt{SuFunctor.QtslSem.}\texttt{\bfseries cellOfSubset}}\hfill{\footnotesize \path{lean/su-functor/SuFunctor/QtslSem/SpecJudg.lean}, line 128}\par\nopagebreak
\begin{leansrc}
def cellOfSubset {El : Type u} {a b : SpecJudgLayer El}
    {M : Pat a.un b.un} {𝒰 𝒱 : Cpl M} (h : 𝒰 ⊆ 𝒱) :
    Cell (C := SpecJudgLayer El)
      (⟨M, 𝒰⟩ : JudgLayer.JLoose (specJudgOrder El) a b) ⟨M, 𝒱⟩
      (𝟙 a) (𝟙 b)
\end{leansrc}
\par\medskip\noindent{\small \texttt{SuFunctor.Site.}\texttt{\bfseries wComp\_mono}}\hfill{\footnotesize \path{lean/su-functor/SuFunctor/Site/Battery.lean}, line 127}\par\nopagebreak
\begin{leansrc}
theorem wComp_mono {M : Pat A B} {N : Pat B C} {𝒰 𝒰' : WCpl U M}
    {𝒱 𝒱' : WCpl U N} (h₁ : ∀ i, 𝒰.1 i ⊆ 𝒰'.1 i)
    (h₂ : ∀ i, 𝒱.1 i ⊆ 𝒱'.1 i) :
    ∀ i, (wComp U 𝒰 𝒱).1 i ⊆ (wComp U 𝒰' 𝒱').1 i
\end{leansrc}
\par\medskip\noindent{\small \texttt{SuFunctor.Site.}\texttt{\bfseries wUnit\_transparent}}\hfill{\footnotesize \path{lean/su-functor/SuFunctor/Site/Battery.lean}, line 87}\par\nopagebreak
\begin{leansrc}
theorem wUnit_transparent {M : Pat A B} (𝒰 : WCpl U M) :
    wCellMap U (PseudoDoubleCategoryStruct.lLeftUnitor (C := WireSpan El) M)
      (wComp U (wUnit U A) 𝒰) = 𝒰
\end{leansrc}
\par\medskip\noindent{\small \texttt{CategoryTheory.}\texttt{\bfseries PseudoDoubleCategoryStruct}}\hfill{\footnotesize \path{lean/double-cat/DoubleCat/Basic.lean}, line 101}\par\nopagebreak
\begin{leansrc}
class PseudoDoubleCategoryStruct (C : Type u) [Category.{v} C] where …
\end{leansrc}
\par\medskip\noindent{\small \texttt{CategoryTheory.}\texttt{\bfseries PseudoDoubleCategory}}\hfill{\footnotesize \path{lean/double-cat/DoubleCat/Basic.lean}, line 322}\par\nopagebreak
\begin{leansrc}
class PseudoDoubleCategory (C : Type u) [Category.{v} C]
    [PseudoDoubleCategoryStruct C] : Prop where …
\end{leansrc}
\par\medskip\noindent{\small \texttt{CategoryTheory.PseudoDoubleCategoryStruct.Cell.}\texttt{\bfseries cast}}\hfill{\footnotesize \path{lean/double-cat/DoubleCat/Basic.lean}, line 178}\par\nopagebreak
\begin{leansrc}
def Cell.cast {a b c d : C} {m : Loose a b} {n : Loose c d} {f f' : a ⟶ c}
    {g g' : b ⟶ d} (hf : f = f') (hg : g = g') (α : Cell m n f g) :
    Cell m n f' g'
\end{leansrc}
\par\medskip\noindent{\small \texttt{CategoryTheory.}\texttt{\bfseries LaxDoubleFunctor}}\hfill{\footnotesize \path{lean/double-cat/DoubleCat/Functor/Lax.lean}, line 48}\par\nopagebreak
\begin{leansrc}
structure LaxDoubleFunctor extends C ⥤ D where …
\end{leansrc}
\par\medskip\noindent{\small \texttt{SuFunctor.Plain.MeasJoint.}\texttt{\bfseries push\_id}}\hfill{\footnotesize \path{lean/su-functor/SuFunctor/Plain/MeasJoint.lean}, line 156}\par\nopagebreak
\begin{leansrc}
@[simp] theorem push_id (μ : MeasJoint X) : μ.push id = μ
\end{leansrc}
\par\medskip\noindent{\small \texttt{SuFunctor.Plain.MeasJoint.}\texttt{\bfseries push\_push}}\hfill{\footnotesize \path{lean/su-functor/SuFunctor/Plain/MeasJoint.lean}, line 163}\par\nopagebreak
\begin{leansrc}
theorem push_push (h₁ : X → Y) (h₂ : Y → Z) (hm₁ : Measurable h₁)
    (hm₂ : Measurable h₂) (μ : MeasJoint X) :
    (μ.push h₁).push h₂ = μ.push (h₂ ∘ h₁)
\end{leansrc}
\par\medskip\noindent{\small \texttt{SuFunctor.Plain.MeasJoint.}\texttt{\bfseries mass\_push}}\hfill{\footnotesize \path{lean/su-functor/SuFunctor/Plain/MeasJoint.lean}, line 170}\par\nopagebreak
\begin{leansrc}
theorem mass_push (h : X → Y) (hm : Measurable h) (μ : MeasJoint X) :
    (μ.push h).1.mass = μ.1.mass
\end{leansrc}
\par\medskip\noindent{\small \texttt{SuFunctor.Plain.}\texttt{\bfseries continuous\_push}}\hfill{\footnotesize \path{lean/su-functor/SuFunctor/Plain/MeasJoint.lean}, line 336}\par\nopagebreak
\begin{leansrc}
theorem continuous_push {h : X → Y} (h_cont : Continuous h) :
    Continuous (fun μ : MeasJoint X => μ.push h)
\end{leansrc}
\par\medskip\noindent{\small \texttt{SuFunctor.Plain.MeasJoint.}\texttt{\bfseries push\_toMeasure}}\hfill{\footnotesize \path{lean/su-functor/SuFunctor/Plain/MeasJoint.lean}, line 153}\par\nopagebreak
\begin{leansrc}
@[simp] theorem push_toMeasure (h : X → Y) (μ : MeasJoint X) :
    ((μ.push h).1 : Measure Y) = (μ.1 : Measure X).map h
\end{leansrc}
\par\medskip\noindent{\small \texttt{SuFunctor.Universe.}\texttt{\bfseries exists\_lowerSemicontinuous\_not\_lowerSemicontinuous\_infProj}}\hfill{\footnotesize \path{lean/su-functor/SuFunctor/Universe/Marginalization.lean}, line 558}\par\nopagebreak
\begin{leansrc}
theorem exists_lowerSemicontinuous_not_lowerSemicontinuous_infProj :
    ∃ σ : ProbabilityMeasure ℕ → ℝ≥0∞, LowerSemicontinuous σ ∧
      ¬ LowerSemicontinuous (infProj liveMass σ)
\end{leansrc}
\par\medskip\noindent{\small \texttt{SuFunctor.Universe.}\texttt{\bfseries le\_lscEnv\_iff}}\hfill{\footnotesize \path{lean/su-functor/SuFunctor/Universe/Marginalization.lean}, line 173}\par\nopagebreak
\begin{leansrc}
theorem le_lscEnv_iff {g φ : Z → ℝ≥0∞} (hφ : LowerSemicontinuous φ) :
    φ ≤ lscEnv g ↔ φ ≤ g
\end{leansrc}
\par\medskip\noindent{\small \texttt{SuFunctor.Universe.}\texttt{\bfseries isCompact\_sublevel\_of\_isTightMeasureSet}}\hfill{\footnotesize \path{lean/su-functor/SuFunctor/Universe/Marginalization.lean}, line 317}\par\nopagebreak
\begin{leansrc}
theorem isCompact_sublevel_of_isTightMeasureSet
    {σ : ProbabilityMeasure 𝓧 → ℝ≥0∞} (hσ : LowerSemicontinuous σ)
    {c : ℝ≥0∞}
    (ht : IsTightMeasureSet
      {((μ : ProbabilityMeasure 𝓧) : Measure 𝓧) | μ ∈ {ν | σ ν ≤ c}}) :
    IsCompact {ν | σ ν ≤ c}
\end{leansrc}
\par\medskip\noindent{\small \texttt{SuFunctor.Universe.}\texttt{\bfseries sublevel\_infProj\_eq}}\hfill{\footnotesize \path{lean/su-functor/SuFunctor/Universe/Marginalization.lean}, line 213}\par\nopagebreak
\begin{leansrc}
theorem sublevel_infProj_eq [T1Space W] {p : Z → W} (hp : Continuous p)
    {σ : Z → ℝ≥0∞} (hσ : LowerSemicontinuous σ)
    (hcpt : ∀ c : ℝ≥0∞, c ≠ ⊤ → IsCompact {z | σ z ≤ c})
    {c : ℝ≥0∞} (hc : c ≠ ⊤) :
    {w | infProj p σ w ≤ c} = p '' {z | σ z ≤ c}
\end{leansrc}
\par\medskip\noindent{\small \texttt{SuFunctor.Universe.}\texttt{\bfseries exists\_infProj\_attained}}\hfill{\footnotesize \path{lean/su-functor/SuFunctor/Universe/Marginalization.lean}, line 270}\par\nopagebreak
\begin{leansrc}
theorem exists_infProj_attained [T1Space W] {p : Z → W} (hp : Continuous p)
    {σ : Z → ℝ≥0∞} (hσ : LowerSemicontinuous σ)
    (hcpt : ∀ c : ℝ≥0∞, c ≠ ⊤ → IsCompact {z | σ z ≤ c})
    {ν : W} (hν : infProj p σ ν ≠ ⊤) :
    ∃ z, p z = ν ∧ σ z = infProj p σ ν
\end{leansrc}
\par\medskip\noindent{\small \texttt{SuFunctor.Universe.}\texttt{\bfseries isCompact\_sublevel\_infProj}}\hfill{\footnotesize \path{lean/su-functor/SuFunctor/Universe/Marginalization.lean}, line 283}\par\nopagebreak
\begin{leansrc}
theorem isCompact_sublevel_infProj [T1Space W] {p : Z → W} (hp : Continuous p)
    {σ : Z → ℝ≥0∞} (hσ : LowerSemicontinuous σ)
    (hcpt : ∀ c : ℝ≥0∞, c ≠ ⊤ → IsCompact {z | σ z ≤ c})
    {c : ℝ≥0∞} (hc : c ≠ ⊤) :
    IsCompact {w | infProj p σ w ≤ c}
\end{leansrc}
\par\medskip\noindent{\small \texttt{SuFunctor.Universe.}\texttt{\bfseries lowerSemicontinuous\_infProj}}\hfill{\footnotesize \path{lean/su-functor/SuFunctor/Universe/Marginalization.lean}, line 292}\par\nopagebreak
\begin{leansrc}
theorem lowerSemicontinuous_infProj [T2Space W] {p : Z → W} (hp : Continuous p)
    {σ : Z → ℝ≥0∞} (hσ : LowerSemicontinuous σ)
    (hcpt : ∀ c : ℝ≥0∞, c ≠ ⊤ → IsCompact {z | σ z ≤ c}) :
    LowerSemicontinuous (infProj p σ)
\end{leansrc}
\par\medskip\noindent{\small \texttt{SuFunctor.Universe.}\texttt{\bfseries sublevel\_infProj\_margFst\_eq}}\hfill{\footnotesize \path{lean/su-functor/SuFunctor/Universe/Marginalization.lean}, line 347}\par\nopagebreak
\begin{leansrc}
theorem sublevel_infProj_margFst_eq
    {σ : ProbabilityMeasure (𝓧 × 𝓨) → ℝ≥0∞} (hσ : LowerSemicontinuous σ)
    (ht : ∀ c : ℝ≥0∞, c ≠ ⊤ → IsTightMeasureSet
      {((ρ : ProbabilityMeasure (𝓧 × 𝓨)) : Measure (𝓧 × 𝓨)) | ρ ∈ {ν | σ ν ≤ c}})
    {c : ℝ≥0∞} (hc : c ≠ ⊤) :
    {w | infProj margFst σ w ≤ c} = margFst '' {ρ | σ ρ ≤ c}
\end{leansrc}
\par\medskip\noindent{\small \texttt{SuFunctor.Plain.Pen.}\texttt{\bfseries pull}}\hfill{\footnotesize \path{lean/su-functor/SuFunctor/Plain/Penalty.lean}, line 100}\par\nopagebreak
\begin{leansrc}
def pull (h : X → Y) (f : Pen Y) : Pen X
\end{leansrc}
\par\medskip\noindent{\small \texttt{SuFunctor.Site.}\texttt{\bfseries siteSpecFunctor}}\hfill{\footnotesize \path{lean/su-functor/SuFunctor/Site/SiteFunctor.lean}, line 128}\par\nopagebreak
\begin{leansrc}
def siteSpecFunctor (U : ι → Set T) (El : Type u) :
    SpanValued.{u + 1, u, u, u + 1, u} (WireSpan El) where …
\end{leansrc}
\par\medskip\noindent{\small \texttt{SuFunctor.Plain.}\texttt{\bfseries measSpecObj}}\hfill{\footnotesize \path{lean/su-functor/SuFunctor/Plain/GeneralFunctor.lean}, line 56}\par\nopagebreak
\begin{leansrc}
def measSpecObj : MeasCat.{u}ᵒᵖ ⥤ Type u where …
\end{leansrc}
\par\medskip\noindent{\small \texttt{SuFunctor.Site.}\texttt{\bfseries SiteSpecTotal}}\hfill{\footnotesize \path{lean/su-functor/SuFunctor/Site/SiteFunctor.lean}, line 279}\par\nopagebreak
\begin{leansrc}
abbrev SiteSpecTotal (U : ι → Set T) (El : Type u)
\end{leansrc}
\par\medskip\noindent{\small \texttt{SuFunctor.Site.}\texttt{\bfseries siteSpecTotal\_pseudoDoubleCategory}}\hfill{\footnotesize \path{lean/su-functor/SuFunctor/Site/SiteFunctor.lean}, line 285}\par\nopagebreak
\begin{leansrc}
theorem siteSpecTotal_pseudoDoubleCategory (U : ι → Set T) (El : Type u) :
    ∃ _ : PseudoDoubleCategory (SiteSpecTotal U El), True
\end{leansrc}
\par\medskip\noindent{\small \texttt{SuFunctor.Site.}\texttt{\bfseries siteSpecProj\_isPseudo}}\hfill{\footnotesize \path{lean/su-functor/SuFunctor/Site/SiteFunctor.lean}, line 295}\par\nopagebreak
\begin{leansrc}
theorem siteSpecProj_isPseudo (U : ι → Set T) (El : Type u) :
    (siteSpecProj U El).IsPseudo
\end{leansrc}
\par\medskip\noindent{\small \texttt{SuFunctor.Site.}\texttt{\bfseries siteSpecProj\_discreteOpfibration}}\hfill{\footnotesize \path{lean/su-functor/SuFunctor/Site/SiteFunctor.lean}, line 304}\par\nopagebreak
\begin{leansrc}
theorem siteSpecProj_discreteOpfibration (U : ι → Set T) (El : Type u) :
    IsDiscreteDoubleOpfibration (SpanValued.proj (siteSpecFunctor U El))
\end{leansrc}
\par\medskip\noindent{\small \texttt{SuFunctor.}\texttt{\bfseries SingleSiteSemanticUniverse}}\hfill{\footnotesize \path{lean/su-functor/SuFunctor/SingleSiteSemanticUniverse.lean}, line 489}\par\nopagebreak
\begin{leansrc}
abbrev SingleSiteSemanticUniverse (El : Type u)
\end{leansrc}
\par\medskip\noindent{\small \texttt{SuFunctor.}\texttt{\bfseries dnfSpecFunctor}}\hfill{\footnotesize \path{lean/su-functor/SuFunctor/SingleSiteSemanticUniverse.lean}, line 365}\par\nopagebreak
\begin{leansrc}
def dnfSpecFunctor (El : Type u) :
    SpanValued.{u + 1, u, u, u + 1, u} (WireSpan El) where …
\end{leansrc}
\par\medskip\noindent{\small \texttt{SuFunctor.}\texttt{\bfseries singleSiteUniverseProj}}\hfill{\footnotesize \path{lean/su-functor/SuFunctor/SingleSiteSemanticUniverse.lean}, line 500}\par\nopagebreak
\begin{leansrc}
def singleSiteUniverseProj (El : Type u) :
    LaxDoubleFunctor (SingleSiteSemanticUniverse El) (WireSpan El)
\end{leansrc}
\par\medskip\noindent{\small \texttt{SuFunctor.}\texttt{\bfseries singleSiteUniverseProj\_isPseudo}}\hfill{\footnotesize \path{lean/su-functor/SuFunctor/SingleSiteSemanticUniverse.lean}, line 505}\par\nopagebreak
\begin{leansrc}
theorem singleSiteUniverseProj_isPseudo (El : Type u) :
    (singleSiteUniverseProj El).IsPseudo
\end{leansrc}
\par\medskip\noindent{\small \texttt{SuFunctor.}\texttt{\bfseries singleSiteUniverseProj\_discreteOpfibration}}\hfill{\footnotesize \path{lean/su-functor/SuFunctor/SingleSiteSemanticUniverse.lean}, line 514}\par\nopagebreak
\begin{leansrc}
theorem singleSiteUniverseProj_discreteOpfibration (El : Type u) :
    IsDiscreteDoubleOpfibration (singleSiteUniverseProj El)
\end{leansrc}
\par\medskip\noindent{\small \texttt{SuFunctor.}\texttt{\bfseries specMap\_modelClass\_eq}}\hfill{\footnotesize \path{lean/su-functor/SuFunctor/SingleSiteSemanticUniverse.lean}, line 533}\par\nopagebreak
\begin{leansrc}
theorem specMap_modelClass_eq (φ : A ⟶ B)
    (hφ : Function.Surjective (res φ)) (D : DnfSpec A) :
    specMap φ D.modelClass
      = (dnfPull (res φ) D).modelClass ∩ Set.range (Pen.pull (res φ))
\end{leansrc}
\par\medskip\noindent{\small \texttt{SuFunctor.}\texttt{\bfseries specMap\_modelClass\_subset}}\hfill{\footnotesize \path{lean/su-functor/SuFunctor/SingleSiteSemanticUniverse.lean}, line 540}\par\nopagebreak
\begin{leansrc}
theorem specMap_modelClass_subset (φ : A ⟶ B)
    (hφ : Function.Surjective (res φ)) (D : DnfSpec A) :
    specMap φ D.modelClass ⊆ (dnfPull (res φ) D).modelClass
\end{leansrc}
\par\medskip\noindent{\small \texttt{SuFunctor.Universe.SpecDnf.}\texttt{\bfseries isConstructibleUnion\_probRead}}\hfill{\footnotesize \path{lean/su-functor/SuFunctor/SingleSiteSemanticUniverse.lean}, line 594}\par\nopagebreak
\begin{leansrc}
theorem Universe.SpecDnf.isConstructibleUnion_probRead
    (D : SpecDnf (Joint V))
    (hfloor : ∀ a ∈ D.probRead.atoms, LowerSemicontinuous a.floor)
    (hmarks : ∀ a ∈ D.probRead.atoms, ∀ v ∈ a.marks, LowerSemicontinuous v) :
    IsConstructibleUnion D.probRead.crescentClass
\end{leansrc}
\par\medskip\noindent{\small \texttt{SuFunctor.Site.}\texttt{\bfseries singleton\_boxplusW}}\hfill{\footnotesize \path{lean/su-functor/SuFunctor/Site/Battery.lean}, line 103}\par\nopagebreak
\begin{leansrc}
theorem singleton_boxplusW {W : Set T} {M : Pat A B} {N : Pat B C}
    (u₁ : Pen (Traj W (Pat.Tot M))) (u₂ : Pen (Traj W (Pat.Tot N))) :
    Set.image2 boxplusW ({u₁} : Set (Pen (Traj W (Pat.Tot M))))
        ({u₂} : Set (Pen (Traj W (Pat.Tot N)))) =
      {boxplusW u₁ u₂}
\end{leansrc}
\par\medskip\noindent{\small \texttt{SuFunctor.Site.}\texttt{\bfseries transport\_singletonW}}\hfill{\footnotesize \path{lean/su-functor/SuFunctor/Site/Battery.lean}, line 113}\par\nopagebreak
\begin{leansrc}
theorem transport_singletonW {W : Set T} (φ : A ⟶ B)
    (σ : Pen (Traj W (states A))) :
    Pen.pull (trajMap (res φ)) '' ({σ} : Set (Pen (Traj W (states A)))) =
      {Pen.pull (trajMap (res φ)) σ}
\end{leansrc}
\par\medskip\noindent{\small \texttt{SuFunctor.Site.}\texttt{\bfseries wIdealPoint\_le\_iff}}\hfill{\footnotesize \path{lean/su-functor/SuFunctor/Site/Battery.lean}, line 121}\par\nopagebreak
\begin{leansrc}
theorem wIdealPoint_le_iff {W : Set T} (σ : Pen (Traj W (states A)))
    (S : Set (Pen (Traj W (states A)))) :
    ({σ} : Set (Pen (Traj W (states A)))) ⊆ S ↔ σ ∈ S
\end{leansrc}
\par\medskip\noindent{\small \texttt{SuFunctor.Universe.}\texttt{\bfseries exists\_coupledSat\_add\_fails}}\hfill{\footnotesize \path{lean/su-functor/SuFunctor/Universe/SpecDnf.lean}, line 420}\par\nopagebreak
\begin{leansrc}
theorem exists_coupledSat_add_fails :
    ∃ (σ₁ σ₂ u₁ u₂ : Unit → ℝ≥0∞) (𝒱₁ 𝒱₂ : Set (Unit → ℝ≥0∞)),
      ⊤ ∈ 𝒱₁ ∧ ⊤ ∈ 𝒱₂ ∧
      CoupledSat σ₁ u₁ 𝒱₁ ∧ CoupledSat σ₂ u₂ 𝒱₂ ∧
      ¬ CoupledSat (σ₁ + σ₂) (u₁ + u₂) (𝒱₁ ∪ 𝒱₂)
\end{leansrc}
\par\medskip\noindent{\small \texttt{SuFunctor.Universe.}\texttt{\bfseries add\_comp\_safety\_le}}\hfill{\footnotesize \path{lean/su-functor/SuFunctor/Universe/SpecDnf.lean}, line 405}\par\nopagebreak
\begin{leansrc}
theorem add_comp_safety_le (g₁ : K → K₁) (g₂ : K → K₂)
    {u₁ σ₁ : K₁ → ℝ≥0∞} {u₂ σ₂ : K₂ → ℝ≥0∞}
    (h₁ : u₁ ≤ σ₁) (h₂ : u₂ ≤ σ₂) :
    (u₁ ∘ g₁) + (u₂ ∘ g₂) ≤ (σ₁ ∘ g₁) + (σ₂ ∘ g₂)
\end{leansrc}
\par\medskip\noindent{\small \texttt{CategoryTheory.ImplDouble.}\texttt{\bfseries cell\_subsingleton}}\hfill{\footnotesize \path{lean/double-cat/DoubleCat/ImplDouble.lean}, line 68}\par\nopagebreak
\begin{leansrc}
theorem cell_subsingleton {a b c d : ImplDouble C} (m : Loose a b)
    (n : Loose c d) (f : a ⟶ c) (g : b ⟶ d) :
    Subsingleton (Cell m n f g)
\end{leansrc}
\par\medskip\noindent{\small \texttt{SuFunctor.Site.}\texttt{\bfseries mem\_wComp\_iff}}\hfill{\footnotesize \path{lean/su-functor/SuFunctor/Site/Battery.lean}, line 69}\par\nopagebreak
\begin{leansrc}
theorem mem_wComp_iff {M : Pat A B} {N : Pat B C} {𝒰 : WCpl U M}
    {𝒱 : WCpl U N} {i : ι}
    {w : Pen (Traj (U i) (Pat.Tot (Pat.comp M N)))} :
    w ∈ (wComp U 𝒰 𝒱).1 i ↔
      ∃ u₁ ∈ 𝒰.1 i, ∃ u₂ ∈ 𝒱.1 i, boxplusW u₁ u₂ = w
\end{leansrc}
\par\medskip\noindent{\small \texttt{SuFunctor.Site.}\texttt{\bfseries wComp\_nonempty}}\hfill{\footnotesize \path{lean/su-functor/SuFunctor/Site/Battery.lean}, line 78}\par\nopagebreak
\begin{leansrc}
theorem wComp_nonempty {M : Pat A B} {N : Pat B C} {𝒰 : WCpl U M}
    {𝒱 : WCpl U N} (h₁ : ∀ i, (𝒰.1 i).Nonempty)
    (h₂ : ∀ i, (𝒱.1 i).Nonempty) :
    ∀ i, ((wComp U 𝒰 𝒱).1 i).Nonempty
\end{leansrc}
\par\medskip\noindent{\small \texttt{SuFunctor.Site.}\texttt{\bfseries taxedBoxplusW\_breaks\_unit\_transparency}}\hfill{\footnotesize \path{lean/su-functor/SuFunctor/Site/Battery.lean}, line 152}\par\nopagebreak
\begin{leansrc}
theorem taxedBoxplusW_breaks_unit_transparency {W : Set T} (M : Pat A B)
    {c : ℝ≥0∞} (hc : c ≠ 0) :
    ¬ ∀ u : Pen (Traj W (Pat.Tot M)),
      Pen.pull (trajMap (PatCell.tot
          (PseudoDoubleCategoryStruct.lLeftUnitor (C := WireSpan El) M)))
          (taxedBoxplusW c (unitPenW A W) u) = u
\end{leansrc}
\par\medskip\noindent{\small \texttt{SuFunctor.Site.}\texttt{\bfseries funnelCompW\_breaks\_unit\_coherence}}\hfill{\footnotesize \path{lean/su-functor/SuFunctor/Site/Battery.lean}, line 176}\par\nopagebreak
\begin{leansrc}
theorem funnelCompW_breaks_unit_coherence {W : Set T} (M : Pat A B) :
    ¬ ∀ 𝒰 : Set (Pen (Traj W (Pat.Tot M))),
      Pen.pull (trajMap (PatCell.tot
          (PseudoDoubleCategoryStruct.lLeftUnitor (C := WireSpan El) M))) ''
          funnelCompW ({unitPenW A W} :
            Set (Pen (Traj W (Pat.Tot (Pat.id A))))) 𝒰 = 𝒰
\end{leansrc}
\par\medskip\noindent{\small \texttt{Qtsl.Imprecise.}\texttt{\bfseries sublevel\_compositeG\_subset\_glueSetG}}\hfill{\footnotesize \path{lean/qtsl-aux/QtslAux/Imprecise/ContainmentGeneral.lean}, line 111}\par\nopagebreak
\begin{leansrc}
theorem sublevel_compositeG_subset_glueSetG (m : ∀ i, J → Mi i)
    (f : ∀ i, Mi i → ENNRat) (F : ENNRat) :
    {ν | compositePenaltyG m f ν ≤ F} ⊆ glueSetG m f F
\end{leansrc}
\par\medskip\noindent{\small \texttt{Qtsl.Imprecise.}\texttt{\bfseries certificate\_transfer}}\hfill{\footnotesize \path{lean/qtsl-aux/QtslAux/Imprecise/ContainmentGeneral.lean}, line 121}\par\nopagebreak
\begin{leansrc}
theorem certificate_transfer (m : ∀ i, J → Mi i) (f : ∀ i, Mi i → ENNRat)
    (F : ENNRat) {P : J → Prop} (h : ∀ ν ∈ glueSetG m f F, P ν) :
    ∀ ν, compositePenaltyG m f ν ≤ F → P ν
\end{leansrc}
\par\medskip\noindent{\small \texttt{Qtsl.Imprecise.}\texttt{\bfseries sublevel\_pmComposite\_subset\_glue}}\hfill{\footnotesize \path{lean/qtsl-aux/QtslAux/Imprecise/ContainmentGeneral.lean}, line 307}\par\nopagebreak
\begin{leansrc}
theorem sublevel_pmComposite_subset_glue {proj : ∀ i, X → Xi i}
    (hproj : ∀ i, Continuous (proj i))
    (f : ∀ i, ProbabilityMeasure (Xi i) → ENNRat) (F : ENNRat) :
    {ν | pmCompositePenalty proj (fun i => (hproj i).measurable) f ν ≤ F}
      ⊆ glueSetG (pmMarginal proj fun i => (hproj i).measurable) f F
\end{leansrc}
\par\medskip\noindent{\small \texttt{Qtsl.Imprecise.}\texttt{\bfseries isClosed\_pmGlueSet}}\hfill{\footnotesize \path{lean/qtsl-aux/QtslAux/Imprecise/ContainmentGeneral.lean}, line 300}\par\nopagebreak
\begin{leansrc}
theorem isClosed_pmGlueSet {proj : ∀ i, X → Xi i}
    (hproj : ∀ i, Continuous (proj i))
    {f : ∀ i, ProbabilityMeasure (Xi i) → ENNRat}
    (hf : ∀ i, LowerSemicontinuous (f i)) (F : ENNRat) :
    IsClosed (glueSetG (pmMarginal proj fun i => (hproj i).measurable) f F)
\end{leansrc}
\par\medskip\noindent{\small \texttt{Qtsl.Imprecise.}\texttt{\bfseries forall\_sublevel\_eq\_iff}}\hfill{\footnotesize \path{lean/qtsl-aux/QtslAux/Imprecise/ContainmentGeneral.lean}, line 162}\par\nopagebreak
\begin{leansrc}
theorem forall_sublevel_eq_iff {g h : J → ENNRat} :
    (∀ F, {ν | g ν ≤ F} = {ν | h ν ≤ F}) ↔ g = h
\end{leansrc}
\par\medskip\noindent{\small \texttt{Qtsl.Imprecise.ENNRat.}\texttt{\bfseries add\_eq\_max\_iff}}\hfill{\footnotesize \path{lean/qtsl-aux/QtslAux/Imprecise/ContainmentGeneral.lean}, line 81}\par\nopagebreak
\begin{leansrc}
theorem ENNRat.add_eq_max_iff (a b : ENNRat) :
    a + b = max a b ↔ min a b = 0 ∨ max a b = ⊤
\end{leansrc}
\par\medskip\noindent{\small \texttt{Qtsl.Imprecise.}\texttt{\bfseries binaryComposite\_ne\_glue}}\hfill{\footnotesize \path{lean/qtsl-aux/QtslAux/Imprecise/ContainmentGeneral.lean}, line 191}\par\nopagebreak
\begin{leansrc}
theorem binaryComposite_ne_glue :
    (fun ν : Unit => (fun _ : Unit => (1 : ENNRat)) ν
        + (fun _ : Unit => (1 : ENNRat)) ν)
      ≠ fun ν : Unit => max ((fun _ : Unit => (1 : ENNRat)) ν)
          ((fun _ : Unit => (1 : ENNRat)) ν)
\end{leansrc}
\par\medskip\noindent{\small \texttt{Qtsl.Imprecise.}\texttt{\bfseries sublevel\_ssubset\_glue\_witness}}\hfill{\footnotesize \path{lean/qtsl-aux/QtslAux/Imprecise/ContainmentGeneral.lean}, line 206}\par\nopagebreak
\begin{leansrc}
theorem sublevel_ssubset_glue_witness :
    {ν : Unit | compositePenaltyG (fun _ : Bool => id)
        (fun _ : Bool => fun _ : Unit => (1 : ENNRat)) ν ≤ 1}
      ⊂ glueSetG (fun _ : Bool => id)
        (fun _ : Bool => fun _ : Unit => (1 : ENNRat)) 1
\end{leansrc}
\par\medskip\noindent{\small \texttt{Qtsl.Imprecise.}\texttt{\bfseries signed\_relaxation\_fails}}\hfill{\footnotesize \path{lean/qtsl-aux/QtslAux/Imprecise/MarginalCompat.lean}, line 179}\par\nopagebreak
\begin{leansrc}
theorem signed_relaxation_fails :
    ∃ f : Fin 2 → FinPMF (Fin 2) → ℚ,
      ¬ ∀ (i : Fin 2) (ν : FinPMF (Fin 2)),
          f i ν ≤ compositePenaltyQ (Xi := fun _ => Fin 2)
            (fun _ ν => ν) f ν
\end{leansrc}
\par\medskip\noindent{\small \texttt{Qtsl.Imprecise.}\texttt{\bfseries subunit\_weight\_fails}}\hfill{\footnotesize \path{lean/qtsl-aux/QtslAux/Imprecise/MarginalCompat.lean}, line 189}\par\nopagebreak
\begin{leansrc}
theorem subunit_weight_fails :
    ∃ (c : Fin 1 → ℚ) (f : Fin 1 → FinPMF (Fin 2) → ℚ),
      (∀ i ν, 0 ≤ f i ν) ∧ (∀ i, 0 ≤ c i) ∧
      ¬ ∀ (i : Fin 1) (ν : FinPMF (Fin 2)),
          f i ν ≤ compositePenaltyWQ (Xi := fun _ => Fin 2) c
            (fun _ ν => ν) f ν
\end{leansrc}
\par\medskip\noindent{\small \texttt{SuFunctor.Plain.}\texttt{\bfseries jointEquivBel}}\hfill{\footnotesize \path{lean/su-functor/SuFunctor/Plain/Seam.lean}, line 100}\par\nopagebreak
\begin{leansrc}
def jointEquivBel : Joint V ≃ Universe.BelState V where …
\end{leansrc}
\par\medskip\noindent{\small \texttt{SuFunctor.Plain.}\texttt{\bfseries jointEquivProb}}\hfill{\footnotesize \path{lean/su-functor/SuFunctor/Plain/Seam.lean}, line 147}\par\nopagebreak
\begin{leansrc}
def jointEquivProb :
    {μ : Joint V // μ.wt.sum (fun _ v => v) = 1} ≃ Universe.ProbState V where …
\end{leansrc}
\par\medskip\noindent{\small \texttt{SuFunctor.Plain.}\texttt{\bfseries penEquiv}}\hfill{\footnotesize \path{lean/su-functor/SuFunctor/Plain/Seam.lean}, line 200}\par\nopagebreak
\begin{leansrc}
def penEquiv : Pen V ≃o (Universe.BelState V → ℝ≥0∞) where …
\end{leansrc}
\par\medskip\noindent{\small \texttt{SuFunctor.Plain.Joint.}\texttt{\bfseries toMeas}}\hfill{\footnotesize \path{lean/su-functor/SuFunctor/Plain/GeneralSeam.lean}, line 76}\par\nopagebreak
\begin{leansrc}
def toMeas (μ : Joint X) : MeasJoint X
\end{leansrc}
\par\medskip\noindent{\small \texttt{SuFunctor.Plain.Joint.}\texttt{\bfseries pull\_toMeas}}\hfill{\footnotesize \path{lean/su-functor/SuFunctor/Plain/GeneralSeam.lean}, line 150}\par\nopagebreak
\begin{leansrc}
theorem pull_toMeas (h : X → Y) (hm : Measurable h) (F : MeasPen Y) :
    (MeasPen.pull h F) ∘ toMeas = Pen.pull h (F ∘ toMeas)
\end{leansrc}
\par\medskip\noindent{\small \texttt{SuFunctor.Plain.Joint.}\texttt{\bfseries image\_pull\_toMeas}}\hfill{\footnotesize \path{lean/su-functor/SuFunctor/Plain/GeneralSeam.lean}, line 159}\par\nopagebreak
\begin{leansrc}
theorem image_pull_toMeas (h : X → Y) (hm : Measurable h)
    (𝒮 : Set (MeasPen Y)) :
    (fun F => F ∘ toMeas) '' (MeasPen.pull h '' 𝒮) =
      Pen.pull h '' ((fun F => F ∘ toMeas) '' 𝒮)
\end{leansrc}
\par\medskip\noindent{\small \texttt{SuFunctor.Universe.General.}\texttt{\bfseries coupledSat\_pullPen\_iff}}\hfill{\footnotesize \path{lean/su-functor/SuFunctor/Universe/General/TransportGeneral.lean}, line 239}\par\nopagebreak
\begin{leansrc}
theorem coupledSat_pullPen_iff (hθ : Continuous θ)
    (hsurj : Function.Surjective (probPush hθ))
    {σ u : Penalty Y} {𝒱 : Set (Penalty Y)} :
    CoupledSat (pullPen hθ σ) (pullPen hθ u) (pullPen hθ '' 𝒱)
      ↔ CoupledSat σ u 𝒱
\end{leansrc}
\par\medskip\noindent{\small \texttt{SuFunctor.Universe.General.}\texttt{\bfseries transportClass\_crescentClass}}\hfill{\footnotesize \path{lean/su-functor/SuFunctor/Universe/General/TransportGeneral.lean}, line 418}\par\nopagebreak
\begin{leansrc}
theorem transportClass_crescentClass [Nonempty X] (hθ : Continuous θ)
    (hsurj : Function.Surjective
      (push θ : SubProbabilityMeasure X → SubProbabilityMeasure Y))
    (hmap : IsClosedMap
      (push θ : SubProbabilityMeasure X → SubProbabilityMeasure Y))
    (u : Penalty Y) (𝒱 : Set (Penalty Y)) :
    transportClass θ (crescentClass u 𝒱)
      = crescentClass (pullPen hθ u) (pullPen hθ '' 𝒱)
        ∩ {L | Saturated θ (L : Set (SubProbabilityMeasure X))}
\end{leansrc}
\par\medskip\noindent{\small \texttt{SuFunctor.Universe.General.}\texttt{\bfseries exists\_transport\_not\_subset\_pulled\_crescent}}\hfill{\footnotesize \path{lean/su-functor/SuFunctor/Universe/General/TransportGeneral.lean}, line 665}\par\nopagebreak
\begin{leansrc}
theorem exists_transport_not_subset_pulled_crescent :
    ∃ (θ : Unit → Bool) (hθ : Continuous θ) (u : Penalty Bool)
      (𝒱 : Set (Penalty Bool)),
      LowerSemicontinuous u ∧ (∀ v ∈ 𝒱, LowerSemicontinuous v) ∧ 𝒱.Finite ∧
      ¬ (transportClass θ (crescentClass u 𝒱)
          ⊆ crescentClass (pullPen hθ u) (pullPen hθ '' 𝒱))
\end{leansrc}
\par\medskip\noindent{\small \texttt{SuFunctor.Universe.General.}\texttt{\bfseries compile\_add\_radProd}}\hfill{\footnotesize \path{lean/su-functor/SuFunctor/Universe/General/TransportGeneral.lean}, line 478}\par\nopagebreak
\begin{leansrc}
theorem compile_add_radProd [Nonempty X] (f g : Penalty X) :
    compile (f + g) = radProd (compile f) (compile g)
\end{leansrc}
\par\medskip\noindent{\small \texttt{SuFunctor.Universe.General.}\texttt{\bfseries probPush\_surjective\_of\_section}}\hfill{\footnotesize \path{lean/su-functor/SuFunctor/Universe/General/TransportGeneral.lean}, line 208}\par\nopagebreak
\begin{leansrc}
theorem probPush_surjective_of_section (hθ : Continuous θ) {s : Y → X}
    (hs_mble : Measurable s) (hs : Function.RightInverse s θ) :
    Function.Surjective (probPush hθ)
\end{leansrc}
\par\medskip\noindent{\small \texttt{SuFunctor.Universe.}\texttt{\bfseries cross\_site\_union\_refinement\_iff}}\hfill{\footnotesize \path{lean/su-functor/SuFunctor/Universe/DemonicClock.lean}, line 262}\par\nopagebreak
\begin{leansrc}
theorem cross_site_union_refinement_iff {P : (K → ℝ≥0∞) → Prop}
    {u w : K → ℝ≥0∞} {𝒱 : Set (K → ℝ≥0∞)} {u' : Fin m → K → ℝ≥0∞}
    {𝒱' : Fin m → Set (K → ℝ≥0∞)}
    (hPc : ∀ c, SelValid 𝒱' c → P (selSup u c ⊔ w)) :
    (∀ σ, P σ → CoupledSat σ (u ⊔ w) 𝒱 → ∃ j, CoupledSat σ (u' j) (𝒱' j))
      ↔ ∀ c, SelValid 𝒱' c → Blocked (u ⊔ w) 𝒱 u' c
\end{leansrc}
\par\medskip\noindent{\small \texttt{SuFunctor.Universe.}\texttt{\bfseries lowerSemicontinuous\_selSup}}\hfill{\footnotesize \path{lean/su-functor/SuFunctor/Universe/DemonicClock.lean}, line 279}\par\nopagebreak
\begin{leansrc}
theorem lowerSemicontinuous_selSup {u : K → ℝ≥0∞}
    {c : Fin m → Option (K → ℝ≥0∞)} (hu : LowerSemicontinuous u)
    (hc : ∀ j w, c j = some w → LowerSemicontinuous w) :
    LowerSemicontinuous (selSup u c)
\end{leansrc}
\par\medskip\noindent{\small \texttt{SuFunctor.Universe.General.}\texttt{\bfseries probWall}}\hfill{\footnotesize \path{lean/su-functor/SuFunctor/Universe/General/DemonicGeneral.lean}, line 49}\par\nopagebreak
\begin{leansrc}
noncomputable def probWall (B : Set (ProbabilityMeasure X)) : Penalty X
\end{leansrc}
\par\medskip\noindent{\small \texttt{SuFunctor.Universe.}\texttt{\bfseries crescentClass\_inter\_wall}}\hfill{\footnotesize \path{lean/su-functor/SuFunctor/Universe/DemonicClock.lean}, line 323}\par\nopagebreak
\begin{leansrc}
theorem crescentClass_inter_wall (u w : ProbPenalty V)
    (𝒱 : Set (ProbPenalty V)) :
    crescentClass u 𝒱 ∩ crescentClass w ∅ = crescentClass (u ⊔ w) 𝒱
\end{leansrc}
\par\medskip\noindent{\small \texttt{SuFunctor.Universe.}\texttt{\bfseries cross\_site\_union\_refinement\_iff\_probWall}}\hfill{\footnotesize \path{lean/su-functor/SuFunctor/Universe/DemonicClock.lean}, line 332}\par\nopagebreak
\begin{leansrc}
theorem cross_site_union_refinement_iff_probWall {B : Set (ProbState V)}
    (hB : IsClosed B) {u : ProbPenalty V} {𝒱 : Set (ProbPenalty V)}
    {u' : Fin m → ProbPenalty V} {𝒱' : Fin m → Set (ProbPenalty V)}
    (hu : LowerSemicontinuous u)
    (h𝒱' : ∀ j, ∀ v ∈ 𝒱' j, LowerSemicontinuous v) :
    (∀ σ, LowerSemicontinuous σ → CoupledSat σ (u ⊔ probWall B) 𝒱 →
        ∃ j, CoupledSat σ (u' j) (𝒱' j))
      ↔ ∀ c, SelValid 𝒱' c → Blocked (u ⊔ probWall B) 𝒱 u' c
\end{leansrc}
\par\medskip\noindent{\small \texttt{SuFunctor.Universe.}\texttt{\bfseries exists\_no\_functional\_transport}}\hfill{\footnotesize \path{lean/su-functor/SuFunctor/Universe/DemonicClock.lean}, line 186}\par\nopagebreak
\begin{leansrc}
theorem exists_no_functional_transport :
    ∃ r : Bool → Bool → Prop, ¬ ∃ f : Bool → Bool, ∀ x y, r x y ↔ f x = y
\end{leansrc}
\par\medskip\noindent{\small \texttt{SuFunctor.Universe.}\texttt{\bfseries emLift\_comp}}\hfill{\footnotesize \path{lean/su-functor/SuFunctor/Universe/DemonicClock.lean}, line 143}\par\nopagebreak
\begin{leansrc}
theorem emLift_comp (r : X → Y → Prop) (s : Y → Z → Prop)
    (A : Set X) (C : Set Z) :
    emLift (relComp r s) A C ↔ ∃ B, emLift r A B ∧ emLift s B C
\end{leansrc}
\par\medskip\noindent{\small \texttt{SuFunctor.Spacetime.}\texttt{\bfseries spacetimeSpecFunctor\_fibre\_eq}}\hfill{\footnotesize \path{lean/su-functor/SuFunctor/Spacetime/Universe.lean}, line 141}\par\nopagebreak
\begin{leansrc}
theorem spacetimeSpecFunctor_fibre_eq (A : SpacetimeBase Clk Rec) :
    (spacetimeSpecFunctor Clk Rec).toFunctor.obj A
      = SpecDnf (Joint (A.wires → ClockedEl Clk Rec))
\end{leansrc}
\par\medskip\noindent{\small \texttt{SuFunctor.Spacetime.}\texttt{\bfseries profileLiftHom}}\hfill{\footnotesize \path{lean/su-functor/SuFunctor/Spacetime/Base.lean}, line 364}\par\nopagebreak
\begin{leansrc}
def profileLiftHom (S : Type u) (φ : S → B.wires) :
    profileLift (Rec := Rec) S φ ⟶ B
\end{leansrc}
\par\medskip\noindent{\small \texttt{SuFunctor.Spacetime.}\texttt{\bfseries profile\_unique\_of\_hom}}\hfill{\footnotesize \path{lean/su-functor/SuFunctor/Spacetime/Base.lean}, line 373}\par\nopagebreak
\begin{leansrc}
theorem profile_unique_of_hom {S : Type u} {κ : S → Clk}
    {φ : S → B.wires} (h : ∀ w, B.profile (φ w) = κ w) :
    κ = fun w => B.profile (φ w)
\end{leansrc}
\par\medskip\noindent{\small \texttt{SuFunctor.Spacetime.}\texttt{\bfseries mixedInterface}}\hfill{\footnotesize \path{lean/su-functor/SuFunctor/Spacetime/Base.lean}, line 382}\par\nopagebreak
\begin{leansrc}
def mixedInterface (c c' : Clk) : SpacetimeBase Clk Rec
\end{leansrc}
\par\medskip\noindent{\small \texttt{SuFunctor.Spacetime.}\texttt{\bfseries profiledEmbed}}\hfill{\footnotesize \path{lean/su-functor/SuFunctor/Spacetime/Universe.lean}, line 94}\par\nopagebreak
\begin{leansrc}
def profiledEmbed (Clk : Type u) (Rec : Clk → Type u) :
    LaxDoubleFunctor (SpacetimeBase Clk Rec)
      (WireSpan (ClockedEl Clk Rec)) where …
\end{leansrc}
\par\medskip\noindent{\small \texttt{SuFunctor.Spacetime.}\texttt{\bfseries SemanticUniverse}}\hfill{\footnotesize \path{lean/su-functor/SuFunctor/Spacetime/Universe.lean}, line 159}\par\nopagebreak
\begin{leansrc}
abbrev SemanticUniverse (Clk : Type u) (Rec : Clk → Type u)
\end{leansrc}
\par\medskip\noindent{\small \texttt{SuFunctor.Spacetime.}\texttt{\bfseries universeProj\_discreteOpfibration}}\hfill{\footnotesize \path{lean/su-functor/SuFunctor/Spacetime/Universe.lean}, line 185}\par\nopagebreak
\begin{leansrc}
theorem universeProj_discreteOpfibration
    (Clk : Type u) (Rec : Clk → Type u) :
    IsDiscreteDoubleOpfibration (universeProj Clk Rec)
\end{leansrc}
\par\medskip\noindent{\small \texttt{SuFunctor.Spacetime.}\texttt{\bfseries constHomEquiv}}\hfill{\footnotesize \path{lean/su-functor/SuFunctor/Spacetime/Universe.lean}, line 216}\par\nopagebreak
\begin{leansrc}
def constHomEquiv (S T : WireSpan (Rec c)) :
    (S ⟶ T) ≃ (constObj (Rec := Rec) c S.wires ⟶ constObj c T.wires) where …
\end{leansrc}
\par\medskip\noindent{\small \texttt{SuFunctor.Universe.}\texttt{\bfseries fibInf}}\hfill{\footnotesize \path{lean/su-functor/SuFunctor/Universe/CertificateLayer.lean}, line 79}\par\nopagebreak
\begin{leansrc}
def fibInf (g : K₁ → K₂) (τ : K₁ → ℝ≥0∞) : K₂ → ℝ≥0∞
\end{leansrc}
\par\medskip\noindent{\small \texttt{SuFunctor.Universe.}\texttt{\bfseries coupledSat\_comap\_iff\_fibInf}}\hfill{\footnotesize \path{lean/su-functor/SuFunctor/Universe/CertificateLayer.lean}, line 128}\par\nopagebreak
\begin{leansrc}
theorem coupledSat_comap_iff_fibInf {g : K₁ → K₂} {τ : K₁ → ℝ≥0∞}
    {u : K₂ → ℝ≥0∞} {𝒱 : Set (K₂ → ℝ≥0∞)} :
    CoupledSat τ (u ∘ g) ((· ∘ g) '' 𝒱) ↔ CoupledSat (fibInf g τ) u 𝒱
\end{leansrc}
\par\medskip\noindent{\small \texttt{SuFunctor.Universe.SpecDnf.}\texttt{\bfseries comap\_modelClass\_mono}}\hfill{\footnotesize \path{lean/su-functor/SuFunctor/Universe/CertificateLayer.lean}, line 181}\par\nopagebreak
\begin{leansrc}
theorem SpecDnf.comap_modelClass_mono {g : K₁ → K₂} {D E : SpecDnf K₂}
    (h : D.modelClass ⊆ E.modelClass) :
    (D.comap g).modelClass ⊆ (E.comap g).modelClass
\end{leansrc}
\par\medskip\noindent{\small \texttt{SuFunctor.Universe.}\texttt{\bfseries exists\_comp\_notMem\_modelClass\_comap}}\hfill{\footnotesize \path{lean/su-functor/SuFunctor/Universe/CertificateLayer.lean}, line 213}\par\nopagebreak
\begin{leansrc}
theorem exists_comp_notMem_modelClass_comap :
    ∃ (g : Unit → Bool) (a : SpecAtom Bool) (σ : Bool → ℝ≥0∞),
      σ ∈ a.modelClass ∧ σ ∘ g ∉ (a.comap g).modelClass
\end{leansrc}
\par\medskip\noindent{\small \texttt{SuFunctor.Universe.}\texttt{\bfseries coupledSat\_comp\_add\_of\_dips}}\hfill{\footnotesize \path{lean/su-functor/SuFunctor/Universe/CertificateLayer.lean}, line 258}\par\nopagebreak
\begin{leansrc}
theorem coupledSat_comp_add_of_dips (g₁ : K → K₁) (g₂ : K → K₂)
    {σ₁ u₁ : K₁ → ℝ≥0∞} {σ₂ u₂ : K₂ → ℝ≥0∞}
    {𝒱₁ : Set (K₁ → ℝ≥0∞)} {𝒱₂ : Set (K₂ → ℝ≥0∞)}
    (h₁ : CoupledSat σ₁ u₁ 𝒱₁) (h₂ : CoupledSat σ₂ u₂ 𝒱₂)
    (hcon : ∀ w ∈ (· ∘ g₁) '' 𝒱₁ ∪ (· ∘ g₂) '' 𝒱₂,
      Dip ((σ₁ ∘ g₁) + (σ₂ ∘ g₂)) w) :
    CoupledSat ((σ₁ ∘ g₁) + (σ₂ ∘ g₂)) ((u₁ ∘ g₁) + (u₂ ∘ g₂))
      ((· ∘ g₁) '' 𝒱₁ ∪ (· ∘ g₂) '' 𝒱₂)
\end{leansrc}
\par\medskip\noindent{\small \texttt{SuFunctor.Universe.}\texttt{\bfseries exists\_boxplus\_not\_semantically\_monotone}}\hfill{\footnotesize \path{lean/su-functor/SuFunctor/Universe/CertificateLayer.lean}, line 285}\par\nopagebreak
\begin{leansrc}
theorem exists_boxplus_not_semantically_monotone :
    ∃ (a e e₁ e₂ : SpecAtom Unit) (τ : Unit → ℝ≥0∞),
      e.modelClass = e₁.modelClass ∪ e₂.modelClass ∧
      CoupledSat τ (a.floor + e.floor) (a.marks ∪ e.marks) ∧
      ¬ CoupledSat τ (a.floor + e₁.floor) (a.marks ∪ e₁.marks) ∧
      ¬ CoupledSat τ (a.floor + e₂.floor) (a.marks ∪ e₂.marks)
\end{leansrc}
\par\medskip\noindent{\small \texttt{SuFunctor.Universe.SuFunctor.}\texttt{\bfseries dnfBoxplus\_mono}}\hfill{\footnotesize \path{lean/su-functor/SuFunctor/Universe/CertificateLayer.lean}, line 405}\par\nopagebreak
\begin{leansrc}
theorem dnfBoxplus_mono {A B C : WireSpan El} {M : Pat A B} {N : Pat B C}
    {D₁ E₁ : DnfCpl M} {D₂ E₂ : DnfCpl N}
    (h₁ : D₁.atoms ⊆ E₁.atoms) (h₂ : D₂.atoms ⊆ E₂.atoms) :
    (dnfBoxplus D₁ D₂).atoms ⊆ (dnfBoxplus E₁ E₂).atoms
\end{leansrc}
\par\medskip\noindent{\small \texttt{SuFunctor.Spacetime.}\texttt{\bfseries spacetimeSpecFunctor}}\hfill{\footnotesize \path{lean/su-functor/SuFunctor/Spacetime/Universe.lean}, line 134}\par\nopagebreak
\begin{leansrc}
noncomputable def spacetimeSpecFunctor (Clk : Type u) (Rec : Clk → Type u) :
    SpanValued.{u + 1, u, u, u + 1, u} (SpacetimeBase Clk Rec)
\end{leansrc}
\par\medskip\noindent{\small \texttt{SuFunctor.Spacetime.}\texttt{\bfseries semanticUniverse\_pseudoDoubleCategory}}\hfill{\footnotesize \path{lean/su-functor/SuFunctor/Spacetime/Universe.lean}, line 164}\par\nopagebreak
\begin{leansrc}
theorem semanticUniverse_pseudoDoubleCategory
    (Clk : Type u) (Rec : Clk → Type u) :
    ∃ _ : PseudoDoubleCategory (SemanticUniverse Clk Rec), True
\end{leansrc}
\par\medskip\noindent{\small \texttt{SuFunctor.Spacetime.}\texttt{\bfseries universeProj\_isPseudo}}\hfill{\footnotesize \path{lean/su-functor/SuFunctor/Spacetime/Universe.lean}, line 177}\par\nopagebreak
\begin{leansrc}
theorem universeProj_isPseudo (Clk : Type u) (Rec : Clk → Type u) :
    (universeProj Clk Rec).IsPseudo
\end{leansrc}
\par\medskip\noindent{\small \texttt{SuFunctor.Spacetime.}\texttt{\bfseries profiledEmbed\_faithful\_tight}}\hfill{\footnotesize \path{lean/su-functor/SuFunctor/Spacetime/Universe.lean}, line 122}\par\nopagebreak
\begin{leansrc}
theorem profiledEmbed_faithful_tight {A B : SpacetimeBase Clk Rec}
    {φ ψ : A ⟶ B}
    (h : (profiledEmbed Clk Rec).map φ = (profiledEmbed Clk Rec).map ψ) :
    φ = ψ
\end{leansrc}
\par\medskip\noindent{\small \texttt{SuFunctor.Spacetime.}\texttt{\bfseries constProfile}}\hfill{\footnotesize \path{lean/su-functor/SuFunctor/Spacetime/Universe.lean}, line 298}\par\nopagebreak
\begin{leansrc}
def constProfile :
    LaxDoubleFunctor (WireSpan (Rec c)) (SpacetimeBase Clk Rec) where …
\end{leansrc}
\par\medskip\noindent{\small \texttt{SuFunctor.Spacetime.}\texttt{\bfseries constProfile\_isPseudo}}\hfill{\footnotesize \path{lean/su-functor/SuFunctor/Spacetime/Universe.lean}, line 340}\par\nopagebreak
\begin{leansrc}
theorem constProfile_isPseudo :
    (constProfile (Rec := Rec) c).IsPseudo where
  laxComp_isInvertible {_ T _} M N
\end{leansrc}
\par\medskip\noindent{\small \texttt{SuFunctor.Spacetime.}\texttt{\bfseries siteRead}}\hfill{\footnotesize \path{lean/su-functor/SuFunctor/Spacetime/Universe.lean}, line 273}\par\nopagebreak
\begin{leansrc}
noncomputable def siteRead (c : Clk) {S : Type u}
    (D : SpecDnf (Joint (S → ClockedEl Clk Rec))) :
    SpecDnf (Joint (S → Rec c))
\end{leansrc}
\par\medskip\noindent{\small \texttt{SuFunctor.Spacetime.SuFunctor.Universe.SpecDnf.}\texttt{\bfseries mem\_modelClass\_withWall}}\hfill{\footnotesize \path{lean/su-functor/SuFunctor/Spacetime/Universe.lean}, line 408}\par\nopagebreak
\begin{leansrc}
theorem SpecDnf.mem_modelClass_withWall {D : SpecDnf K}
    {w σ : K → ℝ≥0∞} :
    σ ∈ (D.withWall w).modelClass ↔ w ≤ σ ∧ σ ∈ D.modelClass
\end{leansrc}
\par\medskip\noindent{\small \texttt{SuFunctor.Adequacy.}\texttt{\bfseries rwall\_singletonC}}\hfill{\footnotesize \path{lean/su-functor/SuFunctor/Universe/RadialUniverse.lean}, line 303}\par\nopagebreak
\begin{leansrc}
theorem rwall_singletonC (μ : SubPMF V) :
    rwall (singletonC μ) = rayPoint μ
\end{leansrc}
\par\medskip\noindent{\small \texttt{SuFunctor.Adequacy.}\texttt{\bfseries idealPoint}}\hfill{\footnotesize \path{lean/su-functor/SuFunctor/Adequacy/SemanticUniverseWiring.lean}, line 213}\par\nopagebreak
\begin{leansrc}
noncomputable def idealPoint (μ : SubPMF V) : Plain.Pen V
\end{leansrc}
\par\medskip\noindent{\small \texttt{SuFunctor.SubPMF.}\texttt{\bfseries toJoint}}\hfill{\footnotesize \path{lean/su-functor/SuFunctor/Adequacy/SemanticUniverseWiring.lean}, line 99}\par\nopagebreak
\begin{leansrc}
noncomputable def toJoint (μ : SubPMF V) : Plain.Joint V where …
\end{leansrc}
\par\medskip\noindent{\small \texttt{SuFunctor.Adequacy.}\texttt{\bfseries diagramDnf}}\hfill{\footnotesize \path{lean/su-functor/SuFunctor/Adequacy/SemanticUniverseWiring.lean}, line 385}\par\nopagebreak
\begin{leansrc}
noncomputable def diagramDnf (L : List (SubPMF V → Prop)) :
    SpecDnf (Plain.Joint V) where …
\end{leansrc}
\par\medskip\noindent{\small \texttt{SuFunctor.Adequacy.}\texttt{\bfseries coupledSat\_idealPoint\_cwall\_iff}}\hfill{\footnotesize \path{lean/su-functor/SuFunctor/Adequacy/SemanticUniverseWiring.lean}, line 295}\par\nopagebreak
\begin{leansrc}
theorem coupledSat_idealPoint_cwall_iff {c : SubPMF V → Prop}
    {μ : SubPMF V} :
    CoupledSat (idealPoint μ) (cwall c) {(⊤ : Plain.Pen V)} ↔ c μ
\end{leansrc}
\par\medskip\noindent{\small \texttt{SuFunctor.Adequacy.}\texttt{\bfseries idealPoint\_mem\_diagramDnf\_iff}}\hfill{\footnotesize \path{lean/su-functor/SuFunctor/Adequacy/SemanticUniverseWiring.lean}, line 394}\par\nopagebreak
\begin{leansrc}
theorem idealPoint_mem_diagramDnf_iff (L : List (SubPMF V → Prop))
    (μ : SubPMF V) :
    idealPoint μ ∈ (diagramDnf L).modelClass ↔ looseConj L μ
\end{leansrc}
\par\medskip\noindent{\small \texttt{SuFunctor.Plain.}\texttt{\bfseries semanticallyValid\_iff\_subset}}\hfill{\footnotesize \path{lean/su-functor/SuFunctor/Plain/CheckerFace.lean}, line 215}\par\nopagebreak
\begin{leansrc}
theorem semanticallyValid_iff_subset (H : Hypernet) (seq : Sequent) :
    seq.semanticallyValid H ↔ admCyl H seq ⊆ postEvt H seq
\end{leansrc}
\par\medskip\noindent{\small \texttt{SuFunctor.Plain.}\texttt{\bfseries upperProb\_eq\_zero\_iff}}\hfill{\footnotesize \path{lean/su-functor/SuFunctor/Plain/CheckerFace.lean}, line 177}\par\nopagebreak
\begin{leansrc}
theorem upperProb_eq_zero_iff (B E : Set X) :
    upperProb B E = 0 ↔ E ∩ B = ∅
\end{leansrc}
\par\medskip\noindent{\small \texttt{SuFunctor.Plain.}\texttt{\bfseries semanticallyValid\_iff\_face}}\hfill{\footnotesize \path{lean/su-functor/SuFunctor/Plain/CheckerFace.lean}, line 224}\par\nopagebreak
\begin{leansrc}
theorem semanticallyValid_iff_face (H : Hypernet) (seq : Sequent) :
    seq.semanticallyValid H ↔
      upperProb (admCyl H seq) (postEvt H seq)ᶜ = 0
\end{leansrc}
\par\medskip\noindent{\small \texttt{SuFunctor.Plain.}\texttt{\bfseries suppIn\_empty}}\hfill{\footnotesize \path{lean/su-functor/SuFunctor/Plain/CheckerFace.lean}, line 78}\par\nopagebreak
\begin{leansrc}
theorem suppIn_empty : suppIn (∅ : Set X) = {Joint.zero X}
\end{leansrc}
\par\medskip\noindent{\small \texttt{SuFunctor.Plain.}\texttt{\bfseries face\_ex\_falso}}\hfill{\footnotesize \path{lean/su-functor/SuFunctor/Plain/CheckerFace.lean}, line 243}\par\nopagebreak
\begin{leansrc}
theorem face_ex_falso (H : Hypernet) (seq : Sequent)
    (h : admCyl H seq = ∅) :
    upperProb (admCyl H seq) (postEvt H seq)ᶜ = 0
\end{leansrc}
\par\medskip\noindent{\small \texttt{SuFunctor.}\texttt{\bfseries hnetInclusion}}\hfill{\footnotesize \path{lean/su-functor/SuFunctor/HypernetDouble.lean}, line 559}\par\nopagebreak
\begin{leansrc}
def hnetInclusion : LaxDoubleFunctor (HNet ty) (SingleSiteSemanticUniverse ElVal)
\end{leansrc}
\par\medskip\noindent{\small \texttt{SuFunctor.Universe.General.}\texttt{\bfseries sat\_wall\_iff}}\hfill{\footnotesize \path{lean/su-functor/SuFunctor/Universe/General/CheckerAnchor.lean}, line 82}\par\nopagebreak
\begin{leansrc}
theorem sat_wall_iff {σ : StatePenalty X}
    {C : Set (SubProbabilityMeasure X)} :
    Sat σ (wall C) ↔ ∀ ν, σ ν ≠ ⊤ → ν ∈ C
\end{leansrc}
\par\medskip\noindent{\small \texttt{SuFunctor.Universe.General.}\texttt{\bfseries sat\_wall\_wall\_iff}}\hfill{\footnotesize \path{lean/su-functor/SuFunctor/Universe/General/CheckerAnchor.lean}, line 92}\par\nopagebreak
\begin{leansrc}
theorem sat_wall_wall_iff {S C : Set (SubProbabilityMeasure X)} :
    Sat (wall S) (wall C) ↔ S ⊆ C
\end{leansrc}
\par\medskip\noindent{\small \texttt{SuFunctor.Universe.General.}\texttt{\bfseries sat\_iff\_sublevel\_subset}}\hfill{\footnotesize \path{lean/su-functor/SuFunctor/Universe/General/CheckerAnchor.lean}, line 111}\par\nopagebreak
\begin{leansrc}
theorem sat_iff_sublevel_subset {σ φ : StatePenalty X} :
    Sat σ φ ↔ ∀ s : ℝ≥0∞,
      {ν : SubProbabilityMeasure X | σ ν ≤ s}
        ⊆ {ν : SubProbabilityMeasure X | φ ν ≤ s}
\end{leansrc}
\par\medskip\noindent{\small \texttt{PPLC.}\texttt{\bfseries OpenHypernet}}\hfill{\footnotesize \path{lean/qtsl/Qtsl/HypernetCospan.lean}, line 222}\par\nopagebreak
\begin{leansrc}
structure OpenHypernet where …
\end{leansrc}
\par\medskip\noindent{\small \texttt{PPLC.OpenHypernet.}\texttt{\bfseries glue\_admissible\_iff}}\hfill{\footnotesize \path{lean/qtsl/Qtsl/HypernetCospan.lean}, line 444}\par\nopagebreak
\begin{leansrc}
theorem glue_admissible_iff (σ : (glueNet hA).Valuation) :
    σ.Admissible ↔
      (∀ a ∈ H.net.applications, a.SemOk σ) ∧
      (∀ a ∈ G.net.applications, a.SemOk σ)
\end{leansrc}
\par\medskip\noindent{\small \texttt{PPLC.OpenHypernet.}\texttt{\bfseries unitOpen}}\hfill{\footnotesize \path{lean/qtsl/Qtsl/HypernetCospan.lean}, line 465}\par\nopagebreak
\begin{leansrc}
def unitOpen (X : List Wire) (wt₀ : Wire → Option TensorType)
    (hX : X.Nodup) : OpenHypernet where …
\end{leansrc}
\par\medskip\noindent{\small \texttt{PPLC.OpenHypernet.}\texttt{\bfseries glue\_assoc}}\hfill{\footnotesize \path{lean/qtsl/Qtsl/HypernetCospan.lean}, line 558}\par\nopagebreak
\begin{leansrc}
theorem glue_assoc {H G K : OpenHypernet}
    (hHG : Aligned H G) (hGK : Aligned G K)
    (hL : Aligned (glue hHG) K) (hR : Aligned H (glue hGK)) :
    glue hL = glue hR
\end{leansrc}
\par\medskip\noindent{\small \texttt{SuFunctor.}\texttt{\bfseries theta\_rename}}\hfill{\footnotesize \path{lean/su-functor/SuFunctor/HypernetInclusion.lean}, line 912}\par\nopagebreak
\begin{leansrc}
theorem theta_rename {m : AlignedLoose X Y} {m' : AlignedLoose X' Y'}
    (c : AlignedCell ty m m') (T : List PrimApp)
    (hTC : ∀ a ∈ T, a.TypeConsistent (tableOn ty m.apex))
    (hw : ∀ a ∈ T, ∀ w ∈ a.inputWires ++ a.outputWires, w ∈ m.apex) :
    dnfPull (renameTot ty c) (austereDnf ty m T)
      = austereDnf ty m' (T.map (PrimApp.rename c.map))
\end{leansrc}
\par\medskip\noindent{\small \texttt{PPLC.ProofObject.}\texttt{\bfseries sound}}\hfill{\footnotesize \path{lean/qtsl/Qtsl/Soundness.lean}, line 8330}\par\nopagebreak
\begin{leansrc}
theorem ProofObject.sound {H : Hypernet} {seq : Sequent}
    (pf : ProofObject H seq) : seq.semanticallyValid H
\end{leansrc}
\par\medskip\noindent{\small \texttt{PPLC.}\texttt{\bfseries check\_sound}}\hfill{\footnotesize \path{lean/qtsl/Qtsl/Soundness.lean}, line 8439}\par\nopagebreak
\begin{leansrc}
theorem check_sound {H : Hypernet} {seq : Sequent} (pf : ProofObject H seq)
    (_hcheck : check pf = true) : seq.semanticallyValid H
\end{leansrc}
\par\medskip\noindent Modules cited: \texttt{QtslAux.DiamondSite}, \texttt{QtslAux.DiamondTopology}, \texttt{SuFunctor.Site.TemporalTypeTheory}, \texttt{SuFunctor.Site.BehaviorSpec}, \texttt{SuFunctor.Site.CausalConcordance}, \texttt{SuFunctor.StateSpace}, \texttt{SuFunctor.StateSpace.FamCat}, \texttt{QtslAux.QuantumCovariance}, \texttt{QtslAux.OpenQuantum}, \texttt{QtslAux.QuantumRecords}, \texttt{QtslAux.QuantumDescent}, \texttt{QtslAux.QuantumStates}, \texttt{QtslAux.QuantumAssembly}, \texttt{QtslAux.QuantumDeterminacy}, \texttt{QtslAux.QuantumOuterStorey}, \texttt{SuFunctor.Universe.FellTopology}, \texttt{SuFunctor.Universe.FellRestrict}, \texttt{SuFunctor.Universe.FellSeamWitness}, \texttt{SuFunctor.Universe.FellDescent}, \texttt{SuFunctor.Universe.FellRecordLocus}, \texttt{SuFunctor.Universe.FellRecordWitness}, \texttt{SuFunctor.Universe.FellOnePoint}, \texttt{SuFunctor.Universe.FellCapacity}, \texttt{Pointfree.Frame.Presentation}, \texttt{Pointfree.Frame.HyperspacePresentation}, \texttt{Pointfree.Valuation.Basic}, \texttt{Pointfree.Valuation.Topology}, \texttt{Pointfree.Valuation.EscapeWitness}, \texttt{Pointfree.Valuation.Bridge}, \texttt{Pointfree.Measure.DaggerChart}, \texttt{SuFunctor.Universe.General}, \texttt{SuFunctor.Universe.Judgment}, \texttt{SuFunctor.Universe.Closure}, \texttt{SuFunctor.Universe.Generation}, \texttt{SuFunctor.Universe.Tents}, \texttt{QtslAux.Imprecise.MarginalCompat}, \texttt{QtslAux.Imprecise.MarginalCompatGeneral}, \texttt{SuFunctor.Universe.RadialUniverse}, \texttt{SuFunctor.StateSpace.FamCompletion}, \texttt{SuFunctor.Universe.VietorisStratum}, \texttt{SuFunctor.Universe.ImageChar}, \texttt{QtslAux.Imprecise.PenaltyCredalIso}, \texttt{SuFunctor.Universe.EpistemicStates}, \texttt{SuFunctor.Universe.General.CredalPoset}, \texttt{SuFunctor.Universe.General.Richardson}, \texttt{SuFunctor.Universe.CredalPoset}, \texttt{SuFunctor.Universe.HausdorffKL}, \texttt{SuFunctor.Adequacy.Envelope}, \texttt{SuFunctor.Adequacy.EnvelopeUnion}, \texttt{QtslAux.Imprecise.Credal}, \texttt{QtslAux.Imprecise.NaturalExtensionGeneral}, \texttt{QtslAux.Imprecise.NaturalExtension}, \texttt{QtslAux.LeakAware}, \texttt{SuFunctor.Universe.General.CrescentStratum}, \texttt{SuFunctor.Universe.HullClosure}, \texttt{SuFunctor.Site.GeneralWindow}, \texttt{SuFunctor.Plain.SeamDensity}, \texttt{Pointfree.MeasureAlgebra.Basic}, \texttt{Pointfree.Prevision.Basic}, \texttt{SuFunctor.QtslSem.SpecJudg}, \texttt{DoubleCat.Basic}, \texttt{DoubleCat.Functor.Lax}, \texttt{DoubleCat.DepSpan}, \texttt{SuFunctor.Universe.Marginalization}, \texttt{DoubleCat.WireSpan}, \texttt{SuFunctor.Plain.Penalty}, \texttt{SuFunctor.Plain.SpecFunctor}, \texttt{SuFunctor.Plain.Battery}, \texttt{SuFunctor.Site.WindowPenalty}, \texttt{SuFunctor.Site.SiteFunctor}, \texttt{SuFunctor.Site.Battery}, \texttt{SuFunctor.Site.MatchingFork}, \texttt{SuFunctor.Plain.GeneralFunctor}, \texttt{DoubleCat.Grothendieck}, \texttt{DoubleCat.Opfibration}, \texttt{DoubleCat.JudgLayer}, \texttt{QtslAux.Imprecise.Containment}, \texttt{QtslAux.Imprecise.ContainmentGeneral}, \texttt{QtslAux.Imprecise.CouplingGlue}, \texttt{QtslAux.Imprecise.FeasibilityGlue}, \texttt{SuFunctor.Plain.Seam}, \texttt{SuFunctor.Plain.GeneralSeam}, \texttt{SuFunctor.Universe.TransportCitizenship}, \texttt{SuFunctor.Universe.General.TransportGeneral}, \texttt{SuFunctor.Universe.DemonicClock}, \texttt{SuFunctor.Universe.General.DemonicGeneral}, \texttt{DoubleCat.CrossSite}, \texttt{SuFunctor.Spacetime.Base}, \texttt{SuFunctor.Universe.CertificateLayer}, \texttt{SuFunctor.Spacetime.Universe}, \texttt{SuFunctor.Target}, \texttt{SuFunctor.Adequacy.Inclusions}, \texttt{SuFunctor.Adequacy.Dtmc}, \texttt{SuFunctor.Adequacy.Pgm}, \texttt{SuFunctor.Adequacy.Mdp}, \texttt{SuFunctor.Adequacy.MdpDemonic}, \texttt{SuFunctor.Adequacy.DelayLine}, \texttt{SuFunctor.Adequacy.HeatTorus}, \texttt{SuFunctor.Adequacy.MarginalGlue}, \texttt{SuFunctor.Adequacy.HeatTorusInit}, \texttt{SuFunctor.Adequacy.SolutionSetMesh}, \texttt{SuFunctor.Plain.CheckerFace}, \texttt{SuFunctor.HypernetDouble}, \texttt{SuFunctor.Universe.General.CheckerAnchor}, \texttt{Qtsl.HypernetCospan}, \texttt{SuFunctor.HypernetInclusion}, \texttt{Qtsl.Soundness}.

\section{Hypernets: Primitive Functions}
Citations from Chapter~\ref{ch:qtsl-hypernets}.

\par\medskip\noindent{\small \texttt{PPLC.OpenHypernet.}\texttt{\bfseries glue\_unit\_left}}\hfill{\footnotesize \path{lean/qtsl/Qtsl/HypernetCospan.lean}, line 517}\par\nopagebreak
\begin{leansrc}
theorem glue_unit_left (H : OpenHypernet)
    (wt₀ : Wire → Option TensorType)
    (hsub : ∀ w τ, wt₀ w = some τ → H.net.wireType w = some τ)
    (hA : Aligned (unitOpen H.inputs wt₀ H.inputs_nodup) H) :
    glue hA = H
\end{leansrc}
\par\medskip\noindent{\small \texttt{PPLC.OpenHypernet.}\texttt{\bfseries glue\_unit\_right}}\hfill{\footnotesize \path{lean/qtsl/Qtsl/HypernetCospan.lean}, line 532}\par\nopagebreak
\begin{leansrc}
theorem glue_unit_right (H : OpenHypernet)
    (wt₀ : Wire → Option TensorType)
    (hsub : ∀ w τ, wt₀ w = some τ → H.net.wireType w = some τ)
    (hA : Aligned H (unitOpen H.outputs wt₀ H.outputs_nodup)) :
    glue hA = H
\end{leansrc}
\par\medskip\noindent{\small \texttt{SuFunctor.}\texttt{\bfseries theta\_unit}}\hfill{\footnotesize \path{lean/su-functor/SuFunctor/HypernetInclusion.lean}, line 842}\par\nopagebreak
\begin{leansrc}
theorem theta_unit (X : Finset Wire) :
    dnfPull (idTot ty X) (dnfUnit (ctxIface X))
      = austereDnf ty (AlignedLoose.unit X) []
\end{leansrc}
\par\medskip\noindent{\small \texttt{SuFunctor.}\texttt{\bfseries theta\_comp}}\hfill{\footnotesize \path{lean/su-functor/SuFunctor/HypernetInclusion.lean}, line 864}\par\nopagebreak
\begin{leansrc}
theorem theta_comp (m : AlignedLoose X Y) (n : AlignedLoose Y Z)
    (T₁ T₂ : List PrimApp)
    (hTC₁ : ∀ a ∈ T₁, a.TypeConsistent (tableOn ty m.apex))
    (hw₁ : ∀ a ∈ T₁, ∀ w ∈ a.inputWires ++ a.outputWires, w ∈ m.apex)
    (hTC₂ : ∀ a ∈ T₂, a.TypeConsistent (tableOn ty n.apex))
    (hw₂ : ∀ a ∈ T₂, ∀ w ∈ a.inputWires ++ a.outputWires, w ∈ n.apex) :
    dnfPull (splitTot ty m n)
        (dnfBoxplus (austereDnf ty m T₁) (austereDnf ty n T₂))
      = austereDnf ty (m.seq n) (T₁ ++ T₂)
\end{leansrc}
\par\medskip\noindent{\small \texttt{SuFunctor.}\texttt{\bfseries vacDnf\_dnfPull}}\hfill{\footnotesize \path{lean/su-functor/SuFunctor/HypernetInclusion.lean}, line 811}\par\nopagebreak
\begin{leansrc}
theorem vacDnf_dnfPull {A B : Type} (h : A → B) :
    dnfPull h (vacDnf (Plain.Joint B)) = vacDnf (Plain.Joint A)
\end{leansrc}
\par\medskip\noindent{\small \texttt{SuFunctor.}\texttt{\bfseries hnetInclusion\_pushGLoose}}\hfill{\footnotesize \path{lean/su-functor/SuFunctor/HypernetDouble.lean}, line 574}\par\nopagebreak
\begin{leansrc}
theorem hnetInclusion_pushGLoose {p q : HNet ty}
    (M : SpanValued.GLoose p q) {c d : AlignedCtx} {n : Loose c d}
    {f' : p.1 ⟶ c} {g' : q.1 ⟶ d} (γ : Cell M.base n f' g') :
    SpanValued.bmapLoose (stateBase ty) (hnetTheta ty)
        (SpanValued.pushGLoose M γ) ≍
      SpanValued.pushGLoose
        (SpanValued.bmapLoose (stateBase ty) (hnetTheta ty) M)
        ((stateBase ty).mapCell γ)
\end{leansrc}
\par\medskip\noindent{\small \texttt{SuFunctor.}\texttt{\bfseries austereDnf\_eq\_iff}}\hfill{\footnotesize \path{lean/su-functor/SuFunctor/HypernetInclusion.lean}, line 934}\par\nopagebreak
\begin{leansrc}
theorem austereDnf_eq_iff (m : AlignedLoose X Y) (T T' : List PrimApp) :
    austereDnf ty m T = austereDnf ty m T'
      ↔ AdmRegion ty m T = AdmRegion ty m T'
\end{leansrc}
\par\medskip\noindent{\small \texttt{PPLC.}\texttt{\bfseries substituteRaw\_eq\_glue}}\hfill{\footnotesize \path{lean/qtsl/Qtsl/SubstitutionBridge.lean}, line 1960}\par\nopagebreak
\begin{leansrc}
theorem substituteRaw_eq_glue :
    substituteRaw outer φ E inner =
      { wireTypes := bridgeTable outer φ E inner
        apps := (substituteOpen outer φ E inner
            hwfO hwfI hpo hti hto).net.applications
        ins := (substituteOpen outer φ E inner
            hwfO hwfI hpo hti hto).inputs
        outs := (substituteOpen outer φ E inner
            hwfO hwfI hpo hti hto).outputs }
\end{leansrc}
\par\medskip\noindent{\small \texttt{PPLC.}\texttt{\bfseries substitute\_eq\_glue\_renumber}}\hfill{\footnotesize \path{lean/qtsl/Qtsl/SubstitutionBridge.lean}, line 1985}\par\nopagebreak
\begin{leansrc}
theorem substitute_eq_glue_renumber :
    substitute outer φ E inner =
      (HypernetP.mk (bridgeTable outer φ E inner)
        ((substituteOpen outer φ E inner
            hwfO hwfI hpo hti hto).net.applications)
        ((substituteOpen outer φ E inner hwfO hwfI hpo hti hto).inputs)
        ((substituteOpen outer φ E inner
            hwfO hwfI hpo hti hto).outputs)).renumber
\end{leansrc}
\par\medskip\noindent{\small \texttt{PPLC.}\texttt{\bfseries substituteOpen\_admissible\_iff}}\hfill{\footnotesize \path{lean/qtsl/Qtsl/SubstitutionBridge.lean}, line 2001}\par\nopagebreak
\begin{leansrc}
theorem substituteOpen_admissible_iff
    (σ : (substituteOpen outer φ E inner
        hwfO hwfI hpo hti hto).net.Valuation) :
    σ.Admissible ↔
      ((∀ a ∈ inner.apps, a.SemOk σ) ∧
        (∀ a ∈ (substSeam outer φ E inner).repairs, a.SemOk σ)) ∧
      (∀ a ∈ outer.apps.map
          (PrimApp.remapWires (substOuterWire outer φ E inner)),
        a.SemOk σ)
\end{leansrc}
\par\medskip\noindent Modules cited: \texttt{Qtsl.SubstitutionBridge}, \texttt{Qtsl.PackingGap}.
\par\medskip\noindent Also cited in this chapter, listed under an earlier one: \texttt{Qtsl.HypernetCospan}, \texttt{SuFunctor.HypernetDouble}, \texttt{SuFunctor.HypernetInclusion}, \texttt{glue\_admissible\_iff}, \texttt{glue\_assoc}, \texttt{hnetInclusion}, \texttt{theta\_rename}.

\section{Hypernets: Markov Kernels}
Citations from Chapter~\ref{ch:qtsl-markov-kernels}.

\par\medskip\noindent{\small \texttt{SuFunctor.Stochastic.}\texttt{\bfseries sum\_noiseDenotFun}}\hfill{\footnotesize \path{lean/su-functor/SuFunctor/Stochastic/NoiseSemantics.lean}, line 117}\par\nopagebreak
\begin{leansrc}
theorem sum_noiseDenotFun (g : (Fin N → ℝ) → V) (hg : Measurable g) :
    ∑ v, noiseDenotFun g v = 1
\end{leansrc}
\par\medskip\noindent{\small \texttt{SuFunctor.Stochastic.}\texttt{\bfseries noiseDenot\_comp}}\hfill{\footnotesize \path{lean/su-functor/SuFunctor/Stochastic/NoiseSemantics.lean}, line 162}\par\nopagebreak
\begin{leansrc}
theorem noiseDenot_comp (g : (Fin N → ℝ) → V) (π : V → W)
    (hg : Measurable g) (hπg : Measurable (π ∘ g)) :
    noiseDenot (π ∘ g) hπg = ProbState.push π (noiseDenot g hg)
\end{leansrc}
\par\medskip\noindent{\small \texttt{SuFunctor.Stochastic.}\texttt{\bfseries noiseDenot\_seqComp}}\hfill{\footnotesize \path{lean/su-functor/SuFunctor/Stochastic/NoiseKleisli.lean}, line 232}\par\nopagebreak
\begin{leansrc}
theorem noiseDenot_seqComp (g : (Fin N₁ → ℝ) → V) (k : V → (Fin N₂ → ℝ) → W)
    (hg : Measurable g) (hk : ∀ v, Measurable (k v)) :
    noiseDenot (seqComp g k) (measurable_seqComp hg hk)
      = (noiseDenot g hg).bind fun v => noiseDenot (k v) (hk v)
\end{leansrc}
\par\medskip\noindent{\small \texttt{SuFunctor.Stochastic.}\texttt{\bfseries copy\_ne\_kleisli}}\hfill{\footnotesize \path{lean/su-functor/SuFunctor/Stochastic/NoiseKleisli.lean}, line 332}\par\nopagebreak
\begin{leansrc}
theorem copy_ne_kleisli :
    noiseDenot copyPair measurable_copyPair
      ≠ (noiseDenot bitStage measurable_bitStage).bind
          (fun v => noiseDenot (pairStage v) (measurable_pairStage v))
\end{leansrc}
\par\medskip\noindent{\small \texttt{SuFunctor.Stochastic.}\texttt{\bfseries ThinningForm}}\hfill{\footnotesize \path{lean/su-functor/SuFunctor/Stochastic/JumpMeasure.lean}, line 297}\par\nopagebreak
\begin{leansrc}
structure ThinningForm (V : Type u) [Fintype V] : Type u where …
\end{leansrc}
\par\medskip\noindent{\small \texttt{SuFunctor.Stochastic.JumpTerm.}\texttt{\bfseries elaborate\_mass\_partition}}\hfill{\footnotesize \path{lean/su-functor/SuFunctor/Stochastic/JumpMeasure.lean}, line 541}\par\nopagebreak
\begin{leansrc}
theorem elaborate_mass_partition (J : JumpTerm V) (x : V) :
    ((elaborate J).rate : ℝ) *
        ((∑ y, ((elaborate J).kernel x : Step V → ℝ) (Step.next y))
          + ((elaborate J).kernel x : Step V → ℝ) Step.stay
          + ((elaborate J).kernel x : Step V → ℝ) Step.cemetery)
      = (J.intensity : ℝ)
\end{leansrc}
\par\medskip\noindent{\small \texttt{SuFunctor.Stochastic.JumpTerm.}\texttt{\bfseries outflow\_le\_intensity}}\hfill{\footnotesize \path{lean/su-functor/SuFunctor/Stochastic/JumpMeasure.lean}, line 557}\par\nopagebreak
\begin{leansrc}
theorem outflow_le_intensity (J : JumpTerm V) (x : V) :
    ((elaborate J).rate : ℝ) *
        ((∑ y, ((elaborate J).kernel x : Step V → ℝ) (Step.next y))
          + ((elaborate J).kernel x : Step V → ℝ) Step.cemetery)
      ≤ (J.intensity : ℝ)
\end{leansrc}
\par\medskip\noindent{\small \texttt{SuFunctor.Stochastic.JumpTerm.}\texttt{\bfseries generator\_elaborate}}\hfill{\footnotesize \path{lean/su-functor/SuFunctor/Stochastic/JumpMeasure.lean}, line 381}\par\nopagebreak
\begin{leansrc}
theorem generator_elaborate (J : JumpTerm V) :
    (elaborate J).generator = J.generator
\end{leansrc}
\par\medskip\noindent{\small \texttt{SuFunctor.Stochastic.ThinningForm.}\texttt{\bfseries generator\_dilate}}\hfill{\footnotesize \path{lean/su-functor/SuFunctor/Stochastic/JumpMeasure.lean}, line 592}\par\nopagebreak
\begin{leansrc}
theorem generator_dilate (c : ℚ≥0) (hc : 1 ≤ c) (T : ThinningForm V) :
    (T.dilate c hc).generator = T.generator
\end{leansrc}
\par\medskip\noindent{\small \texttt{SuFunctor.Stochastic.ThinningForm.}\texttt{\bfseries generator\_absorbDiagonal}}\hfill{\footnotesize \path{lean/su-functor/SuFunctor/Stochastic/JumpMeasure.lean}, line 681}\par\nopagebreak
\begin{leansrc}
theorem generator_absorbDiagonal (T : ThinningForm V) :
    T.absorbDiagonal.generator = T.generator
\end{leansrc}
\par\medskip\noindent{\small \texttt{SuFunctor.Stochastic.ThinningForm.JumpTerm.}\texttt{\bfseries elaborate\_superpose}}\hfill{\footnotesize \path{lean/su-functor/SuFunctor/Stochastic/JumpAlgebra.lean}, line 534}\par\nopagebreak
\begin{leansrc}
theorem elaborate_superpose (J₁ J₂ : JumpTerm V) :
    (superpose J₁ J₂).elaborate
      = (J₁.elaborate).superpose (J₂.elaborate)
\end{leansrc}
\par\medskip\noindent{\small \texttt{SuFunctor.Stochastic.ThinningForm.JumpTerm.}\texttt{\bfseries generator\_elaborate\_interleave}}\hfill{\footnotesize \path{lean/su-functor/SuFunctor/Stochastic/JumpAlgebra.lean}, line 561}\par\nopagebreak
\begin{leansrc}
theorem generator_elaborate_interleave (J₁ : JumpTerm V₁)
    (J₂ : JumpTerm V₂) :
    (interleave J₁ J₂).elaborate.generator
      = ((J₁.elaborate).interleave (J₂.elaborate)).generator
\end{leansrc}
\par\medskip\noindent{\small \texttt{SuFunctor.Stochastic.}\texttt{\bfseries generator\_eq\_stepMix}}\hfill{\footnotesize \path{lean/su-functor/SuFunctor/Stochastic/JumpDiffusionThinning.lean}, line 187}\par\nopagebreak
\begin{leansrc}
theorem generator_eq_stepMix (P : JdProfile n m) (V : State n → ℝ)
    (gradV : State n → State n) (hessV : State n → Fin n → Fin n → ℝ)
    (z : State n) :
    Qtsl.JumpDiffusion.generator P V gradV hessV z
      = (∑ i, gradV z i * P.f z i)
        + (1 / 2) * (∑ i, ∑ j, P.gram z i j * hessV z i j)
        + P.kill z * stepMixGenerator (killMixture n) V z
        + (P.legs.map fun lm =>
            lm.1 z * stepMixGenerator (deficitToCemetery lm.2) V z).sum
\end{leansrc}
\par\medskip\noindent Modules cited: \texttt{SuFunctor.Stochastic}, \texttt{SuFunctor.Stochastic.JumpAlgebra}.

\section{Systems: Mathematical Models}
Citations from Chapter~\ref{ch:qtsl-mathematical-models}.

\par\medskip\noindent{\small \texttt{SuFunctor.Adequacy.}\texttt{\bfseries isOdeSol\_iff\_integralCurve}}\hfill{\footnotesize \path{lean/su-functor/SuFunctor/Adequacy/Ode.lean}, line 209}\par\nopagebreak
\begin{leansrc}
theorem isOdeSol_iff_integralCurve {f : E → E} {x₀ : E} {y : ℝ → E}
    (hf : Continuous f) :
    IsOdeSol f x₀ y ↔
      y 0 = x₀ ∧ IsIntegralCurveOn y (fun _ => f) (Ici 0)
\end{leansrc}
\par\medskip\noindent{\small \texttt{SuFunctor.Adequacy.}\texttt{\bfseries isOdeSol\_iff\_integral}}\hfill{\footnotesize \path{lean/su-functor/SuFunctor/Adequacy/Ode.lean}, line 180}\par\nopagebreak
\begin{leansrc}
theorem isOdeSol_iff_integral {f : E → E} {x₀ : E} {y : ℝ → E}
    (hf : Continuous f) :
    IsOdeSol f x₀ y ↔ ContinuousOn y (Ici 0) ∧
      ∀ t ∈ Ici (0 : ℝ), y t = x₀ + ∫ s in (0 : ℝ)..t, f (y s)
\end{leansrc}
\par\medskip\noindent{\small \texttt{SuFunctor.Adequacy.}\texttt{\bfseries ode\_adequacy}}\hfill{\footnotesize \path{lean/su-functor/SuFunctor/Adequacy/Ode.lean}, line 319}\par\nopagebreak
\begin{leansrc}
theorem ode_adequacy [CompleteSpace E] {f : E → E} (hf : Continuous f)
    (x₀ : E) {a : Crd (ℝ → E)} :
    ∃ α : Cell (looseConj (odeAtoms f x₀) : Loose a a)
        (suppC (integralCurves f x₀)) (𝟙 a) (𝟙 a),
      Cell.IsInvertible α
\end{leansrc}
\par\medskip\noindent{\small \texttt{SuFunctor.Adequacy.}\texttt{\bfseries ode\_box\_adequacy}}\hfill{\footnotesize \path{lean/su-functor/SuFunctor/Adequacy/Ode.lean}, line 414}\par\nopagebreak
\begin{leansrc}
theorem ode_box_adequacy [CompleteSpace E] {f : E → E} (hf : Continuous f)
    (B : Set E) {a : Crd (ℝ → E)} :
    ∃ α : Cell (looseConj (odeBoxAtoms f B) : Loose a a)
        (suppC (⋃ x₀ ∈ B, integralCurves f x₀)) (𝟙 a) (𝟙 a),
      Cell.IsInvertible α
\end{leansrc}
\par\medskip\noindent{\small \texttt{SuFunctor.Adequacy.}\texttt{\bfseries odeSols\_eqOn\_of\_lipschitzOnWith}}\hfill{\footnotesize \path{lean/su-functor/SuFunctor/Adequacy/Ode.lean}, line 234}\par\nopagebreak
\begin{leansrc}
theorem odeSols_eqOn_of_lipschitzOnWith {f : E → E} {K : ℝ≥0} {s : Set E}
    (hL : LipschitzOnWith K f s) {x₀ : E} {y₁ y₂ : ℝ → E}
    (h₁ : IsOdeSol f x₀ y₁) (h₂ : IsOdeSol f x₀ y₂)
    (hs₁ : ∀ t ∈ Ici (0 : ℝ), y₁ t ∈ s)
    (hs₂ : ∀ t ∈ Ici (0 : ℝ), y₂ t ∈ s) :
    EqOn y₁ y₂ (Ici 0)
\end{leansrc}
\par\medskip\noindent{\small \texttt{SuFunctor.Adequacy.}\texttt{\bfseries su\_ode\_adequacy}}\hfill{\footnotesize \path{lean/su-functor/SuFunctor/Adequacy/Ode.lean}, line 453}\par\nopagebreak
\begin{leansrc}
theorem su_ode_adequacy [CompleteSpace E] {f : E → E} (hf : Continuous f)
    (x₀ : E) (μ : SubPMF (ℝ → E)) :
    idealPoint μ ∈ (odeDnf f x₀).modelClass ↔
      suppC (integralCurves f x₀) μ
\end{leansrc}
\par\medskip\noindent{\small \texttt{SuFunctor.Adequacy.}\texttt{\bfseries heat\_mesh\_glues}}\hfill{\footnotesize \path{lean/su-functor/SuFunctor/Adequacy/HeatTorus.lean}, line 695}\par\nopagebreak
\begin{leansrc}
theorem heat_mesh_glues (hT : IsOpen T)
    (h1 : IsHeatSol D M.I1 T c1) (h2 : IsHeatSol D M.I2 T c2)
    (h3 : IsHeatSol D M.I3 T c3)
    (hs12 : ∀ x ∈ M.I1 ∩ M.I2, ∀ t ∈ T, c1 x t = c2 x t)
    (hs23 : ∀ x ∈ M.I2 ∩ M.I3, ∀ t ∈ T, c2 x t = c3 x t)
    (hw : WrapWall M T c1 c3) :
    ∃ u : ℝ → ℝ → E, IsHeatTorusSol D M.L T u ∧
      (∀ x ∈ M.I1, ∀ t ∈ T, u x t = c1 x t) ∧
      (∀ x ∈ M.I2, ∀ t ∈ T, u x t = c2 x t) ∧
      (∀ x ∈ M.I3, ∀ t ∈ T, u x t = c3 x t)
\end{leansrc}
\par\medskip\noindent{\small \texttt{SuFunctor.Adequacy.}\texttt{\bfseries heat\_adequacy}}\hfill{\footnotesize \path{lean/su-functor/SuFunctor/Adequacy/HeatTorus.lean}, line 1122}\par\nopagebreak
\begin{leansrc}
theorem heat_adequacy (M : RingMesh) (D : ℝ) {T : Set ℝ}
    (hT : IsOpen T) {a : Crd (HeatRec E)} :
    ∃ α : Cell (looseConj (heatAtoms M D T) : Loose a a)
        (suppC (heatClassicalRecs M D T)) (𝟙 a) (𝟙 a),
      Cell.IsInvertible α
\end{leansrc}
\par\medskip\noindent{\small \texttt{SuFunctor.Adequacy.}\texttt{\bfseries chainAtoms\_dirac\_lin}}\hfill{\footnotesize \path{lean/su-functor/SuFunctor/Adequacy/HeatTorus.lean}, line 919}\par\nopagebreak
\begin{leansrc}
theorem chainAtoms_dirac_lin (e : E) :
    looseConj (heatChainAtoms M D T) (SubPMF.dirac (linRec e))
\end{leansrc}
\par\medskip\noindent{\small \texttt{SuFunctor.Adequacy.}\texttt{\bfseries not\_wrap\_lin}}\hfill{\footnotesize \path{lean/su-functor/SuFunctor/Adequacy/HeatTorus.lean}, line 932}\par\nopagebreak
\begin{leansrc}
theorem not_wrap_lin {e : E} (he : e ≠ 0) {t₀ : ℝ} (ht₀ : t₀ ∈ T) :
    ¬ suppC (wrapC M T) (SubPMF.dirac (linRec e))
\end{leansrc}
\par\medskip\noindent{\small \texttt{SuFunctor.Adequacy.}\texttt{\bfseries heat\_chain\_glues}}\hfill{\footnotesize \path{lean/su-functor/SuFunctor/Adequacy/HeatTorus.lean}, line 484}\par\nopagebreak
\begin{leansrc}
theorem heat_chain_glues {M : RingMesh} {D : ℝ} {T : Set ℝ}
    (hT : IsOpen T) {c1 c2 c3 : ℝ → ℝ → E}
    (h1 : IsHeatSol D M.I1 T c1) (h2 : IsHeatSol D M.I2 T c2)
    (h3 : IsHeatSol D M.I3 T c3)
    (hs12 : ∀ x ∈ M.I1 ∩ M.I2, ∀ t ∈ T, c1 x t = c2 x t)
    (hs23 : ∀ x ∈ M.I2 ∩ M.I3, ∀ t ∈ T, c2 x t = c3 x t) :
    ∃ w : ℝ → ℝ → E, IsHeatSol D M.chart T w ∧
      (∀ x ∈ M.I1, ∀ t ∈ T, w x t = c1 x t) ∧
      (∀ x ∈ M.I2, ∀ t ∈ T, w x t = c2 x t) ∧
      (∀ x ∈ M.I3, ∀ t ∈ T, w x t = c3 x t)
\end{leansrc}
\par\medskip\noindent{\small \texttt{SuFunctor.}\texttt{\bfseries vorobev\_no\_joint}}\hfill{\footnotesize \path{lean/su-functor/SuFunctor/Adequacy/MarginalGlue.lean}, line 246}\par\nopagebreak
\begin{leansrc}
theorem vorobev_no_joint (μ : SubPMF (Bool × Bool × Bool)) :
    ¬ (μ.push proj12 = pairCorr ∧ μ.push proj23 = pairCorr ∧
        μ.push proj13 = pairIndep)
\end{leansrc}
\par\medskip\noindent{\small \texttt{SuFunctor.}\texttt{\bfseries credalCompile\_vorobevGlue}}\hfill{\footnotesize \path{lean/su-functor/SuFunctor/Adequacy/MarginalGlue.lean}, line 415}\par\nopagebreak
\begin{leansrc}
theorem credalCompile_vorobevGlue :
    credalCompile vorobevGlue = rayC (SubPMF.zero (Bool × Bool × Bool))
\end{leansrc}
\par\medskip\noindent{\small \texttt{SuFunctor.Adequacy.}\texttt{\bfseries heat\_ring\_glues\_vorobev\_ring\_not}}\hfill{\footnotesize \path{lean/su-functor/SuFunctor/Adequacy/HeatTorus.lean}, line 1038}\par\nopagebreak
\begin{leansrc}
theorem heat_ring_glues_vorobev_ring_not :
    (∃ μ : SubPMF (HeatRec E), looseConj (heatAtoms M D T) μ ∧
        μ ≠ SubPMF.zero (HeatRec E)) ∧
      ¬ ∃ μ : SubPMF (Bool × Bool × Bool), vorobevGlue μ
\end{leansrc}
\par\medskip\noindent{\small \texttt{SuFunctor.Adequacy.}\texttt{\bfseries heat\_init\_adequacy}}\hfill{\footnotesize \path{lean/su-functor/SuFunctor/Adequacy/HeatTorusInit.lean}, line 291}\par\nopagebreak
\begin{leansrc}
theorem heat_init_adequacy (M : RingMesh) (D : ℝ) {T : Set ℝ}
    (hT : IsOpen T) (g : ℝ → E) {a : Crd (HeatRec E)} :
    ∃ α : Cell (looseConj (heatInitAtoms M D T g) : Loose a a)
        (suppC (heatInitClassicalRecs M D T g)) (𝟙 a) (𝟙 a),
      Cell.IsInvertible α
\end{leansrc}
\par\medskip\noindent{\small \texttt{SuFunctor.Adequacy.}\texttt{\bfseries heatInitAtoms\_consistent\_of\_feed}}\hfill{\footnotesize \path{lean/su-functor/SuFunctor/Adequacy/HeatTorusInit.lean}, line 323}\par\nopagebreak
\begin{leansrc}
theorem heatInitAtoms_consistent_of_feed {M : RingMesh} {D : ℝ}
    {T : Set ℝ} (hT : IsOpen T) {g : ℝ → E}
    (hfeed : PortFeed M D T g) :
    ∃ μ : SubPMF (HeatRec E), looseConj (heatInitAtoms M D T g) μ ∧
      μ ≠ SubPMF.zero (HeatRec E)
\end{leansrc}
\par\medskip\noindent{\small \texttt{SuFunctor.Adequacy.}\texttt{\bfseries initFace\_wrap\_compat}}\hfill{\footnotesize \path{lean/su-functor/SuFunctor/Adequacy/HeatTorusInit.lean}, line 178}\par\nopagebreak
\begin{leansrc}
theorem initFace_wrap_compat {M : RingMesh} {D : ℝ} {T : Set ℝ}
    {g : ℝ → E} {p : HeatRec E} (h : IsInitMeshSol M D T g p)
    (hne : (𝓝[T] (0 : ℝ)).NeBot) {x : ℝ} (hx3 : x ∈ M.I3)
    (hx1 : x - M.L ∈ M.I1) : g x = g (x - M.L)
\end{leansrc}
\par\medskip\noindent{\small \texttt{SuFunctor.Adequacy.IsHeatTorusSolInit.}\texttt{\bfseries trace\_wrap\_compat}}\hfill{\footnotesize \path{lean/su-functor/SuFunctor/Adequacy/HeatTorusInit.lean}, line 208}\par\nopagebreak
\begin{leansrc}
theorem IsHeatTorusSolInit.trace_wrap_compat {M : RingMesh} {D : ℝ}
    {T : Set ℝ} {u : ℝ → ℝ → E} {g : ℝ → E}
    (h : IsHeatTorusSolInit M D T u g) (hne : (𝓝[T] (0 : ℝ)).NeBot)
    {x : ℝ} (hx : x ∈ M.chart) (hxL : x + M.L ∈ M.chart) :
    g (x + M.L) = g x
\end{leansrc}
\par\medskip\noindent{\small \texttt{SuFunctor.Adequacy.}\texttt{\bfseries portFeed\_const}}\hfill{\footnotesize \path{lean/su-functor/SuFunctor/Adequacy/HeatTorusInit.lean}, line 335}\par\nopagebreak
\begin{leansrc}
theorem portFeed_const (M : RingMesh) (D : ℝ) (T : Set ℝ) (k : E) :
    PortFeed M D T (fun _ => k)
\end{leansrc}
\par\medskip\noindent{\small \texttt{SuFunctor.Adequacy.}\texttt{\bfseries portFeed\_mode}}\hfill{\footnotesize \path{lean/su-functor/SuFunctor/Adequacy/HeatTorusInit.lean}, line 436}\par\nopagebreak
\begin{leansrc}
theorem portFeed_mode (M : RingMesh) (D : ℝ) (T : Set ℝ) (e : E)
    {ω : ℝ} (k : ℤ) (hω : ω * M.L = k * (2 * Real.pi)) :
    PortFeed M D T (modeDatum ω e)
\end{leansrc}
\par\medskip\noindent{\small \texttt{SuFunctor.Adequacy.}\texttt{\bfseries modeDatum\_ne\_const}}\hfill{\footnotesize \path{lean/su-functor/SuFunctor/Adequacy/HeatTorusInit.lean}, line 447}\par\nopagebreak
\begin{leansrc}
theorem modeDatum_ne_const {e : E} (he : e ≠ 0) {ω : ℝ} (hω : ω ≠ 0) :
    modeDatum ω e (ω⁻¹ * (Real.pi / 2)) ≠ modeDatum ω e 0
\end{leansrc}
\par\medskip\noindent{\small \texttt{SuFunctor.Adequacy.}\texttt{\bfseries initFace\_neBot}}\hfill{\footnotesize \path{lean/su-functor/SuFunctor/Adequacy/HeatTorusInit.lean}, line 99}\par\nopagebreak
\begin{leansrc}
theorem initFace_neBot {τ : ℝ} (hτ : 0 < τ) :
    (𝓝[Ioo (0 : ℝ) τ] (0 : ℝ)).NeBot
\end{leansrc}
\par\medskip\noindent{\small \texttt{SuFunctor.Universe.FellCloseds.}\texttt{\bfseries equalizerHomeomorph}}\hfill{\footnotesize \path{lean/su-functor/SuFunctor/Universe/FellDescentEqualizer.lean}, line 287}\par\nopagebreak
\begin{leansrc}
def equalizerHomeomorph [R1Space α] (U : ι → Set α) (hUo : ∀ i, IsOpen (U i))
    (hcov : ⋃ i, U i = univ) :
    FellCloseds α ≃ₜ Equalizer U where …
\end{leansrc}
\par\medskip\noindent{\small \texttt{SuFunctor.Adequacy.}\texttt{\bfseries meshDescent}}\hfill{\footnotesize \path{lean/su-functor/SuFunctor/Adequacy/SolutionSetMesh.lean}, line 235}\par\nopagebreak
\begin{leansrc}
noncomputable def meshDescent (M : RingMesh) (T : Set ℝ) :
    FellCloseds (MeshPt M T E) ≃ₜ
      FellCloseds.Equalizer (arcRegion M T E)
\end{leansrc}
\par\medskip\noindent{\small \texttt{SuFunctor.Adequacy.}\texttt{\bfseries setStratum\_holonomy}}\hfill{\footnotesize \path{lean/su-functor/SuFunctor/Adequacy/SolutionSetMesh.lean}, line 688}\par\nopagebreak
\begin{leansrc}
theorem setStratum_holonomy (D : ℝ) {t₀ : ℝ} (ht₀ : t₀ ∈ T) :
    ∃ p : FellCloseds.Equalizer (arcRegion M T ℝ),
      (∀ i, IsCellTube M D T i (p.1 i : Set ↥(arcRegion M T ℝ i))) ∧
      TubeWrapWall M T (p.1 0 : Set ↥(arcRegion M T ℝ 0))
        (p.1 2 : Set ↥(arcRegion M T ℝ 2)) ∧
      IsChartTube M D T (meshGlue M T p : Set (MeshPt M T ℝ)) ∧
      ¬ IsTorusTube M D T (meshGlue M T p : Set (MeshPt M T ℝ))
\end{leansrc}
\par\medskip\noindent{\small \texttt{SuFunctor.Adequacy.}\texttt{\bfseries orbitTube\_not\_torusTube}}\hfill{\footnotesize \path{lean/su-functor/SuFunctor/Adequacy/SolutionSetMesh.lean}, line 640}\par\nopagebreak
\begin{leansrc}
theorem orbitTube_not_torusTube (D : ℝ) {t₀ : ℝ} (ht₀ : t₀ ∈ T) :
    ¬ IsTorusTube M D T (orbitTube M T)
\end{leansrc}
\par\medskip\noindent{\small \texttt{SuFunctor.Adequacy.}\texttt{\bfseries not\_isLocalLocusEq\_torusTube}}\hfill{\footnotesize \path{lean/su-functor/SuFunctor/Adequacy/SolutionSetMesh.lean}, line 723}\par\nopagebreak
\begin{leansrc}
theorem not_isLocalLocusEq_torusTube (D : ℝ) {t₀ : ℝ} (ht₀ : t₀ ∈ T) :
    ¬ FellCloseds.IsLocalLocusEq (arcRegion M T ℝ)
        (torusTubeLocus M D T) (cellTubeLocus M D T)
\end{leansrc}
\par\medskip\noindent{\small \texttt{SuFunctor.Adequacy.}\texttt{\bfseries tubeAtoms\_not\_classical}}\hfill{\footnotesize \path{lean/su-functor/SuFunctor/Adequacy/SolutionSetMesh.lean}, line 863}\par\nopagebreak
\begin{leansrc}
theorem tubeAtoms_not_classical {t₀ : ℝ} (ht₀ : t₀ ∈ T) (D : ℝ) :
    ¬ ∀ μ : SubPMF (FellCloseds.Equalizer (arcRegion M T ℝ)),
        looseConj (tubeAtoms M D T) μ ↔ suppC (torusTubeEqs M D T) μ
\end{leansrc}
\par\medskip\noindent{\small \texttt{SuFunctor.Adequacy.}\texttt{\bfseries solset\_forward\_cell}}\hfill{\footnotesize \path{lean/su-functor/SuFunctor/Adequacy/SolutionSetMesh.lean}, line 818}\par\nopagebreak
\begin{leansrc}
theorem solset_forward_cell
    {a : Crd (FellCloseds.Equalizer (arcRegion M T E))} :
    ∃ _ : Cell (suppC (torusTubeEqs M D T) : Loose a a)
        (looseConj (tubeAtoms M D T)) (𝟙 a) (𝟙 a), True
\end{leansrc}
\par\medskip\noindent{\small \texttt{SuFunctor.Adequacy.}\texttt{\bfseries isTorusTube\_univ}}\hfill{\footnotesize \path{lean/su-functor/SuFunctor/Adequacy/SolutionSetMesh.lean}, line 372}\par\nopagebreak
\begin{leansrc}
theorem isTorusTube_univ :
    IsTorusTube M D T (univ : Set (MeshPt M T E))
\end{leansrc}
\par\medskip\noindent{\small \texttt{SuFunctor.Adequacy.}\texttt{\bfseries labelled\_ensembles\_glue}}\hfill{\footnotesize \path{lean/su-functor/SuFunctor/Adequacy/SolutionSetMesh.lean}, line 385}\par\nopagebreak
\begin{leansrc}
theorem labelled_ensembles_glue {ι' : Type*} (hT : IsOpen T)
    {c1 c2 c3 : ι' → ℝ → ℝ → E}
    (h1 : ∀ j, IsHeatSol D M.I1 T (c1 j))
    (h2 : ∀ j, IsHeatSol D M.I2 T (c2 j))
    (h3 : ∀ j, IsHeatSol D M.I3 T (c3 j))
    (hs12 : ∀ j, ∀ x ∈ M.I1 ∩ M.I2, ∀ t ∈ T, c1 j x t = c2 j x t)
    (hs23 : ∀ j, ∀ x ∈ M.I2 ∩ M.I3, ∀ t ∈ T, c2 j x t = c3 j x t)
    (hw : ∀ j, WrapWall M T (c1 j) (c3 j)) :
    ∃ u : ι' → ℝ → ℝ → E, ∀ j,
      IsHeatTorusSol D M.L T (u j) ∧
      (∀ x ∈ M.I1, ∀ t ∈ T, u j x t = c1 j x t) ∧
      (∀ x ∈ M.I2, ∀ t ∈ T, u j x t = c2 j x t) ∧
      (∀ x ∈ M.I3, ∀ t ∈ T, u j x t = c3 j x t)
\end{leansrc}
\par\medskip\noindent{\small \texttt{SuFunctor.Adequacy.}\texttt{\bfseries orbit\_enclosure\_forces\_zero}}\hfill{\footnotesize \path{lean/su-functor/SuFunctor/Adequacy/SolutionSetMesh.lean}, line 991}\par\nopagebreak
\begin{leansrc}
theorem orbit_enclosure_forces_zero (D : ℝ) {t₀ : ℝ} (ht₀ : t₀ ∈ T) :
    ∀ μ : SubPMF (HeatRec ℝ),
      looseConj (heatAtoms M D T (E := ℝ)
          ++ [suppC (enclose1C M T (orbitSet M)),
              suppC (enclose2C M T (orbitSet M)),
              suppC (enclose3C M T (orbitSet M))]) μ →
      μ = SubPMF.zero (HeatRec ℝ)
\end{leansrc}
\par\medskip\noindent{\small \texttt{SuFunctor.Adequacy.LocalRows.}\texttt{\bfseries mesh\_adequacy}}\hfill{\footnotesize \path{lean/su-functor/SuFunctor/Adequacy/Pde.lean}, line 400}\par\nopagebreak
\begin{leansrc}
theorem mesh_adequacy (R : LocalRows E) (M : RingMesh) {T : Set ℝ}
    (hT : IsOpen T) {a : Crd (HeatRec E)} :
    ∃ α : Cell (looseConj (R.meshAtoms M T) : Loose a a)
        (suppC (R.classicalRecs M T)) (𝟙 a) (𝟙 a),
      Cell.IsInvertible α
\end{leansrc}
\par\medskip\noindent{\small \texttt{SuFunctor.Adequacy.LocalRows.PdeAt.}\texttt{\bfseries congr}}\hfill{\footnotesize \path{lean/su-functor/SuFunctor/Adequacy/Pde.lean}, line 644}\par\nopagebreak
\begin{leansrc}
theorem PdeAt.congr {u v : ℝ → ℝ → ι → ℝ} {x t : ℝ}
    {Sx St : Set ℝ} (hSx : IsOpen Sx) (hSt : IsOpen St)
    (hx : x ∈ Sx) (ht : t ∈ St)
    (hagree : ∀ y ∈ Sx, ∀ s ∈ St, u y s = v y s)
    (h : PdeAt P v x t) : PdeAt P u x t
\end{leansrc}
\par\medskip\noindent{\small \texttt{SuFunctor.Adequacy.LocalRows.PdeAt.}\texttt{\bfseries comp\_sub\_const}}\hfill{\footnotesize \path{lean/su-functor/SuFunctor/Adequacy/Pde.lean}, line 683}\par\nopagebreak
\begin{leansrc}
theorem PdeAt.comp_sub_const {w : ℝ → ℝ → ι → ℝ} {x t c : ℝ}
    (h : PdeAt P w (x - c) t) :
    PdeAt P (fun y s => w (y - c) s) x t
\end{leansrc}
\par\medskip\noindent{\small \texttt{SuFunctor.Adequacy.LocalRows.}\texttt{\bfseries pde\_adequacy}}\hfill{\footnotesize \path{lean/su-functor/SuFunctor/Adequacy/Pde.lean}, line 768}\par\nopagebreak
\begin{leansrc}
theorem pde_adequacy (P : PdeMechanism ι k) (M : RingMesh) {T : Set ℝ}
    (hT : IsOpen T) {a : Crd (HeatRec (ι → ℝ))} :
    ∃ α : Cell (looseConj (P.rows.meshAtoms M T) : Loose a a)
        (suppC (P.rows.classicalRecs M T)) (𝟙 a) (𝟙 a),
      Cell.IsInvertible α
\end{leansrc}
\par\medskip\noindent{\small \texttt{SuFunctor.Adequacy.LocalRows.}\texttt{\bfseries heat\_adequacy\_of\_mesh}}\hfill{\footnotesize \path{lean/su-functor/SuFunctor/Adequacy/Pde.lean}, line 510}\par\nopagebreak
\begin{leansrc}
theorem heat_adequacy_of_mesh (M : RingMesh) (D : ℝ) {T : Set ℝ}
    (hT : IsOpen T) {a : Crd (HeatRec E)} :
    ∃ α : Cell (looseConj (heatAtoms M D T) : Loose a a)
        (suppC (heatClassicalRecs M D T)) (𝟙 a) (𝟙 a),
      Cell.IsInvertible α
\end{leansrc}
\par\medskip\noindent{\small \texttt{SuFunctor.Adequacy.LocalRows.}\texttt{\bfseries heatAt\_iff\_pdeAt}}\hfill{\footnotesize \path{lean/su-functor/SuFunctor/Adequacy/Pde.lean}, line 835}\par\nopagebreak
\begin{leansrc}
theorem heatAt_iff_pdeAt {D : ℝ} {u : ℝ → ℝ → ℝ} {x t : ℝ} :
    HeatAt D u x t ↔ PdeAt (heatMechanism D) (fun y s _ => u y s) x t
\end{leansrc}
\par\medskip\noindent{\small \texttt{SuFunctor.Adequacy.LocalRows.}\texttt{\bfseries burgers\_adequacy}}\hfill{\footnotesize \path{lean/su-functor/SuFunctor/Adequacy/Pde.lean}, line 918}\par\nopagebreak
\begin{leansrc}
theorem burgers_adequacy (ν : ℝ) (M : RingMesh) {T : Set ℝ}
    (hT : IsOpen T) {a : Crd (HeatRec (PUnit → ℝ))} :
    ∃ α : Cell
        (looseConj ((burgersMechanism ν).rows.meshAtoms M T) : Loose a a)
        (suppC ((burgersMechanism ν).rows.classicalRecs M T))
        (𝟙 a) (𝟙 a),
      Cell.IsInvertible α
\end{leansrc}
\par\medskip\noindent{\small \texttt{SuFunctor.Adequacy.LocalRows.}\texttt{\bfseries kpp\_adequacy}}\hfill{\footnotesize \path{lean/su-functor/SuFunctor/Adequacy/Pde.lean}, line 928}\par\nopagebreak
\begin{leansrc}
theorem kpp_adequacy (M : RingMesh) {T : Set ℝ}
    (hT : IsOpen T) {a : Crd (HeatRec (PUnit → ℝ))} :
    ∃ α : Cell
        (looseConj (kppMechanism.rows.meshAtoms M T) : Loose a a)
        (suppC (kppMechanism.rows.classicalRecs M T)) (𝟙 a) (𝟙 a),
      Cell.IsInvertible α
\end{leansrc}
\par\medskip\noindent{\small \texttt{SuFunctor.Adequacy.LocalRows.}\texttt{\bfseries coupled\_adequacy}}\hfill{\footnotesize \path{lean/su-functor/SuFunctor/Adequacy/Pde.lean}, line 937}\par\nopagebreak
\begin{leansrc}
theorem coupled_adequacy (M : RingMesh) {T : Set ℝ}
    (hT : IsOpen T) {a : Crd (HeatRec (Fin 2 → ℝ))} :
    ∃ α : Cell
        (looseConj (coupledMechanism.rows.meshAtoms M T) : Loose a a)
        (suppC (coupledMechanism.rows.classicalRecs M T)) (𝟙 a) (𝟙 a),
      Cell.IsInvertible α
\end{leansrc}
\par\medskip\noindent{\small \texttt{SuFunctor.Adequacy.LocalRows.}\texttt{\bfseries PdeRung}}\hfill{\footnotesize \path{lean/su-functor/SuFunctor/Adequacy/Pde.lean}, line 1010}\par\nopagebreak
\begin{leansrc}
structure PdeRung {k : ℕ} (P : PdeMechanism ι k) where …
\end{leansrc}
\par\medskip\noindent{\small \texttt{SuFunctor.Adequacy.LocalRows.}\texttt{\bfseries rung\_adequacy}}\hfill{\footnotesize \path{lean/su-functor/SuFunctor/Adequacy/Pde.lean}, line 1045}\par\nopagebreak
\begin{leansrc}
theorem rung_adequacy {k : ℕ} {P : PdeMechanism ι k} (B : PdeRung P)
    (M : RingMesh) {T : Set ℝ} (hT : IsOpen T)
    {a : Crd (HeatRec (ι → ℝ))} :
    ∃ α : Cell (looseConj (B.rows.meshAtoms M T) : Loose a a)
        (suppC (B.rows.classicalRecs M T)) (𝟙 a) (𝟙 a),
      Cell.IsInvertible α
\end{leansrc}
\par\medskip\noindent{\small \texttt{SuFunctor.Adequacy.LocalRows.}\texttt{\bfseries pde\_rung\_transfer}}\hfill{\footnotesize \path{lean/su-functor/SuFunctor/Adequacy/Pde.lean}, line 1070}\par\nopagebreak
\begin{leansrc}
theorem pde_rung_transfer {k : ℕ} {P : PdeMechanism ι k} (B : PdeRung P)
    (M : RingMesh) (T : Set ℝ) {a : Crd (HeatRec (ι → ℝ))} :
    ∃ _ : Cell (looseConj (P.rows.meshAtoms M T) : Loose a a)
        (looseConj (B.rows.meshAtoms M T)) (𝟙 a) (𝟙 a), True
\end{leansrc}
\par\medskip\noindent{\small \texttt{SuFunctor.Adequacy.LocalRows.}\texttt{\bfseries pdeAtoms\_consistent\_of\_equilibrium}}\hfill{\footnotesize \path{lean/su-functor/SuFunctor/Adequacy/Pde.lean}, line 795}\par\nopagebreak
\begin{leansrc}
theorem pdeAtoms_consistent_of_equilibrium {M : RingMesh} {T : Set ℝ}
    {e : ι → ℝ} (he : P.IsEquilibriumValue e) :
    ∃ μ : SubPMF (HeatRec (ι → ℝ)),
      looseConj (P.rows.meshAtoms M T) μ ∧
        μ ≠ SubPMF.zero (HeatRec (ι → ℝ))
\end{leansrc}
\par\medskip\noindent{\small \texttt{SuFunctor.Adequacy.LocalRows.}\texttt{\bfseries kpp\_not\_equilibrium\_half}}\hfill{\footnotesize \path{lean/su-functor/SuFunctor/Adequacy/Pde.lean}, line 975}\par\nopagebreak
\begin{leansrc}
theorem kpp_not_equilibrium_half :
    ¬ kppMechanism.IsEquilibriumValue (fun _ => (1 / 2 : ℝ))
\end{leansrc}
\par\medskip\noindent{\small \texttt{SuFunctor.Adequacy.LocalRows.}\texttt{\bfseries su\_pde\_adequacy}}\hfill{\footnotesize \path{lean/su-functor/SuFunctor/Adequacy/Pde.lean}, line 1132}\par\nopagebreak
\begin{leansrc}
theorem su_pde_adequacy {k : ℕ} (P : PdeMechanism ι k) (M : RingMesh)
    {T : Set ℝ} (hT : IsOpen T) (μ : SubPMF (HeatRec (ι → ℝ))) :
    idealPoint μ ∈ (P.rows.meshDnf M T).modelClass ↔
      suppC (P.rows.classicalRecs M T) μ
\end{leansrc}
\par\medskip\noindent{\small \texttt{SuFunctor.Adequacy.PdeRowSpec.}\texttt{\bfseries toMechanism}}\hfill{\footnotesize \path{lean/su-functor/SuFunctor/Adequacy/PdeRows.lean}, line 173}\par\nopagebreak
\begin{leansrc}
noncomputable def toMechanism [Inhabited ι] (S : PdeRowSpec ι k) :
    PdeMechanism ι k where …
\end{leansrc}
\par\medskip\noindent{\small \texttt{SuFunctor.Adequacy.}\texttt{\bfseries toMv\_congr}}\hfill{\footnotesize \path{lean/su-functor/SuFunctor/Adequacy/PdeRows.lean}, line 154}\par\nopagebreak
\begin{leansrc}
theorem toMv_congr {σ : Type} {b₁ b₂ : ℕ → σ} :
    ∀ rows : List RowTerm,
      (∀ t ∈ rows, ∀ j, j < t.exps.length → b₁ j = b₂ j) →
      toMv b₁ rows = toMv b₂ rows
  | [], _ => rfl
  | t :: ts, h => by
    have hterm : termProd b₁ 0 t.exps = termProd b₂ 0 t.exps
\end{leansrc}
\par\medskip\noindent{\small \texttt{SuFunctor.Adequacy.PdeRowSpec.}\texttt{\bfseries toMechanism\_congr\_inhabited}}\hfill{\footnotesize \path{lean/su-functor/SuFunctor/Adequacy/PdeRows.lean}, line 181}\par\nopagebreak
\begin{leansrc}
theorem toMechanism_congr_inhabited (S : PdeRowSpec ι k) (hWF : S.WF)
    (d₁ d₂ : Inhabited ι) :
    @toMechanism ι k d₁ S = @toMechanism ι k d₂ S
\end{leansrc}
\par\medskip\noindent{\small \texttt{SuFunctor.Adequacy.}\texttt{\bfseries heatRowSpec\_toMechanism}}\hfill{\footnotesize \path{lean/su-functor/SuFunctor/Adequacy/PdeRows.lean}, line 220}\par\nopagebreak
\begin{leansrc}
theorem heatRowSpec_toMechanism (c : ℚ) :
    (heatRowSpec c).toMechanism = heatMechanism (c : ℝ)
\end{leansrc}
\par\medskip\noindent{\small \texttt{SuFunctor.Adequacy.}\texttt{\bfseries heat\_rows\_adequacy}}\hfill{\footnotesize \path{lean/su-functor/SuFunctor/Adequacy/PdeRows.lean}, line 239}\par\nopagebreak
\begin{leansrc}
theorem heat_rows_adequacy (c : ℚ) (M : RingMesh) {T : Set ℝ}
    (hT : IsOpen T) {a : Crd (HeatRec (PUnit → ℝ))} :
    ∃ α : Cell
        (looseConj (((heatRowSpec c).toMechanism).rows.meshAtoms M T) : Loose a a)
        (suppC (((heatRowSpec c).toMechanism).rows.classicalRecs M T))
        (𝟙 a) (𝟙 a),
      Cell.IsInvertible α
\end{leansrc}
\par\medskip\noindent{\small \texttt{SuFunctor.Adequacy.}\texttt{\bfseries su\_pde\_rows\_adequacy}}\hfill{\footnotesize \path{lean/su-functor/SuFunctor/Adequacy/PdeRows.lean}, line 310}\par\nopagebreak
\begin{leansrc}
theorem su_pde_rows_adequacy {ι : Type} [Fintype ι] [Inhabited ι] {k : ℕ}
    (S : PdeRowSpec ι k) (M : RingMesh) {T : Set ℝ} (hT : IsOpen T)
    (μ : SubPMF (HeatRec (ι → ℝ))) :
    idealPoint μ ∈ ((S.toMechanism).rows.meshDnf M T).modelClass ↔
      suppC ((S.toMechanism).rows.classicalRecs M T) μ
\end{leansrc}
\par\medskip\noindent{\small \texttt{SuFunctor.Adequacy.}\texttt{\bfseries composite\_iff\_dde}}\hfill{\footnotesize \path{lean/su-functor/SuFunctor/Adequacy/DelayLine.lean}, line 442}\par\nopagebreak
\begin{leansrc}
theorem composite_iff_dde {f : E → E → E} {c τ : ℝ} {φ : ℝ → E}
    (hc : 0 < c) (hτ : 0 < τ) (hφ : ContinuousOn φ (Icc (-τ) 0))
    {y z : ℝ → E} {u : ℝ → ℝ → E} :
    IsCompositeSol f c τ φ y z u ↔
      ∃ Y, IsDdeSol f τ φ Y ∧ RecMatches c τ Y y z u
\end{leansrc}
\par\medskip\noindent{\small \texttt{SuFunctor.Adequacy.}\texttt{\bfseries delayline\_adequacy}}\hfill{\footnotesize \path{lean/su-functor/SuFunctor/Adequacy/DelayLine.lean}, line 592}\par\nopagebreak
\begin{leansrc}
theorem delayline_adequacy (f : E → E → E) {c τ : ℝ} (hc : 0 < c)
    (hτ : 0 < τ) {φ : ℝ → E} (hφ : ContinuousOn φ (Icc (-τ) 0))
    {a : Crd (DelayRec E)} :
    ∃ α : Cell (looseConj (delayAtoms f c τ φ) : Loose a a)
        (suppC (classicalRecs f c τ φ)) (𝟙 a) (𝟙 a),
      Cell.IsInvertible α
\end{leansrc}
\par\medskip\noindent{\small \texttt{SuFunctor.Adequacy.IsCompositeSol.}\texttt{\bfseries outflow\_delay}}\hfill{\footnotesize \path{lean/su-functor/SuFunctor/Adequacy/DelayLine.lean}, line 295}\par\nopagebreak
\begin{leansrc}
theorem outflow_delay (hc : 0 < c) (hτ : 0 < τ)
    (h : IsCompositeSol f c τ φ y z u) :
    ∀ t ∈ Ici (0 : ℝ), z t = concat φ y (t - τ)
\end{leansrc}
\par\medskip\noindent{\small \texttt{SuFunctor.Adequacy.IsCompositeSol.}\texttt{\bfseries field\_slice\_is\_history}}\hfill{\footnotesize \path{lean/su-functor/SuFunctor/Adequacy/DelayLine.lean}, line 310}\par\nopagebreak
\begin{leansrc}
theorem field_slice_is_history (hc : 0 < c)
    (h : IsCompositeSol f c τ φ y z u) {t : ℝ} (ht : t ∈ Ici (0 : ℝ)) :
    ∀ s ∈ Icc (t - τ) t, u (c * (t - s)) t = concat φ y s
\end{leansrc}
\par\medskip\noindent{\small \texttt{Qtsl.JumpDiffusion.}\texttt{\bfseries MpSolution}}\hfill{\footnotesize \path{lean/qtsl-aux/QtslAux/JumpDiffusionMartingaleProblem.lean}, line 191}\par\nopagebreak
\begin{leansrc}
structure MpSolution (P : JdProfile n m) (D : Set (State n))
    (Ω : Type*) [MeasureSpace Ω] where …
\end{leansrc}
\par\medskip\noindent{\small \texttt{Qtsl.JumpDiffusion.MpSolution.}\texttt{\bfseries supermartingale\_of\_gen\_nonpos}}\hfill{\footnotesize \path{lean/qtsl-aux/QtslAux/JumpDiffusionMartingaleProblem.lean}, line 234}\par\nopagebreak
\begin{leansrc}
theorem supermartingale_of_gen_nonpos
    [IsProbabilityMeasure (volume : Measure Ω)]
    (mp : MpSolution P D Ω) (hC2 : IsC2Fn V gradV hessV)
    (hgint : ∀ (ω : Ω) (t : ℝ), IntervalIntegrable
      (fun s => generator P V gradV hessV (mp.X (Real.toNNReal s) ω))
      volume 0 t)
    (hVint : ∀ t : ℝ≥0, Integrable (fun ω => V (mp.X t ω))
      (volume : Measure Ω))
    (hrow : ∀ z ∈ D, generator P V gradV hessV z ≤ 0) :
    Supermartingale (fun (t : ℝ≥0) (ω : Ω) => V (mp.X t ω)) mp.ℱ
      (volume : Measure Ω)
\end{leansrc}
\par\medskip\noindent{\small \texttt{Qtsl.SdeMartingaleProblem.MpSolution.}\texttt{\bfseries toSdeDynkinVille}}\hfill{\footnotesize \path{lean/qtsl-aux/QtslAux/SdeMartingaleProblem.lean}, line 409}\par\nopagebreak
\begin{leansrc}
noncomputable def toSdeDynkinVille
    [IsProbabilityMeasure (volume : Measure Ω)]
    (mp : MpSolution P D Ω) (τ : ℕ → ℝ≥0) (hτ : Monotone τ)
    (hC2 : IsC2Fn V gradV hessV)
    (hgint : ∀ (ω : Ω) (t : ℝ), IntervalIntegrable
      (fun s => generator P V gradV hessV (mp.X (Real.toNNReal s) ω))
      volume 0 t)
    (hVint : ∀ t : ℝ≥0, Integrable (fun ω => V (mp.X t ω))
      (volume : Measure Ω)) :
    SdeDynkinVille P D V gradV hessV Ω where …
\end{leansrc}
\par\medskip\noindent{\small \texttt{Qtsl.JumpDiffusion.}\texttt{\bfseries genPoly\_eval\_eq\_generator}}\hfill{\footnotesize \path{lean/qtsl-aux/QtslAux/JumpDiffusion.lean}, line 269}\par\nopagebreak
\begin{leansrc}
theorem genPoly_eval_eq_generator (P : PolyProfile n m)
    (V : MvPolynomial (Fin n) ℚ) (z : State n) :
    (aeval z (genPoly P V) : ℝ)
      = generator P.toJd
          (fun w => (aeval w V : ℝ))
          (fun w i => (aeval w (pderiv i V) : ℝ))
          (fun w i j => (aeval w (pderiv i (pderiv j V)) : ℝ)) z
\end{leansrc}
\par\medskip\noindent{\small \texttt{Qtsl.SdeMartingaleProblem.}\texttt{\bfseries isC2Fn\_poly}}\hfill{\footnotesize \path{lean/qtsl-aux/QtslAux/SdeMartingaleProblem.lean}, line 285}\par\nopagebreak
\begin{leansrc}
theorem isC2Fn_poly (V : MvPolynomial (Fin n) ℚ) :
    IsC2Fn (fun z : State n => (aeval z V : ℝ))
      (fun z i => (aeval z (pderiv i V) : ℝ))
      (fun z i j => (aeval z (pderiv i (pderiv j V)) : ℝ)) where
  hasFDerivAt_V z
\end{leansrc}
\par\medskip\noindent{\small \texttt{Qtsl.JumpDiffusion.}\texttt{\bfseries mp\_avoid\_row\_of\_genPoly}}\hfill{\footnotesize \path{lean/qtsl-aux/QtslAux/JumpDiffusionMartingaleProblem.lean}, line 326}\par\nopagebreak
\begin{leansrc}
theorem mp_avoid_row_of_genPoly
    [IsProbabilityMeasure (volume : Measure Ω)]
    {P₀ : PolyProfile n m} {D : Set (State n)}
    (V₀ : MvPolynomial (Fin n) ℚ)
    (mp : MpSolution P₀.toJd D Ω)
    (τ : ℕ → ℝ≥0) (hτ : Monotone τ)
    (hgint : ∀ (ω : Ω) (t : ℝ), IntervalIntegrable
      (fun s => (aeval (mp.X (Real.toNNReal s) ω) (genPoly P₀ V₀) : ℝ))
      volume 0 t)
    (hVint : ∀ t : ℝ≥0, Integrable
      (fun ω => (aeval (mp.X t ω) V₀ : ℝ)) (volume : Measure Ω))
    (hrow : ∀ z ∈ D, (aeval z (genPoly P₀ V₀) : ℝ) ≤ 0)
    (hnonneg : ∀ z ∈ D, 0 ≤ (aeval z V₀ : ℝ))
    {Bad : Set (State n)} {a : ℝ}
    (hfloor : ∀ z ∈ Bad, a ≤ (aeval z V₀ : ℝ)) :
    ENNReal.ofReal a * volume {ω | ∃ k, mp.X (τ k) ω ∈ Bad}
      ≤ ENNReal.ofReal (∫ ω, (aeval (mp.X (τ 0) ω) V₀ : ℝ))
\end{leansrc}
\par\medskip\noindent{\small \texttt{Qtsl.JumpDiffusion.}\texttt{\bfseries mp\_avoid\_init\_of\_genPoly}}\hfill{\footnotesize \path{lean/qtsl-aux/QtslAux/JumpDiffusionMartingaleProblem.lean}, line 370}\par\nopagebreak
\begin{leansrc}
theorem mp_avoid_init_of_genPoly
    [IsProbabilityMeasure (volume : Measure Ω)]
    {P₀ : PolyProfile n m} {D : Set (State n)}
    (V₀ : MvPolynomial (Fin n) ℚ)
    (mp : MpSolution P₀.toJd D Ω)
    (τ : ℕ → ℝ≥0) (hτ : Monotone τ)
    (hgint : ∀ (ω : Ω) (t : ℝ), IntervalIntegrable
      (fun s => (aeval (mp.X (Real.toNNReal s) ω) (genPoly P₀ V₀) : ℝ))
      volume 0 t)
    (hVint : ∀ t : ℝ≥0, Integrable
      (fun ω => (aeval (mp.X t ω) V₀ : ℝ)) (volume : Measure Ω))
    (hrow : ∀ z ∈ D, (aeval z (genPoly P₀ V₀) : ℝ) ≤ 0)
    (hnonneg : ∀ z ∈ D, 0 ≤ (aeval z V₀ : ℝ))
    {Bad : Set (State n)} {z0 : State n}
    (hstart : ∀ ω : Ω, mp.X (τ 0) ω = z0)
    (hfloor : ∀ z ∈ Bad, 1 ≤ (aeval z V₀ : ℝ)) :
    volume {ω | ∃ k, mp.X (τ k) ω ∈ Bad}
      ≤ ENNReal.ofReal (aeval z0 V₀ : ℝ)
\end{leansrc}
\par\medskip\noindent{\small \texttt{Qtsl.JumpDiffusion.}\texttt{\bfseries mp\_control}}\hfill{\footnotesize \path{lean/qtsl-aux/QtslAux/JumpDiffusionMartingaleProblem.lean}, line 543}\par\nopagebreak
\begin{leansrc}
theorem mp_control [IsProbabilityMeasure (volume : Measure Ω)]
    (x₀ : State 1) {Bad : Set (State 1)}
    (hfloor : ∀ z ∈ Bad,
      1 ≤ (aeval z (X 0 * X 0 : MvPolynomial (Fin 1) ℚ) : ℝ)) :
    volume {_ω : Ω | ∃ _k : ℕ, x₀ ∈ Bad}
      ≤ ENNReal.ofReal (aeval x₀ (X 0 * X 0 : MvPolynomial (Fin 1) ℚ) : ℝ)
\end{leansrc}
\par\medskip\noindent{\small \texttt{SuFunctor.Adequacy.SpectrumDecl.}\texttt{\bfseries plainTrace\_divergent\_none}}\hfill{\footnotesize \path{lean/su-functor/SuFunctor/Adequacy/PdeRows.lean}, line 549}\par\nopagebreak
\begin{leansrc}
theorem plainTrace_divergent_none (ds : List DriverDecl) :
    (SpectrumDecl.mk ds .divergent).plainTrace = none
\end{leansrc}
\par\medskip\noindent{\small \texttt{SuFunctor.Adequacy.}\texttt{\bfseries isRodeSol\_iff\_integralCurve}}\hfill{\footnotesize \path{lean/su-functor/SuFunctor/Adequacy/Rode.lean}, line 63}\par\nopagebreak
\begin{leansrc}
theorem isRodeSol_iff_integralCurve [CompleteSpace E] {f : Θ → E → E}
    {x₀ : E} {θ : Θ} {y : ℝ → E} (hf : Continuous (f θ)) :
    IsRodeSol f x₀ θ y ↔
      y 0 = x₀ ∧ IsIntegralCurveOn y (fun _ => f θ) (Ici 0)
\end{leansrc}
\par\medskip\noindent{\small \texttt{SuFunctor.Adequacy.}\texttt{\bfseries rode\_draw\_adequacy}}\hfill{\footnotesize \path{lean/su-functor/SuFunctor/Adequacy/Rode.lean}, line 71}\par\nopagebreak
\begin{leansrc}
theorem rode_draw_adequacy [CompleteSpace E] {f : Θ → E → E} {θ : Θ}
    (hf : Continuous (f θ)) (x₀ : E) {a : Crd (ℝ → E)} :
    ∃ α : Cell (looseConj (odeAtoms (f θ) x₀) : Loose a a)
        (suppC (integralCurves (f θ) x₀)) (𝟙 a) (𝟙 a),
      Cell.IsInvertible α
\end{leansrc}
\par\medskip\noindent{\small \texttt{SuFunctor.Adequacy.}\texttt{\bfseries rode\_adequacy}}\hfill{\footnotesize \path{lean/su-functor/SuFunctor/Adequacy/Rode.lean}, line 160}\par\nopagebreak
\begin{leansrc}
theorem rode_adequacy [CompleteSpace E] {f : Θ → E → E} {P : Set Θ}
    (hf : ∀ θ ∈ P, Continuous (f θ)) (x₀ : E) {a : Crd (ℝ → E)} :
    ∃ α : Cell (looseConj (rodeAtoms f P x₀) : Loose a a)
        (suppC (⋃ θ ∈ P, integralCurves (f θ) x₀)) (𝟙 a) (𝟙 a),
      Cell.IsInvertible α
\end{leansrc}
\par\medskip\noindent{\small \texttt{SuFunctor.Adequacy.}\texttt{\bfseries rode\_box\_adequacy}}\hfill{\footnotesize \path{lean/su-functor/SuFunctor/Adequacy/Rode.lean}, line 241}\par\nopagebreak
\begin{leansrc}
theorem rode_box_adequacy [CompleteSpace E] {f : Θ → E → E} {P : Set Θ}
    (hf : ∀ θ ∈ P, Continuous (f θ)) (B : Set E) {a : Crd (ℝ → E)} :
    ∃ α : Cell (looseConj (rodeBoxAtoms f P B) : Loose a a)
        (suppC (⋃ θ ∈ P, ⋃ x₀ ∈ B, integralCurves (f θ) x₀))
        (𝟙 a) (𝟙 a),
      Cell.IsInvertible α
\end{leansrc}
\par\medskip\noindent{\small \texttt{SuFunctor.Adequacy.}\texttt{\bfseries rodeSols\_eqOn\_of\_lipschitzOnWith}}\hfill{\footnotesize \path{lean/su-functor/SuFunctor/Adequacy/Rode.lean}, line 81}\par\nopagebreak
\begin{leansrc}
theorem rodeSols_eqOn_of_lipschitzOnWith {f : Θ → E → E} {K : ℝ≥0}
    {s : Set E} {P : Set Θ} (hL : ∀ θ ∈ P, LipschitzOnWith K (f θ) s)
    {x₀ : E} {θ : Θ} (hθ : θ ∈ P) {y₁ y₂ : ℝ → E}
    (h₁ : IsRodeSol f x₀ θ y₁) (h₂ : IsRodeSol f x₀ θ y₂)
    (hs₁ : ∀ t ∈ Ici (0 : ℝ), y₁ t ∈ s)
    (hs₂ : ∀ t ∈ Ici (0 : ℝ), y₂ t ∈ s) :
    EqOn y₁ y₂ (Ici 0)
\end{leansrc}
\par\medskip\noindent{\small \texttt{SuFunctor.Adequacy.}\texttt{\bfseries su\_rode\_adequacy}}\hfill{\footnotesize \path{lean/su-functor/SuFunctor/Adequacy/Rode.lean}, line 281}\par\nopagebreak
\begin{leansrc}
theorem su_rode_adequacy [CompleteSpace E] {f : Θ → E → E} {P : Set Θ}
    (hf : ∀ θ ∈ P, Continuous (f θ)) (x₀ : E) (μ : SubPMF (ℝ → E)) :
    idealPoint μ ∈ (rodeDnf f P x₀).modelClass ↔
      suppC (⋃ θ ∈ P, integralCurves (f θ) x₀) μ
\end{leansrc}
\par\medskip\noindent{\small \texttt{Qtsl.JumpDiffusion.}\texttt{\bfseries JdProfile}}\hfill{\footnotesize \path{lean/qtsl-aux/QtslAux/JumpDiffusion.lean}, line 121}\par\nopagebreak
\begin{leansrc}
structure JdProfile (n m : ℕ) where …
\end{leansrc}
\par\medskip\noindent{\small \texttt{Qtsl.JumpDiffusion.}\texttt{\bfseries generator\_one\_dirac}}\hfill{\footnotesize \path{lean/qtsl-aux/QtslAux/JumpDiffusion.lean}, line 178}\par\nopagebreak
\begin{leansrc}
theorem generator_one_dirac (P : JdProfile n m) (lam : State n → ℝ)
    (g : State n → State n) (hkill : ∀ z, P.kill z = 0)
    (hlegs : P.legs = [(lam, [((1 : ℝ), g)])])
    (V : State n → ℝ) (gradV : State n → State n)
    (hessV : State n → Fin n → Fin n → ℝ) (z : State n) :
    generator P V gradV hessV z
      = SdeGenerator.generator ⟨P.f, P.σ, lam, g⟩ V gradV hessV z
\end{leansrc}
\par\medskip\noindent{\small \texttt{Qtsl.JumpDiffusion.}\texttt{\bfseries toJd\_generator\_one\_dirac}}\hfill{\footnotesize \path{lean/qtsl-aux/QtslAux/JumpDiffusion.lean}, line 293}\par\nopagebreak
\begin{leansrc}
theorem toJd_generator_one_dirac (P : PolyProfile n m)
    (lam : MvPolynomial (Fin n) ℚ) (g : Fin n → MvPolynomial (Fin n) ℚ)
    (hkill : P.kill = 0) (hlegs : P.legs = [(lam, [((1 : ℚ), g)])])
    (V : State n → ℝ) (gradV : State n → State n)
    (hessV : State n → Fin n → Fin n → ℝ) (z : State n) :
    generator P.toJd V gradV hessV z
      = SdeGenerator.generator (P.toSde lam g) V gradV hessV z
\end{leansrc}
\par\medskip\noindent{\small \texttt{Qtsl.JumpDiffusion.}\texttt{\bfseries transferSum\_sub\_self}}\hfill{\footnotesize \path{lean/qtsl-aux/QtslAux/JumpDiffusion.lean}, line 101}\par\nopagebreak
\begin{leansrc}
theorem transferSum_sub_self (μ : MixtureR n) (V : State n → ℝ)
    (z : State n) :
    transferSum μ V z - V z
      = (μ.map fun wg => wg.1 * (V (wg.2 z) - V z)).sum
        - (1 - (μ.map fun wg => wg.1).sum) * V z
\end{leansrc}
\par\medskip\noindent{\small \texttt{Qtsl.JumpDiffusion.}\texttt{\bfseries two\_leg\_control}}\hfill{\footnotesize \path{lean/qtsl-aux/QtslAux/JumpDiffusion.lean}, line 329}\par\nopagebreak
\begin{leansrc}
theorem two_leg_control (z : State 1) :
    4 * generator twoLeg.toJd (fun w => (w 0) ^ 2)
        (fun w _ => 2 * w 0) (fun _ _ _ => 2) z
      = 4 + (z 0) ^ 2 - 3 * (z 0) ^ 4
\end{leansrc}
\par\medskip\noindent{\small \texttt{Qtsl.JumpDiffusion.}\texttt{\bfseries JdDynkinVille}}\hfill{\footnotesize \path{lean/qtsl-aux/QtslAux/JumpDiffusionMartingaleProblem.lean}, line 82}\par\nopagebreak
\begin{leansrc}
structure JdDynkinVille (P : JdProfile n m) (D : Set (State n))
    (V : State n → ℝ) (gradV : State n → State n)
    (hessV : State n → Fin n → Fin n → ℝ)
    (Ω : Type*) [MeasureSpace Ω] where …
\end{leansrc}
\par\medskip\noindent{\small \texttt{Qtsl.JumpDiffusion.}\texttt{\bfseries generator\_toJd}}\hfill{\footnotesize \path{lean/qtsl-aux/QtslAux/JumpDiffusionMartingaleProblem.lean}, line 431}\par\nopagebreak
\begin{leansrc}
theorem generator_toJd (P : SdeProfile n m) (V : State n → ℝ)
    (gradV : State n → State n) (hessV : State n → Fin n → Fin n → ℝ)
    (z : State n) :
    generator P.toJd V gradV hessV z
      = SdeGenerator.generator P V gradV hessV z
\end{leansrc}
\par\medskip\noindent{\small \texttt{Qtsl.JumpDiffusion.\_root\_.Qtsl.SdeMartingaleProblem.MpSolution.}\texttt{\bfseries toJd}}\hfill{\footnotesize \path{lean/qtsl-aux/QtslAux/JumpDiffusionMartingaleProblem.lean}, line 445}\par\nopagebreak
\begin{leansrc}
noncomputable def _root_.Qtsl.SdeMartingaleProblem.MpSolution.toJd
    {P : SdeProfile n m} {D : Set (State n)} {Ω : Type*} [MeasureSpace Ω]
    (mp : SdeMartingaleProblem.MpSolution P D Ω) :
    MpSolution P.toJd D Ω where …
\end{leansrc}
\par\medskip\noindent{\small \texttt{SuFunctor.Adequacy.}\texttt{\bfseries satAtoms\_iff}}\hfill{\footnotesize \path{lean/su-functor/SuFunctor/Adequacy/Dtmc.lean}, line 598}\par\nopagebreak
\begin{leansrc}
theorem satAtoms_iff (hR : R.WF) (μ : SubPMF (Traj S k)) :
    looseConj (atoms R k) μ ↔ μ = rowsLaw R hR k
\end{leansrc}
\par\medskip\noindent{\small \texttt{SuFunctor.Adequacy.}\texttt{\bfseries dtmc\_adequacy}}\hfill{\footnotesize \path{lean/su-functor/SuFunctor/Adequacy/Dtmc.lean}, line 622}\par\nopagebreak
\begin{leansrc}
theorem dtmc_adequacy (M : FinDTMC S) (k : ℕ) {a : Crd (Traj S k)} :
    ∃ α : Cell (looseConj (atoms (compile M) k) : Loose a a)
        (singletonC (M.pathLaw k)) (𝟙 a) (𝟙 a),
      Cell.IsInvertible α
\end{leansrc}
\par\medskip\noindent{\small \texttt{SuFunctor.Adequacy.FinDTMC.}\texttt{\bfseries pathLaw\_mass}}\hfill{\footnotesize \path{lean/su-functor/SuFunctor/Adequacy/Dtmc.lean}, line 259}\par\nopagebreak
\begin{leansrc}
theorem FinDTMC.pathLaw_mass (M : FinDTMC S) (k : ℕ) :
    (M.pathLaw k).mass = 1
\end{leansrc}
\par\medskip\noindent{\small \texttt{SuFunctor.Adequacy.FinDTMC.}\texttt{\bfseries trajPMF\_eq}}\hfill{\footnotesize \path{lean/su-functor/SuFunctor/Adequacy/Dtmc.lean}, line 743}\par\nopagebreak
\begin{leansrc}
theorem FinDTMC.trajPMF_eq :
    ∀ (k : ℕ) (ω : Traj S k),
      (M.trajPMF k) ω = q2e ((M.pathLaw k).wt ω)
\end{leansrc}
\par\medskip\noindent{\small \texttt{SuFunctor.Adequacy.}\texttt{\bfseries ctmc\_jump\_adequacy}}\hfill{\footnotesize \path{lean/su-functor/SuFunctor/Adequacy/CtmcJumpChain.lean}, line 371}\par\nopagebreak
\begin{leansrc}
theorem ctmc_jump_adequacy (M : FinCTMC S) (k : ℕ) {a : Crd (Traj S k)} :
    ∃ α : Cell
        (looseConj (atoms (jumpRows (compileCtmc M)) k) : Loose a a)
        (singletonC (M.jumpLaw k)) (𝟙 a) (𝟙 a),
      Cell.IsInvertible α
\end{leansrc}
\par\medskip\noindent{\small \texttt{SuFunctor.Adequacy.}\texttt{\bfseries compile\_jump\_comm}}\hfill{\footnotesize \path{lean/su-functor/SuFunctor/Adequacy/CtmcJumpChain.lean}, line 352}\par\nopagebreak
\begin{leansrc}
theorem compile_jump_comm (M : FinCTMC S) :
    jumpRows (compileCtmc M) = compile M.jumpChain
\end{leansrc}
\par\medskip\noindent{\small \texttt{SuFunctor.Adequacy.}\texttt{\bfseries jumpRows\_wf}}\hfill{\footnotesize \path{lean/su-functor/SuFunctor/Adequacy/CtmcJumpChain.lean}, line 338}\par\nopagebreak
\begin{leansrc}
theorem jumpRows_wf {R : CtmcRows S} (hR : R.WF) : (jumpRows R).WF where
  outcome_nonneg p
\end{leansrc}
\par\medskip\noindent{\small \texttt{SuFunctor.Adequacy.FinCTMC.}\texttt{\bfseries jumpChain\_step\_mul\_exit}}\hfill{\footnotesize \path{lean/su-functor/SuFunctor/Adequacy/CtmcJumpChain.lean}, line 210}\par\nopagebreak
\begin{leansrc}
theorem FinCTMC.jumpChain_step_mul_exit (M : FinCTMC S) (s s' : S) :
    M.jumpChain.step s s' * M.exitRate s = M.offRate s s'
\end{leansrc}
\par\medskip\noindent{\small \texttt{SuFunctor.Adequacy.}\texttt{\bfseries compileCtmc\_diag\_blind}}\hfill{\footnotesize \path{lean/su-functor/SuFunctor/Adequacy/CtmcJumpChain.lean}, line 385}\par\nopagebreak
\begin{leansrc}
theorem compileCtmc_diag_blind (M M' : FinCTMC S)
    (hoff : ∀ s s', s ≠ s' → M.rate s s' = M'.rate s s')
    (hinit : M.init = M'.init) : compileCtmc M = compileCtmc M'
\end{leansrc}
\par\medskip\noindent{\small \texttt{SuFunctor.Adequacy.FinCTMC.}\texttt{\bfseries jumpLaw\_diag\_blind}}\hfill{\footnotesize \path{lean/su-functor/SuFunctor/Adequacy/CtmcJumpChain.lean}, line 407}\par\nopagebreak
\begin{leansrc}
theorem FinCTMC.jumpLaw_diag_blind (M M' : FinCTMC S)
    (hoff : ∀ s s', s ≠ s' → M.rate s s' = M'.rate s s')
    (hinit : M.init = M'.init) (k : ℕ) :
    M.jumpLaw k = M'.jumpLaw k
\end{leansrc}
\par\medskip\noindent{\small \texttt{SuFunctor.Adequacy.FinCTMC.}\texttt{\bfseries jumpTrajPMF\_eq}}\hfill{\footnotesize \path{lean/su-functor/SuFunctor/Adequacy/CtmcJumpChain.lean}, line 440}\par\nopagebreak
\begin{leansrc}
theorem FinCTMC.jumpTrajPMF_eq (M : FinCTMC S) (k : ℕ) (ω : Traj S k) :
    M.jumpTrajPMF k ω = q2e ((M.jumpLaw k).wt ω)
\end{leansrc}
\par\medskip\noindent{\small \texttt{SuFunctor.Adequacy.}\texttt{\bfseries mdp\_adequacy}}\hfill{\footnotesize \path{lean/su-functor/SuFunctor/Adequacy/Mdp.lean}, line 284}\par\nopagebreak
\begin{leansrc}
theorem mdp_adequacy (M : FinMDP S A) (π : Policy S A) (k : ℕ)
    {a : Crd (Traj S k)} :
    ∃ α : Cell
        (looseConj (atoms (closeRows (compileMdp M) (compilePolicy π)) k) :
          Loose a a)
        (singletonC (M.pathLaw π k)) (𝟙 a) (𝟙 a),
      Cell.IsInvertible α
\end{leansrc}
\par\medskip\noindent{\small \texttt{SuFunctor.Adequacy.}\texttt{\bfseries compile\_close\_comm}}\hfill{\footnotesize \path{lean/su-functor/SuFunctor/Adequacy/Mdp.lean}, line 266}\par\nopagebreak
\begin{leansrc}
theorem compile_close_comm (M : FinMDP S A) (π : Policy S A) :
    closeRows (compileMdp M) (compilePolicy π) =
      compile (M.closedLoop π)
\end{leansrc}
\par\medskip\noindent{\small \texttt{SuFunctor.Adequacy.FinMDP.}\texttt{\bfseries closedLoop\_stepPMF}}\hfill{\footnotesize \path{lean/su-functor/SuFunctor/Adequacy/Mdp.lean}, line 334}\par\nopagebreak
\begin{leansrc}
theorem FinMDP.closedLoop_stepPMF (M : FinMDP S A) (π : Policy S A)
    (s : S) :
    (M.closedLoop π).stepPMF s =
      (π.actPMF s).bind (fun a => M.stepPMF s a)
\end{leansrc}
\par\medskip\noindent{\small \texttt{SuFunctor.}\texttt{\bfseries unionSingletons}}\hfill{\footnotesize \path{lean/su-functor/SuFunctor/Adequacy/EnvelopeUnion.lean}, line 104}\par\nopagebreak
\begin{leansrc}
def unionSingletons (μs : ι → SubPMF V) : SubPMF V → Prop
\end{leansrc}
\par\medskip\noindent{\small \texttt{SuFunctor.Adequacy.FinMDP.}\texttt{\bfseries demonicCredal}}\hfill{\footnotesize \path{lean/su-functor/SuFunctor/Adequacy/MdpDemonic.lean}, line 81}\par\nopagebreak
\begin{leansrc}
def FinMDP.demonicCredal (M : FinMDP S A) (k : ℕ) :
    SubPMF (Traj S k) → Prop
\end{leansrc}
\par\medskip\noindent{\small \texttt{SuFunctor.Adequacy.}\texttt{\bfseries mdp\_demonic\_adequacy}}\hfill{\footnotesize \path{lean/su-functor/SuFunctor/Adequacy/MdpDemonic.lean}, line 157}\par\nopagebreak
\begin{leansrc}
theorem mdp_demonic_adequacy (M : FinMDP S A) (k : ℕ)
    {a : Crd (Traj S k)} :
    ∃ α : Cell (demonicC (compileMdp M) k : Loose a a)
        (M.demonicLaws k) (𝟙 a) (𝟙 a),
      Cell.IsInvertible α
\end{leansrc}
\par\medskip\noindent{\small \texttt{SuFunctor.Adequacy.}\texttt{\bfseries demonicC}}\hfill{\footnotesize \path{lean/su-functor/SuFunctor/Adequacy/MdpDemonic.lean}, line 92}\par\nopagebreak
\begin{leansrc}
def demonicC (R : MdpRows S A) (k : ℕ) : SubPMF (Traj S k) → Prop
\end{leansrc}
\par\medskip\noindent{\small \texttt{SuFunctor.Adequacy.}\texttt{\bfseries demonicC\_iff}}\hfill{\footnotesize \path{lean/su-functor/SuFunctor/Adequacy/MdpDemonic.lean}, line 128}\par\nopagebreak
\begin{leansrc}
theorem demonicC_iff (M : FinMDP S A) (k : ℕ) (μ : SubPMF (Traj S k)) :
    demonicC (compileMdp M) k μ ↔ M.demonicLaws k μ
\end{leansrc}
\par\medskip\noindent{\small \texttt{SuFunctor.Adequacy.}\texttt{\bfseries demonic\_maxSlice}}\hfill{\footnotesize \path{lean/su-functor/SuFunctor/Adequacy/MdpDemonic.lean}, line 184}\par\nopagebreak
\begin{leansrc}
theorem demonic_maxSlice (M : FinMDP S A) (k : ℕ)
    (ν : SubPMF (Traj S k)) :
    maxSlice (M.demonicCredal k) ν ↔ M.demonicLaws k ν
\end{leansrc}
\par\medskip\noindent{\small \texttt{SuFunctor.Adequacy.}\texttt{\bfseries satCredal\_lower\_demonicCredal}}\hfill{\footnotesize \path{lean/su-functor/SuFunctor/Adequacy/MdpDemonic.lean}, line 205}\par\nopagebreak
\begin{leansrc}
theorem satCredal_lower_demonicCredal (M : FinMDP S A) (k : ℕ)
    (c : SubPMF (Traj S k) → Prop) :
    SatCredal .lower (M.demonicCredal k) c ↔
      ∀ π : Policy S A, c (M.pathLaw π k)
\end{leansrc}
\par\medskip\noindent{\small \texttt{SuFunctor.Adequacy.}\texttt{\bfseries satCredal\_upper\_demonicCredal}}\hfill{\footnotesize \path{lean/su-functor/SuFunctor/Adequacy/MdpDemonic.lean}, line 214}\par\nopagebreak
\begin{leansrc}
theorem satCredal_upper_demonicCredal (M : FinMDP S A) (k : ℕ)
    {c : SubPMF (Traj S k) → Prop} (hc : DownClosed c) :
    SatCredal .upper (M.demonicCredal k) c ↔
      ∀ π : Policy S A, c (M.pathLaw π k)
\end{leansrc}
\par\medskip\noindent{\small \texttt{SuFunctor.Adequacy.}\texttt{\bfseries not\_satCredal\_upper\_demonicCredal\_geC}}\hfill{\footnotesize \path{lean/su-functor/SuFunctor/Adequacy/MdpDemonic.lean}, line 265}\par\nopagebreak
\begin{leansrc}
theorem not_satCredal_upper_demonicCredal_geC (M : FinMDP S A) (k : ℕ)
    [Nonempty A] {E : Set (Traj S k)} {q : ℚ} (hq : 0 < q) :
    ¬ SatCredal .upper (M.demonicCredal k) (geC E q)
\end{leansrc}
\par\medskip\noindent{\small \texttt{SuFunctor.Adequacy.}\texttt{\bfseries isEmpty\_policy\_iff}}\hfill{\footnotesize \path{lean/su-functor/SuFunctor/Adequacy/MdpDemonic.lean}, line 290}\par\nopagebreak
\begin{leansrc}
theorem isEmpty_policy_iff (M : FinMDP S A) :
    IsEmpty (Policy S A) ↔ IsEmpty A
\end{leansrc}
\par\medskip\noindent{\small \texttt{SuFunctor.}\texttt{\bfseries not\_maxSlice\_unionRays\_halfPair}}\hfill{\footnotesize \path{lean/su-functor/SuFunctor/Adequacy/EnvelopeUnion.lean}, line 316}\par\nopagebreak
\begin{leansrc}
theorem not_maxSlice_unionRays_halfPair (μ : SubPMF V) :
    ¬ maxSlice (unionRays (halfPair μ)) (halfPair μ false)
\end{leansrc}
\par\medskip\noindent{\small \texttt{SuFunctor.Adequacy.}\texttt{\bfseries ma\_adequacy}}\hfill{\footnotesize \path{lean/su-functor/SuFunctor/Adequacy/MarkovAutomaton.lean}, line 839}\par\nopagebreak
\begin{leansrc}
theorem ma_adequacy (M : FinMA S I A) (σ : Policy I A) (r : I → ℕ)
    (hr : ZenoFree (M.imStep σ) r) (k : ℕ) {a : Crd (Traj S k)} :
    ∃ α : Cell
        (looseConj (atoms (jumpRows
          (closeMaRows (compileMa M) (compilePolicy σ) (absFuel r))) k) :
          Loose a a)
        (singletonC (M.pathLaw σ r hr k)) (𝟙 a) (𝟙 a),
      Cell.IsInvertible α
\end{leansrc}
\par\medskip\noindent{\small \texttt{SuFunctor.Adequacy.}\texttt{\bfseries compile\_closeMa\_comm}}\hfill{\footnotesize \path{lean/su-functor/SuFunctor/Adequacy/MarkovAutomaton.lean}, line 819}\par\nopagebreak
\begin{leansrc}
theorem compile_closeMa_comm (M : FinMA S I A) (σ : Policy I A)
    (r : I → ℕ) (hr : ZenoFree (M.imStep σ) r) :
    closeMaRows (compileMa M) (compilePolicy σ) (absFuel r)
      = compileCtmc (M.close σ r hr)
\end{leansrc}
\par\medskip\noindent{\small \texttt{SuFunctor.Adequacy.}\texttt{\bfseries absorbIter\_mass}}\hfill{\footnotesize \path{lean/su-functor/SuFunctor/Adequacy/MarkovAutomaton.lean}, line 246}\par\nopagebreak
\begin{leansrc}
theorem absorbIter_mass (f : I → S ⊕ I → ℚ) (r : I → ℕ)
    (hr : ZenoFree f r) (hrow : ∀ i, ∑ y, f i y = 1) :
    ∀ (n : ℕ) (i : I), r i < n →
      ∑ m, absorbIter f n (Sum.inr i) m = 1
  | 0, _, hn => absurd hn (Nat.not_lt_zero _)
  | n + 1, i, hn => by
    simp_rw [absorbIter_succ_inr]
    rw [Finset.sum_comm]
    simp_rw [← Finset.mul_sum]
    have hterm : ∀ y, f i y * (∑ m, absorbIter f n y m) = f i y * 1
\end{leansrc}
\par\medskip\noindent{\small \texttt{SuFunctor.Adequacy.}\texttt{\bfseries absorbIter\_unique}}\hfill{\footnotesize \path{lean/su-functor/SuFunctor/Adequacy/MarkovAutomaton.lean}, line 276}\par\nopagebreak
\begin{leansrc}
theorem absorbIter_unique (f : I → S ⊕ I → ℚ) (r : I → ℕ)
    (hr : ZenoFree f r) (g : S ⊕ I → S → ℚ)
    (hg_inl : ∀ m' m, g (Sum.inl m') m = hitMass (Sum.inl m' : S ⊕ I) m)
    (hg_step : ∀ i m, g (Sum.inr i) m = ∑ y, f i y * g y m) :
    ∀ (n : ℕ) (i : I) (m : S), r i < n →
      g (Sum.inr i) m = absorbIter f n (Sum.inr i) m
  | 0, _, _, hn => absurd hn (Nat.not_lt_zero _)
  | n + 1, i, m, hn => by
    rw [hg_step, absorbIter_succ_inr]
    refine Finset.sum_congr rfl fun y _ => ?_
    by_cases hz : f i y = 0
    · rw [hz, zero_mul, zero_mul]
    · congr 1
      cases y with
      | inl m' => rw [hg_inl, absorbIter_inl]
      | inr j =>
        exact absorbIter_unique f r hr g hg_inl hg_step n j m
          (lt_of_lt_of_le (hr i j hz) (Nat.lt_succ_iff.mp hn))

def absFuel (r : I → ℕ) : ℕ
\end{leansrc}
\par\medskip\noindent{\small \texttt{SuFunctor.Adequacy.}\texttt{\bfseries absorbIter\_stab}}\hfill{\footnotesize \path{lean/su-functor/SuFunctor/Adequacy/MarkovAutomaton.lean}, line 223}\par\nopagebreak
\begin{leansrc}
theorem absorbIter_stab (f : I → S ⊕ I → ℚ) (r : I → ℕ)
    (hr : ZenoFree f r) (i : I) (m : S) {n n' : ℕ}
    (h : r i < n) (h' : r i < n') :
    absorbIter f n (Sum.inr i) m = absorbIter f n' (Sum.inr i) m
\end{leansrc}
\par\medskip\noindent{\small \texttt{SuFunctor.Adequacy.FinMA.}\texttt{\bfseries close\_offRate\_factor}}\hfill{\footnotesize \path{lean/su-functor/SuFunctor/Adequacy/MarkovAutomaton.lean}, line 540}\par\nopagebreak
\begin{leansrc}
theorem FinMA.close_offRate_factor (M : FinMA S I A) (σ : Policy I A)
    (r : I → ℕ) (hr : ZenoFree (M.imStep σ) r) {s s' : S}
    (hne : s ≠ s') :
    (M.close σ r hr).offRate s s'
      = M.exitRate s * ∑ x, M.preChain.jumpChain.step (Sum.inl s) x
          * M.absorb σ r x s'
\end{leansrc}
\par\medskip\noindent{\small \texttt{SuFunctor.Adequacy.FinMA.}\texttt{\bfseries close\_exitRate\_thinning}}\hfill{\footnotesize \path{lean/su-functor/SuFunctor/Adequacy/MarkovAutomaton.lean}, line 562}\par\nopagebreak
\begin{leansrc}
theorem FinMA.close_exitRate_thinning (M : FinMA S I A)
    (σ : Policy I A) (r : I → ℕ) (hr : ZenoFree (M.imStep σ) r)
    (s : S) :
    (M.close σ r hr).exitRate s
      = M.exitRate s - ∑ x, M.offRate s x * M.absorb σ r x s
\end{leansrc}
\par\medskip\noindent{\small \texttt{SuFunctor.Adequacy.FinMA.}\texttt{\bfseries close\_exitRate\_le}}\hfill{\footnotesize \path{lean/su-functor/SuFunctor/Adequacy/MarkovAutomaton.lean}, line 594}\par\nopagebreak
\begin{leansrc}
theorem FinMA.close_exitRate_le (M : FinMA S I A) (σ : Policy I A)
    (r : I → ℕ) (hr : ZenoFree (M.imStep σ) r) (s : S) :
    (M.close σ r hr).exitRate s ≤ M.exitRate s
\end{leansrc}
\par\medskip\noindent{\small \texttt{SuFunctor.Adequacy.}\texttt{\bfseries compileMa\_diag\_blind}}\hfill{\footnotesize \path{lean/su-functor/SuFunctor/Adequacy/MarkovAutomaton.lean}, line 713}\par\nopagebreak
\begin{leansrc}
theorem compileMa_diag_blind (M M' : FinMA S I A)
    (hoff : ∀ s x, Sum.inl s ≠ x → M.rate s x = M'.rate s x)
    (hbr : M.branch = M'.branch) (hinit : M.init = M'.init) :
    compileMa M = compileMa M'
\end{leansrc}
\par\medskip\noindent{\small \texttt{SuFunctor.Adequacy.FinMA.}\texttt{\bfseries trajPMF\_eq}}\hfill{\footnotesize \path{lean/su-functor/SuFunctor/Adequacy/MarkovAutomaton.lean}, line 869}\par\nopagebreak
\begin{leansrc}
theorem FinMA.trajPMF_eq (M : FinMA S I A) (σ : Policy I A)
    (r : I → ℕ) (hr : ZenoFree (M.imStep σ) r) (k : ℕ)
    (ω : Traj S k) :
    M.trajPMF σ r hr k ω = q2e ((M.pathLaw σ r hr k).wt ω)
\end{leansrc}
\par\medskip\noindent{\small \texttt{SuFunctor.Adequacy.}\texttt{\bfseries ma\_demonic\_adequacy}}\hfill{\footnotesize \path{lean/su-functor/SuFunctor/Adequacy/MaDemonic.lean}, line 299}\par\nopagebreak
\begin{leansrc}
theorem ma_demonic_adequacy (M : FinMA S I A) {r : I → ℕ}
    (hb : M.BranchZenoFree r) (k : ℕ) {a : Crd (Traj S k)} :
    ∃ α : Cell (maDemonicC (compileMa M) (absFuel r) k : Loose a a)
        (M.demonicLaws r hb k) (𝟙 a) (𝟙 a),
      Cell.IsInvertible α
\end{leansrc}
\par\medskip\noindent{\small \texttt{SuFunctor.Adequacy.}\texttt{\bfseries maDemonicC}}\hfill{\footnotesize \path{lean/su-functor/SuFunctor/Adequacy/MaDemonic.lean}, line 238}\par\nopagebreak
\begin{leansrc}
def maDemonicC (R : MaRows S I A) (n k : ℕ) :
    SubPMF (Traj S k) → Prop
\end{leansrc}
\par\medskip\noindent{\small \texttt{SuFunctor.Adequacy.}\texttt{\bfseries maDemonicC\_iff}}\hfill{\footnotesize \path{lean/su-functor/SuFunctor/Adequacy/MaDemonic.lean}, line 262}\par\nopagebreak
\begin{leansrc}
theorem maDemonicC_iff (M : FinMA S I A) {r : I → ℕ}
    (hb : M.BranchZenoFree r) (k : ℕ) (μ : SubPMF (Traj S k)) :
    maDemonicC (compileMa M) (absFuel r) k μ ↔
      M.demonicLaws r hb k μ
\end{leansrc}
\par\medskip\noindent{\small \texttt{SuFunctor.Adequacy.FinMA.}\texttt{\bfseries BranchZenoFree}}\hfill{\footnotesize \path{lean/su-functor/SuFunctor/Adequacy/MaDemonic.lean}, line 122}\par\nopagebreak
\begin{leansrc}
def FinMA.BranchZenoFree (M : FinMA S I A) (r : I → ℕ) : Prop
\end{leansrc}
\par\medskip\noindent{\small \texttt{SuFunctor.Adequacy.}\texttt{\bfseries zenoFree\_of\_branch}}\hfill{\footnotesize \path{lean/su-functor/SuFunctor/Adequacy/MarkovAutomaton.lean}, line 411}\par\nopagebreak
\begin{leansrc}
theorem zenoFree_of_branch (M : FinMA S I A) (r : I → ℕ)
    (hb : ∀ i a j, M.branch i a (Sum.inr j) ≠ 0 → r j < r i)
    (σ : Policy I A) : ZenoFree (M.imStep σ) r
\end{leansrc}
\par\medskip\noindent{\small \texttt{SuFunctor.Adequacy.FinMA.}\texttt{\bfseries pathLaw\_rank\_irrel}}\hfill{\footnotesize \path{lean/su-functor/SuFunctor/Adequacy/MaDemonic.lean}, line 169}\par\nopagebreak
\begin{leansrc}
theorem FinMA.pathLaw_rank_irrel (M : FinMA S I A) (σ : Policy I A)
    {r r' : I → ℕ} (hr : ZenoFree (M.imStep σ) r)
    (hr' : ZenoFree (M.imStep σ) r') (k : ℕ) :
    M.pathLaw σ r hr k = M.pathLaw σ r' hr' k
\end{leansrc}
\par\medskip\noindent{\small \texttt{SuFunctor.Adequacy.FinMA.}\texttt{\bfseries demonicLaws\_rank\_irrel}}\hfill{\footnotesize \path{lean/su-functor/SuFunctor/Adequacy/MaDemonic.lean}, line 207}\par\nopagebreak
\begin{leansrc}
theorem FinMA.demonicLaws_rank_irrel (M : FinMA S I A) {r r' : I → ℕ}
    (hb : M.BranchZenoFree r) (hb' : M.BranchZenoFree r') (k : ℕ) :
    M.demonicLaws r hb k = M.demonicLaws r' hb' k
\end{leansrc}
\par\medskip\noindent{\small \texttt{SuFunctor.Adequacy.}\texttt{\bfseries ma\_demonic\_maxSlice}}\hfill{\footnotesize \path{lean/su-functor/SuFunctor/Adequacy/MaDemonic.lean}, line 326}\par\nopagebreak
\begin{leansrc}
theorem ma_demonic_maxSlice (M : FinMA S I A) {r : I → ℕ}
    (hb : M.BranchZenoFree r) (k : ℕ) (ν : SubPMF (Traj S k)) :
    maxSlice (M.demonicCredal r hb k) ν ↔ M.demonicLaws r hb k ν
\end{leansrc}
\par\medskip\noindent{\small \texttt{SuFunctor.Adequacy.}\texttt{\bfseries satCredal\_lower\_maDemonicCredal}}\hfill{\footnotesize \path{lean/su-functor/SuFunctor/Adequacy/MaDemonic.lean}, line 349}\par\nopagebreak
\begin{leansrc}
theorem satCredal_lower_maDemonicCredal (M : FinMA S I A) {r : I → ℕ}
    (hb : M.BranchZenoFree r) (k : ℕ) (c : SubPMF (Traj S k) → Prop) :
    SatCredal .lower (M.demonicCredal r hb k) c ↔
      ∀ σ : Policy I A, c (M.pathLaw σ r (hb.zenoFree σ) k)
\end{leansrc}
\par\medskip\noindent{\small \texttt{SuFunctor.Adequacy.}\texttt{\bfseries satCredal\_upper\_maDemonicCredal}}\hfill{\footnotesize \path{lean/su-functor/SuFunctor/Adequacy/MaDemonic.lean}, line 359}\par\nopagebreak
\begin{leansrc}
theorem satCredal_upper_maDemonicCredal (M : FinMA S I A) {r : I → ℕ}
    (hb : M.BranchZenoFree r) (k : ℕ) {c : SubPMF (Traj S k) → Prop}
    (hc : DownClosed c) :
    SatCredal .upper (M.demonicCredal r hb k) c ↔
      ∀ σ : Policy I A, c (M.pathLaw σ r (hb.zenoFree σ) k)
\end{leansrc}
\par\medskip\noindent{\small \texttt{SuFunctor.Adequacy.Policy.}\texttt{\bfseries nonempty\_of\_isEmpty}}\hfill{\footnotesize \path{lean/su-functor/SuFunctor/Adequacy/MaDemonic.lean}, line 406}\par\nopagebreak
\begin{leansrc}
theorem Policy.nonempty_of_isEmpty {X : Type} [IsEmpty X] :
    Nonempty (Policy X A)
\end{leansrc}
\par\medskip\noindent{\small \texttt{SuFunctor.Adequacy.Policy.}\texttt{\bfseries isEmpty\_iff}}\hfill{\footnotesize \path{lean/su-functor/SuFunctor/Adequacy/MaDemonic.lean}, line 418}\par\nopagebreak
\begin{leansrc}
theorem Policy.isEmpty_iff {X : Type} [Nonempty X] [DecidableEq A] :
    IsEmpty (Policy X A) ↔ IsEmpty A
\end{leansrc}
\par\medskip\noindent{\small \texttt{SuFunctor.Adequacy.}\texttt{\bfseries not\_satCredal\_upper\_maDemonicCredal\_geC}}\hfill{\footnotesize \path{lean/su-functor/SuFunctor/Adequacy/MaDemonic.lean}, line 438}\par\nopagebreak
\begin{leansrc}
theorem not_satCredal_upper_maDemonicCredal_geC (M : FinMA S I A)
    {r : I → ℕ} (hb : M.BranchZenoFree r) (k : ℕ)
    [hp : Nonempty (Policy I A)] {E : Set (Traj S k)} {q : ℚ}
    (hq : 0 < q) :
    ¬ SatCredal .upper (M.demonicCredal r hb k) (geC E q)
\end{leansrc}
\par\medskip\noindent{\small \texttt{SuFunctor.Adequacy.FinMA.}\texttt{\bfseries nonempty\_markovian}}\hfill{\footnotesize \path{lean/su-functor/SuFunctor/Adequacy/MaDemonic.lean}, line 474}\par\nopagebreak
\begin{leansrc}
theorem FinMA.nonempty_markovian (M : FinMA S I A) [Nonempty I]
    [Nonempty A] {r : I → ℕ} (hb : M.BranchZenoFree r) : Nonempty S
\end{leansrc}
\par\medskip\noindent{\small \texttt{SuFunctor.Adequacy.FinMA.}\texttt{\bfseries CertScheduler}}\hfill{\footnotesize \path{lean/su-functor/SuFunctor/Adequacy/MaDemonic.lean}, line 504}\par\nopagebreak
\begin{leansrc}
structure FinMA.CertScheduler (M : FinMA S I A) where …
\end{leansrc}
\par\medskip\noindent{\small \texttt{SuFunctor.Adequacy.FinMA.}\texttt{\bfseries demonicLawsCert}}\hfill{\footnotesize \path{lean/su-functor/SuFunctor/Adequacy/MaDemonic.lean}, line 515}\par\nopagebreak
\begin{leansrc}
def FinMA.demonicLawsCert (M : FinMA S I A) (k : ℕ) :
    SubPMF (Traj S k) → Prop
\end{leansrc}
\par\medskip\noindent{\small \texttt{SuFunctor.Adequacy.FinMA.}\texttt{\bfseries demonicLawsCert\_iff}}\hfill{\footnotesize \path{lean/su-functor/SuFunctor/Adequacy/MaDemonic.lean}, line 531}\par\nopagebreak
\begin{leansrc}
theorem FinMA.demonicLawsCert_iff (M : FinMA S I A) {r : I → ℕ}
    (hb : M.BranchZenoFree r) (k : ℕ) (ν : SubPMF (Traj S k)) :
    M.demonicLawsCert k ν ↔ M.demonicLaws r hb k ν
\end{leansrc}
\par\medskip\noindent{\small \texttt{SuFunctor.Adequacy.Bag.}\texttt{\bfseries sortTuple\_apply\_eq\_inf\_sup}}\hfill{\footnotesize \path{lean/su-functor/SuFunctor/Adequacy/Bag.lean}, line 128}\par\nopagebreak
\begin{leansrc}
theorem sortTuple_apply_eq_inf_sup (x : Fin n → B) (k : Fin n) :
    sortTuple x k =
      ((univ : Finset (Fin n)).powersetCard (k.val + 1)).attach.inf'
        ((powersetCard_succ_nonempty k).attach)
        (fun S => S.1.sup' (powersetCard_mem_nonempty S.2) x)
\end{leansrc}
\par\medskip\noindent{\small \texttt{SuFunctor.Adequacy.Bag.}\texttt{\bfseries measurable\_symSec}}\hfill{\footnotesize \path{lean/su-functor/SuFunctor/Adequacy/Bag.lean}, line 377}\par\nopagebreak
\begin{leansrc}
theorem measurable_symSec
    (hle : MeasurableSet {p : B × B | p.1 ≤ p.2}) :
    Measurable (symSec : Sym B n → (Fin n → B))
\end{leansrc}
\par\medskip\noindent{\small \texttt{SuFunctor.Adequacy.Bag.}\texttt{\bfseries monotone\_eq\_of\_toSym\_eq}}\hfill{\footnotesize \path{lean/su-functor/SuFunctor/Adequacy/Bag.lean}, line 335}\par\nopagebreak
\begin{leansrc}
theorem monotone_eq_of_toSym_eq {x y : Fin n → B}
    (hx : Monotone x) (hy : Monotone y) (h : toSym x = toSym y) :
    x = y
\end{leansrc}
\par\medskip\noindent{\small \texttt{SuFunctor.Adequacy.Bag.}\texttt{\bfseries quotKernel\_comp\_eq}}\hfill{\footnotesize \path{lean/su-functor/SuFunctor/Adequacy/Bag.lean}, line 411}\par\nopagebreak
\begin{leansrc}
theorem quotKernel_comp_eq (L : Kernel α T) (hs : Measurable s)
    (hL : ∀ x, L (s (q x)) = L x) (x : α) :
    quotKernel L hs (q x) = L x
\end{leansrc}
\par\medskip\noindent{\small \texttt{SuFunctor.Adequacy.Bag.}\texttt{\bfseries iter\_intertwine}}\hfill{\footnotesize \path{lean/su-functor/SuFunctor/Adequacy/Bag.lean}, line 451}\par\nopagebreak
\begin{leansrc}
theorem iter_intertwine (hq : Measurable q) (hs : Measurable s)
    (K : Kernel α α)
    (hK : ∀ x, (K.map q) (s (q x)) = (K.map q) x) (k : ℕ) :
    Kernel.deterministic q hq ∘ₖ iterK K k =
      iterK (quotKernel (K.map q) hs) k ∘ₖ Kernel.deterministic q hq
\end{leansrc}
\par\medskip\noindent{\small \texttt{SuFunctor.Adequacy.Bag.}\texttt{\bfseries comap\_section\_congr}}\hfill{\footnotesize \path{lean/su-functor/SuFunctor/Adequacy/Bag.lean}, line 419}\par\nopagebreak
\begin{leansrc}
theorem comap_section_congr (L : Kernel α T)
    {s' : ᾱ → α} (hs : Measurable s) (hs' : Measurable s')
    (hsec' : ∀ m, q (s' m) = m)
    (hL : ∀ x, L (s (q x)) = L x) :
    L.comap s' hs' = L.comap s hs
\end{leansrc}
\par\medskip\noindent{\small \texttt{SuFunctor.Adequacy.Bag.}\texttt{\bfseries bag\_transport}}\hfill{\footnotesize \path{lean/su-functor/SuFunctor/Adequacy/Bag.lean}, line 598}\par\nopagebreak
\begin{leansrc}
theorem bag_transport [LinearOrder B] [MeasurableSpace B]
    (hle : MeasurableSet {p : B × B | p.1 ≤ p.2})
    (K : Kernel (BagPoints B) (BagPoints B))
    (hK : ∀ (n : ℕ) (x : Fin n → B) (σ : Equiv.Perm (Fin n)),
      (K.map toSymSigma) ⟨n, x ∘ σ⟩ = (K.map toSymSigma) ⟨n, x⟩)
    (k : ℕ) (p : BagPoints B) :
    (iterK K k p).map toSymSigma =
      iterK (quotKernel (K.map toSymSigma) (measurable_symSecSigma hle)) k
        (toSymSigma p)
\end{leansrc}
\par\medskip\noindent{\small \texttt{SuFunctor.Adequacy.Bag.}\texttt{\bfseries invariance\_iff\_perm}}\hfill{\footnotesize \path{lean/su-functor/SuFunctor/Adequacy/Bag.lean}, line 572}\par\nopagebreak
\begin{leansrc}
theorem invariance_iff_perm [LinearOrder B] [MeasurableSpace B]
    {T : Type*} [MeasurableSpace T]
    (L : Kernel (BagPoints B) T) :
    (∀ p : BagPoints B, L (symSecSigma (toSymSigma p)) = L p) ↔
      (∀ (n : ℕ) (x : Fin n → B) (σ : Equiv.Perm (Fin n)),
        L ⟨n, x ∘ σ⟩ = L ⟨n, x⟩)
\end{leansrc}
\par\medskip\noindent{\small \texttt{SuFunctor.Adequacy.Bag.}\texttt{\bfseries measurableSet\_le\_lex}}\hfill{\footnotesize \path{lean/su-functor/SuFunctor/Adequacy/Bag.lean}, line 627}\par\nopagebreak
\begin{leansrc}
theorem measurableSet_le_lex (d : ℕ) :
    MeasurableSet {q : (Fin d → ℝ) × (Fin d → ℝ) |
      toLex q.1 ≤ toLex q.2}
\end{leansrc}
\par\medskip\noindent{\small \texttt{SuFunctor.Adequacy.FinISPN.}\texttt{\bfseries not\_branchZenoFree\_of\_disabled}}\hfill{\footnotesize \path{lean/su-functor/SuFunctor/Adequacy/StochasticPetriDemonic.lean}, line 105}\par\nopagebreak
\begin{leansrc}
theorem not_branchZenoFree_of_disabled {i : N.Vanishing} {u : U}
    (h : ¬ N.IEnabled u i.1) (r : N.Vanishing → ℕ) :
    ¬ N.toFinMA.BranchZenoFree r
\end{leansrc}
\par\medskip\noindent{\small \texttt{SuFunctor.Adequacy.FinISPN.}\texttt{\bfseries RespScheduler}}\hfill{\footnotesize \path{lean/su-functor/SuFunctor/Adequacy/StochasticPetriDemonic.lean}, line 121}\par\nopagebreak
\begin{leansrc}
def RespScheduler : Type
\end{leansrc}
\par\medskip\noindent{\small \texttt{SuFunctor.Adequacy.SpnDemonicControls.}\texttt{\bfseries arena\_plural}}\hfill{\footnotesize \path{lean/su-functor/SuFunctor/Adequacy/StochasticPetriDemonic.lean}, line 532}\par\nopagebreak
\begin{leansrc}
theorem arena_plural :
    (net₁.uniformResp : net₁.RespScheduler).1.choice i10 1
      ≠ discardResp.1.choice i10 1
\end{leansrc}
\par\medskip\noindent{\small \texttt{SuFunctor.Adequacy.FinISPN.}\texttt{\bfseries RespectingRows}}\hfill{\footnotesize \path{lean/su-functor/SuFunctor/Adequacy/StochasticPetriDemonic.lean}, line 193}\par\nopagebreak
\begin{leansrc}
structure RespectingRows where …
\end{leansrc}
\par\medskip\noindent{\small \texttt{SuFunctor.Adequacy.FinISPN.}\texttt{\bfseries ispn\_demonic\_adequacy}}\hfill{\footnotesize \path{lean/su-functor/SuFunctor/Adequacy/StochasticPetriDemonic.lean}, line 295}\par\nopagebreak
\begin{leansrc}
theorem ispn_demonic_adequacy {r : N.Vanishing → ℕ}
    (hr : N.CascadeRank r) (k : ℕ) {a : Crd (Traj N.Tangible k)} :
    ∃ α : Cell (N.ispnDemonicC (absFuel r) k : Loose a a)
        (N.demonicTangibleLaws hr k) (𝟙 a) (𝟙 a),
      Cell.IsInvertible α
\end{leansrc}
\par\medskip\noindent{\small \texttt{SuFunctor.Adequacy.FinISPN.}\texttt{\bfseries ispnDemonicC\_iff}}\hfill{\footnotesize \path{lean/su-functor/SuFunctor/Adequacy/StochasticPetriDemonic.lean}, line 250}\par\nopagebreak
\begin{leansrc}
theorem ispnDemonicC_iff {r : N.Vanishing → ℕ} (hr : N.CascadeRank r)
    (k : ℕ) (μ : SubPMF (Traj N.Tangible k)) :
    N.ispnDemonicC (absFuel r) k μ ↔ N.demonicTangibleLaws hr k μ
\end{leansrc}
\par\medskip\noindent{\small \texttt{SuFunctor.Adequacy.FinISPN.}\texttt{\bfseries demonicTangibleLaws\_rank\_irrel}}\hfill{\footnotesize \path{lean/su-functor/SuFunctor/Adequacy/StochasticPetriDemonic.lean}, line 158}\par\nopagebreak
\begin{leansrc}
theorem demonicTangibleLaws_rank_irrel {r r' : N.Vanishing → ℕ}
    (hr : N.CascadeRank r) (hr' : N.CascadeRank r') (k : ℕ) :
    N.demonicTangibleLaws hr k = N.demonicTangibleLaws hr' k
\end{leansrc}
\par\medskip\noindent{\small \texttt{SuFunctor.Adequacy.FinISPN.}\texttt{\bfseries demonicTangibleCredal\_rank\_irrel}}\hfill{\footnotesize \path{lean/su-functor/SuFunctor/Adequacy/StochasticPetriDemonic.lean}, line 172}\par\nopagebreak
\begin{leansrc}
theorem demonicTangibleCredal_rank_irrel {r r' : N.Vanishing → ℕ}
    (hr : N.CascadeRank r) (hr' : N.CascadeRank r') (k : ℕ) :
    N.demonicTangibleCredal hr k = N.demonicTangibleCredal hr' k
\end{leansrc}
\par\medskip\noindent{\small \texttt{SuFunctor.Adequacy.FinISPN.}\texttt{\bfseries ispn\_demonic\_maxSlice}}\hfill{\footnotesize \path{lean/su-functor/SuFunctor/Adequacy/StochasticPetriDemonic.lean}, line 321}\par\nopagebreak
\begin{leansrc}
theorem ispn_demonic_maxSlice {r : N.Vanishing → ℕ}
    (hr : N.CascadeRank r) (k : ℕ) (ν : SubPMF (Traj N.Tangible k)) :
    maxSlice (N.demonicTangibleCredal hr k) ν ↔
      N.demonicTangibleLaws hr k ν
\end{leansrc}
\par\medskip\noindent{\small \texttt{SuFunctor.Adequacy.FinISPN.}\texttt{\bfseries satCredal\_lower\_ispnDemonicCredal}}\hfill{\footnotesize \path{lean/su-functor/SuFunctor/Adequacy/StochasticPetriDemonic.lean}, line 346}\par\nopagebreak
\begin{leansrc}
theorem satCredal_lower_ispnDemonicCredal {r : N.Vanishing → ℕ}
    (hr : N.CascadeRank r) (k : ℕ)
    (c : SubPMF (Traj N.Tangible k) → Prop) :
    SatCredal .lower (N.demonicTangibleCredal hr k) c ↔
      ∀ σ : N.RespScheduler,
        c (N.tangiblePathLaw σ.1 r (N.respZenoFree hr σ) k)
\end{leansrc}
\par\medskip\noindent{\small \texttt{SuFunctor.Adequacy.FinISPN.}\texttt{\bfseries su\_ispn\_demonic\_adequacy}}\hfill{\footnotesize \path{lean/su-functor/SuFunctor/Adequacy/StochasticPetriDemonic.lean}, line 412}\par\nopagebreak
\begin{leansrc}
theorem su_ispn_demonic_adequacy {r : N.Vanishing → ℕ}
    (hr : N.CascadeRank r) (k : ℕ) (μ : SubPMF (Traj N.Tangible k)) :
    idealPoint μ ∈ (N.ispnDemonicDnf (absFuel r) k).modelClass
      ↔ N.demonicTangibleLaws hr k μ
\end{leansrc}
\par\medskip\noindent{\small \texttt{SuFunctor.Adequacy.FinISPN.}\texttt{\bfseries chainFire\_val}}\hfill{\footnotesize \path{lean/su-functor/SuFunctor/Adequacy/StochasticPetriSetFiring.lean}, line 196}\par\nopagebreak
\begin{leansrc}
theorem chainFire_val : ∀ (l : List U) (m : BMarking P B),
    N.ChainEnabled l m → ∀ p : P,
    (N.chainFire l m p : ℕ) + N.preSum ↑l p
      = (m p : ℕ) + N.postSum ↑l p
  | [], m, _, p => by simp
  | u :: l, m, h, p => by
    obtain ⟨hu, hl⟩
\end{leansrc}
\par\medskip\noindent{\small \texttt{SuFunctor.Adequacy.FinISPN.}\texttt{\bfseries chainFire\_eq\_of\_enabled}}\hfill{\footnotesize \path{lean/su-functor/SuFunctor/Adequacy/StochasticPetriSetFiring.lean}, line 213}\par\nopagebreak
\begin{leansrc}
theorem chainFire_eq_of_enabled {l₁ l₂ : List U} {m : BMarking P B}
    (h₁ : N.ChainEnabled l₁ m) (h₂ : N.ChainEnabled l₂ m)
    (hml : (↑l₁ : Multiset U) = ↑l₂) :
    N.chainFire l₁ m = N.chainFire l₂ m
\end{leansrc}
\par\medskip\noindent{\small \texttt{SuFunctor.Adequacy.FinISPN.}\texttt{\bfseries chainFire\_eq\_compositeFire}}\hfill{\footnotesize \path{lean/su-functor/SuFunctor/Adequacy/StochasticPetriSetFiring.lean}, line 262}\par\nopagebreak
\begin{leansrc}
theorem chainFire_eq_compositeFire {l : List U} {m : BMarking P B}
    (h : N.ChainEnabled l m) (hsup : N.Supplied ↑l m) :
    N.chainFire l m = N.compositeFire ↑l m
\end{leansrc}
\par\medskip\noindent{\small \texttt{SuFunctor.Adequacy.FinISPN.}\texttt{\bfseries JointReservation}}\hfill{\footnotesize \path{lean/su-functor/SuFunctor/Adequacy/StochasticPetriSetFiring.lean}, line 296}\par\nopagebreak
\begin{leansrc}
def JointReservation (S : Multiset U) (m : BMarking P B) : Prop
\end{leansrc}
\par\medskip\noindent{\small \texttt{SuFunctor.Adequacy.FinISPN.}\texttt{\bfseries setNonConflicting\_of\_headroom}}\hfill{\footnotesize \path{lean/su-functor/SuFunctor/Adequacy/StochasticPetriSetFiring.lean}, line 377}\par\nopagebreak
\begin{leansrc}
theorem setNonConflicting_of_headroom {S : Multiset U} {m : BMarking P B}
    (hJ : N.JointReservation S m)
    (hcap : ∀ p, (m p : ℕ) + N.postSum S p ≤ B) :
    N.SetNonConflicting S m
\end{leansrc}
\par\medskip\noindent{\small \texttt{SuFunctor.Adequacy.FinISPN.}\texttt{\bfseries r1r3\_composite\_adequacy}}\hfill{\footnotesize \path{lean/su-functor/SuFunctor/Adequacy/StochasticPetriSetFiring.lean}, line 391}\par\nopagebreak
\begin{leansrc}
theorem r1r3_composite_adequacy {S : Multiset U} {m : BMarking P B}
    (hJ : N.JointReservation S m)
    (hcap : ∀ p, (m p : ℕ) + N.postSum S p ≤ B) :
    ∀ l : List U, (↑l : Multiset U) = S →
      N.ChainEnabled l m ∧ N.chainFire l m = N.compositeFire S m
\end{leansrc}
\par\medskip\noindent{\small \texttt{SuFunctor.Adequacy.FinISPN.}\texttt{\bfseries flatten\_triangle\_of\_agree}}\hfill{\footnotesize \path{lean/su-functor/SuFunctor/Adequacy/StochasticPetriFlatten.lean}, line 324}\par\nopagebreak
\begin{leansrc}
theorem flatten_triangle_of_agree (σ : Policy N.Vanishing U)
    (r : N.Vanishing → ℕ) (hr : ZenoFree (N.toFinMA.imStep σ) r)
    {Mode : Type} [Fintype Mode] (e : Mode ≃ N.Tangible)
    (flat : FinCTMC Mode)
    (hoff : ∀ s s' : N.Tangible, s ≠ s' →
      (flat.relabel e).rate s s'
        = (N.toFinMA.close σ r hr).rate s s')
    (hinit : (flat.relabel e).init = (N.toFinMA.close σ r hr).init)
    (k : ℕ) {a : Crd (Traj N.Tangible k)} :
    ∃ α : Cell
        (looseConj (atoms (jumpRows
          (closeMaRows N.compileIspn (compilePolicy σ) (absFuel r))) k) :
          Loose a a)
        (looseConj (atoms (jumpRows (compileCtmc (flat.relabel e))) k))
        (𝟙 a) (𝟙 a),
      Cell.IsInvertible α
\end{leansrc}
\par\medskip\noindent{\small \texttt{SuFunctor.}\texttt{\bfseries SatAngelic}}\hfill{\footnotesize \path{lean/su-functor/SuFunctor/Adequacy/StochasticPetriAngelic.lean}, line 121}\par\nopagebreak
\begin{leansrc}
def SatAngelic (C c : SubPMF V → Prop) : Prop
\end{leansrc}
\par\medskip\noindent{\small \texttt{SuFunctor.Adequacy.FinISPN.}\texttt{\bfseries satAngelic\_ispnDemonicCredal}}\hfill{\footnotesize \path{lean/su-functor/SuFunctor/Adequacy/StochasticPetriAngelic.lean}, line 196}\par\nopagebreak
\begin{leansrc}
theorem satAngelic_ispnDemonicCredal {r : N.Vanishing → ℕ}
    (hr : N.CascadeRank r) (k : ℕ)
    (c : SubPMF (Traj N.Tangible k) → Prop) :
    SatAngelic (N.demonicTangibleCredal hr k) c ↔
      ∃ σ : N.RespScheduler,
        c (N.tangiblePathLaw σ.1 r (N.respZenoFree hr σ) k)
\end{leansrc}
\par\medskip\noindent{\small \texttt{SuFunctor.Adequacy.FinISPN.}\texttt{\bfseries ispn\_demonic\_verdict\_to\_member}}\hfill{\footnotesize \path{lean/su-functor/SuFunctor/Adequacy/StochasticPetriAngelic.lean}, line 220}\par\nopagebreak
\begin{leansrc}
theorem ispn_demonic_verdict_to_member {r : N.Vanishing → ℕ}
    (hr : N.CascadeRank r) (σ : N.RespScheduler) (k : ℕ)
    {c : SubPMF (Traj N.Tangible k) → Prop}
    (h : SatCredal .lower (N.demonicTangibleCredal hr k) c) :
    c (N.tangiblePathLaw σ.1 r (N.respZenoFree hr σ) k)
\end{leansrc}
\par\medskip\noindent{\small \texttt{SuFunctor.Adequacy.FinISPN.}\texttt{\bfseries flatten\_angelic\_of\_agree}}\hfill{\footnotesize \path{lean/su-functor/SuFunctor/Adequacy/StochasticPetriAngelic.lean}, line 264}\par\nopagebreak
\begin{leansrc}
theorem flatten_angelic_of_agree {r : N.Vanishing → ℕ}
    (hrk : N.CascadeRank r) (σ : N.RespScheduler)
    {Mode : Type} [Fintype Mode] (e : Mode ≃ N.Tangible)
    (flat : FinCTMC Mode)
    (hoff : ∀ s s' : N.Tangible, s ≠ s' →
      (flat.relabel e).rate s s'
        = (N.toFinMA.close σ.1 r (N.respZenoFree hrk σ)).rate s s')
    (hinit : (flat.relabel e).init
        = (N.toFinMA.close σ.1 r (N.respZenoFree hrk σ)).init)
    (k : ℕ) {c : SubPMF (Traj N.Tangible k) → Prop}
    (hc : c ((flat.relabel e).jumpLaw k)) :
    SatAngelic (N.demonicTangibleCredal hrk k) c
\end{leansrc}
\par\medskip\noindent{\small \texttt{SuFunctor.Adequacy.PetriVanilla.}\texttt{\bfseries face\_eq\_reach}}\hfill{\footnotesize \path{lean/su-functor/SuFunctor/Adequacy/PetriVanilla.lean}, line 409}\par\nopagebreak
\begin{leansrc}
theorem face_eq_reach (hd : ∀ t, 0 < d t) :
    face N d m0 = {m | Reach N m0 m}
\end{leansrc}
\par\medskip\noindent{\small \texttt{SuFunctor.Adequacy.PetriVanilla.}\texttt{\bfseries face\_dwell\_uniform}}\hfill{\footnotesize \path{lean/su-functor/SuFunctor/Adequacy/PetriVanilla.lean}, line 418}\par\nopagebreak
\begin{leansrc}
theorem face_dwell_uniform {d' : T → ℚ} (hd : ∀ t, 0 < d t)
    (hd' : ∀ t, 0 < d' t) : face N d m0 = face N d' m0
\end{leansrc}
\par\medskip\noindent{\small \texttt{SuFunctor.Adequacy.PetriVanilla.}\texttt{\bfseries su\_vanilla\_demonic\_adequacy}}\hfill{\footnotesize \path{lean/su-functor/SuFunctor/Adequacy/PetriVanilla.lean}, line 570}\par\nopagebreak
\begin{leansrc}
theorem su_vanilla_demonic_adequacy (k : ℕ)
    (μ : SubPMF (Traj (Marking P) k)) :
    idealPoint μ ∈ (vanillaDemonicDnf N d m0 k).modelClass
      ↔ vanillaDemonicLaws N d m0 k μ
\end{leansrc}
\par\medskip\noindent{\small \texttt{SuFunctor.Adequacy.PetriVanilla.}\texttt{\bfseries vanillaDemonicC\_iff}}\hfill{\footnotesize \path{lean/su-functor/SuFunctor/Adequacy/PetriVanilla.lean}, line 515}\par\nopagebreak
\begin{leansrc}
theorem vanillaDemonicC_iff (k : ℕ) (μ : SubPMF (Traj (Marking P) k)) :
    vanillaDemonicC N d m0 k μ ↔ vanillaDemonicLaws N d m0 k μ
\end{leansrc}
\par\medskip\noindent{\small \texttt{SuFunctor.Adequacy.PetriVanilla.}\texttt{\bfseries recordVisited\_eq\_face}}\hfill{\footnotesize \path{lean/su-functor/SuFunctor/Adequacy/PetriVanilla.lean}, line 622}\par\nopagebreak
\begin{leansrc}
theorem recordVisited_eq_face :
    recordVisited N d m0 = face N d m0
\end{leansrc}
\par\medskip\noindent{\small \texttt{SuFunctor.Adequacy.PetriVanilla.}\texttt{\bfseries su\_vanilla\_reach\_iff}}\hfill{\footnotesize \path{lean/su-functor/SuFunctor/Adequacy/PetriVanilla.lean}, line 650}\par\nopagebreak
\begin{leansrc}
theorem su_vanilla_reach_iff (hd : ∀ t, 0 < d t) (m : Marking P) :
    (∃ (k : ℕ) (μ : SubPMF (Traj (Marking P) k)),
        idealPoint μ ∈ (vanillaDemonicDnf N d m0 k).modelClass
          ∧ ∃ ω, μ.wt ω ≠ 0 ∧ ∃ i, ω i = m)
      ↔ Reach N m0 m
\end{leansrc}
\par\medskip\noindent{\small \texttt{SuFunctor.Adequacy.PetriVanilla.}\texttt{\bfseries su\_vanilla\_dwell\_uniform}}\hfill{\footnotesize \path{lean/su-functor/SuFunctor/Adequacy/PetriVanilla.lean}, line 671}\par\nopagebreak
\begin{leansrc}
theorem su_vanilla_dwell_uniform {d' : T → ℚ} (hd : ∀ t, 0 < d t)
    (hd' : ∀ t, 0 < d' t) :
    recordVisited N d m0 = recordVisited N d' m0
\end{leansrc}
\par\medskip\noindent{\small \texttt{SuFunctor.Adequacy.PetriVanilla.}\texttt{\bfseries su\_vanilla\_safety}}\hfill{\footnotesize \path{lean/su-functor/SuFunctor/Adequacy/PetriVanilla.lean}, line 681}\par\nopagebreak
\begin{leansrc}
theorem su_vanilla_safety (hd : ∀ t, 0 < d t) (Φ : Marking P → Prop) :
    (∀ m ∈ recordVisited N d m0, Φ m) ↔
      ∀ m, Reach N m0 m → Φ m
\end{leansrc}
\par\medskip\noindent{\small \texttt{SuFunctor.Adequacy.Pgm.}\texttt{\bfseries pgm\_adequacy}}\hfill{\footnotesize \path{lean/su-functor/SuFunctor/Adequacy/Pgm.lean}, line 727}\par\nopagebreak
\begin{leansrc}
theorem pgm_adequacy (B : FinBN n Val) {a : Crd (Asn Val)} :
    ∃ α : Cell (looseConj (atoms (compile B)) : Loose a a)
        (singletonC B.jointLaw) (𝟙 a) (𝟙 a),
      Cell.IsInvertible α
\end{leansrc}
\par\medskip\noindent{\small \texttt{SuFunctor.Adequacy.Pgm.FinBN.}\texttt{\bfseries jointLaw\_mass}}\hfill{\footnotesize \path{lean/su-functor/SuFunctor/Adequacy/Pgm.lean}, line 374}\par\nopagebreak
\begin{leansrc}
theorem FinBN.jointLaw_mass (B : FinBN n Val) : (B.jointLaw).mass = 1
\end{leansrc}
\par\medskip\noindent{\small \texttt{SuFunctor.Adequacy.Pgm.FinBN.}\texttt{\bfseries ancestralPMF\_eq}}\hfill{\footnotesize \path{lean/su-functor/SuFunctor/Adequacy/Pgm.lean}, line 921}\par\nopagebreak
\begin{leansrc}
theorem FinBN.ancestralPMF_eq (σ : Asn Val) :
    (B.ancestralPMF) σ = q2e ((B.jointLaw).wt σ)
\end{leansrc}
\par\medskip\noindent{\small \texttt{SuFunctor.Adequacy.QNet.}\texttt{\bfseries quantum\_net\_adequacy}}\hfill{\footnotesize \path{lean/su-functor/SuFunctor/Adequacy/QuantumNet.lean}, line 426}\par\nopagebreak
\begin{leansrc}
theorem quantum_net_adequacy (P : QNetSystem X n ι T κ) {a : Crd κ} :
    ∃ α : Cell (looseConj (QNet.atoms (QNet.compile P)) : Loose a a)
        (singletonC P.shadowLaw) (𝟙 a) (𝟙 a),
      Cell.IsInvertible α
\end{leansrc}
\par\medskip\noindent{\small \texttt{SuFunctor.Adequacy.}\texttt{\bfseries QNetSystem}}\hfill{\footnotesize \path{lean/su-functor/SuFunctor/Adequacy/QuantumNet.lean}, line 156}\par\nopagebreak
\begin{leansrc}
structure QNetSystem (X : Type*) [PartialOrder X]
    (n : Type*) [Fintype n] [DecidableEq n]
    (ι T κ : Type*) [Fintype κ] where …
\end{leansrc}
\par\medskip\noindent{\small \texttt{SuFunctor.Adequacy.QNetSystem.}\texttt{\bfseries shadowPMF\_prob\_eq\_born}}\hfill{\footnotesize \path{lean/su-functor/SuFunctor/Adequacy/QuantumNet.lean}, line 204}\par\nopagebreak
\begin{leansrc}
theorem shadowPMF_prob_eq_born (α : κ) :
    ((P.shadowPMF.prob α : ℚ) : ℝ) =
      ((P.state * P.frag.records α).trace).re
\end{leansrc}
\par\medskip\noindent{\small \texttt{SuFunctor.Adequacy.QNetSystem.}\texttt{\bfseries shadowPMF\_eq\_fragment\_shadow}}\hfill{\footnotesize \path{lean/su-functor/SuFunctor/Adequacy/QuantumNet.lean}, line 213}\par\nopagebreak
\begin{leansrc}
theorem shadowPMF_eq_fragment_shadow (α : κ) :
    ((P.shadowPMF.prob α : ℚ) : ℝ) =
      P.frag.shadow
        (stateClassMk P.frag.net P.frag.ambient
          ⟨P.state, P.state_isState⟩) α
\end{leansrc}
\par\medskip\noindent{\small \texttt{SuFunctor.Adequacy.}\texttt{\bfseries shadowPMF\_class\_invariant}}\hfill{\footnotesize \path{lean/su-functor/SuFunctor/Adequacy/QuantumNet.lean}, line 227}\par\nopagebreak
\begin{leansrc}
theorem shadowPMF_class_invariant (F : LCQFTFragment X n ι T κ)
    {ρ σ : Matrix n n ℂ} (hρ : IsState ρ) (hσ : IsState σ)
    {q₁ q₂ : κ → ℚ}
    (h₁ : ∀ α, ((q₁ α : ℚ) : ℝ) = recordShadow F.records ρ α)
    (h₂ : ∀ α, ((q₂ α : ℚ) : ℝ) = recordShadow F.records σ α)
    (hcls : stateClassMk F.net F.ambient ⟨ρ, hρ⟩ =
      stateClassMk F.net F.ambient ⟨σ, hσ⟩) :
    (QNetSystem.mk F ρ hρ q₁ h₁).shadowPMF =
      (QNetSystem.mk F σ hσ q₂ h₂).shadowPMF
\end{leansrc}
\par\medskip\noindent{\small \texttt{SuFunctor.Adequacy.QNetSystem.}\texttt{\bfseries shadowLaw\_mass}}\hfill{\footnotesize \path{lean/su-functor/SuFunctor/Adequacy/QuantumNet.lean}, line 198}\par\nopagebreak
\begin{leansrc}
theorem shadowLaw_mass : P.shadowLaw.mass = 1
\end{leansrc}
\par\medskip\noindent{\small \texttt{SuFunctor.Adequacy.QNet.}\texttt{\bfseries rowsLaw\_mass\_of\_complete}}\hfill{\footnotesize \path{lean/su-functor/SuFunctor/Adequacy/QuantumNet.lean}, line 313}\par\nopagebreak
\begin{leansrc}
theorem rowsLaw_mass_of_complete (R : QNetRows κ) (hR : R.WF)
    (hC : R.Complete) : (rowsLaw R hR).mass = 1
\end{leansrc}
\par\medskip\noindent{\small \texttt{SuFunctor.Adequacy.QNet.}\texttt{\bfseries bellQNet\_shadowPMF}}\hfill{\footnotesize \path{lean/su-functor/SuFunctor/Adequacy/QuantumNet.lean}, line 714}\par\nopagebreak
\begin{leansrc}
theorem bellQNet_shadowPMF : bellQNet.shadowPMF = bellShadowPMF
\end{leansrc}
\par\medskip\noindent{\small \texttt{SuFunctor.Adequacy.QNet.}\texttt{\bfseries mixedQNet\_shadowPMF}}\hfill{\footnotesize \path{lean/su-functor/SuFunctor/Adequacy/QuantumNet.lean}, line 718}\par\nopagebreak
\begin{leansrc}
theorem mixedQNet_shadowPMF : mixedQNet.shadowPMF = mixedShadowPMF
\end{leansrc}
\par\medskip\noindent{\small \texttt{SuFunctor.Adequacy.QNet.}\texttt{\bfseries exists\_incomplete\_mass\_lt\_one}}\hfill{\footnotesize \path{lean/su-functor/SuFunctor/Adequacy/QuantumNet.lean}, line 726}\par\nopagebreak
\begin{leansrc}
theorem exists_incomplete_mass_lt_one :
    ∃ (P : QNetSystem (Mink ℝ) (Fin 2 × Fin 2) (Fin 2) ℝ
        (Fin 2 × Fin 2)) (J : Finset (Fin 2 × Fin 2)),
      (QNet.rowsLaw (QNet.compileSub P J)
        (QNet.compileSub_wf P J)).mass < 1
\end{leansrc}
\par\medskip\noindent{\small \texttt{SuFunctor.Adequacy.QNet.}\texttt{\bfseries exists\_nonPSD\_family\_negative\_shadow}}\hfill{\footnotesize \path{lean/su-functor/SuFunctor/Adequacy/QuantumNet.lean}, line 822}\par\nopagebreak
\begin{leansrc}
theorem exists_nonPSD_family_negative_shadow :
    ∃ (R : Fin 2 → Matrix (Fin 1) (Fin 1) ℂ)
      (ρ : Matrix (Fin 1) (Fin 1) ℂ),
      IsState ρ ∧ (∑ α, R α = 1) ∧ ∃ α, recordShadow R ρ α < 0
\end{leansrc}
\par\medskip\noindent{\small \texttt{SuFunctor.Adequacy.QNet.}\texttt{\bfseries su\_qnet\_adequacy}}\hfill{\footnotesize \path{lean/su-functor/SuFunctor/Adequacy/QuantumNet.lean}, line 895}\par\nopagebreak
\begin{leansrc}
theorem su_qnet_adequacy (P : QNetSystem X n ι T κ) (μ : SubPMF κ) :
    idealPoint (μ.map recState)
        ∈ (qnetSpecW (QNet.compile P)).modelClass
      ↔ μ = P.shadowLaw
\end{leansrc}
\par\medskip\noindent{\small \texttt{SuFunctor.Adequacy.QNet.}\texttt{\bfseries quantum\_credal\_adequacy}}\hfill{\footnotesize \path{lean/su-functor/SuFunctor/Adequacy/QuantumNet.lean}, line 546}\par\nopagebreak
\begin{leansrc}
theorem quantum_credal_adequacy (F : LCQFTFragment X n ι T κ)
    (S : Set {ρ : Matrix n n ℂ // IsState ρ}) {a : Crd κ} :
    ∃ α : Cell (QNet.qnetCredalC F S : Loose a a)
        (QNet.quantumLaws F S) (𝟙 a) (𝟙 a),
      Cell.IsInvertible α
\end{leansrc}
\par\medskip\noindent{\small \texttt{SuFunctor.Adequacy.QNet.QNet.}\texttt{\bfseries qnetCredalC}}\hfill{\footnotesize \path{lean/su-functor/SuFunctor/Adequacy/QuantumNet.lean}, line 514}\par\nopagebreak
\begin{leansrc}
def qnetCredalC (F : LCQFTFragment X n ι T κ)
    (S : Set {ρ : Matrix n n ℂ // IsState ρ}) : SubPMF κ → Prop
\end{leansrc}
\par\medskip\noindent{\small \texttt{SuFunctor.Adequacy.QNet.QNet.}\texttt{\bfseries qnetCredalC\_iff}}\hfill{\footnotesize \path{lean/su-functor/SuFunctor/Adequacy/QuantumNet.lean}, line 522}\par\nopagebreak
\begin{leansrc}
theorem qnetCredalC_iff (F : LCQFTFragment X n ι T κ)
    (S : Set {ρ : Matrix n n ℂ // IsState ρ}) (μ : SubPMF κ) :
    qnetCredalC F S μ ↔ quantumLaws F S μ
\end{leansrc}
\par\medskip\noindent{\small \texttt{SuFunctor.Adequacy.QNet.QNet.}\texttt{\bfseries RatShadowMember}}\hfill{\footnotesize \path{lean/su-functor/SuFunctor/Adequacy/QuantumNet.lean}, line 442}\par\nopagebreak
\begin{leansrc}
structure RatShadowMember (F : LCQFTFragment X n ι T κ)
    (S : Set {ρ : Matrix n n ℂ // IsState ρ}) where …
\end{leansrc}
\par\medskip\noindent{\small \texttt{SuFunctor.Adequacy.QNet.}\texttt{\bfseries quantum\_maxSlice}}\hfill{\footnotesize \path{lean/su-functor/SuFunctor/Adequacy/QuantumNet.lean}, line 556}\par\nopagebreak
\begin{leansrc}
theorem quantum_maxSlice (F : LCQFTFragment X n ι T κ)
    (S : Set {ρ : Matrix n n ℂ // IsState ρ}) (ν : SubPMF κ) :
    maxSlice (QNet.quantumCredal F S) ν ↔ QNet.quantumLaws F S ν
\end{leansrc}
\par\medskip\noindent{\small \texttt{SuFunctor.Adequacy.QNet.}\texttt{\bfseries quantum\_ray\_singleton\_reconciliation}}\hfill{\footnotesize \path{lean/su-functor/SuFunctor/Adequacy/QuantumNet.lean}, line 562}\par\nopagebreak
\begin{leansrc}
theorem quantum_ray_singleton_reconciliation
    (F : LCQFTFragment X n ι T κ)
    (S : Set {ρ : Matrix n n ℂ // IsState ρ}) {a : Crd κ} :
    ∃ α : Cell (maxSlice (QNet.quantumCredal F S) : Loose a a)
        (QNet.quantumLaws F S) (𝟙 a) (𝟙 a),
      Cell.IsInvertible α
\end{leansrc}
\par\medskip\noindent{\small \texttt{SuFunctor.Adequacy.QNet.}\texttt{\bfseries satCredal\_lower\_quantumCredal}}\hfill{\footnotesize \path{lean/su-functor/SuFunctor/Adequacy/QuantumNet.lean}, line 573}\par\nopagebreak
\begin{leansrc}
theorem satCredal_lower_quantumCredal (F : LCQFTFragment X n ι T κ)
    (S : Set {ρ : Matrix n n ℂ // IsState ρ})
    (c : SubPMF κ → Prop) :
    SatCredal .lower (QNet.quantumCredal F S) c ↔
      ∀ w : QNet.RatShadowMember F S, c w.law
\end{leansrc}
\par\medskip\noindent{\small \texttt{SuFunctor.Adequacy.QNet.}\texttt{\bfseries satCredal\_upper\_quantumCredal}}\hfill{\footnotesize \path{lean/su-functor/SuFunctor/Adequacy/QuantumNet.lean}, line 582}\par\nopagebreak
\begin{leansrc}
theorem satCredal_upper_quantumCredal (F : LCQFTFragment X n ι T κ)
    (S : Set {ρ : Matrix n n ℂ // IsState ρ})
    {c : SubPMF κ → Prop} (hc : DownClosed c) :
    SatCredal .upper (QNet.quantumCredal F S) c ↔
      ∀ w : QNet.RatShadowMember F S, c w.law
\end{leansrc}
\par\medskip\noindent{\small \texttt{SuFunctor.Adequacy.QNet.QNet.}\texttt{\bfseries quantumLaws\_mix}}\hfill{\footnotesize \path{lean/su-functor/SuFunctor/Adequacy/QuantumNet.lean}, line 660}\par\nopagebreak
\begin{leansrc}
theorem quantumLaws_mix (hS : MixClosed S)
    (w₁ w₂ : RatShadowMember F S) {t : ℚ} (h₀ : 0 ≤ t) (h₁ : t ≤ 1) :
    quantumLaws F S (FinPMF.mix w₁.pmf w₂.pmf t h₀ h₁).toSubPMF
\end{leansrc}
\par\medskip\noindent{\small \texttt{Qtsl.QuantumStates.}\texttt{\bfseries recordShadowPMF\_mix}}\hfill{\footnotesize \path{lean/qtsl-aux/QtslAux/QuantumStates.lean}, line 678}\par\nopagebreak
\begin{leansrc}
theorem recordShadowPMF_mix [Fintype ι] (R : ι → Matrix n n ℂ)
    (ρ σ : Matrix n n ℂ) (q₁ q₂ : ι → ℚ)
    (hq₁ : ∀ α, (q₁ α : ℝ) = recordShadow R ρ α)
    (hq₂ : ∀ α, (q₂ α : ℝ) = recordShadow R σ α)
    (hR : ∀ α, (R α).PosSemidef) (hexh : ∑ α, R α = 1)
    (hρ : IsState ρ) (hσ : IsState σ) {t : ℚ} (h₀ : 0 ≤ t) (h₁ : t ≤ 1) :
    recordShadowPMF R (((t : ℚ) : ℂ) • ρ + (((1 - t : ℚ)) : ℂ) • σ)
        (fun α => t * q₁ α + (1 - t) * q₂ α)
        (fun α => recordShadow_mix_rat_witness hq₁ hq₂ α)
        hR hexh (isState_mix hρ hσ h₀ h₁) =
      FinPMF.mix (recordShadowPMF R ρ q₁ hq₁ hR hexh hρ)
        (recordShadowPMF R σ q₂ hq₂ hR hexh hσ) t h₀ h₁
\end{leansrc}
\par\medskip\noindent{\small \texttt{SuFunctor.Adequacy.QNet.}\texttt{\bfseries mixLaw\_not\_mem\_pair}}\hfill{\footnotesize \path{lean/su-functor/SuFunctor/Adequacy/QuantumNet.lean}, line 775}\par\nopagebreak
\begin{leansrc}
theorem mixLaw_not_mem_pair :
    ¬ QNet.quantumLaws isingQubitFragment bellMixedSet
        (FinPMF.mix bellMember.pmf mixedMember.pmf (1 / 2)
          (by norm_num) (by norm_num)).toSubPMF
\end{leansrc}
\par\medskip\noindent{\small \texttt{SuFunctor.Adequacy.QNet.}\texttt{\bfseries not\_satCredal\_upper\_quantumCredal\_geC}}\hfill{\footnotesize \path{lean/su-functor/SuFunctor/Adequacy/QuantumNet.lean}, line 609}\par\nopagebreak
\begin{leansrc}
theorem not_satCredal_upper_quantumCredal_geC
    (F : LCQFTFragment X n ι T κ)
    (S : Set {ρ : Matrix n n ℂ // IsState ρ})
    [hne : Nonempty (QNet.RatShadowMember F S)]
    {E : Set κ} {q : ℚ} (hq : 0 < q) :
    ¬ SatCredal .upper (QNet.quantumCredal F S) (geC E q)
\end{leansrc}
\par\medskip\noindent{\small \texttt{SuFunctor.Adequacy.QNet.}\texttt{\bfseries satCredal\_quantum\_of\_isEmpty}}\hfill{\footnotesize \path{lean/su-functor/SuFunctor/Adequacy/QuantumNet.lean}, line 625}\par\nopagebreak
\begin{leansrc}
theorem satCredal_quantum_of_isEmpty (F : LCQFTFragment X n ι T κ)
    (S : Set {ρ : Matrix n n ℂ // IsState ρ})
    [h : IsEmpty (QNet.RatShadowMember F S)] (kind : EnvelopeKind)
    (c : SubPMF κ → Prop) :
    SatCredal kind (QNet.quantumCredal F S) c
\end{leansrc}
\par\medskip\noindent{\small \texttt{Qtsl.QuantumStates.}\texttt{\bfseries existsUnique\_isState\_traceRepr}}\hfill{\footnotesize \path{lean/qtsl-aux/QtslAux/QuantumStates.lean}, line 518}\par\nopagebreak
\begin{leansrc}
theorem existsUnique_isState_traceRepr
    (hpos : ∀ v : n → ℂ, 0 ≤ φ (Matrix.vecMulVec v (star v)))
    (hone : φ 1 = 1) :
    ∃! ρ : Matrix n n ℂ, IsState ρ ∧ ∀ a, φ a = (ρ * a).trace
\end{leansrc}
\par\medskip\noindent{\small \texttt{Qtsl.QuantumStates.}\texttt{\bfseries exists\_shadow\_lowerE\_lt\_upperE}}\hfill{\footnotesize \path{lean/qtsl-aux/QtslAux/QuantumStates.lean}, line 953}\par\nopagebreak
\begin{leansrc}
theorem exists_shadow_lowerE_lt_upperE :
    ∃ f : Fin 2 × Fin 2 → ℚ,
      lowerE ![bellShadowPMF, mixedShadowPMF] f <
        upperE ![bellShadowPMF, mixedShadowPMF] f
\end{leansrc}
\par\medskip\noindent{\small \texttt{Qtsl.QuantumCoarseGraining.}\texttt{\bfseries pushforward\_ofPresentation}}\hfill{\footnotesize \path{lean/qtsl-aux/QtslAux/QuantumCoarseGraining.lean}, line 474}\par\nopagebreak
\begin{leansrc}
theorem pushforward_ofPresentation (f : ι → κ) (C : ι → Matrix d d ℂ)
    (ψ : d → ℂ)
    (hnorm : ∑ α, ∑ β, decoherenceFunctional C ψ α β = 1) :
    (QuantumMeasure.ofPresentation C ψ hnorm).pushforward f
      = QuantumMeasure.ofPresentation (coarsenFamily f C) ψ
          (by
            simp only [decoherenceFunctional_coarsenFamily]
            have h := sum_sum_fiberEvent f Finset.univ Finset.univ
              fun a b => decoherenceFunctional C ψ a b
            rw [h, fiberEvent_univ]
            exact hnorm)
\end{leansrc}
\par\medskip\noindent{\small \texttt{SuFunctor.Adequacy.}\texttt{\bfseries tvDistQ}}\hfill{\footnotesize \path{lean/su-functor/SuFunctor/Adequacy/QuantumContainment.lean}, line 135}\par\nopagebreak
\begin{leansrc}
def tvDistQ {κ : Type u} [Fintype κ] (μ ν : SubPMF κ) : ℚ
\end{leansrc}
\par\medskip\noindent{\small \texttt{SuFunctor.Adequacy.}\texttt{\bfseries tvDistQ\_eq\_zero\_iff}}\hfill{\footnotesize \path{lean/su-functor/SuFunctor/Adequacy/QuantumContainment.lean}, line 166}\par\nopagebreak
\begin{leansrc}
theorem tvDistQ_eq_zero_iff {κ : Type u} [Fintype κ] {μ ν : SubPMF κ} :
    tvDistQ μ ν = 0 ↔ μ = ν
\end{leansrc}
\par\medskip\noindent{\small \texttt{SuFunctor.Adequacy.}\texttt{\bfseries credalInflate}}\hfill{\footnotesize \path{lean/su-functor/SuFunctor/Adequacy/QuantumContainment.lean}, line 214}\par\nopagebreak
\begin{leansrc}
def credalInflate {κ : Type u} [Fintype κ] (ε : ℚ)
    (L : SubPMF κ → Prop) : SubPMF κ → Prop
\end{leansrc}
\par\medskip\noindent{\small \texttt{SuFunctor.Adequacy.Containment.}\texttt{\bfseries mono}}\hfill{\footnotesize \path{lean/su-functor/SuFunctor/Adequacy/QuantumContainment.lean}, line 282}\par\nopagebreak
\begin{leansrc}
theorem Containment.mono {ι κ : Type u} [Fintype ι] [Fintype κ]
    {LUV : SubPMF ι → Prop} {f : ι → κ} {LIR : SubPMF κ → Prop}
    {ε ε' : ℚ} (h : ε ≤ ε')
    (hc : Containment LUV f LIR ε) : Containment LUV f LIR ε'
\end{leansrc}
\par\medskip\noindent{\small \texttt{SuFunctor.Adequacy.Containment.}\texttt{\bfseries comp}}\hfill{\footnotesize \path{lean/su-functor/SuFunctor/Adequacy/QuantumContainment.lean}, line 292}\par\nopagebreak
\begin{leansrc}
theorem Containment.comp {ι κ σ : Type u}
    [Fintype ι] [Fintype κ] [Fintype σ]
    {L₁ : SubPMF ι → Prop} {f : ι → κ} {L₂ : SubPMF κ → Prop} {ε₁ : ℚ}
    {g : κ → σ} {L₃ : SubPMF σ → Prop} {ε₂ : ℚ}
    (h₁ : Containment L₁ f L₂ ε₁) (h₂ : Containment L₂ g L₃ ε₂) :
    Containment L₁ (g ∘ f) L₃ (ε₁ + ε₂)
\end{leansrc}
\par\medskip\noindent{\small \texttt{SuFunctor.SubPMF.}\texttt{\bfseries map\_wt\_apply}}\hfill{\footnotesize \path{lean/su-functor/SuFunctor/Adequacy/QuantumContainment.lean}, line 102}\par\nopagebreak
\begin{leansrc}
theorem map_wt_apply {X Y : Type u} [Fintype X] [DecidableEq Y]
    (f : X → Y) (μ : SubPMF X) (c : Y) :
    (μ.map f).wt c
      = ∑ a ∈ Finset.univ.filter (fun a => f a = c), μ.wt a
\end{leansrc}
\par\medskip\noindent{\small \texttt{SuFunctor.Adequacy.}\texttt{\bfseries tvDistQ\_map\_le}}\hfill{\footnotesize \path{lean/su-functor/SuFunctor/Adequacy/QuantumContainment.lean}, line 193}\par\nopagebreak
\begin{leansrc}
theorem tvDistQ_map_le {X Y : Type u} [Fintype X] [Fintype Y]
    (f : X → Y) (μ ν : SubPMF X) :
    tvDistQ (μ.map f) (ν.map f) ≤ tvDistQ μ ν
\end{leansrc}
\par\medskip\noindent{\small \texttt{SuFunctor.Adequacy.}\texttt{\bfseries recordShadow\_coarsenFamily}}\hfill{\footnotesize \path{lean/su-functor/SuFunctor/Adequacy/QuantumContainment.lean}, line 327}\par\nopagebreak
\begin{leansrc}
theorem recordShadow_coarsenFamily {ι' κ' : Type*} [Fintype ι']
    [Fintype κ'] [DecidableEq ι'] [DecidableEq κ']
    {d : Type*} [Fintype d]
    (f : ι' → κ') (R : ι' → Matrix d d ℂ) (ρ : Matrix d d ℂ) (c : κ') :
    recordShadow (coarsenFamily f R) ρ c
      = ∑ a ∈ fiber f c, recordShadow R ρ a
\end{leansrc}
\par\medskip\noindent{\small \texttt{SuFunctor.Adequacy.}\texttt{\bfseries uv\_ir\_containment}}\hfill{\footnotesize \path{lean/su-functor/SuFunctor/Adequacy/QuantumContainment.lean}, line 763}\par\nopagebreak
\begin{leansrc}
theorem uv_ir_containment :
    Containment (QNet.quantumLaws isingQubitFragment uvSeg) Prod.fst
      (QNet.quantumLaws irFragment irSet) (49 / 625)
\end{leansrc}
\par\medskip\noindent{\small \texttt{SuFunctor.Adequacy.}\texttt{\bfseries uv\_ir\_containment\_lower}}\hfill{\footnotesize \path{lean/su-functor/SuFunctor/Adequacy/QuantumContainment.lean}, line 788}\par\nopagebreak
\begin{leansrc}
theorem uv_ir_containment_lower (ε : ℚ)
    (h : Containment (QNet.quantumLaws isingQubitFragment uvSeg)
      Prod.fst (QNet.quantumLaws irFragment irSet) ε) :
    49 / 625 ≤ ε
\end{leansrc}
\par\medskip\noindent{\small \texttt{SuFunctor.Adequacy.}\texttt{\bfseries not\_uv\_ir\_containment\_zero}}\hfill{\footnotesize \path{lean/su-functor/SuFunctor/Adequacy/QuantumContainment.lean}, line 817}\par\nopagebreak
\begin{leansrc}
theorem not_uv_ir_containment_zero :
    ¬ Containment (QNet.quantumLaws isingQubitFragment uvSeg) Prod.fst
      (QNet.quantumLaws irFragment irSet) 0
\end{leansrc}
\par\medskip\noindent{\small \texttt{SuFunctor.Adequacy.}\texttt{\bfseries onestep\_defect}}\hfill{\footnotesize \path{lean/su-functor/SuFunctor/Adequacy/QuantumContainment.lean}, line 1056}\par\nopagebreak
\begin{leansrc}
theorem onestep_defect (x : Fin 2) :
    tvDistQ (coarseLaw (x + 1)) (irStep (coarseLaw x)) ≤ 49 / 625
\end{leansrc}
\par\medskip\noindent{\small \texttt{SuFunctor.Adequacy.}\texttt{\bfseries horizon\_bound}}\hfill{\footnotesize \path{lean/su-functor/SuFunctor/Adequacy/QuantumContainment.lean}, line 1073}\par\nopagebreak
\begin{leansrc}
theorem horizon_bound (x : Fin 2) (k : ℕ) :
    tvDistQ (coarseLaw ((· + 1)^[k] x)) (irStep^[k] (coarseLaw x))
      ≤ (k : ℚ) * (49 / 625)
\end{leansrc}
\par\medskip\noindent{\small \texttt{SuFunctor.Adequacy.}\texttt{\bfseries lax\_horizon}}\hfill{\footnotesize \path{lean/su-functor/SuFunctor/Adequacy/QuantumContainment.lean}, line 395}\par\nopagebreak
\begin{leansrc}
theorem lax_horizon {A : Type*} {B : Type*}
    (push : A → B) (stepA : A → A) (stepB : B → B)
    (d : B → B → ℚ) (ε : ℚ) (Reach : A → Prop)
    (hself : ∀ x, d x x = 0)
    (htri : ∀ x y z, d x z ≤ d x y + d y z)
    (hnonexp : ∀ x y, d (stepB x) (stepB y) ≤ d x y)
    (hinv : ∀ a, Reach a → Reach (stepA a))
    (hdef : ∀ a, Reach a → d (push (stepA a)) (stepB (push a)) ≤ ε) :
    ∀ a, Reach a → ∀ k : ℕ,
      d (push (stepA^[k] a)) (stepB^[k] (push a)) ≤ (k : ℚ) * ε
\end{leansrc}
\par\medskip\noindent{\small \texttt{SuFunctor.Adequacy.}\texttt{\bfseries tvDistQ\_stochAct\_le}}\hfill{\footnotesize \path{lean/su-functor/SuFunctor/Adequacy/QuantumContainment.lean}, line 368}\par\nopagebreak
\begin{leansrc}
theorem tvDistQ_stochAct_le {κ : Type u} [Fintype κ] (Q : κ → κ → ℚ)
    (hQ0 : ∀ c c', 0 ≤ Q c c') (hQ1 : ∀ c, ∑ c', Q c c' = 1)
    (μ ν : SubPMF κ) :
    tvDistQ (stochAct Q hQ0 hQ1 μ) (stochAct Q hQ0 hQ1 ν)
      ≤ tvDistQ μ ν
\end{leansrc}
\par\medskip\noindent{\small \texttt{SuFunctor.Adequacy.}\texttt{\bfseries ir\_tot\_containment}}\hfill{\footnotesize \path{lean/su-functor/SuFunctor/Adequacy/QuantumContainment.lean}, line 1135}\par\nopagebreak
\begin{leansrc}
theorem ir_tot_containment :
    Containment (QNet.quantumLaws irFragment irSet) totalMerge
      (QNet.quantumLaws totFragment totSet) 0
\end{leansrc}
\par\medskip\noindent{\small \texttt{SuFunctor.Adequacy.}\texttt{\bfseries uv\_tot\_containment}}\hfill{\footnotesize \path{lean/su-functor/SuFunctor/Adequacy/QuantumContainment.lean}, line 1159}\par\nopagebreak
\begin{leansrc}
theorem uv_tot_containment :
    Containment (QNet.quantumLaws isingQubitFragment uvSeg)
      (totalMerge ∘ Prod.fst)
      (QNet.quantumLaws totFragment totSet) (49 / 625 + 0)
\end{leansrc}
\par\medskip\noindent{\small \texttt{SuFunctor.Adequacy.}\texttt{\bfseries uvSeg\_mixClosed}}\hfill{\footnotesize \path{lean/su-functor/SuFunctor/Adequacy/QuantumContainment.lean}, line 570}\par\nopagebreak
\begin{leansrc}
theorem uvSeg_mixClosed : QNet.MixClosed uvSeg
\end{leansrc}
\par\medskip\noindent{\small \texttt{SuFunctor.Adequacy.}\texttt{\bfseries uvLaws\_mix\_exercised}}\hfill{\footnotesize \path{lean/su-functor/SuFunctor/Adequacy/QuantumContainment.lean}, line 620}\par\nopagebreak
\begin{leansrc}
theorem uvLaws_mix_exercised :
    QNet.quantumLaws isingQubitFragment uvSeg
      ((FinPMF.mix upMember.pmf entMember.pmf (1 / 2)
        (by norm_num) (by norm_num)).toSubPMF)
\end{leansrc}
\par\medskip\noindent{\small \texttt{SuFunctor.Adequacy.}\texttt{\bfseries pushforward\_probRow\_bridge}}\hfill{\footnotesize \path{lean/su-functor/SuFunctor/Adequacy/QuantumContainment.lean}, line 1302}\par\nopagebreak
\begin{leansrc}
theorem pushforward_probRow_bridge (c : Fin 2) :
    (uvMeasure.pushforward Prod.fst).probRow c
      = (((entMember.law.map Prod.fst).wt c : ℚ) : ℝ)
\end{leansrc}
\par\medskip\noindent{\small \texttt{SuFunctor.Adequacy.}\texttt{\bfseries eftTower\_offDiagMass\_antitone}}\hfill{\footnotesize \path{lean/su-functor/SuFunctor/Adequacy/QuantumContainment.lean}, line 1276}\par\nopagebreak
\begin{leansrc}
theorem eftTower_offDiagMass_antitone :
    Antitone fun n => (eftStage n).offDiagMass
\end{leansrc}
\par\medskip\noindent{\small \texttt{SuFunctor.Adequacy.}\texttt{\bfseries eftStage\_classical}}\hfill{\footnotesize \path{lean/su-functor/SuFunctor/Adequacy/QuantumContainment.lean}, line 1268}\par\nopagebreak
\begin{leansrc}
theorem eftStage_classical (n : ℕ) :
    ∃ (p : eftTowerX n → ℝ) (hp : ∀ c, 0 ≤ p c) (hs : ∑ c, p c = 1),
      eftStage n = QuantumMeasure.classical p hp hs
\end{leansrc}
\par\medskip\noindent{\small \texttt{SuFunctor.Adequacy.}\texttt{\bfseries recordShadow\_kchanApply}}\hfill{\footnotesize \path{lean/su-functor/SuFunctor/Adequacy/QuantumCompression.lean}, line 317}\par\nopagebreak
\begin{leansrc}
theorem recordShadow_kchanApply (A : J → Matrix e d ℂ)
    (S : κ' → Matrix e e ℂ) (ρ : Matrix d d ℂ) (c : κ') :
    recordShadow S (kchanApply A ρ) c
      = recordShadow (fun c => kchanDual A (S c)) ρ c
\end{leansrc}
\par\medskip\noindent{\small \texttt{SuFunctor.Adequacy.}\texttt{\bfseries kchanApply\_ptrKraus\_eq\_ptR}}\hfill{\footnotesize \path{lean/su-functor/SuFunctor/Adequacy/QuantumCompression.lean}, line 450}\par\nopagebreak
\begin{leansrc}
theorem kchanApply_ptrKraus_eq_ptR (ρ : TwoQubits) :
    kchanApply ptrKraus ρ = ptR ρ
\end{leansrc}
\par\medskip\noindent{\small \texttt{SuFunctor.Adequacy.}\texttt{\bfseries uv\_smallIR\_containment}}\hfill{\footnotesize \path{lean/su-functor/SuFunctor/Adequacy/QuantumCompression.lean}, line 718}\par\nopagebreak
\begin{leansrc}
theorem uv_smallIR_containment :
    Containment (QNet.quantumLaws isingQubitFragment uvSeg) Prod.fst
      (QNet.quantumLaws singleQubitFragment smallSet) (49 / 625)
\end{leansrc}
\par\medskip\noindent{\small \texttt{SuFunctor.Adequacy.QNet.}\texttt{\bfseries QNetRows}}\hfill{\footnotesize \path{lean/su-functor/SuFunctor/Adequacy/QuantumNet.lean}, line 251}\par\nopagebreak
\begin{leansrc}
structure QNetRows (κ : Type*) where …
\end{leansrc}
\par\medskip\noindent Modules cited: \texttt{SuFunctor.Adequacy.QuantumNet}, \texttt{SuFunctor.Adequacy.Ode}, \texttt{SuFunctor.Adequacy.Pde}, \texttt{SuFunctor.Adequacy.PdeRows}, \texttt{QtslAux.SdeGenerator}, \texttt{QtslAux.SdeGeneratorPoly}, \texttt{QtslAux.SdeMartingaleProblem}, \texttt{SuFunctor.Adequacy.Rode}, \texttt{QtslAux.JumpDiffusion}, \texttt{QtslAux.JumpDiffusionMartingaleProblem}, \texttt{SuFunctor.Adequacy.CtmcJumpChain}, \texttt{SuFunctor.Adequacy.MarkovAutomaton}, \texttt{SuFunctor.Adequacy.MaDemonic}, \texttt{SuFunctor.Adequacy.Bag}, \texttt{SuFunctor.TokenGame}, \texttt{SuFunctor.Adequacy.StochasticPetri}, \texttt{SuFunctor.Adequacy.StochasticPetriCascade}, \texttt{SuFunctor.Adequacy.StochasticPetriDemonic}, \texttt{SuFunctor.Adequacy.StochasticPetriSetFiring}, \texttt{SuFunctor.Adequacy.StochasticPetriFlatten}, \texttt{SuFunctor.Adequacy.StochasticPetriFlattenAgreement}, \texttt{SuFunctor.Adequacy.StochasticPetriComposed}, \texttt{QtslAux.KilledTimeIndexed}, \texttt{SuFunctor.Adequacy.StochasticPetriAngelic}, \texttt{SuFunctor.Adequacy.StochasticPetriColoured}, \texttt{SuFunctor.Adequacy.StochasticPetriColouredCascade}, \texttt{SuFunctor.Adequacy.PetriVanilla}, \texttt{Qtsl.ScmTwin}, \texttt{SuFunctor.Stochastic.NoiseKleisli}, \texttt{QtslAux.QuantumCoarseGraining}, \texttt{SuFunctor.Adequacy.QuantumContainment}.
\par\medskip\noindent Also cited in this chapter, listed under an earlier one: \texttt{QtslAux.QuantumAssembly}, \texttt{QtslAux.QuantumRecords}, \texttt{SuFunctor.Adequacy.DelayLine}, \texttt{SuFunctor.Adequacy.Dtmc}, \texttt{SuFunctor.Adequacy.EnvelopeUnion}, \texttt{SuFunctor.Adequacy.HeatTorus}, \texttt{SuFunctor.Adequacy.HeatTorusInit}, \texttt{SuFunctor.Adequacy.Inclusions}, \texttt{SuFunctor.Adequacy.MarginalGlue}, \texttt{SuFunctor.Adequacy.Mdp}, \texttt{SuFunctor.Adequacy.MdpDemonic}, \texttt{SuFunctor.Adequacy.Pgm}, \texttt{SuFunctor.Adequacy.SolutionSetMesh}, \texttt{copy\_ne\_kleisli}, \texttt{maxSlice\_unionRays}, \texttt{maxSlice\_unionRays\_subset}, \texttt{noiseDenot\_seqComp}.

\section{Systems: Composition Patterns}
Citations from Chapter~\ref{ch:qtsl-composition-patterns}.

\par\medskip\noindent{\small \texttt{SuFunctor.AlignedLoose.}\texttt{\bfseries seq\_assoc}}\hfill{\footnotesize \path{lean/su-functor/SuFunctor/HypernetComposition.lean}, line 219}\par\nopagebreak
\begin{leansrc}
theorem AlignedLoose.seq_assoc (m : AlignedLoose X Y) (n : AlignedLoose Y Z)
    (p : AlignedLoose Z V) : (m.seq n).seq p = m.seq (n.seq p)
\end{leansrc}
\par\medskip\noindent{\small \texttt{SuFunctor.}\texttt{\bfseries austereDnf\_subst}}\hfill{\footnotesize \path{lean/su-functor/SuFunctor/HypernetComposition.lean}, line 350}\par\nopagebreak
\begin{leansrc}
theorem austereDnf_subst (m : AlignedLoose X Y) (n : AlignedLoose Y Z)
    {T₁ T₁' T₂ T₂' : List PrimApp}
    (h₁ : austereDnf ty m T₁ = austereDnf ty m T₁')
    (h₂ : austereDnf ty n T₂ = austereDnf ty n T₂')
    (hTC₁ : ∀ a ∈ T₁, a.TypeConsistent (tableOn ty m.apex))
    (hw₁ : ∀ a ∈ T₁, ∀ w ∈ a.inputWires ++ a.outputWires, w ∈ m.apex)
    (hTC₂ : ∀ a ∈ T₂, a.TypeConsistent (tableOn ty n.apex))
    (hw₂ : ∀ a ∈ T₂, ∀ w ∈ a.inputWires ++ a.outputWires, w ∈ n.apex)
    (hTC₁' : ∀ a ∈ T₁', a.TypeConsistent (tableOn ty m.apex))
    (hw₁' : ∀ a ∈ T₁', ∀ w ∈ a.inputWires ++ a.outputWires, w ∈ m.apex)
    (hTC₂' : ∀ a ∈ T₂', a.TypeConsistent (tableOn ty n.apex))
    (hw₂' : ∀ a ∈ T₂', ∀ w ∈ a.inputWires ++ a.outputWires, w ∈ n.apex) :
    austereDnf ty (m.seq n) (T₁ ++ T₂)
      = austereDnf ty (m.seq n) (T₁' ++ T₂')
\end{leansrc}
\par\medskip\noindent{\small \texttt{SuFunctor.}\texttt{\bfseries admRegion\_append\_empty\_left}}\hfill{\footnotesize \path{lean/su-functor/SuFunctor/HypernetComposition.lean}, line 373}\par\nopagebreak
\begin{leansrc}
theorem admRegion_append_empty_left (m : AlignedLoose X Y)
    (n : AlignedLoose Y Z) {T₁ T₂ : List PrimApp}
    (h : AdmRegion ty m T₁ = ∅)
    (hTC₁ : ∀ a ∈ T₁, a.TypeConsistent (tableOn ty m.apex))
    (hw₁ : ∀ a ∈ T₁, ∀ w ∈ a.inputWires ++ a.outputWires, w ∈ m.apex)
    (hTC₂ : ∀ a ∈ T₂, a.TypeConsistent (tableOn ty n.apex))
    (hw₂ : ∀ a ∈ T₂, ∀ w ∈ a.inputWires ++ a.outputWires, w ∈ n.apex) :
    AdmRegion ty (m.seq n) (T₁ ++ T₂) = ∅
\end{leansrc}
\par\medskip\noindent{\small \texttt{SuFunctor.}\texttt{\bfseries share\_contradiction}}\hfill{\footnotesize \path{lean/su-functor/SuFunctor/HypernetComposition.lean}, line 465}\par\nopagebreak
\begin{leansrc}
theorem share_contradiction {w : Wire} {τ : TensorType}
    (c₁ c₂ : TensorValue τ) {i₀ : ℕ}
    (hne : c₁.entryE i₀ ≠ c₂.entryE i₀) :
    AdmRegion ty
        ((AlignedLoose.unit {w}).seq (AlignedLoose.unit {w}))
        ([PrimApp.const τ c₁ w] ++ [PrimApp.const τ c₂ w]) = ∅
\end{leansrc}
\par\medskip\noindent{\small \texttt{SuFunctor.}\texttt{\bfseries theta\_rename\_seq}}\hfill{\footnotesize \path{lean/su-functor/SuFunctor/HypernetComposition.lean}, line 564}\par\nopagebreak
\begin{leansrc}
theorem theta_rename_seq {m : AlignedLoose X Y} {n : AlignedLoose Y Z}
    {m' : AlignedLoose X' Y'} {n' : AlignedLoose Y' Z'}
    (c₁ : AlignedCell ty m m') (c₂ : AlignedCell ty n n')
    (hmap : c₁.map = c₂.map) (T₁ T₂ : List PrimApp)
    (hTC₁ : ∀ a ∈ T₁, a.TypeConsistent (tableOn ty m.apex))
    (hw₁ : ∀ a ∈ T₁, ∀ w ∈ a.inputWires ++ a.outputWires, w ∈ m.apex)
    (hTC₂ : ∀ a ∈ T₂, a.TypeConsistent (tableOn ty n.apex))
    (hw₂ : ∀ a ∈ T₂, ∀ w ∈ a.inputWires ++ a.outputWires, w ∈ n.apex)
    (hTC₁' : ∀ a ∈ T₁.map (PrimApp.rename c₁.map),
      a.TypeConsistent (tableOn ty m'.apex))
    (hw₁' : ∀ a ∈ T₁.map (PrimApp.rename c₁.map),
      ∀ w ∈ a.inputWires ++ a.outputWires, w ∈ m'.apex)
    (hTC₂' : ∀ a ∈ T₂.map (PrimApp.rename c₂.map),
      a.TypeConsistent (tableOn ty n'.apex))
    (hw₂' : ∀ a ∈ T₂.map (PrimApp.rename c₂.map),
      ∀ w ∈ a.inputWires ++ a.outputWires, w ∈ n'.apex) :
    dnfPull (renameTot ty (AlignedCell.seqCell ty c₁ c₂ hmap))
        (austereDnf ty (m.seq n) (T₁ ++ T₂))
      = dnfPull (splitTot ty m' n')
          (dnfBoxplus
            (dnfPull (renameTot ty c₁) (austereDnf ty m T₁))
            (dnfPull (renameTot ty c₂) (austereDnf ty n T₂)))
\end{leansrc}
\par\medskip\noindent{\small \texttt{SuFunctor.}\texttt{\bfseries admRegion\_append}}\hfill{\footnotesize \path{lean/su-functor/SuFunctor/HypernetInclusion.lean}, line 728}\par\nopagebreak
\begin{leansrc}
theorem admRegion_append (m : AlignedLoose X Y) (n : AlignedLoose Y Z)
    (T₁ T₂ : List PrimApp)
    (hTC₁ : ∀ a ∈ T₁, a.TypeConsistent (tableOn ty m.apex))
    (hw₁ : ∀ a ∈ T₁, ∀ w ∈ a.inputWires ++ a.outputWires, w ∈ m.apex)
    (hTC₂ : ∀ a ∈ T₂, a.TypeConsistent (tableOn ty n.apex))
    (hw₂ : ∀ a ∈ T₂, ∀ w ∈ a.inputWires ++ a.outputWires, w ∈ n.apex) :
    AdmRegion ty (m.seq n) (T₁ ++ T₂)
      = splitTot ty m n ⁻¹'
          (Pat.compFst (alignedPat ty m) (alignedPat ty n) ⁻¹'
              AdmRegion ty m T₁
            ∩ Pat.compSnd (alignedPat ty m) (alignedPat ty n) ⁻¹'
              AdmRegion ty n T₂)
\end{leansrc}
\par\medskip\noindent{\small \texttt{SuFunctor.}\texttt{\bfseries share\_feasible}}\hfill{\footnotesize \path{lean/su-functor/SuFunctor/HypernetComposition.lean}, line 442}\par\nopagebreak
\begin{leansrc}
theorem share_feasible {w : Wire} {τ : TensorType} (hty : ty w = some τ)
    (c : TensorValue τ) :
    (AdmRegion ty (AlignedLoose.unit {w})
      [PrimApp.const τ c w]).Nonempty
\end{leansrc}
\par\medskip\noindent{\small \texttt{SuFunctor.Spacetime.}\texttt{\bfseries hybridTime\_square\_subset}}\hfill{\footnotesize \path{lean/su-functor/SuFunctor/Adequacy/HybridTime.lean}, line 489}\par\nopagebreak
\begin{leansrc}
theorem hybridTime_square_subset
    (D : SpecDnf (Joint (W → ClockedEl Clk (SiteRec κ X)))) :
    ((D.withWall (sampleWireWall κ X ports i i' φ)).withWall
        (possWall (perTickRowsSet κ X ports i i' φ Pc))).modelClass ⊆
      ((D.withWall (sampleWireWall κ X ports i i' φ)).withWall
        (possWall (coarseRowsSet κ X ports i Pc))).modelClass
\end{leansrc}
\par\medskip\noindent{\small \texttt{SuFunctor.Spacetime.}\texttt{\bfseries tickSeam\_hybridTime\_square\_eq}}\hfill{\footnotesize \path{lean/su-functor/SuFunctor/Adequacy/HybridTime.lean}, line 552}\par\nopagebreak
\begin{leansrc}
theorem tickSeam_hybridTime_square_eq {X : Type u}
    (Pc : SpacetimeRec (tickClkOf.{u} ⟨false⟩) X → Prop)
    (D : SpecDnf (Joint (ULift.{u} Bool →
      ClockedEl (ULift.{u} Bool) (SiteRec tickClkOf X)))) :
    ((D.withWall (sampleWireWall tickClkOf X tickPorts ⟨false⟩ ⟨true⟩
          tickSeamHom)).withWall
        (possWall (perTickRowsSet tickClkOf X tickPorts ⟨false⟩ ⟨true⟩
          tickSeamHom Pc))).modelClass =
      ((D.withWall (sampleWireWall tickClkOf X tickPorts ⟨false⟩ ⟨true⟩
          tickSeamHom)).withWall
        (possWall (coarseRowsSet tickClkOf X tickPorts ⟨false⟩
          Pc))).modelClass
\end{leansrc}
\par\medskip\noindent{\small \texttt{SuFunctor.Plain.}\texttt{\bfseries funnelComp\_breaks\_unit\_coherence}}\hfill{\footnotesize \path{lean/su-functor/SuFunctor/Plain/Battery.lean}, line 177}\par\nopagebreak
\begin{leansrc}
theorem funnelComp_breaks_unit_coherence (M : Pat A B) :
    ¬ ∀ 𝒰 : Cpl M,
      Pen.pull (PatCell.tot
          (PseudoDoubleCategoryStruct.lLeftUnitor (C := WireSpan El) M)) ''
          funnelComp ({unitPen A} : Cpl (Pat.id A)) 𝒰 = 𝒰
\end{leansrc}
\par\medskip\noindent{\small \texttt{SuFunctor.Plain.}\texttt{\bfseries unitClass\_transparent}}\hfill{\footnotesize \path{lean/su-functor/SuFunctor/Plain/Battery.lean}, line 86}\par\nopagebreak
\begin{leansrc}
theorem unitClass_transparent (M : Pat A B) (𝒰 : Cpl M) :
    Pen.pull (PatCell.tot
        (PseudoDoubleCategoryStruct.lLeftUnitor (C := WireSpan El) M)) ''
        Set.image2 boxplus ({unitPen A} : Cpl (Pat.id A)) 𝒰 = 𝒰
\end{leansrc}
\par\medskip\noindent{\small \texttt{SuFunctor.Adequacy.}\texttt{\bfseries funnel\_coupledSat\_of\_parts}}\hfill{\footnotesize \path{lean/su-functor/SuFunctor/Adequacy/Funnel.lean}, line 195}\par\nopagebreak
\begin{leansrc}
theorem funnel_coupledSat_of_parts (F : FunnelData K) {σ₁ σ₂ : K → ℝ≥0∞}
    (h₁ : σ₁ ∈ F.transit.modelClass) (h₂ : σ₂ ∈ F.hold.modelClass)
    (hcon : ∀ w ∈ F.composite.marks, Dip (σ₁ + σ₂) w) :
    σ₁ + σ₂ ∈ F.composite.modelClass
\end{leansrc}
\par\medskip\noindent{\small \texttt{SuFunctor.Adequacy.FunnelParts.}\texttt{\bfseries assemble\_isSome\_iff}}\hfill{\footnotesize \path{lean/su-functor/SuFunctor/Adequacy/Funnel.lean}, line 138}\par\nopagebreak
\begin{leansrc}
theorem assemble_isSome_iff {p : FunnelParts K} :
    p.assemble.isSome ↔ p.transitTarget = p.holdInvariant ∧ 0 < p.budget
\end{leansrc}
\par\medskip\noindent{\small \texttt{SuFunctor.}\texttt{\bfseries vorobev\_record\_glue}}\hfill{\footnotesize \path{lean/su-functor/SuFunctor/Adequacy/MarginalGlue.lean}, line 274}\par\nopagebreak
\begin{leansrc}
theorem vorobev_record_glue (v : Bool × Bool × Bool) :
    (0 < pairCorr.wt (proj12 v) ∧ 0 < pairCorr.wt (proj23 v) ∧
        0 < pairIndep.wt (proj13 v)) ↔
      v = (false, false, false) ∨ v = (true, true, true)
\end{leansrc}
\par\medskip\noindent Modules cited: \texttt{SuFunctor.HypernetComposition}, \texttt{SuFunctor.Adequacy.HybridTime}, \texttt{SuFunctor.Adequacy.Funnel}, \texttt{SuFunctor.Adequacy.StochasticPetriSupervise}.
\par\medskip\noindent Also cited in this chapter, listed under an earlier one: \texttt{Qtsl.ScmTwin}, \texttt{SuFunctor.Adequacy.MarginalGlue}, \texttt{SuFunctor.HypernetInclusion}, \texttt{SuFunctor.Plain.Battery}, \texttt{SuFunctor.SingleSiteSemanticUniverse}, \texttt{coupledSat\_comp\_add\_of\_dips}, \texttt{credalCompile\_vorobevGlue}, \texttt{glue\_admissible\_iff}, \texttt{theta\_comp}, \texttt{vorobev\_no\_joint}.

\section{Specification Language}
Citations from Chapter~\ref{ch:qtsl-specification-language}.

\par\medskip\noindent{\small \texttt{SuFunctor.Universe.}\texttt{\bfseries Sat}}\hfill{\footnotesize \path{lean/su-functor/SuFunctor/Universe/Judgment.lean}, line 54}\par\nopagebreak
\begin{leansrc}
def Sat (σ φ : Penalty V) : Prop
\end{leansrc}
\par\medskip\noindent{\small \texttt{SuFunctor.Universe.}\texttt{\bfseries wallOn\_le\_iff}}\hfill{\footnotesize \path{lean/su-functor/SuFunctor/Universe/SpecElab.lean}, line 120}\par\nopagebreak
\begin{leansrc}
theorem wallOn_le_iff {C : Set K} {σ : K → ℝ≥0∞} :
    wallOn C ≤ σ ↔ permitted σ ⊆ C
\end{leansrc}
\par\medskip\noindent{\small \texttt{SuFunctor.Universe.}\texttt{\bfseries dip\_top\_iff}}\hfill{\footnotesize \path{lean/su-functor/SuFunctor/Universe/SpecElab.lean}, line 137}\par\nopagebreak
\begin{leansrc}
theorem dip_top_iff {σ : K → ℝ≥0∞} :
    Dip σ (⊤ : K → ℝ≥0∞) ↔ (permitted σ).Nonempty
\end{leansrc}
\par\medskip\noindent{\small \texttt{SuFunctor.}\texttt{\bfseries SatCredal}}\hfill{\footnotesize \path{lean/su-functor/SuFunctor/Adequacy/Envelope.lean}, line 289}\par\nopagebreak
\begin{leansrc}
def SatCredal (k : EnvelopeKind) (C c : SubPMF V → Prop) : Prop
\end{leansrc}
\par\medskip\noindent{\small \texttt{SuFunctor.Universe.SpecRecord.}\texttt{\bfseries elabSpec\_top\_mem}}\hfill{\footnotesize \path{lean/su-functor/SuFunctor/Universe/SpecElab.lean}, line 935}\par\nopagebreak
\begin{leansrc}
theorem SpecRecord.elabSpec_top_mem {width : α → ℕ} {A : α → Set K}
    {r : SpecRecord α} {D : SpecDnf K}
    (_h : r.elabSpec width A = .ok D) :
    ∀ a ∈ D.atoms, (⊤ : K → ℝ≥0∞) ∈ a.marks
\end{leansrc}
\par\medskip\noindent{\small \texttt{SuFunctor.Universe.SpecRecord.}\texttt{\bfseries elabSpec}}\hfill{\footnotesize \path{lean/su-functor/SuFunctor/Universe/SpecElab.lean}, line 902}\par\nopagebreak
\begin{leansrc}
noncomputable def SpecRecord.elabSpec (width : α → ℕ) (A : α → Set K)
    (r : SpecRecord α) : Except (ConformanceError α) (SpecDnf K)
\end{leansrc}
\par\medskip\noindent{\small \texttt{SuFunctor.Universe.SpecRecord.}\texttt{\bfseries resolveArity\_isOk\_iff}}\hfill{\footnotesize \path{lean/su-functor/SuFunctor/Universe/SpecElab.lean}, line 742}\par\nopagebreak
\begin{leansrc}
theorem SpecRecord.resolveArity_isOk_iff {width : α → ℕ}
    {r : SpecRecord α} :
    IsOk (r.resolveArity width) ↔ r.Conformant width
\end{leansrc}
\par\medskip\noindent{\small \texttt{SuFunctor.Universe.SpecRecord.}\texttt{\bfseries resolveArity\_error\_names\_offender}}\hfill{\footnotesize \path{lean/su-functor/SuFunctor/Universe/SpecElab.lean}, line 835}\par\nopagebreak
\begin{leansrc}
theorem SpecRecord.resolveArity_error_names_offender {width : α → ℕ}
    {r : SpecRecord α} {a : α} {f e : ℕ} {b : Bool}
    (h : r.resolveArity width = .error (.widthMismatch a f e b)) :
    width a = f ∧ f ≠ e ∧
      ((b = false ∧ a ∈ r.root.atoms) ∨
        (b = true ∧ ∃ I, r.iface? = some I ∧
          ∃ φ, I.assumption? = some φ ∧ a ∈ φ.atoms))
\end{leansrc}
\par\medskip\noindent{\small \texttt{SuFunctor.Universe.SpecRecord.}\texttt{\bfseries elabSpec\_isOk\_iff}}\hfill{\footnotesize \path{lean/su-functor/SuFunctor/Universe/SpecElab.lean}, line 909}\par\nopagebreak
\begin{leansrc}
theorem SpecRecord.elabSpec_isOk_iff {width : α → ℕ} {A : α → Set K}
    {r : SpecRecord α} :
    IsOk (r.elabSpec width A) ↔ r.Conformant width
\end{leansrc}
\par\medskip\noindent{\small \texttt{SuFunctor.Universe.SpecRecord.}\texttt{\bfseries elabSpec\_ok\_modelClass}}\hfill{\footnotesize \path{lean/su-functor/SuFunctor/Universe/SpecElab.lean}, line 918}\par\nopagebreak
\begin{leansrc}
theorem SpecRecord.elabSpec_ok_modelClass {width : α → ℕ}
    {A : α → Set K} {r : SpecRecord α} {D : SpecDnf K}
    (h : r.elabSpec width A = .ok D) :
    D.modelClass = NnfForm.classSem A r.root.nnf
\end{leansrc}
\par\medskip\noindent{\small \texttt{SuFunctor.Universe.SpecDnf.}\texttt{\bfseries modelClass\_top}}\hfill{\footnotesize \path{lean/su-functor/SuFunctor/Universe/SpecElab.lean}, line 323}\par\nopagebreak
\begin{leansrc}
theorem SpecDnf.modelClass_top :
    (SpecDnf.top : SpecDnf K).modelClass = safeProper Set.univ
\end{leansrc}
\par\medskip\noindent{\small \texttt{SuFunctor.Universe.SpecDnf.}\texttt{\bfseries modelClass\_bot}}\hfill{\footnotesize \path{lean/su-functor/SuFunctor/Universe/SpecElab.lean}, line 328}\par\nopagebreak
\begin{leansrc}
theorem SpecDnf.modelClass_bot :
    (SpecDnf.bot : SpecDnf K).modelClass = ∅
\end{leansrc}
\par\medskip\noindent{\small \texttt{SuFunctor.Universe.}\texttt{\bfseries den\_compl\_le}}\hfill{\footnotesize \path{lean/su-functor/SuFunctor/Universe/SpecElab.lean}, line 601}\par\nopagebreak
\begin{leansrc}
theorem den_compl_le (f : K → ℝ) (c : ℚ) :
    {μ | f μ ≤ (c : ℝ)}ᶜ = {μ | (c : ℝ) < f μ}
\end{leansrc}
\par\medskip\noindent{\small \texttt{SuFunctor.Universe.}\texttt{\bfseries safetyAtom}}\hfill{\footnotesize \path{lean/su-functor/SuFunctor/Universe/SpecElab.lean}, line 176}\par\nopagebreak
\begin{leansrc}
noncomputable def safetyAtom (C : Set K) : SpecAtom K where …
\end{leansrc}
\par\medskip\noindent{\small \texttt{SuFunctor.Universe.}\texttt{\bfseries coupledSat\_safetyAtom\_iff}}\hfill{\footnotesize \path{lean/su-functor/SuFunctor/Universe/SpecElab.lean}, line 191}\par\nopagebreak
\begin{leansrc}
theorem coupledSat_safetyAtom_iff {C : Set K} {σ : K → ℝ≥0∞} :
    CoupledSat σ (wallOn C) {⊤} ↔
      permitted σ ⊆ C ∧ (permitted σ).Nonempty
\end{leansrc}
\par\medskip\noindent{\small \texttt{SuFunctor.Universe.}\texttt{\bfseries safetyAtom\_modelClass}}\hfill{\footnotesize \path{lean/su-functor/SuFunctor/Universe/SpecElab.lean}, line 205}\par\nopagebreak
\begin{leansrc}
theorem safetyAtom_modelClass (C : Set K) :
    (safetyAtom C).modelClass = safeProper C
\end{leansrc}
\par\medskip\noindent{\small \texttt{SuFunctor.Universe.SpecAtom.}\texttt{\bfseries modelClass\_merge}}\hfill{\footnotesize \path{lean/su-functor/SuFunctor/Universe/SpecElab.lean}, line 243}\par\nopagebreak
\begin{leansrc}
theorem SpecAtom.modelClass_merge (a b : SpecAtom K) :
    (a.merge b).modelClass = a.modelClass ∩ b.modelClass
\end{leansrc}
\par\medskip\noindent{\small \texttt{SuFunctor.Universe.SpecDnf.}\texttt{\bfseries modelClass\_conj}}\hfill{\footnotesize \path{lean/su-functor/SuFunctor/Universe/SpecElab.lean}, line 292}\par\nopagebreak
\begin{leansrc}
theorem SpecDnf.modelClass_conj (D E : SpecDnf K) :
    (D.conj E).modelClass = D.modelClass ∩ E.modelClass
\end{leansrc}
\par\medskip\noindent{\small \texttt{SuFunctor.Universe.SpecDnf.}\texttt{\bfseries modelClass\_disj}}\hfill{\footnotesize \path{lean/su-functor/SuFunctor/Universe/SpecElab.lean}, line 308}\par\nopagebreak
\begin{leansrc}
theorem SpecDnf.modelClass_disj (D E : SpecDnf K) :
    (D.disj E).modelClass = D.modelClass ∪ E.modelClass
\end{leansrc}
\par\medskip\noindent{\small \texttt{SuFunctor.Universe.NnfForm.}\texttt{\bfseries exists\_hedging\_gap}}\hfill{\footnotesize \path{lean/su-functor/SuFunctor/Universe/SpecElab.lean}, line 557}\par\nopagebreak
\begin{leansrc}
theorem exists_hedging_gap :
    ∃ (A : Bool → Set Bool) (ψ χ : NnfForm Bool) (σ : Bool → ℝ≥0∞),
      σ ∈ safeProper ((NnfForm.disj ψ χ).den A) ∧
        σ ∉ classSem A (NnfForm.disj ψ χ)
\end{leansrc}
\par\medskip\noindent{\small \texttt{SuFunctor.Universe.NnfForm.}\texttt{\bfseries elabNnf\_modelClass\_disjFree}}\hfill{\footnotesize \path{lean/su-functor/SuFunctor/Universe/SpecElab.lean}, line 546}\par\nopagebreak
\begin{leansrc}
theorem elabNnf_modelClass_disjFree (A : α → Set K) (ψ : NnfForm α)
    (h : ψ.DisjFree) :
    (elabNnf A ψ).modelClass = safeProper (ψ.den A)
\end{leansrc}
\par\medskip\noindent{\small \texttt{SuFunctor.Universe.StateFormula.}\texttt{\bfseries den\_nnf}}\hfill{\footnotesize \path{lean/su-functor/SuFunctor/Universe/SpecElab.lean}, line 462}\par\nopagebreak
\begin{leansrc}
theorem den_nnf (A : α → Set K) (φ : StateFormula α) :
    (φ.nnf).den A = φ.den A
\end{leansrc}
\par\medskip\noindent{\small \texttt{SuFunctor.Universe.}\texttt{\bfseries recordCon\_iff\_stages}}\hfill{\footnotesize \path{lean/su-functor/SuFunctor/Adequacy/SpecPathLens.lean}, line 720}\par\nopagebreak
\begin{leansrc}
theorem recordCon_iff_stages {p : NonemptyCompacts X × ProbState X} :
    p ∈ recordCon ↔ stages .dwell p ⊆ stages .touch p
\end{leansrc}
\par\medskip\noindent{\small \texttt{SuFunctor.Universe.}\texttt{\bfseries exists\_untilOn\_species\_incomparable}}\hfill{\footnotesize \path{lean/su-functor/SuFunctor/Adequacy/SpecPathLens.lean}, line 840}\par\nopagebreak
\begin{leansrc}
theorem exists_untilOn_species_incomparable :
    ∃ p ∈ recordCon (X := Bool),
      (untilOn (fun b => cond b 1 0) 1 ∅ {false} (stages .touch p) ∧
        ¬ untilOn (fun b => cond b 1 0) 1 ∅ {false} (stages .dwell p)) ∧
      (untilOn (fun b => cond b 1 0) 1 {true} {true} (stages .dwell p) ∧
        ¬ untilOn (fun b => cond b 1 0) 1 {true} {true}
          (stages .touch p))
\end{leansrc}
\par\medskip\noindent{\small \texttt{SuFunctor.Universe.}\texttt{\bfseries exists\_streettOn\_species\_incomparable}}\hfill{\footnotesize \path{lean/su-functor/SuFunctor/Adequacy/SpecPathLens.lean}, line 888}\par\nopagebreak
\begin{leansrc}
theorem exists_streettOn_species_incomparable :
    ∃ p ∈ recordCon (X := Bool),
      (streettOn [(({true} : Set Bool), ({false} : Set Bool))]
          (stages .touch p) ∧
        ¬ streettOn [(({true} : Set Bool), ({false} : Set Bool))]
          (stages .dwell p)) ∧
      (streettOn [(({false} : Set Bool), (∅ : Set Bool))]
          (stages .dwell p) ∧
        ¬ streettOn [(({false} : Set Bool), (∅ : Set Bool))]
          (stages .touch p))
\end{leansrc}
\par\medskip\noindent{\small \texttt{SuFunctor.Universe.}\texttt{\bfseries touchWall\_eq\_wallOn}}\hfill{\footnotesize \path{lean/su-functor/SuFunctor/Adequacy/SpecPathLens.lean}, line 744}\par\nopagebreak
\begin{leansrc}
theorem touchWall_eq_wallOn (Bad : Set X) :
    touchWall Bad
      = wallOn {p : NonemptyCompacts X × ProbState X |
          alwaysOn Badᶜ (stages .touch p)}
\end{leansrc}
\par\medskip\noindent{\small \texttt{SuFunctor.Universe.}\texttt{\bfseries dwellWall\_eq\_wallOn}}\hfill{\footnotesize \path{lean/su-functor/SuFunctor/Adequacy/SpecPathLens.lean}, line 761}\par\nopagebreak
\begin{leansrc}
theorem dwellWall_eq_wallOn (Bad : Set X) :
    dwellWall Bad
      = wallOn {p : NonemptyCompacts X × ProbState X |
          alwaysOn Badᶜ (stages .dwell p)}
\end{leansrc}
\par\medskip\noindent{\small \texttt{SuFunctor.Universe.}\texttt{\bfseries markOn}}\hfill{\footnotesize \path{lean/su-functor/SuFunctor/Adequacy/SpecPathLens.lean}, line 100}\par\nopagebreak
\begin{leansrc}
noncomputable def markOn (C : Set K) : K → ℝ≥0∞
\end{leansrc}
\par\medskip\noindent{\small \texttt{SuFunctor.Universe.}\texttt{\bfseries markOn\_univ}}\hfill{\footnotesize \path{lean/su-functor/SuFunctor/Adequacy/SpecPathLens.lean}, line 112}\par\nopagebreak
\begin{leansrc}
@[simp] theorem markOn_univ : markOn (Set.univ : Set K) = ⊤
\end{leansrc}
\par\medskip\noindent{\small \texttt{SuFunctor.Universe.}\texttt{\bfseries coupledSat\_pathAtom\_iff}}\hfill{\footnotesize \path{lean/su-functor/SuFunctor/Adequacy/SpecPathLens.lean}, line 167}\par\nopagebreak
\begin{leansrc}
theorem coupledSat_pathAtom_iff {C G : Set K} {σ : K → ℝ≥0∞} :
    CoupledSat σ (pathAtom C G).floor (pathAtom C G).marks ↔
      permitted σ ⊆ C ∧ (permitted σ ∩ G).Nonempty
\end{leansrc}
\par\medskip\noindent{\small \texttt{SuFunctor.Universe.}\texttt{\bfseries stages\_uniformLift}}\hfill{\footnotesize \path{lean/su-functor/SuFunctor/Adequacy/SpecPathLens.lean}, line 732}\par\nopagebreak
\begin{leansrc}
theorem stages_uniformLift (sp : RecordSpecies) (G : NonemptyCompacts X) :
    stages sp (G, uniformLift G) = (G : Set X)
\end{leansrc}
\par\medskip\noindent{\small \texttt{SuFunctor.Universe.}\texttt{\bfseries alwaysOn}}\hfill{\footnotesize \path{lean/su-functor/SuFunctor/Adequacy/SpecPathLens.lean}, line 266}\par\nopagebreak
\begin{leansrc}
def alwaysOn (Good : Set X) (S : Set X) : Prop
\end{leansrc}
\par\medskip\noindent{\small \texttt{SuFunctor.Universe.}\texttt{\bfseries isClosed\_recordClass\_always}}\hfill{\footnotesize \path{lean/su-functor/SuFunctor/Adequacy/SpecPathLens.lean}, line 1010}\par\nopagebreak
\begin{leansrc}
theorem isClosed_recordClass_always (sp : RecordSpecies) (τ : X → ℕ)
    (A : α → Set X) (φ : StateFormula α) :
    IsClosed (recordClass sp τ A (.always φ))
\end{leansrc}
\par\medskip\noindent{\small \texttt{SuFunctor.Universe.}\texttt{\bfseries windowAlwaysOn}}\hfill{\footnotesize \path{lean/su-functor/SuFunctor/Adequacy/SpecPathLens.lean}, line 270}\par\nopagebreak
\begin{leansrc}
def windowAlwaysOn (τ : X → ℕ) (n : ℕ) (Good : Set X) (S : Set X) :
    Prop
\end{leansrc}
\par\medskip\noindent{\small \texttt{SuFunctor.Universe.}\texttt{\bfseries isClosed\_recordClass\_alwaysT}}\hfill{\footnotesize \path{lean/su-functor/SuFunctor/Adequacy/SpecPathLens.lean}, line 1017}\par\nopagebreak
\begin{leansrc}
theorem isClosed_recordClass_alwaysT (sp : RecordSpecies) (τ : X → ℕ)
    (A : α → Set X) (n : ℕ) (φ : StateFormula α) :
    IsClosed (recordClass sp τ A (.alwaysT n φ))
\end{leansrc}
\par\medskip\noindent{\small \texttt{SuFunctor.Universe.}\texttt{\bfseries windowReachOn}}\hfill{\footnotesize \path{lean/su-functor/SuFunctor/Adequacy/SpecPathLens.lean}, line 276}\par\nopagebreak
\begin{leansrc}
def windowReachOn (τ : X → ℕ) (n : ℕ) (Goal : Set X) (S : Set X) :
    Prop
\end{leansrc}
\par\medskip\noindent{\small \texttt{SuFunctor.Universe.}\texttt{\bfseries isOpen\_recordClass\_deadline}}\hfill{\footnotesize \path{lean/su-functor/SuFunctor/Adequacy/SpecPathLens.lean}, line 1030}\par\nopagebreak
\begin{leansrc}
theorem isOpen_recordClass_deadline (sp : RecordSpecies) (τ : X → ℕ)
    (A : α → Set X) (n : ℕ) (φ : StateFormula α) (b : Bool) :
    IsOpen (recordClass sp τ A (.deadline n φ b))
\end{leansrc}
\par\medskip\noindent{\small \texttt{SuFunctor.Universe.}\texttt{\bfseries untilOn}}\hfill{\footnotesize \path{lean/su-functor/SuFunctor/Adequacy/SpecPathLens.lean}, line 283}\par\nopagebreak
\begin{leansrc}
def untilOn (τ : X → ℕ) (n : ℕ) (Safe Target : Set X) (S : Set X) :
    Prop
\end{leansrc}
\par\medskip\noindent{\small \texttt{SuFunctor.Universe.}\texttt{\bfseries recordClass\_untilF\_eq}}\hfill{\footnotesize \path{lean/su-functor/SuFunctor/Adequacy/SpecPathLens.lean}, line 1045}\par\nopagebreak
\begin{leansrc}
theorem recordClass_untilF_eq (sp : RecordSpecies) (τ : X → ℕ)
    (A : α → Set X) (n : ℕ) (sf tg : StateFormula α) :
    recordClass sp τ A (.untilF n sf tg)
      = ⋃ t : Fin (n + 1),
          ({p : NonemptyCompacts X × ProbState X |
              (stages sp p ∩ (tg.den A ∩ {x | τ x = (t : ℕ)})).Nonempty}
            ∩ {p : NonemptyCompacts X × ProbState X |
                stages sp p ⊆ sf.den A ∪ {x | ¬ τ x < (t : ℕ)}})
\end{leansrc}
\par\medskip\noindent{\small \texttt{SuFunctor.Universe.}\texttt{\bfseries isConstructibleUnion\_recordClass\_untilF}}\hfill{\footnotesize \path{lean/su-functor/SuFunctor/Adequacy/SpecPathLens.lean}, line 1062}\par\nopagebreak
\begin{leansrc}
theorem isConstructibleUnion_recordClass_untilF (sp : RecordSpecies)
    (τ : X → ℕ) (A : α → Set X) (n : ℕ) (sf tg : StateFormula α) :
    IsConstructibleUnion (recordClass sp τ A (.untilF n sf tg))
\end{leansrc}
\par\medskip\noindent{\small \texttt{Qtsl.QStreett.}\texttt{\bfseries sound\_QStreettBundle}}\hfill{\footnotesize \path{lean/qtsl-aux/QtslAux/QStreett.lean}, line 347}\par\nopagebreak
\begin{leansrc}
theorem sound_QStreettBundle {Rq G : ι → Set S} {W : ι → S → ℝ} {ε C : ι → ℝ}
    (hD : MeasurableSet D) (hDbar : MeasurableSet Dbar) (hU : MeasurableSet U)
    (hV₀ : Measurable V₀)
    (hRq : ∀ i, MeasurableSet (Rq i)) (hG : ∀ i, MeasurableSet (G i))
    (hW : ∀ i, Measurable (W i))
    (hΘ : p + thr ≤ 1)
    (hDsub : D ⊆ Dbar)
    (hR3 : ∀ y ∈ D, y ∉ U → K y Dbar = 1)
    (hR4 : ∀ y ∈ D, y ∉ U → K y (D ∪ U) = 1)
    (hV₀int : ∀ y ∈ D, y ∉ U → Integrable V₀ (K y))
    (hV₀drift : ∀ y ∈ D, y ∉ U → ∫ z, V₀ z ∂K y ≤ V₀ y)
    (hV₀floor : ∀ y ∈ U, y ∈ Dbar → 1 ≤ V₀ y)
    (hV₀nonneg : ∀ y ∈ Dbar, 0 ≤ V₀ y)
    (hε : ∀ i, 0 < ε i)
    (hW0 : ∀ i, ∀ y ∈ Dbar, 0 ≤ W i y) (hWC : ∀ i, ∀ y ∈ Dbar, W i y ≤ C i)
    (hWint : ∀ i, ∀ y ∈ D, y ∉ U → y ∉ G i → Integrable (W i) (K y))
    (hWdrift : ∀ i, ∀ y ∈ D, y ∉ U → y ∉ G i → ∫ z, W i z ∂K y ≤ W i y)
    (hWdriftR : ∀ i, ∀ y ∈ D, y ∉ U → y ∈ Rq i → y ∉ G i →
      ∫ z, W i z ∂K y ≤ W i y - ε i)
    (hx : x ∈ D) (hthr : V₀ x ≤ thr) :
    ENNReal.ofReal p ≤ pathMeasure K x
      {ω | ∀ i, (∃ᶠ n in atTop, ω n ∈ Rq i) → ∃ᶠ n in atTop, ω n ∈ G i}
\end{leansrc}
\par\medskip\noindent{\small \texttt{SuFunctor.Universe.}\texttt{\bfseries streettOn}}\hfill{\footnotesize \path{lean/su-functor/SuFunctor/Adequacy/SpecPathLens.lean}, line 292}\par\nopagebreak
\begin{leansrc}
def streettOn (prs : List (Set X × Set X)) (S : Set X) : Prop
\end{leansrc}
\par\medskip\noindent{\small \texttt{SuFunctor.Universe.}\texttt{\bfseries pathNnf\_notPath\_alwaysT}}\hfill{\footnotesize \path{lean/su-functor/SuFunctor/Adequacy/SpecPathLens.lean}, line 677}\par\nopagebreak
\begin{leansrc}
@[simp] theorem pathNnf_notPath_alwaysT (n : ℕ) (φ : StateFormula α) :
    pathNnf (PathFormula.notPath (.alwaysT n φ))
      = .ok (.deadline n φ.neg false)
\end{leansrc}
\par\medskip\noindent{\small \texttt{SuFunctor.Universe.}\texttt{\bfseries pathNnf\_notPath\_notPath\_deadline}}\hfill{\footnotesize \path{lean/su-functor/SuFunctor/Adequacy/SpecPathLens.lean}, line 683}\par\nopagebreak
\begin{leansrc}
@[simp] theorem pathNnf_notPath_notPath_deadline (n : ℕ)
    (φ : StateFormula α) (b : Bool) :
    pathNnf (PathFormula.notPath (.notPath (.deadline n φ b)))
      = .ok (.deadline n φ b)
\end{leansrc}
\par\medskip\noindent{\small \texttt{SuFunctor.Universe.}\texttt{\bfseries pathNnf\_notPath\_always}}\hfill{\footnotesize \path{lean/su-functor/SuFunctor/Adequacy/SpecPathLens.lean}, line 659}\par\nopagebreak
\begin{leansrc}
@[simp] theorem pathNnf_notPath_always (φ : StateFormula α) :
    pathNnf (PathFormula.notPath (.always φ))
      = .error .negatedAlways
\end{leansrc}
\par\medskip\noindent{\small \texttt{SuFunctor.Universe.}\texttt{\bfseries pathNnf\_notPath\_untilF}}\hfill{\footnotesize \path{lean/su-functor/SuFunctor/Adequacy/SpecPathLens.lean}, line 664}\par\nopagebreak
\begin{leansrc}
@[simp] theorem pathNnf_notPath_untilF (n : ℕ)
    (sf tg : StateFormula α) :
    pathNnf (PathFormula.notPath (.untilF n sf tg))
      = .error .negatedUntil
\end{leansrc}
\par\medskip\noindent{\small \texttt{SuFunctor.Universe.}\texttt{\bfseries pathNnf\_notPath\_streett}}\hfill{\footnotesize \path{lean/su-functor/SuFunctor/Adequacy/SpecPathLens.lean}, line 671}\par\nopagebreak
\begin{leansrc}
@[simp] theorem pathNnf_notPath_streett (prs : List (StreettPair α)) :
    pathNnf (PathFormula.notPath (.streett prs))
      = .error .negatedStreett
\end{leansrc}
\par\medskip\noindent{\small \texttt{SuFunctor.Universe.}\texttt{\bfseries not\_alwaysOn\_iff}}\hfill{\footnotesize \path{lean/su-functor/SuFunctor/Adequacy/SpecPathLens.lean}, line 422}\par\nopagebreak
\begin{leansrc}
theorem not_alwaysOn_iff {Good S : Set X} :
    ¬ alwaysOn Good S ↔ ∃ x ∈ S, x ∈ Goodᶜ
\end{leansrc}
\par\medskip\noindent{\small \texttt{SuFunctor.Universe.}\texttt{\bfseries satCredal\_lower\_singletonC\_zero}}\hfill{\footnotesize \path{lean/su-functor/SuFunctor/Adequacy/SpecBehavioral.lean}, line 143}\par\nopagebreak
\begin{leansrc}
theorem satCredal_lower_singletonC_zero (c : SubPMF V → Prop) :
    SatCredal .lower (singletonC (SubPMF.zero V)) c
\end{leansrc}
\par\medskip\noindent{\small \texttt{SuFunctor.Universe.}\texttt{\bfseries probThresholdSet}}\hfill{\footnotesize \path{lean/su-functor/SuFunctor/Adequacy/SpecBehavioral.lean}, line 168}\par\nopagebreak
\begin{leansrc}
def probThresholdSet (K : Kernel S S) [IsMarkovKernel K] (p : ℝ)
    (E : Set (ℕ → S)) : Set S
\end{leansrc}
\par\medskip\noindent{\small \texttt{SuFunctor.Universe.}\texttt{\bfseries safetyAtom\_modelClass\_mono}}\hfill{\footnotesize \path{lean/su-functor/SuFunctor/Adequacy/SpecBehavioral.lean}, line 92}\par\nopagebreak
\begin{leansrc}
theorem safetyAtom_modelClass_mono {C D : Set X} (h : C ⊆ D) :
    (safetyAtom C).modelClass ⊆ (safetyAtom D).modelClass
\end{leansrc}
\par\medskip\noindent{\small \texttt{SuFunctor.Universe.}\texttt{\bfseries sprobAtom\_modelClass\_mono}}\hfill{\footnotesize \path{lean/su-functor/SuFunctor/Adequacy/SpecBehavioral.lean}, line 210}\par\nopagebreak
\begin{leansrc}
theorem sprobAtom_modelClass_mono {p : ℝ} {E E' : Set (ℕ → S)}
    (hE : E ⊆ E') :
    (safetyAtom (probThresholdSet K p E)).modelClass ⊆
      (safetyAtom (probThresholdSet K p E')).modelClass
\end{leansrc}
\par\medskip\noindent{\small \texttt{SuFunctor.Universe.}\texttt{\bfseries sprob\_sound\_of\_streettBundle}}\hfill{\footnotesize \path{lean/su-functor/SuFunctor/Adequacy/SpecBehavioral.lean}, line 224}\par\nopagebreak
\begin{leansrc}
theorem sprob_sound_of_streettBundle
    {K : Kernel S S} [IsMarkovKernel K]
    {D Dbar U : Set S} {V₀ : S → ℝ} {p thr : ℝ} {x : S}
    {ι : Type*} [Countable ι]
    {Rq G : ι → Set S} {W : ι → S → ℝ} {ε C : ι → ℝ}
    (hD : MeasurableSet D) (hDbar : MeasurableSet Dbar)
    (hU : MeasurableSet U) (hV₀ : Measurable V₀)
    (hRq : ∀ i, MeasurableSet (Rq i)) (hG : ∀ i, MeasurableSet (G i))
    (hW : ∀ i, Measurable (W i))
    (hΘ : p + thr ≤ 1)
    (hDsub : D ⊆ Dbar)
    (hR3 : ∀ y ∈ D, y ∉ U → K y Dbar = 1)
    (hR4 : ∀ y ∈ D, y ∉ U → K y (D ∪ U) = 1)
    (hV₀int : ∀ y ∈ D, y ∉ U → Integrable V₀ (K y))
    (hV₀drift : ∀ y ∈ D, y ∉ U → ∫ z, V₀ z ∂K y ≤ V₀ y)
    (hV₀floor : ∀ y ∈ U, y ∈ Dbar → 1 ≤ V₀ y)
    (hV₀nonneg : ∀ y ∈ Dbar, 0 ≤ V₀ y)
    (hε : ∀ i, 0 < ε i)
    (hW0 : ∀ i, ∀ y ∈ Dbar, 0 ≤ W i y)
    (hWC : ∀ i, ∀ y ∈ Dbar, W i y ≤ C i)
    (hWint : ∀ i, ∀ y ∈ D, y ∉ U → y ∉ G i → Integrable (W i) (K y))
    (hWdrift : ∀ i, ∀ y ∈ D, y ∉ U → y ∉ G i →
      ∫ z, W i z ∂K y ≤ W i y)
    (hWdriftR : ∀ i, ∀ y ∈ D, y ∉ U → y ∈ Rq i → y ∉ G i →
      ∫ z, W i z ∂K y ≤ W i y - ε i)
    (hx : x ∈ D) (hthr : V₀ x ≤ thr)
    {E : Set (ℕ → S)}
    (hE : {ω | ∀ i, (∃ᶠ n in atTop, ω n ∈ Rq i) →
      ∃ᶠ n in atTop, ω n ∈ G i} ⊆ E) :
    x ∈ probThresholdSet K p E
\end{leansrc}
\par\medskip\noindent{\small \texttt{SuFunctor.Universe.}\texttt{\bfseries satCredal\_lower\_singletonC\_geC}}\hfill{\footnotesize \path{lean/su-functor/SuFunctor/Adequacy/SpecBehavioral.lean}, line 115}\par\nopagebreak
\begin{leansrc}
theorem satCredal_lower_singletonC_geC {μ : SubPMF V}
    (hμ : μ ≠ SubPMF.zero V) (A : Set V) (q : ℚ) :
    SatCredal .lower (singletonC μ) (geC A q) ↔ q ≤ μ.prob A
\end{leansrc}
\par\medskip\noindent{\small \texttt{SuFunctor.Universe.}\texttt{\bfseries returnGESet}}\hfill{\footnotesize \path{lean/su-functor/SuFunctor/Adequacy/SpecBehavioral.lean}, line 284}\par\nopagebreak
\begin{leansrc}
def returnGESet (K : Kernel S S) [IsMarkovKernel K] (r : S → ℝ)
    (b : ℝ) (n : ℕ) : Set S
\end{leansrc}
\par\medskip\noindent{\small \texttt{SuFunctor.Universe.}\texttt{\bfseries returnGESet\_mono\_reward}}\hfill{\footnotesize \path{lean/su-functor/SuFunctor/Adequacy/SpecBehavioral.lean}, line 299}\par\nopagebreak
\begin{leansrc}
theorem returnGESet_mono_reward {r r' : S → ℝ} {b : ℝ} {n : ℕ}
    (hr : ∀ s, r s ≤ r' s)
    (hint : ∀ x, Integrable (finReturn r n) (pathMeasure K x))
    (hint' : ∀ x, Integrable (finReturn r' n) (pathMeasure K x)) :
    returnGESet K r b n ⊆ returnGESet K r' b n
\end{leansrc}
\par\medskip\noindent{\small \texttt{SuFunctor.Universe.}\texttt{\bfseries returnGEAtom\_modelClass\_mono\_reward}}\hfill{\footnotesize \path{lean/su-functor/SuFunctor/Adequacy/SpecBehavioral.lean}, line 334}\par\nopagebreak
\begin{leansrc}
theorem returnGEAtom_modelClass_mono_reward {r r' : S → ℝ} {b : ℝ}
    {n : ℕ} (hr : ∀ s, r s ≤ r' s)
    (hint : ∀ x, Integrable (finReturn r n) (pathMeasure K x))
    (hint' : ∀ x, Integrable (finReturn r' n) (pathMeasure K x)) :
    (safetyAtom (returnGESet K r b n)).modelClass ⊆
      (safetyAtom (returnGESet K r' b n)).modelClass
\end{leansrc}
\par\medskip\noindent{\small \texttt{SuFunctor.Universe.}\texttt{\bfseries returnGE\_sound\_of\_deadline}}\hfill{\footnotesize \path{lean/su-functor/SuFunctor/Adequacy/SpecBehavioral.lean}, line 441}\par\nopagebreak
\begin{leansrc}
theorem returnGE_sound_of_deadline [IsProbabilityMeasure μ]
    (hβ : 0 ≤ β) (ha : 0 < a)
    (hint : ∀ n, Integrable (L n) μ) (hnonneg : 0 ≤ L)
    (hmB : ∀ n, MeasurableSet[𝒢 n] {ω | X n ω ∈ B})
    (hdrift : ∀ n, ∀ᵐ ω ∂μ, X n ω ∉ B →
      (μ[L (n + 1) | 𝒢 n]) ω ≤ β * L n ω)
    (hbarrier : ∀ n ω, X n ω ∉ B → a ≤ L n ω) (n : ℕ) :
    1 - β ^ n * (∫ ω, L 0 ω ∂μ) / a ≤
      ∫ ω, finReturn (B.indicator fun _ => (1 : ℝ)) (n + 1)
        (fun t => X t ω) ∂μ
\end{leansrc}
\par\medskip\noindent{\small \texttt{SuFunctor.Universe.}\texttt{\bfseries discountedReturn\_le\_of\_drift}}\hfill{\footnotesize \path{lean/su-functor/SuFunctor/Adequacy/SpecBehavioral.lean}, line 543}\par\nopagebreak
\begin{leansrc}
theorem discountedReturn_le_of_drift [IsFiniteMeasure μ]
    {γ : ℝ} {r : ℕ → Ω → ℝ}
    (hβ : 0 ≤ β) (hγ : 0 ≤ γ)
    (hint : ∀ n, Integrable (L n) μ) (hnonneg : 0 ≤ L)
    (hdrift : ∀ n, ∀ᵐ ω ∂μ, (μ[L (n + 1) | 𝒢 n]) ω ≤ β * L n ω)
    (hrint : ∀ n, Integrable (r n) μ)
    (hr : ∀ n ω, r n ω ≤ L n ω) (n : ℕ) :
    ∫ ω, discountedPartialReturn γ r n ω ∂μ ≤
      (∑ t ∈ Finset.range n, (γ * β) ^ t) * ∫ ω, L 0 ω ∂μ
\end{leansrc}
\par\medskip\noindent{\small \texttt{SuFunctor.Universe.}\texttt{\bfseries integral\_le\_pow\_mul}}\hfill{\footnotesize \path{lean/su-functor/SuFunctor/Adequacy/SpecBehavioral.lean}, line 511}\par\nopagebreak
\begin{leansrc}
theorem integral_le_pow_mul [IsFiniteMeasure μ] (hβ : 0 ≤ β)
    (hint : ∀ n, Integrable (L n) μ) (hnonneg : 0 ≤ L)
    (hdrift : ∀ n, ∀ᵐ ω ∂μ, (μ[L (n + 1) | 𝒢 n]) ω ≤ β * L n ω)
    (n : ℕ) :
    ∫ ω, L n ω ∂μ ≤ β ^ n * ∫ ω, L 0 ω ∂μ
\end{leansrc}
\par\medskip\noindent{\small \texttt{SuFunctor.}\texttt{\bfseries livenessTilt\_top}}\hfill{\footnotesize \path{lean/su-functor-legacy/SuFunctorLegacy/Threshold.lean}, line 275}\par\nopagebreak
\begin{leansrc}
theorem livenessTilt_top (goal : Set V) (p : ℚ) :
    livenessTilt goal p ⊤ = livenessWall goal p
\end{leansrc}
\par\medskip\noindent{\small \texttt{SuFunctor.Universe.}\texttt{\bfseries untilOn\_univ\_safe}}\hfill{\footnotesize \path{lean/su-functor/SuFunctor/Adequacy/SpecPathLens.lean}, line 392}\par\nopagebreak
\begin{leansrc}
theorem untilOn_univ_safe {τ : X → ℕ} {n : ℕ} {Target S : Set X} :
    untilOn τ n Set.univ Target S ↔ windowReachOn τ n Target S
\end{leansrc}
\par\medskip\noindent{\small \texttt{SuFunctor.Universe.}\texttt{\bfseries streettOn\_singleton\_univ\_iff}}\hfill{\footnotesize \path{lean/su-functor/SuFunctor/Universe/SpecRewrite.lean}, line 599}\par\nopagebreak
\begin{leansrc}
theorem streettOn_singleton_univ_iff {Y : Type*} {G S : Set Y} :
    streettOn [(Set.univ, G)] S ↔ (S.Nonempty → (S ∩ G).Nonempty)
\end{leansrc}
\par\medskip\noindent{\small \texttt{SuFunctor.Universe.}\texttt{\bfseries streettOn\_singleton\_univ\_empty}}\hfill{\footnotesize \path{lean/su-functor/SuFunctor/Universe/SpecRewrite.lean}, line 625}\par\nopagebreak
\begin{leansrc}
theorem streettOn_singleton_univ_empty {Y : Type*} {G : Set Y} :
    streettOn [(Set.univ, G)] (∅ : Set Y) ∧ ¬ ((∅ : Set Y) ∩ G).Nonempty
\end{leansrc}
\par\medskip\noindent{\small \texttt{Qtsl.TemporalRewrite.}\texttt{\bfseries no\_f\_and\_fusion}}\hfill{\footnotesize \path{lean/qtsl-aux/QtslAux/TemporalRewrite.lean}, line 274}\par\nopagebreak
\begin{leansrc}
theorem no_f_and_fusion :
    ∃ V : Nat → Nat → Prop,
      sat V 0 (.and (.F (.atom 0)) (.F (.atom 1))) ∧
      ¬ sat V 0 (.F (.and (.atom 0) (.atom 1)))
\end{leansrc}
\par\medskip\noindent{\small \texttt{Qtsl.TemporalRewrite.}\texttt{\bfseries no\_g\_or\_fusion}}\hfill{\footnotesize \path{lean/qtsl-aux/QtslAux/TemporalRewrite.lean}, line 285}\par\nopagebreak
\begin{leansrc}
theorem no_g_or_fusion :
    ∃ V : Nat → Nat → Prop,
      sat V 0 (.G (.or (.atom 0) (.atom 1))) ∧
      ¬ sat V 0 (.or (.G (.atom 0)) (.G (.atom 1)))
\end{leansrc}
\par\medskip\noindent{\small \texttt{Qtsl.TemporalRewrite.}\texttt{\bfseries no\_gf\_and\_fusion}}\hfill{\footnotesize \path{lean/qtsl-aux/QtslAux/TemporalRewrite.lean}, line 298}\par\nopagebreak
\begin{leansrc}
theorem no_gf_and_fusion :
    ∃ V : Nat → Nat → Prop,
      sat V 0 (.and (.G (.F (.atom 0))) (.G (.F (.atom 1)))) ∧
      ¬ sat V 0 (.G (.F (.and (.atom 0) (.atom 1))))
\end{leansrc}
\par\medskip\noindent{\small \texttt{Qtsl.TemporalRewrite.}\texttt{\bfseries no\_fg\_or\_fusion}}\hfill{\footnotesize \path{lean/qtsl-aux/QtslAux/TemporalRewrite.lean}, line 311}\par\nopagebreak
\begin{leansrc}
theorem no_fg_or_fusion :
    ∃ V : Nat → Nat → Prop,
      sat V 0 (.F (.G (.or (.atom 0) (.atom 1)))) ∧
      ¬ sat V 0 (.or (.F (.G (.atom 0))) (.F (.G (.atom 1))))
\end{leansrc}
\par\medskip\noindent{\small \texttt{SuFunctor.Universe.}\texttt{\bfseries exists\_denEq\_not\_classEq}}\hfill{\footnotesize \path{lean/su-functor/SuFunctor/Universe/SpecRewrite.lean}, line 206}\par\nopagebreak
\begin{leansrc}
theorem exists_denEq_not_classEq :
    ∃ (A : Fin 3 → Set Bool) (φ ψ : StateFormula (Fin 3)),
      DenEq A φ ψ ∧ ¬ ClassEq A φ ψ
\end{leansrc}
\par\medskip\noindent{\small \texttt{SuFunctor.Universe.}\texttt{\bfseries classEquiv\_of\_nnf\_eq}}\hfill{\footnotesize \path{lean/su-functor/SuFunctor/Universe/SpecRewrite.lean}, line 139}\par\nopagebreak
\begin{leansrc}
theorem classEquiv_of_nnf_eq {φ ψ : StateFormula α} (h : φ.nnf = ψ.nnf) :
    ClassEquiv.{u} φ ψ
\end{leansrc}
\par\medskip\noindent Modules cited: \texttt{SuFunctor.Universe.SpecElab}, \texttt{SuFunctor.Universe.SpecRewrite}, \texttt{SuFunctor.Universe.PairSequent}.
\par\medskip\noindent Also cited in this chapter, listed under an earlier one: \texttt{General.iSup\_gambleRows\_eq\_self}, \texttt{iSup\_boxRows\_eq\_self}, \texttt{sat\_sup\_iff}, \texttt{top\_notMem\_modelClass}.

\section{Rule Catalogue}
Citations from Chapter~\ref{ch:qtsl-manual}.

\par\medskip\noindent{\small \texttt{PPLC.}\texttt{\bfseries PAxData}}\hfill{\footnotesize \path{lean/qtsl/Qtsl/Core/Catalogue.lean}, line 1353}\par\nopagebreak
\begin{leansrc}
inductive PAxData : Hypernet → Sequent → Type where …
\end{leansrc}
\par\medskip\noindent{\small \texttt{SuFunctor.StateSpace.FamObj.}\texttt{\bfseries famEq}}\hfill{\footnotesize \path{lean/su-functor/SuFunctor/StateSpace/FamLimits.lean}, line 169}\par\nopagebreak
\begin{leansrc}
def famEq : FamObj C
\end{leansrc}
\par\medskip\noindent{\small \texttt{PPLC.}\texttt{\bfseries sound\_ax}}\hfill{\footnotesize \path{lean/qtsl/Qtsl/Soundness.lean}, line 4165}\par\nopagebreak
\begin{leansrc}
theorem sound_ax {H : Hypernet} {seq : Sequent}
    (h : seq.pre = seq.post) : seq.semanticallyValid H
\end{leansrc}
\par\medskip\noindent{\small \texttt{PPLC.Formula.}\texttt{\bfseries holds}}\hfill{\footnotesize \path{lean/qtsl/Qtsl/Catalogue.lean}, line 402}\par\nopagebreak
\begin{leansrc}
def Formula.holds {H : Hypernet} (σ : H.Valuation) : Formula → Prop
  | .atom b l r => ∃ a a' : QTensor,
      Term.evalQ σ l = some a ∧ Term.evalQ σ r = some a' ∧
      QTensor.rel2 (BowTie.holdsQ b) a a'
  | .inRegion t simplex S =>
      simplex = S.simplex ∧
      ∃ q : QTensor, Term.evalQ σ t = some q ∧
        q.shape = [S.simplex.complex.ambDim] ∧
        MaskRegion S.simplex.complex S.simplex.contains S.contains
          (fun j => q.entries.getD j.val 0)
  | .and l r => Formula.holds σ l ∧ Formula.holds σ r
  | .or l r => Formula.holds σ l ∨ Formula.holds σ r
  | .tt => True
  | .ff => False

namespace ContractPattern

variable {inputRanks : List Nat} {outputRank : Nat}

def idxShape (pat : ContractPattern inputRanks outputRank) : Shape
\end{leansrc}
\par\medskip\noindent{\small \texttt{PPLC.Sequent.}\texttt{\bfseries semanticallyValid}}\hfill{\footnotesize \path{lean/qtsl/Qtsl/Catalogue.lean}, line 527}\par\nopagebreak
\begin{leansrc}
def Sequent.semanticallyValid (H : Hypernet) (seq : Sequent) : Prop
\end{leansrc}
\par\medskip\noindent{\small \texttt{PPLC.M3.}\texttt{\bfseries le\_refl}}\hfill{\footnotesize \path{lean/qtsl/Qtsl/Admissible.lean}, line 112}\par\nopagebreak
\begin{leansrc}
theorem le_refl : ∀ x : M3, le x x = true
\end{leansrc}
\par\medskip\noindent{\small \texttt{PPLC.}\texttt{\bfseries PWeakenData}}\hfill{\footnotesize \path{lean/qtsl/Qtsl/Core/Catalogue.lean}, line 1357}\par\nopagebreak
\begin{leansrc}
inductive PWeakenData : Hypernet → Sequent → Type where …
\end{leansrc}
\par\medskip\noindent{\small \texttt{PPLC.}\texttt{\bfseries ProofObject}}\hfill{\footnotesize \path{lean/qtsl/Qtsl/Core/Catalogue.lean}, line 1316}\par\nopagebreak
\begin{leansrc}
inductive ProofObject : Hypernet → Sequent → Type where …
\end{leansrc}
\par\medskip\noindent{\small \texttt{PPLC.}\texttt{\bfseries sound\_weaken}}\hfill{\footnotesize \path{lean/qtsl/Qtsl/Soundness.lean}, line 4171}\par\nopagebreak
\begin{leansrc}
theorem sound_weaken {H : Hypernet} {seq premise : Sequent} {extra : Formula}
    (ih : premise.semanticallyValid H)
    (hpre : seq.pre = .and premise.pre extra)
    (hpost : seq.post = premise.post) : seq.semanticallyValid H
\end{leansrc}
\par\medskip\noindent{\small \texttt{PPLC.M3.}\texttt{\bfseries weaken\_kit}}\hfill{\footnotesize \path{lean/qtsl/Qtsl/Admissible.lean}, line 121}\par\nopagebreak
\begin{leansrc}
theorem weaken_kit : ∀ x y z : M3,
    le x z = true → le (inf x y) z = true
\end{leansrc}
\par\medskip\noindent{\small \texttt{PPLC.}\texttt{\bfseries sound\_cut}}\hfill{\footnotesize \path{lean/qtsl/Qtsl/Soundness.lean}, line 4180}\par\nopagebreak
\begin{leansrc}
theorem sound_cut {H : Hypernet} {seq leftPremise rightPremise : Sequent}
    {cutFormula : Formula}
    (ihl : leftPremise.semanticallyValid H)
    (ihr : rightPremise.semanticallyValid H)
    (hpre : seq.pre = leftPremise.pre)
    (hpost : seq.post = rightPremise.post)
    (hcutL : leftPremise.post = cutFormula)
    (hcutR : rightPremise.pre = .and leftPremise.pre cutFormula) :
    seq.semanticallyValid H
\end{leansrc}
\par\medskip\noindent{\small \texttt{PPLC.}\texttt{\bfseries PCutData}}\hfill{\footnotesize \path{lean/qtsl/Qtsl/Core/Catalogue.lean}, line 1365}\par\nopagebreak
\begin{leansrc}
inductive PCutData : Hypernet → Sequent → Type where …
\end{leansrc}
\par\medskip\noindent{\small \texttt{PPLC.M3.}\texttt{\bfseries cut\_kit}}\hfill{\footnotesize \path{lean/qtsl/Qtsl/Admissible.lean}, line 125}\par\nopagebreak
\begin{leansrc}
theorem cut_kit : ∀ x m y : M3,
    le x m = true → le (inf x m) y = true → le x y = true
\end{leansrc}
\par\medskip\noindent{\small \texttt{PPLC.}\texttt{\bfseries sound\_andI}}\hfill{\footnotesize \path{lean/qtsl/Qtsl/Soundness.lean}, line 4198}\par\nopagebreak
\begin{leansrc}
theorem sound_andI {H : Hypernet} {seq leftPremise rightPremise : Sequent}
    (ihl : leftPremise.semanticallyValid H)
    (ihr : rightPremise.semanticallyValid H)
    (hpreMatch : leftPremise.pre = rightPremise.pre)
    (hpre : seq.pre = leftPremise.pre)
    (hpost : seq.post = .and leftPremise.post rightPremise.post) :
    seq.semanticallyValid H
\end{leansrc}
\par\medskip\noindent{\small \texttt{PPLC.}\texttt{\bfseries PAndIData}}\hfill{\footnotesize \path{lean/qtsl/Qtsl/Core/Catalogue.lean}, line 1376}\par\nopagebreak
\begin{leansrc}
inductive PAndIData : Hypernet → Sequent → Type where …
\end{leansrc}
\par\medskip\noindent{\small \texttt{PPLC.M3.}\texttt{\bfseries andI\_kit}}\hfill{\footnotesize \path{lean/qtsl/Qtsl/Admissible.lean}, line 129}\par\nopagebreak
\begin{leansrc}
theorem andI_kit : ∀ x y z : M3,
    le x y = true → le x z = true → le x (inf y z) = true
\end{leansrc}
\par\medskip\noindent{\small \texttt{PPLC.}\texttt{\bfseries sound\_andE\_left}}\hfill{\footnotesize \path{lean/qtsl/Qtsl/Soundness.lean}, line 4212}\par\nopagebreak
\begin{leansrc}
theorem sound_andE_left {H : Hypernet} {seq premise : Sequent} {other : Formula}
    (ih : premise.semanticallyValid H)
    (hpre : seq.pre = premise.pre)
    (hpost : premise.post = .and seq.post other) :
    seq.semanticallyValid H
\end{leansrc}
\par\medskip\noindent{\small \texttt{PPLC.}\texttt{\bfseries sound\_andE\_right}}\hfill{\footnotesize \path{lean/qtsl/Qtsl/Soundness.lean}, line 4223}\par\nopagebreak
\begin{leansrc}
theorem sound_andE_right {H : Hypernet} {seq premise : Sequent} {other : Formula}
    (ih : premise.semanticallyValid H)
    (hpre : seq.pre = premise.pre)
    (hpost : premise.post = .and other seq.post) :
    seq.semanticallyValid H
\end{leansrc}
\par\medskip\noindent{\small \texttt{PPLC.M3.}\texttt{\bfseries andE\_l\_kit}}\hfill{\footnotesize \path{lean/qtsl/Qtsl/Admissible.lean}, line 133}\par\nopagebreak
\begin{leansrc}
theorem andE_l_kit : ∀ x y z : M3,
    le x (inf y z) = true → le x y = true
\end{leansrc}
\par\medskip\noindent{\small \texttt{PPLC.M3.}\texttt{\bfseries andE\_r\_kit}}\hfill{\footnotesize \path{lean/qtsl/Qtsl/Admissible.lean}, line 137}\par\nopagebreak
\begin{leansrc}
theorem andE_r_kit : ∀ x y z : M3,
    le x (inf y z) = true → le x z = true
\end{leansrc}
\par\medskip\noindent{\small \texttt{PPLC.}\texttt{\bfseries sound\_orI\_left}}\hfill{\footnotesize \path{lean/qtsl/Qtsl/Soundness.lean}, line 4234}\par\nopagebreak
\begin{leansrc}
theorem sound_orI_left {H : Hypernet} {seq premise : Sequent} {other : Formula}
    (ih : premise.semanticallyValid H)
    (hpre : seq.pre = premise.pre)
    (hpost : seq.post = .or premise.post other) :
    seq.semanticallyValid H
\end{leansrc}
\par\medskip\noindent{\small \texttt{PPLC.}\texttt{\bfseries sound\_orI\_right}}\hfill{\footnotesize \path{lean/qtsl/Qtsl/Soundness.lean}, line 4244}\par\nopagebreak
\begin{leansrc}
theorem sound_orI_right {H : Hypernet} {seq premise : Sequent} {other : Formula}
    (ih : premise.semanticallyValid H)
    (hpre : seq.pre = premise.pre)
    (hpost : seq.post = .or other premise.post) :
    seq.semanticallyValid H
\end{leansrc}
\par\medskip\noindent{\small \texttt{PPLC.M3.}\texttt{\bfseries orI\_l\_kit}}\hfill{\footnotesize \path{lean/qtsl/Qtsl/Admissible.lean}, line 141}\par\nopagebreak
\begin{leansrc}
theorem orI_l_kit : ∀ x y z : M3,
    le x y = true → le x (sup y z) = true
\end{leansrc}
\par\medskip\noindent{\small \texttt{PPLC.M3.}\texttt{\bfseries orI\_r\_kit}}\hfill{\footnotesize \path{lean/qtsl/Qtsl/Admissible.lean}, line 145}\par\nopagebreak
\begin{leansrc}
theorem orI_r_kit : ∀ x y z : M3,
    le x z = true → le x (sup y z) = true
\end{leansrc}
\par\medskip\noindent{\small \texttt{PPLC.}\texttt{\bfseries sound\_orE}}\hfill{\footnotesize \path{lean/qtsl/Qtsl/Soundness.lean}, line 4254}\par\nopagebreak
\begin{leansrc}
theorem sound_orE {H : Hypernet} {seq leftPremise rightPremise : Sequent}
    (ihl : leftPremise.semanticallyValid H)
    (ihr : rightPremise.semanticallyValid H)
    (hpostMatch : leftPremise.post = rightPremise.post)
    (hpre : seq.pre = .or leftPremise.pre rightPremise.pre)
    (hpost : seq.post = leftPremise.post) :
    seq.semanticallyValid H
\end{leansrc}
\par\medskip\noindent{\small \texttt{PPLC.}\texttt{\bfseries OrECtxStep}}\hfill{\footnotesize \path{lean/qtsl/Qtsl/Admissible.lean}, line 20}\par\nopagebreak
\begin{leansrc}
structure OrECtxStep where …
\end{leansrc}
\par\medskip\noindent{\small \texttt{PPLC.}\texttt{\bfseries sound\_orECtx}}\hfill{\footnotesize \path{lean/qtsl/Qtsl/Admissible.lean}, line 33}\par\nopagebreak
\begin{leansrc}
theorem sound_orECtx {H : Hypernet} (st : OrECtxStep)
    (ihl : st.leftPremise.semanticallyValid H)
    (ihr : st.rightPremise.semanticallyValid H) :
    st.seq.semanticallyValid H
\end{leansrc}
\par\medskip\noindent{\small \texttt{PPLC.}\texttt{\bfseries cx\_not\_derivable}}\hfill{\footnotesize \path{lean/qtsl/Qtsl/Admissible.lean}, line 657}\par\nopagebreak
\begin{leansrc}
theorem cx_not_derivable : IsEmpty (ProofObject cxH cxSeq)
\end{leansrc}
\par\medskip\noindent{\small \texttt{PPLC.}\texttt{\bfseries orECtx\_not\_elaborable}}\hfill{\footnotesize \path{lean/qtsl/Qtsl/Admissible.lean}, line 666}\par\nopagebreak
\begin{leansrc}
theorem orECtx_not_elaborable :
    ¬ (∀ (st : OrECtxStep) (H : Hypernet),
        Nonempty (ProofObject H st.leftPremise) →
        Nonempty (ProofObject H st.rightPremise) →
        Nonempty (ProofObject H st.seq))
\end{leansrc}
\par\medskip\noindent{\small \texttt{PPLC.ProofObject.}\texttt{\bfseries m3sound}}\hfill{\footnotesize \path{lean/qtsl/Qtsl/Admissible.lean}, line 406}\par\nopagebreak
\begin{leansrc}
theorem ProofObject.m3sound : ∀ {H : Hypernet} {seq : Sequent},
    (∀ w, H.wireType w = none) → ProofObject H seq →
    M3.le (mval seq.pre) (mval seq.post) = true
  | _, _, _, .mk _ _ _ (.ax _ _ (.mk _ _ hpp)) => by
      rw [hpp]; exact M3.le_refl _
  | _, _, hE, .mk _ _ _ (.weaken _ _ (.mk _ _ _ pproof _ hpre hpost)) => by
      rw [hpre, hpost]
      exact M3.weaken_kit _ _ _ (pproof.m3sound hE)
  | _, _, hE, .mk _ _ _ (.cut _ _
      (.mk _ _ _ _ lproof _ rproof hpre hpost hcl hcr)) => by
      have hl
\end{leansrc}
\par\medskip\noindent{\small \texttt{PPLC.}\texttt{\bfseries cxStep}}\hfill{\footnotesize \path{lean/qtsl/Qtsl/Admissible.lean}, line 632}\par\nopagebreak
\begin{leansrc}
def cxStep : OrECtxStep where …
\end{leansrc}
\par\medskip\noindent{\small \texttt{PPLC.OrECtxStep.}\texttt{\bfseries elabHere}}\hfill{\footnotesize \path{lean/qtsl/Qtsl/Admissible.lean}, line 58}\par\nopagebreak
\begin{leansrc}
def OrECtxStep.elabHere {H : Hypernet} (st : OrECtxStep)
    (hpath : st.path = .here)
    (lpf : ProofObject H st.leftPremise) (rpf : ProofObject H st.rightPremise) :
    ProofObject H st.seq
\end{leansrc}
\par\medskip\noindent{\small \texttt{PPLC.}\texttt{\bfseries POrEData}}\hfill{\footnotesize \path{lean/qtsl/Qtsl/Core/Catalogue.lean}, line 1413}\par\nopagebreak
\begin{leansrc}
inductive POrEData : Hypernet → Sequent → Type where …
\end{leansrc}
\par\medskip\noindent{\small \texttt{PPLC.M3.}\texttt{\bfseries orE\_kit}}\hfill{\footnotesize \path{lean/qtsl/Qtsl/Admissible.lean}, line 149}\par\nopagebreak
\begin{leansrc}
theorem orE_kit : ∀ x y z : M3,
    le x z = true → le y z = true → le (sup x y) z = true
\end{leansrc}
\par\medskip\noindent{\small \texttt{PPLC.}\texttt{\bfseries substituteInFormula}}\hfill{\footnotesize \path{lean/qtsl/Qtsl/Core/Catalogue.lean}, line 581}\par\nopagebreak
\begin{leansrc}
def substituteInFormula (F : Formula) (w : Wire) (t : Term) (F' : Formula) : Prop
\end{leansrc}
\par\medskip\noindent{\small \texttt{PPLC.}\texttt{\bfseries sound\_eqElim}}\hfill{\footnotesize \path{lean/qtsl/Qtsl/Soundness.lean}, line 4269}\par\nopagebreak
\begin{leansrc}
theorem sound_eqElim {H : Hypernet} {seq eqPremise bodyPremise : Sequent}
    {wire : Wire} {termT termU : Term} {bodyForm : Formula}
    (ihEq : eqPremise.semanticallyValid H)
    (ihBody : bodyPremise.semanticallyValid H)
    (hpre : seq.pre = eqPremise.pre)
    (hpreMatch : eqPremise.pre = bodyPremise.pre)
    (heqPost : eqPremise.post = .atom .eq termT termU)
    (htyEq : ∃ s, termT.bshape H = some s ∧ termU.bshape H = some s)
    (hbodySubT : substituteInFormula bodyForm wire termT bodyPremise.post)
    (hconclSubU : substituteInFormula bodyForm wire termU seq.post) :
    seq.semanticallyValid H
\end{leansrc}
\par\medskip\noindent{\small \texttt{PPLC.Term.}\texttt{\bfseries evalQ\_subst\_congr}}\hfill{\footnotesize \path{lean/qtsl/Qtsl/Soundness.lean}, line 270}\par\nopagebreak
\begin{leansrc}
theorem Term.evalQ_subst_congr {H : Hypernet} (σ : H.Valuation)
    {w : Wire} {u₁ u₂ : Term} (hu : Term.evalQ σ u₁ = Term.evalQ σ u₂) :
    ∀ (t : Term),
      Term.evalQ σ (Term.subst w u₁ t) = Term.evalQ σ (Term.subst w u₂ t)
\end{leansrc}
\par\medskip\noindent{\small \texttt{PPLC.Formula.}\texttt{\bfseries holds\_subst\_congr}}\hfill{\footnotesize \path{lean/qtsl/Qtsl/Soundness.lean}, line 288}\par\nopagebreak
\begin{leansrc}
theorem Formula.holds_subst_congr {H : Hypernet} (σ : H.Valuation)
    {w : Wire} {u₁ u₂ : Term} (hu : Term.evalQ σ u₁ = Term.evalQ σ u₂) :
    ∀ (F : Formula), (Formula.holds σ (F.subst w u₁) ↔ Formula.holds σ (F.subst w
      u₂))
\end{leansrc}
\par\medskip\noindent{\small \texttt{PPLC.Term.}\texttt{\bfseries evalQ\_bshape}}\hfill{\footnotesize \path{lean/qtsl/Qtsl/Soundness.lean}, line 479}\par\nopagebreak
\begin{leansrc}
theorem Term.evalQ_bshape {H : Hypernet} (σ : H.Valuation) :
    ∀ (t : Term) (s : Shape) (q : QTensor),
    t.bshape H = some s → Term.evalQ σ t = some q → q.shape = s
\end{leansrc}
\par\medskip\noindent{\small \texttt{PPLC.Hypernet.}\texttt{\bfseries Valuation}}\hfill{\footnotesize \path{lean/qtsl/Qtsl/Catalogue.lean}, line 268}\par\nopagebreak
\begin{leansrc}
def Hypernet.Valuation (H : Hypernet) : Type
\end{leansrc}
\par\medskip\noindent{\small \texttt{PPLC.Term.}\texttt{\bfseries evalQ}}\hfill{\footnotesize \path{lean/qtsl/Qtsl/Catalogue.lean}, line 352}\par\nopagebreak
\begin{leansrc}
def Term.evalQ {H : Hypernet} (σ : H.Valuation) : Term → Option QTensor
  | .var w => σ.wireQ w
  | .const c => some ⟨[], [c], rfl⟩
  | .add l r => do QTensor.map2 (· + ·) (← Term.evalQ σ l) (← Term.evalQ σ r)
  | .mul l r => do QTensor.map2 (· * ·) (← Term.evalQ σ l) (← Term.evalQ σ r)
  | .sub l r => do QTensor.map2 (· - ·) (← Term.evalQ σ l) (← Term.evalQ σ r)
  | .neg t => (Term.evalQ σ t).map (QTensor.map1 Neg.neg)
  | .idx t i => do

    let q ← Term.evalQ σ t
    match q.shape with
    | [e] => if i < e then some ⟨[], [q.entries.getD i 0], rfl⟩ else none
    | _ => none

def BowTie.holdsQ : BowTie → ℚ → ℚ → Prop
  | .lt => (· < ·)
  | .le => (· ≤ ·)
  | .eq => (· = ·)
  | .ge => fun a b => b ≤ a
  | .gt => fun a b => b < a
  | .ne => (· ≠ ·)

def finSum (n : ℕ) (f : Fin n → ℝ) : ℝ
\end{leansrc}
\par\medskip\noindent{\small \texttt{PPLC.QTensor.}\texttt{\bfseries rel2}}\hfill{\footnotesize \path{lean/qtsl/Qtsl/Catalogue.lean}, line 332}\par\nopagebreak
\begin{leansrc}
def rel2 (r : ℚ → ℚ → Prop) (a b : QTensor) : Prop
\end{leansrc}
\par\medskip\noindent{\small \texttt{PPLC.QTensor.}\texttt{\bfseries eq\_of\_rel2}}\hfill{\footnotesize \path{lean/qtsl/Qtsl/Soundness.lean}, line 2431}\par\nopagebreak
\begin{leansrc}
theorem QTensor.eq_of_rel2 {r : ℚ → ℚ → Prop} {a b : QTensor}
    (hr : ∀ x y, r x y → x = y) (hs : a.shape = b.shape)
    (h : QTensor.rel2 r a b) : a = b
\end{leansrc}
\par\medskip\noindent{\small \texttt{PPLC.Formula.}\texttt{\bfseries subst}}\hfill{\footnotesize \path{lean/qtsl/Qtsl/Core/Catalogue.lean}, line 573}\par\nopagebreak
\begin{leansrc}
def Formula.subst (w : Wire) (t : Term) : Formula → Formula
  | .atom b l r => .atom b (Term.subst w t l) (Term.subst w t r)
  | .inRegion u simplex S => .inRegion (Term.subst w t u) simplex S
  | .and l r => .and (Formula.subst w t l) (Formula.subst w t r)
  | .or l r => .or (Formula.subst w t l) (Formula.subst w t r)
  | .tt => .tt
  | .ff => .ff

def substituteInFormula (F : Formula) (w : Wire) (t : Term) (F' : Formula) : Prop
\end{leansrc}
\par\medskip\noindent{\small \texttt{PPLC.Term.}\texttt{\bfseries bshape}}\hfill{\footnotesize \path{lean/qtsl/Qtsl/Core/Catalogue.lean}, line 517}\par\nopagebreak
\begin{leansrc}
def Term.bshape (H : Hypernet) : Term → Option Shape
  | .var w => (H.wireType w).map (·.shape)
  | .const _ => some []
  | .add l r | .mul l r | .sub l r =>
    match l.bshape H, r.bshape H with
    | some s₁, some s₂ =>
      if s₁ = s₂ then some s₁
      else if s₁ = [] then some s₂
      else if s₂ = [] then some s₁
      else none
    | _, _ => none
  | .neg t => t.bshape H
  | .idx t i =>

    match t.bshape H with
    | some [e] => if i < e then some [] else none
    | _ => none

def Term.finWires (H : Hypernet) (t : Term) : Prop
\end{leansrc}
\par\medskip\noindent{\small \texttt{PPLC.Term.}\texttt{\bfseries finWires}}\hfill{\footnotesize \path{lean/qtsl/Qtsl/Core/Catalogue.lean}, line 535}\par\nopagebreak
\begin{leansrc}
def Term.finWires (H : Hypernet) (t : Term) : Prop
\end{leansrc}
\par\medskip\noindent{\small \texttt{PPLC.Term.}\texttt{\bfseries evalQ\_defined\_of\_bshape\_fin}}\hfill{\footnotesize \path{lean/qtsl/Qtsl/Soundness.lean}, line 3488}\par\nopagebreak
\begin{leansrc}
theorem Term.evalQ_defined_of_bshape_fin {H : Hypernet} (σ : H.Valuation) :
    ∀ (t : Term) (s : Shape), t.bshape H = some s → t.finWires H →
    ∃ q, Term.evalQ σ t = some q ∧ q.shape = s
\end{leansrc}
\par\medskip\noindent{\small \texttt{PPLC.QTensor.}\texttt{\bfseries rel2\_refl}}\hfill{\footnotesize \path{lean/qtsl/Qtsl/Soundness.lean}, line 6773}\par\nopagebreak
\begin{leansrc}
theorem QTensor.rel2_refl {r : ℚ → ℚ → Prop} (hr : ∀ x, r x x)
    (a : QTensor) : QTensor.rel2 r a a
\end{leansrc}
\par\medskip\noindent{\small \texttt{PPLC.}\texttt{\bfseries sound\_eqRefl}}\hfill{\footnotesize \path{lean/qtsl/Qtsl/Soundness.lean}, line 6853}\par\nopagebreak
\begin{leansrc}
theorem sound_eqRefl {H : Hypernet} {seq : Sequent} {t : Term}
    (htyDef : ∃ s, t.bshape H = some s) (hfin : t.finWires H)
    (hpost : seq.post = .atom .eq t t) :
    seq.semanticallyValid H
\end{leansrc}
\par\medskip\noindent{\small \texttt{PPLC.Hypernet.Valuation.}\texttt{\bfseries wireQ\_isSome}}\hfill{\footnotesize \path{lean/qtsl/Qtsl/Soundness.lean}, line 70}\par\nopagebreak
\begin{leansrc}
theorem Hypernet.Valuation.wireQ_isSome {H : Hypernet} (σ : H.Valuation)
    {w : Wire} {τ : TensorType} (hτ : H.wireType w = some τ)
    (hfin : τ.dtype.ext = Ext.fin) : ∃ q, σ.wireQ w = some q
\end{leansrc}
\par\medskip\noindent{\small \texttt{PPLC.QTensor.}\texttt{\bfseries rel2\_symm}}\hfill{\footnotesize \path{lean/qtsl/Qtsl/Soundness.lean}, line 6789}\par\nopagebreak
\begin{leansrc}
theorem QTensor.rel2_symm {r : ℚ → ℚ → Prop}
    (hr : ∀ x y, r x y → r y x) {a b : QTensor}
    (h : QTensor.rel2 r a b) : QTensor.rel2 r b a
\end{leansrc}
\par\medskip\noindent{\small \texttt{PPLC.}\texttt{\bfseries sound\_eqSym}}\hfill{\footnotesize \path{lean/qtsl/Qtsl/Soundness.lean}, line 6863}\par\nopagebreak
\begin{leansrc}
theorem sound_eqSym {H : Hypernet} {seq premise : Sequent} {tT tU : Term}
    (ih : premise.semanticallyValid H)
    (hpre : seq.pre = premise.pre)
    (hpremPost : premise.post = .atom .eq tT tU)
    (hpost : seq.post = .atom .eq tU tT) :
    seq.semanticallyValid H
\end{leansrc}
\par\medskip\noindent{\small \texttt{PPLC.}\texttt{\bfseries PEqSymData}}\hfill{\footnotesize \path{lean/qtsl/Qtsl/Core/Catalogue.lean}, line 1443}\par\nopagebreak
\begin{leansrc}
inductive PEqSymData : Hypernet → Sequent → Type where …
\end{leansrc}
\par\medskip\noindent{\small \texttt{PPLC.QTensor.}\texttt{\bfseries rel2\_conj}}\hfill{\footnotesize \path{lean/qtsl/Qtsl/Soundness.lean}, line 6834}\par\nopagebreak
\begin{leansrc}
theorem QTensor.rel2_conj {r s t : ℚ → ℚ → Prop}
    (hrst : ∀ x y, r x y → s x y → t x y) {a b : QTensor}
    (hr : QTensor.rel2 r a b) (hs : QTensor.rel2 s a b) :
    QTensor.rel2 t a b
\end{leansrc}
\par\medskip\noindent{\small \texttt{PPLC.}\texttt{\bfseries sound\_antisym}}\hfill{\footnotesize \path{lean/qtsl/Qtsl/Soundness.lean}, line 6888}\par\nopagebreak
\begin{leansrc}
theorem sound_antisym {H : Hypernet} {seq premiseLe premiseGe : Sequent}
    {tT tU : Term}
    (ihLe : premiseLe.semanticallyValid H)
    (ihGe : premiseGe.semanticallyValid H)
    (hpreLe : seq.pre = premiseLe.pre)
    (hpreGe : seq.pre = premiseGe.pre)
    (hpremLePost : premiseLe.post = .atom .le tT tU)
    (hpremGePost : premiseGe.post = .atom .ge tT tU)
    (hpost : seq.post = .atom .eq tT tU) :
    seq.semanticallyValid H
\end{leansrc}
\par\medskip\noindent{\small \texttt{PPLC.}\texttt{\bfseries baryCoeffs1D}}\hfill{\footnotesize \path{lean/qtsl/Qtsl/Core/Catalogue.lean}, line 584}\par\nopagebreak
\begin{leansrc}
def baryCoeffs1D (cx : SimpC) (dContains : Fin cx.nVertices → Bool) :
    Option (List (Fin cx.nVertices × Rat × Rat))
\end{leansrc}
\par\medskip\noindent{\small \texttt{PPLC.}\texttt{\bfseries primAffine}}\hfill{\footnotesize \path{lean/qtsl/Qtsl/Core/Catalogue.lean}, line 599}\par\nopagebreak
\begin{leansrc}
def primAffine (s c : Rat) : Rat × Rat
\end{leansrc}
\par\medskip\noindent{\small \texttt{PPLC.}\texttt{\bfseries f4VertexAtom}}\hfill{\footnotesize \path{lean/qtsl/Qtsl/Core/Catalogue.lean}, line 612}\par\nopagebreak
\begin{leansrc}
def f4VertexAtom (inS : Bool) (off slp : Rat) (t : Term) : Formula
\end{leansrc}
\par\medskip\noindent{\small \texttt{PPLC.}\texttt{\bfseries barycentricHalfspaceBuild}}\hfill{\footnotesize \path{lean/qtsl/Qtsl/Core/Catalogue.lean}, line 616}\par\nopagebreak
\begin{leansrc}
def barycentricHalfspaceBuild (t : Term) (simplex : Simplex)
    (S : ExtensionSet) : Option Formula
\end{leansrc}
\par\medskip\noindent{\small \texttt{PPLC.}\texttt{\bfseries barycentricHalfspaceForm}}\hfill{\footnotesize \path{lean/qtsl/Qtsl/Core/Catalogue.lean}, line 625}\par\nopagebreak
\begin{leansrc}
def barycentricHalfspaceForm (t : Term) (simplex : Simplex)
    (S : ExtensionSet) (F : Formula) : Prop
\end{leansrc}
\par\medskip\noindent{\small \texttt{PPLC.}\texttt{\bfseries barycentricHalfspaceBuild\_spec}}\hfill{\footnotesize \path{lean/qtsl/Qtsl/Soundness.lean}, line 634}\par\nopagebreak
\begin{leansrc}
theorem barycentricHalfspaceBuild_spec {t : Term} {simplex : Simplex}
    {S : ExtensionSet} {F : Formula}
    (h : barycentricHalfspaceBuild t simplex S = some F) :
    S.simplex = simplex ∧ simplex.complex.ambDim = 1 ∧
    ∃ (a b : Fin simplex.complex.nVertices) (V₀ V₁ : ℚ),
      (List.finRange simplex.complex.nVertices).filter
        (fun v => simplex.contains v) = [a, b] ∧
      (∀ j, simplex.complex.coord a j = V₀) ∧
      (∀ j, simplex.complex.coord b j = V₁) ∧
      V₀ ≠ V₁ ∧
      F = .and
        (f4VertexAtom (S.containsNat a.val)
          (V₁ / (V₁ - V₀)) (-1 / (V₁ - V₀)) t)
        (f4VertexAtom (S.containsNat b.val)
          (-V₀ / (V₁ - V₀)) (1 / (V₁ - V₀)) t)
\end{leansrc}
\par\medskip\noindent{\small \texttt{PPLC.}\texttt{\bfseries sound\_convElim}}\hfill{\footnotesize \path{lean/qtsl/Qtsl/Soundness.lean}, line 4672}\par\nopagebreak
\begin{leansrc}
theorem sound_convElim {H : Hypernet} {seq premise : Sequent}
    {term : Term} {simplex : Simplex} {S : ExtensionSet}
    (ih : premise.semanticallyValid H)
    (hpre : seq.pre = premise.pre)
    (hpremiseForm : premise.post = .inRegion term simplex S)
    (hhalf : barycentricHalfspaceForm term simplex S seq.post) :
    seq.semanticallyValid H
\end{leansrc}
\par\medskip\noindent{\small \texttt{PPLC.}\texttt{\bfseries f4VertexAtom\_holds\_iff}}\hfill{\footnotesize \path{lean/qtsl/Qtsl/Soundness.lean}, line 1393}\par\nopagebreak
\begin{leansrc}
theorem f4VertexAtom_holds_iff {H : Hypernet} {σ : H.Valuation} {t : Term}
    {q : QTensor} (hq : Term.evalQ σ t = some q) (hsq : SingletonShape q.shape)
    (inS : Bool) (off slp : ℚ) :
    Formula.holds σ (f4VertexAtom inS off slp t) ↔
      (if inS then off + slp * q.scalar ≤ 0
       else 0 ≤ off + slp * q.scalar)
\end{leansrc}
\par\medskip\noindent{\small \texttt{PPLC.}\texttt{\bfseries maskRegion\_pair\_elim}}\hfill{\footnotesize \path{lean/qtsl/Qtsl/Soundness.lean}, line 2526}\par\nopagebreak
\begin{leansrc}
theorem maskRegion_pair_elim {C : SimpC} {mask ext : Fin C.nVertices → Bool}
    {x : Fin C.ambDim → ℚ} {a b : Fin C.nVertices} {V₀ V₁ : ℚ}
    (hfil : (List.finRange C.nVertices).filter (fun v => mask v) = [a, b])
    (hVa : ∀ j, C.coord a j = V₀) (hVb : ∀ j, C.coord b j = V₁)
    (hne : V₀ ≠ V₁) (j0 : Fin C.ambDim)
    (h : MaskRegion C mask ext x) :
    (ext a = true → (V₁ - x j0) / (V₁ - V₀) ≤ 0) ∧
    (ext a = false → 0 ≤ (V₁ - x j0) / (V₁ - V₀)) ∧
    (ext b = true → (x j0 - V₀) / (V₁ - V₀) ≤ 0) ∧
    (ext b = false → 0 ≤ (x j0 - V₀) / (V₁ - V₀))
\end{leansrc}
\par\medskip\noindent{\small \texttt{PPLC.}\texttt{\bfseries sound\_convIntro}}\hfill{\footnotesize \path{lean/qtsl/Qtsl/Soundness.lean}, line 4593}\par\nopagebreak
\begin{leansrc}
theorem sound_convIntro {H : Hypernet} {seq premise : Sequent}
    {term : Term} {simplex : Simplex} {S : ExtensionSet}
    (ih : premise.semanticallyValid H)
    (hpre : seq.pre = premise.pre)
    (htermShape : term.bshape H = some [simplex.complex.ambDim])
    (hhalf : barycentricHalfspaceForm term simplex S premise.post)
    (hpost : seq.post = .inRegion term simplex S) :
    seq.semanticallyValid H
\end{leansrc}
\par\medskip\noindent{\small \texttt{PPLC.}\texttt{\bfseries f4VertexAtom\_holds\_inv}}\hfill{\footnotesize \path{lean/qtsl/Qtsl/Soundness.lean}, line 1442}\par\nopagebreak
\begin{leansrc}
theorem f4VertexAtom_holds_inv {H : Hypernet} {σ : H.Valuation} {t : Term}
    {inS : Bool} {off slp : ℚ}
    (h : Formula.holds σ (f4VertexAtom inS off slp t)) :
    ∃ q, Term.evalQ σ t = some q
\end{leansrc}
\par\medskip\noindent{\small \texttt{PPLC.}\texttt{\bfseries maskRegion\_pair\_intro}}\hfill{\footnotesize \path{lean/qtsl/Qtsl/Soundness.lean}, line 2575}\par\nopagebreak
\begin{leansrc}
theorem maskRegion_pair_intro {C : SimpC} {mask ext : Fin C.nVertices → Bool}
    {x : Fin C.ambDim → ℚ} {a b : Fin C.nVertices} {V₀ V₁ : ℚ}
    (hfil : (List.finRange C.nVertices).filter (fun v => mask v) = [a, b])
    (hVa : ∀ j, C.coord a j = V₀) (hVb : ∀ j, C.coord b j = V₁)
    (hne : V₀ ≠ V₁) (hdim : C.ambDim = 1) (j0 : Fin C.ambDim)
    (hw0t : ext a = true → (V₁ - x j0) / (V₁ - V₀) ≤ 0)
    (hw0f : ext a = false → 0 ≤ (V₁ - x j0) / (V₁ - V₀))
    (hw1t : ext b = true → (x j0 - V₀) / (V₁ - V₀) ≤ 0)
    (hw1f : ext b = false → 0 ≤ (x j0 - V₀) / (V₁ - V₀)) :
    MaskRegion C mask ext x
\end{leansrc}
\par\medskip\noindent{\small \texttt{PPLC.}\texttt{\bfseries sound\_arithAtom\_fold}}\hfill{\footnotesize \path{lean/qtsl/Qtsl/Soundness.lean}, line 4297}\par\nopagebreak
\begin{leansrc}
theorem sound_arithAtom_fold {H : Hypernet} {seq premise : Sequent}
    {termL termR : Term} {b : BowTie}
    (ih : premise.semanticallyValid H)
    (hpre : seq.pre = premise.pre)
    (hprem : atomicCompareZero b termL termR premise.post)
    (hconcl : atomicCompareForm b termL termR seq.post) :
    seq.semanticallyValid H
\end{leansrc}
\par\medskip\noindent{\small \texttt{PPLC.}\texttt{\bfseries sound\_arithAtom\_unfold}}\hfill{\footnotesize \path{lean/qtsl/Qtsl/Soundness.lean}, line 4318}\par\nopagebreak
\begin{leansrc}
theorem sound_arithAtom_unfold {H : Hypernet} {seq premise : Sequent}
    {termL termR : Term} {b : BowTie}
    (ih : premise.semanticallyValid H)
    (hpre : seq.pre = premise.pre)
    (hprem : atomicCompareForm b termL termR premise.post)
    (hconcl : atomicCompareZero b termL termR seq.post) :
    seq.semanticallyValid H
\end{leansrc}
\par\medskip\noindent{\small \texttt{PPLC.}\texttt{\bfseries atomicCompareForm}}\hfill{\footnotesize \path{lean/qtsl/Qtsl/Core/Catalogue.lean}, line 553}\par\nopagebreak
\begin{leansrc}
def atomicCompareForm (b : BowTie) (l r : Term) (F : Formula) : Prop
\end{leansrc}
\par\medskip\noindent{\small \texttt{PPLC.}\texttt{\bfseries atomicCompareZero}}\hfill{\footnotesize \path{lean/qtsl/Qtsl/Core/Catalogue.lean}, line 556}\par\nopagebreak
\begin{leansrc}
def atomicCompareZero (b : BowTie) (l r : Term) (F : Formula) : Prop
\end{leansrc}
\par\medskip\noindent{\small \texttt{PPLC.QTensor.}\texttt{\bfseries map2\_defined\_of\_rel2}}\hfill{\footnotesize \path{lean/qtsl/Qtsl/Soundness.lean}, line 3462}\par\nopagebreak
\begin{leansrc}
theorem QTensor.map2_defined_of_rel2 {r : ℚ → ℚ → Prop} {a b : QTensor}
    (h : QTensor.rel2 r a b) (f : ℚ → ℚ → ℚ) :
    ∃ c, QTensor.map2 f a b = some c
\end{leansrc}
\par\medskip\noindent{\small \texttt{PPLC.}\texttt{\bfseries BowTie}}\hfill{\footnotesize \path{lean/qtsl/Qtsl/Core/Catalogue.lean}, line 346}\par\nopagebreak
\begin{leansrc}
inductive BowTie where …
\end{leansrc}
\par\medskip\noindent{\small \texttt{PPLC.QTensor.}\texttt{\bfseries rel2\_sub\_zero\_iff}}\hfill{\footnotesize \path{lean/qtsl/Qtsl/Soundness.lean}, line 2466}\par\nopagebreak
\begin{leansrc}
theorem QTensor.rel2_sub_zero_iff {b : BowTie} {al ar d : QTensor}
    (hd : QTensor.map2 (· - ·) al ar = some d) :
    QTensor.rel2 (BowTie.holdsQ b) d ⟨[], [0], rfl⟩
      ↔ QTensor.rel2 (BowTie.holdsQ b) al ar
\end{leansrc}
\par\medskip\noindent{\small \texttt{PPLC.Term.}\texttt{\bfseries evalQ\_sub\_inv}}\hfill{\footnotesize \path{lean/qtsl/Qtsl/Soundness.lean}, line 1306}\par\nopagebreak
\begin{leansrc}
theorem Term.evalQ_sub_inv {H : Hypernet} (σ : H.Valuation) {l r : Term}
    {q : QTensor} (h : Term.evalQ σ (.sub l r) = some q) :
    ∃ a b, Term.evalQ σ l = some a ∧ Term.evalQ σ r = some b ∧
      QTensor.map2 (· - ·) a b = some q
\end{leansrc}
\par\medskip\noindent{\small \texttt{PPLC.}\texttt{\bfseries PArithAtomData}}\hfill{\footnotesize \path{lean/qtsl/Qtsl/Core/Catalogue.lean}, line 1471}\par\nopagebreak
\begin{leansrc}
inductive PArithAtomData : Hypernet → Sequent → Type where …
\end{leansrc}
\par\medskip\noindent{\small \texttt{PPLC.}\texttt{\bfseries allForwardDiffBernCoeffsSatisfyP}}\hfill{\footnotesize \path{lean/qtsl/Qtsl/Core/Catalogue.lean}, line 760}\par\nopagebreak
\begin{leansrc}
def allForwardDiffBernCoeffsSatisfyP
    (simplex : Simplex) (S : ExtensionSet) (degree : Nat)
    (polynomial : Term) (predicate : Formula) : Prop
\end{leansrc}
\par\medskip\noindent{\small \texttt{PPLC.}\texttt{\bfseries satisfiesPredQ}}\hfill{\footnotesize \path{lean/qtsl/Qtsl/Core/Catalogue.lean}, line 701}\par\nopagebreak
\begin{leansrc}
def satisfiesPredQ (P : Formula) (c : Rat) : Bool
\end{leansrc}
\par\medskip\noindent{\small \texttt{PPLC.}\texttt{\bfseries bernConeCoeffCheck}}\hfill{\footnotesize \path{lean/qtsl/Qtsl/Core/Catalogue.lean}, line 727}\par\nopagebreak
\begin{leansrc}
def bernConeCoeffCheck (simplex : Simplex) (S : ExtensionSet) (degree : Nat)
    (polynomial : Term) (predicate : Formula) : Bool
\end{leansrc}
\par\medskip\noindent{\small \texttt{PPLC.}\texttt{\bfseries sound\_bernCone}}\hfill{\footnotesize \path{lean/qtsl/Qtsl/Soundness.lean}, line 5498}\par\nopagebreak
\begin{leansrc}
theorem sound_bernCone {H : Hypernet} {seq regionPremise : Sequent}
    {termX : Term} {simplex : Simplex} {S : ExtensionSet} {degree : Nat}
    {polynomial : Term} {predicate : Formula}
    (ih : regionPremise.semanticallyValid H)
    (hpre : seq.pre = regionPremise.pre)
    (hregionPost : regionPremise.post = .inRegion termX simplex S)
    (hpolyTie : polynomial.freeWires = [] ∨
      ∃ w, termX = Term.var w ∧ polynomial.freeWires ⊆ [w])
    (hpredShape : (¬ isExtensionSetEmpty S → isConicPredicate predicate)
      ∧ (isExtensionSetEmpty S → isConvexPredicate predicate))
    (hcoeff : allForwardDiffBernCoeffsSatisfyP simplex S degree polynomial
      predicate)
    (hpostApply : applyPredicateToTerm predicate polynomial seq.post) :
    seq.semanticallyValid H
\end{leansrc}
\par\medskip\noindent{\small \texttt{PPLC.}\texttt{\bfseries sound\_bernCone\_empty}}\hfill{\footnotesize \path{lean/qtsl/Qtsl/Soundness.lean}, line 5058}\par\nopagebreak
\begin{leansrc}
theorem sound_bernCone_empty {H : Hypernet} {seq regionPremise : Sequent}
    {termX : Term} {simplex : Simplex} {S : ExtensionSet} {degree : Nat}
    {polynomial : Term} {predicate : Formula}
    (ih : regionPremise.semanticallyValid H)
    (hpre : seq.pre = regionPremise.pre)
    (hregionPost : regionPremise.post = .inRegion termX simplex S)
    (hpolyTie : polynomial.freeWires = [] ∨
      ∃ w, termX = Term.var w ∧ polynomial.freeWires ⊆ [w])
    (hconvex : isExtensionSetEmpty S → isConvexPredicate predicate)
    (hSempty : isExtensionSetEmpty S)
    (hcoeff : allForwardDiffBernCoeffsSatisfyP simplex S degree polynomial
      predicate)
    (hpostApply : applyPredicateToTerm predicate polynomial seq.post) :
    seq.semanticallyValid H
\end{leansrc}
\par\medskip\noindent{\small \texttt{PPLC.}\texttt{\bfseries MaskRegion}}\hfill{\footnotesize \path{lean/qtsl/Qtsl/Catalogue.lean}, line 376}\par\nopagebreak
\begin{leansrc}
def MaskRegion (C : SimpC) (mask ext : Fin C.nVertices → Bool)
    (x : Fin C.ambDim → ℚ) : Prop
\end{leansrc}
\par\medskip\noindent{\small \texttt{PPLC.Term.}\texttt{\bfseries evalQ\_brRead}}\hfill{\footnotesize \path{lean/qtsl/Qtsl/Soundness.lean}, line 1578}\par\nopagebreak
\begin{leansrc}
theorem Term.evalQ_brRead {H : Hypernet} (σ : H.Valuation) (s : Shape) :
    ∀ (t : Term),
    (∀ w ∈ t.freeWires, ∃ qw, σ.wireQ w = some qw ∧
      (qw.shape = [] ∨ qw.shape = s)) →
    t.idxFree = true →
    ∃ q, Term.evalQ σ t = some q ∧ (q.shape = [] ∨ q.shape = s) ∧
      (q.shape = [] → ∀ w ∈ t.freeWires, ∀ qw, σ.wireQ w = some qw →
        qw.shape = []) ∧
      ∀ i, i < Shape.totalSize s →
        q.brRead i = QArith.Poly.eval (σ.entryVal i) t.toPoly
\end{leansrc}
\par\medskip\noindent{\small \texttt{PPLC.}\texttt{\bfseries maskRegion\_rat\_weights}}\hfill{\footnotesize \path{lean/qtsl/Qtsl/Soundness.lean}, line 4798}\par\nopagebreak
\begin{leansrc}
theorem maskRegion_rat_weights {C : SimpC}
    {mask ext : Fin C.nVertices → Bool} {x : Fin C.ambDim → ℚ}
    (hmask : MaskRegion C mask ext x)
    (hext0 : ∀ v, ext v = false) :
    ∃ qw : Fin ((List.finRange C.nVertices).filter
        (fun v => mask v)).length → ℚ,
      (∀ i, 0 ≤ qw i) ∧ (∑ i, qw i) = 1 ∧
      ∀ j : Fin C.ambDim,
        (∑ i, qw i * C.coord (((List.finRange C.nVertices).filter
          (fun v => mask v)).get i) j) = x j
\end{leansrc}
\par\medskip\noindent{\small \texttt{PPLC.ConvexRat.}\texttt{\bfseries exists\_nonneg\_rat\_solution}}\hfill{\footnotesize \path{lean/qtsl/Qtsl/ConvexRat.lean}, line 269}\par\nopagebreak
\begin{leansrc}
theorem exists_nonneg_rat_solution (A : Fin k → Fin m → ℚ)
    (b : Fin m → ℚ) (c : Fin k → ℝ) (hc0 : ∀ i, 0 ≤ c i)
    (hcsol : ∀ j, (∑ i, c i * (A i j : ℝ)) = (b j : ℝ)) :
    ∃ q : Fin k → ℚ, (∀ i, 0 ≤ q i) ∧ ∀ j, (∑ i, q i * A i j) = b j
\end{leansrc}
\par\medskip\noindent{\small \texttt{PPLC.}\texttt{\bfseries classifyPred}}\hfill{\footnotesize \path{lean/qtsl/Qtsl/Core/Catalogue.lean}, line 667}\par\nopagebreak
\begin{leansrc}
def classifyPred : Formula → Option PredClass
  | .and l r => do
    let a ← classifyPred l
    let b ← classifyPred r
    pure ⟨a.conic && b.conic, a.upClosed && b.upClosed⟩
  | f => classifyPredAtom f

def isConicPredicate (P : Formula) : Prop
\end{leansrc}
\par\medskip\noindent{\small \texttt{PPLC.}\texttt{\bfseries satisfiesPredQ\_convex}}\hfill{\footnotesize \path{lean/qtsl/Qtsl/Soundness.lean}, line 2871}\par\nopagebreak
\begin{leansrc}
theorem satisfiesPredQ_convex {P : Formula} {cls : PredClass}
    (hP : classifyPred P = some cls) :
    ∀ a₁ b₁ l : ℚ, 0 ≤ l → l ≤ 1 →
      satisfiesPredQ P a₁ = true → satisfiesPredQ P b₁ = true →
      satisfiesPredQ P (l * a₁ + (1 - l) * b₁) = true
\end{leansrc}
\par\medskip\noindent{\small \texttt{PPLC.QArith.}\texttt{\bfseries fdBernstein\_sound\_dim1}}\hfill{\footnotesize \path{lean/qtsl/Qtsl/QArithKernel.lean}, line 2219}\par\nopagebreak
\begin{leansrc}
theorem fdBernstein_sound_dim1
    (verts : List (List ℚ)) (smask : List Bool) (d : Nat) (p : Poly)
    (t : FDTable) (pred : ℚ → Bool) (u₀ : ℚ)
    (hmask : smask = [true, false] ∨ smask = [false, true])
    (hfd : fdBernstein verts smask d p = some t)
    (hgrid : gridAll pred t = true)
    (hu : 0 ≤ u₀)
    (hzero : pred 0 = true)
    (hscale : ∀ a c : ℚ, 0 ≤ c → pred a = true → pred (c * a) = true)
    (hconv : ∀ a b l : ℚ, 0 ≤ l → l ≤ 1 → pred a = true → pred b = true →
      pred (l * a + (1 - l) * b) = true) :
    pred (Poly.eval (fun j =>
      ((List.range verts.length).map
        fun i => (verts.getD i []).getD j 0 *
          (if smask.getD i false then -u₀ else 1 + u₀)).sum) p) = true
\end{leansrc}
\par\medskip\noindent{\small \texttt{PPLC.QArith.}\texttt{\bfseries fdBernstein\_sound\_const}}\hfill{\footnotesize \path{lean/qtsl/Qtsl/QArithKernel.lean}, line 1733}\par\nopagebreak
\begin{leansrc}
theorem fdBernstein_sound_const
    (verts : List (List ℚ)) (smask : List Bool) (d : Nat) (c₀ : ℚ)
    (t : FDTable) (pred : ℚ → Bool)
    (hfd : fdBernstein verts smask d (Poly.const c₀) = some t)
    (hgrid : gridAll pred t = true) :
    pred c₀ = true
\end{leansrc}
\par\medskip\noindent{\small \texttt{PPLC.}\texttt{\bfseries satisfiesPredQ\_conic\_zero}}\hfill{\footnotesize \path{lean/qtsl/Qtsl/Soundness.lean}, line 3066}\par\nopagebreak
\begin{leansrc}
theorem satisfiesPredQ_conic_zero {P : Formula} {cls : PredClass}
    (hP : classifyPred P = some cls) (hc : cls.conic = true) :
    satisfiesPredQ P 0 = true
\end{leansrc}
\par\medskip\noindent{\small \texttt{PPLC.}\texttt{\bfseries satisfiesPredQ\_conic\_scale}}\hfill{\footnotesize \path{lean/qtsl/Qtsl/Soundness.lean}, line 3102}\par\nopagebreak
\begin{leansrc}
theorem satisfiesPredQ_conic_scale {P : Formula} {cls : PredClass}
    (hP : classifyPred P = some cls) (hc : cls.conic = true) :
    ∀ a c : ℚ, 0 ≤ c → satisfiesPredQ P a = true →
      satisfiesPredQ P (c * a) = true
\end{leansrc}
\par\medskip\noindent{\small \texttt{PPLC.QArith.}\texttt{\bfseries fdBernstein}}\hfill{\footnotesize \path{lean/qtsl/Qtsl/Core/QArith.lean}, line 158}\par\nopagebreak
\begin{leansrc}
def fdBernstein (verts : List (List Rat)) (smask : List Bool) (d : Nat)
    (p : Poly) : Option FDTable
\end{leansrc}
\par\medskip\noindent{\small \texttt{PPLC.QArith.}\texttt{\bfseries gridAll}}\hfill{\footnotesize \path{lean/qtsl/Qtsl/Core/QArith.lean}, line 190}\par\nopagebreak
\begin{leansrc}
def gridAll (pred : Rat → Bool) (t : FDTable) : Bool
\end{leansrc}
\par\medskip\noindent{\small \texttt{PPLC.Term.}\texttt{\bfseries idxBounded}}\hfill{\footnotesize \path{lean/qtsl/Qtsl/Core/Catalogue.lean}, line 478}\par\nopagebreak
\begin{leansrc}
def Term.idxBounded (n : Nat) : Term → Bool
  | .var _ | .const _ => true
  | .add l r | .mul l r | .sub l r => l.idxBounded n && r.idxBounded n
  | .neg t => t.idxBounded n
  | .idx t i => decide (i < n) && t.idxBounded n

def Formula.freeWires : Formula → List Wire
  | .atom _ l r => l.freeWires ++ r.freeWires
  | .inRegion t _ _ => t.freeWires
  | .and l r | .or l r => l.freeWires ++ r.freeWires
  | .tt | .ff => []

def Term.type (H : Hypernet) : Term → Option TensorType
  | .var w => H.wireType w
  | .const _ => some ⟨[], Dtype.Q⟩
  | .idx t i =>

    match Term.type H t with
    | some τ =>
      match τ.shape with
      | [e] => if i < e then some ⟨[], τ.dtype⟩ else none
      | _ => none
    | none => none
  | .add l r | .mul l r | .sub l r =>
    match Term.type H l, Term.type H r with
    | some τ₁, some τ₂ => if τ₁ = τ₂ then some τ₁ else none
    | _, _ => none
  | .neg t => Term.type H t

def Formula.WellTyped (H : Hypernet) : Formula → Prop
  | .atom _ l r => ∃ τ, l.type H = some τ ∧ r.type H = some τ
  | .inRegion t simplex S =>
      (∃ τ, t.type H = some τ ∧ τ.shape = [simplex.complex.ambDim]) ∧
      S.simplex = simplex
  | .and l r => l.WellTyped H ∧ r.WellTyped H
  | .or l r  => l.WellTyped H ∧ r.WellTyped H
  | .tt      => True
  | .ff      => True

def Term.bshape (H : Hypernet) : Term → Option Shape
  …
\end{leansrc}
\par\medskip\noindent{\small \texttt{PPLC.}\texttt{\bfseries applyPredicateToTerm}}\hfill{\footnotesize \path{lean/qtsl/Qtsl/Core/Catalogue.lean}, line 631}\par\nopagebreak
\begin{leansrc}
def applyPredicateToTerm (P : Formula) (arg : Term) (F : Formula) : Prop
\end{leansrc}
\par\medskip\noindent{\small \texttt{PPLC.}\texttt{\bfseries holds\_substHole\_of\_satisfies}}\hfill{\footnotesize \path{lean/qtsl/Qtsl/Soundness.lean}, line 3144}\par\nopagebreak
\begin{leansrc}
theorem holds_substHole_of_satisfies {H : Hypernet} (σ : H.Valuation)
    {P : Formula} {cls : PredClass} (hP : classifyPred P = some cls)
    {arg : Term} {qa : QTensor}
    (ha : Term.evalQ σ arg = some qa) (hsa : SingletonShape qa.shape)
    (hsat : satisfiesPredQ P qa.scalar = true) :
    Formula.holds σ (Formula.subst predicateHole arg P)
\end{leansrc}
\par\medskip\noindent{\small \texttt{PPLC.ConeLemma.}\texttt{\bfseries StarIdentity}}\hfill{\footnotesize \path{lean/qtsl/Qtsl/ConeLemma.lean}, line 18}\par\nopagebreak
\begin{leansrc}
def StarIdentity : Prop
\end{leansrc}
\par\medskip\noindent{\small \texttt{PPLC.}\texttt{\bfseries PNormData}}\hfill{\footnotesize \path{lean/qtsl/Qtsl/Core/Catalogue.lean}, line 1505}\par\nopagebreak
\begin{leansrc}
inductive PNormData : Hypernet → Sequent → Type where …
\end{leansrc}
\par\medskip\noindent{\small \texttt{PPLC.}\texttt{\bfseries sound\_norm}}\hfill{\footnotesize \path{lean/qtsl/Qtsl/Soundness.lean}, line 4494}\par\nopagebreak
\begin{leansrc}
theorem sound_norm {H : Hypernet} {seq : Sequent} {termL termR : Term}
    (hcanon : polynomialCanonicalEq termL termR)
    (htyEq : termL.bshape H = termR.bshape H)
    (htyDef : ∃ s, termL.bshape H = some s)
    (hfinL : termL.finWires H) (hfinR : termR.finWires H)
    (hpost : seq.post = .atom .eq termL termR) :
    seq.semanticallyValid H
\end{leansrc}
\par\medskip\noindent{\small \texttt{PPLC.}\texttt{\bfseries polynomialCanonicalEq}}\hfill{\footnotesize \path{lean/qtsl/Qtsl/Core/Catalogue.lean}, line 698}\par\nopagebreak
\begin{leansrc}
def polynomialCanonicalEq (l r : Term) : Prop
\end{leansrc}
\par\medskip\noindent{\small \texttt{PPLC.Term.}\texttt{\bfseries toPoly}}\hfill{\footnotesize \path{lean/qtsl/Qtsl/Core/Catalogue.lean}, line 414}\par\nopagebreak
\begin{leansrc}
def Term.toPoly : Term → QArith.Poly
  | .var w => QArith.Poly.var w
  | .const c => QArith.Poly.const c
  | .add l r => QArith.Poly.add l.toPoly r.toPoly
  | .mul l r => QArith.Poly.mul l.toPoly r.toPoly
  | .sub l r => QArith.Poly.sub l.toPoly r.toPoly
  | .neg t => QArith.Poly.neg t.toPoly

  | .idx _ _ => QArith.Poly.const 0

def Term.evalQ? : Term → Option Rat
  | .var _ => none
  | .const c => some c
  | .add l r => do pure ((← l.evalQ?) + (← r.evalQ?))
  | .mul l r => do pure ((← l.evalQ?) * (← r.evalQ?))
  | .sub l r => do pure ((← l.evalQ?) - (← r.evalQ?))
  | .neg t => do pure (-(← t.evalQ?))
  | .idx _ _ => none

def BowTie.evalQ : BowTie → Rat → Rat → Bool
  | .lt, x, y => decide (x < y)
  | .le, x, y => decide (x ≤ y)
  | .eq, x, y => decide (x = y)
  | .ge, x, y => decide (y ≤ x)
  | .gt, x, y => decide (y < x)
  | .ne, x, y => decide (¬ x = y)

def Formula.evalClosed : Formula → Option Bool
  | .atom b l r => do pure (b.evalQ (← l.evalQ?) (← r.evalQ?))
  | .inRegion _ _ _ => none
  | .and l r => do pure ((← l.evalClosed) && (← r.evalClosed))
  | .or l r => do pure ((← l.evalClosed) || (← r.evalClosed))
  | .tt => some true
  | .ff => some false

def baryLambdas1D (cx : SimpC) (dContains : Fin cx.nVertices → Bool)
    (t : Term) : Option (List (Fin cx.nVertices × Term))
\end{leansrc}
\par\medskip\noindent{\small \texttt{PPLC.Term.}\texttt{\bfseries idxFree}}\hfill{\footnotesize \path{lean/qtsl/Qtsl/Core/Catalogue.lean}, line 472}\par\nopagebreak
\begin{leansrc}
def Term.idxFree : Term → Bool
  | .var _ | .const _ => true
  | .add l r | .mul l r | .sub l r => l.idxFree && r.idxFree
  | .neg t => t.idxFree
  | .idx _ _ => false

def Term.idxBounded (n : Nat) : Term → Bool
  | .var _ | .const _ => true
  | .add l r | .mul l r | .sub l r => l.idxBounded n && r.idxBounded n
  | .neg t => t.idxBounded n
  | .idx t i => decide (i < n) && t.idxBounded n

def Formula.freeWires : Formula → List Wire
  | .atom _ l r => l.freeWires ++ r.freeWires
  | .inRegion t _ _ => t.freeWires
  | .and l r | .or l r => l.freeWires ++ r.freeWires
  | .tt | .ff => []

def Term.type (H : Hypernet) : Term → Option TensorType
  | .var w => H.wireType w
  | .const _ => some ⟨[], Dtype.Q⟩
  | .idx t i =>

    match Term.type H t with
    | some τ =>
      match τ.shape with
      | [e] => if i < e then some ⟨[], τ.dtype⟩ else none
      | _ => none
    | none => none
  | .add l r | .mul l r | .sub l r =>
    match Term.type H l, Term.type H r with
    | some τ₁, some τ₂ => if τ₁ = τ₂ then some τ₁ else none
    | _, _ => none
  | .neg t => Term.type H t

def Formula.WellTyped (H : Hypernet) : Formula → Prop
  | .atom _ l r => ∃ τ, l.type H = some τ ∧ r.type H = some τ
  | .inRegion t simplex S =>
      (∃ τ, t.type H = some τ ∧ τ.shape = [simplex.complex.ambDim]) ∧
      S.simplex = simplex
  …
\end{leansrc}
\par\medskip\noindent{\small \texttt{PPLC.}\texttt{\bfseries farkasCombinationValid}}\hfill{\footnotesize \path{lean/qtsl/Qtsl/Core/Catalogue.lean}, line 918}\par\nopagebreak
\begin{leansrc}
def farkasCombinationValid (cc : Term) (premise : Formula)
    (conclusion : Term) : Prop
\end{leansrc}
\par\medskip\noindent{\small \texttt{PPLC.}\texttt{\bfseries farkasEntries}}\hfill{\footnotesize \path{lean/qtsl/Qtsl/Core/Catalogue.lean}, line 778}\par\nopagebreak
\begin{leansrc}
def farkasEntries : Term → Option (List (Rat × Term))
  | .const c => if c = 0 then some [] else none
  | .add (.mul (.const l) e) rest =>
    (farkasEntries rest).map ((l, e) :: ·)
  | _ => none

def wireCoordKey (w i : Nat) : Nat
\end{leansrc}
\par\medskip\noindent{\small \texttt{PPLC.}\texttt{\bfseries atomLeZeroPolys}}\hfill{\footnotesize \path{lean/qtsl/Qtsl/Core/Catalogue.lean}, line 769}\par\nopagebreak
\begin{leansrc}
def atomLeZeroPolys : Formula → List QArith.Poly
  | .atom .le l r => [QArith.Poly.sub l.toPoly r.toPoly]
  | .atom .lt l r => [QArith.Poly.sub l.toPoly r.toPoly]
  | .atom .ge l r => [QArith.Poly.sub r.toPoly l.toPoly]
  | .atom .gt l r => [QArith.Poly.sub r.toPoly l.toPoly]
  | .atom .eq l r => [QArith.Poly.sub l.toPoly r.toPoly,
                      QArith.Poly.sub r.toPoly l.toPoly]
  | _ => []

def farkasEntries : Term → Option (List (Rat × Term))
  | .const c => if c = 0 then some [] else none
  | .add (.mul (.const l) e) rest =>
    (farkasEntries rest).map ((l, e) :: ·)
  | _ => none

def wireCoordKey (w i : Nat) : Nat
\end{leansrc}
\par\medskip\noindent{\small \texttt{PPLC.}\texttt{\bfseries farkasCheckCore}}\hfill{\footnotesize \path{lean/qtsl/Qtsl/Core/Catalogue.lean}, line 899}\par\nopagebreak
\begin{leansrc}
def farkasCheckCore (cc : Term) (premise : Formula) (conclusion : Term) : Bool
\end{leansrc}
\par\medskip\noindent{\small \texttt{PPLC.}\texttt{\bfseries farkasCheck}}\hfill{\footnotesize \path{lean/qtsl/Qtsl/Core/Catalogue.lean}, line 914}\par\nopagebreak
\begin{leansrc}
def farkasCheck (cc : Term) (premise : Formula) (conclusion : Term) : Bool
\end{leansrc}
\par\medskip\noindent{\small \texttt{PPLC.}\texttt{\bfseries sound\_farkas}}\hfill{\footnotesize \path{lean/qtsl/Qtsl/Soundness.lean}, line 4337}\par\nopagebreak
\begin{leansrc}
theorem sound_farkas {H : Hypernet} {seq : Sequent}
    {combinationCoeffs conclusion : Term}
    (hshape : ∃ s, conclusion.bshape H = some s)
    (hfin : conclusion.finWires H)
    (hwpos : ∀ w ∈ seq.pre.freeWires ++ conclusion.freeWires,
      ∃ τ, H.wireType w = some τ ∧ 0 < Shape.totalSize τ.shape)
    (hvalid : farkasCombinationValid combinationCoeffs seq.pre conclusion)
    (hpost : seq.post = .atom .le conclusion (.const 0)) :
    seq.semanticallyValid H
\end{leansrc}
\par\medskip\noindent{\small \texttt{PPLC.}\texttt{\bfseries atomLeZeroPolysW}}\hfill{\footnotesize \path{lean/qtsl/Qtsl/Core/Catalogue.lean}, line 800}\par\nopagebreak
\begin{leansrc}
def atomLeZeroPolysW : Formula → List QArith.Poly
  | .atom op l r =>
    match l.asCoordPolyIdxW? with
    | none => []
    | some pl =>
      match r.asCoordPolyIdxW? with
      | none => []
      | some pr =>
        match op with
        | .le | .lt => [QArith.Poly.sub pl pr]
        | .ge | .gt => [QArith.Poly.sub pr pl]
        | .eq => [QArith.Poly.sub pl pr, QArith.Poly.sub pr pl]
        | .ne => []
  | _ => []

def farkasEntriesW : List (Rat × Term) → Option (List (Rat × QArith.Poly))
  | [] => some []
  | (l, e) :: rest =>
    match e.asCoordPolyIdxW?, farkasEntriesW rest with
    | some pe, some prest => some ((l, pe) :: prest)
    | _, _ => none

def farkasCheckW (cc : Term) (premise : Formula) (conclusion : Term) : Bool
\end{leansrc}
\par\medskip\noindent{\small \texttt{PPLC.}\texttt{\bfseries atomLeZeroPolysM}}\hfill{\footnotesize \path{lean/qtsl/Qtsl/Core/Catalogue.lean}, line 854}\par\nopagebreak
\begin{leansrc}
def atomLeZeroPolysM : Formula → List QArith.Poly
  | .atom op l r =>
    match l.asCoordPolyMixed? with
    | none => []
    | some pl =>
      match r.asCoordPolyMixed? with
      | none => []
      | some pr =>
        match op with
        | .le | .lt => [QArith.Poly.sub pl pr]
        | .ge | .gt => [QArith.Poly.sub pr pl]
        | .eq => [QArith.Poly.sub pl pr, QArith.Poly.sub pr pl]
        | .ne => []
  | _ => []

def farkasEntriesM : List (Rat × Term) → Option (List (Rat × QArith.Poly))
  | [] => some []
  | (l, e) :: rest =>
    match e.asCoordPolyMixed?, farkasEntriesM rest with
    | some pe, some prest => some ((l, pe) :: prest)
    | _, _ => none

def farkasCheckM (cc : Term) (premise : Formula) (conclusion : Term) : Bool
\end{leansrc}
\par\medskip\noindent{\small \texttt{PPLC.}\texttt{\bfseries farkasCheckW}}\hfill{\footnotesize \path{lean/qtsl/Qtsl/Core/Catalogue.lean}, line 822}\par\nopagebreak
\begin{leansrc}
def farkasCheckW (cc : Term) (premise : Formula) (conclusion : Term) : Bool
\end{leansrc}
\par\medskip\noindent{\small \texttt{PPLC.}\texttt{\bfseries farkasCheckM}}\hfill{\footnotesize \path{lean/qtsl/Qtsl/Core/Catalogue.lean}, line 876}\par\nopagebreak
\begin{leansrc}
def farkasCheckM (cc : Term) (premise : Formula) (conclusion : Term) : Bool
\end{leansrc}
\par\medskip\noindent{\small \texttt{PPLC.}\texttt{\bfseries wireCoordKey}}\hfill{\footnotesize \path{lean/qtsl/Qtsl/Core/Catalogue.lean}, line 784}\par\nopagebreak
\begin{leansrc}
def wireCoordKey (w i : Nat) : Nat
\end{leansrc}
\par\medskip\noindent{\small \texttt{PPLC.}\texttt{\bfseries farkas\_fold\_eval\_nonpos}}\hfill{\footnotesize \path{lean/qtsl/Qtsl/Soundness.lean}, line 3445}\par\nopagebreak
\begin{leansrc}
theorem farkas_fold_eval_nonpos (ρ : Nat → ℚ) :
    ∀ (entries : List (ℚ × Term)),
    (∀ e ∈ entries, 0 ≤ e.1 ∧ QArith.Poly.eval ρ e.2.toPoly ≤ 0) →
    QArith.Poly.eval ρ
      (entries.foldr
        (fun (le : ℚ × Term) acc =>
          QArith.Poly.add (QArith.Poly.scale le.1 le.2.toPoly) acc) []) ≤ 0
\end{leansrc}
\par\medskip\noindent{\small \texttt{PPLC.}\texttt{\bfseries farkas\_fold\_eval\_nonpos\_polys}}\hfill{\footnotesize \path{lean/qtsl/Qtsl/Soundness.lean}, line 3830}\par\nopagebreak
\begin{leansrc}
theorem farkas_fold_eval_nonpos_polys (ρ : Nat → ℚ) :
    ∀ (pentries : List (ℚ × QArith.Poly)),
    (∀ e ∈ pentries, 0 ≤ e.1 ∧ QArith.Poly.eval ρ e.2 ≤ 0) →
    QArith.Poly.eval ρ
      (pentries.foldr
        (fun (le : ℚ × QArith.Poly) acc =>
          QArith.Poly.add (QArith.Poly.scale le.1 le.2) acc) []) ≤ 0
\end{leansrc}
\par\medskip\noindent{\small \texttt{PPLC.}\texttt{\bfseries PFarkasData}}\hfill{\footnotesize \path{lean/qtsl/Qtsl/Core/Catalogue.lean}, line 1517}\par\nopagebreak
\begin{leansrc}
inductive PFarkasData : Hypernet → Sequent → Type where …
\end{leansrc}
\par\medskip\noindent{\small \texttt{PPLC.}\texttt{\bfseries coverInfeasibilityValid}}\hfill{\footnotesize \path{lean/qtsl/Qtsl/Core/Catalogue.lean}, line 1151}\par\nopagebreak
\begin{leansrc}
def coverInfeasibilityValid
    (_termX : Term) (simplex : Simplex) (S : ExtensionSet)
    (nBranches : Nat)
    (branchSimplex : Fin nBranches → Simplex)
    (branchS : Fin nBranches → ExtensionSet)
    (_coverCert : Term) : Prop
\end{leansrc}
\par\medskip\noindent{\small \texttt{PPLC.}\texttt{\bfseries coverExtendPass}}\hfill{\footnotesize \path{lean/qtsl/Qtsl/Core/Catalogue.lean}, line 1127}\par\nopagebreak
\begin{leansrc}
def coverExtendPass (ivs : List (Rat × Rat)) (r : Rat) : Rat
\end{leansrc}
\par\medskip\noindent{\small \texttt{PPLC.}\texttt{\bfseries coverReach}}\hfill{\footnotesize \path{lean/qtsl/Qtsl/Core/Catalogue.lean}, line 1130}\par\nopagebreak
\begin{leansrc}
def coverReach : Nat → List (Rat × Rat) → Rat → Rat
  | 0, _, r => r
  | n + 1, ivs, r => coverReach n ivs (coverExtendPass ivs r)

def coverChainCheck (ivs : List (Rat × Rat)) (lo hi : Rat) : Bool
\end{leansrc}
\par\medskip\noindent{\small \texttt{PPLC.}\texttt{\bfseries coverChainCheck}}\hfill{\footnotesize \path{lean/qtsl/Qtsl/Core/Catalogue.lean}, line 1134}\par\nopagebreak
\begin{leansrc}
def coverChainCheck (ivs : List (Rat × Rat)) (lo hi : Rat) : Bool
\end{leansrc}
\par\medskip\noindent{\small \texttt{PPLC.}\texttt{\bfseries coverCheck}}\hfill{\footnotesize \path{lean/qtsl/Qtsl/Core/Catalogue.lean}, line 1137}\par\nopagebreak
\begin{leansrc}
def coverCheck (simplex : Simplex) (S : ExtensionSet) (nBranches : Nat)
    (branchSimplex : Fin nBranches → Simplex)
    (branchS : Fin nBranches → ExtensionSet) : Bool
\end{leansrc}
\par\medskip\noindent{\small \texttt{PPLC.}\texttt{\bfseries sound\_subdivide}}\hfill{\footnotesize \path{lean/qtsl/Qtsl/Soundness.lean}, line 7758}\par\nopagebreak
\begin{leansrc}
theorem sound_subdivide {H : Hypernet} {seq regionPremise : Sequent}
    {nBranches : Nat} {branchPremise : Fin nBranches → Sequent}
    {termX : Term} {simplex : Simplex} {S : ExtensionSet}
    {branchSimplex : Fin nBranches → Simplex}
    {branchS : Fin nBranches → ExtensionSet}
    {coverCertificate : Term}
    (ihRegion : regionPremise.semanticallyValid H)
    (ihBranch : ∀ k, (branchPremise k).semanticallyValid H)
    (hpre : seq.pre = regionPremise.pre)
    (hregionPost : regionPremise.post = .inRegion termX simplex S)
    (hbranchPre : ∀ k, (branchPremise k).pre
                   = .and regionPremise.pre
                          (.inRegion termX (branchSimplex k) (branchS k)))
    (hbranchPost : ∀ k, (branchPremise k).post = seq.post)
    (hcover : coverInfeasibilityValid termX simplex S nBranches
                branchSimplex branchS coverCertificate) :
    seq.semanticallyValid H
\end{leansrc}
\par\medskip\noindent{\small \texttt{PPLC.}\texttt{\bfseries interval1D\_elim}}\hfill{\footnotesize \path{lean/qtsl/Qtsl/Soundness.lean}, line 7642}\par\nopagebreak
\begin{leansrc}
theorem interval1D_elim {S : ExtensionSet} {lo hi : ℚ}
    (hint : S.interval1D? = some (lo, hi))
    {x : Fin S.simplex.complex.ambDim → ℚ}
    (hmask : MaskRegion S.simplex.complex S.simplex.contains S.contains x)
    (j0 : Fin S.simplex.complex.ambDim) :
    lo ≤ x j0 ∧ x j0 ≤ hi ∧ lo < hi
\end{leansrc}
\par\medskip\noindent{\small \texttt{PPLC.}\texttt{\bfseries interval1D\_intro}}\hfill{\footnotesize \path{lean/qtsl/Qtsl/Soundness.lean}, line 7676}\par\nopagebreak
\begin{leansrc}
theorem interval1D_intro {S : ExtensionSet} {lo hi : ℚ}
    (hint : S.interval1D? = some (lo, hi))
    {x : Fin S.simplex.complex.ambDim → ℚ}
    (j0 : Fin S.simplex.complex.ambDim)
    (hxlo : lo ≤ x j0) (hxhi : x j0 ≤ hi) :
    MaskRegion S.simplex.complex S.simplex.contains S.contains x
\end{leansrc}
\par\medskip\noindent{\small \texttt{PPLC.}\texttt{\bfseries ChainInv}}\hfill{\footnotesize \path{lean/qtsl/Qtsl/Soundness.lean}, line 7700}\par\nopagebreak
\begin{leansrc}
def ChainInv (ivs : List (ℚ × ℚ)) (lo r : ℚ) : Prop
\end{leansrc}
\par\medskip\noindent{\small \texttt{PPLC.}\texttt{\bfseries chainInv\_extendPass}}\hfill{\footnotesize \path{lean/qtsl/Qtsl/Soundness.lean}, line 7704}\par\nopagebreak
\begin{leansrc}
theorem chainInv_extendPass (ivs : List (ℚ × ℚ)) (lo : ℚ) :
    ∀ (sub : List (ℚ × ℚ)), (∀ iv ∈ sub, iv ∈ ivs) →
    ∀ (r : ℚ), ChainInv ivs lo r →
    ChainInv ivs lo (coverExtendPass sub r)
\end{leansrc}
\par\medskip\noindent{\small \texttt{PPLC.}\texttt{\bfseries chainInv\_reach}}\hfill{\footnotesize \path{lean/qtsl/Qtsl/Soundness.lean}, line 7735}\par\nopagebreak
\begin{leansrc}
theorem chainInv_reach (ivs : List (ℚ × ℚ)) (lo : ℚ) :
    ∀ (n : Nat) (r : ℚ), ChainInv ivs lo r →
    ChainInv ivs lo (coverReach n ivs r)
\end{leansrc}
\par\medskip\noindent{\small \texttt{PPLC.}\texttt{\bfseries coverChain\_covers}}\hfill{\footnotesize \path{lean/qtsl/Qtsl/Soundness.lean}, line 7745}\par\nopagebreak
\begin{leansrc}
theorem coverChain_covers {ivs : List (ℚ × ℚ)} {lo hi : ℚ} (hlt : lo < hi)
    (hcheck : coverChainCheck ivs lo hi = true) :
    ∀ x : ℚ, lo ≤ x → x ≤ hi → ∃ iv ∈ ivs, iv.1 ≤ x ∧ x ≤ iv.2
\end{leansrc}
\par\medskip\noindent{\small \texttt{PPLC.}\texttt{\bfseries PSubdivideData}}\hfill{\footnotesize \path{lean/qtsl/Qtsl/Core/Catalogue.lean}, line 1533}\par\nopagebreak
\begin{leansrc}
inductive PSubdivideData : Hypernet → Sequent → Type where …
\end{leansrc}
\par\medskip\noindent{\small \texttt{PPLC.}\texttt{\bfseries reluLinked}}\hfill{\footnotesize \path{lean/qtsl/Qtsl/Core/Catalogue.lean}, line 1223}\par\nopagebreak
\begin{leansrc}
def reluLinked (H : Hypernet) (preact post : Wire) : Prop
\end{leansrc}
\par\medskip\noindent{\small \texttt{PPLC.PrimApp.}\texttt{\bfseries isReluAt}}\hfill{\footnotesize \path{lean/qtsl/Qtsl/Core/Catalogue.lean}, line 1217}\par\nopagebreak
\begin{leansrc}
def PrimApp.isReluAt (preact post : Wire) : PrimApp → Bool
  | .relu s i o =>
      decide (i = preact) && decide (o = post) &&
      decide (0 < Shape.totalSize s)
  | _ => false

def reluLinked (H : Hypernet) (preact post : Wire) : Prop
\end{leansrc}
\par\medskip\noindent{\small \texttt{PPLC.}\texttt{\bfseries affineWireTerm}}\hfill{\footnotesize \path{lean/qtsl/Qtsl/Core/Catalogue.lean}, line 1159}\par\nopagebreak
\begin{leansrc}
def affineWireTerm (offset slope : Rat) (w : Wire) : Term
\end{leansrc}
\par\medskip\noindent{\small \texttt{PPLC.}\texttt{\bfseries lirpaConclusionForm}}\hfill{\footnotesize \path{lean/qtsl/Qtsl/Core/Catalogue.lean}, line 1162}\par\nopagebreak
\begin{leansrc}
def lirpaConclusionForm
    (preact post : Wire) (α β γ δ : Rat) (F : Formula) : Prop
\end{leansrc}
\par\medskip\noindent{\small \texttt{PPLC.}\texttt{\bfseries reluLinked\_facts}}\hfill{\footnotesize \path{lean/qtsl/Qtsl/Soundness.lean}, line 8057}\par\nopagebreak
\begin{leansrc}
theorem reluLinked_facts {H : Hypernet} {preact post : Wire}
    (hlink : reluLinked H preact post) (σ : H.Valuation)
    (hadm : σ.Admissible) :
    ∃ qp qy : QTensor,
      σ.wireQ preact = some qp ∧ σ.wireQ post = some qy ∧
      qp.shape = qy.shape ∧ 0 < qp.entries.length ∧
      qp.entries.length = qy.entries.length ∧
      ∀ i, i < qp.entries.length →
        (qp.entries.getD i 0 ≤ 0 ∧ qy.entries.getD i 0 = 0) ∨
        (0 ≤ qp.entries.getD i 0 ∧
          qy.entries.getD i 0 = qp.entries.getD i 0)
\end{leansrc}
\par\medskip\noindent{\small \texttt{PPLC.}\texttt{\bfseries sound\_lirpa}}\hfill{\footnotesize \path{lean/qtsl/Qtsl/Soundness.lean}, line 8235}\par\nopagebreak
\begin{leansrc}
theorem sound_lirpa {H : Hypernet}
    {seq lowerPremise upperPremise : Sequent}
    {preact post : Wire} {a b α β γ δ : ℚ}
    (ihl : lowerPremise.semanticallyValid H)
    (ihu : upperPremise.semanticallyValid H)
    (hlink : reluLinked H preact post)
    (hAneg : a < 0) (hBpos : 0 < b)
    (hLower1 : α * a + β ≤ 0) (hLower2 : β ≤ 0) (hLower3 : α * b + β ≤ b)
    (hUpper1 : 0 ≤ γ * a + δ) (hUpper2 : 0 ≤ δ) (hUpper3 : b ≤ γ * b + δ)
    (hpre : seq.pre = lowerPremise.pre)
    (hpreMatch : lowerPremise.pre = upperPremise.pre)
    (hlowerPost : lowerPremise.post = .atom .le (.const a) (.var preact))
    (hupperPost : upperPremise.post = .atom .le (.var preact) (.const b))
    (hpostForm : lirpaConclusionForm preact post α β γ δ seq.post) :
    seq.semanticallyValid H
\end{leansrc}
\par\medskip\noindent{\small \texttt{PPLC.ReluExpansion.}\texttt{\bfseries reluExpH}}\hfill{\footnotesize \path{lean/qtsl/Qtsl/ReluExpansion.lean}, line 156}\par\nopagebreak
\begin{leansrc}
def reluExpH : Hypernet where …
\end{leansrc}
\par\medskip\noindent{\small \texttt{PPLC.ReluExpansion.}\texttt{\bfseries relu\_projection}}\hfill{\footnotesize \path{lean/qtsl/Qtsl/ReluExpansion.lean}, line 462}\par\nopagebreak
\begin{leansrc}
theorem relu_projection (p y : ℚ) :
    (∃ σ : reluExpH.Valuation, σ.Admissible ∧
        σ.entryFlat 0 0 = some (.fin p) ∧
        σ.entryFlat 8 0 = some (.fin y)) ↔
    ((p ≤ 0 ∧ y = 0) ∨ (0 ≤ p ∧ y = p))
\end{leansrc}
\par\medskip\noindent{\small \texttt{PPLC.}\texttt{\bfseries lirpa\_lower\_neg}}\hfill{\footnotesize \path{lean/qtsl/Qtsl/Soundness.lean}, line 8014}\par\nopagebreak
\begin{leansrc}
theorem lirpa_lower_neg {a α β p : ℚ} (ha : a < 0)
    (h1 : α * a + β ≤ 0) (h2 : β ≤ 0) (hap : a ≤ p) (hp0 : p ≤ 0) :
    α * p + β ≤ 0
\end{leansrc}
\par\medskip\noindent{\small \texttt{PPLC.}\texttt{\bfseries lirpa\_lower\_pos}}\hfill{\footnotesize \path{lean/qtsl/Qtsl/Soundness.lean}, line 8025}\par\nopagebreak
\begin{leansrc}
theorem lirpa_lower_pos {b α β p : ℚ} (hb : 0 < b)
    (h2 : β ≤ 0) (h3 : α * b + β ≤ b) (hp0 : 0 ≤ p) (hpb : p ≤ b) :
    α * p + β ≤ p
\end{leansrc}
\par\medskip\noindent{\small \texttt{PPLC.}\texttt{\bfseries lirpa\_upper\_neg}}\hfill{\footnotesize \path{lean/qtsl/Qtsl/Soundness.lean}, line 8036}\par\nopagebreak
\begin{leansrc}
theorem lirpa_upper_neg {a γ δ p : ℚ} (ha : a < 0)
    (h1 : 0 ≤ γ * a + δ) (h2 : 0 ≤ δ) (hap : a ≤ p) (hp0 : p ≤ 0) :
    0 ≤ γ * p + δ
\end{leansrc}
\par\medskip\noindent{\small \texttt{PPLC.}\texttt{\bfseries lirpa\_upper\_pos}}\hfill{\footnotesize \path{lean/qtsl/Qtsl/Soundness.lean}, line 8047}\par\nopagebreak
\begin{leansrc}
theorem lirpa_upper_pos {b γ δ p : ℚ} (hb : 0 < b)
    (h2 : 0 ≤ δ) (h3 : b ≤ γ * b + δ) (hp0 : 0 ≤ p) (hpb : p ≤ b) :
    p ≤ γ * p + δ
\end{leansrc}
\par\medskip\noindent{\small \texttt{PPLC.}\texttt{\bfseries affineWireTerm\_evalQ}}\hfill{\footnotesize \path{lean/qtsl/Qtsl/Soundness.lean}, line 7994}\par\nopagebreak
\begin{leansrc}
theorem affineWireTerm_evalQ {H : Hypernet} (σ : H.Valuation) {w : Wire}
    {qw : QTensor} (h : σ.wireQ w = some qw) (off slope : ℚ) :
    ∃ qa, Term.evalQ σ (affineWireTerm off slope w) = some qa ∧
      qa.shape = qw.shape ∧
      qa.entries.length = qw.entries.length ∧
      ∀ i, i < qw.entries.length →
        qa.entries.getD i 0 = off + slope * qw.entries.getD i 0
\end{leansrc}
\par\medskip\noindent{\small \texttt{PPLC.}\texttt{\bfseries sound\_reluPos}}\hfill{\footnotesize \path{lean/qtsl/Qtsl/Soundness.lean}, line 8156}\par\nopagebreak
\begin{leansrc}
theorem sound_reluPos {H : Hypernet} {seq premise : Sequent}
    {preact post : Wire}
    (ih : premise.semanticallyValid H)
    (hlink : reluLinked H preact post)
    (hpre : seq.pre = premise.pre)
    (hpremPost : premise.post = .atom .ge (.var preact) (.const 0))
    (hpost : seq.post = .atom .eq (.var post) (.var preact)) :
    seq.semanticallyValid H
\end{leansrc}
\par\medskip\noindent{\small \texttt{PPLC.}\texttt{\bfseries PReLUPosData}}\hfill{\footnotesize \path{lean/qtsl/Qtsl/Core/Catalogue.lean}, line 1578}\par\nopagebreak
\begin{leansrc}
inductive PReLUPosData : Hypernet → Sequent → Type where …
\end{leansrc}
\par\medskip\noindent{\small \texttt{PPLC.QTensor.}\texttt{\bfseries rel2\_const\_right\_elim}}\hfill{\footnotesize \path{lean/qtsl/Qtsl/Soundness.lean}, line 7914}\par\nopagebreak
\begin{leansrc}
theorem QTensor.rel2_const_right_elim {r : ℚ → ℚ → Prop} {q : QTensor}
    {c : ℚ} (h : QTensor.rel2 r q ⟨[], [c], rfl⟩) :
    ∀ x ∈ q.entries, r x c
\end{leansrc}
\par\medskip\noindent{\small \texttt{PPLC.}\texttt{\bfseries sound\_reluNeg}}\hfill{\footnotesize \path{lean/qtsl/Qtsl/Soundness.lean}, line 8193}\par\nopagebreak
\begin{leansrc}
theorem sound_reluNeg {H : Hypernet} {seq premise : Sequent}
    {preact post : Wire}
    (ih : premise.semanticallyValid H)
    (hlink : reluLinked H preact post)
    (hpre : seq.pre = premise.pre)
    (hpremPost : premise.post = .atom .le (.var preact) (.const 0))
    (hpost : seq.post = .atom .eq (.var post) (.const 0)) :
    seq.semanticallyValid H
\end{leansrc}
\par\medskip\noindent{\small \texttt{PPLC.}\texttt{\bfseries PReLUNegData}}\hfill{\footnotesize \path{lean/qtsl/Qtsl/Core/Catalogue.lean}, line 1588}\par\nopagebreak
\begin{leansrc}
inductive PReLUNegData : Hypernet → Sequent → Type where …
\end{leansrc}
\par\medskip\noindent{\small \texttt{PPLC.}\texttt{\bfseries PReLUNNData}}\hfill{\footnotesize \path{lean/qtsl/Qtsl/Core/Catalogue.lean}, line 1598}\par\nopagebreak
\begin{leansrc}
inductive PReLUNNData : Hypernet → Sequent → Type where …
\end{leansrc}
\par\medskip\noindent{\small \texttt{PPLC.}\texttt{\bfseries sound\_reluNN}}\hfill{\footnotesize \path{lean/qtsl/Qtsl/Soundness.lean}, line 8116}\par\nopagebreak
\begin{leansrc}
theorem sound_reluNN {H : Hypernet} {seq : Sequent} {preact post : Wire}
    (hlink : reluLinked H preact post)
    (hpost : seq.post = .atom .ge (.var post) (.const 0)) :
    seq.semanticallyValid H
\end{leansrc}
\par\medskip\noindent{\small \texttt{PPLC.}\texttt{\bfseries sound\_pConstExact}}\hfill{\footnotesize \path{lean/qtsl/Qtsl/Soundness.lean}, line 7403}\par\nopagebreak
\begin{leansrc}
theorem sound_pConstExact {H : Hypernet} {seq : Sequent} {e₁ e₂ : Term}
    {b : BowTie}
    (hEval : (Formula.atom b e₁ e₂).evalClosed = some true)
    (hpost : seq.post = .atom b e₁ e₂) :
    seq.semanticallyValid H
\end{leansrc}
\par\medskip\noindent{\small \texttt{PPLC.}\texttt{\bfseries PReLUAboveData}}\hfill{\footnotesize \path{lean/qtsl/Qtsl/Core/Catalogue.lean}, line 1605}\par\nopagebreak
\begin{leansrc}
inductive PReLUAboveData : Hypernet → Sequent → Type where …
\end{leansrc}
\par\medskip\noindent{\small \texttt{PPLC.}\texttt{\bfseries sound\_reluAbove}}\hfill{\footnotesize \path{lean/qtsl/Qtsl/Soundness.lean}, line 8138}\par\nopagebreak
\begin{leansrc}
theorem sound_reluAbove {H : Hypernet} {seq : Sequent} {preact post : Wire}
    (hlink : reluLinked H preact post)
    (hpost : seq.post = .atom .ge (.var post) (.var preact)) :
    seq.semanticallyValid H
\end{leansrc}
\par\medskip\noindent{\small \texttt{PPLC.}\texttt{\bfseries sound\_pBound}}\hfill{\footnotesize \path{lean/qtsl/Qtsl/Soundness.lean}, line 7412}\par\nopagebreak
\begin{leansrc}
theorem sound_pBound {H : Hypernet} {seq : Sequent} {components : List Sequent}
    {combination conclusion : Term}
    (hcomps : ∀ c ∈ components, c.semanticallyValid H)
    (hpreC : ∀ c ∈ components, premiseContains seq.pre c.pre)
    (hshape : ∃ s, conclusion.bshape H = some s)
    (hfin : conclusion.finWires H)
    (hwpos : ∀ w ∈ seq.pre.freeWires ++ conclusion.freeWires,
      ∃ τ, H.wireType w = some τ ∧ 0 < Shape.totalSize τ.shape)
    (hcompPos : ∀ c ∈ components, ∀ w ∈ c.post.freeWires,
      ∃ τ, H.wireType w = some τ ∧ 0 < Shape.totalSize τ.shape)
    (hvalid : syzygyCombinationValid combination components seq.pre conclusion)
    (hform : syzygyConclusionForm conclusion seq.post) :
    seq.semanticallyValid H
\end{leansrc}
\par\medskip\noindent{\small \texttt{PPLC.}\texttt{\bfseries baryLambdas1D}}\hfill{\footnotesize \path{lean/qtsl/Qtsl/Core/Catalogue.lean}, line 449}\par\nopagebreak
\begin{leansrc}
def baryLambdas1D (cx : SimpC) (dContains : Fin cx.nVertices → Bool)
    (t : Term) : Option (List (Fin cx.nVertices × Term))
\end{leansrc}
\par\medskip\noindent{\small \texttt{PPLC.}\texttt{\bfseries MaskRegionStrict}}\hfill{\footnotesize \path{lean/qtsl/Qtsl/Catalogue.lean}, line 386}\par\nopagebreak
\begin{leansrc}
def MaskRegionStrict (C : SimpC) (mask ext : Fin C.nVertices → Bool)
    (x : Fin C.ambDim → ℚ) : Prop
\end{leansrc}
\par\medskip\noindent{\small \texttt{PPLC.SimpC.}\texttt{\bfseries CellIntDisjoint}}\hfill{\footnotesize \path{lean/qtsl/Qtsl/Catalogue.lean}, line 396}\par\nopagebreak
\begin{leansrc}
def SimpC.CellIntDisjoint (C : SimpC) : Prop
\end{leansrc}
\par\medskip\noindent{\small \texttt{PPLC.}\texttt{\bfseries PCellInIntData}}\hfill{\footnotesize \path{lean/qtsl/Qtsl/Core/Catalogue.lean}, line 1612}\par\nopagebreak
\begin{leansrc}
inductive PCellInIntData : Hypernet → Sequent → Type where …
\end{leansrc}
\par\medskip\noindent{\small \texttt{PPLC.}\texttt{\bfseries sound\_cellInInterior}}\hfill{\footnotesize \path{lean/qtsl/Qtsl/Soundness.lean}, line 6531}\par\nopagebreak
\begin{leansrc}
theorem sound_cellInInterior {H : Hypernet} {seq premise : Sequent}
    {schema : SimpC} {input : Wire} {outputs : List Wire}
    {whichI : Fin schema.cells.length}
    (ih : premise.semanticallyValid H)
    (happ : PrimApp.cell schema input outputs ∈ H.applications)
    (hinshape : ∃ τ, H.wireType input = some τ ∧
      τ.shape = [schema.ambDim])
    (houts : ∀ o ∈ outputs, ∃ τ, H.wireType o = some τ ∧
      Shape.totalSize τ.shape = 1)
    (hpre : seq.pre = premise.pre)
    (hint : cellStrictInteriorForm schema (.var input) whichI premise.post)
    (hpost : cellOneHotForm schema outputs whichI seq.post) :
    seq.semanticallyValid H
\end{leansrc}
\par\medskip\noindent{\small \texttt{PPLC.}\texttt{\bfseries cellStrictInteriorBuild}}\hfill{\footnotesize \path{lean/qtsl/Qtsl/Core/Catalogue.lean}, line 1167}\par\nopagebreak
\begin{leansrc}
def cellStrictInteriorBuild (schema : SimpC) (t : Term)
    (k : Fin schema.cells.length) : Option Formula
\end{leansrc}
\par\medskip\noindent{\small \texttt{PPLC.}\texttt{\bfseries cellOneHotBuild}}\hfill{\footnotesize \path{lean/qtsl/Qtsl/Core/Catalogue.lean}, line 1180}\par\nopagebreak
\begin{leansrc}
def cellOneHotBuild (schema : SimpC) (outputs : List Wire)
    (k : Fin schema.cells.length) : Option Formula
\end{leansrc}
\par\medskip\noindent{\small \texttt{PPLC.}\texttt{\bfseries maskRegionStrict\_pair\_intro}}\hfill{\footnotesize \path{lean/qtsl/Qtsl/Soundness.lean}, line 6218}\par\nopagebreak
\begin{leansrc}
theorem maskRegionStrict_pair_intro {C : SimpC}
    {mask ext : Fin C.nVertices → Bool}
    {x : Fin C.ambDim → ℚ} {a b : Fin C.nVertices} {V₀ V₁ : ℚ}
    (hfil : (List.finRange C.nVertices).filter (fun v => mask v) = [a, b])
    (hscope : ∀ v, ext v = true → mask v = true)
    (hVa : ∀ j, C.coord a j = V₀) (hVb : ∀ j, C.coord b j = V₁)
    (hne : V₀ ≠ V₁) (hdim : C.ambDim = 1) (j0 : Fin C.ambDim)
    (hw0t : ext a = true → (V₁ - x j0) / (V₁ - V₀) < 0)
    (hw0f : ext a = false → 0 < (V₁ - x j0) / (V₁ - V₀))
    (hw1t : ext b = true → (x j0 - V₀) / (V₁ - V₀) < 0)
    (hw1f : ext b = false → 0 < (x j0 - V₀) / (V₁ - V₀)) :
    MaskRegionStrict C mask ext x
\end{leansrc}
\par\medskip\noindent{\small \texttt{PPLC.}\texttt{\bfseries holds\_cellOneHot}}\hfill{\footnotesize \path{lean/qtsl/Qtsl/Soundness.lean}, line 6480}\par\nopagebreak
\begin{leansrc}
theorem holds_cellOneHot {H : Hypernet} (σ : H.Valuation) {schema : SimpC}
    {outputs : List Wire} {k : Fin schema.cells.length} {F : Formula}
    (hbuild : cellOneHotForm schema outputs k F)
    (houts : ∀ o ∈ outputs, ∃ τ, H.wireType o = some τ ∧
      Shape.totalSize τ.shape = 1)
    (hsem : ∀ j : Fin outputs.length,
      σ.entryFlat (outputs.get j) 0
        = some (.fin (if j.val = k.val then 1 else 0))) :
    Formula.holds σ F
\end{leansrc}
\par\medskip\noindent{\small \texttt{PPLC.}\texttt{\bfseries cellStrictInteriorBuild\_spec}}\hfill{\footnotesize \path{lean/qtsl/Qtsl/Soundness.lean}, line 6091}\par\nopagebreak
\begin{leansrc}
theorem cellStrictInteriorBuild_spec {schema : SimpC} {t : Term}
    {k : Fin schema.cells.length} {F : Formula}
    (h : cellStrictInteriorForm schema t k F) :
    schema.ambDim = 1 ∧
    ∃ (a b : Fin schema.nVertices) (V₀ V₁ : ℚ),
      (List.finRange schema.nVertices).filter
        (fun v => (schema.cells.get k).1 v) = [a, b] ∧
      (∀ j, schema.coord a j = V₀) ∧
      (∀ j, schema.coord b j = V₁) ∧
      V₀ ≠ V₁ ∧
      F = .and
        (.atom (if (schema.cells.get k).2 a then .lt else .gt)
          (.add (.const (V₁ / (V₁ - V₀))) (.mul (.const (-1 / (V₁ - V₀))) t))
          (.const 0))
        (.atom (if (schema.cells.get k).2 b then .lt else .gt)
          (.add (.const (-V₀ / (V₁ - V₀))) (.mul (.const (1 / (V₁ - V₀))) t))
          (.const 0))
\end{leansrc}
\par\medskip\noindent{\small \texttt{PPLC.PrimApp.}\texttt{\bfseries SemOk}}\hfill{\footnotesize \path{lean/qtsl/Qtsl/Catalogue.lean}, line 445}\par\nopagebreak
\begin{leansrc}
def PrimApp.SemOk {H : Hypernet} (σ : H.Valuation) : PrimApp → Prop
  | .const _ty c out =>
      ∀ i : ℕ, σ.entryFlat out i = c.entryE i
  | .contract δ r inputs output _inputRanks _outputRank pat =>
      ∀ ν₀ ∈ Shape.positions pat.idxShape,
        ∃ eOut eExact,
          σ.entryAt output (pat.outCoords ν₀) = some eOut ∧
          contractExact σ inputs pat (pat.outCoords ν₀) = some eExact ∧
          δ.RoundOkE r eExact eOut
  | .cast _δfrom δto r shape input output =>
      ∀ i : Fin (Shape.totalSize shape),
        ∃ eIn eOut, σ.entryFlat input i.val = some eIn ∧
          σ.entryFlat output i.val = some eOut ∧ δto.RoundOkE r eIn eOut
  | .cell schema input outputs =>

      schema.CellIntDisjoint ∧
      ∃ x : Fin schema.ambDim → ℚ,
        (∀ j : Fin schema.ambDim, σ.entryAt input [j.val] = some (.fin (x j))) ∧
        ∃ k : Fin schema.cells.length,
          MaskRegion schema (schema.cells.get k).1 (schema.cells.get k).2 x ∧
          ∀ j : Fin outputs.length,
            σ.entryFlat (outputs.get j) 0 =
              some (.fin (if j.val = k.val then 1 else 0))
  | .reshape _newShape input output =>
      ∀ i : ℕ, σ.entryFlat output i = σ.entryFlat input i
  | .broadcast toShape input output =>
      ∃ τ, H.wireType input = some τ ∧
        ∀ c ∈ Shape.positions toShape,
          σ.entryAt output c = σ.entryAt input (Shape.bcastBack τ.shape toShape c)
  | .copy _k input outputs =>
      ∀ o ∈ outputs, ∀ i : ℕ, σ.entryFlat o i = σ.entryFlat input i
  | .delete _ => True
  | .swap in1 in2 out1 out2 =>
      (∀ i : ℕ, σ.entryFlat out1 i = σ.entryFlat in2 i) ∧
      (∀ i : ℕ, σ.entryFlat out2 i = σ.entryFlat in1 i)
  | .stack _k shape inputs output =>

      ∀ i : Fin inputs.length, ∀ j : Fin (Shape.totalSize shape),
        σ.entryFlat output (i.val * Shape.totalSize shape + j.val)
          = σ.entryFlat inputs[i] j.val
  …
\end{leansrc}
\par\medskip\noindent{\small \texttt{PPLC.}\texttt{\bfseries PCellOutData}}\hfill{\footnotesize \path{lean/qtsl/Qtsl/Core/Catalogue.lean}, line 1628}\par\nopagebreak
\begin{leansrc}
inductive PCellOutData : Hypernet → Sequent → Type where …
\end{leansrc}
\par\medskip\noindent{\small \texttt{PPLC.}\texttt{\bfseries cellOutZeroBuild}}\hfill{\footnotesize \path{lean/qtsl/Qtsl/Core/Catalogue.lean}, line 1207}\par\nopagebreak
\begin{leansrc}
def cellOutZeroBuild (schema : SimpC) (outputs : List Wire)
    (k : Fin schema.cells.length) : Option Formula
\end{leansrc}
\par\medskip\noindent{\small \texttt{PPLC.}\texttt{\bfseries sound\_cellOut}}\hfill{\footnotesize \path{lean/qtsl/Qtsl/Soundness.lean}, line 6632}\par\nopagebreak
\begin{leansrc}
theorem sound_cellOut {H : Hypernet} {seq premise : Sequent}
    {schema : SimpC} {input : Wire} {outputs : List Wire}
    {whichJ : Fin schema.cells.length} {wv : Fin schema.nVertices}
    (ih : premise.semanticallyValid H)
    (happ : PrimApp.cell schema input outputs ∈ H.applications)
    (hinshape : ∃ τ, H.wireType input = some τ ∧
      τ.shape = [schema.ambDim])
    (houts : ∀ o ∈ outputs, ∃ τ, H.wireType o = some τ ∧
      Shape.totalSize τ.shape = 1)
    (hpre : seq.pre = premise.pre)
    (hwit : cellOutWitnessForm schema (.var input) whichJ wv premise.post)
    (hpost : cellOutZeroForm schema outputs whichJ seq.post) :
    seq.semanticallyValid H
\end{leansrc}
\par\medskip\noindent{\small \texttt{PPLC.}\texttt{\bfseries cellOutWitnessBuild}}\hfill{\footnotesize \path{lean/qtsl/Qtsl/Core/Catalogue.lean}, line 1193}\par\nopagebreak
\begin{leansrc}
def cellOutWitnessBuild (schema : SimpC) (t : Term)
    (k : Fin schema.cells.length) (w : Fin schema.nVertices) :
    Option Formula
\end{leansrc}
\par\medskip\noindent{\small \texttt{PPLC.}\texttt{\bfseries cellOutWitnessBuild\_spec}}\hfill{\footnotesize \path{lean/qtsl/Qtsl/Soundness.lean}, line 6145}\par\nopagebreak
\begin{leansrc}
theorem cellOutWitnessBuild_spec {schema : SimpC} {t : Term}
    {k : Fin schema.cells.length} {w : Fin schema.nVertices} {F : Formula}
    (h : cellOutWitnessForm schema t k w F) :
    schema.ambDim = 1 ∧
    (schema.cells.get k).1 w = true ∧
    ∃ (a b : Fin schema.nVertices) (V₀ V₁ : ℚ),
      (List.finRange schema.nVertices).filter
        (fun v => (schema.cells.get k).1 v) = [a, b] ∧
      (∀ j, schema.coord a j = V₀) ∧
      (∀ j, schema.coord b j = V₁) ∧
      V₀ ≠ V₁ ∧
      ((w = a ∧ F = .atom (if (schema.cells.get k).2 a then .gt else .lt)
          (.add (.const (V₁ / (V₁ - V₀))) (.mul (.const (-1 / (V₁ - V₀))) t))
          (.const 0)) ∨
       (w = b ∧ F = .atom (if (schema.cells.get k).2 b then .gt else .lt)
          (.add (.const (-V₀ / (V₁ - V₀))) (.mul (.const (1 / (V₁ - V₀))) t))
          (.const 0)))
\end{leansrc}
\par\medskip\noindent{\small \texttt{PPLC.Formula.}\texttt{\bfseries evalClosed}}\hfill{\footnotesize \path{lean/qtsl/Qtsl/Core/Catalogue.lean}, line 441}\par\nopagebreak
\begin{leansrc}
def Formula.evalClosed : Formula → Option Bool
  | .atom b l r => do pure (b.evalQ (← l.evalQ?) (← r.evalQ?))
  | .inRegion _ _ _ => none
  | .and l r => do pure ((← l.evalClosed) && (← r.evalClosed))
  | .or l r => do pure ((← l.evalClosed) || (← r.evalClosed))
  | .tt => some true
  | .ff => some false

def baryLambdas1D (cx : SimpC) (dContains : Fin cx.nVertices → Bool)
    (t : Term) : Option (List (Fin cx.nVertices × Term))
\end{leansrc}
\par\medskip\noindent{\small \texttt{PPLC.Formula.}\texttt{\bfseries holds\_of\_evalClosed\_true}}\hfill{\footnotesize \path{lean/qtsl/Qtsl/Soundness.lean}, line 2728}\par\nopagebreak
\begin{leansrc}
theorem Formula.holds_of_evalClosed_true {H : Hypernet} (σ : H.Valuation) :
    ∀ (F : Formula), F.evalClosed = some true → F.holds σ
\end{leansrc}
\par\medskip\noindent{\small \texttt{PPLC.Formula.}\texttt{\bfseries holds\_of\_evalClosed\_true}}\hfill{\footnotesize \path{lean/qtsl/Qtsl/Soundness.lean}, line 2728}\par\nopagebreak
\begin{leansrc}
theorem Formula.holds_of_evalClosed_true {H : Hypernet} (σ : H.Valuation) :
    ∀ (F : Formula), F.evalClosed = some true → F.holds σ
\end{leansrc}
\par\medskip\noindent{\small \texttt{PPLC.Formula.}\texttt{\bfseries evalClosed}}\hfill{\footnotesize \path{lean/qtsl/Qtsl/Core/Catalogue.lean}, line 441}\par\nopagebreak
\begin{leansrc}
def Formula.evalClosed : Formula → Option Bool
  | .atom b l r => do pure (b.evalQ (← l.evalQ?) (← r.evalQ?))
  | .inRegion _ _ _ => none
  | .and l r => do pure ((← l.evalClosed) && (← r.evalClosed))
  | .or l r => do pure ((← l.evalClosed) || (← r.evalClosed))
  | .tt => some true
  | .ff => some false

def baryLambdas1D (cx : SimpC) (dContains : Fin cx.nVertices → Bool)
    (t : Term) : Option (List (Fin cx.nVertices × Term))
\end{leansrc}
\par\medskip\noindent{\small \texttt{PPLC.}\texttt{\bfseries SyzygyStep}}\hfill{\footnotesize \path{lean/qtsl/Qtsl/Soundness.lean}, line 7150}\par\nopagebreak
\begin{leansrc}
structure SyzygyStep (H : Hypernet) where …
\end{leansrc}
\par\medskip\noindent{\small \texttt{PPLC.}\texttt{\bfseries sound\_syzygy}}\hfill{\footnotesize \path{lean/qtsl/Qtsl/Soundness.lean}, line 7205}\par\nopagebreak
\begin{leansrc}
theorem sound_syzygy {H : Hypernet} (st : SyzygyStep H)
    (hcomps : ∀ c ∈ st.components, c.semanticallyValid H) :
    st.seq.semanticallyValid H
\end{leansrc}
\par\medskip\noindent{\small \texttt{PPLC.}\texttt{\bfseries syzygyConclusionForm}}\hfill{\footnotesize \path{lean/qtsl/Qtsl/Core/Catalogue.lean}, line 559}\par\nopagebreak
\begin{leansrc}
def syzygyConclusionForm (conclusion : Term) (F : Formula) : Prop
\end{leansrc}
\par\medskip\noindent{\small \texttt{PPLC.}\texttt{\bfseries syzygyCheck}}\hfill{\footnotesize \path{lean/qtsl/Qtsl/Core/Catalogue.lean}, line 1026}\par\nopagebreak
\begin{leansrc}
def syzygyCheck (cc : Term) (components : List Sequent) (premise : Formula)
    (conclusion : Term) : Bool
\end{leansrc}
\par\medskip\noindent{\small \texttt{PPLC.}\texttt{\bfseries PBoundData}}\hfill{\footnotesize \path{lean/qtsl/Qtsl/Core/Catalogue.lean}, line 1690}\par\nopagebreak
\begin{leansrc}
inductive PBoundData : Hypernet → Sequent → Type where …
\end{leansrc}
\par\medskip\noindent{\small \texttt{PPLC.}\texttt{\bfseries syzygyCombinationValid}}\hfill{\footnotesize \path{lean/qtsl/Qtsl/Core/Catalogue.lean}, line 1032}\par\nopagebreak
\begin{leansrc}
def syzygyCombinationValid (cc : Term) (components : List Sequent)
    (premise : Formula) (conclusion : Term) : Prop
\end{leansrc}
\par\medskip\noindent{\small \texttt{PPLC.}\texttt{\bfseries holds\_of\_syzygyConclusionForm}}\hfill{\footnotesize \path{lean/qtsl/Qtsl/Soundness.lean}, line 7165}\par\nopagebreak
\begin{leansrc}
theorem holds_of_syzygyConclusionForm {H : Hypernet} {σ : H.Valuation}
    {conclusion : Term} {F : Formula}
    (hle : Formula.holds σ (.atom .le conclusion (.const 0)))
    (hform : syzygyConclusionForm conclusion F) :
    Formula.holds σ F
\end{leansrc}
\par\medskip\noindent{\small \texttt{PPLC.PsdCert.}\texttt{\bfseries conclCheck}}\hfill{\footnotesize \path{lean/qtsl/Qtsl/Core/Catalogue.lean}, line 1105}\par\nopagebreak
\begin{leansrc}
def conclCheck (n : Nat) (M : Fin n → Fin n → Rat) (ts : List Term)
    (conclusion : Term) : Bool
\end{leansrc}
\par\medskip\noindent{\small \texttt{PPLC.Term.}\texttt{\bfseries toPoly}}\hfill{\footnotesize \path{lean/qtsl/Qtsl/Core/Catalogue.lean}, line 414}\par\nopagebreak
\begin{leansrc}
def Term.toPoly : Term → QArith.Poly
  | .var w => QArith.Poly.var w
  | .const c => QArith.Poly.const c
  | .add l r => QArith.Poly.add l.toPoly r.toPoly
  | .mul l r => QArith.Poly.mul l.toPoly r.toPoly
  | .sub l r => QArith.Poly.sub l.toPoly r.toPoly
  | .neg t => QArith.Poly.neg t.toPoly

  | .idx _ _ => QArith.Poly.const 0

def Term.evalQ? : Term → Option Rat
  | .var _ => none
  | .const c => some c
  | .add l r => do pure ((← l.evalQ?) + (← r.evalQ?))
  | .mul l r => do pure ((← l.evalQ?) * (← r.evalQ?))
  | .sub l r => do pure ((← l.evalQ?) - (← r.evalQ?))
  | .neg t => do pure (-(← t.evalQ?))
  | .idx _ _ => none

def BowTie.evalQ : BowTie → Rat → Rat → Bool
  | .lt, x, y => decide (x < y)
  | .le, x, y => decide (x ≤ y)
  | .eq, x, y => decide (x = y)
  | .ge, x, y => decide (y ≤ x)
  | .gt, x, y => decide (y < x)
  | .ne, x, y => decide (¬ x = y)

def Formula.evalClosed : Formula → Option Bool
  | .atom b l r => do pure (b.evalQ (← l.evalQ?) (← r.evalQ?))
  | .inRegion _ _ _ => none
  | .and l r => do pure ((← l.evalClosed) && (← r.evalClosed))
  | .or l r => do pure ((← l.evalClosed) || (← r.evalClosed))
  | .tt => some true
  | .ff => some false

def baryLambdas1D (cx : SimpC) (dContains : Fin cx.nVertices → Bool)
    (t : Term) : Option (List (Fin cx.nVertices × Term))
\end{leansrc}
\par\medskip\noindent{\small \texttt{PPLC.PsdCert.}\texttt{\bfseries conclCheckCore}}\hfill{\footnotesize \path{lean/qtsl/Qtsl/Core/Catalogue.lean}, line 1092}\par\nopagebreak
\begin{leansrc}
def conclCheckCore (n : Nat) (M : Fin n → Fin n → Rat) (ts : List Term)
    (conclusion : Term) : Bool
\end{leansrc}
\par\medskip\noindent{\small \texttt{PPLC.PsdCert.}\texttt{\bfseries conclCheckW}}\hfill{\footnotesize \path{lean/qtsl/Qtsl/Core/Catalogue.lean}, line 1098}\par\nopagebreak
\begin{leansrc}
def conclCheckW (n : Nat) (M : Fin n → Fin n → Rat) (ts : List Term)
    (conclusion : Term) : Bool
\end{leansrc}
\par\medskip\noindent{\small \texttt{PPLC.}\texttt{\bfseries sound\_psdCertOf}}\hfill{\footnotesize \path{lean/qtsl/Qtsl/Soundness.lean}, line 7466}\par\nopagebreak
\begin{leansrc}
theorem sound_psdCertOf {H : Hypernet} {seq : Sequent} {a : Nat}
    {M : Fin a → Fin a → ℚ} {ts : List Term} {conclusion : Term}
    (hpsd : ∀ x : Fin a → ℚ, 0 ≤ PsdCert.finQSum a fun i =>
      PsdCert.finQSum a fun j => M i j * x i * x j)
    (hconcl : PsdCert.conclCheck a M ts conclusion = true)
    (hshape : ∃ s, conclusion.bshape H = some s)
    (hfin : conclusion.finWires H)
    (hwpos : ∀ w ∈ conclusion.freeWires,
      ∃ τ, H.wireType w = some τ ∧ 0 < Shape.totalSize τ.shape)
    (hform : syzygyConclusionForm conclusion seq.post) :
    seq.semanticallyValid H
\end{leansrc}
\par\medskip\noindent{\small \texttt{PPLC.PsdCert.}\texttt{\bfseries congruence\_eq\_diagonal}}\hfill{\footnotesize \path{lean/qtsl/Qtsl/PsdCert.lean}, line 41}\par\nopagebreak
\begin{leansrc}
theorem congruence_eq_diagonal {A Z : Matrix n n ℝ}
    (hA : ∀ i j, A i j = A j i)
    (hZu : ∀ ⦃i j⦄, j < i → Z i j = 0)
    (hMl : ∀ ⦃i j⦄, i < j → (A * Z) i j = 0) :
    Zᴴ * (A * Z) = diagonal fun i => Z i i * (A * Z) i i
\end{leansrc}
\par\medskip\noindent{\small \texttt{PPLC.PsdCert.}\texttt{\bfseries posDef\_of\_congruence\_posDef}}\hfill{\footnotesize \path{lean/qtsl/Qtsl/PsdCert.lean}, line 77}\par\nopagebreak
\begin{leansrc}
theorem posDef_of_congruence_posDef {A Z : Matrix n n ℝ}
    (hZu : ∀ ⦃i j⦄, j < i → Z i j = 0)
    (hZd : ∀ i, 0 < Z i i)
    (hBd : (Zᴴ * (A * Z)).PosDef) :
    A.PosDef
\end{leansrc}
\par\medskip\noindent{\small \texttt{PPLC.PsdCert.}\texttt{\bfseries posDef\_of\_triangular\_witness}}\hfill{\footnotesize \path{lean/qtsl/Qtsl/PsdCert.lean}, line 98}\par\nopagebreak
\begin{leansrc}
theorem posDef_of_triangular_witness {A Z : Matrix n n ℝ}
    (hA : ∀ i j, A i j = A j i)
    (hZu : ∀ ⦃i j⦄, j < i → Z i j = 0)
    (hZd : ∀ i, 0 < Z i i)
    (hMl : ∀ ⦃i j⦄, i < j → (A * Z) i j = 0)
    (hMd : ∀ i, 0 < (A * Z) i i) :
    A.PosDef
\end{leansrc}
\par\medskip\noindent{\small \texttt{PPLC.PsdCert.}\texttt{\bfseries quadForm\_nonneg\_of\_posSemidef}}\hfill{\footnotesize \path{lean/qtsl/Qtsl/PsdCert.lean}, line 431}\par\nopagebreak
\begin{leansrc}
theorem quadForm_nonneg_of_posSemidef {a : Nat} {M : Fin a → Fin a → ℚ}
    (h : (matR M).PosSemidef) (x : Fin a → ℚ) :
    0 ≤ finQSum a (fun i => finQSum a (fun j => M i j * x i * x j))
\end{leansrc}
\par\medskip\noindent{\small \texttt{PPLC.PsdCert.}\texttt{\bfseries symB}}\hfill{\footnotesize \path{lean/qtsl/Qtsl/Core/Catalogue.lean}, line 1047}\par\nopagebreak
\begin{leansrc}
def symB (n : Nat) (A : Fin n → Fin n → Rat) : Bool
\end{leansrc}
\par\medskip\noindent{\small \texttt{PPLC.PsdCert.}\texttt{\bfseries upperTriPosDiagB}}\hfill{\footnotesize \path{lean/qtsl/Qtsl/Core/Catalogue.lean}, line 1051}\par\nopagebreak
\begin{leansrc}
def upperTriPosDiagB (n : Nat) (Z : Fin n → Fin n → Rat) : Bool
\end{leansrc}
\par\medskip\noindent{\small \texttt{PPLC.PsdCert.}\texttt{\bfseries mulEntry}}\hfill{\footnotesize \path{lean/qtsl/Qtsl/Core/Catalogue.lean}, line 1044}\par\nopagebreak
\begin{leansrc}
def mulEntry (n : Nat) (A B : Fin n → Fin n → Rat) (i j : Fin n) : Rat
\end{leansrc}
\par\medskip\noindent{\small \texttt{PPLC.PsdCert.}\texttt{\bfseries conclCheck}}\hfill{\footnotesize \path{lean/qtsl/Qtsl/Core/Catalogue.lean}, line 1105}\par\nopagebreak
\begin{leansrc}
def conclCheck (n : Nat) (M : Fin n → Fin n → Rat) (ts : List Term)
    (conclusion : Term) : Bool
\end{leansrc}
\par\medskip\noindent{\small \texttt{PPLC.PsdCert.}\texttt{\bfseries negQuadFormPoly}}\hfill{\footnotesize \path{lean/qtsl/Qtsl/Core/Catalogue.lean}, line 1085}\par\nopagebreak
\begin{leansrc}
def negQuadFormPoly (n : Nat) (M : Fin n → Fin n → Rat)
    (ps : List QArith.Poly) : QArith.Poly
\end{leansrc}
\par\medskip\noindent{\small \texttt{PPLC.PsdCert.}\texttt{\bfseries posDef\_of\_dominant}}\hfill{\footnotesize \path{lean/qtsl/Qtsl/PsdCert.lean}, line 126}\par\nopagebreak
\begin{leansrc}
theorem posDef_of_dominant [DecidableEq n] {B : Matrix n n ℝ}
    (hsym : ∀ i j, B i j = B j i)
    (hdom : ∀ i, (∑ j ∈ Finset.univ.erase i, |B i j|) < B i i) :
    B.PosDef
\end{leansrc}
\par\medskip\noindent{\small \texttt{PPLC.PsdCert.}\texttt{\bfseries posDef\_of\_dd\_witness}}\hfill{\footnotesize \path{lean/qtsl/Qtsl/PsdCert.lean}, line 213}\par\nopagebreak
\begin{leansrc}
theorem posDef_of_dd_witness [DecidableEq n] {A Z : Matrix n n ℝ}
    (hA : ∀ i j, A i j = A j i)
    (hZu : ∀ ⦃i j⦄, j < i → Z i j = 0)
    (hZd : ∀ i, 0 < Z i i)
    (hdom : ∀ i, (∑ j ∈ Finset.univ.erase i, |(Zᴴ * (A * Z)) i j|)
              < (Zᴴ * (A * Z)) i i) :
    A.PosDef
\end{leansrc}
\par\medskip\noindent{\small \texttt{PPLC.PsdCert.}\texttt{\bfseries finQSum}}\hfill{\footnotesize \path{lean/qtsl/Qtsl/Core/Catalogue.lean}, line 1042}\par\nopagebreak
\begin{leansrc}
def finQSum (n : Nat) (f : Fin n → Rat) : Rat
\end{leansrc}
\par\medskip\noindent{\small \texttt{PPLC.PsdCert.}\texttt{\bfseries posDef\_of\_dominant}}\hfill{\footnotesize \path{lean/qtsl/Qtsl/PsdCert.lean}, line 126}\par\nopagebreak
\begin{leansrc}
theorem posDef_of_dominant [DecidableEq n] {B : Matrix n n ℝ}
    (hsym : ∀ i j, B i j = B j i)
    (hdom : ∀ i, (∑ j ∈ Finset.univ.erase i, |B i j|) < B i i) :
    B.PosDef
\end{leansrc}
\par\medskip\noindent{\small \texttt{PPLC.PsdCert.}\texttt{\bfseries congEntry}}\hfill{\footnotesize \path{lean/qtsl/Qtsl/Core/Catalogue.lean}, line 1062}\par\nopagebreak
\begin{leansrc}
def congEntry (n : Nat) (A Z : Fin n → Fin n → Rat) (i j : Fin n) : Rat
\end{leansrc}
\par\medskip\noindent{\small \texttt{PPLC.PsdCert.}\texttt{\bfseries posSemidef\_of\_factor}}\hfill{\footnotesize \path{lean/qtsl/Qtsl/PsdCert.lean}, line 233}\par\nopagebreak
\begin{leansrc}
theorem posSemidef_of_factor {m' n' : Type*} [Fintype m'] [Fintype n']
    {A : Matrix n' n' ℝ} {B : Matrix m' m' ℝ} {L : Matrix m' n' ℝ}
    (hB : B.PosSemidef) (hA : A = Lᴴ * B * L) : A.PosSemidef
\end{leansrc}
\par\medskip\noindent{\small \texttt{PPLC.PsdCert.}\texttt{\bfseries kernelExt}}\hfill{\footnotesize \path{lean/qtsl/Qtsl/Core/Catalogue.lean}, line 1071}\par\nopagebreak
\begin{leansrc}
def kernelExt (m k : Nat) (N : Fin m → Fin k → Rat) (a : Fin m)
    (j : Fin (m + k)) : Rat
\end{leansrc}
\par\medskip\noindent{\small \texttt{PPLC.PsdCert.}\texttt{\bfseries posSemidef\_of\_factor}}\hfill{\footnotesize \path{lean/qtsl/Qtsl/PsdCert.lean}, line 233}\par\nopagebreak
\begin{leansrc}
theorem posSemidef_of_factor {m' n' : Type*} [Fintype m'] [Fintype n']
    {A : Matrix n' n' ℝ} {B : Matrix m' m' ℝ} {L : Matrix m' n' ℝ}
    (hB : B.PosSemidef) (hA : A = Lᴴ * B * L) : A.PosSemidef
\end{leansrc}
\par\medskip\noindent{\small \texttt{PPLC.}\texttt{\bfseries sound\_reluSplit}}\hfill{\footnotesize \path{lean/qtsl/Qtsl/ReluSplit.lean}, line 15}\par\nopagebreak
\begin{leansrc}
theorem sound_reluSplit {H : Hypernet}
    {seq activePremise inactivePremise : Sequent}
    {preact post marker : Wire} {bodyForm : Formula}
    (ihActive : activePremise.semanticallyValid H)
    (ihInactive : inactivePremise.semanticallyValid H)
    (hlink : reluLinked H preact post)
    (hrank0 : ∃ τ, H.wireType post = some τ ∧ τ.shape = [])
    (hActPre : activePremise.pre
      = .and seq.pre (.atom .ge (.var preact) (.const 0)))
    (hInactPre : inactivePremise.pre
      = .and seq.pre (.atom .le (.var preact) (.const 0)))
    (hActPost : substituteInFormula bodyForm marker (.var preact)
      activePremise.post)
    (hInactPost : substituteInFormula bodyForm marker (.const 0)
      inactivePremise.post)
    (hconcl : substituteInFormula bodyForm marker (.var post) seq.post) :
    seq.semanticallyValid H
\end{leansrc}
\par\medskip\noindent{\small \texttt{PPLC.}\texttt{\bfseries reluSplit\_rank0\_necessary}}\hfill{\footnotesize \path{lean/qtsl/Qtsl/ReluSplit.lean}, line 153}\par\nopagebreak
\begin{leansrc}
theorem reluSplit_rank0_necessary :
    ¬ (∀ (H : Hypernet) (seq activePremise inactivePremise : Sequent)
        (preact post marker : Wire) (bodyForm : Formula),
        activePremise.semanticallyValid H →
        inactivePremise.semanticallyValid H →
        reluLinked H preact post →
        activePremise.pre
          = .and seq.pre (.atom .ge (.var preact) (.const 0)) →
        inactivePremise.pre
          = .and seq.pre (.atom .le (.var preact) (.const 0)) →
        substituteInFormula bodyForm marker (.var preact)
          activePremise.post →
        substituteInFormula bodyForm marker (.const 0)
          inactivePremise.post →
        substituteInFormula bodyForm marker (.var post) seq.post →
        seq.semanticallyValid H)
\end{leansrc}
\par\medskip\noindent{\small \texttt{PPLC.ReluSplitScope.}\texttt{\bfseries HSplit}}\hfill{\footnotesize \path{lean/qtsl/Qtsl/ReluSplit.lean}, line 95}\par\nopagebreak
\begin{leansrc}
def HSplit : Hypernet where …
\end{leansrc}
\par\medskip\noindent{\small \texttt{PPLC.}\texttt{\bfseries PromiseStep}}\hfill{\footnotesize \path{lean/qtsl/Qtsl/Soundness.lean}, line 7116}\par\nopagebreak
\begin{leansrc}
structure PromiseStep where …
\end{leansrc}
\par\medskip\noindent{\small \texttt{PPLC.}\texttt{\bfseries sound\_promise}}\hfill{\footnotesize \path{lean/qtsl/Qtsl/Soundness.lean}, line 7132}\par\nopagebreak
\begin{leansrc}
theorem sound_promise {H : Hypernet} (st : PromiseStep)
    (ihs : st.supplier.semanticallyValid H)
    (ihc : st.consumer.semanticallyValid H) :
    st.seq.semanticallyValid H
\end{leansrc}
\par\medskip\noindent{\small \texttt{PPLC.Formula.}\texttt{\bfseries conjGet\_of\_conjReplace}}\hfill{\footnotesize \path{lean/qtsl/Qtsl/Soundness.lean}, line 6994}\par\nopagebreak
\begin{leansrc}
theorem Formula.conjGet_of_conjReplace :
    ∀ {π : ConjPath} {Γ Γ' g : Formula},
    Γ.conjReplace? π g = some Γ' → Γ'.conjGet? π = some g
\end{leansrc}
\par\medskip\noindent{\small \texttt{PPLC.Formula.}\texttt{\bfseries holds\_of\_conjGet}}\hfill{\footnotesize \path{lean/qtsl/Qtsl/Soundness.lean}, line 6916}\par\nopagebreak
\begin{leansrc}
theorem Formula.holds_of_conjGet {H : Hypernet} {σ : H.Valuation} :
    ∀ {π : ConjPath} {Γ d : Formula},
    Γ.conjGet? π = some d → Γ.holds σ → d.holds σ
\end{leansrc}
\par\medskip\noindent{\small \texttt{PPLC.Formula.}\texttt{\bfseries conjReplace\_back}}\hfill{\footnotesize \path{lean/qtsl/Qtsl/Soundness.lean}, line 7042}\par\nopagebreak
\begin{leansrc}
theorem Formula.conjReplace_back :
    ∀ {π : ConjPath} {Γ Γ' q g : Formula},
    Γ.conjGet? π = some q → Γ.conjReplace? π g = some Γ' →
    Γ'.conjReplace? π q = some Γ
\end{leansrc}
\par\medskip\noindent{\small \texttt{PPLC.Formula.}\texttt{\bfseries holds\_conjReplace}}\hfill{\footnotesize \path{lean/qtsl/Qtsl/Soundness.lean}, line 6949}\par\nopagebreak
\begin{leansrc}
theorem Formula.holds_conjReplace {H : Hypernet} {σ : H.Valuation} :
    ∀ {π : ConjPath} {Γ Γ' g : Formula},
    Γ.conjReplace? π g = some Γ' → Γ.holds σ → g.holds σ → Γ'.holds σ
\end{leansrc}
\par\medskip\noindent{\small \texttt{PPLC.}\texttt{\bfseries formulaWiresWithin}}\hfill{\footnotesize \path{lean/qtsl/Qtsl/Core/Catalogue.lean}, line 953}\par\nopagebreak
\begin{leansrc}
def formulaWiresWithin (f : Formula) (iface : List Wire) : Prop
\end{leansrc}
\par\medskip\noindent{\small \texttt{PPLC.}\texttt{\bfseries wiresSeparated}}\hfill{\footnotesize \path{lean/qtsl/Qtsl/Core/Catalogue.lean}, line 956}\par\nopagebreak
\begin{leansrc}
def wiresSeparated (left ctx post : Formula) (iface : List Wire) : Prop
\end{leansrc}
\par\medskip\noindent{\small \texttt{PPLC.}\texttt{\bfseries syzygyCheckCore}}\hfill{\footnotesize \path{lean/qtsl/Qtsl/Core/Catalogue.lean}, line 964}\par\nopagebreak
\begin{leansrc}
def syzygyCheckCore (cc : Term) (components : List Sequent) (premise : Formula)
    (conclusion : Term) : Bool
\end{leansrc}
\par\medskip\noindent{\small \texttt{PPLC.}\texttt{\bfseries syzygyCheckW}}\hfill{\footnotesize \path{lean/qtsl/Qtsl/Core/Catalogue.lean}, line 982}\par\nopagebreak
\begin{leansrc}
def syzygyCheckW (cc : Term) (components : List Sequent) (premise : Formula)
    (conclusion : Term) : Bool
\end{leansrc}
\par\medskip\noindent{\small \texttt{PPLC.}\texttt{\bfseries syzygyCheckM}}\hfill{\footnotesize \path{lean/qtsl/Qtsl/Core/Catalogue.lean}, line 1004}\par\nopagebreak
\begin{leansrc}
def syzygyCheckM (cc : Term) (components : List Sequent) (premise : Formula)
    (conclusion : Term) : Bool
\end{leansrc}
\par\medskip\noindent{\small \texttt{PPLC.}\texttt{\bfseries premiseContains}}\hfill{\footnotesize \path{lean/qtsl/Qtsl/Core/Catalogue.lean}, line 950}\par\nopagebreak
\begin{leansrc}
def premiseContains (pre sub : Formula) : Prop
\end{leansrc}
\par\medskip\noindent{\small \texttt{PPLC.}\texttt{\bfseries atomLeZeroPolys\_sound\_entryC}}\hfill{\footnotesize \path{lean/qtsl/Qtsl/Soundness.lean}, line 3380}\par\nopagebreak
\begin{leansrc}
theorem atomLeZeroPolys_sound_entryC {H : Hypernet} (σ : H.Valuation)
    {F : Formula} (hF : Formula.holds σ F)
    (hw : ∀ w ∈ F.freeWires, ∃ τ, H.wireType w = some τ ∧
      0 < Shape.totalSize τ.shape)
    (hix : F.atomSidesIdxFree = true) (i : ℕ) :
    ∀ pa ∈ atomLeZeroPolys F, QArith.Poly.eval (σ.entryValC i) pa ≤ 0
\end{leansrc}
\par\medskip\noindent{\small \texttt{PPLC.}\texttt{\bfseries wire\_discipline\_of\_bshape\_fin}}\hfill{\footnotesize \path{lean/qtsl/Qtsl/Soundness.lean}, line 2265}\par\nopagebreak
\begin{leansrc}
theorem wire_discipline_of_bshape_fin {H : Hypernet} (σ : H.Valuation)
    {t : Term} {s : Shape} (hb : t.bshape H = some s)
    (hf : t.finWires H) (hix : t.idxFree = true) :
    ∀ w ∈ t.freeWires, ∃ qw, σ.wireQ w = some qw ∧
      (qw.shape = [] ∨ qw.shape = s)
\end{leansrc}
\par\medskip\noindent{\small \texttt{PPLC.Term.}\texttt{\bfseries evalQ\_total\_singleton}}\hfill{\footnotesize \path{lean/qtsl/Qtsl/Soundness.lean}, line 1009}\par\nopagebreak
\begin{leansrc}
theorem Term.evalQ_total_singleton {H : Hypernet} (σ : H.Valuation) :
    ∀ (t : Term),
    (∀ w ∈ t.freeWires, ∃ qw, σ.wireQ w = some qw ∧ SingletonShape qw.shape) →
    t.idxFree = true →
    ∃ q, Term.evalQ σ t = some q ∧ SingletonShape q.shape ∧
      q.scalar = QArith.Poly.eval σ.rho t.toPoly
\end{leansrc}
\par\medskip\noindent{\small \texttt{PPLC.SmallGainStep.}\texttt{\bfseries toSyzygy}}\hfill{\footnotesize \path{lean/qtsl/Qtsl/Composition.lean}, line 409}\par\nopagebreak
\begin{leansrc}
def toSyzygy (st : SmallGainStep H) : SyzygyStep H where …
\end{leansrc}
\par\medskip\noindent{\small \texttt{PPLC.}\texttt{\bfseries sound\_smallGain}}\hfill{\footnotesize \path{lean/qtsl/Qtsl/Composition.lean}, line 435}\par\nopagebreak
\begin{leansrc}
theorem sound_smallGain {H : Hypernet} (st : SmallGainStep H)
    (h1 : st.diss1.semanticallyValid H)
    (h2 : st.diss2.semanticallyValid H)
    (hm1 : st.mem1.semanticallyValid H)
    (hm2 : st.mem2.semanticallyValid H) :
    st.seq.semanticallyValid H
\end{leansrc}
\par\medskip\noindent{\small \texttt{PPLC.SmallGainStep.}\texttt{\bfseries toSyzygy}}\hfill{\footnotesize \path{lean/qtsl/Qtsl/Composition.lean}, line 409}\par\nopagebreak
\begin{leansrc}
def toSyzygy (st : SmallGainStep H) : SyzygyStep H where …
\end{leansrc}
\par\medskip\noindent{\small \texttt{PPLC.}\texttt{\bfseries RuleTag}}\hfill{\footnotesize \path{lean/qtsl/Qtsl/Core/Catalogue.lean}, line 1291}\par\nopagebreak
\begin{leansrc}
inductive RuleTag where …
\end{leansrc}
\par\medskip\noindent{\small \texttt{PPLC.}\texttt{\bfseries SmallGainStep}}\hfill{\footnotesize \path{lean/qtsl/Qtsl/Composition.lean}, line 253}\par\nopagebreak
\begin{leansrc}
structure SmallGainStep (H : Hypernet) where …
\end{leansrc}
\par\medskip\noindent{\small \texttt{PPLC.SmallGainStep.}\texttt{\bfseries combTerm}}\hfill{\footnotesize \path{lean/qtsl/Qtsl/Composition.lean}, line 322}\par\nopagebreak
\begin{leansrc}
def combTerm (st : SmallGainStep H) : Term
\end{leansrc}
\par\medskip\noindent{\small \texttt{PPLC.SmallGainStep.}\texttt{\bfseries combValid}}\hfill{\footnotesize \path{lean/qtsl/Qtsl/Composition.lean}, line 338}\par\nopagebreak
\begin{leansrc}
theorem combValid (st : SmallGainStep H) :
    syzygyCombinationValid st.combTerm st.components st.seq.pre st.concl
\end{leansrc}
\par\medskip\noindent{\small \texttt{PPLC.SmallGainStep.}\texttt{\bfseries toSyzygy\_components}}\hfill{\footnotesize \path{lean/qtsl/Qtsl/Composition.lean}, line 430}\par\nopagebreak
\begin{leansrc}
theorem toSyzygy_components (st : SmallGainStep H) :
    st.toSyzygy.components = st.components
\end{leansrc}
\par\medskip\noindent{\small \texttt{PPLC.}\texttt{\bfseries smallGainN\_is\_syzygy}}\hfill{\footnotesize \path{lean/qtsl/Qtsl/Composition.lean}, line 582}\par\nopagebreak
\begin{leansrc}
theorem smallGainN_is_syzygy (kk : ℕ) (LV A Sg : Fin kk → Poly)
    (lam : Fin kk → ℚ) (B : Fin kk → Fin kk → ℚ) (th : ℚ)
    (hcol : ∀ j, ((List.finRange kk).map (fun i => B i j)).sum = lam j) :
    (((List.finRange kk).map
        (fun i => (lam i, Poly.sub (Poly.add (LV i) (A i)) (Sg i)))
      ++ (List.finRange kk).map
        (fun i => ((1 : ℚ), smallGainNMemPoly kk A Sg lam B th i))).foldr
      (fun pr acc => Poly.add (Poly.scale pr.1 pr.2) acc) [])
    = Poly.add
        ((List.finRange kk).foldr
          (fun i acc => Poly.add (Poly.scale (lam i) (LV i)) acc) [])
        (Poly.scale th
          ((List.finRange kk).foldr
            (fun i acc => Poly.add (Poly.scale (lam i) (A i)) acc) []))
\end{leansrc}
\par\medskip\noindent{\small \texttt{PPLC.SmallGainNStep.}\texttt{\bfseries toSyzygy}}\hfill{\footnotesize \path{lean/qtsl/Qtsl/Composition.lean}, line 796}\par\nopagebreak
\begin{leansrc}
def toSyzygy (st : SmallGainNStep H kk) : SyzygyStep H where …
\end{leansrc}
\par\medskip\noindent{\small \texttt{PPLC.}\texttt{\bfseries sound\_smallGainN}}\hfill{\footnotesize \path{lean/qtsl/Qtsl/Composition.lean}, line 825}\par\nopagebreak
\begin{leansrc}
theorem sound_smallGainN {H : Hypernet} {kk : ℕ} (st : SmallGainNStep H kk)
    (hdiss : ∀ i, (st.diss i).semanticallyValid H)
    (hmem : ∀ i, (st.mem i).semanticallyValid H) :
    st.seq.semanticallyValid H
\end{leansrc}
\par\medskip\noindent{\small \texttt{PPLC.}\texttt{\bfseries smallGainN\_colSum\_mutant\_refuted}}\hfill{\footnotesize \path{lean/qtsl/Qtsl/Composition.lean}, line 857}\par\nopagebreak
\begin{leansrc}
theorem smallGainN_colSum_mutant_refuted :
    ¬ (∀ (kk : ℕ) (LV A Sg : Fin kk → Poly) (lam : Fin kk → ℚ)
        (B : Fin kk → Fin kk → ℚ) (th : ℚ),
        (((List.finRange kk).map
            (fun i => (lam i, Poly.sub (Poly.add (LV i) (A i)) (Sg i)))
          ++ (List.finRange kk).map
            (fun i => ((1 : ℚ), smallGainNMemPoly kk A Sg lam B th i))).foldr
          (fun pr acc => Poly.add (Poly.scale pr.1 pr.2) acc) [])
        = Poly.add
            ((List.finRange kk).foldr
              (fun i acc => Poly.add (Poly.scale (lam i) (LV i)) acc) [])
            (Poly.scale th
              ((List.finRange kk).foldr
                (fun i acc => Poly.add (Poly.scale (lam i) (A i)) acc) [])))
\end{leansrc}
\par\medskip\noindent{\small \texttt{Qtsl.QtslCompose.Discrete.}\texttt{\bfseries smallGainN\_drift}}\hfill{\footnotesize \path{lean/qtsl-aux/QtslAux/QtslCompose/AGComposeDiscrete.lean}, line 372}\par\nopagebreak
\begin{leansrc}
theorem smallGainN_drift (hlam : ∀ i, 0 ≤ lam i)
    (hcol : ∀ j, ∑ i, B i j = lam j)
    (hint : ∀ i, ∀ y, y ∉ A → Integrable (V i) (K y))
    (hd : ∀ i, ∀ y, y ∉ A →
        ∫ z, V i z ∂K y ≤ V i y - alpha i y + sigma i y)
    (hm : ∀ i, ∀ y, y ∉ A →
        lam i * sigma i y ≤ (1 - theta) * ∑ j, B i j * alpha j y) :
    ∀ y, y ∉ A → ∫ z, (∑ i, lam i * V i z) ∂K y
      ≤ (∑ i, lam i * V i y) - theta * ∑ j, lam j * alpha j y
\end{leansrc}
\par\medskip\noindent{\small \texttt{Qtsl.QtslCompose.Discrete.}\texttt{\bfseries smallGain\_drift\_of\_N}}\hfill{\footnotesize \path{lean/qtsl-aux/QtslAux/QtslCompose/AGComposeDiscrete.lean}, line 416}\par\nopagebreak
\begin{leansrc}
theorem smallGain_drift_of_N {V₁ V₂ alpha₁ alpha₂ sigma₁ sigma₂ : S → ℝ}
    {lam theta : ℝ} (hlam : 0 ≤ lam)
    (hint₁ : ∀ y, y ∉ A → Integrable V₁ (K y))
    (hint₂ : ∀ y, y ∉ A → Integrable V₂ (K y))
    (hd₁ : ∀ y, y ∉ A → ∫ z, V₁ z ∂K y ≤ V₁ y - alpha₁ y + sigma₁ y)
    (hd₂ : ∀ y, y ∉ A → ∫ z, V₂ z ∂K y ≤ V₂ y - alpha₂ y + sigma₂ y)
    (hm₁ : ∀ y, y ∉ A → lam * sigma₂ y ≤ (1 - theta) * alpha₁ y)
    (hm₂ : ∀ y, y ∉ A → sigma₁ y ≤ lam * ((1 - theta) * alpha₂ y)) :
    ∀ y, y ∉ A → ∫ z, (V₁ z + lam * V₂ z) ∂K y
      ≤ (V₁ y + lam * V₂ y) - theta * (alpha₁ y + lam * alpha₂ y)
\end{leansrc}
\par\medskip\noindent{\small \texttt{PPLC.}\texttt{\bfseries citeFacet}}\hfill{\footnotesize \path{lean/qtsl/Qtsl/ProductLattice.lean}, line 675}\par\nopagebreak
\begin{leansrc}
def citeFacet (s : Sequent) (rider : Formula) : Sequent
\end{leansrc}
\par\medskip\noindent{\small \texttt{PPLC.}\texttt{\bfseries prodBlock}}\hfill{\footnotesize \path{lean/qtsl/Qtsl/ProductLattice.lean}, line 678}\par\nopagebreak
\begin{leansrc}
def prodBlock (b1 b2 : BlockCert) : BlockCert
\end{leansrc}
\par\medskip\noindent{\small \texttt{PPLC.PProductLatticeData.}\texttt{\bfseries prodBlocks}}\hfill{\footnotesize \path{lean/qtsl/Qtsl/ProductLattice.lean}, line 688}\par\nopagebreak
\begin{leansrc}
def PProductLatticeData.prodBlocks {H1 H2 : Hypernet}
    (d : PProductLatticeData H1 H2) : List BlockCert
\end{leansrc}
\par\medskip\noindent{\small \texttt{PPLC.}\texttt{\bfseries CertifiedBlockLattice}}\hfill{\footnotesize \path{lean/qtsl/Qtsl/ProductLattice.lean}, line 658}\par\nopagebreak
\begin{leansrc}
structure CertifiedBlockLattice (H : Hypernet) where …
\end{leansrc}
\par\medskip\noindent{\small \texttt{PPLC.Hypernet.}\texttt{\bfseries disjUnion}}\hfill{\footnotesize \path{lean/qtsl/Qtsl/ProductLattice.lean}, line 252}\par\nopagebreak
\begin{leansrc}
def Hypernet.disjUnion (H1 H2 : Hypernet) (hd : H1.WireDisjoint H2) : Hypernet
      where …
\end{leansrc}
\par\medskip\noindent{\small \texttt{PPLC.}\texttt{\bfseries sound\_productLattice}}\hfill{\footnotesize \path{lean/qtsl/Qtsl/ProductLattice.lean}, line 712}\par\nopagebreak
\begin{leansrc}
theorem sound_productLattice {H1 H2 : Hypernet}
    (d : PProductLatticeData H1 H2)
    (hL : d.left.Valid) (hR : d.right.Valid) :
    ∀ b ∈ d.prodBlocks, ∀ s ∈ b.facets,
      s.semanticallyValid (H1.disjUnion H2 d.disjoint)
\end{leansrc}
\par\medskip\noindent{\small \texttt{PPLC.Sequent.}\texttt{\bfseries semanticallyValid\_extend}}\hfill{\footnotesize \path{lean/qtsl/Qtsl/ProductLattice.lean}, line 636}\par\nopagebreak
\begin{leansrc}
theorem Sequent.semanticallyValid_extend {H H' : Hypernet} {seq : Sequent}
    (hext : H.WireExtends H')
    (happs : ∀ app ∈ H.applications, app ∈ H'.applications)
    (hws : seq.wellScoped H) (hval : seq.semanticallyValid H) :
    seq.semanticallyValid H'
\end{leansrc}
\par\medskip\noindent{\small \texttt{PPLC.Sequent.}\texttt{\bfseries semanticallyValid\_weaken}}\hfill{\footnotesize \path{lean/qtsl/Qtsl/ProductLattice.lean}, line 648}\par\nopagebreak
\begin{leansrc}
theorem Sequent.semanticallyValid_weaken {H : Hypernet}
    {pre extra post : Formula}
    (h : Sequent.semanticallyValid H ⟨pre, post⟩) :
    Sequent.semanticallyValid H ⟨.and pre extra, post⟩
\end{leansrc}
\par\medskip\noindent{\small \texttt{PPLC.}\texttt{\bfseries sound\_productLattice\_citeLeft}}\hfill{\footnotesize \path{lean/qtsl/Qtsl/ProductLattice.lean}, line 692}\par\nopagebreak
\begin{leansrc}
theorem sound_productLattice_citeLeft {H1 H2 : Hypernet}
    (hd : H1.WireDisjoint H2) {s0 : Sequent} (rider : Formula)
    (hws : s0.wellScoped H1) (hval : s0.semanticallyValid H1) :
    (citeFacet s0 rider).semanticallyValid (H1.disjUnion H2 hd)
\end{leansrc}
\par\medskip\noindent{\small \texttt{PPLC.PProductLatticeData.}\texttt{\bfseries prodBlocks}}\hfill{\footnotesize \path{lean/qtsl/Qtsl/ProductLattice.lean}, line 688}\par\nopagebreak
\begin{leansrc}
def PProductLatticeData.prodBlocks {H1 H2 : Hypernet}
    (d : PProductLatticeData H1 H2) : List BlockCert
\end{leansrc}
\par\medskip\noindent{\small \texttt{PPLC.Hypernet.Valuation.}\texttt{\bfseries restrict}}\hfill{\footnotesize \path{lean/qtsl/Qtsl/ProductLattice.lean}, line 320}\par\nopagebreak
\begin{leansrc}
def Hypernet.Valuation.restrict {H H' : Hypernet} (σ : H'.Valuation)
    (hext : H.WireExtends H') : H.Valuation
\end{leansrc}
\par\medskip\noindent{\small \texttt{PPLC.Formula.}\texttt{\bfseries holds\_restrict}}\hfill{\footnotesize \path{lean/qtsl/Qtsl/ProductLattice.lean}, line 429}\par\nopagebreak
\begin{leansrc}
theorem Formula.holds_restrict {H H' : Hypernet} (σ : H'.Valuation)
    (hext : H.WireExtends H') :
    ∀ (F : Formula), (∀ w ∈ F.freeWires, (H.wireType w).isSome) →
      (Formula.holds (σ.restrict hext) F ↔ Formula.holds σ F)
\end{leansrc}
\par\medskip\noindent{\small \texttt{PPLC.PrimApp.}\texttt{\bfseries semOk\_restrict}}\hfill{\footnotesize \path{lean/qtsl/Qtsl/ProductLattice.lean}, line 485}\par\nopagebreak
\begin{leansrc}
theorem PrimApp.semOk_restrict {H H' : Hypernet} (σ : H'.Valuation)
    (hext : H.WireExtends H') {app : PrimApp}
    (hTC : app.TypeConsistent H.wireType) (h : app.SemOk σ) :
    app.SemOk (σ.restrict hext)
\end{leansrc}
\par\medskip\noindent{\small \texttt{PPLC.Hypernet.}\texttt{\bfseries WireDisjoint}}\hfill{\footnotesize \path{lean/qtsl/Qtsl/ProductLattice.lean}, line 17}\par\nopagebreak
\begin{leansrc}
def Hypernet.WireDisjoint (H1 H2 : Hypernet) : Prop
\end{leansrc}
\par\medskip\noindent{\small \texttt{PPLC.}\texttt{\bfseries prodBlocks\_confined}}\hfill{\footnotesize \path{lean/qtsl/Qtsl/ProductLattice.lean}, line 729}\par\nopagebreak
\begin{leansrc}
theorem prodBlocks_confined {H1 H2 : Hypernet}
    (d : PProductLatticeData H1 H2) :
    ∀ b ∈ d.prodBlocks,
      (∀ w ∈ b.region.freeWires,
        ((H1.disjUnion H2 d.disjoint).wireType w).isSome) ∧
      ∀ s ∈ b.facets, s.wellScoped (H1.disjUnion H2 d.disjoint)
\end{leansrc}
\par\medskip\noindent{\small \texttt{PPLC.CertifiedBlockLattice.}\texttt{\bfseries Valid}}\hfill{\footnotesize \path{lean/qtsl/Qtsl/ProductLattice.lean}, line 671}\par\nopagebreak
\begin{leansrc}
def CertifiedBlockLattice.Valid {H : Hypernet}
    (L : CertifiedBlockLattice H) : Prop
\end{leansrc}
\par\medskip\noindent{\small \texttt{Qtsl.Aliases.}\texttt{\bfseries sound\_QSatFromRegion}}\hfill{\footnotesize \path{lean/qtsl-aux/QtslAux/QtslAliases.lean}, line 35}\par\nopagebreak
\begin{leansrc}
theorem sound_QSatFromRegion {X : Type*} {R : Set X} {sat : X → Prop} {x₀ : X}
    (hchild : ∀ x ∈ R, sat x) (heval : x₀ ∈ R) : sat x₀
\end{leansrc}
\par\medskip\noindent{\small \texttt{Qtsl.Aliases.}\texttt{\bfseries sound\_QSatFromRegion\_sublevel}}\hfill{\footnotesize \path{lean/qtsl-aux/QtslAux/QtslAliases.lean}, line 58}\par\nopagebreak
\begin{leansrc}
theorem sound_QSatFromRegion_sublevel {σ : Type*} {sat : (σ → ℝ) → Prop}
    {P : MvPolynomial σ ℚ} {thr : ℚ} {xq : σ → ℚ}
    (hchild : ∀ x ∈ {x : σ → ℝ |
      MvPolynomial.eval x (MvPolynomial.map (Rat.castHom ℝ) P) ≤ (thr : ℝ)}, sat x)
    (heval : MvPolynomial.eval xq P ≤ thr) :
    sat (fun i => (xq i : ℝ))
\end{leansrc}
\par\medskip\noindent{\small \texttt{PPLC.}\texttt{\bfseries rootEvalFormula\_true\_iff}}\hfill{\footnotesize \path{lean/qtsl/Qtsl/Elaborators.lean}, line 137}\par\nopagebreak
\begin{leansrc}
theorem rootEvalFormula_true_iff (x0 : Nat → Rat) (V : Poly) (thr : Rat) :
    (rootEvalFormula x0 V thr).evalClosed = some true
      ↔ Poly.eval x0 V ≤ thr
\end{leansrc}
\par\medskip\noindent{\small \texttt{Qtsl.Guard.}\texttt{\bfseries GuardSat}}\hfill{\footnotesize \path{lean/qtsl-aux/QtslAux/QGuard.lean}, line 57}\par\nopagebreak
\begin{leansrc}
def GuardSat (init R : Set X) (sat : X → Prop) : Prop
\end{leansrc}
\par\medskip\noindent{\small \texttt{Qtsl.Guard.}\texttt{\bfseries sound\_QGuardIntro}}\hfill{\footnotesize \path{lean/qtsl-aux/QtslAux/QGuard.lean}, line 74}\par\nopagebreak
\begin{leansrc}
theorem sound_QGuardIntro {init R : Set X} {sat : X → Prop}
    (hchild : ∀ x ∈ R, sat x) : GuardSat init R sat
\end{leansrc}
\par\medskip\noindent{\small \texttt{Qtsl.Guard.}\texttt{\bfseries sound\_QGuardIntro\_sublevel}}\hfill{\footnotesize \path{lean/qtsl-aux/QtslAux/QGuard.lean}, line 171}\par\nopagebreak
\begin{leansrc}
theorem sound_QGuardIntro_sublevel {σ : Type*} {init : Set (σ → ℝ)}
    {sat : (σ → ℝ) → Prop} {P : MvPolynomial σ ℚ} {thr : ℚ}
    (hchild : ∀ x ∈ {x : σ → ℝ |
        MvPolynomial.eval x (MvPolynomial.map (Rat.castHom ℝ) P) ≤ (thr : ℝ)},
      sat x) :
    GuardSat init
      {x : σ → ℝ |
        MvPolynomial.eval x (MvPolynomial.map (Rat.castHom ℝ) P) ≤ (thr : ℝ)}
      sat
\end{leansrc}
\par\medskip\noindent{\small \texttt{Qtsl.Guard.}\texttt{\bfseries guard\_eval\_sublevel}}\hfill{\footnotesize \path{lean/qtsl-aux/QtslAux/QGuard.lean}, line 185}\par\nopagebreak
\begin{leansrc}
theorem guard_eval_sublevel {σ : Type*} {P : MvPolynomial σ ℚ} {thr : ℚ}
    {xq : σ → ℚ} (heval : MvPolynomial.eval xq P ≤ thr) :
    (fun i => (xq i : ℝ)) ∈ {x : σ → ℝ |
      MvPolynomial.eval x (MvPolynomial.map (Rat.castHom ℝ) P) ≤ (thr : ℝ)}
\end{leansrc}
\par\medskip\noindent{\small \texttt{Qtsl.Guard.}\texttt{\bfseries sound\_QGuardIntro\_prob}}\hfill{\footnotesize \path{lean/qtsl-aux/QtslAux/QGuard.lean}, line 193}\par\nopagebreak
\begin{leansrc}
theorem sound_QGuardIntro_prob [MeasurableSpace X] {init R : Set X}
    {μ : X → Measure (ℕ → X)} {p : ENNReal} {A : Set (ℕ → X)}
    (hchild : ∀ x ∈ R, p ≤ μ x A) :
    GuardSat init R fun x => p ≤ μ x A
\end{leansrc}
\par\medskip\noindent{\small \texttt{Qtsl.Guard.}\texttt{\bfseries exists\_blanket\_not\_guardSat}}\hfill{\footnotesize \path{lean/qtsl-aux/QtslAux/QGuard.lean}, line 137}\par\nopagebreak
\begin{leansrc}
theorem exists_blanket_not_guardSat :
    ∃ (init R : Set Bool) (sat : Bool → Prop),
      GuardSatBlanket init R sat ∧ ¬ GuardSat init R sat
\end{leansrc}
\par\medskip\noindent{\small \texttt{Qtsl.Guard.}\texttt{\bfseries exists\_not\_guardSat\_wrong\_region}}\hfill{\footnotesize \path{lean/qtsl-aux/QtslAux/QGuard.lean}, line 160}\par\nopagebreak
\begin{leansrc}
theorem exists_not_guardSat_wrong_region :
    ∃ (init R R' : Set Bool) (sat : Bool → Prop),
      (∀ x ∈ R, sat x) ∧ ¬ GuardSat init R' sat
\end{leansrc}
\par\medskip\noindent{\small \texttt{Qtsl.Guard.}\texttt{\bfseries sound\_QGuardElim}}\hfill{\footnotesize \path{lean/qtsl-aux/QtslAux/QGuard.lean}, line 81}\par\nopagebreak
\begin{leansrc}
theorem sound_QGuardElim {init R : Set X} {sat : X → Prop}
    (hg : GuardSat init R sat) (heval : init ⊆ R) : ∀ x₀ ∈ init, sat x₀
\end{leansrc}
\par\medskip\noindent{\small \texttt{Qtsl.Guard.}\texttt{\bfseries exists\_guardSat\_not\_unguarded}}\hfill{\footnotesize \path{lean/qtsl-aux/QtslAux/QGuard.lean}, line 150}\par\nopagebreak
\begin{leansrc}
theorem exists_guardSat_not_unguarded :
    ∃ (init R : Set Bool) (sat : Bool → Prop),
      GuardSat init R sat ∧ ¬ ∀ x₀ ∈ init, sat x₀
\end{leansrc}
\par\medskip\noindent{\small \texttt{Qtsl.Guard.}\texttt{\bfseries qsatFromRegion\_derived}}\hfill{\footnotesize \path{lean/qtsl-aux/QtslAux/QGuard.lean}, line 88}\par\nopagebreak
\begin{leansrc}
theorem qsatFromRegion_derived {init R : Set X} {sat : X → Prop}
    (hchild : ∀ x ∈ R, sat x) (heval : init ⊆ R) : ∀ x₀ ∈ init, sat x₀
\end{leansrc}
\par\medskip\noindent{\small \texttt{Qtsl.LayerD.}\texttt{\bfseries sound\_GSHS\_safety\_mp}}\hfill{\footnotesize \path{lean/qtsl-aux/QtslAux/LayerD/MPDischarge.lean}, line 511}\par\nopagebreak
\begin{leansrc}
theorem sound_GSHS_safety_mp [IsMarkovKernel M.seg] [IsMarkovKernel M.jump]
    {V : State → ℝ} {U : Set State} (hV : Measurable V) (hU : MeasurableSet U)
    (hmp : ∀ x : State, Nonempty (MPFlow M 0 V x))
    (hForce : ForcedNonexpansion M V)
    (hNonneg : NonnegV V) {C : ℝ} (hC0 : 0 ≤ C) (hVC : ∀ y, V y ≤ C)
    (hFloor : ∀ y ∈ U, 1 ≤ V y) (x : State) :
    Battery.pathMeasure (skel M) (SkelState.alive x)
        {ω | ∃ n, ω n ∈ skelBad U}
      ≤ ENNReal.ofReal (V x)
\end{leansrc}
\par\medskip\noindent{\small \texttt{Qtsl.LayerD.}\texttt{\bfseries sound\_GSHS\_safety\_battery\_mp}}\hfill{\footnotesize \path{lean/qtsl-aux/QtslAux/LayerD/MPDischarge.lean}, line 524}\par\nopagebreak
\begin{leansrc}
theorem sound_GSHS_safety_battery_mp
    [IsMarkovKernel M.seg] [IsMarkovKernel M.jump]
    {V : State → ℝ} {U : Set State} {p thr : ℝ}
    (hV : Measurable V) (hU : MeasurableSet U)
    (hmp : ∀ x : State, Nonempty (MPFlow M 0 V x))
    (hForce : ForcedNonexpansion M V)
    (hNonneg : NonnegV V) {C : ℝ} (hC0 : 0 ≤ C) (hVC : ∀ y, V y ≤ C)
    (hFloor : ∀ y ∈ U, 1 ≤ V y)
    (hΘ : p + thr ≤ 1) (x : State) (hthr : V x ≤ thr) :
    ENNReal.ofReal p ≤ Battery.pathMeasure (skel M) (SkelState.alive x)
      {ω | ∀ n, ω n ∉ skelBad U}
\end{leansrc}
\par\medskip\noindent{\small \texttt{Qtsl.Battery.}\texttt{\bfseries sound\_QSafetyBattery\_staged}}\hfill{\footnotesize \path{lean/qtsl-aux/QtslAux/QtslBattery.lean}, line 556}\par\nopagebreak
\begin{leansrc}
theorem sound_QSafetyBattery_staged
    {K : Kernel (Md × X) (Md × X)} [IsMarkovKernel K]
    {Dm Dbarm Um : Md → Set X} {p thr : ℚ} {sat : Md × X → Prop}
    (hDm : ∀ m, MeasurableSet (Dm m)) (hDbarm : ∀ m, MeasurableSet (Dbarm m))
    (hUm : ∀ m, MeasurableSet (Um m)) (hVmeas : ∀ m, Measurable (V m))
    (hF : ∀ m z, ¬ sat (m, z) → z ∈ Um m)
    (hΘ : p + thr ≤ 1)
    (hDsub : ∀ m, Dm m ⊆ Dbarm m)
    (hR3 : ∀ m z, z ∈ Dm m → z ∉ Um m → K (m, z) {y | y.2 ∈ Dbarm y.1} = 1)
    (hR4 : ∀ m z, z ∈ Dm m → z ∉ Um m →
      K (m, z) ({y | y.2 ∈ Dm y.1} ∪ {y | y.2 ∈ Um y.1}) = 1)
    (hint : ∀ m z, z ∈ Dm m → z ∉ Um m → Integrable (V m) ((flow m) z))
    (hB1 : ∀ m z, z ∈ Dm m → z ∉ Um m → ∫ ζ, V m ζ ∂(flow m) z ≤ V m z)
    (hB2 : ∀ j, ∀ z ∈ G j, V (dst j) (g j z) ≤ V (src j) z)
    (hB3 : ∀ m z, z ∈ Um m → z ∈ Dbarm m → 1 ≤ V m z)
    (hB4 : ∀ m z, z ∈ Dbarm m → 0 ≤ V m z)
    (hmix : ∀ m z, z ∈ Dm m → z ∉ Um m →
      ∃ (w₀ : ℝ≥0) (w : J → ℝ≥0), w₀ + ∑ j, w j = 1 ∧
        (∀ j, w j ≠ 0 → src j = m ∧ z ∈ G j) ∧
        K (m, z) = (w₀ : ℝ≥0∞) • ((flow m) z).map (Prod.mk m)
          + ∑ j, (w j : ℝ≥0∞) • Measure.dirac (dst j, g j z)) :
    ∀ m z, z ∈ Dm m → V m z ≤ (thr : ℝ) →
      ENNReal.ofReal (p : ℝ)
        ≤ pathMeasure K (m, z) {ω | ∀ n, sat (ω n)}
\end{leansrc}
\par\medskip\noindent{\small \texttt{SuFunctor.Site.}\texttt{\bfseries measurableSet\_confinedEvent\_cylinder}}\hfill{\footnotesize \path{lean/su-functor/SuFunctor/Site/FellAcceptance.lean}, line 413}\par\nopagebreak
\begin{leansrc}
theorem measurableSet_confinedEvent_cylinder {D : Set X}
    (hD : IsClosed D) :
    MeasurableSet
      (confinedEvent ((Set.univ : Set ↥W) ×ˢ D) : Set (FTraj W X))
\end{leansrc}
\par\medskip\noindent{\small \texttt{SuFunctor.Site.}\texttt{\bfseries satAtM\_safetySpecF\_iff}}\hfill{\footnotesize \path{lean/su-functor/SuFunctor/Site/FellBehaviorSpec.lean}, line 254}\par\nopagebreak
\begin{leansrc}
theorem satAtM_safetySpecF_iff (U : ι → Set T)
    (hU : ∀ i, IsOpen (U i)) {D : Set X} (hD : IsClosed D) (j : ι)
    (μ : MeasJoint (FTraj (U j) X)) :
    SatAtM μ (specOfBehaviorF U (safetyBehaviorF U D) j) ↔
      (μ.1 : Measure (FTraj (U j) X))
        (confinedEvent ((Set.univ : Set ↥(U j)) ×ˢ D))ᶜ = 0
\end{leansrc}
\par\medskip\noindent{\small \texttt{SuFunctor.Site.}\texttt{\bfseries pull\_indPenM}}\hfill{\footnotesize \path{lean/su-functor/SuFunctor/Site/FellBehaviorSpec.lean}, line 110}\par\nopagebreak
\begin{leansrc}
theorem pull_indPenM (h : K → L) (hm : Measurable h) {S : Set L}
    (hS : MeasurableSet S) :
    MeasPen.pull h (indPenM S) = indPenM (h ⁻¹' S)
\end{leansrc}
\par\medskip\noindent{\small \texttt{SuFunctor.Site.}\texttt{\bfseries graphCl\_mem\_confinedEvent\_iff}}\hfill{\footnotesize \path{lean/su-functor/SuFunctor/Site/FellBehaviorSpec.lean}, line 340}\par\nopagebreak
\begin{leansrc}
theorem graphCl_mem_confinedEvent_iff {D : Set X} (hD : IsClosed D)
    (f : Traj W X) :
    graphCl f ∈ confinedEvent ((Set.univ : Set ↥W) ×ˢ D) ↔
      ∀ w, f w ∈ D
\end{leansrc}
\par\medskip\noindent{\small \texttt{SuFunctor.Site.}\texttt{\bfseries mboxplusWF\_indPenM}}\hfill{\footnotesize \path{lean/su-functor/SuFunctor/Site/FellBehaviorSpec.lean}, line 280}\par\nopagebreak
\begin{leansrc}
theorem mboxplusWF_indPenM {W : Set T} {π₁ : Z → U'} {π₂ : Z → V'}
    (hp₁ : IsProperMap π₁) (hp₂ : IsProperMap π₂)
    {S₁ : Set (FTraj W U')} {S₂ : Set (FTraj W V')}
    (hS₁ : MeasurableSet S₁) (hS₂ : MeasurableSet S₂) :
    mboxplusWF π₁ π₂ (indPenM S₁) (indPenM S₂)
      = indPenM (fellMap (W := W) π₁ ⁻¹' S₁ ∩
          fellMap (W := W) π₂ ⁻¹' S₂)
\end{leansrc}
\par\medskip\noindent{\small \texttt{Qtsl.Battery.}\texttt{\bfseries sound\_QSafetyBattery}}\hfill{\footnotesize \path{lean/qtsl-aux/QtslAux/Battery.lean}, line 433}\par\nopagebreak
\begin{leansrc}
theorem sound_QSafetyBattery {sat : S → Prop}
    (hD : MeasurableSet D) (hDbar : MeasurableSet Dbar) (hU : MeasurableSet U)
    (hV : Measurable V)
    (hF : ∀ y, ¬ sat y → y ∈ U)
    (hΘ : p + thr ≤ 1)
    (hDsub : D ⊆ Dbar)
    (hR3 : ∀ y ∈ D, y ∉ U → K y Dbar = 1)
    (hR4 : ∀ y ∈ D, y ∉ U → K y (D ∪ U) = 1)
    (hB12int : ∀ y ∈ D, y ∉ U → Integrable V (K y))
    (hB12 : ∀ y ∈ D, y ∉ U → ∫ z, V z ∂K y ≤ V y)
    (hB3 : ∀ y ∈ U, y ∈ Dbar → 1 ≤ V y)
    (hB4 : ∀ y ∈ Dbar, 0 ≤ V y)
    (hx : x ∈ D) (hthr : V x ≤ thr) :
    ENNReal.ofReal p ≤ pathMeasure K x {ω | ∀ n, sat (ω n)}
\end{leansrc}
\par\medskip\noindent{\small \texttt{Qtsl.Battery.}\texttt{\bfseries pathMeasure}}\hfill{\footnotesize \path{lean/qtsl-aux/QtslAux/Battery.lean}, line 69}\par\nopagebreak
\begin{leansrc}
noncomputable def pathMeasure (K : Kernel S S) [IsMarkovKernel K] (x : S) :
    Measure (ℕ → S)
\end{leansrc}
\par\medskip\noindent{\small \texttt{Qtsl.LayerD.}\texttt{\bfseries sound\_GSHS\_safety\_battery}}\hfill{\footnotesize \path{lean/qtsl-aux/QtslAux/LayerD/Skeleton.lean}, line 253}\par\nopagebreak
\begin{leansrc}
theorem sound_GSHS_safety_battery [IsMarkovKernel M.seg] [IsMarkovKernel M.jump]
    {V : State → ℝ} {U : Set State} {p thr : ℝ}
    (hV : Measurable V) (hU : MeasurableSet U)
    (hSegD : SegmentDynkin M 0 V) (hForce : ForcedNonexpansion M V)
    (hNonneg : NonnegV V) {C : ℝ} (hC0 : 0 ≤ C) (hVC : ∀ y, V y ≤ C)
    (hFloor : ∀ y ∈ U, 1 ≤ V y)
    (hΘ : p + thr ≤ 1) (x : State) (hthr : V x ≤ thr) :
    ENNReal.ofReal p ≤ Battery.pathMeasure (skel M) (SkelState.alive x)
      {ω | ∀ n, ω n ∉ skelBad U}
\end{leansrc}
\par\medskip\noindent{\small \texttt{Qtsl.Battery.}\texttt{\bfseries sound\_QSafetyBattery\_failClosed}}\hfill{\footnotesize \path{lean/qtsl-aux/QtslAux/Battery.lean}, line 311}\par\nopagebreak
\begin{leansrc}
theorem sound_QSafetyBattery_failClosed
    (hD : MeasurableSet D) (hDbar : MeasurableSet Dbar) (hU : MeasurableSet U)
    (hV : Measurable V)
    (hΘ : p + thr ≤ 1)
    (hDsub : D ⊆ Dbar)
    (hR3 : ∀ y ∈ D, y ∉ U → K y Dbar = 1)
    (hR4 : ∀ y ∈ D, y ∉ U → K y (D ∪ U) = 1)
    (hB12int : ∀ y ∈ D, y ∉ U → Integrable V (K y))
    (hB12 : ∀ y ∈ D, y ∉ U → ∫ z, V z ∂K y ≤ V y)
    (hB3 : ∀ y ∈ U, y ∈ Dbar → 1 ≤ V y)
    (hB4 : ∀ y ∈ Dbar, 0 ≤ V y)
    (hx : x ∈ D) (hthr : V x ≤ thr) :
    ENNReal.ofReal p ≤ pathMeasure K x {ω | ∀ n, ω n ∈ D \ U}
\end{leansrc}
\par\medskip\noindent{\small \texttt{OptionalSamplingRC.}\texttt{\bfseries supermartingale\_of\_compensator\_nonpos}}\hfill{\footnotesize \path{lean/optional-sampling/OptionalSamplingRC/Basic.lean}, line 1150}\par\nopagebreak
\begin{leansrc}
theorem supermartingale_of_compensator_nonpos
    [IsProbabilityMeasure μ]
    (W : ℝ≥0 → Ω → ℝ) (g : ℝ → Ω → ℝ)
    (hmart : Martingale
      (fun (t : ℝ≥0) (ω : Ω) => W t ω - ∫ s in (0:ℝ)..(t:ℝ), g s ω) ℱ μ)
    (hGadp : Adapted ℱ
      (fun (t : ℝ≥0) (ω : Ω) => ∫ s in (0:ℝ)..(t:ℝ), g s ω))
    (hWint : ∀ t : ℝ≥0, Integrable (W t) μ)
    (hgint : ∀ ω, ∀ t : ℝ, IntervalIntegrable (fun s => g s ω) volume 0 t)
    (hg0 : ∀ ω, ∀ s : ℝ, 0 ≤ s → g s ω ≤ 0) :
    Supermartingale W ℱ μ
\end{leansrc}
\par\medskip\noindent{\small \texttt{Qtsl.Battery.}\texttt{\bfseries mixture\_drift}}\hfill{\footnotesize \path{lean/qtsl-aux/QtslAux/QtslBattery.lean}, line 488}\par\nopagebreak
\begin{leansrc}
lemma mixture_drift {m : Md} {z : X} {κ : Measure (Md × X)}
    {w₀ : ℝ≥0} {w : J → ℝ≥0}
    (hVmeas : ∀ m', Measurable (V m'))
    (hκ : κ = (w₀ : ℝ≥0∞) • ((flow m) z).map (Prod.mk m)
      + ∑ j, (w j : ℝ≥0∞) • Measure.dirac (dst j, g j z))
    (hsum : w₀ + ∑ j, w j = 1)
    (hguard : ∀ j, w j ≠ 0 → V (dst j) (g j z) ≤ V m z)
    (hflowint : Integrable (fun ζ => V m ζ) ((flow m) z))
    (hflowdrift : ∫ ζ, V m ζ ∂(flow m) z ≤ V m z) :
    ∫ y, V y.1 y.2 ∂κ ≤ V m z
\end{leansrc}
\par\medskip\noindent{\small \texttt{Qtsl.Battery.}\texttt{\bfseries stronglyAdapted\_stopped}}\hfill{\footnotesize \path{lean/qtsl-aux/QtslAux/Battery.lean}, line 187}\par\nopagebreak
\begin{leansrc}
lemma stronglyAdapted_stopped {f : S → ℝ} {A : Set S} (hA : MeasurableSet A)
    (hf : Measurable f) :
    StronglyAdapted (Filtration.piLE (X := fun _ : ℕ => S)) (stopped f A)
\end{leansrc}
\par\medskip\noindent{\small \texttt{Qtsl.Battery.}\texttt{\bfseries one\_le\_stopped}}\hfill{\footnotesize \path{lean/qtsl-aux/QtslAux/Battery.lean}, line 133}\par\nopagebreak
\begin{leansrc}
lemma one_le_stopped {f : S → ℝ} {A : Set S} (hf1 : ∀ y ∈ A, 1 ≤ f y) :
    ∀ (n : ℕ) (ω : ℕ → S), (∃ k ≤ n, ω k ∈ A) → 1 ≤ stopped f A n ω
\end{leansrc}
\par\medskip\noindent{\small \texttt{Qtsl.Battery.}\texttt{\bfseries stopped\_congr}}\hfill{\footnotesize \path{lean/qtsl-aux/QtslAux/Battery.lean}, line 154}\par\nopagebreak
\begin{leansrc}
lemma stopped_congr {f : S → ℝ} {A : Set S} : ∀ {n : ℕ} {ω ω' : ℕ → S},
    (∀ k ≤ n, ω k = ω' k) → stopped f A n ω = stopped f A n ω'
\end{leansrc}
\par\medskip\noindent{\small \texttt{Qtsl.Battery.}\texttt{\bfseries supermartingale\_stopped}}\hfill{\footnotesize \path{lean/qtsl-aux/QtslAux/Battery.lean}, line 263}\par\nopagebreak
\begin{leansrc}
theorem supermartingale_stopped {f : S → ℝ} {A : Set S} {C : ℝ}
    (hA : MeasurableSet A) (hf : Measurable f) (hf0 : ∀ y, 0 ≤ f y) (hfC : ∀ y, f y
      ≤ C)
    (hdrift : ∀ y, y ∉ A → ∫ z, f z ∂K y ≤ f y) (x : S) :
    Supermartingale (stopped f A) (Filtration.piLE (X := fun _ : ℕ => S))
      (pathMeasure K x)
\end{leansrc}
\par\medskip\noindent{\small \texttt{Qtsl.Battery.}\texttt{\bfseries measure\_reach\_le}}\hfill{\footnotesize \path{lean/qtsl-aux/QtslAux/Battery.lean}, line 279}\par\nopagebreak
\begin{leansrc}
theorem measure_reach_le {f : S → ℝ} {A : Set S} {C : ℝ}
    (hA : MeasurableSet A) (hf : Measurable f) (hf0 : ∀ y, 0 ≤ f y) (hfC : ∀ y, f y
      ≤ C)
    (hf1 : ∀ y ∈ A, 1 ≤ f y) (hdrift : ∀ y, y ∉ A → ∫ z, f z ∂K y ≤ f y) (x : S) :
    pathMeasure K x {ω | ∃ n, ω n ∈ A} ≤ ENNReal.ofReal (f x)
\end{leansrc}
\par\medskip\noindent{\small \texttt{Qtsl.Battery.}\texttt{\bfseries mixture\_integrable}}\hfill{\footnotesize \path{lean/qtsl-aux/QtslAux/QtslBattery.lean}, line 468}\par\nopagebreak
\begin{leansrc}
lemma mixture_integrable {m : Md} {z : X} {κ : Measure (Md × X)}
    {w₀ : ℝ≥0} {w : J → ℝ≥0}
    (hVmeas : ∀ m', Measurable (V m'))
    (hκ : κ = (w₀ : ℝ≥0∞) • ((flow m) z).map (Prod.mk m)
      + ∑ j, (w j : ℝ≥0∞) • Measure.dirac (dst j, g j z))
    (hflowint : Integrable (fun ζ => V m ζ) ((flow m) z)) :
    Integrable (fun y : Md × X => V y.1 y.2) κ
\end{leansrc}
\par\medskip\noindent{\small \texttt{Qtsl.LayerD.}\texttt{\bfseries sound\_GSHS\_deadline\_mp}}\hfill{\footnotesize \path{lean/qtsl-aux/QtslAux/LayerD/MPDischarge.lean}, line 539}\par\nopagebreak
\begin{leansrc}
theorem sound_GSHS_deadline_mp [IsMarkovKernel M.seg] [IsMarkovKernel M.jump]
    {V : State → ℝ} {θ ε : ℝ} {T : ℝ≥0}
    (hV : Measurable V) (hθ : 0 ≤ θ)
    (hmp : ∀ x : State, Nonempty (MPFlow M θ V x))
    (hForce : ForcedNonexpansion M V)
    (hNonneg : NonnegV V) (hε : 0 < ε) (hFloor : ∀ y, ε ≤ V y) (x : State) :
    Battery.pathMeasure (tskel M) (TSkelState.alive (0, x))
        {ω | ∃ n, T < tclock (ω n)}
      ≤ ENNReal.ofReal (Real.exp (-(θ * T)) * V x / ε)
\end{leansrc}
\par\medskip\noindent{\small \texttt{Qtsl.LayerD.}\texttt{\bfseries sound\_GSHS\_deadline\_battery\_mp}}\hfill{\footnotesize \path{lean/qtsl-aux/QtslAux/LayerD/MPDischarge.lean}, line 552}\par\nopagebreak
\begin{leansrc}
theorem sound_GSHS_deadline_battery_mp
    [IsMarkovKernel M.seg] [IsMarkovKernel M.jump]
    {V : State → ℝ} {θ ε p thr : ℝ} {T : ℝ≥0}
    (hV : Measurable V) (hθ : 0 ≤ θ)
    (hmp : ∀ x : State, Nonempty (MPFlow M θ V x))
    (hForce : ForcedNonexpansion M V)
    (hNonneg : NonnegV V) (hε : 0 < ε) (hFloor : ∀ y, ε ≤ V y)
    (hΘ : p + Real.exp (-(θ * T)) * thr / ε ≤ 1) (x : State)
    (hthr : V x ≤ thr) :
    ENNReal.ofReal p ≤ Battery.pathMeasure (tskel M) (TSkelState.alive (0, x))
      {ω | ∀ n, tclock (ω n) ≤ T}
\end{leansrc}
\par\medskip\noindent{\small \texttt{SuFunctor.Site.}\texttt{\bfseries satAtM\_acceptanceSpecF\_ae}}\hfill{\footnotesize \path{lean/su-functor/SuFunctor/Site/FellBehaviorSpec.lean}, line 180}\par\nopagebreak
\begin{leansrc}
theorem satAtM_acceptanceSpecF_ae (U : ι → Set T)
    (E : ∀ i, Set (FTraj (U i) X)) (j : ι)
    {μ : MeasJoint (FTraj (U j) X)}
    (h : SatAtM μ (acceptanceSpecF U E j)) :
    (μ.1 : Measure (FTraj (U j) X)) (E j)ᶜ = 0
\end{leansrc}
\par\medskip\noindent{\small \texttt{SuFunctor.Site.}\texttt{\bfseries hitEvent\_iUnion}}\hfill{\footnotesize \path{lean/su-functor/SuFunctor/Site/FellAcceptance.lean}, line 123}\par\nopagebreak
\begin{leansrc}
theorem hitEvent_iUnion {ι : Sort*} (A : ι → Set α) :
    hitEvent (⋃ i, A i) = ⋃ i, hitEvent (A i)
\end{leansrc}
\par\medskip\noindent{\small \texttt{SuFunctor.Site.}\texttt{\bfseries measurableSet\_hitEvent\_univ\_prod\_of\_isCompact}}\hfill{\footnotesize \path{lean/su-functor/SuFunctor/Site/FellAcceptance.lean}, line 397}\par\nopagebreak
\begin{leansrc}
theorem measurableSet_hitEvent_univ_prod_of_isCompact
    [SigmaCompactSpace ↥W] {R : Set X} (hR : IsCompact R) :
    MeasurableSet
      (hitEvent ((Set.univ : Set ↥W) ×ˢ R) : Set (FTraj W X))
\end{leansrc}
\par\medskip\noindent{\small \texttt{Qtsl.Battery.}\texttt{\bfseries sound\_QDeadline}}\hfill{\footnotesize \path{lean/qtsl-aux/QtslAux/Battery.lean}, line 749}\par\nopagebreak
\begin{leansrc}
theorem sound_QDeadline {T : ℕ} {M : ℝ} {sat : S → Prop}
    (hD : MeasurableSet D) (hDbar : MeasurableSet Dbar) (hB : MeasurableSet B)
    (hV : Measurable V)
    (hF : ∀ y ∈ B, sat y)
    (hβ : 0 ≤ β)
    (hΘF : β ^ T * thr ≤ 1 - p)
    (hDsub : D ⊆ Dbar)
    (hR3 : ∀ y ∈ D, y ∉ B → K y Dbar = 1)
    (hR4 : ∀ y ∈ D, y ∉ B → K y (D ∪ B) = 1)
    (hint : ∀ y ∈ D, y ∉ B → Integrable V (K y))
    (hdecay : ∀ y ∈ D, y ∉ B → ∫ z, V z ∂K y ≤ β * V y)
    (hfloor : ∀ y ∈ Dbar, y ∉ B → 1 ≤ V y)
    (hnonneg : ∀ y ∈ Dbar, 0 ≤ V y)
    (hM : ∀ y ∈ Dbar, V y ≤ M)
    (hx : x ∈ D) (hthr : V x ≤ thr) :
    ENNReal.ofReal p ≤ pathMeasure K x {ω | ∃ k ≤ T, sat (ω k)}
\end{leansrc}
\par\medskip\noindent{\small \texttt{Qtsl.Battery.}\texttt{\bfseries sound\_QDeadlineBattery}}\hfill{\footnotesize \path{lean/qtsl-aux/QtslAux/QtslBattery.lean}, line 413}\par\nopagebreak
\begin{leansrc}
theorem sound_QDeadlineBattery {n : ℕ} {sat : S → Prop}
    (hD : MeasurableSet D) (hDbar : MeasurableSet Dbar) (hB : MeasurableSet B)
    (hV : Measurable V)
    (hF : ∀ y ∈ B, sat y)
    (hβ : 0 ≤ β) (hβθ : β ≤ Real.exp (-θ)) (hθ : 0 ≤ θ)
    (hq0 : 0 ≤ q)
    (hq : q ≤ ∑ k ∈ range (n + 1), (θ * T) ^ k / (Nat.factorial k : ℝ))
    (hp : p ≤ 1)
    (hDsub : D ⊆ Dbar)
    (hR3 : ∀ y ∈ D, y ∉ B → K y Dbar = 1)
    (hR4 : ∀ y ∈ D, y ∉ B → K y (D ∪ B) = 1)
    (hint : ∀ y ∈ D, y ∉ B → Integrable V (K y))
    (hdecay : ∀ y ∈ D, y ∉ B → ∫ z, V z ∂K y ≤ β * V y)
    (hfloor : ∀ y ∈ Dbar, y ∉ B → 1 ≤ V y)
    (hnonneg : ∀ y ∈ Dbar, 0 ≤ V y)
    (hM : ∀ y ∈ Dbar, V y ≤ M)
    (hx : x ∈ D) (hthr : V x ≤ (1 - p) * q) :
    ENNReal.ofReal p ≤ pathMeasure K x {ω | ∃ k ≤ T, sat (ω k)}
\end{leansrc}
\par\medskip\noindent{\small \texttt{Qtsl.Battery.}\texttt{\bfseries sound\_QDeadline\_hit}}\hfill{\footnotesize \path{lean/qtsl-aux/QtslAux/Battery.lean}, line 568}\par\nopagebreak
\begin{leansrc}
theorem sound_QDeadline_hit {T : ℕ} {M : ℝ}
    (hD : MeasurableSet D) (hDbar : MeasurableSet Dbar) (hB : MeasurableSet B)
    (hV : Measurable V)
    (hβ : 0 ≤ β)
    (hΘF : β ^ T * thr ≤ 1 - p)
    (hDsub : D ⊆ Dbar)
    (hR3 : ∀ y ∈ D, y ∉ B → K y Dbar = 1)
    (hR4 : ∀ y ∈ D, y ∉ B → K y (D ∪ B) = 1)
    (hint : ∀ y ∈ D, y ∉ B → Integrable V (K y))
    (hdecay : ∀ y ∈ D, y ∉ B → ∫ z, V z ∂K y ≤ β * V y)
    (hfloor : ∀ y ∈ Dbar, y ∉ B → 1 ≤ V y)
    (hnonneg : ∀ y ∈ Dbar, 0 ≤ V y)
    (hM : ∀ y ∈ Dbar, V y ≤ M)
    (hx : x ∈ D) (hthr : V x ≤ thr) :
    ENNReal.ofReal p ≤ pathMeasure K x {ω | ∃ k ≤ T, ω k ∈ B}
\end{leansrc}
\par\medskip\noindent{\small \texttt{MeasureTheory.}\texttt{\bfseries supermartingale\_inv\_pow\_mul}}\hfill{\footnotesize \path{lean/qtsl-aux/QtslAux/Deadline.lean}, line 34}\par\nopagebreak
\begin{leansrc}
theorem supermartingale_inv_pow_mul [IsFiniteMeasure μ] {f : ℕ → Ω → ℝ} {β : ℝ}
    (hadp : StronglyAdapted 𝒢 f) (hint : ∀ n, Integrable (f n) μ) (hβ : 0 < β)
    (hdrift : ∀ n, μ[f (n + 1) | 𝒢 n] ≤ᵐ[μ] fun ω => β * f n ω) :
    Supermartingale (fun n ω => β⁻¹ ^ n * f n ω) 𝒢 μ
\end{leansrc}
\par\medskip\noindent{\small \texttt{Qtsl.Deadline.}\texttt{\bfseries expMoment\_hitting}}\hfill{\footnotesize \path{lean/qtsl-aux/QtslAux/Deadline.lean}, line 128}\par\nopagebreak
\begin{leansrc}
theorem expMoment_hitting (hβ : 0 < β)
    (hint : ∀ n, Integrable (V n) μ) (hnonneg : 0 ≤ V)
    (hmB : ∀ n, MeasurableSet[𝒢 n] {ω | X n ω ∈ B})
    (hdrift : ∀ n, ∀ᵐ ω ∂μ, X n ω ∉ B → (μ[V (n + 1) | 𝒢 n]) ω ≤ β * V n ω)
    (hbarrier : ∀ n ω, X n ω ∈ B → a ≤ V n ω) :
    ENNReal.ofReal a *
        ∑' k, ENNReal.ofReal (β⁻¹ ^ k) * μ {ω | X k ω ∈ B ∧ ∀ j < k, X j ω ∉ B}
      ≤ ENNReal.ofReal (∫ ω, V 0 ω ∂μ)
\end{leansrc}
\par\medskip\noindent{\small \texttt{Qtsl.Battery.}\texttt{\bfseries ae\_no\_escape}}\hfill{\footnotesize \path{lean/qtsl-aux/QtslAux/Battery.lean}, line 502}\par\nopagebreak
\begin{leansrc}
private lemma ae_no_escape (hD : MeasurableSet D) (hB : MeasurableSet B)
    (hR4 : ∀ y ∈ D, y ∉ B → K y (D ∪ B) = 1) (hx : x ∈ D) :
    pathMeasure K x (⋃ n, {ω : ℕ → S | (∀ k ≤ n, ω k ∉ B) ∧ ω n ∉ D}) = 0
\end{leansrc}
\par\medskip\noindent{\small \texttt{Qtsl.Battery.}\texttt{\bfseries ae\_coord\_zero\_mem}}\hfill{\footnotesize \path{lean/qtsl-aux/QtslAux/Battery.lean}, line 484}\par\nopagebreak
\begin{leansrc}
private lemma ae_coord_zero_mem {E : Set S} (hE : MeasurableSet E) (hxE : x ∈ E) :
    ∀ᵐ ω ∂pathMeasure K x, ω 0 ∈ E
\end{leansrc}
\par\medskip\noindent{\small \texttt{Qtsl.Battery.}\texttt{\bfseries compl\_firstExit\_subset}}\hfill{\footnotesize \path{lean/qtsl-aux/QtslAux/QtslBattery.lean}, line 70}\par\nopagebreak
\begin{leansrc}
theorem compl_firstExit_subset {S : Type*} (Stop Good : Set S) (T : ℕ) :
    {ω : ℕ → S | ∃ n ≤ T, ω n ∈ Good ∧ ∀ j < n, ω j ∉ Stop}ᶜ ⊆
      {ω : ℕ → S | ∃ n, ω n ∈ Stop \ Good ∧ ∀ j < n, ω j ∉ Stop} ∪
        {ω : ℕ → S | ∀ k ≤ T, ω k ∉ Stop}
\end{leansrc}
\par\medskip\noindent{\small \texttt{Qtsl.Battery.}\texttt{\bfseries sound\_QUnionBound}}\hfill{\footnotesize \path{lean/qtsl-aux/QtslAux/QtslBattery.lean}, line 50}\par\nopagebreak
\begin{leansrc}
theorem sound_QUnionBound {Ω : Type*} [MeasurableSpace Ω] (μ : Measure Ω)
    {E₁ E₂ F : Set Ω} {a b : ℝ≥0∞}
    (h₁ : μ E₁ ≤ a) (h₂ : μ E₂ ≤ b) (hsub : F ⊆ E₁ ∪ E₂) :
    μ F ≤ a + b
\end{leansrc}
\par\medskip\noindent{\small \texttt{Qtsl.Battery.}\texttt{\bfseries sound\_QUnionBound\_ofReal}}\hfill{\footnotesize \path{lean/qtsl-aux/QtslAux/QtslBattery.lean}, line 58}\par\nopagebreak
\begin{leansrc}
theorem sound_QUnionBound_ofReal {Ω : Type*} [MeasurableSpace Ω] (μ : Measure Ω)
    {E₁ E₂ F : Set Ω} {a b : ℝ} (ha : 0 ≤ a) (hb : 0 ≤ b)
    (h₁ : μ E₁ ≤ ENNReal.ofReal a) (h₂ : μ E₂ ≤ ENNReal.ofReal b)
    (hsub : F ⊆ E₁ ∪ E₂) :
    μ F ≤ ENNReal.ofReal (a + b)
\end{leansrc}
\par\medskip\noindent{\small \texttt{Qtsl.Battery.}\texttt{\bfseries sound\_QUntilBattery\_firstHit}}\hfill{\footnotesize \path{lean/qtsl-aux/QtslAux/QtslBattery.lean}, line 206}\par\nopagebreak
\begin{leansrc}
theorem sound_QUntilBattery_firstHit
    (hD : MeasurableSet D) (hDbar : MeasurableSet Dbar)
    (hA : MeasurableSet A) (hR : MeasurableSet R)
    (hVa : Measurable Va) (hVd : Measurable Vd)
    (hβ : 0 ≤ β) (hc : β ^ T ≤ c)
    (hDsub : D ⊆ Dbar)
    (hR3 : ∀ y ∈ D, y ∉ A ∪ R → K y Dbar = 1)
    (hR4 : ∀ y ∈ D, y ∉ A ∪ R → K y (D ∪ (A ∪ R)) = 1)
    (hVaint : ∀ y ∈ D, y ∉ A ∪ R → Integrable Va (K y))
    (hVadrift : ∀ y ∈ D, y ∉ A ∪ R → ∫ z, Va z ∂K y ≤ Va y)
    (hVafloor : ∀ y ∈ A, y ∈ Dbar → 1 ≤ Va y)
    (hVanonneg : ∀ y ∈ Dbar, 0 ≤ Va y)
    (hVdint : ∀ y ∈ D, y ∉ A ∪ R → Integrable Vd (K y))
    (hVddecay : ∀ y ∈ D, y ∉ A ∪ R → ∫ z, Vd z ∂K y ≤ β * Vd y)
    (hVdfloor : ∀ y ∈ Dbar, y ∉ A ∪ R → 1 ≤ Vd y)
    (hVdnonneg : ∀ y ∈ Dbar, 0 ≤ Vd y)
    (hVdM : ∀ y ∈ Dbar, Vd y ≤ M)
    (hx : x ∈ D) (hthr : Va x + c * Vd x ≤ 1 - p) :
    ENNReal.ofReal p ≤ pathMeasure K x
      {ω | ∃ n ≤ T, ω n ∈ R ∧ ∀ j < n, ω j ∉ (A ∪ R) ∪ Dᶜ}
\end{leansrc}
\par\medskip\noindent{\small \texttt{Qtsl.LayerD.}\texttt{\bfseries sound\_GSHS\_until\_mp}}\hfill{\footnotesize \path{lean/qtsl-aux/QtslAux/LayerD/MPDischarge.lean}, line 584}\par\nopagebreak
\begin{leansrc}
theorem sound_GSHS_until_mp [IsMarkovKernel M.seg] [IsMarkovKernel M.jump]
    {Va Vd : State → ℝ} {Us : Set State} {θ ε : ℝ} {T : ℝ≥0}
    (hVa : Measurable Va) (hUs : MeasurableSet Us)
    (hmpa : ∀ x : State, Nonempty (MPFlow M 0 Va x))
    (hForcea : ForcedNonexpansion M Va)
    (hNonnega : NonnegV Va) {C : ℝ} (hC0 : 0 ≤ C) (hVC : ∀ y, Va y ≤ C)
    (hFloora : ∀ y ∈ Us, 1 ≤ Va y)
    (hVd : Measurable Vd) (hθ : 0 ≤ θ)
    (hmpd : ∀ x : State, Nonempty (MPFlow M θ Vd x))
    (hForced : ForcedNonexpansion M Vd)
    (hNonnegd : NonnegV Vd) (hε : 0 < ε) (hFloord : ∀ y, ε ≤ Vd y)
    (x : State) :
    Battery.pathMeasure (tskel M) (TSkelState.alive (0, x))
        {ω | ∃ n, ω n ∈ tbadIn Us ∨ T < tclock (ω n)}
      ≤ ENNReal.ofReal (Va x)
        + ENNReal.ofReal (Real.exp (-(θ * T)) * Vd x / ε)
\end{leansrc}
\par\medskip\noindent{\small \texttt{Qtsl.LayerD.}\texttt{\bfseries sound\_GSHS\_until\_battery\_mp}}\hfill{\footnotesize \path{lean/qtsl-aux/QtslAux/LayerD/MPDischarge.lean}, line 644}\par\nopagebreak
\begin{leansrc}
theorem sound_GSHS_until_battery_mp
    [IsMarkovKernel M.seg] [IsMarkovKernel M.jump]
    {Va Vd : State → ℝ} {Us : Set State} {θ ε p thra thrd : ℝ} {T : ℝ≥0}
    (hVa : Measurable Va) (hUs : MeasurableSet Us)
    (hmpa : ∀ x : State, Nonempty (MPFlow M 0 Va x))
    (hForcea : ForcedNonexpansion M Va)
    (hNonnega : NonnegV Va) {C : ℝ} (hC0 : 0 ≤ C) (hVC : ∀ y, Va y ≤ C)
    (hFloora : ∀ y ∈ Us, 1 ≤ Va y)
    (hVd : Measurable Vd) (hθ : 0 ≤ θ)
    (hmpd : ∀ x : State, Nonempty (MPFlow M θ Vd x))
    (hForced : ForcedNonexpansion M Vd)
    (hNonnegd : NonnegV Vd) (hε : 0 < ε) (hFloord : ∀ y, ε ≤ Vd y)
    (hΘ : p + (thra + Real.exp (-(θ * T)) * thrd / ε) ≤ 1)
    (x : State) (hthra : Va x ≤ thra) (hthrd : Vd x ≤ thrd) :
    ENNReal.ofReal p ≤ Battery.pathMeasure (tskel M) (TSkelState.alive (0, x))
      {ω | ∀ n, ω n ∉ tbadIn Us ∧ tclock (ω n) ≤ T}
\end{leansrc}
\par\medskip\noindent{\small \texttt{SuFunctor.Site.}\texttt{\bfseries untilEvent}}\hfill{\footnotesize \path{lean/su-functor/SuFunctor/Site/FellAcceptance.lean}, line 343}\par\nopagebreak
\begin{leansrc}
def untilEvent (Reach Avoid : ℕ → Set α) : Set (FellCloseds α)
\end{leansrc}
\par\medskip\noindent{\small \texttt{SuFunctor.Site.}\texttt{\bfseries measurableSet\_untilEvent}}\hfill{\footnotesize \path{lean/su-functor/SuFunctor/Site/FellAcceptance.lean}, line 353}\par\nopagebreak
\begin{leansrc}
theorem measurableSet_untilEvent {Reach Avoid : ℕ → Set α}
    (h₁ : ∀ n, MeasurableSet (hitEvent (Reach n)))
    (h₂ : ∀ n, MeasurableSet (missEvent (Avoid n))) :
    MeasurableSet (untilEvent Reach Avoid)
\end{leansrc}
\par\medskip\noindent{\small \texttt{Qtsl.Battery.}\texttt{\bfseries sound\_QUntilBattery}}\hfill{\footnotesize \path{lean/qtsl-aux/QtslAux/QtslBattery.lean}, line 371}\par\nopagebreak
\begin{leansrc}
theorem sound_QUntilBattery {sat_s sat : S → Prop}
    (hD : MeasurableSet D) (hDbar : MeasurableSet Dbar)
    (hA : MeasurableSet A) (hR : MeasurableSet R)
    (hVa : Measurable Va) (hVd : Measurable Vd)
    (hFs : ∀ y, ¬ sat_s y → y ∈ A) (hFt : ∀ y ∈ R, sat y)
    (hβ : 0 ≤ β) (hc : β ^ T ≤ c)
    (hDsub : D ⊆ Dbar)
    (hR3 : ∀ y ∈ D, y ∉ A ∪ R → K y Dbar = 1)
    (hR4 : ∀ y ∈ D, y ∉ A ∪ R → K y (D ∪ (A ∪ R)) = 1)
    (hVaint : ∀ y ∈ D, y ∉ A ∪ R → Integrable Va (K y))
    (hVadrift : ∀ y ∈ D, y ∉ A ∪ R → ∫ z, Va z ∂K y ≤ Va y)
    (hVafloor : ∀ y ∈ A, y ∈ Dbar → 1 ≤ Va y)
    (hVanonneg : ∀ y ∈ Dbar, 0 ≤ Va y)
    (hVdint : ∀ y ∈ D, y ∉ A ∪ R → Integrable Vd (K y))
    (hVddecay : ∀ y ∈ D, y ∉ A ∪ R → ∫ z, Vd z ∂K y ≤ β * Vd y)
    (hVdfloor : ∀ y ∈ Dbar, y ∉ A ∪ R → 1 ≤ Vd y)
    (hVdnonneg : ∀ y ∈ Dbar, 0 ≤ Vd y)
    (hVdM : ∀ y ∈ Dbar, Vd y ≤ M)
    (hx : x ∈ D) (hthr : Va x + c * Vd x ≤ 1 - p) :
    ENNReal.ofReal p ≤ pathMeasure K x
      {ω | ∃ n ≤ T, sat (ω n) ∧ ∀ j < n, sat_s (ω j)}
\end{leansrc}
\par\medskip\noindent{\small \texttt{Qtsl.Battery.}\texttt{\bfseries measure\_firstBad\_le}}\hfill{\footnotesize \path{lean/qtsl-aux/QtslAux/QtslBattery.lean}, line 110}\par\nopagebreak
\begin{leansrc}
theorem measure_firstBad_le {K : Kernel S S} [IsMarkovKernel K]
    {f : S → ℝ} {Stop Bad : Set S} {C : ℝ}
    (hStop : MeasurableSet Stop) (hf : Measurable f)
    (hf0 : ∀ y, 0 ≤ f y) (hfC : ∀ y, f y ≤ C)
    (hf1 : ∀ y ∈ Bad, 1 ≤ f y)
    (hdrift : ∀ y, y ∉ Stop → ∫ z, f z ∂K y ≤ f y) (x : S) :
    pathMeasure K x {ω | ∃ n, ω n ∈ Bad ∧ ∀ j < n, ω j ∉ Stop}
      ≤ ENNReal.ofReal (f x)
\end{leansrc}
\par\medskip\noindent{\small \texttt{Qtsl.Battery.}\texttt{\bfseries measure\_missTarget\_le}}\hfill{\footnotesize \path{lean/qtsl-aux/QtslAux/QtslBattery.lean}, line 143}\par\nopagebreak
\begin{leansrc}
theorem measure_missTarget_le {K : Kernel S S} [IsMarkovKernel K]
    {D Dbar B : Set S} {V : S → ℝ} {β : ℝ} {x : S} {T : ℕ} {M : ℝ}
    (hD : MeasurableSet D) (hDbar : MeasurableSet Dbar) (hB : MeasurableSet B)
    (hV : Measurable V) (hβ : 0 ≤ β)
    (hDsub : D ⊆ Dbar)
    (hR3 : ∀ y ∈ D, y ∉ B → K y Dbar = 1)
    (hR4 : ∀ y ∈ D, y ∉ B → K y (D ∪ B) = 1)
    (hint : ∀ y ∈ D, y ∉ B → Integrable V (K y))
    (hdecay : ∀ y ∈ D, y ∉ B → ∫ z, V z ∂K y ≤ β * V y)
    (hfloor : ∀ y ∈ Dbar, y ∉ B → 1 ≤ V y)
    (hnonneg : ∀ y ∈ Dbar, 0 ≤ V y)
    (hM : ∀ y ∈ Dbar, V y ≤ M)
    (hx : x ∈ D) :
    pathMeasure K x {ω | ∀ k ≤ T, ω k ∉ B} ≤ ENNReal.ofReal (β ^ T * V x)
\end{leansrc}
\par\medskip\noindent{\small \texttt{MeasureTheory.Supermartingale.}\texttt{\bfseries ville\_ineq}}\hfill{\footnotesize \path{lean/qtsl-aux/QtslAux/Ville.lean}, line 133}\par\nopagebreak
\begin{leansrc}
theorem Supermartingale.ville_ineq [IsFiniteMeasure μ] (hsup : Supermartingale f 𝒢
      μ)
    (hnonneg : 0 ≤ f) (ε : ℝ≥0) :
    ε * μ {ω | ∃ k, (ε : ℝ) ≤ f k ω} ≤ ENNReal.ofReal (∫ ω, f 0 ω ∂μ)
\end{leansrc}
\par\medskip\noindent{\small \texttt{Qtsl.QStreett.}\texttt{\bfseries absorptionBound\_of\_kernel}}\hfill{\footnotesize \path{lean/qtsl-aux/QtslAux/QStreett.lean}, line 804}\par\nopagebreak
\begin{leansrc}
theorem absorptionBound_of_kernel {A J Rq : Set S} {V : S → ℝ} {C γ : ℝ} {x : S}
    (hRq : MeasurableSet Rq) (hV : Measurable V)
    (hV0 : ∀ y, 0 ≤ V y) (hVC : ∀ y, V y ≤ C)
    (hVfloor : ∀ y ∈ Rq, 1 ≤ V y)
    (hVdrift : ∀ y, y ∉ Rq → ∫ z, V z ∂K y ≤ V y)
    (hVJ : ∀ y ∈ J, V y ≤ 1 - γ) :
    AbsorptionBound K x A J Rq γ
\end{leansrc}
\par\medskip\noindent{\small \texttt{Qtsl.LayerD.}\texttt{\bfseries sound\_GSHS\_streett\_bundle\_mp}}\hfill{\footnotesize \path{lean/qtsl-aux/QtslAux/LayerD/MPDischarge.lean}, line 612}\par\nopagebreak
\begin{leansrc}
theorem sound_GSHS_streett_bundle_mp
    {ι : Type*} [Countable ι]
    {D Dbar U : Set State} {V₀ : State → ℝ} {Rq G : ι → Set State}
    {W : ι → State → ℝ} {p thr C₀ : ℝ} {ε C : ι → ℝ} {x : State}
    [IsMarkovKernel M.seg] [IsMarkovKernel M.jump]
    (hD : MeasurableSet D) (hDbar : MeasurableSet Dbar) (hU : MeasurableSet U)
    (hV₀ : Measurable V₀)
    (hRq : ∀ i, MeasurableSet (Rq i)) (hG : ∀ i, MeasurableSet (G i))
    (hW : ∀ i, Measurable (W i))
    (hΘ : p + thr ≤ 1) (hDsub : D ⊆ Dbar)
    (hLive : NoVanish M (D \ U)) (hHalt : HaltsCaught M (D \ U) (U ∩ Dbar))
    (hReset : ResetsInto M (D \ U) ((D ∪ U) ∩ Dbar))
    (hmp₀ : ∀ x : State, Nonempty (MPFlow M 0 V₀ x))
    (hForce₀ : ForcedNonexpansion M V₀) (hNonneg₀ : NonnegV V₀)
    (hC₀0 : 0 ≤ C₀) (hV₀C : ∀ y, V₀ y ≤ C₀)
    (hFloor₀ : ∀ y ∈ U, y ∈ Dbar → 1 ≤ V₀ y)
    (hε : ∀ i, 0 < ε i)
    (hmpW : ∀ i, ∀ x : State, Nonempty (MPFlow M 0 (W i) x))
    (hSegDWR : ∀ i, RankingSegmentDynkin M (((D \ U) ∩ Rq i) \ G i) (W i) (ε i))
    (hForceW : ∀ i, ForcedNonexpansion M (W i)) (hNonnegW : ∀ i, NonnegV (W i))
    (hC0 : ∀ i, 0 ≤ C i) (hWC : ∀ i, ∀ y, W i y ≤ C i)
    (hx : x ∈ D) (hthr : V₀ x ≤ thr) :
    ENNReal.ofReal p ≤ Battery.pathMeasure (skel M) (SkelState.alive x)
      {ω | ∀ i, (∃ᶠ n in atTop, ω n ∈ skelIn (Rq i)) →
        ∃ᶠ n in atTop, ω n ∈ skelIn (G i)}
\end{leansrc}
\par\medskip\noindent{\small \texttt{Qtsl.LayerD.}\texttt{\bfseries sound\_GSHS\_streett\_bundle\_mp\_strict}}\hfill{\footnotesize \path{lean/qtsl-aux/QtslAux/LayerD/MPDischarge.lean}, line 1074}\par\nopagebreak
\begin{leansrc}
theorem sound_GSHS_streett_bundle_mp_strict
    {ι : Type*} [Countable ι]
    {D Dbar U : Set State} {V₀ : State → ℝ} {Rq G : ι → Set State}
    {W : ι → State → ℝ} {p thr C₀ : ℝ} {ε C : ι → ℝ} {x : State}
    [IsMarkovKernel M.seg] [IsMarkovKernel M.jump]
    (hD : MeasurableSet D) (hDbar : MeasurableSet Dbar) (hU : MeasurableSet U)
    (hV₀ : Measurable V₀)
    (hRq : ∀ i, MeasurableSet (Rq i)) (hG : ∀ i, MeasurableSet (G i))
    (hW : ∀ i, Measurable (W i))
    (hΘ : p + thr ≤ 1) (hDsub : D ⊆ Dbar)
    (hLive : NoVanish M (D \ U)) (hHalt : HaltsCaught M (D \ U) (U ∩ Dbar))
    (hReset : ResetsInto M (D \ U) ((D ∪ U) ∩ Dbar))
    (hmp₀ : ∀ x : State, Nonempty (MPFlow M 0 V₀ x))
    (hForce₀ : ForcedNonexpansion M V₀) (hNonneg₀ : NonnegV V₀)
    (hC₀0 : 0 ≤ C₀) (hV₀C : ∀ y, V₀ y ≤ C₀)
    (hFloor₀ : ∀ y ∈ U, y ∈ Dbar → 1 ≤ V₀ y)
    (hε : ∀ i, 0 < ε i)
    (hmpW : ∀ i, ∀ x : State, Nonempty (MPFlow M 0 (W i) x))
    (hmpWR : ∀ i, ∀ x ∈ ((D \ U) ∩ Rq i) \ G i,
      Nonempty (MPFlowStrict M 0 (W i) x (ε i)))
    (hForceW : ∀ i, ForcedNonexpansion M (W i)) (hNonnegW : ∀ i, NonnegV (W i))
    (hC0 : ∀ i, 0 ≤ C i) (hWC : ∀ i, ∀ y, W i y ≤ C i)
    (hx : x ∈ D) (hthr : V₀ x ≤ thr) :
    ENNReal.ofReal p ≤ Battery.pathMeasure (skel M) (SkelState.alive x)
      {ω | ∀ i, (∃ᶠ n in atTop, ω n ∈ skelIn (Rq i)) →
        ∃ᶠ n in atTop, ω n ∈ skelIn (G i)}
\end{leansrc}
\par\medskip\noindent{\small \texttt{Qtsl.LayerD.}\texttt{\bfseries rankingSegmentDynkin\_of\_mp\_strict}}\hfill{\footnotesize \path{lean/qtsl-aux/QtslAux/LayerD/MPDischarge.lean}, line 1060}\par\nopagebreak
\begin{leansrc}
theorem rankingSegmentDynkin_of_mp_strict (M : SegmentModel State)
    (V : State → ℝ) [IsSFiniteKernel M.jump] (hV : Measurable V)
    {Sc : Set State} {ε : ℝ} (hε : 0 ≤ ε)
    (h : ∀ x ∈ Sc, Nonempty (MPFlowStrict M 0 V x ε)) :
    RankingSegmentDynkin M Sc V ε
\end{leansrc}
\par\medskip\noindent{\small \texttt{Qtsl.LayerD.}\texttt{\bfseries MPFlowStrict}}\hfill{\footnotesize \path{lean/qtsl-aux/QtslAux/LayerD/MPDischarge.lean}, line 1047}\par\nopagebreak
\begin{leansrc}
structure MPFlowStrict (M : SegmentModel State) (θ : ℝ) (V : State → ℝ)
    (x : State) (ε : ℝ) extends MPFlow M θ V x where …
\end{leansrc}
\par\medskip\noindent{\small \texttt{Qtsl.Excursions.}\texttt{\bfseries streett\_ae\_frequently\_response}}\hfill{\footnotesize \path{lean/qtsl-aux/QtslAux/Excursions.lean}, line 365}\par\nopagebreak
\begin{leansrc}
theorem streett_ae_frequently_response (hε : 0 < ε₀)
    (hint : ∀ n, Integrable (V n) μ) (hnonneg : 0 ≤ V)
    (hmR : ∀ n, MeasurableSet[𝒢 n] {ω | X n ω ∈ R})
    (hmG : ∀ n, MeasurableSet[𝒢 n] {ω | X n ω ∈ G})
    (hdrift : ∀ n, ∀ᵐ ω ∂μ, X n ω ∉ G → (μ[V (n + 1) | 𝒢 n]) ω ≤ V n ω)
    (hdriftR : ∀ n, ∀ᵐ ω ∂μ,
      X n ω ∈ R → X n ω ∉ G → (μ[V (n + 1) | 𝒢 n]) ω ≤ V n ω - ε₀) :
    ∀ᵐ ω ∂μ, (∃ᶠ n in atTop, X n ω ∈ R) → ∃ᶠ n in atTop, X n ω ∈ G
\end{leansrc}
\par\medskip\noindent{\small \texttt{Qtsl.Excursions.}\texttt{\bfseries streett\_window\_sum\_le}}\hfill{\footnotesize \path{lean/qtsl-aux/QtslAux/Excursions.lean}, line 319}\par\nopagebreak
\begin{leansrc}
theorem streett_window_sum_le (hε : 0 ≤ ε₀)
    (hint : ∀ n, Integrable (V n) μ) (hnonneg : 0 ≤ V)
    (hmR : ∀ n, MeasurableSet[𝒢 n] {ω | X n ω ∈ R})
    (hmG : ∀ n, MeasurableSet[𝒢 n] {ω | X n ω ∈ G})
    (hdrift : ∀ n, ∀ᵐ ω ∂μ, X n ω ∉ G → (μ[V (n + 1) | 𝒢 n]) ω ≤ V n ω)
    (hdriftR : ∀ n, ∀ᵐ ω ∂μ,
      X n ω ∈ R → X n ω ∉ G → (μ[V (n + 1) | 𝒢 n]) ω ≤ V n ω - ε₀)
    (N : ℕ) :
    ENNReal.ofReal ε₀ * ∑' j : ℕ,
        μ ({ω | ∀ k, N ≤ k → k < N + j + 1 → X k ω ∉ G} ∩ {ω | X (N + j) ω ∈ R})
      ≤ ENNReal.ofReal (∫ ω, V N ω ∂μ)
\end{leansrc}
\par\medskip\noindent{\small \texttt{Qtsl.Excursions.}\texttt{\bfseries ranking\_hitting\_sum\_le}}\hfill{\footnotesize \path{lean/qtsl-aux/QtslAux/Excursions.lean}, line 147}\par\nopagebreak
\begin{leansrc}
theorem ranking_hitting_sum_le (hε : 0 ≤ ε₀)
    (hint : ∀ n, Integrable (V n) μ) (hnonneg : 0 ≤ V)
    (hmB : ∀ n, MeasurableSet[𝒢 n] {ω | X n ω ∈ B})
    (hdrift : ∀ n, ∀ᵐ ω ∂μ, X n ω ∉ B → (μ[V (n + 1) | 𝒢 n]) ω ≤ V n ω - ε₀) :
    ENNReal.ofReal ε₀ * ∑' k, μ {ω | ∀ j < k + 1, X j ω ∉ B}
      ≤ ENNReal.ofReal (∫ ω, V 0 ω ∂μ)
\end{leansrc}
\par\medskip\noindent{\small \texttt{Qtsl.Excursions.}\texttt{\bfseries ae\_hits\_of\_ranking}}\hfill{\footnotesize \path{lean/qtsl-aux/QtslAux/Excursions.lean}, line 174}\par\nopagebreak
\begin{leansrc}
theorem ae_hits_of_ranking (hε : 0 < ε₀)
    (hint : ∀ n, Integrable (V n) μ) (hnonneg : 0 ≤ V)
    (hmB : ∀ n, MeasurableSet[𝒢 n] {ω | X n ω ∈ B})
    (hdrift : ∀ n, ∀ᵐ ω ∂μ, X n ω ∉ B → (μ[V (n + 1) | 𝒢 n]) ω ≤ V n ω - ε₀) :
    ∀ᵐ ω ∂μ, ∃ n, X n ω ∈ B
\end{leansrc}
\par\medskip\noindent{\small \texttt{MeasureTheory.Supermartingale.}\texttt{\bfseries exists\_ae\_tendsto\_of\_nonneg}}\hfill{\footnotesize \path{lean/qtsl-aux/QtslAux/Excursions.lean}, line 42}\par\nopagebreak
\begin{leansrc}
theorem Supermartingale.exists_ae_tendsto_of_nonneg [IsFiniteMeasure μ] {f : ℕ → Ω
      → ℝ}
    (hf : Supermartingale f 𝒢 μ) (hnonneg : 0 ≤ f) :
    ∀ᵐ ω ∂μ, ∃ c, Tendsto (fun n => f n ω) atTop (𝓝 c)
\end{leansrc}
\par\medskip\noindent{\small \texttt{Qtsl.QStreett.}\texttt{\bfseries streett\_event\_mono}}\hfill{\footnotesize \path{lean/qtsl-aux/QtslAux/QStreett.lean}, line 418}\par\nopagebreak
\begin{leansrc}
theorem streett_event_mono {Rq Rqhat G Ghat : Set S} (hR : Rq ⊆ Rqhat)
    (hG : Ghat ⊆ G) (ω : ℕ → S)
    (h : (∃ᶠ n in atTop, ω n ∈ Rqhat) → ∃ᶠ n in atTop, ω n ∈ Ghat) :
    (∃ᶠ n in atTop, ω n ∈ Rq) → ∃ᶠ n in atTop, ω n ∈ G
\end{leansrc}
\par\medskip\noindent{\small \texttt{Qtsl.QStreett.}\texttt{\bfseries streett\_pair\_localized}}\hfill{\footnotesize \path{lean/qtsl-aux/QtslAux/QStreett.lean}, line 244}\par\nopagebreak
\begin{leansrc}
theorem streett_pair_localized {D Dbar U Rq G : Set S} {W : S → ℝ} {ε C : ℝ} {x : S}
    (hD : MeasurableSet D) (hDbar : MeasurableSet Dbar) (hU : MeasurableSet U)
    (hRq : MeasurableSet Rq) (hG : MeasurableSet G) (hW : Measurable W)
    (hε : 0 < ε) (hDsub : D ⊆ Dbar)
    (hR3 : ∀ y ∈ D, y ∉ U → K y Dbar = 1)
    (hW0 : ∀ y ∈ Dbar, 0 ≤ W y) (hWC : ∀ y ∈ Dbar, W y ≤ C)
    (hWint : ∀ y ∈ D, y ∉ U → y ∉ G → Integrable W (K y))
    (hWdrift : ∀ y ∈ D, y ∉ U → y ∉ G → ∫ z, W z ∂K y ≤ W y)
    (hWdriftR : ∀ y ∈ D, y ∉ U → y ∈ Rq → y ∉ G → ∫ z, W z ∂K y ≤ W y - ε) :
    ∀ᵐ ω ∂pathMeasure K x, (∀ n, ω n ∈ D \ U) →
      (∃ᶠ n in atTop, ω n ∈ Rq) → ∃ᶠ n in atTop, ω n ∈ G
\end{leansrc}
\par\medskip\noindent{\small \texttt{Qtsl.Excursions.}\texttt{\bfseries streett\_window\_invariant}}\hfill{\footnotesize \path{lean/qtsl-aux/QtslAux/Excursions.lean}, line 201}\par\nopagebreak
\begin{leansrc}
theorem streett_window_invariant
    (hint : ∀ n, Integrable (V n) μ) (hnonneg : 0 ≤ V)
    (hmR : ∀ n, MeasurableSet[𝒢 n] {ω | X n ω ∈ R})
    (hmG : ∀ n, MeasurableSet[𝒢 n] {ω | X n ω ∈ G})
    (hdrift : ∀ n, ∀ᵐ ω ∂μ, X n ω ∉ G → (μ[V (n + 1) | 𝒢 n]) ω ≤ V n ω)
    (hdriftR : ∀ n, ∀ᵐ ω ∂μ,
      X n ω ∈ R → X n ω ∉ G → (μ[V (n + 1) | 𝒢 n]) ω ≤ V n ω - ε₀)
    (N : ℕ) : ∀ n, N ≤ n →
    ∫ ω in {ω | ∀ k, N ≤ k → k < n → X k ω ∉ G}, V n ω ∂μ
        + ε₀ * ∑ k ∈ Ico N n,
            μ.real ({ω | ∀ j, N ≤ j → j < k + 1 → X j ω ∉ G} ∩ {ω | X k ω ∈ R})
      ≤ ∫ ω, V N ω ∂μ
\end{leansrc}
\par\medskip\noindent{\small \texttt{Qtsl.QStreett.}\texttt{\bfseries AbsorptionBound}}\hfill{\footnotesize \path{lean/qtsl-aux/QtslAux/QStreett.lean}, line 479}\par\nopagebreak
\begin{leansrc}
def AbsorptionBound (K : Kernel S S) [IsMarkovKernel K] (x : S)
    (A J Rq : Set S) (γ : ℝ) : Prop
\end{leansrc}
\par\medskip\noindent{\small \texttt{Qtsl.QStreett.}\texttt{\bfseries ae\_not\_frequently\_of\_absorption}}\hfill{\footnotesize \path{lean/qtsl-aux/QtslAux/QStreett.lean}, line 490}\par\nopagebreak
\begin{leansrc}
theorem ae_not_frequently_of_absorption {A J Rq : Set S} {γ : ℝ} {x : S}
    (hRq : MeasurableSet Rq) (hγ : 0 < γ)
    (habs : AbsorptionBound K x A J Rq γ) :
    ∀ᵐ ω ∂pathMeasure K x,
      (∃ᶠ k in atTop, ω ∈ survive A k ∧ ω k ∈ J) →
        ¬ ∃ᶠ n in atTop, ω n ∈ Rq
\end{leansrc}
\par\medskip\noindent{\small \texttt{Qtsl.QStreett.}\texttt{\bfseries supermartingale\_stopped\_shift}}\hfill{\footnotesize \path{lean/qtsl-aux/QtslAux/QStreett.lean}, line 740}\par\nopagebreak
\begin{leansrc}
theorem supermartingale_stopped_shift {f : S → ℝ} {R : Set S} {C : ℝ}
    (hR : MeasurableSet R) (hf : Measurable f) (hf0 : ∀ y, 0 ≤ f y) (hfC : ∀ y, f y
      ≤ C)
    (hdrift : ∀ y, y ∉ R → ∫ z, f z ∂K y ≤ f y) (k : ℕ) (u : Π _i : Iic k, S) :
    Supermartingale (fun n ω => stopped f R n (shift k ω)) (piLEShift k)
      (Kernel.traj (X := fun _ : ℕ => S) (κ := mchain K) k u)
\end{leansrc}
\par\medskip\noindent{\small \texttt{Qtsl.QStreett.}\texttt{\bfseries measure\_reach\_after\_le}}\hfill{\footnotesize \path{lean/qtsl-aux/QtslAux/QStreett.lean}, line 760}\par\nopagebreak
\begin{leansrc}
theorem measure_reach_after_le {f : S → ℝ} {R : Set S} {C : ℝ}
    (hR : MeasurableSet R) (hf : Measurable f) (hf0 : ∀ y, 0 ≤ f y) (hfC : ∀ y, f y
      ≤ C)
    (hf1 : ∀ y ∈ R, 1 ≤ f y) (hdrift : ∀ y, y ∉ R → ∫ z, f z ∂K y ≤ f y)
    (k : ℕ) (u : Π _i : Iic k, S) :
    Kernel.traj (X := fun _ : ℕ => S) (κ := mchain K) k u {ω | ∃ n, k ≤ n ∧ ω n ∈ R}
      ≤ ENNReal.ofReal (f (u ⟨k, mem_Iic.2 le_rfl⟩))
\end{leansrc}
\par\medskip\noindent{\small \texttt{PPLC.Qtsl.Behavior.}\texttt{\bfseries AService}}\hfill{\footnotesize \path{lean/qtsl/Qtsl/TimedCore.lean}, line 32}\par\nopagebreak
\begin{leansrc}
def Behavior.AService (B : Behavior X) (T : ℝ≥0) (φ : Set X) (x : X) : Prop
\end{leansrc}
\par\medskip\noindent{\small \texttt{PPLC.Qtsl.Behavior.}\texttt{\bfseries BoxInv}}\hfill{\footnotesize \path{lean/qtsl/Qtsl/TimedCore.lean}, line 35}\par\nopagebreak
\begin{leansrc}
def Behavior.BoxInv (B : Behavior X) (φ : Set X) (x : X) : Prop
\end{leansrc}
\par\medskip\noindent{\small \texttt{PPLC.Qtsl.Behavior.}\texttt{\bfseries FusionClosed}}\hfill{\footnotesize \path{lean/qtsl/Qtsl/TimedCore.lean}, line 28}\par\nopagebreak
\begin{leansrc}
def Behavior.FusionClosed (B : Behavior X) : Prop
\end{leansrc}
\par\medskip\noindent{\small \texttt{PPLC.Qtsl.Behavior.}\texttt{\bfseries FirstVisits}}\hfill{\footnotesize \path{lean/qtsl/Qtsl/TimedCore.lean}, line 38}\par\nopagebreak
\begin{leansrc}
def Behavior.FirstVisits (B : Behavior X) (φ : Set X) : Prop
\end{leansrc}
\par\medskip\noindent{\small \texttt{PPLC.Qtsl.Behavior.}\texttt{\bfseries fsTarget}}\hfill{\footnotesize \path{lean/qtsl/Qtsl/TimedCore.lean}, line 41}\par\nopagebreak
\begin{leansrc}
def Behavior.fsTarget (B : Behavior X) {n : ℕ}
    (α : Fin n → Set X) (T : Fin n → ℝ≥0) : Set X
\end{leansrc}
\par\medskip\noindent{\small \texttt{PPLC.Qtsl.Behavior.}\texttt{\bfseries FirstVisits}}\hfill{\footnotesize \path{lean/qtsl/Qtsl/TimedCore.lean}, line 38}\par\nopagebreak
\begin{leansrc}
def Behavior.FirstVisits (B : Behavior X) (φ : Set X) : Prop
\end{leansrc}
\par\medskip\noindent{\small \texttt{PPLC.Qtsl.Behavior.}\texttt{\bfseries residualCell}}\hfill{\footnotesize \path{lean/qtsl/Qtsl/TimedCore.lean}, line 45}\par\nopagebreak
\begin{leansrc}
def Behavior.residualCell (B : Behavior X) {n : ℕ}
    (α : Fin n → Set X) (T : Fin n → ℝ≥0) (c : ℝ≥0) : Set X
\end{leansrc}
\par\medskip\noindent{\small \texttt{Qtsl.NNF.}\texttt{\bfseries denot\_elimNeg}}\hfill{\footnotesize \path{lean/qtsl-aux/QtslAux/NNFDuality.lean}, line 198}\par\nopagebreak
\begin{leansrc}
theorem denot_elimNeg {σ : Type*} (den : α → Set σ) (d : α → α)
    (hd : ∀ a, den (d a) = (den a)ᶜ) :
    ∀ φ : NForm α, (φ.elimNeg d).denot den = φ.denot den
\end{leansrc}
\par\medskip\noindent{\small \texttt{Qtsl.QStreett.}\texttt{\bfseries ranking\_guard\_collapse}}\hfill{\footnotesize \path{lean/qtsl-aux/QtslAux/QStreett.lean}, line 464}\par\nopagebreak
\begin{leansrc}
theorem ranking_guard_collapse {W : S → ℝ} {T slab : Set S} {g : S → S}
    (hzero : ∀ y ∈ T, W y = 0) (himg : ∀ y ∈ slab, g y ∈ T) :
    ∀ y ∈ slab, 0 ≤ W y → W (g y) ≤ W y
\end{leansrc}
\par\medskip\noindent{\small \texttt{Qtsl.NNF.}\texttt{\bfseries compl\_le\_atom}}\hfill{\footnotesize \path{lean/qtsl-aux/QtslAux/NNFDuality.lean}, line 36}\par\nopagebreak
\begin{leansrc}
theorem compl_le_atom : {x | f x ≤ c}ᶜ = {x | c < f x}
\end{leansrc}
\par\medskip\noindent{\small \texttt{Qtsl.NNF.}\texttt{\bfseries compl\_lt\_atom}}\hfill{\footnotesize \path{lean/qtsl-aux/QtslAux/NNFDuality.lean}, line 41}\par\nopagebreak
\begin{leansrc}
theorem compl_lt_atom : {x | f x < c}ᶜ = {x | c ≤ f x}
\end{leansrc}
\par\medskip\noindent{\small \texttt{Qtsl.NNF.}\texttt{\bfseries compl\_ge\_atom}}\hfill{\footnotesize \path{lean/qtsl-aux/QtslAux/NNFDuality.lean}, line 46}\par\nopagebreak
\begin{leansrc}
theorem compl_ge_atom : {x | c ≤ f x}ᶜ = {x | f x < c}
\end{leansrc}
\par\medskip\noindent{\small \texttt{Qtsl.NNF.}\texttt{\bfseries compl\_gt\_atom}}\hfill{\footnotesize \path{lean/qtsl-aux/QtslAux/NNFDuality.lean}, line 51}\par\nopagebreak
\begin{leansrc}
theorem compl_gt_atom : {x | c < f x}ᶜ = {x | f x ≤ c}
\end{leansrc}
\par\medskip\noindent{\small \texttt{Qtsl.NNF.}\texttt{\bfseries denot\_nnfAux}}\hfill{\footnotesize \path{lean/qtsl-aux/QtslAux/NNFDuality.lean}, line 111}\par\nopagebreak
\begin{leansrc}
theorem denot_nnfAux : ∀ (φ : Form α) (b : Bool),
    (φ.nnfAux b).denot den = if b then φ.denot den else (φ.denot den)ᶜ
\end{leansrc}
\par\medskip\noindent{\small \texttt{Qtsl.NNF.Form.}\texttt{\bfseries nnf}}\hfill{\footnotesize \path{lean/qtsl-aux/QtslAux/NNFDuality.lean}, line 87}\par\nopagebreak
\begin{leansrc}
def Form.nnf (φ : Form α) : NForm α
\end{leansrc}
\par\medskip\noindent{\small \texttt{Qtsl.NNF.}\texttt{\bfseries denot\_nnf}}\hfill{\footnotesize \path{lean/qtsl-aux/QtslAux/NNFDuality.lean}, line 128}\par\nopagebreak
\begin{leansrc}
theorem denot_nnf (φ : Form α) : (φ.nnf).denot den = φ.denot den
\end{leansrc}
\par\medskip\noindent{\small \texttt{Qtsl.NNF.}\texttt{\bfseries negLits\_nnf}}\hfill{\footnotesize \path{lean/qtsl-aux/QtslAux/NNFDuality.lean}, line 170}\par\nopagebreak
\begin{leansrc}
theorem negLits_nnf (φ : Form α) : (φ.nnf).negLits = φ.occs false
\end{leansrc}
\par\medskip\noindent{\small \texttt{Qtsl.NNF.}\texttt{\bfseries nnf\_no\_negLit\_iff}}\hfill{\footnotesize \path{lean/qtsl-aux/QtslAux/NNFDuality.lean}, line 175}\par\nopagebreak
\begin{leansrc}
theorem nnf_no_negLit_iff (φ : Form α) (P : α → Prop) :
    (∀ a ∈ (φ.nnf).negLits, ¬ P a) ↔ ∀ a ∈ φ.occs false, ¬ P a
\end{leansrc}
\par\medskip\noindent{\small \texttt{Qtsl.NNF.}\texttt{\bfseries negLits\_elimNeg}}\hfill{\footnotesize \path{lean/qtsl-aux/QtslAux/NNFDuality.lean}, line 188}\par\nopagebreak
\begin{leansrc}
theorem negLits_elimNeg (d : α → α) : ∀ φ : NForm α, (φ.elimNeg d).negLits = ∅
\end{leansrc}
\par\medskip\noindent{\small \texttt{Qtsl.NNF.}\texttt{\bfseries negLits\_nnfAux}}\hfill{\footnotesize \path{lean/qtsl-aux/QtslAux/NNFDuality.lean}, line 153}\par\nopagebreak
\begin{leansrc}
theorem negLits_nnfAux : ∀ (φ : Form α) (b : Bool),
    (φ.nnfAux b).negLits = φ.occs (!b)
\end{leansrc}
\par\medskip\noindent{\small \texttt{Qtsl.Aliases.}\texttt{\bfseries sugarAvoid\_denot}}\hfill{\footnotesize \path{lean/qtsl-aux/QtslAux/QtslAliases.lean}, line 95}\par\nopagebreak
\begin{leansrc}
theorem sugarAvoid_denot {S : Type*} (B : Set S) :
    {ω : ℕ → S | ∀ n, ω n ∉ B} = {ω : ℕ → S | ∀ n, ω n ∈ Bᶜ}
\end{leansrc}
\par\medskip\noindent{\small \texttt{Qtsl.Aliases.}\texttt{\bfseries sugarGF\_denot}}\hfill{\footnotesize \path{lean/qtsl-aux/QtslAux/QtslAliases.lean}, line 98}\par\nopagebreak
\begin{leansrc}
theorem sugarGF_denot {S : Type*} (G : Set S) (ω : ℕ → S) :
    ((∃ᶠ n in atTop, ω n ∈ (Set.univ : Set S)) → ∃ᶠ n in atTop, ω n ∈ G)
      ↔ ∃ᶠ n in atTop, ω n ∈ G
\end{leansrc}
\par\medskip\noindent{\small \texttt{Qtsl.Aliases.}\texttt{\bfseries sugarFG\_denot}}\hfill{\footnotesize \path{lean/qtsl-aux/QtslAux/QtslAliases.lean}, line 105}\par\nopagebreak
\begin{leansrc}
theorem sugarFG_denot {S : Type*} (B : Set S) (ω : ℕ → S) :
    ((∃ᶠ n in atTop, ω n ∈ B) → ∃ᶠ n in atTop, ω n ∈ (∅ : Set S))
      ↔ ∀ᶠ n in atTop, ω n ∉ B
\end{leansrc}
\par\medskip\noindent{\small \texttt{Qtsl.Aliases.}\texttt{\bfseries negF\_denot\_subset}}\hfill{\footnotesize \path{lean/qtsl-aux/QtslAux/QtslAliases.lean}, line 84}\par\nopagebreak
\begin{leansrc}
theorem negF_denot_subset {S : Type*} (sat : S → Prop) (T : ℕ) :
    {ω : ℕ → S | ∀ n, ¬ sat (ω n)} ⊆ {ω : ℕ → S | ¬ ∃ k ≤ T, sat (ω k)}
\end{leansrc}
\par\medskip\noindent{\small \texttt{Qtsl.Structural.}\texttt{\bfseries hypMultiplier\_sound}}\hfill{\footnotesize \path{lean/qtsl-aux/QtslAux/Structural.lean}, line 75}\par\nopagebreak
\begin{leansrc}
theorem hypMultiplier_sound (cell childAtom : Set X) (N lam g : X → ℝ)
    (hbridge : childAtom = {x | 0 ≤ g x})
    (hlam : ∀ x ∈ cell, 0 ≤ lam x)
    (hcomb : ∀ x ∈ cell, N x + lam x * g x ≤ 0) :
    ∀ x ∈ cell ∩ childAtom, N x ≤ 0
\end{leansrc}
\par\medskip\noindent{\small \texttt{Qtsl.Structural.}\texttt{\bfseries fullDomainWeaken\_sound}}\hfill{\footnotesize \path{lean/qtsl-aux/QtslAux/Structural.lean}, line 88}\par\nopagebreak
\begin{leansrc}
theorem fullDomainWeaken_sound (Dm scope : Set X) (P : X → Prop)
    (hbridge : scope ⊆ Dm)
    (hrow : ∀ x ∈ Dm, P x) :
    ∀ x ∈ scope, P x
\end{leansrc}
\par\medskip\noindent{\small \texttt{Qtsl.Structural.}\texttt{\bfseries fullDomainWeaken\_offScope}}\hfill{\footnotesize \path{lean/qtsl-aux/QtslAux/Structural.lean}, line 94}\par\nopagebreak
\begin{leansrc}
theorem fullDomainWeaken_offScope (Dm Off : Set X) (P : X → Prop)
    (hrow : ∀ x ∈ Dm, P x) :
    ∀ x ∈ Dm \ Off, P x
\end{leansrc}
\par\medskip\noindent{\small \texttt{Qtsl.Structural.}\texttt{\bfseries regionSat\_mono}}\hfill{\footnotesize \path{lean/qtsl-aux/QtslAux/Structural.lean}, line 62}\par\nopagebreak
\begin{leansrc}
theorem regionSat_mono {R S den : Set X} (h : RegionSat R den) (hsub : S ⊆ R) :
    RegionSat S den
\end{leansrc}
\par\medskip\noindent{\small \texttt{Qtsl.Structural.}\texttt{\bfseries verifiedCover\_components\_sound}}\hfill{\footnotesize \path{lean/qtsl-aux/QtslAux/Structural.lean}, line 116}\par\nopagebreak
\begin{leansrc}
theorem verifiedCover_components_sound {ι : Type*} (Dm C Off scope : Set X)
    (K : ι → Set X) (P : X → Prop)
    (hbridge : scope = Dm \ Off)
    (hrowC : ∀ x ∈ C, P x)
    (hdecomp : Dm \ C ⊆ ⋃ i, K i)
    (hcomp : ∀ i, K i ⊆ Off) :
    ∀ x ∈ scope, P x
\end{leansrc}
\par\medskip\noindent{\small \texttt{Qtsl.Structural.}\texttt{\bfseries verifiedCover\_sound}}\hfill{\footnotesize \path{lean/qtsl-aux/QtslAux/Structural.lean}, line 102}\par\nopagebreak
\begin{leansrc}
theorem verifiedCover_sound (Dm C Off scope : Set X) (P : X → Prop)
    (hbridge : scope = Dm \ Off)
    (hrowC : ∀ x ∈ C, P x)
    (hresidual : Dm \ C ⊆ Off) :
    ∀ x ∈ scope, P x
\end{leansrc}
\par\medskip\noindent{\small \texttt{Qtsl.Structural.}\texttt{\bfseries RegionSat}}\hfill{\footnotesize \path{lean/qtsl-aux/QtslAux/Structural.lean}, line 57}\par\nopagebreak
\begin{leansrc}
def RegionSat (R den : Set X) : Prop
\end{leansrc}
\par\medskip\noindent{\small \texttt{Qtsl.Structural.}\texttt{\bfseries positiveCombinationAnd\_sound}}\hfill{\footnotesize \path{lean/qtsl-aux/QtslAux/Structural.lean}, line 129}\par\nopagebreak
\begin{leansrc}
theorem positiveCombinationAnd_sound (R₁ R₂ den₁ den₂ denAnd : Set X)
    (hbridge : denAnd = den₁ ∩ den₂)
    (h₁ : RegionSat R₁ den₁) (h₂ : RegionSat R₂ den₂) :
    RegionSat (R₁ ∩ R₂) denAnd
\end{leansrc}
\par\medskip\noindent{\small \texttt{Qtsl.Structural.}\texttt{\bfseries positiveCombinationOr\_sound}}\hfill{\footnotesize \path{lean/qtsl-aux/QtslAux/Structural.lean}, line 135}\par\nopagebreak
\begin{leansrc}
theorem positiveCombinationOr_sound (R₁ R₂ den₁ den₂ denOr : Set X)
    (hbridge : denOr = den₁ ∪ den₂)
    (h₁ : RegionSat R₁ den₁) (h₂ : RegionSat R₂ den₂) :
    RegionSat (R₁ ∪ R₂) denOr
\end{leansrc}
\par\medskip\noindent{\small \texttt{Qtsl.Structural.}\texttt{\bfseries positiveCombinationInter\_sound}}\hfill{\footnotesize \path{lean/qtsl-aux/QtslAux/Structural.lean}, line 141}\par\nopagebreak
\begin{leansrc}
theorem positiveCombinationInter_sound {ι : Type*} (R den : ι → Set X)
    (denAll : Set X)
    (hbridge : denAll = ⋂ i, den i)
    (h : ∀ i, RegionSat (R i) (den i)) :
    RegionSat (⋂ i, R i) denAll
\end{leansrc}
\par\medskip\noindent{\small \texttt{Qtsl.Structural.}\texttt{\bfseries positiveCombinationUnion\_sound}}\hfill{\footnotesize \path{lean/qtsl-aux/QtslAux/Structural.lean}, line 148}\par\nopagebreak
\begin{leansrc}
theorem positiveCombinationUnion_sound {ι : Type*} (R den : ι → Set X)
    (denAll : Set X)
    (hbridge : denAll = ⋃ i, den i)
    (h : ∀ i, RegionSat (R i) (den i)) :
    RegionSat (⋃ i, R i) denAll
\end{leansrc}
\par\medskip\noindent{\small \texttt{Qtsl.Structural.}\texttt{\bfseries compl\_over\_approx}}\hfill{\footnotesize \path{lean/qtsl-aux/QtslAux/Structural.lean}, line 201}\par\nopagebreak
\begin{leansrc}
theorem compl_over_approx {R den : Set X} (h : RegionSat R den) :
    denᶜ ⊆ Rᶜ
\end{leansrc}
\par\medskip\noindent{\small \texttt{Qtsl.Structural.}\texttt{\bfseries closureExactLt\_sound}}\hfill{\footnotesize \path{lean/qtsl-aux/QtslAux/Structural.lean}, line 157}\par\nopagebreak
\begin{leansrc}
theorem closureExactLt_sound (p : X → ℝ) (c : ℝ) (strictAtom : Set X)
    (hbridge : strictAtom = {x | p x < c}) :
    strictAtomᶜ = {x | c ≤ p x}
\end{leansrc}
\par\medskip\noindent{\small \texttt{Qtsl.Structural.}\texttt{\bfseries closureExact\_isClosed\_ge}}\hfill{\footnotesize \path{lean/qtsl-aux/QtslAux/Structural.lean}, line 171}\par\nopagebreak
\begin{leansrc}
theorem closureExact_isClosed_ge [TopologicalSpace X] {p : X → ℝ}
    (hp : Continuous p) (c : ℝ) : IsClosed {x | c ≤ p x}
\end{leansrc}
\par\medskip\noindent{\small \texttt{Qtsl.Structural.}\texttt{\bfseries closureExactGt\_sound}}\hfill{\footnotesize \path{lean/qtsl-aux/QtslAux/Structural.lean}, line 164}\par\nopagebreak
\begin{leansrc}
theorem closureExactGt_sound (p : X → ℝ) (c : ℝ) (strictAtom : Set X)
    (hbridge : strictAtom = {x | c < p x}) :
    strictAtomᶜ = {x | p x ≤ c}
\end{leansrc}
\par\medskip\noindent{\small \texttt{Qtsl.Structural.}\texttt{\bfseries closureExact\_isClosed\_le}}\hfill{\footnotesize \path{lean/qtsl-aux/QtslAux/Structural.lean}, line 175}\par\nopagebreak
\begin{leansrc}
theorem closureExact_isClosed_le [TopologicalSpace X] {p : X → ℝ}
    (hp : Continuous p) (c : ℝ) : IsClosed {x | p x ≤ c}
\end{leansrc}
\par\medskip\noindent{\small \texttt{Qtsl.Structural.}\texttt{\bfseries closureOutward\_sound}}\hfill{\footnotesize \path{lean/qtsl-aux/QtslAux/Structural.lean}, line 182}\par\nopagebreak
\begin{leansrc}
theorem closureOutward_sound (p : X → ℝ) (c : ℝ) (cmplAtom : Set X)
    (hbridge : cmplAtom = {x | c ≤ p x}) :
    {x | p x ≤ c}ᶜ ⊆ cmplAtom
\end{leansrc}
\par\medskip\noindent{\small \texttt{Qtsl.NNF.NForm.}\texttt{\bfseries elimNeg}}\hfill{\footnotesize \path{lean/qtsl-aux/QtslAux/NNFDuality.lean}, line 181}\par\nopagebreak
\begin{leansrc}
def NForm.elimNeg (d : α → α) : NForm α → NForm α
  | .pos a => .pos a
  | .negLit a => .pos (d a)
  | .and φ ψ => .and (φ.elimNeg d) (ψ.elimNeg d)
  | .or φ ψ => .or (φ.elimNeg d) (ψ.elimNeg d)

theorem negLits_elimNeg (d : α → α) : ∀ φ : NForm α, (φ.elimNeg d).negLits = ∅
\end{leansrc}
\par\medskip\noindent{\small \texttt{Qtsl.Structural.}\texttt{\bfseries closureOutwardGe\_sound}}\hfill{\footnotesize \path{lean/qtsl-aux/QtslAux/Structural.lean}, line 190}\par\nopagebreak
\begin{leansrc}
theorem closureOutwardGe_sound (p : X → ℝ) (c : ℝ) (cmplAtom : Set X)
    (hbridge : cmplAtom = {x | p x ≤ c}) :
    {x | c ≤ p x}ᶜ ⊆ cmplAtom
\end{leansrc}
\par\medskip\noindent{\small \texttt{Qtsl.Structural.}\texttt{\bfseries closureInwardEpsLt\_sound}}\hfill{\footnotesize \path{lean/qtsl-aux/QtslAux/Structural.lean}, line 209}\par\nopagebreak
\begin{leansrc}
theorem closureInwardEpsLt_sound (p : X → ℝ) (c ε : ℝ) (hε : 0 < ε)
    (roundedAtom : Set X)
    (hbridge : roundedAtom = {x | p x ≤ c - ε}) :
    roundedAtom ⊆ {x | p x < c}
\end{leansrc}
\par\medskip\noindent{\small \texttt{Qtsl.Structural.}\texttt{\bfseries closureInwardEps\_isClosed}}\hfill{\footnotesize \path{lean/qtsl-aux/QtslAux/Structural.lean}, line 229}\par\nopagebreak
\begin{leansrc}
theorem closureInwardEps_isClosed [TopologicalSpace X] {p : X → ℝ}
    (hp : Continuous p) (c ε : ℝ) : IsClosed {x | p x ≤ c - ε}
\end{leansrc}
\par\medskip\noindent{\small \texttt{Qtsl.Structural.}\texttt{\bfseries closureInwardEpsGt\_sound}}\hfill{\footnotesize \path{lean/qtsl-aux/QtslAux/Structural.lean}, line 219}\par\nopagebreak
\begin{leansrc}
theorem closureInwardEpsGt_sound (p : X → ℝ) (c ε : ℝ) (hε : 0 < ε)
    (roundedAtom : Set X)
    (hbridge : roundedAtom = {x | c + ε ≤ p x}) :
    roundedAtom ⊆ {x | c < p x}
\end{leansrc}
\par\medskip\noindent{\small \texttt{Qtsl.Aliases.}\texttt{\bfseries taylorLeaf\_le\_exp}}\hfill{\footnotesize \path{lean/qtsl-aux/QtslAux/QtslAliases.lean}, line 66}\par\nopagebreak
\begin{leansrc}
theorem taylorLeaf_le_exp {q x : ℝ} {n : ℕ} (hx : 0 ≤ x)
    (hq : q ≤ ∑ k ∈ range (n + 1), x ^ k / (Nat.factorial k : ℝ)) :
    q ≤ Real.exp x
\end{leansrc}
\par\medskip\noindent{\small \texttt{Qtsl.Aliases.}\texttt{\bfseries taylorLeaf\_rat\_le\_exp}}\hfill{\footnotesize \path{lean/qtsl-aux/QtslAux/QtslAliases.lean}, line 71}\par\nopagebreak
\begin{leansrc}
theorem taylorLeaf_rat_le_exp {qq θ : ℚ} {T n : ℕ} (hθ : 0 ≤ θ)
    (hleaf : qq ≤ ∑ k ∈ range (n + 1), (θ * T) ^ k / (Nat.factorial k : ℚ)) :
    (qq : ℝ) ≤ Real.exp (θ * T)
\end{leansrc}
\par\medskip\noindent{\small \texttt{PPLC.}\texttt{\bfseries taylorElab\_le\_exp}}\hfill{\footnotesize \path{lean/qtsl/Qtsl/Elaborators.lean}, line 224}\par\nopagebreak
\begin{leansrc}
theorem taylorElab_le_exp {q θ : Rat} {T n : Nat} (hθ : 0 ≤ θ)
    (hEval : (taylorFormula q θ T n).evalClosed = some true) :
    (q : ℝ) ≤ Real.exp (θ * T)
\end{leansrc}
\par\medskip\noindent{\small \texttt{PPLC.}\texttt{\bfseries taylorElab}}\hfill{\footnotesize \path{lean/qtsl/Qtsl/Elaborators.lean}, line 199}\par\nopagebreak
\begin{leansrc}
def taylorElab (H : Hypernet) (seq : Sequent) (q θ : Rat) (T n : Nat)
    (hEval : (taylorFormula q θ T n).evalClosed = some true)
    (hpost : seq.post = taylorFormula q θ T n) :
    PConstExactData H seq
\end{leansrc}
\par\medskip\noindent{\small \texttt{PPLC.}\texttt{\bfseries taylorPartialSum}}\hfill{\footnotesize \path{lean/qtsl/Qtsl/Elaborators.lean}, line 180}\par\nopagebreak
\begin{leansrc}
def taylorPartialSum (θ : Rat) (T n : Nat) : Rat
\end{leansrc}
\par\medskip\noindent{\small \texttt{PPLC.}\texttt{\bfseries taylorElab\_sound}}\hfill{\footnotesize \path{lean/qtsl/Qtsl/Elaborators.lean}, line 208}\par\nopagebreak
\begin{leansrc}
theorem taylorElab_sound {H : Hypernet} {seq : Sequent} {q θ : Rat} {T n : Nat}
    (hEval : (taylorFormula q θ T n).evalClosed = some true)
    (hpost : seq.post = taylorFormula q θ T n) :
    seq.semanticallyValid H
\end{leansrc}
\par\medskip\noindent{\small \texttt{PPLC.}\texttt{\bfseries taylorFormula\_true\_iff}}\hfill{\footnotesize \path{lean/qtsl/Qtsl/Elaborators.lean}, line 190}\par\nopagebreak
\begin{leansrc}
theorem taylorFormula_true_iff (q θ : Rat) (T n : Nat) :
    (taylorFormula q θ T n).evalClosed = some true
      ↔ q ≤ taylorPartialSum θ T n
\end{leansrc}
\par\medskip\noindent{\small \texttt{Qtsl.Aliases.}\texttt{\bfseries rootEval\_eval\_cast}}\hfill{\footnotesize \path{lean/qtsl-aux/QtslAux/QtslAliases.lean}, line 39}\par\nopagebreak
\begin{leansrc}
theorem rootEval_eval_cast {σ : Type*} (P : MvPolynomial σ ℚ) (xq : σ → ℚ) :
    ((MvPolynomial.eval xq P : ℚ) : ℝ)
      = MvPolynomial.eval (fun i => (xq i : ℝ)) (MvPolynomial.map (Rat.castHom ℝ) P)
\end{leansrc}
\par\medskip\noindent{\small \texttt{Qtsl.Aliases.}\texttt{\bfseries rootEval\_sublevel\_iff}}\hfill{\footnotesize \path{lean/qtsl-aux/QtslAux/QtslAliases.lean}, line 45}\par\nopagebreak
\begin{leansrc}
theorem rootEval_sublevel_iff {σ : Type*} (P : MvPolynomial σ ℚ) (xq : σ → ℚ)
    (thr : ℚ) :
    MvPolynomial.eval xq P ≤ thr
      ↔ MvPolynomial.eval (fun i => (xq i : ℝ))
          (MvPolynomial.map (Rat.castHom ℝ) P) ≤ (thr : ℝ)
\end{leansrc}
\par\medskip\noindent{\small \texttt{Qtsl.Aliases.}\texttt{\bfseries rootEval\_box\_iff}}\hfill{\footnotesize \path{lean/qtsl-aux/QtslAux/QtslAliases.lean}, line 53}\par\nopagebreak
\begin{leansrc}
theorem rootEval_box_iff {ι : Type*} (aq bq xq : ι → ℚ) :
    (∀ i, aq i ≤ xq i ∧ xq i ≤ bq i)
      ↔ ∀ i, ((aq i : ℝ) ≤ (xq i : ℝ) ∧ (xq i : ℝ) ≤ (bq i : ℝ))
\end{leansrc}
\par\medskip\noindent{\small \texttt{PPLC.}\texttt{\bfseries polyClosedTerm}}\hfill{\footnotesize \path{lean/qtsl/Qtsl/Elaborators.lean}, line 100}\par\nopagebreak
\begin{leansrc}
def polyClosedTerm (x0 : Nat → Rat) (p : Poly) : Term
\end{leansrc}
\par\medskip\noindent{\small \texttt{PPLC.}\texttt{\bfseries rootEvalElab}}\hfill{\footnotesize \path{lean/qtsl/Qtsl/Elaborators.lean}, line 147}\par\nopagebreak
\begin{leansrc}
def rootEvalElab (H : Hypernet) (seq : Sequent) (x0 : Nat → Rat) (V : Poly)
    (thr : Rat)
    (hEval : (rootEvalFormula x0 V thr).evalClosed = some true)
    (hpost : seq.post = rootEvalFormula x0 V thr) :
    PConstExactData H seq
\end{leansrc}
\par\medskip\noindent{\small \texttt{PPLC.}\texttt{\bfseries polyClosedTerm\_evalQ}}\hfill{\footnotesize \path{lean/qtsl/Qtsl/Elaborators.lean}, line 104}\par\nopagebreak
\begin{leansrc}
@[simp] theorem polyClosedTerm_evalQ (x0 : Nat → Rat) :
    ∀ p : Poly, (polyClosedTerm x0 p).evalQ? = some (Poly.eval x0 p)
  | [] => by simp [polyClosedTerm, Poly.eval, Term.evalQ?]
  | mc :: rest => by
    have ih
\end{leansrc}
\par\medskip\noindent{\small \texttt{PPLC.}\texttt{\bfseries rootEvalElab\_sound}}\hfill{\footnotesize \path{lean/qtsl/Qtsl/Elaborators.lean}, line 158}\par\nopagebreak
\begin{leansrc}
theorem rootEvalElab_sound {H : Hypernet} {seq : Sequent} {x0 : Nat → Rat}
    {V : Poly} {thr : Rat}
    (hEval : (rootEvalFormula x0 V thr).evalClosed = some true)
    (hpost : seq.post = rootEvalFormula x0 V thr) :
    seq.semanticallyValid H
\end{leansrc}
\par\medskip\noindent{\small \texttt{SuFunctor.Universe.}\texttt{\bfseries ratBox}}\hfill{\footnotesize \path{lean/su-functor/SuFunctor/Universe/Generation.lean}, line 67}\par\nopagebreak
\begin{leansrc}
def ratBox (l u : V → ℚ) : Set (ProbState V)
\end{leansrc}
\par\medskip\noindent{\small \texttt{Qtsl.GridCx.}\texttt{\bfseries Box}}\hfill{\footnotesize \path{lean/qtsl-aux/QtslAux/GridCx.lean}, line 49}\par\nopagebreak
\begin{leansrc}
structure Box (d : ℕ) where …
\end{leansrc}
\par\medskip\noindent{\small \texttt{Qtsl.Repar.AffineWitness.}\texttt{\bfseries t0\_soundness}}\hfill{\footnotesize \path{lean/qtsl-aux/QtslAux/Repar.lean}, line 196}\par\nopagebreak
\begin{leansrc}
theorem t0_soundness (N Nhat : MvPolynomial (Fin n) ℚ)
    (X D : Set (Fin n → ℝ))
    (hreplay : Nhat = φ.subst N)
    (hcorr : φ.invMapA '' X ⊆ D)
    (hcells : ∀ w ∈ D, aeval w Nhat ≤ 0) :
    ∀ z ∈ X, aeval z N ≤ 0
\end{leansrc}
\par\medskip\noindent{\small \texttt{Qtsl.Repar.AffineWitness.}\texttt{\bfseries t05\_soundness\_gridcx}}\hfill{\footnotesize \path{lean/qtsl-aux/QtslAux/Repar.lean}, line 247}\par\nopagebreak
\begin{leansrc}
theorem t05_soundness_gridcx {n : ℕ} {ι : Type*} (φs : ι → AffineWitness n)
    (N : MvPolynomial (Fin n) ℚ) (Nhat : ι → MvPolynomial (Fin n) ℚ)
    (B : ι → GridCx.Box n) (T : ∀ i, GridCx.PatchTree n (B i))
    (X : Set (Fin n → ℝ)) (Xp : ι → Set (Fin n → ℝ))
    (hX : X ⊆ ⋃ i, Xp i)
    (hreplay : ∀ i, Nhat i = (φs i).subst N)
    (hcorr : ∀ i, (φs i).invMapA '' Xp i ⊆ (B i).set)
    (hcells : ∀ i (c : (T i).CellIdx), ∀ w ∈ (T i).cell c,
      aeval w (Nhat i) ≤ 0) :
    ∀ z ∈ X, aeval z N ≤ 0
\end{leansrc}
\par\medskip\noindent{\small \texttt{Qtsl.Repar.AffineWitness.}\texttt{\bfseries left\_certified}}\hfill{\footnotesize \path{lean/qtsl-aux/QtslAux/Repar.lean}, line 77}\par\nopagebreak
\begin{leansrc}
theorem left_certified : φ.Ainv * φ.A = 1
\end{leansrc}
\par\medskip\noindent{\small \texttt{Qtsl.BoxBern.}\texttt{\bfseries boxBern\_enclosure}}\hfill{\footnotesize \path{lean/qtsl-aux/QtslAux/BoxBern.lean}, line 170}\par\nopagebreak
\begin{leansrc}
theorem boxBern_enclosure (lo hi : Fin n → ℚ) (N : MvPolynomial (Fin n) ℚ)
    (d : Fin n → ℕ) (b : (∀ v, Fin (d v + 1)) → ℚ) (m M : ℚ)
    (hexp : boxShift lo hi N = ∑ α, C (b α) * tensorBern d α)
    (hm : ∀ α, m ≤ b α) (hM : ∀ α, b α ≤ M) :
    ∀ x ∈ (ratBox lo hi).set, (m : ℝ) ≤ aeval x N ∧ aeval x N ≤ (M : ℝ)
\end{leansrc}
\par\medskip\noindent{\small \texttt{Qtsl.BoxBern.}\texttt{\bfseries boxShift}}\hfill{\footnotesize \path{lean/qtsl-aux/QtslAux/BoxBern.lean}, line 122}\par\nopagebreak
\begin{leansrc}
noncomputable def boxShift (lo hi : Fin n → ℚ) :
    MvPolynomial (Fin n) ℚ →ₐ[ℚ] MvPolynomial (Fin n) ℚ
\end{leansrc}
\par\medskip\noindent{\small \texttt{Qtsl.BoxBern.}\texttt{\bfseries aeval\_boxShift}}\hfill{\footnotesize \path{lean/qtsl-aux/QtslAux/BoxBern.lean}, line 126}\par\nopagebreak
\begin{leansrc}
theorem aeval_boxShift (lo hi : Fin n → ℚ) (u : Fin n → ℝ)
    (N : MvPolynomial (Fin n) ℚ) :
    aeval u (boxShift lo hi N)
      = aeval (fun v => (lo v : ℝ) + ((hi v : ℝ) - (lo v : ℝ)) * u v) N
\end{leansrc}
\par\medskip\noindent{\small \texttt{Qtsl.BoxBern.}\texttt{\bfseries boxBern\_nonpos}}\hfill{\footnotesize \path{lean/qtsl-aux/QtslAux/BoxBern.lean}, line 195}\par\nopagebreak
\begin{leansrc}
theorem boxBern_nonpos (lo hi : Fin n → ℚ) (N : MvPolynomial (Fin n) ℚ)
    (d : Fin n → ℕ) (b : (∀ v, Fin (d v + 1)) → ℚ)
    (hexp : boxShift lo hi N = ∑ α, C (b α) * tensorBern d α)
    (hcoeff : ∀ α, b α ≤ 0) :
    ∀ x ∈ (ratBox lo hi).set, aeval x N ≤ 0
\end{leansrc}
\par\medskip\noindent{\small \texttt{Qtsl.BoxBern.}\texttt{\bfseries boxBern\_enclosure\_minmax}}\hfill{\footnotesize \path{lean/qtsl-aux/QtslAux/BoxBern.lean}, line 184}\par\nopagebreak
\begin{leansrc}
theorem boxBern_enclosure_minmax (lo hi : Fin n → ℚ)
    (N : MvPolynomial (Fin n) ℚ)
    (d : Fin n → ℕ) (b : (∀ v, Fin (d v + 1)) → ℚ)
    (hexp : boxShift lo hi N = ∑ α, C (b α) * tensorBern d α) :
    ∀ x ∈ (ratBox lo hi).set,
      ((Finset.univ.inf' Finset.univ_nonempty b : ℚ) : ℝ) ≤ aeval x N ∧
        aeval x N ≤ ((Finset.univ.sup' Finset.univ_nonempty b : ℚ) : ℝ)
\end{leansrc}
\par\medskip\noindent{\small \texttt{Qtsl.BoxBern.}\texttt{\bfseries sum\_aeval\_tensorBern}}\hfill{\footnotesize \path{lean/qtsl-aux/QtslAux/BoxBern.lean}, line 79}\par\nopagebreak
\begin{leansrc}
theorem sum_aeval_tensorBern (d : Fin n → ℕ) (u : Fin n → ℝ) :
    ∑ α : ∀ v, Fin (d v + 1), aeval u (tensorBern d α) = 1
\end{leansrc}
\par\medskip\noindent{\small \texttt{Qtsl.Enclosure.}\texttt{\bfseries encCheck}}\hfill{\footnotesize \path{lean/qtsl-aux/QtslAux/Enclosure.lean}, line 93}\par\nopagebreak
\begin{leansrc}
def encCheck (w : EncWitness) : Bool
\end{leansrc}
\par\medskip\noindent{\small \texttt{Qtsl.Enclosure.}\texttt{\bfseries EncWitness}}\hfill{\footnotesize \path{lean/qtsl-aux/QtslAux/Enclosure.lean}, line 82}\par\nopagebreak
\begin{leansrc}
structure EncWitness where …
\end{leansrc}
\par\medskip\noindent{\small \texttt{Qtsl.Enclosure.}\texttt{\bfseries encCheck\_sound}}\hfill{\footnotesize \path{lean/qtsl-aux/QtslAux/Enclosure.lean}, line 380}\par\nopagebreak
\begin{leansrc}
theorem encCheck_sound (w : EncWitness) (h : encCheck w = true) :
    ∀ x : ℝ, |x| ≤ (w.r : ℝ) →
      |w.tag.eval x - evalPoly w.coeffs x| ≤ (w.R : ℝ)
\end{leansrc}
\par\medskip\noindent{\small \texttt{Qtsl.Enclosure.}\texttt{\bfseries encCheck\_sound\_bracket}}\hfill{\footnotesize \path{lean/qtsl-aux/QtslAux/Enclosure.lean}, line 426}\par\nopagebreak
\begin{leansrc}
theorem encCheck_sound_bracket (w : EncWitness) (h : encCheck w = true) :
    ∀ x : ℝ, |x| ≤ (w.r : ℝ) →
      evalPoly w.coeffs x - (w.R : ℝ) ≤ w.tag.eval x ∧
      w.tag.eval x ≤ evalPoly w.coeffs x + (w.R : ℝ)
\end{leansrc}
\par\medskip\noindent{\small \texttt{Qtsl.Enclosure.}\texttt{\bfseries encCheck\_point}}\hfill{\footnotesize \path{lean/qtsl-aux/QtslAux/Enclosure.lean}, line 437}\par\nopagebreak
\begin{leansrc}
theorem encCheck_point (w : EncWitness) (h : encCheck w = true) (c : ℚ)
    (hc : |c| ≤ w.r) :
    evalPoly w.coeffs (c : ℝ) - (w.R : ℝ) ≤ w.tag.eval (c : ℝ) ∧
    w.tag.eval (c : ℝ) ≤ evalPoly w.coeffs (c : ℝ) + (w.R : ℝ)
\end{leansrc}
\par\medskip\noindent{\small \texttt{Qtsl.Enclosure.}\texttt{\bfseries linear\_interp\_le}}\hfill{\footnotesize \path{lean/qtsl-aux/QtslAux/Enclosure.lean}, line 445}\par\nopagebreak
\begin{leansrc}
theorem linear_interp_le {A C l h w : ℝ} (hlw : l ≤ w) (hwh : w ≤ h)
    (hl : A + C * l ≤ 0) (hh : A + C * h ≤ 0) : A + C * w ≤ 0
\end{leansrc}
\par\medskip\noindent{\small \texttt{Qtsl.Enclosure.}\texttt{\bfseries integral\_exp\_neg\_sq}}\hfill{\footnotesize \path{lean/qtsl-aux/QtslAux/Enclosure.lean}, line 459}\par\nopagebreak
\begin{leansrc}
theorem integral_exp_neg_sq (z : ℝ) :
    ∫ t in (0:ℝ)..z, Real.exp (-(t ^ 2)) = z * FnTag.erfT.eval (z ^ 2)
\end{leansrc}
\par\medskip\noindent{\small \texttt{PPLC.Qtsl.Behavior.}\texttt{\bfseries qMorseDischarge\_sound}}\hfill{\footnotesize \path{lean/qtsl/Qtsl/Behavioral.lean}, line 369}\par\nopagebreak
\begin{leansrc}
theorem qMorseDischarge_sound (B : Behavior X) (D : Set X)
    (ρ : Vtx → Set X) (cells : ι → Set X) (hmap : ι → Option Vtx)
    (hContain : ∀ i v, hmap i = some v → cells i ⊆ ρ v)
    (hCompact : ∀ v, IsCompact (B.MaxInv (fiberUnion cells hmap v)))
    (hDisj : ∀ v w, v ≠ w →
      Disjoint (B.MaxInv (fiberUnion cells hmap v))
        (B.MaxInv (fiberUnion cells hmap w)))
    (hCell : ∀ x ∈ D, B.ChainRecurrentIn D x → ∃ i, x ∈ B.MaxInv (cells i))
    (hTransit : ∀ i, hmap i = none → B.HasTransitCert (cells i))
    (hOrder : ∀ γ ∈ B.cpl, Set.range γ ⊆ D →
      (∀ v, ¬ Set.range γ ⊆ B.MaxInv (fiberUnion cells hmap v)) →
      ∃ v w, w < v ∧ alphaSet γ ⊆ B.MaxInv (fiberUnion cells hmap v) ∧
        omegaSet γ ⊆ B.MaxInv (fiberUnion cells hmap w)) :
    B.MorseSat D ρ
\end{leansrc}
\par\medskip\noindent{\small \texttt{PPLC.Qtsl.Behavior.}\texttt{\bfseries MorseWitness}}\hfill{\footnotesize \path{lean/qtsl/Qtsl/Behavioral.lean}, line 99}\par\nopagebreak
\begin{leansrc}
structure MorseWitness (B : Behavior X) (D : Set X) [PartialOrder Vtx]
    (ρ : Vtx → Set X) where …
\end{leansrc}
\par\medskip\noindent{\small \texttt{PPLC.Qtsl.Behavior.}\texttt{\bfseries HasTransitCert}}\hfill{\footnotesize \path{lean/qtsl/Qtsl/Behavioral.lean}, line 124}\par\nopagebreak
\begin{leansrc}
def HasTransitCert (B : Behavior X) (C : Set X) : Prop
\end{leansrc}
\par\medskip\noindent{\small \texttt{PPLC.Qtsl.Semiflow.}\texttt{\bfseries FlowDerivOn}}\hfill{\footnotesize \path{lean/qtsl/Qtsl/MeanValue.lean}, line 24}\par\nopagebreak
\begin{leansrc}
def FlowDerivOn (φ : Semiflow X) (V W : X → ℝ) (C : Set X) : Prop
\end{leansrc}
\par\medskip\noindent{\small \texttt{PPLC.Qtsl.Semiflow.}\texttt{\bfseries maxInv\_isCompact}}\hfill{\footnotesize \path{lean/qtsl/Qtsl/Discharge.lean}, line 231}\par\nopagebreak
\begin{leansrc}
theorem maxInv_isCompact [T2Space X] {φ : Semiflow X} {C : Set X}
    (hflow : ∀ t, Continuous (φ.flow t)) (hC : IsCompact C) :
    IsCompact (φ.MaxInv C)
\end{leansrc}
\par\medskip\noindent{\small \texttt{PPLC.Qtsl.Semiflow.}\texttt{\bfseries fiberUnion\_isCompact}}\hfill{\footnotesize \path{lean/qtsl/Qtsl/Discharge.lean}, line 261}\par\nopagebreak
\begin{leansrc}
theorem fiberUnion_isCompact {Vtx ι : Type*} [Finite ι]
    {cells : ι → Set X} {hmap : ι → Option Vtx}
    (hcells : ∀ i, IsCompact (cells i)) (v : Vtx) :
    IsCompact (fiberUnion cells hmap v)
\end{leansrc}
\par\medskip\noindent{\small \texttt{PPLC.Qtsl.Semiflow.HasTransitCert.}\texttt{\bfseries disjoint\_maxInv}}\hfill{\footnotesize \path{lean/qtsl/Qtsl/Discharge.lean}, line 81}\par\nopagebreak
\begin{leansrc}
theorem HasTransitCert.disjoint_maxInv {φ : Semiflow X} {C₁ C₂ : Set X}
    (h : φ.HasTransitCert (C₁ ∩ C₂)) :
    Disjoint (φ.MaxInv C₁) (φ.MaxInv C₂)
\end{leansrc}
\par\medskip\noindent{\small \texttt{PPLC.Qtsl.Behavior.}\texttt{\bfseries crCover\_assemble}}\hfill{\footnotesize \path{lean/qtsl/Qtsl/Behavioral.lean}, line 352}\par\nopagebreak
\begin{leansrc}
theorem crCover_assemble {B : Behavior X} {D : Set X}
    {cells : ι → Set X} {hmap : ι → Option Vtx}
    (hCell : ∀ x ∈ D, B.ChainRecurrentIn D x → ∃ i, x ∈ B.MaxInv (cells i))
    (hTransit : ∀ i, hmap i = none → B.HasTransitCert (cells i)) :
    ∀ x ∈ D, B.ChainRecurrentIn D x →
      ∃ v, x ∈ B.MaxInv (fiberUnion cells hmap v)
\end{leansrc}
\par\medskip\noindent{\small \texttt{PPLC.Qtsl.Behavior.HasTransitCert.}\texttt{\bfseries maxInv\_empty}}\hfill{\footnotesize \path{lean/qtsl/Qtsl/Behavioral.lean}, line 144}\par\nopagebreak
\begin{leansrc}
theorem HasTransitCert.maxInv_empty {B : Behavior X} {C : Set X}
    (h : B.HasTransitCert C) : B.MaxInv C = ∅
\end{leansrc}
\par\medskip\noindent{\small \texttt{PPLC.Qtsl.Semiflow.}\texttt{\bfseries qMorseDischarge\_sound\_orderAssembled}}\hfill{\footnotesize \path{lean/qtsl/Qtsl/OrderAssembly.lean}, line 250}\par\nopagebreak
\begin{leansrc}
theorem qMorseDischarge_sound_orderAssembled (φ : Semiflow X) (D : Set X)
    (ρ : Vtx → Set X) (cells : ι → Set X) (hmap : ι → Option Vtx)
    (hD : IsCompact D)
    (hflow : ∀ t, Continuous (φ.flow t))
    (hcells : ∀ i, IsCompact (cells i))
    (hContain : ∀ i v, hmap i = some v → cells i ⊆ ρ v)
    (hOverlap : ∀ v w, v ≠ w →
      φ.HasTransitCert
        (fiberUnion cells hmap v ∩ fiberUnion cells hmap w))
    (hCell : ∀ x ∈ D, ChainRecurrentIn φ D x → ∃ i, x ∈ φ.MaxInv (cells i))
    (hTransitCert : ∀ i, hmap i = none → φ.HasTransitCert (cells i))
    (hConv : ∀ i j k v, i ≤ j → j ≤ k →
      hmap i = some v → hmap k = some v → hmap j = some v)
    (hMono : ∀ i j v w, i ≤ j →
      hmap i = some v → hmap j = some w → v ≤ w)
    (hInterval : ∀ γ : ℝ → X, φ.CompleteOrbit γ → Set.range γ ⊆ D →
      ∃ bα bω : ι, bω ≤ bα ∧
        alphaSet γ ⊆ cells bα ∧ omegaSet γ ⊆ cells bω ∧
        ∀ t : ℝ, ∃ i, bω ≤ i ∧ i ≤ bα ∧ γ t ∈ cells i) :
    MorseSat φ D ρ
\end{leansrc}
\par\medskip\noindent{\small \texttt{PPLC.Qtsl.Semiflow.}\texttt{\bfseries qMorseDischarge\_sound\_eventualMember}}\hfill{\footnotesize \path{lean/qtsl/Qtsl/EventualMember.lean}, line 274}\par\nopagebreak
\begin{leansrc}
theorem qMorseDischarge_sound_eventualMember (φ : Semiflow X) (D : Set X)
    (ρ : Vtx → Set X) (cells : ι → Set X) (hmap : ι → Option Vtx)
    (blocks : β → Set X) (emb : ι → β) (cover : ι → β) (bD : β)
    (hD : IsCompact D)
    (hflow : ∀ t, Continuous (φ.flow t))
    (hcells : ∀ i, IsCompact (cells i))
    (hContain : ∀ i v, hmap i = some v → cells i ⊆ ρ v)
    (hOverlap : ∀ v w, v ≠ w →
      φ.HasTransitCert
        (fiberUnion cells hmap v ∩ fiberUnion cells hmap w))
    (hCell : ∀ x ∈ D, ChainRecurrentIn φ D x → ∃ i, x ∈ φ.MaxInv (cells i))
    (hTransitCert : ∀ i, hmap i = none → φ.HasTransitCert (cells i))
    (hConv : ∀ i j k v, i ≤ j → j ≤ k →
      hmap i = some v → hmap k = some v → hmap j = some v)
    (hMono : ∀ i j v w, i ≤ j →
      hmap i = some v → hmap j = some w → v ≤ w)
    (hbD : D ⊆ blocks bD)
    (hJoin : ∀ b : β, ∀ x ∈ blocks b,
      ∃ i : ι, x ∈ blocks (emb i) ∧ blocks (emb i) ⊆ blocks b)
    (hcover : ∀ i : ι, blocks (cover i) ⊂ blocks (emb i))
    (hcellsSup : ∀ i : ι, blocks (emb i) \ blocks (cover i) ⊆ cells i)
    (horder : ∀ i j : ι, blocks (emb i) ⊆ blocks (emb j) → i ≤ j)
    (hclosed : ∀ i : ι, IsClosed (blocks (emb i)))
    (hInv : ∀ b : β, φ.Invariant (blocks b))
    (hMeet : ∀ a b : β, ∃ c : β, blocks c ⊆ blocks a ∩ blocks b ∧
      ∀ γ : ℝ → X, φ.CompleteOrbit γ → Set.range γ ⊆ D →
        ∀ s t : ℝ, (∀ u, s ≤ u → u ≤ t → γ u ∈ blocks a ∩ blocks b) →
          ∀ u, s ≤ u → u ≤ t → γ u ∈ blocks c)
    (hAbsorbOmega : ∀ γ : ℝ → X, φ.CompleteOrbit γ → Set.range γ ⊆ D →
      ∀ b : β, (omegaSet γ ∩ blocks b).Nonempty →
        ∃ T : ℝ, ∀ t, T ≤ t → γ t ∈ blocks b)
    (hAbsorbAlpha : ∀ γ : ℝ → X, φ.CompleteOrbit γ → Set.range γ ⊆ D →
      ∀ b : β, (alphaSet γ ∩ blocks b).Nonempty →
        ∃ T : ℝ, ∀ t, t ≤ T → γ t ∈ blocks b) :
    MorseSat φ D ρ
\end{leansrc}
\par\medskip\noindent{\small \texttt{PPLC.Qtsl.Semiflow.}\texttt{\bfseries qMorseDischarge\_sound\_cellCover}}\hfill{\footnotesize \path{lean/qtsl/Qtsl/Crossing.lean}, line 120}\par\nopagebreak
\begin{leansrc}
theorem qMorseDischarge_sound_cellCover (φ : Semiflow X) (D : Set X)
    (ρ : Vtx → Set X) (cells : ι → Set X) (hmap : ι → Option Vtx)
    (hflow : ∀ t, Continuous (φ.flow t))
    (htraj : ∀ x, Continuous fun t : ℝ≥0 => φ.flow t x)
    (hcells : ∀ i, IsCompact (cells i))
    (hContain : ∀ i v, hmap i = some v → cells i ⊆ ρ v)
    (hCross : ∀ v w, v ≠ w → ∀ i j, hmap i = some v → hmap j = some w →
      φ.HasCrossingCert (cells i ∩ cells j))
    (hCell : ∀ x ∈ D, ChainRecurrentIn φ D x → ∃ i, x ∈ φ.MaxInv (cells i))
    (hTransitCert : ∀ i, hmap i = none → φ.HasTransitCert (cells i))
    (hOrder : ∀ γ : ℝ → X, φ.CompleteOrbit γ → Set.range γ ⊆ D →
      (∀ v, ¬ Set.range γ ⊆ φ.MaxInv (fiberUnion cells hmap v)) →
      ∃ v w, w < v ∧ alphaSet γ ⊆ φ.MaxInv (fiberUnion cells hmap v) ∧
        omegaSet γ ⊆ φ.MaxInv (fiberUnion cells hmap w)) :
    MorseSat φ D ρ
\end{leansrc}
\par\medskip\noindent{\small \texttt{PPLC.Qtsl.Behavior.}\texttt{\bfseries maxInv\_subset}}\hfill{\footnotesize \path{lean/qtsl/Qtsl/Behavioral.lean}, line 51}\par\nopagebreak
\begin{leansrc}
theorem maxInv_subset (B : Behavior X) (C : Set X) : B.MaxInv C ⊆ C
\end{leansrc}
\par\medskip\noindent{\small \texttt{PPLC.Qtsl.Behavior.}\texttt{\bfseries maxInv\_isInvariant}}\hfill{\footnotesize \path{lean/qtsl/Qtsl/Behavioral.lean}, line 69}\par\nopagebreak
\begin{leansrc}
theorem maxInv_isInvariant (B : Behavior X) (C : Set X) :
    B.IsInvariantIn (B.MaxInv C)
\end{leansrc}
\par\medskip\noindent{\small \texttt{PPLC.Qtsl.Semiflow.}\texttt{\bfseries maxInv\_invariant}}\hfill{\footnotesize \path{lean/qtsl/Qtsl/QtslSoundness.lean}, line 63}\par\nopagebreak
\begin{leansrc}
theorem maxInv_invariant (φ : Semiflow X) (C : Set X) :
    φ.Invariant (φ.MaxInv C)
\end{leansrc}
\par\medskip\noindent{\small \texttt{PPLC.Qtsl.Semiflow.}\texttt{\bfseries fiberUnion\_subset}}\hfill{\footnotesize \path{lean/qtsl/Qtsl/QtslSoundness.lean}, line 195}\par\nopagebreak
\begin{leansrc}
theorem fiberUnion_subset {cells : ι → Set X} {hmap : ι → Option Vtx}
    {ρ : Vtx → Set X} (hContain : ∀ i v, hmap i = some v → cells i ⊆ ρ v)
    (v : Vtx) : fiberUnion cells hmap v ⊆ ρ v
\end{leansrc}
\par\medskip\noindent{\small \texttt{PPLC.Qtsl.Semiflow.}\texttt{\bfseries maxInv\_forward\_confined}}\hfill{\footnotesize \path{lean/qtsl/Qtsl/Discharge.lean}, line 38}\par\nopagebreak
\begin{leansrc}
theorem maxInv_forward_confined {φ : Semiflow X} {C : Set X} {x : X}
    (hx : x ∈ φ.MaxInv C) : ∀ t : ℝ≥0, φ.flow t x ∈ C
\end{leansrc}
\par\medskip\noindent{\small \texttt{PPLC.Qtsl.Behavior.HasTransitCert.}\texttt{\bfseries not\_forward\_confined}}\hfill{\footnotesize \path{lean/qtsl/Qtsl/Behavioral.lean}, line 129}\par\nopagebreak
\begin{leansrc}
theorem HasTransitCert.not_forward_confined {B : Behavior X} {C : Set X}
    (h : B.HasTransitCert C) {γ : ℝ≥0 → X} (hγ : γ ∈ B.fwd) :
    ¬ ∀ t : ℝ≥0, γ t ∈ C
\end{leansrc}
\par\medskip\noindent{\small \texttt{PPLC.Qtsl.Semiflow.}\texttt{\bfseries transit\_maxInv\_empty}}\hfill{\footnotesize \path{lean/qtsl/Qtsl/QtslSoundness.lean}, line 114}\par\nopagebreak
\begin{leansrc}
theorem transit_maxInv_empty {φ : Semiflow X} {C : Set X} {V : X → ℝ}
    {ε : ℝ} (hε : 0 < ε) (hBdd : BddBelow (V '' C))
    (hDec : ∀ x ∈ C, ∀ t : ℝ≥0,
      (∀ s : ℝ≥0, s ≤ t → φ.flow s x ∈ C) → V (φ.flow t x) ≤ V x - ε * t) :
    φ.MaxInv C = ∅
\end{leansrc}
\par\medskip\noindent{\small \texttt{PPLC.Qtsl.Behavior.}\texttt{\bfseries maxInv\_mono}}\hfill{\footnotesize \path{lean/qtsl/Qtsl/Behavioral.lean}, line 55}\par\nopagebreak
\begin{leansrc}
theorem maxInv_mono {B : Behavior X} {C D : Set X} (h : C ⊆ D) :
    B.MaxInv C ⊆ B.MaxInv D
\end{leansrc}
\par\medskip\noindent{\small \texttt{PPLC.Qtsl.Semiflow.}\texttt{\bfseries subset\_fiberUnion}}\hfill{\footnotesize \path{lean/qtsl/Qtsl/Discharge.lean}, line 45}\par\nopagebreak
\begin{leansrc}
theorem subset_fiberUnion {Vtx ι : Type*} {cells : ι → Set X}
    {hmap : ι → Option Vtx} {i : ι} {v : Vtx} (hi : hmap i = some v) :
    cells i ⊆ fiberUnion cells hmap v
\end{leansrc}
\par\medskip\noindent{\small \texttt{PPLC.Qtsl.Semiflow.}\texttt{\bfseries qMorseDischarge\_sound\_modelClass}}\hfill{\footnotesize \path{lean/qtsl/Qtsl/Discharge.lean}, line 297}\par\nopagebreak
\begin{leansrc}
theorem qMorseDischarge_sound_modelClass (φ : Semiflow X) (D : Set X)
    (ρ : Vtx → Set X) (cells : ι → Set X) (hmap : ι → Option Vtx)
    (hflow : ∀ t, Continuous (φ.flow t))
    (hcells : ∀ i, IsCompact (cells i))
    (hContain : ∀ i v, hmap i = some v → cells i ⊆ ρ v)
    (hOverlap : ∀ v w, v ≠ w →
      φ.HasTransitCert
        (fiberUnion cells hmap v ∩ fiberUnion cells hmap w))
    (hCell : ∀ x ∈ D, ChainRecurrentIn φ D x → ∃ i, x ∈ φ.MaxInv (cells i))
    (hTransitCert : ∀ i, hmap i = none → φ.HasTransitCert (cells i))
    (hOrder : ∀ γ : ℝ → X, φ.CompleteOrbit γ → Set.range γ ⊆ D →
      (∀ v, ¬ Set.range γ ⊆ φ.MaxInv (fiberUnion cells hmap v)) →
      ∃ v w, w < v ∧ alphaSet γ ⊆ φ.MaxInv (fiberUnion cells hmap v) ∧
        omegaSet γ ⊆ φ.MaxInv (fiberUnion cells hmap w)) :
    MorseSat φ D ρ
\end{leansrc}
\par\medskip\noindent{\small \texttt{PPLC.Qtsl.Behavior.}\texttt{\bfseries MorseSat}}\hfill{\footnotesize \path{lean/qtsl/Qtsl/Behavioral.lean}, line 118}\par\nopagebreak
\begin{leansrc}
def MorseSat (B : Behavior X) (D : Set X) [PartialOrder Vtx]
    (ρ : Vtx → Set X) : Prop
\end{leansrc}
\par\medskip\noindent{\small \texttt{PPLC.Qtsl.Behavior.}\texttt{\bfseries EpsChainIn}}\hfill{\footnotesize \path{lean/qtsl/Qtsl/Behavioral.lean}, line 80}\par\nopagebreak
\begin{leansrc}
inductive EpsChainIn (B : Behavior X) (D : Set X) (ε : ℝ) (T : ℝ≥0) :
    X → X → Prop
  | step {x y : X} {γ : ℝ≥0 → X} {t : ℝ≥0} (hγ : γ ∈ B.fwd) (hγ0 : γ 0 = x)
      (hx : x ∈ D) (hy : y ∈ D) (ht : T ≤ t) (hjump : dist (γ t) y < ε) :
      EpsChainIn B D ε T x y
  | trans {x y z : X} (h1 : EpsChainIn B D ε T x y)
      (h2 : EpsChainIn B D ε T y z) : EpsChainIn B D ε T x z

def ChainRecurrentIn (B : Behavior X) (D : Set X) (x : X) : Prop
\end{leansrc}
\par\medskip\noindent{\small \texttt{PPLC.Qtsl.Behavior.}\texttt{\bfseries ChainRecurrentIn}}\hfill{\footnotesize \path{lean/qtsl/Qtsl/Behavioral.lean}, line 88}\par\nopagebreak
\begin{leansrc}
def ChainRecurrentIn (B : Behavior X) (D : Set X) (x : X) : Prop
\end{leansrc}
\par\medskip\noindent{\small \texttt{PPLC.Qtsl.Semiflow.}\texttt{\bfseries omegaSet}}\hfill{\footnotesize \path{lean/qtsl/Qtsl/QtslSoundness.lean}, line 142}\par\nopagebreak
\begin{leansrc}
def omegaSet [TopologicalSpace X] (γ : ℝ → X) : Set X
\end{leansrc}
\par\medskip\noindent{\small \texttt{PPLC.Qtsl.Semiflow.}\texttt{\bfseries alphaSet}}\hfill{\footnotesize \path{lean/qtsl/Qtsl/QtslSoundness.lean}, line 145}\par\nopagebreak
\begin{leansrc}
def alphaSet [TopologicalSpace X] (γ : ℝ → X) : Set X
\end{leansrc}
\par\medskip\noindent{\small \texttt{PPLC.Qtsl.Semiflow.}\texttt{\bfseries orderClause\_assemble}}\hfill{\footnotesize \path{lean/qtsl/Qtsl/OrderAssembly.lean}, line 183}\par\nopagebreak
\begin{leansrc}
theorem orderClause_assemble (φ : Semiflow X) (D : Set X)
    (hD : IsCompact D) (hflow : ∀ t, Continuous (φ.flow t))
    (cells : ι → Set X) (hmap : ι → Option Vtx)
    (hConv : ∀ i j k v, i ≤ j → j ≤ k →
      hmap i = some v → hmap k = some v → hmap j = some v)
    (hMono : ∀ i j v w, i ≤ j →
      hmap i = some v → hmap j = some w → v ≤ w)
    (hTransit : ∀ i, hmap i = none → φ.HasTransitCert (cells i))
    (hInterval : ∀ γ : ℝ → X, φ.CompleteOrbit γ → Set.range γ ⊆ D →
      ∃ bα bω : ι, bω ≤ bα ∧
        alphaSet γ ⊆ cells bα ∧ omegaSet γ ⊆ cells bω ∧
        ∀ t : ℝ, ∃ i, bω ≤ i ∧ i ≤ bα ∧ γ t ∈ cells i) :
    ∀ γ : ℝ → X, φ.CompleteOrbit γ → Set.range γ ⊆ D →
      (∀ v, ¬ Set.range γ ⊆ φ.MaxInv (fiberUnion cells hmap v)) →
      ∃ v w, w < v ∧
        alphaSet γ ⊆ φ.MaxInv (fiberUnion cells hmap v) ∧
        omegaSet γ ⊆ φ.MaxInv (fiberUnion cells hmap w)
\end{leansrc}
\par\medskip\noindent{\small \texttt{PPLC.Qtsl.Semiflow.}\texttt{\bfseries eventualMember\_interval}}\hfill{\footnotesize \path{lean/qtsl/Qtsl/EventualMember.lean}, line 145}\par\nopagebreak
\begin{leansrc}
theorem eventualMember_interval (φ : Semiflow X) (D : Set X)
    (blocks : β → Set X) (emb : ι → β) (cover : ι → β) (bD : β)
    (cells : ι → Set X)
    (hbD : D ⊆ blocks bD)
    (hJoin : ∀ b : β, ∀ x ∈ blocks b,
      ∃ i : ι, x ∈ blocks (emb i) ∧ blocks (emb i) ⊆ blocks b)
    (hcover : ∀ i : ι, blocks (cover i) ⊂ blocks (emb i))
    (hcellsSup : ∀ i : ι, blocks (emb i) \ blocks (cover i) ⊆ cells i)
    (horder : ∀ i j : ι, blocks (emb i) ⊆ blocks (emb j) → i ≤ j)
    (hclosed : ∀ i : ι, IsClosed (blocks (emb i)))
    (hInv : ∀ b : β, φ.Invariant (blocks b))
    (hMeet : ∀ a b : β, ∃ c : β, blocks c ⊆ blocks a ∩ blocks b ∧
      ∀ γ : ℝ → X, φ.CompleteOrbit γ → Set.range γ ⊆ D →
        ∀ s t : ℝ, (∀ u, s ≤ u → u ≤ t → γ u ∈ blocks a ∩ blocks b) →
          ∀ u, s ≤ u → u ≤ t → γ u ∈ blocks c)
    (hAbsorbOmega : ∀ γ : ℝ → X, φ.CompleteOrbit γ → Set.range γ ⊆ D →
      ∀ b : β, (omegaSet γ ∩ blocks b).Nonempty →
        ∃ T : ℝ, ∀ t, T ≤ t → γ t ∈ blocks b)
    (hAbsorbAlpha : ∀ γ : ℝ → X, φ.CompleteOrbit γ → Set.range γ ⊆ D →
      ∀ b : β, (alphaSet γ ∩ blocks b).Nonempty →
        ∃ T : ℝ, ∀ t, t ≤ T → γ t ∈ blocks b) :
    ∀ γ : ℝ → X, φ.CompleteOrbit γ → Set.range γ ⊆ D →
      ∃ bα bω : ι, bω ≤ bα ∧
        alphaSet γ ⊆ cells bα ∧ omegaSet γ ⊆ cells bω ∧
        ∀ t : ℝ, ∃ i, bω ≤ i ∧ i ≤ bα ∧ γ t ∈ cells i
\end{leansrc}
\par\medskip\noindent{\small \texttt{PPLC.Qtsl.Semiflow.}\texttt{\bfseries prod}}\hfill{\footnotesize \path{lean/qtsl/Qtsl/Product.lean}, line 21}\par\nopagebreak
\begin{leansrc}
def prod (φ₁ : Semiflow X₁) (φ₂ : Semiflow X₂) : Semiflow (X₁ × X₂) where …
\end{leansrc}
\par\medskip\noindent{\small \texttt{PPLC.Qtsl.Semiflow.MorseWitness.}\texttt{\bfseries prod}}\hfill{\footnotesize \path{lean/qtsl/Qtsl/Product.lean}, line 154}\par\nopagebreak
\begin{leansrc}
def MorseWitness.prod (W₁ : MorseWitness φ₁ D₁ ρ₁)
    (W₂ : MorseWitness φ₂ D₂ ρ₂) :
    MorseWitness (φ₁.prod φ₂) (D₁ ×ˢ D₂)
      (fun p : Vtx₁ × Vtx₂ => ρ₁ p.1 ×ˢ ρ₂ p.2) where …
\end{leansrc}
\par\medskip\noindent{\small \texttt{PPLC.Qtsl.Semiflow.}\texttt{\bfseries qMorseProduct\_sound}}\hfill{\footnotesize \path{lean/qtsl/Qtsl/Product.lean}, line 235}\par\nopagebreak
\begin{leansrc}
theorem qMorseProduct_sound
    (h₁ : MorseSat φ₁ D₁ ρ₁) (h₂ : MorseSat φ₂ D₂ ρ₂) :
    MorseSat (φ₁.prod φ₂) (D₁ ×ˢ D₂)
      (fun p : Vtx₁ × Vtx₂ => ρ₁ p.1 ×ˢ ρ₂ p.2)
\end{leansrc}
\par\medskip\noindent{\small \texttt{PPLC.Qtsl.Semiflow.}\texttt{\bfseries prod\_flow}}\hfill{\footnotesize \path{lean/qtsl/Qtsl/Product.lean}, line 26}\par\nopagebreak
\begin{leansrc}
@[simp] theorem prod_flow (φ₁ : Semiflow X₁) (φ₂ : Semiflow X₂)
    (t : ℝ≥0) (p : X₁ × X₂) :
    (φ₁.prod φ₂).flow t p = (φ₁.flow t p.1, φ₂.flow t p.2)
\end{leansrc}
\par\medskip\noindent{\small \texttt{PPLC.Qtsl.Semiflow.}\texttt{\bfseries completeOrbit\_prod\_iff}}\hfill{\footnotesize \path{lean/qtsl/Qtsl/Product.lean}, line 39}\par\nopagebreak
\begin{leansrc}
theorem completeOrbit_prod_iff {γ : ℝ → X₁ × X₂} :
    (φ₁.prod φ₂).CompleteOrbit γ ↔
      φ₁.CompleteOrbit (fun s => (γ s).1) ∧
      φ₂.CompleteOrbit (fun s => (γ s).2)
\end{leansrc}
\par\medskip\noindent{\small \texttt{PPLC.Qtsl.Semiflow.CompleteOrbit.}\texttt{\bfseries prodMk}}\hfill{\footnotesize \path{lean/qtsl/Qtsl/Product.lean}, line 32}\par\nopagebreak
\begin{leansrc}
theorem CompleteOrbit.prodMk {γ₁ : ℝ → X₁} {γ₂ : ℝ → X₂}
    (h₁ : φ₁.CompleteOrbit γ₁) (h₂ : φ₂.CompleteOrbit γ₂) :
    (φ₁.prod φ₂).CompleteOrbit fun s => (γ₁ s, γ₂ s)
\end{leansrc}
\par\medskip\noindent{\small \texttt{PPLC.Qtsl.Semiflow.}\texttt{\bfseries omegaSet\_prod\_subset}}\hfill{\footnotesize \path{lean/qtsl/Qtsl/Product.lean}, line 55}\par\nopagebreak
\begin{leansrc}
theorem omegaSet_prod_subset (γ : ℝ → X₁ × X₂) :
    omegaSet γ ⊆
      (omegaSet fun s => (γ s).1) ×ˢ (omegaSet fun s => (γ s).2)
\end{leansrc}
\par\medskip\noindent{\small \texttt{PPLC.Qtsl.Semiflow.}\texttt{\bfseries alphaSet\_prod\_subset}}\hfill{\footnotesize \path{lean/qtsl/Qtsl/Product.lean}, line 72}\par\nopagebreak
\begin{leansrc}
theorem alphaSet_prod_subset (γ : ℝ → X₁ × X₂) :
    alphaSet γ ⊆
      (alphaSet fun s => (γ s).1) ×ˢ (alphaSet fun s => (γ s).2)
\end{leansrc}
\par\medskip\noindent{\small \texttt{PPLC.Qtsl.Semiflow.EpsChainIn.}\texttt{\bfseries fst}}\hfill{\footnotesize \path{lean/qtsl/Qtsl/Product.lean}, line 111}\par\nopagebreak
\begin{leansrc}
theorem EpsChainIn.fst {D₁ : Set X₁} {D₂ : Set X₂} {ε : ℝ} {T : ℝ≥0}
    {x y : X₁ × X₂}
    (h : EpsChainIn (φ₁.prod φ₂) (D₁ ×ˢ D₂) ε T x y) :
    EpsChainIn φ₁ D₁ ε T x.1 y.1
\end{leansrc}
\par\medskip\noindent{\small \texttt{PPLC.Qtsl.Semiflow.ChainRecurrentIn.}\texttt{\bfseries fst}}\hfill{\footnotesize \path{lean/qtsl/Qtsl/Product.lean}, line 133}\par\nopagebreak
\begin{leansrc}
theorem ChainRecurrentIn.fst {D₁ : Set X₁} {D₂ : Set X₂} {x : X₁ × X₂}
    (h : ChainRecurrentIn (φ₁.prod φ₂) (D₁ ×ˢ D₂) x) :
    ChainRecurrentIn φ₁ D₁ x.1
\end{leansrc}
\par\medskip\noindent{\small \texttt{PPLC.Qtsl.Semiflow.}\texttt{\bfseries alphaSet\_subset\_closure\_range}}\hfill{\footnotesize \path{lean/qtsl/Qtsl/Product.lean}, line 100}\par\nopagebreak
\begin{leansrc}
theorem alphaSet_subset_closure_range (γ : ℝ → X) :
    alphaSet γ ⊆ closure (Set.range γ)
\end{leansrc}
\par\medskip\noindent{\small \texttt{PPLC.Qtsl.Semiflow.}\texttt{\bfseries omegaSet\_subset\_closure\_range}}\hfill{\footnotesize \path{lean/qtsl/Qtsl/Product.lean}, line 95}\par\nopagebreak
\begin{leansrc}
theorem omegaSet_subset_closure_range (γ : ℝ → X) :
    omegaSet γ ⊆ closure (Set.range γ)
\end{leansrc}
\par\medskip\noindent{\small \texttt{PPLC.Qtsl.Behavior.}\texttt{\bfseries qBlockInv\_sound}}\hfill{\footnotesize \path{lean/qtsl/Qtsl/Behavioral.lean}, line 323}\par\nopagebreak
\begin{leansrc}
theorem qBlockInv_sound {B : Behavior X} {N X0 : Set X}
    (hInv : B.Invariant N) (hInit : X0 ⊆ N) : B.InvScopeSat N X0
\end{leansrc}
\par\medskip\noindent{\small \texttt{PPLC.Qtsl.Behavior.}\texttt{\bfseries invariant\_of\_boundaryDwell}}\hfill{\footnotesize \path{lean/qtsl/Qtsl/Behavioral.lean}, line 213}\par\nopagebreak
\begin{leansrc}
theorem invariant_of_boundaryDwell {B : Behavior X} {N : Set X}
    (hN : IsClosed N)
    (hcont : ∀ γ ∈ B.fwd, Continuous γ)
    (hdwell : ∀ γ ∈ B.fwd, γ 0 ∈ N → γ 0 ∈ frontier N →
      ∃ δ : ℝ≥0, 0 < δ ∧ ∀ s ≤ δ, γ s ∈ N) :
    B.Invariant N
\end{leansrc}
\par\medskip\noindent{\small \texttt{PPLC.Qtsl.Behavior.HasFacetRows.}\texttt{\bfseries boundaryDwell}}\hfill{\footnotesize \path{lean/qtsl/Qtsl/Behavioral.lean}, line 283}\par\nopagebreak
\begin{leansrc}
theorem HasFacetRows.boundaryDwell {B : Behavior X} {D N : Set X}
    (h : B.HasFacetRows D N) (hND : N ⊆ D) (hD : B.Invariant D)
    (hcont : ∀ γ ∈ B.fwd, Continuous γ) :
    ∀ γ ∈ B.fwd, γ 0 ∈ N → γ 0 ∈ frontier N →
      ∃ δ : ℝ≥0, 0 < δ ∧ ∀ s ≤ δ, γ s ∈ N
\end{leansrc}
\par\medskip\noindent{\small \texttt{PPLC.Qtsl.Behavior.}\texttt{\bfseries invariant\_univ}}\hfill{\footnotesize \path{lean/qtsl/Qtsl/Behavioral.lean}, line 320}\par\nopagebreak
\begin{leansrc}
theorem invariant_univ (B : Behavior X) : B.Invariant Set.univ
\end{leansrc}
\par\medskip\noindent{\small \texttt{PPLC.Qtsl.Behavior.}\texttt{\bfseries HasFacetRows}}\hfill{\footnotesize \path{lean/qtsl/Qtsl/Behavioral.lean}, line 200}\par\nopagebreak
\begin{leansrc}
def HasFacetRows (B : Behavior X) [TopologicalSpace X] (D N : Set X) : Prop
\end{leansrc}
\par\medskip\noindent{\small \texttt{PPLC.Qtsl.Semiflow.}\texttt{\bfseries qBlockInv\_sound\_of\_rows}}\hfill{\footnotesize \path{lean/qtsl/Qtsl/MeanValue.lean}, line 166}\par\nopagebreak
\begin{leansrc}
theorem qBlockInv_sound_of_rows {φ : Semiflow X} {D N X0 : Set X}
    {n : ℕ} {ℓ : Fin n → X → ℝ} {ε : Fin n → ℝ} {F : Fin n → Set X}
    (W : Fin n → X → ℝ)
    (hN : IsClosed N)
    (htraj : ∀ x, Continuous fun t : ℝ≥0 => φ.flow t x)
    (hND : N ⊆ D) (hD : φ.Invariant D)
    (hε : ∀ i, 0 < ε i)
    (hℓ : ∀ i, ContinuousOn (ℓ i) (F i))
    (hW : ∀ i, φ.FlowDerivOn (ℓ i) (W i) (F i))
    (hrow : ∀ i, ∀ y ∈ F i, ε i ≤ W i y)
    (hlocal : ∀ x ∈ frontier N, ∃ (A : Set (Fin n)) (U : Set X), U ∈ 𝓝 x ∧
      (∀ i ∈ A, ℓ i x = 0 ∧ F i ∈ 𝓝 x) ∧
      ∀ y ∈ U, y ∈ D → (∀ i ∈ A, 0 ≤ ℓ i y) → y ∈ N)
    (hInit : X0 ⊆ N) : φ.InvScopeSat N X0
\end{leansrc}
\par\medskip\noindent{\small \texttt{PPLC.Qtsl.Grid.}\texttt{\bfseries qBlockInv\_sound\_of\_grid}}\hfill{\footnotesize \path{lean/qtsl/Qtsl/GridGeometry.lean}, line 381}\par\nopagebreak
\begin{leansrc}
theorem qBlockInv_sound_of_grid (φ : Semiflow (Fin d → ℝ)) (B : Finset G.CellIdx)
    {X0 : Set (Fin d → ℝ)}
    (htraj : ∀ x, Continuous fun t : ℝ≥0 => φ.flow t x)
    (hArena : φ.Invariant G.arena)
    (ε : G.FacetData → ℝ) (F : G.FacetData → Set (Fin d → ℝ))
    (hε : ∀ p, G.IsFacet B p → 0 < ε p)
    (hF : ∀ p, G.IsFacet B p → ∀ x ∈ G.facetBox p, F p ∈ 𝓝 x)
    (hrow : ∀ p, G.IsFacet B p → ∀ x ∈ F p, ∀ t : ℝ≥0,
      (∀ s ≤ t, φ.flow s x ∈ F p) →
      G.obs p x + ε p * (t : ℝ) ≤ G.obs p (φ.flow t x))
    (hInit : X0 ⊆ G.blockRegion B) :
    φ.InvScopeSat (G.blockRegion B) X0
\end{leansrc}
\par\medskip\noindent{\small \texttt{PPLC.Qtsl.Grid.}\texttt{\bfseries qBlockInv\_sound\_of\_pointwiseRows}}\hfill{\footnotesize \path{lean/qtsl/Qtsl/GridGeometry.lean}, line 397}\par\nopagebreak
\begin{leansrc}
theorem qBlockInv_sound_of_pointwiseRows (φ : Semiflow (Fin d → ℝ))
    (B : Finset G.CellIdx) {X0 : Set (Fin d → ℝ)}
    (htraj : ∀ x, Continuous fun t : ℝ≥0 => φ.flow t x)
    (hArena : φ.Invariant G.arena)
    (ε : G.FacetData → ℝ) (F : G.FacetData → Set (Fin d → ℝ))
    (W : G.FacetData → (Fin d → ℝ) → ℝ)
    (hε : ∀ p, G.IsFacet B p → 0 < ε p)
    (hF : ∀ p, G.IsFacet B p → ∀ x ∈ G.facetBox p, F p ∈ 𝓝 x)
    (hV : ∀ p, G.IsFacet B p → ContinuousOn (G.obs p) (F p))
    (hW : ∀ p, G.IsFacet B p → φ.FlowDerivOn (G.obs p) (W p) (F p))
    (hrow : ∀ p, G.IsFacet B p → ∀ y ∈ F p, ε p ≤ W p y)
    (hInit : X0 ⊆ G.blockRegion B) :
    φ.InvScopeSat (G.blockRegion B) X0
\end{leansrc}
\par\medskip\noindent{\small \texttt{PPLC.Qtsl.Semiflow.}\texttt{\bfseries qBlockInv\_sound\_of\_dwell}}\hfill{\footnotesize \path{lean/qtsl/Qtsl/Discharge.lean}, line 168}\par\nopagebreak
\begin{leansrc}
theorem qBlockInv_sound_of_dwell {φ : Semiflow X} {N X0 : Set X}
    (hN : IsClosed N)
    (htraj : ∀ x, Continuous fun t : ℝ≥0 => φ.flow t x)
    (hdwell : ∀ x ∈ N, x ∈ frontier N →
      ∃ δ : ℝ≥0, 0 < δ ∧ ∀ s ≤ δ, φ.flow s x ∈ N)
    (hInit : X0 ⊆ N) : φ.InvScopeSat N X0
\end{leansrc}
\par\medskip\noindent{\small \texttt{PPLC.Qtsl.Behavior.}\texttt{\bfseries qBlockInv\_sound\_of\_facetRows}}\hfill{\footnotesize \path{lean/qtsl/Qtsl/Behavioral.lean}, line 310}\par\nopagebreak
\begin{leansrc}
theorem qBlockInv_sound_of_facetRows {B : Behavior X} {D N X0 : Set X}
    (hN : IsClosed N)
    (hcont : ∀ γ ∈ B.fwd, Continuous γ)
    (hND : N ⊆ D) (hD : B.Invariant D) (hrows : B.HasFacetRows D N)
    (hInit : X0 ⊆ N) : B.InvScopeSat N X0
\end{leansrc}
\par\medskip\noindent{\small \texttt{PPLC.Qtsl.Semiflow.}\texttt{\bfseries invariant\_of\_invScopeSat}}\hfill{\footnotesize \path{lean/qtsl/Qtsl/GridGeometry.lean}, line 422}\par\nopagebreak
\begin{leansrc}
theorem invariant_of_invScopeSat {φ : Semiflow X} {D X0 : Set X}
    (h : φ.InvScopeSat D X0) : φ.Invariant D
\end{leansrc}
\par\medskip\noindent{\small \texttt{PPLC.Qtsl.Semiflow.}\texttt{\bfseries integrated\_ge\_of\_flowDeriv\_ge}}\hfill{\footnotesize \path{lean/qtsl/Qtsl/MeanValue.lean}, line 105}\par\nopagebreak
\begin{leansrc}
theorem integrated_ge_of_flowDeriv_ge {φ : Semiflow X} {C : Set X}
    {V W : X → ℝ} {c : ℝ}
    (htraj : ∀ x, Continuous fun t : ℝ≥0 => φ.flow t x)
    (hV : ContinuousOn V C)
    (hW : φ.FlowDerivOn V W C)
    (hrow : ∀ y ∈ C, c ≤ W y)
    {x : X} {t : ℝ≥0} (hconf : ∀ s ≤ t, φ.flow s x ∈ C) :
    V x + c * t ≤ V (φ.flow t x)
\end{leansrc}
\par\medskip\noindent{\small \texttt{PPLC.Qtsl.Grid.}\texttt{\bfseries wedge\_containment}}\hfill{\footnotesize \path{lean/qtsl/Qtsl/GridGeometry.lean}, line 106}\par\nopagebreak
\begin{leansrc}
theorem wedge_containment (B : Finset G.CellIdx) (x : Fin d → ℝ)
    (hx : x ∈ G.blockRegion B) :
    ∃ U ∈ 𝓝 x, ∀ y ∈ U, y ∈ G.arena →
      (∀ p : G.FacetData, G.IsFacet B p → x ∈ G.facetBox p → 0 ≤ G.obs p y) →
      y ∈ G.blockRegion B
\end{leansrc}
\par\medskip\noindent{\small \texttt{PPLC.Qtsl.Grid.}\texttt{\bfseries qBlockInv\_sound\_of\_grid}}\hfill{\footnotesize \path{lean/qtsl/Qtsl/GridGeometry.lean}, line 381}\par\nopagebreak
\begin{leansrc}
theorem qBlockInv_sound_of_grid (φ : Semiflow (Fin d → ℝ)) (B : Finset G.CellIdx)
    {X0 : Set (Fin d → ℝ)}
    (htraj : ∀ x, Continuous fun t : ℝ≥0 => φ.flow t x)
    (hArena : φ.Invariant G.arena)
    (ε : G.FacetData → ℝ) (F : G.FacetData → Set (Fin d → ℝ))
    (hε : ∀ p, G.IsFacet B p → 0 < ε p)
    (hF : ∀ p, G.IsFacet B p → ∀ x ∈ G.facetBox p, F p ∈ 𝓝 x)
    (hrow : ∀ p, G.IsFacet B p → ∀ x ∈ F p, ∀ t : ℝ≥0,
      (∀ s ≤ t, φ.flow s x ∈ F p) →
      G.obs p x + ε p * (t : ℝ) ≤ G.obs p (φ.flow t x))
    (hInit : X0 ⊆ G.blockRegion B) :
    φ.InvScopeSat (G.blockRegion B) X0
\end{leansrc}
\par\medskip\noindent{\small \texttt{PPLC.Qtsl.Grid.}\texttt{\bfseries qBlockInv\_sound\_of\_pointwiseRows}}\hfill{\footnotesize \path{lean/qtsl/Qtsl/GridGeometry.lean}, line 397}\par\nopagebreak
\begin{leansrc}
theorem qBlockInv_sound_of_pointwiseRows (φ : Semiflow (Fin d → ℝ))
    (B : Finset G.CellIdx) {X0 : Set (Fin d → ℝ)}
    (htraj : ∀ x, Continuous fun t : ℝ≥0 => φ.flow t x)
    (hArena : φ.Invariant G.arena)
    (ε : G.FacetData → ℝ) (F : G.FacetData → Set (Fin d → ℝ))
    (W : G.FacetData → (Fin d → ℝ) → ℝ)
    (hε : ∀ p, G.IsFacet B p → 0 < ε p)
    (hF : ∀ p, G.IsFacet B p → ∀ x ∈ G.facetBox p, F p ∈ 𝓝 x)
    (hV : ∀ p, G.IsFacet B p → ContinuousOn (G.obs p) (F p))
    (hW : ∀ p, G.IsFacet B p → φ.FlowDerivOn (G.obs p) (W p) (F p))
    (hrow : ∀ p, G.IsFacet B p → ∀ y ∈ F p, ε p ≤ W p y)
    (hInit : X0 ⊆ G.blockRegion B) :
    φ.InvScopeSat (G.blockRegion B) X0
\end{leansrc}
\par\medskip\noindent{\small \texttt{PPLC.Qtsl.}\texttt{\bfseries qLatticeInv\_exact}}\hfill{\footnotesize \path{lean/qtsl/Qtsl/LatticeInv.lean}, line 28}\par\nopagebreak
\begin{leansrc}
theorem qLatticeInv_exact
    (hdc : IsLowerSet C)
    (hJoin : ∀ b : β, ∀ x ∈ blocks b,
      ∃ i : ι, x ∈ blocks (emb i) ∧ blocks (emb i) ⊆ blocks b)
    (hcover : ∀ i : ι, blocks (cover i) ⊂ blocks (emb i))
    (hcellsSup : ∀ i : ι, blocks (emb i) \ blocks (cover i) ⊆ cells i)
    (horder : ∀ i j : ι, blocks (emb i) ⊆ blocks (emb j) → i ≤ j)
    (hcellSub : ∀ i : ι, cells i ⊆ blocks (emb i)) :
    (⋃ i ∈ C, cells i) = ⋃ i ∈ C, blocks (emb i)
\end{leansrc}
\par\medskip\noindent{\small \texttt{PPLC.Qtsl.Semiflow.}\texttt{\bfseries exists\_cellMem\_subset}}\hfill{\footnotesize \path{lean/qtsl/Qtsl/EventualMember.lean}, line 112}\par\nopagebreak
\begin{leansrc}
theorem exists_cellMem_subset
    {blocks : β → Set X} {emb : ι → β} {cover : ι → β} {cells : ι → Set X}
    (hJoin : ∀ b : β, ∀ x ∈ blocks b,
      ∃ i : ι, x ∈ blocks (emb i) ∧ blocks (emb i) ⊆ blocks b)
    (hcover : ∀ i : ι, blocks (cover i) ⊂ blocks (emb i))
    (hcellsSup : ∀ i : ι, blocks (emb i) \ blocks (cover i) ⊆ cells i)
    {x : X} (i : ι) (hmem : x ∈ blocks (emb i)) :
    ∃ j : ι, blocks (emb j) ⊆ blocks (emb i) ∧
      x ∈ blocks (emb j) ∧ x ∈ cells j
\end{leansrc}
\par\medskip\noindent{\small \texttt{PPLC.Qtsl.Semiflow.}\texttt{\bfseries qLatticeInv\_sound}}\hfill{\footnotesize \path{lean/qtsl/Qtsl/LatticeInv.lean}, line 65}\par\nopagebreak
\begin{leansrc}
theorem Semiflow.qLatticeInv_sound {φ : Semiflow X} {X0 : Set X}
    (hdc : IsLowerSet C)
    (hJoin : ∀ b : β, ∀ x ∈ blocks b,
      ∃ i : ι, x ∈ blocks (emb i) ∧ blocks (emb i) ⊆ blocks b)
    (hcover : ∀ i : ι, blocks (cover i) ⊂ blocks (emb i))
    (hcellsSup : ∀ i : ι, blocks (emb i) \ blocks (cover i) ⊆ cells i)
    (horder : ∀ i j : ι, blocks (emb i) ⊆ blocks (emb j) → i ≤ j)
    (hcellSub : ∀ i : ι, cells i ⊆ blocks (emb i))
    (hMember : φ.InvScopeSat (⋃ i ∈ C, blocks (emb i)) X0) :
    φ.InvScopeSat (⋃ i ∈ C, cells i) X0
\end{leansrc}
\par\medskip\noindent{\small \texttt{PPLC.Qtsl.Behavior.}\texttt{\bfseries qLatticeInv\_sound}}\hfill{\footnotesize \path{lean/qtsl/Qtsl/LatticeInv.lean}, line 93}\par\nopagebreak
\begin{leansrc}
theorem Behavior.qLatticeInv_sound {B : Behavior X} {X0 : Set X}
    (hdc : IsLowerSet C)
    (hJoin : ∀ b : β, ∀ x ∈ blocks b,
      ∃ i : ι, x ∈ blocks (emb i) ∧ blocks (emb i) ⊆ blocks b)
    (hcover : ∀ i : ι, blocks (cover i) ⊂ blocks (emb i))
    (hcellsSup : ∀ i : ι, blocks (emb i) \ blocks (cover i) ⊆ cells i)
    (horder : ∀ i j : ι, blocks (emb i) ⊆ blocks (emb j) → i ≤ j)
    (hcellSub : ∀ i : ι, cells i ⊆ blocks (emb i))
    (hMember : B.InvScopeSat (⋃ i ∈ C, blocks (emb i)) X0) :
    B.InvScopeSat (⋃ i ∈ C, cells i) X0
\end{leansrc}
\par\medskip\noindent{\small \texttt{PPLC.Qtsl.Behavior.}\texttt{\bfseries qBasin\_sound}}\hfill{\footnotesize \path{lean/qtsl/Qtsl/Behavioral.lean}, line 327}\par\nopagebreak
\begin{leansrc}
theorem qBasin_sound {B : Behavior X} {S X0 Bad : Set X}
    (hscope : B.InvScopeSat S X0) (hAvoid : S ∩ Bad = ∅) :
    B.SafetySat X0 Bad
\end{leansrc}
\par\medskip\noindent{\small \texttt{PPLC.Qtsl.}\texttt{\bfseries member\_subset\_cells\_le}}\hfill{\footnotesize \path{lean/qtsl/Qtsl/LatticeInv.lean}, line 44}\par\nopagebreak
\begin{leansrc}
theorem member_subset_cells_le
    (hJoin : ∀ b : β, ∀ x ∈ blocks b,
      ∃ i : ι, x ∈ blocks (emb i) ∧ blocks (emb i) ⊆ blocks b)
    (hcover : ∀ i : ι, blocks (cover i) ⊂ blocks (emb i))
    (hcellsSup : ∀ i : ι, blocks (emb i) \ blocks (cover i) ⊆ cells i)
    (horder : ∀ i j : ι, blocks (emb i) ⊆ blocks (emb j) → i ≤ j)
    (i : ι) :
    blocks (emb i) ⊆ ⋃ j ∈ Set.Iic i, cells j
\end{leansrc}
\par\medskip\noindent{\small \texttt{PPLC.Qtsl.Semiflow.}\texttt{\bfseries qLatticeInv\_sound}}\hfill{\footnotesize \path{lean/qtsl/Qtsl/LatticeInv.lean}, line 65}\par\nopagebreak
\begin{leansrc}
theorem Semiflow.qLatticeInv_sound {φ : Semiflow X} {X0 : Set X}
    (hdc : IsLowerSet C)
    (hJoin : ∀ b : β, ∀ x ∈ blocks b,
      ∃ i : ι, x ∈ blocks (emb i) ∧ blocks (emb i) ⊆ blocks b)
    (hcover : ∀ i : ι, blocks (cover i) ⊂ blocks (emb i))
    (hcellsSup : ∀ i : ι, blocks (emb i) \ blocks (cover i) ⊆ cells i)
    (horder : ∀ i j : ι, blocks (emb i) ⊆ blocks (emb j) → i ≤ j)
    (hcellSub : ∀ i : ι, cells i ⊆ blocks (emb i))
    (hMember : φ.InvScopeSat (⋃ i ∈ C, blocks (emb i)) X0) :
    φ.InvScopeSat (⋃ i ∈ C, cells i) X0
\end{leansrc}
\par\medskip\noindent{\small \texttt{PPLC.Qtsl.Behavior.}\texttt{\bfseries qBasin\_batch}}\hfill{\footnotesize \path{lean/qtsl/Qtsl/Behavioral.lean}, line 336}\par\nopagebreak
\begin{leansrc}
theorem qBasin_batch {B : Behavior X} {S X0 : Set X} {J : Type*}
    {Bad : J → Set X} (hscope : B.InvScopeSat S X0)
    (hAvoid : ∀ j, S ∩ Bad j = ∅) :
    B.SafetySat X0 (⋃ j, Bad j)
\end{leansrc}
\par\medskip\noindent{\small \texttt{PPLC.Qtsl.Behavior.}\texttt{\bfseries InvScopeSat}}\hfill{\footnotesize \path{lean/qtsl/Qtsl/Behavioral.lean}, line 45}\par\nopagebreak
\begin{leansrc}
def InvScopeSat (B : Behavior X) (S X0 : Set X) : Prop
\end{leansrc}
\par\medskip\noindent{\small \texttt{PPLC.Qtsl.Behavior.}\texttt{\bfseries SafetySat}}\hfill{\footnotesize \path{lean/qtsl/Qtsl/Behavioral.lean}, line 42}\par\nopagebreak
\begin{leansrc}
def SafetySat (B : Behavior X) (X0 Bad : Set X) : Prop
\end{leansrc}
\par\medskip\noindent{\small \texttt{PPLC.Qtsl.Semiflow.}\texttt{\bfseries InvScopeSat}}\hfill{\footnotesize \path{lean/qtsl/Qtsl/QtslSoundness.lean}, line 84}\par\nopagebreak
\begin{leansrc}
def InvScopeSat (φ : Semiflow X) (S X0 : Set X) : Prop
\end{leansrc}
\par\medskip\noindent{\small \texttt{PPLC.Qtsl.Semiflow.}\texttt{\bfseries SafetySat}}\hfill{\footnotesize \path{lean/qtsl/Qtsl/QtslSoundness.lean}, line 87}\par\nopagebreak
\begin{leansrc}
def SafetySat (φ : Semiflow X) (X0 B : Set X) : Prop
\end{leansrc}
\par\medskip\noindent{\small \texttt{Qtsl.QtslCompose.System.}\texttt{\bfseries qagTier}}\hfill{\footnotesize \path{lean/qtsl-aux/QtslAux/QtslCompose/Presentation.lean}, line 111}\par\nopagebreak
\begin{leansrc}
def System.qagTier {n : ℕ} (S : System n) : Tier
\end{leansrc}
\par\medskip\noindent{\small \texttt{Qtsl.QtslCompose.}\texttt{\bfseries admissible\_constructs}}\hfill{\footnotesize \path{lean/qtsl-aux/QtslAux/QtslCompose/Presentation.lean}, line 163}\par\nopagebreak
\begin{leansrc}
theorem admissible_constructs {ι κIdx X : Type*} [Fintype ι] [Fintype κIdx]
    {R : Set X} (lam : ι → ℝ) (e : ι → X → ℝ) (nu : κIdx → ℝ) (g : κIdx → X → ℝ)
    (hlam : ∀ i, 0 ≤ lam i)
    (he : ∀ i, ∀ x ∈ R, e i x ≤ 0)
    (hg : ∀ k, ∀ x ∈ R, g k x = 0) :
    ∀ x ∈ R, (∑ i, lam i * e i x) + (∑ k, nu k * g k x) ≤ 0
\end{leansrc}
\par\medskip\noindent{\small \texttt{Qtsl.AGCompose.Parallel.}\texttt{\bfseries fail\_bound}}\hfill{\footnotesize \path{lean/qtsl-aux/QtslAux/AGCompose/GshsParallel.lean}, line 128}\par\nopagebreak
\begin{leansrc}
theorem fail_bound {ι : Type*} (s : Finset ι) {sat : ι → Set Ω} {satOut : Set Ω}
    {ε : ι → ℝ≥0∞}
    (h : ∀ i ∈ s, μ (sat i)ᶜ ≤ ε i)
    (hout : (⋂ i ∈ s, sat i) ⊆ satOut) :
    μ satOutᶜ ≤ ∑ i ∈ s, ε i
\end{leansrc}
\par\medskip\noindent{\small \texttt{Qtsl.QtslCompose.}\texttt{\bfseries fail\_bound\_out}}\hfill{\footnotesize \path{lean/qtsl-aux/QtslAux/QtslCompose/AGCompose.lean}, line 131}\par\nopagebreak
\begin{leansrc}
theorem fail_bound_out [MeasurableSpace Ω] {μ : Measure Ω} {ε : Fin n → ℝ≥0∞}
    {satOut : Set Ω}
    (hw : ∀ j, conformExt ∩ (⋂ i ∈ S.upstream j, sat i) ⊆ conform j)
    (hout : conformExt ∩ (⋂ j, sat j) ⊆ satOut)
    (hε : ∀ j, μ (conform j ∩ (sat j)ᶜ) ≤ ε j) :
    μ (conformExt ∩ satOutᶜ) ≤ ∑ j, ε j
\end{leansrc}
\par\medskip\noindent{\small \texttt{Qtsl.QtslCompose.}\texttt{\bfseries sat\_of\_conform}}\hfill{\footnotesize \path{lean/qtsl-aux/QtslAux/QtslCompose/AGCompose.lean}, line 48}\par\nopagebreak
\begin{leansrc}
theorem sat_of_conform
    (hw : ∀ j, conformExt ∩ (⋂ i ∈ S.upstream j, sat i) ⊆ conform j)
    {ω : Ω} (h : ∀ j, ω ∈ conform j → ω ∈ sat j)
    (hext : ω ∈ conformExt) : ∀ j, ω ∈ sat j
\end{leansrc}
\par\medskip\noindent{\small \texttt{Qtsl.AGCompose.Series.}\texttt{\bfseries fail\_subset}}\hfill{\footnotesize \path{lean/qtsl-aux/QtslAux/AGCompose/GshsSeries.lean}, line 45}\par\nopagebreak
\begin{leansrc}
theorem fail_subset
    (hw : conform₁ ∩ sat₁ ⊆ conform₂)
    (hout : conform₁ ∩ sat₁ ∩ sat₂ ⊆ satOut) :
    conform₁ ∩ satOutᶜ ⊆ (conform₁ ∩ sat₁ᶜ) ∪ (conform₂ ∩ sat₂ᶜ)
\end{leansrc}
\par\medskip\noindent{\small \texttt{Qtsl.QtslCompose.System.}\texttt{\bfseries qagTier\_admissible\_iff}}\hfill{\footnotesize \path{lean/qtsl-aux/QtslAux/QtslCompose/Presentation.lean}, line 116}\par\nopagebreak
\begin{leansrc}
theorem System.qagTier_admissible_iff {n : ℕ} (S : System n) :
    S.qagTier = .admissible ↔
      ∀ i, (S.charts i).profile = .ode ∨ (S.charts i).profile = .discrete
\end{leansrc}
\par\medskip\noindent{\small \texttt{Qtsl.QtslCompose.System.}\texttt{\bfseries qagTier\_core\_iff}}\hfill{\footnotesize \path{lean/qtsl-aux/QtslAux/QtslCompose/Presentation.lean}, line 135}\par\nopagebreak
\begin{leansrc}
theorem System.qagTier_core_iff {n : ℕ} (S : System n) :
    S.qagTier = .core ↔
      ∃ i, (S.charts i).profile = .sde ∨ (S.charts i).profile = .gshs
\end{leansrc}
\par\medskip\noindent{\small \texttt{Qtsl.AGCompose.Series.}\texttt{\bfseries certified}}\hfill{\footnotesize \path{lean/qtsl-aux/QtslAux/AGCompose/GshsSeries.lean}, line 132}\par\nopagebreak
\begin{leansrc}
theorem certified [IsFiniteMeasure μ]
    {Y₁ Y₂ : ℕ → Ω → ℝ} {X₁ : ℕ → Ω → S₁} {X₂ : ℕ → Ω → S₂}
    {B₁ : Set S₁} {B₂ : Set S₂} {a₁ a₂ : ℝ}

    (hY₁ : Supermartingale Y₁ 𝒢 μ) (hY₁0 : 0 ≤ Y₁)
    (hY₂ : Supermartingale Y₂ 𝒢 μ) (hY₂0 : 0 ≤ Y₂)

    (ha₁ : 0 < a₁) (ha₂ : 0 < a₂)
    (hbar₁ : ∀ n ω, X₁ n ω ∈ B₁ → a₁ ≤ Y₁ n ω)
    (hbar₂ : ∀ n ω, X₂ n ω ∈ B₂ → a₂ ≤ Y₂ n ω)

    (hwit₁ : conform₁ ∩ sat₁ᶜ ⊆ {ω | ∃ n, X₁ n ω ∈ B₁})
    (hwit₂ : conform₂ ∩ sat₂ᶜ ⊆ {ω | ∃ n, X₂ n ω ∈ B₂})

    (hw : conform₁ ∩ sat₁ ⊆ conform₂)
    (hout : conform₁ ∩ sat₁ ∩ sat₂ ⊆ satOut) :
    μ (conform₁ ∩ satOutᶜ) ≤
      ENNReal.ofReal (∫ ω, Y₁ 0 ω ∂μ) / ENNReal.ofReal a₁ +
      ENNReal.ofReal (∫ ω, Y₂ 0 ω ∂μ) / ENNReal.ofReal a₂
\end{leansrc}
\par\medskip\noindent{\small \texttt{Qtsl.AGCompose.Parallel.}\texttt{\bfseries certified}}\hfill{\footnotesize \path{lean/qtsl-aux/QtslAux/AGCompose/GshsParallel.lean}, line 58}\par\nopagebreak
\begin{leansrc}
theorem certified (s : Finset ι) [IsFiniteMeasure μ]
    {Y : ι → ℕ → Ω → ℝ} {lam : ι → ℝ} {X : ℕ → Ω → St} {Bad : ι → Set St} {a : ℝ}
    (hY : ∀ i ∈ s, Supermartingale (Y i) 𝒢 μ)
    (hY0 : ∀ i ∈ s, 0 ≤ Y i)
    (hlam : ∀ i ∈ s, 0 ≤ lam i)
    (hbar : ∀ i ∈ s, ∀ n ω, X n ω ∈ Bad i → a ≤ lam i * Y i n ω) :
    ENNReal.ofReal a * μ {ω | ∃ n, X n ω ∈ ⋃ i ∈ s, Bad i} ≤
      ENNReal.ofReal (∑ i ∈ s, lam i * ∫ ω, Y i 0 ω ∂μ)
\end{leansrc}
\par\medskip\noindent{\small \texttt{Qtsl.AGCompose.Parallel.}\texttt{\bfseries certified\_pair}}\hfill{\footnotesize \path{lean/qtsl-aux/QtslAux/AGCompose/GshsParallel.lean}, line 87}\par\nopagebreak
\begin{leansrc}
theorem certified_pair [IsFiniteMeasure μ]
    {Y₁ Y₂ : ℕ → Ω → ℝ} {X : ℕ → Ω → St} {B₁ B₂ : Set St} {lam a : ℝ}
    (hY₁ : Supermartingale Y₁ 𝒢 μ) (hY₁0 : 0 ≤ Y₁)
    (hY₂ : Supermartingale Y₂ 𝒢 μ) (hY₂0 : 0 ≤ Y₂) (hlam : 0 ≤ lam)
    (hbar₁ : ∀ n ω, X n ω ∈ B₁ → a ≤ Y₁ n ω)
    (hbar₂ : ∀ n ω, X n ω ∈ B₂ → a ≤ lam * Y₂ n ω) :
    ENNReal.ofReal a * μ {ω | ∃ n, X n ω ∈ B₁ ∪ B₂} ≤
      ENNReal.ofReal (∫ ω, Y₁ 0 ω ∂μ + lam * ∫ ω, Y₂ 0 ω ∂μ)
\end{leansrc}
\par\medskip\noindent{\small \texttt{Qtsl.AGCompose.SumFormFeedback.}\texttt{\bfseries avoid\_bound}}\hfill{\footnotesize \path{lean/qtsl-aux/QtslAux/GshsFeedbackSum.lean}, line 166}\par\nopagebreak
\begin{leansrc}
theorem avoid_bound {S : Type*} {X : ℕ → Ω → S} {B : Set S} {a : ℝ}
    (hbarrier : ∀ k ω, X k ω ∈ B → a ≤ F.W k ω) :
    ENNReal.ofReal a * μ {ω | ∃ k, X k ω ∈ B} ≤ ENNReal.ofReal (∫ ω, F.W 0 ω ∂μ)
\end{leansrc}
\par\medskip\noindent{\small \texttt{Qtsl.AGCompose.}\texttt{\bfseries SumFormFeedback}}\hfill{\footnotesize \path{lean/qtsl-aux/QtslAux/GshsFeedbackSum.lean}, line 79}\par\nopagebreak
\begin{leansrc}
structure SumFormFeedback (μ : Measure Ω) (𝒢 : Filtration ℕ m0) where …
\end{leansrc}
\par\medskip\noindent{\small \texttt{Qtsl.LayerD.}\texttt{\bfseries sound\_GSHS\_safety\_barrier}}\hfill{\footnotesize \path{lean/qtsl-aux/QtslAux/MaxFormPDMP.lean}, line 608}\par\nopagebreak
\begin{leansrc}
theorem sound_GSHS_safety_barrier [IsMarkovKernel M.seg] [IsMarkovKernel M.jump]
    {W : State → ℝ} {U : Set State} {m : ℝ} (hm : 0 < m)
    (hW : Measurable W) (hU : MeasurableSet U)
    (hSegD : SegmentDynkin M 0 W) (hForce : ForcedNonexpansion M W)
    (hNonneg : NonnegV W) {C : ℝ} (hC0 : 0 ≤ C) (hWC : ∀ y, W y ≤ C)
    (hFloor : ∀ y ∈ U, m ≤ W y) (x : State) :
    Battery.pathMeasure (skel M) (SkelState.alive x) {ω | ∃ n, ω n ∈ skelBad U}
      ≤ ENNReal.ofReal (W x / m)
\end{leansrc}
\par\medskip\noindent{\small \texttt{Qtsl.QtslCompose.Discrete.}\texttt{\bfseries firstFail}}\hfill{\footnotesize \path{lean/qtsl-aux/QtslAux/QtslCompose/AGComposeDiscrete.lean}, line 52}\par\nopagebreak
\begin{leansrc}
def firstFail (bad : Fin m → Set S) (j : Fin m) : Set (ℕ → S)
\end{leansrc}
\par\medskip\noindent{\small \texttt{Qtsl.QtslCompose.Discrete.}\texttt{\bfseries measure\_firstFail\_le}}\hfill{\footnotesize \path{lean/qtsl-aux/QtslAux/QtslCompose/AGComposeDiscrete.lean}, line 86}\par\nopagebreak
\begin{leansrc}
theorem measure_firstFail_le {bad : Fin m → Set S} {V : S → ℝ} {C : ℝ} {j : Fin m}
    (hA : MeasurableSet (⋃ i, bad i)) (hV : Measurable V)
    (hV0 : ∀ y, 0 ≤ V y) (hVC : ∀ y, V y ≤ C)
    (hbar : ∀ y ∈ bad j, 1 ≤ V y)
    (hdrift : ∀ y, y ∉ (⋃ i, bad i) → ∫ z, V z ∂K y ≤ V y) (x : S) :
    pathMeasure K x (firstFail bad j) ≤ ENNReal.ofReal (V x)
\end{leansrc}
\par\medskip\noindent{\small \texttt{Qtsl.QtslCompose.Discrete.}\texttt{\bfseries fail\_bound\_wired}}\hfill{\footnotesize \path{lean/qtsl-aux/QtslAux/QtslCompose/AGComposeDiscrete.lean}, line 156}\par\nopagebreak
\begin{leansrc}
theorem fail_bound_wired (Sys : System m) {conf bad : Fin m → Set S}
    {V : Fin m → S → ℝ} {C : Fin m → ℝ}
    (hA : MeasurableSet (⋃ i, bad i))
    (hw : ∀ j, (⋂ i ∈ Sys.upstream j, (bad i)ᶜ) ⊆ conf j)
    (hV : ∀ j, Measurable (V j)) (hV0 : ∀ j y, 0 ≤ V j y)
    (hVC : ∀ j y, V j y ≤ C j)
    (hbar : ∀ j, ∀ y ∈ bad j, 1 ≤ V j y)
    (hdrift : ∀ j, ∀ y ∈ conf j, y ∉ bad j → ∫ z, V j z ∂K y ≤ V j y) (x : S) :
    pathMeasure K x {ω | ∃ n, ω n ∈ ⋃ i, bad i}
      ≤ ∑ j, ENNReal.ofReal (V j x)
\end{leansrc}
\par\medskip\noindent{\small \texttt{MeasureTheory.Supermartingale.}\texttt{\bfseries ville\_ineq}}\hfill{\footnotesize \path{lean/qtsl-aux/QtslAux/Ville.lean}, line 133}\par\nopagebreak
\begin{leansrc}
theorem Supermartingale.ville_ineq [IsFiniteMeasure μ] (hsup : Supermartingale f 𝒢
      μ)
    (hnonneg : 0 ≤ f) (ε : ℝ≥0) :
    ε * μ {ω | ∃ k, (ε : ℝ) ≤ f k ω} ≤ ENNReal.ofReal (∫ ω, f 0 ω ∂μ)
\end{leansrc}
\par\medskip\noindent{\small \texttt{Qtsl.QtslCompose.Discrete.}\texttt{\bfseries drift\_parallelComp}}\hfill{\footnotesize \path{lean/qtsl-aux/QtslAux/QtslCompose/AGComposeDiscrete.lean}, line 207}\par\nopagebreak
\begin{leansrc}
theorem drift_parallelComp {bad₁ : Set S₁} {bad₂ : Set S₂}
    {V₁ : S₁ → ℝ} {V₂ : S₂ → ℝ} {lam : ℝ} (hlam : 0 ≤ lam)
    (hint₁ : ∀ y, y ∉ bad₁ → Integrable V₁ (K₁ y))
    (hint₂ : ∀ y, y ∉ bad₂ → Integrable V₂ (K₂ y))
    (hd₁ : ∀ y, y ∉ bad₁ → ∫ z, V₁ z ∂K₁ y ≤ V₁ y)
    (hd₂ : ∀ y, y ∉ bad₂ → ∫ z, V₂ z ∂K₂ y ≤ V₂ y) :
    ∀ p : S₁ × S₂, p ∉ (Prod.fst ⁻¹' bad₁ ∪ Prod.snd ⁻¹' bad₂) →
      ∫ q, (V₁ q.1 + lam * V₂ q.2) ∂(K₁ ∥ₖ K₂) p
        ≤ V₁ p.1 + lam * V₂ p.2
\end{leansrc}
\par\medskip\noindent{\small \texttt{Qtsl.QtslCompose.Discrete.}\texttt{\bfseries avoid\_parallelComp}}\hfill{\footnotesize \path{lean/qtsl-aux/QtslAux/QtslCompose/AGComposeDiscrete.lean}, line 233}\par\nopagebreak
\begin{leansrc}
theorem avoid_parallelComp {bad₁ : Set S₁} {bad₂ : Set S₂}
    {V₁ : S₁ → ℝ} {V₂ : S₂ → ℝ} {C₁ C₂ lam : ℝ} (hlam : 0 ≤ lam)
    (hbad₁ : MeasurableSet bad₁) (hbad₂ : MeasurableSet bad₂)
    (hV₁ : Measurable V₁) (hV₂ : Measurable V₂)
    (hV₁0 : ∀ y, 0 ≤ V₁ y) (hV₂0 : ∀ y, 0 ≤ V₂ y)
    (hV₁C : ∀ y, V₁ y ≤ C₁) (hV₂C : ∀ y, V₂ y ≤ C₂)
    (hbar₁ : ∀ y ∈ bad₁, 1 ≤ V₁ y) (hbar₂ : ∀ y ∈ bad₂, 1 ≤ lam * V₂ y)
    (hd₁ : ∀ y, y ∉ bad₁ → ∫ z, V₁ z ∂K₁ y ≤ V₁ y)
    (hd₂ : ∀ y, y ∉ bad₂ → ∫ z, V₂ z ∂K₂ y ≤ V₂ y) (x : S₁ × S₂) :
    pathMeasure (K₁ ∥ₖ K₂) x
        {ω | ∃ n, ω n ∈ (Prod.fst ⁻¹' bad₁ ∪ Prod.snd ⁻¹' bad₂ : Set (S₁ × S₂))}
      ≤ ENNReal.ofReal (V₁ x.1 + lam * V₂ x.2)
\end{leansrc}
\par\medskip\noindent{\small \texttt{Qtsl.QtslCompose.Discrete.}\texttt{\bfseries smallGain\_drift}}\hfill{\footnotesize \path{lean/qtsl-aux/QtslAux/QtslCompose/AGComposeDiscrete.lean}, line 278}\par\nopagebreak
\begin{leansrc}
theorem smallGain_drift (hlam : 0 ≤ lam)
    (hint₁ : ∀ y, y ∉ A → Integrable V₁ (K y))
    (hint₂ : ∀ y, y ∉ A → Integrable V₂ (K y))
    (hd₁ : ∀ y, y ∉ A → ∫ z, V₁ z ∂K y ≤ V₁ y - alpha₁ y + sigma₁ y)
    (hd₂ : ∀ y, y ∉ A → ∫ z, V₂ z ∂K y ≤ V₂ y - alpha₂ y + sigma₂ y)
    (hm₁ : ∀ y, y ∉ A → lam * sigma₂ y ≤ (1 - theta) * alpha₁ y)
    (hm₂ : ∀ y, y ∉ A → sigma₁ y ≤ lam * ((1 - theta) * alpha₂ y)) :
    ∀ y, y ∉ A → ∫ z, (V₁ z + lam * V₂ z) ∂K y
      ≤ (V₁ y + lam * V₂ y) - theta * (alpha₁ y + lam * alpha₂ y)
\end{leansrc}
\par\medskip\noindent{\small \texttt{Qtsl.QtslCompose.Discrete.}\texttt{\bfseries avoid\_smallGain}}\hfill{\footnotesize \path{lean/qtsl-aux/QtslAux/QtslCompose/AGComposeDiscrete.lean}, line 310}\par\nopagebreak
\begin{leansrc}
theorem avoid_smallGain (hlam : 0 ≤ lam) (htheta : 0 ≤ theta)
    {C₁ C₂ : ℝ} (hA : MeasurableSet A)
    (hV₁ : Measurable V₁) (hV₂ : Measurable V₂)
    (hV₁0 : ∀ y, 0 ≤ V₁ y) (hV₂0 : ∀ y, 0 ≤ V₂ y)
    (hV₁C : ∀ y, V₁ y ≤ C₁) (hV₂C : ∀ y, V₂ y ≤ C₂)
    (hbar : ∀ y ∈ A, 1 ≤ V₁ y + lam * V₂ y)
    (halpha₁ : ∀ y, y ∉ A → 0 ≤ alpha₁ y)
    (halpha₂ : ∀ y, y ∉ A → 0 ≤ alpha₂ y)
    (hd₁ : ∀ y, y ∉ A → ∫ z, V₁ z ∂K y ≤ V₁ y - alpha₁ y + sigma₁ y)
    (hd₂ : ∀ y, y ∉ A → ∫ z, V₂ z ∂K y ≤ V₂ y - alpha₂ y + sigma₂ y)
    (hm₁ : ∀ y, y ∉ A → lam * sigma₂ y ≤ (1 - theta) * alpha₁ y)
    (hm₂ : ∀ y, y ∉ A → sigma₁ y ≤ lam * ((1 - theta) * alpha₂ y)) (x : S) :
    pathMeasure K x {ω | ∃ n, ω n ∈ A}
      ≤ ENNReal.ofReal (V₁ x + lam * V₂ x)
\end{leansrc}
\par\medskip\noindent{\small \texttt{Qtsl.AGCompose.Series.}\texttt{\bfseries tower\_step}}\hfill{\footnotesize \path{lean/qtsl-aux/QtslAux/AGCompose/GshsSeries.lean}, line 81}\par\nopagebreak
\begin{leansrc}
theorem tower_step [IsFiniteMeasure μ] {mG : MeasurableSpace Ω} (hG : mG ≤ m0)
    {S : Set Ω} (hS : MeasurableSet[mG] S) {E : Set Ω} (hE : MeasurableSet[m0] E)
    {p : ℝ}
    (hcond : ∀ᵐ ω ∂μ, ω ∈ S → p ≤ (μ[E.indicator (fun _ => (1 : ℝ)) | mG]) ω) :
    ENNReal.ofReal p * μ S ≤ μ (S ∩ E)
\end{leansrc}
\par\medskip\noindent{\small \texttt{Qtsl.AGCompose.Parallel.}\texttt{\bfseries supermartingale\_weighted\_sum}}\hfill{\footnotesize \path{lean/qtsl-aux/QtslAux/AGCompose/GshsParallel.lean}, line 31}\par\nopagebreak
\begin{leansrc}
theorem supermartingale_weighted_sum {ι : Type*} (s : Finset ι) [IsFiniteMeasure μ]
    {Y : ι → ℕ → Ω → ℝ} {lam : ι → ℝ}
    (hY : ∀ i ∈ s, Supermartingale (Y i) 𝒢 μ) (hlam : ∀ i ∈ s, 0 ≤ lam i) :
    Supermartingale (fun n ω => ∑ i ∈ s, lam i * Y i n ω) 𝒢 μ
\end{leansrc}
\par\medskip\noindent{\small \texttt{Qtsl.AGCompose.SumFormFeedback.}\texttt{\bfseries debut\_avoid\_bound}}\hfill{\footnotesize \path{lean/qtsl-aux/QtslAux/GshsFeedbackSum.lean}, line 171}\par\nopagebreak
\begin{leansrc}
theorem debut_avoid_bound {S : Type*} {X : ℕ → Ω → S} {B : Set S} {a : ℝ}
    (hbarrier : ∀ k ω, X k ω ∈ B → a ≤ F.W k ω) {E : Set Ω}
    (hdebut : E ⊆ {ω | ∃ k, X k ω ∈ B}) :
    ENNReal.ofReal a * μ E ≤ ENNReal.ofReal (∫ ω, F.W 0 ω ∂μ)
\end{leansrc}
\par\medskip\noindent{\small \texttt{Qtsl.QtslCompose.}\texttt{\bfseries qualitative}}\hfill{\footnotesize \path{lean/qtsl-aux/QtslAux/QtslCompose/AGCompose.lean}, line 69}\par\nopagebreak
\begin{leansrc}
theorem qualitative [MeasurableSpace Ω] {μ : Measure Ω}
    (hw : ∀ j, conformExt ∩ (⋂ i ∈ S.upstream j, sat i) ⊆ conform j)
    (h : ∀ j, ∀ᵐ ω ∂μ, ω ∈ conform j → ω ∈ sat j)
    {satOut : Set Ω} (hout : conformExt ∩ (⋂ j, sat j) ⊆ satOut) :
    ∀ᵐ ω ∂μ, ω ∈ conformExt → ω ∈ satOut
\end{leansrc}
\par\medskip\noindent{\small \texttt{Qtsl.AGCompose.Series.}\texttt{\bfseries qualitative}}\hfill{\footnotesize \path{lean/qtsl-aux/QtslAux/AGCompose/GshsSeries.lean}, line 35}\par\nopagebreak
\begin{leansrc}
theorem qualitative
    (h₁ : ∀ᵐ ω ∂μ, ω ∈ conform₁ → ω ∈ sat₁)
    (hw : conform₁ ∩ sat₁ ⊆ conform₂)
    (h₂ : ∀ᵐ ω ∂μ, ω ∈ conform₂ → ω ∈ sat₂)
    (hout : conform₁ ∩ sat₁ ∩ sat₂ ⊆ satOut) :
    ∀ᵐ ω ∂μ, ω ∈ conform₁ → ω ∈ satOut
\end{leansrc}
\par\medskip\noindent{\small \texttt{Qtsl.AGCompose.Parallel.}\texttt{\bfseries supermartingale\_weighted\_sum}}\hfill{\footnotesize \path{lean/qtsl-aux/QtslAux/AGCompose/GshsParallel.lean}, line 31}\par\nopagebreak
\begin{leansrc}
theorem supermartingale_weighted_sum {ι : Type*} (s : Finset ι) [IsFiniteMeasure μ]
    {Y : ι → ℕ → Ω → ℝ} {lam : ι → ℝ}
    (hY : ∀ i ∈ s, Supermartingale (Y i) 𝒢 μ) (hlam : ∀ i ∈ s, 0 ≤ lam i) :
    Supermartingale (fun n ω => ∑ i ∈ s, lam i * Y i n ω) 𝒢 μ
\end{leansrc}
\par\medskip\noindent{\small \texttt{Qtsl.AGCompose.}\texttt{\bfseries SkeletonCertificate}}\hfill{\footnotesize \path{lean/qtsl-aux/QtslAux/GshsFeedbackSum.lean}, line 39}\par\nopagebreak
\begin{leansrc}
structure SkeletonCertificate (μ : Measure Ω) (𝒢 : Filtration ℕ m0) where …
\end{leansrc}
\par\medskip\noindent{\small \texttt{Qtsl.AGCompose.SkeletonCertificate.}\texttt{\bfseries step}}\hfill{\footnotesize \path{lean/qtsl-aux/QtslAux/GshsFeedbackSum.lean}, line 65}\par\nopagebreak
\begin{leansrc}
theorem step (C : SkeletonCertificate μ 𝒢) (k : ℕ) :
    μ[C.val (k + 1) + C.alphaInt k - C.sigmaInt k | 𝒢 k] ≤ᵐ[μ] C.val k
\end{leansrc}
\par\medskip\noindent{\small \texttt{Qtsl.AGCompose.SumFormFeedback.}\texttt{\bfseries drift}}\hfill{\footnotesize \path{lean/qtsl-aux/QtslAux/GshsFeedbackSum.lean}, line 124}\par\nopagebreak
\begin{leansrc}
theorem drift (k : ℕ) : μ[F.W (k + 1) + F.margin k | 𝒢 k] ≤ᵐ[μ] F.W k
\end{leansrc}
\par\medskip\noindent{\small \texttt{Qtsl.AGCompose.SumFormFeedback.}\texttt{\bfseries compensated\_supermartingale}}\hfill{\footnotesize \path{lean/qtsl-aux/QtslAux/GshsFeedbackSum.lean}, line 196}\par\nopagebreak
\begin{leansrc}
theorem compensated_supermartingale
    (hα₁ : ∀ k, StronglyMeasurable[𝒢 (k + 1)] (F.cert₁.alphaInt k))
    (hα₂ : ∀ k, StronglyMeasurable[𝒢 (k + 1)] (F.cert₂.alphaInt k)) :
    Supermartingale F.compensated 𝒢 μ
\end{leansrc}
\par\medskip\noindent{\small \texttt{Qtsl.AGCompose.SumFormFeedback.}\texttt{\bfseries margin\_budget}}\hfill{\footnotesize \path{lean/qtsl-aux/QtslAux/GshsFeedbackSum.lean}, line 226}\par\nopagebreak
\begin{leansrc}
theorem margin_budget
    (hα₁ : ∀ k, StronglyMeasurable[𝒢 (k + 1)] (F.cert₁.alphaInt k))
    (hα₂ : ∀ k, StronglyMeasurable[𝒢 (k + 1)] (F.cert₂.alphaInt k)) (n : ℕ) :
    ∫ ω, (∑ j ∈ Finset.range n, F.margin j) ω ∂μ ≤ ∫ ω, F.W 0 ω ∂μ
\end{leansrc}
\par\medskip\noindent{\small \texttt{Qtsl.QtslCompose.ChartShape.}\texttt{\bfseries profile}}\hfill{\footnotesize \path{lean/qtsl-aux/QtslAux/QtslCompose/Presentation.lean}, line 57}\par\nopagebreak
\begin{leansrc}
def ChartShape.profile (c : ChartShape) : Profile
\end{leansrc}
\par\medskip\noindent{\small \texttt{Qtsl.QtslCompose.}\texttt{\bfseries IsClassK0Inf}}\hfill{\footnotesize \path{lean/qtsl-aux/QtslAux/QtslCompose/WireCertDischarge.lean}, line 38}\par\nopagebreak
\begin{leansrc}
structure IsClassK0Inf (κ : ℝ → ℝ) : Prop where …
\end{leansrc}
\par\medskip\noindent{\small \texttt{Qtsl.QtslCompose.}\texttt{\bfseries polynomial\_discharge}}\hfill{\footnotesize \path{lean/qtsl-aux/QtslAux/QtslCompose/WireCertDischarge.lean}, line 75}\par\nopagebreak
\begin{leansrc}
theorem polynomial_discharge (p : Polynomial ℝ)
    (h0 : p.coeff 0 = 0)
    (hderiv : ∀ x : ℝ, 0 ≤ x → 0 ≤ p.derivative.eval x)
    (hdeg : 0 < p.degree) (hlead : 0 ≤ p.leadingCoeff) :
    IsClassK0Inf (fun x => p.eval x) where
  zero_at_zero
\end{leansrc}
\par\medskip\noindent{\small \texttt{Qtsl.QtslCompose.}\texttt{\bfseries slack\_triangle}}\hfill{\footnotesize \path{lean/qtsl-aux/QtslAux/QtslCompose/WireCertDischarge.lean}, line 94}\par\nopagebreak
\begin{leansrc}
theorem slack_triangle {X : Type*} {f g h : X → ℝ} {c₁ c₂ : ℝ}
    (h₁ : ∀ x, f x ≤ g x + c₁) (h₂ : ∀ x, g x ≤ h x + c₂) :
    ∀ x, f x ≤ h x + (c₁ + c₂)
\end{leansrc}
\par\medskip\noindent{\small \texttt{Qtsl.QtslCompose.}\texttt{\bfseries guarantee\_leg\_comp}}\hfill{\footnotesize \path{lean/qtsl-aux/QtslAux/QtslCompose/WireCertDischarge.lean}, line 100}\par\nopagebreak
\begin{leansrc}
theorem guarantee_leg_comp {A B C : Type*} {w₁ : A → B} {w₂ : B → C}
    {γ₁ : A → ℝ} {γ₂ : B → ℝ} {γ₃ : C → ℝ}
    (h₁ : ∀ a, γ₂ (w₁ a) ≤ γ₁ a) (h₂ : ∀ b, γ₃ (w₂ b) ≤ γ₂ b) :
    ∀ a, γ₃ (w₂ (w₁ a)) ≤ γ₁ a
\end{leansrc}
\par\medskip\noindent{\small \texttt{Qtsl.QtslCompose.IsClassK0Inf.}\texttt{\bfseries add}}\hfill{\footnotesize \path{lean/qtsl-aux/QtslAux/QtslCompose/WireCertDischarge.lean}, line 52}\par\nopagebreak
\begin{leansrc}
theorem add (h₁ : IsClassK0Inf κ₁) (h₂ : IsClassK0Inf κ₂) :
    IsClassK0Inf (fun x => κ₁ x + κ₂ x) where
  zero_at_zero
\end{leansrc}
\par\medskip\noindent{\small \texttt{Qtsl.QtslCompose.IsClassK0Inf.}\texttt{\bfseries comp}}\hfill{\footnotesize \path{lean/qtsl-aux/QtslAux/QtslCompose/WireCertDischarge.lean}, line 63}\par\nopagebreak
\begin{leansrc}
theorem comp (h₁ : IsClassK0Inf κ₁) (h₂ : IsClassK0Inf κ₂) :
    IsClassK0Inf (κ₁ ∘ κ₂) where
  zero_at_zero
\end{leansrc}
\par\medskip\noindent{\small \texttt{Qtsl.QtslCompose.IsClassK0Inf.}\texttt{\bfseries nonneg}}\hfill{\footnotesize \path{lean/qtsl-aux/QtslAux/QtslCompose/WireCertDischarge.lean}, line 47}\par\nopagebreak
\begin{leansrc}
theorem nonneg (h : IsClassK0Inf κ) {x : ℝ} (hx : 0 ≤ x) : 0 ≤ κ x
\end{leansrc}
\par\medskip\noindent{\small \texttt{BernsteinLogic.QtslCompose.Marks.}\texttt{\bfseries ConditionalDelivery}}\hfill{\footnotesize \path{lean/qtsl-aux/QtslAux/QtslCompose/AGComposeMarks.lean}, line 157}\par\nopagebreak
\begin{leansrc}
def ConditionalDelivery (N : MarkNet n V) (mark : Fin n → Set V) (i : Fin n) : Prop
\end{leansrc}
\par\medskip\noindent{\small \texttt{BernsteinLogic.QtslCompose.Marks.}\texttt{\bfseries DepLocal}}\hfill{\footnotesize \path{lean/qtsl-aux/QtslAux/QtslCompose/AGComposeMarks.lean}, line 109}\par\nopagebreak
\begin{leansrc}
def DepLocal (N : MarkNet n V) : Prop
\end{leansrc}
\par\medskip\noindent{\small \texttt{BernsteinLogic.QtslCompose.Marks.}\texttt{\bfseries MarkFlowFeasible}}\hfill{\footnotesize \path{lean/qtsl-aux/QtslAux/QtslCompose/AGComposeMarks.lean}, line 123}\par\nopagebreak
\begin{leansrc}
def MarkFlowFeasible (N : MarkNet n V) (t : Fin n → ℕ) : Prop
\end{leansrc}
\par\medskip\noindent{\small \texttt{BernsteinLogic.QtslCompose.Marks.}\texttt{\bfseries CarvedDelivery}}\hfill{\footnotesize \path{lean/qtsl-aux/QtslAux/QtslCompose/AGComposeCarves.lean}, line 69}\par\nopagebreak
\begin{leansrc}
def CarvedDelivery (N : MarkNet n V) (mark P : Fin n → Set V) : Prop
\end{leansrc}
\par\medskip\noindent{\small \texttt{BernsteinLogic.QtslCompose.Marks.}\texttt{\bfseries CarveClosure}}\hfill{\footnotesize \path{lean/qtsl-aux/QtslAux/QtslCompose/AGComposeCarves.lean}, line 76}\par\nopagebreak
\begin{leansrc}
def CarveClosure (N : MarkNet n V) (P : Fin n → Set V) (t : Fin n → ℕ) : Prop
\end{leansrc}
\par\medskip\noindent{\small \texttt{BernsteinLogic.QtslCompose.Marks.}\texttt{\bfseries PinnedDelivery}}\hfill{\footnotesize \path{lean/qtsl-aux/QtslAux/QtslCompose/AGComposeMarks.lean}, line 130}\par\nopagebreak
\begin{leansrc}
def PinnedDelivery (N : MarkNet n V) (mark : Fin n → Set V) (a : Fin n → V) : Prop
\end{leansrc}
\par\medskip\noindent{\small \texttt{BernsteinLogic.QtslCompose.Marks.}\texttt{\bfseries InstantClosure}}\hfill{\footnotesize \path{lean/qtsl-aux/QtslAux/QtslCompose/AGComposeMarks.lean}, line 148}\par\nopagebreak
\begin{leansrc}
def InstantClosure (N : MarkNet n V) (a : Fin n → V) (t : Fin n → ℕ) : Prop
\end{leansrc}
\par\medskip\noindent{\small \texttt{BernsteinLogic.QtslCompose.Marks.}\texttt{\bfseries carvedDelivery\_of\_pinned}}\hfill{\footnotesize \path{lean/qtsl-aux/QtslAux/QtslCompose/AGComposeCarves.lean}, line 276}\par\nopagebreak
\begin{leansrc}
theorem carvedDelivery_of_pinned (hpin : PinnedDelivery N mark a) :
    CarvedDelivery N mark (fun i => {a i})
\end{leansrc}
\par\medskip\noindent{\small \texttt{BernsteinLogic.QtslCompose.Marks.}\texttt{\bfseries carveClosure\_of\_pinned}}\hfill{\footnotesize \path{lean/qtsl-aux/QtslAux/QtslCompose/AGComposeCarves.lean}, line 290}\par\nopagebreak
\begin{leansrc}
theorem carveClosure_of_pinned (hfeas : MarkFlowFeasible N t)
    (hpin : PinnedDelivery N mark a) (hclose : InstantClosure N a t) :
    CarveClosure N (fun i => {a i}) t
\end{leansrc}
\par\medskip\noindent{\small \texttt{BernsteinLogic.QtslCompose.Marks.}\texttt{\bfseries marks\_of\_pinned\_delivery\_via\_carve}}\hfill{\footnotesize \path{lean/qtsl-aux/QtslAux/QtslCompose/AGComposeCarves.lean}, line 315}\par\nopagebreak
\begin{leansrc}
theorem marks_of_pinned_delivery_via_carve (hser : Serial N)
    (hfeas : MarkFlowFeasible N t) (hpin : PinnedDelivery N mark a)
    (hclose : InstantClosure N a t) :
    ∃ ω : ℕ → Fin n → V, JointRun N ω ∧ ∀ i, ω (t i) i ∈ mark i
\end{leansrc}
\par\medskip\noindent{\small \texttt{BernsteinLogic.QtslCompose.Marks.}\texttt{\bfseries carvedDelivery\_mark\_iff\_conditional}}\hfill{\footnotesize \path{lean/qtsl-aux/QtslAux/QtslCompose/AGComposeCarves.lean}, line 332}\par\nopagebreak
\begin{leansrc}
theorem carvedDelivery_mark_iff_conditional :
    CarvedDelivery N mark mark ↔ ∀ i, ConditionalDelivery N mark i
\end{leansrc}
\par\medskip\noindent{\small \texttt{BernsteinLogic.QtslCompose.Marks.}\texttt{\bfseries marks\_of\_carved\_delivery}}\hfill{\footnotesize \path{lean/qtsl-aux/QtslAux/QtslCompose/AGComposeCarves.lean}, line 260}\par\nopagebreak
\begin{leansrc}
theorem marks_of_carved_delivery (hser : Serial N) (hfeas : MarkFlowFeasible N t)
    (hpin : CarvedDelivery N mark P) (hclose : CarveClosure N P t) :
    ∃ ω : ℕ → Fin n → V, JointRun N ω ∧ ∀ i, ω (t i) i ∈ mark i
\end{leansrc}
\par\medskip\noindent{\small \texttt{BernsteinLogic.QtslCompose.Marks.}\texttt{\bfseries marks\_of\_conditional\_delivery\_acyclic}}\hfill{\footnotesize \path{lean/qtsl-aux/QtslAux/QtslCompose/AGComposeMarks.lean}, line 421}\par\nopagebreak
\begin{leansrc}
theorem marks_of_conditional_delivery_acyclic (hser : Serial N) (hloc : DepLocal N)
    (ord : Fin n → ℕ) (hord : ∀ i, ∀ j ∈ N.dep i, ord j < ord i)
    (hfeas : MarkFlowFeasible N t) (hdel : ∀ i, ConditionalDelivery N mark i) :
    ∃ ω : ℕ → Fin n → V, JointRun N ω ∧ ∀ i, ω (t i) i ∈ mark i
\end{leansrc}
\par\medskip\noindent{\small \texttt{BernsteinLogic.QtslCompose.Marks.}\texttt{\bfseries marks\_acyclic\_via\_maximal\_carve}}\hfill{\footnotesize \path{lean/qtsl-aux/QtslAux/QtslCompose/AGComposeCarves.lean}, line 475}\par\nopagebreak
\begin{leansrc}
theorem marks_acyclic_via_maximal_carve (hser : Serial N) (hloc : DepLocal N)
    (ord : Fin n → ℕ) (hord : ∀ i, ∀ j ∈ N.dep i, ord j < ord i)
    (hfeas : MarkFlowFeasible N t)
    (hdel : ∀ i, ConditionalDelivery N mark i) :
    ∃ ω : ℕ → Fin n → V, JointRun N ω ∧ ∀ i, ω (t i) i ∈ mark i
\end{leansrc}
\par\medskip\noindent{\small \texttt{BernsteinLogic.QtslCompose.Marks.}\texttt{\bfseries JointRun}}\hfill{\footnotesize \path{lean/qtsl-aux/QtslAux/QtslCompose/AGComposeMarks.lean}, line 104}\par\nopagebreak
\begin{leansrc}
def JointRun (N : MarkNet n V) (ω : ℕ → Fin n → V) : Prop
\end{leansrc}
\par\medskip\noindent{\small \texttt{BernsteinLogic.QtslCompose.Marks.}\texttt{\bfseries z1\_refutes\_walkless}}\hfill{\footnotesize \path{lean/qtsl-aux/QtslAux/QtslCompose/AGComposeMarks.lean}, line 505}\par\nopagebreak
\begin{leansrc}
theorem z1_refutes_walkless :
    ¬ (∀ (m : ℕ) (W : Type) (_ : Nonempty W) (N : MarkNet m W)
        (mark : Fin m → Set W) (a : Fin m → W) (t : Fin m → ℕ),
        Serial N → PinnedDelivery N mark a → InstantClosure N a t →
        ∃ ω : ℕ → Fin m → W, JointRun N ω ∧ ∀ i, ∃ k, ω k i ∈ mark i)
\end{leansrc}
\par\medskip\noindent{\small \texttt{BernsteinLogic.QtslCompose.Marks.}\texttt{\bfseries z1\_refutes\_walkless\_carve}}\hfill{\footnotesize \path{lean/qtsl-aux/QtslAux/QtslCompose/AGComposeCarves.lean}, line 590}\par\nopagebreak
\begin{leansrc}
theorem z1_refutes_walkless_carve :
    ¬ (∀ (m : ℕ) (W : Type) (_ : Nonempty W) (N : MarkNet m W)
        (mark P : Fin m → Set W) (t : Fin m → ℕ),
        Serial N → CarvedDelivery N mark P → CarveClosure N P t →
        ∃ ω : ℕ → Fin m → W, JointRun N ω ∧ ∀ i, ∃ k, ω k i ∈ mark i)
\end{leansrc}
\par\medskip\noindent{\small \texttt{BernsteinLogic.QtslCompose.Marks.}\texttt{\bfseries z5\_refutes\_pinless}}\hfill{\footnotesize \path{lean/qtsl-aux/QtslAux/QtslCompose/AGComposeMarks.lean}, line 597}\par\nopagebreak
\begin{leansrc}
theorem z5_refutes_pinless :
    ¬ (∀ (m : ℕ) (W : Type) (_ : Nonempty W) (N : MarkNet m W)
        (mark : Fin m → Set W) (t : Fin m → ℕ),
        Serial N → DepLocal N → (∀ i, ConditionalDelivery N mark i) →
        MarkFlowFeasible N t → W0Serial N → (∃ ω, JointRun N ω) →
        ∃ ω : ℕ → Fin m → W, JointRun N ω ∧ ∀ i, ∃ k, ω k i ∈ mark i)
\end{leansrc}
\par\medskip\noindent{\small \texttt{BernsteinLogic.QtslCompose.Marks.}\texttt{\bfseries z5\_refutes\_closureless}}\hfill{\footnotesize \path{lean/qtsl-aux/QtslAux/QtslCompose/AGComposeCarves.lean}, line 530}\par\nopagebreak
\begin{leansrc}
theorem z5_refutes_closureless :
    ¬ (∀ (m : ℕ) (W : Type) (_ : Nonempty W) (N : MarkNet m W)
        (mark P : Fin m → Set W) (t : Fin m → ℕ),
        Serial N → DepLocal N → CarvedDelivery N mark P →
        MarkFlowFeasible N t → W0Serial N → (∃ ω, JointRun N ω) →
        ∃ ω : ℕ → Fin m → W, JointRun N ω ∧ ∀ i, ∃ k, ω k i ∈ mark i)
\end{leansrc}
\par\medskip\noindent{\small \texttt{BernsteinLogic.QtslCompose.Marks.}\texttt{\bfseries z5\_carvedDelivery\_mark}}\hfill{\footnotesize \path{lean/qtsl-aux/QtslAux/QtslCompose/AGComposeCarves.lean}, line 495}\par\nopagebreak
\begin{leansrc}
theorem z5_carvedDelivery_mark : CarvedDelivery z5Net z5Mark z5Mark
\end{leansrc}
\par\medskip\noindent{\small \texttt{BernsteinLogic.QtslCompose.Marks.}\texttt{\bfseries z5\_no\_carve\_closure}}\hfill{\footnotesize \path{lean/qtsl-aux/QtslAux/QtslCompose/AGComposeCarves.lean}, line 509}\par\nopagebreak
\begin{leansrc}
theorem z5_no_carve_closure :
    ¬ ∃ P : Fin 2 → Set (Fin 3), (∀ i, P i ⊆ z5Mark i) ∧
      CarveClosure z5Net P ![0, 0]
\end{leansrc}
\par\medskip\noindent{\small \texttt{BernsteinLogic.QtslCompose.Marks.}\texttt{\bfseries z5\_no\_pin}}\hfill{\footnotesize \path{lean/qtsl-aux/QtslAux/QtslCompose/AGComposeMarks.lean}, line 663}\par\nopagebreak
\begin{leansrc}
theorem z5_no_pin : ¬ ∃ a : Fin 2 → Fin 3, PinnedDelivery z5Net z5Mark a
\end{leansrc}
\par\medskip\noindent{\small \texttt{BernsteinLogic.QtslCompose.Marks.}\texttt{\bfseries z1\_walk\_infeasible}}\hfill{\footnotesize \path{lean/qtsl-aux/QtslAux/QtslCompose/AGComposeMarks.lean}, line 460}\par\nopagebreak
\begin{leansrc}
theorem z1_walk_infeasible : ¬ ∃ t : Fin 2 → ℕ, MarkFlowFeasible z1Net t
\end{leansrc}
\par\medskip\noindent{\small \texttt{BernsteinLogic.QtslCompose.Marks.}\texttt{\bfseries allOnes\_walk\_feasible}}\hfill{\footnotesize \path{lean/qtsl-aux/QtslAux/QtslCompose/AGComposeMarks.lean}, line 469}\par\nopagebreak
\begin{leansrc}
theorem allOnes_walk_feasible : MarkFlowFeasible allOnesNet ![0, 0]
\end{leansrc}
\par\medskip\noindent{\small \texttt{BernsteinLogic.QtslCompose.Marks.}\texttt{\bfseries allOnes\_marks\_realized}}\hfill{\footnotesize \path{lean/qtsl-aux/QtslAux/QtslCompose/AGComposeMarks.lean}, line 677}\par\nopagebreak
\begin{leansrc}
theorem allOnes_marks_realized :
    ∃ ω : ℕ → Fin 2 → Bool, JointRun allOnesNet ω ∧
      ∀ i, ω ((![0, 0] : Fin 2 → ℕ) i) i ∈ boolMark i
\end{leansrc}
\par\medskip\noindent{\small \texttt{BernsteinLogic.QtslCompose.Marks.}\texttt{\bfseries allOnes\_direct}}\hfill{\footnotesize \path{lean/qtsl-aux/QtslAux/QtslCompose/AGComposeMarks.lean}, line 711}\par\nopagebreak
\begin{leansrc}
theorem allOnes_direct : JointRun allOnesNet fun _ => ![true, true]
\end{leansrc}
\par\medskip\noindent{\small \texttt{BernsteinLogic.QtslCompose.Marks.}\texttt{\bfseries fatCarve\_marks\_realized}}\hfill{\footnotesize \path{lean/qtsl-aux/QtslAux/QtslCompose/AGComposeCarves.lean}, line 666}\par\nopagebreak
\begin{leansrc}
theorem fatCarve_marks_realized :
    ∃ ω : ℕ → Fin 2 → Fin 3, JointRun relay3Net ω ∧
      ∀ i, ω ((![0, 0] : Fin 2 → ℕ) i) i ∈ relay3Mark i
\end{leansrc}
\par\medskip\noindent{\small \texttt{BernsteinLogic.QtslCompose.Marks.}\texttt{\bfseries marks\_of\_pinned\_delivery}}\hfill{\footnotesize \path{lean/qtsl-aux/QtslAux/QtslCompose/AGComposeMarks.lean}, line 287}\par\nopagebreak
\begin{leansrc}
theorem marks_of_pinned_delivery (hser : Serial N) (hfeas : MarkFlowFeasible N t)
    (hpin : PinnedDelivery N mark a) (hclose : InstantClosure N a t) :
    ∃ ω : ℕ → Fin n → V, JointRun N ω ∧ ∀ i, ω (t i) i ∈ mark i
\end{leansrc}
\par\medskip\noindent{\small \texttt{BernsteinLogic.QtslCompose.Marks.}\texttt{\bfseries pins\_of\_conditional\_delivery}}\hfill{\footnotesize \path{lean/qtsl-aux/QtslAux/QtslCompose/AGComposeMarks.lean}, line 308}\par\nopagebreak
\begin{leansrc}
theorem pins_of_conditional_delivery (hloc : DepLocal N) (ord : Fin n → ℕ)
    (hord : ∀ i, ∀ j ∈ N.dep i, ord j < ord i)
    (hdel : ∀ i, ConditionalDelivery N mark i) :
    ∃ a : Fin n → V, PinnedDelivery N mark a
\end{leansrc}
\par\medskip\noindent{\small \texttt{BernsteinLogic.QtslCompose.Marks.}\texttt{\bfseries instantClosure\_of\_acyclic}}\hfill{\footnotesize \path{lean/qtsl-aux/QtslAux/QtslCompose/AGComposeMarks.lean}, line 357}\par\nopagebreak
\begin{leansrc}
theorem instantClosure_of_acyclic (hser : Serial N) (hloc : DepLocal N)
    (ord : Fin n → ℕ) (hord : ∀ i, ∀ j ∈ N.dep i, ord j < ord i)
    (a : Fin n → V) (t : Fin n → ℕ) :
    InstantClosure N a t
\end{leansrc}
\par\medskip\noindent{\small \texttt{BernsteinLogic.QtslCompose.Marks.}\texttt{\bfseries carveClosure\_of\_acyclic}}\hfill{\footnotesize \path{lean/qtsl-aux/QtslAux/QtslCompose/AGComposeCarves.lean}, line 357}\par\nopagebreak
\begin{leansrc}
theorem carveClosure_of_acyclic (hser : Serial N) (hloc : DepLocal N)
    (ord : Fin n → ℕ) (hord : ∀ i, ∀ j ∈ N.dep i, ord j < ord i)
    (hfeas : MarkFlowFeasible N t) (hpin : CarvedDelivery N mark P) :
    CarveClosure N P t
\end{leansrc}
\par\medskip\noindent{\small \texttt{Qtsl.QtslCompose.}\texttt{\bfseries sim\_expect}}\hfill{\footnotesize \path{lean/qtsl-aux/QtslAux/QtslCompose/SimTransport.lean}, line 61}\par\nopagebreak
\begin{leansrc}
theorem sim_expect {σ : α → β} {K : α → α → ℝ} {Q : β → β → ℝ}
    (hσ : PushKernel σ K Q) (μ₀ : α → ℝ) (g : β → ℝ) (n : ℕ) :
    ∑ s, law K μ₀ n s * g (σ s) = ∑ b, law Q (push σ μ₀) n b * g b
\end{leansrc}
\par\medskip\noindent{\small \texttt{Qtsl.QtslCompose.}\texttt{\bfseries sim\_expect\_bound}}\hfill{\footnotesize \path{lean/qtsl-aux/QtslAux/QtslCompose/SimTransport.lean}, line 68}\par\nopagebreak
\begin{leansrc}
theorem sim_expect_bound {σ : α → β} {K : α → α → ℝ} {Q : β → β → ℝ}
    (hσ : PushKernel σ K Q) (μ₀ : α → ℝ) (g : β → ℝ) {n : ℕ} {b₀ : ℝ}
    (hb : ∑ b, law Q (push σ μ₀) n b * g b ≤ b₀) :
    ∑ s, law K μ₀ n s * g (σ s) ≤ b₀
\end{leansrc}
\par\medskip\noindent{\small \texttt{Qtsl.QtslCompose.}\texttt{\bfseries sim\_region}}\hfill{\footnotesize \path{lean/qtsl-aux/QtslAux/QtslCompose/SimTransport.lean}, line 77}\par\nopagebreak
\begin{leansrc}
theorem sim_region {σ : α → β} {K : α → α → ℝ} {Q : β → β → ℝ}
    (hσ : PushKernel σ K Q) {μ₀ : α → ℝ}
    (hK : ∀ s t, 0 ≤ K s t) (hμ₀ : ∀ s, 0 ≤ μ₀ s) {n : ℕ} {R' : Set β}
    (hR' : ∀ b, b ∉ R' → law Q (push σ μ₀) n b = 0) :
    ∀ s, σ s ∉ R' → law K μ₀ n s = 0
\end{leansrc}
\par\medskip\noindent{\small \texttt{Qtsl.QtslCompose.}\texttt{\bfseries expect\_push}}\hfill{\footnotesize \path{lean/qtsl-aux/QtslAux/QtslCompose/SimTransport.lean}, line 46}\par\nopagebreak
\begin{leansrc}
theorem expect_push (σ : α → β) (μ : α → ℝ) (g : β → ℝ) :
    ∑ b, push σ μ b * g b = ∑ s, μ s * g (σ s)
\end{leansrc}
\par\medskip\noindent{\small \texttt{Qtsl.QRelChain.}\texttt{\bfseries push\_law}}\hfill{\footnotesize \path{lean/qtsl-aux/QtslAux/QRelChain.lean}, line 152}\par\nopagebreak
\begin{leansrc}
theorem push_law [Fintype β] {π : α → β} {K : α → α → ℝ} {Q : β → β → ℝ}
    (h : PushKernel π K Q) (μ₀ : α → ℝ) :
    ∀ n, push π (law K μ₀ n) = law Q (push π μ₀) n
  | 0 => rfl
  | n + 1 => by rw [law_succ, push_step h, push_law h μ₀ n, law_succ]

end Push

section Drift

variable [Fintype α] {K : α → α → ℝ} {V μ₀ μ : α → ℝ}

theorem expect_step_le (hμ : ∀ s, 0 ≤ μ s)
    (hdrift : ∀ s, ∑ t, K s t * V t ≤ V s) :
    ∑ t, step K μ t * V t ≤ ∑ s, μ s * V s
\end{leansrc}
\par\medskip\noindent{\small \texttt{Qtsl.QtslCompose.Meas.}\texttt{\bfseries sim\_expect\_lintegral}}\hfill{\footnotesize \path{lean/qtsl-aux/QtslAux/QtslCompose/SimTransportMeas.lean}, line 91}\par\nopagebreak
\begin{leansrc}
theorem sim_expect_lintegral {σ : α → β} {K : Kernel α α} {Q : Kernel β β}
    (hσ : Measurable σ) (h : PushKernel σ K Q) (μ : Measure α)
    {g : β → ℝ≥0∞} (hg : Measurable g) (n : ℕ) :
    ∫⁻ s, g (σ s) ∂law K μ n = ∫⁻ b, g b ∂law Q (μ.map σ) n
\end{leansrc}
\par\medskip\noindent{\small \texttt{Qtsl.QtslCompose.Meas.}\texttt{\bfseries pushKernel\_iff\_comp}}\hfill{\footnotesize \path{lean/qtsl-aux/QtslAux/QtslCompose/SimTransportMeas.lean}, line 132}\par\nopagebreak
\begin{leansrc}
theorem pushKernel_iff_comp {σ : α → β} (hσ : Measurable σ)
    {K : Kernel α α} {Q : Kernel β β} :
    PushKernel σ K Q ↔
      Kernel.deterministic σ hσ ∘ₖ K = Q ∘ₖ Kernel.deterministic σ hσ
\end{leansrc}
\par\medskip\noindent{\small \texttt{Qtsl.QRelChain.}\texttt{\bfseries law\_nonneg}}\hfill{\footnotesize \path{lean/qtsl-aux/QtslAux/QRelChain.lean}, line 85}\par\nopagebreak
\begin{leansrc}
theorem law_nonneg (hK : ∀ s t, 0 ≤ K s t) (hμ₀ : ∀ s, 0 ≤ μ₀ s) :
    ∀ n s, 0 ≤ law K μ₀ n s
  | 0, s => hμ₀ s
  | n + 1, t => step_nonneg hK (law_nonneg hK hμ₀ n) t

theorem sum_step (hrow : ∀ s, ∑ t, K s t = 1) : ∑ t, step K μ t = ∑ s, μ s
\end{leansrc}
\par\medskip\noindent{\small \texttt{Qtsl.QRelChain.}\texttt{\bfseries PushKernel}}\hfill{\footnotesize \path{lean/qtsl-aux/QtslAux/QRelChain.lean}, line 130}\par\nopagebreak
\begin{leansrc}
def PushKernel (π : α → β) (K : α → α → ℝ) (Q : β → β → ℝ) : Prop
\end{leansrc}
\par\medskip\noindent{\small \texttt{Qtsl.QtslCompose.}\texttt{\bfseries law\_prodKernel}}\hfill{\footnotesize \path{lean/qtsl-aux/QtslAux/QtslCompose/Flatten.lean}, line 94}\par\nopagebreak
\begin{leansrc}
theorem law_prodKernel (K₁ : α → α → ℝ) (K₂ : β → β → ℝ)
    (μ₁ : α → ℝ) (μ₂ : β → ℝ) :
    ∀ n, law (prodKernel K₁ K₂) (prodDist μ₁ μ₂) n
      = prodDist (law K₁ μ₁ n) (law K₂ μ₂ n)
  | 0 => rfl
  | n + 1 => by
    rw [law_succ, law_prodKernel K₁ K₂ μ₁ μ₂ n, step_prodKernel, law_succ,
      law_succ]

section Projections

variable [DecidableEq α] [DecidableEq β]

theorem pushKernel_fst {K₁ : α → α → ℝ} {K₂ : β → β → ℝ}
    (h₂ : ∀ b, ∑ b', K₂ b b' = 1) :
    PushKernel Prod.fst (prodKernel K₁ K₂) K₁
\end{leansrc}
\par\medskip\noindent{\small \texttt{Qtsl.QtslCompose.}\texttt{\bfseries pushKernel\_fst}}\hfill{\footnotesize \path{lean/qtsl-aux/QtslAux/QtslCompose/Flatten.lean}, line 109}\par\nopagebreak
\begin{leansrc}
theorem pushKernel_fst {K₁ : α → α → ℝ} {K₂ : β → β → ℝ}
    (h₂ : ∀ b, ∑ b', K₂ b b' = 1) :
    PushKernel Prod.fst (prodKernel K₁ K₂) K₁
\end{leansrc}
\par\medskip\noindent{\small \texttt{Qtsl.QtslCompose.}\texttt{\bfseries prodKernel}}\hfill{\footnotesize \path{lean/qtsl-aux/QtslAux/QtslCompose/Flatten.lean}, line 34}\par\nopagebreak
\begin{leansrc}
def prodKernel (K₁ : α → α → ℝ) (K₂ : β → β → ℝ) :
    α × β → α × β → ℝ
\end{leansrc}
\par\medskip\noindent{\small \texttt{Qtsl.QtslCompose.Meas.}\texttt{\bfseries law\_parallelComp}}\hfill{\footnotesize \path{lean/qtsl-aux/QtslAux/QtslCompose/FlattenMeas.lean}, line 75}\par\nopagebreak
\begin{leansrc}
theorem law_parallelComp (μ : Measure α) (ν : Measure β)
    (K : Kernel α α) (Q : Kernel β β)
    [SFinite μ] [SFinite ν] [IsSFiniteKernel K] [IsSFiniteKernel Q] :
    ∀ n, law (K ∥ₖ Q) (μ.prod ν) n = (law K μ n).prod (law Q ν n)
  | 0 => rfl
  | n + 1 => by
    rw [law_succ, law_parallelComp μ ν K Q n, bind_parallelComp,
      ← law_succ, ← law_succ]

theorem pushKernel_fst {K : Kernel α α} {Q : Kernel β β}
    [IsSFiniteKernel K] [IsMarkovKernel Q] :
    PushKernel Prod.fst (K ∥ₖ Q) K
\end{leansrc}
\par\medskip\noindent{\small \texttt{Qtsl.QtslCompose.Meas.}\texttt{\bfseries PushKernel}}\hfill{\footnotesize \path{lean/qtsl-aux/QtslAux/QtslCompose/SimTransportMeas.lean}, line 43}\par\nopagebreak
\begin{leansrc}
def PushKernel (σ : α → β) (K : Kernel α α) (Q : Kernel β β) : Prop
\end{leansrc}
\par\medskip\noindent{\small \texttt{Qtsl.QtslCompose.Meas.}\texttt{\bfseries bind\_parallelComp}}\hfill{\footnotesize \path{lean/qtsl-aux/QtslAux/QtslCompose/FlattenMeas.lean}, line 58}\par\nopagebreak
\begin{leansrc}
theorem bind_parallelComp (μ : Measure α) (ν : Measure β)
    (K : Kernel α α) (Q : Kernel β β)
    [SFinite μ] [SFinite ν] [IsSFiniteKernel K] [IsSFiniteKernel Q] :
    (μ.prod ν).bind (K ∥ₖ Q) = (μ.bind K).prod (ν.bind Q)
\end{leansrc}
\par\medskip\noindent{\small \texttt{Qtsl.QtslCompose.}\texttt{\bfseries isMarkov\_prodKernel}}\hfill{\footnotesize \path{lean/qtsl-aux/QtslAux/QtslCompose/Flatten.lean}, line 45}\par\nopagebreak
\begin{leansrc}
theorem isMarkov_prodKernel {K₁ : α → α → ℝ} {K₂ : β → β → ℝ}
    (h₁ : IsMarkov K₁) (h₂ : IsMarkov K₂) :
    IsMarkov (prodKernel K₁ K₂) where
  nonneg
\end{leansrc}
\par\medskip\noindent{\small \texttt{Qtsl.QtslCompose.}\texttt{\bfseries isDist\_prodDist}}\hfill{\footnotesize \path{lean/qtsl-aux/QtslAux/QtslCompose/Flatten.lean}, line 61}\par\nopagebreak
\begin{leansrc}
theorem isDist_prodDist {μ₁ : α → ℝ} {μ₂ : β → ℝ}
    (h₁ : IsDist μ₁) (h₂ : IsDist μ₂) : IsDist (prodDist μ₁ μ₂) where
  nonneg
\end{leansrc}
\par\medskip\noindent{\small \texttt{Qtsl.QtslCompose.}\texttt{\bfseries push\_fst\_prodDist}}\hfill{\footnotesize \path{lean/qtsl-aux/QtslAux/QtslCompose/Flatten.lean}, line 135}\par\nopagebreak
\begin{leansrc}
theorem push_fst_prodDist {μ₁ : α → ℝ} {μ₂ : β → ℝ} (h₂ : IsDist μ₂) :
    push Prod.fst (prodDist μ₁ μ₂) = μ₁
\end{leansrc}
\par\medskip\noindent{\small \texttt{Qtsl.QtslCompose.Meas.}\texttt{\bfseries sfinite\_law}}\hfill{\footnotesize \path{lean/qtsl-aux/QtslAux/QtslCompose/FlattenMeas.lean}, line 38}\par\nopagebreak
\begin{leansrc}
instance sfinite_law (K : Kernel α α) [IsSFiniteKernel K]
    (μ : Measure α) [SFinite μ] : ∀ n, SFinite (law K μ n)
  | 0 => inferInstanceAs (SFinite μ)
  | n + 1 => by
    have
\end{leansrc}
\par\medskip\noindent{\small \texttt{Qtsl.QtslCompose.Meas.}\texttt{\bfseries isProbabilityMeasure\_law}}\hfill{\footnotesize \path{lean/qtsl-aux/QtslAux/QtslCompose/FlattenMeas.lean}, line 47}\par\nopagebreak
\begin{leansrc}
instance isProbabilityMeasure_law (K : Kernel α α) [IsMarkovKernel K]
    (μ : Measure α) [IsProbabilityMeasure μ] :
    ∀ n, IsProbabilityMeasure (law K μ n)
  | 0 => inferInstanceAs (IsProbabilityMeasure μ)
  | n + 1 => by
    have
\end{leansrc}
\par\medskip\noindent{\small \texttt{Qtsl.QtslCompose.Meas.}\texttt{\bfseries sim\_expect\_fst}}\hfill{\footnotesize \path{lean/qtsl-aux/QtslAux/QtslCompose/FlattenMeas.lean}, line 106}\par\nopagebreak
\begin{leansrc}
theorem sim_expect_fst {K : Kernel α α} {Q : Kernel β β}
    [IsSFiniteKernel K] [IsMarkovKernel Q]
    (μ : Measure α) (ν : Measure β) [IsProbabilityMeasure ν]
    {g : α → ℝ≥0∞} (hg : Measurable g) (n : ℕ) :
    ∫⁻ p, g p.1 ∂law (K ∥ₖ Q) (μ.prod ν) n = ∫⁻ a, g a ∂law K μ n
\end{leansrc}
\par\medskip\noindent{\small \texttt{Qtsl.QtslCompose.Meas.}\texttt{\bfseries pushKernel\_iff\_stochSim}}\hfill{\footnotesize \path{lean/qtsl-aux/QtslAux/QtslCompose/FlattenMeas.lean}, line 128}\par\nopagebreak
\begin{leansrc}
theorem pushKernel_iff_stochSim {σ : α → β} {K : Kernel α α}
    {Q : Kernel β β} : PushKernel σ K Q ↔ StochSim σ K Q
\end{leansrc}
\par\medskip\noindent{\small \texttt{Qtsl.QtslCompose.}\texttt{\bfseries prodDist}}\hfill{\footnotesize \path{lean/qtsl-aux/QtslAux/QtslCompose/Flatten.lean}, line 40}\par\nopagebreak
\begin{leansrc}
def prodDist (μ₁ : α → ℝ) (μ₂ : β → ℝ) : α × β → ℝ
\end{leansrc}
\par\medskip\noindent{\small \texttt{Qtsl.QtslCompose.}\texttt{\bfseries step\_prodKernel}}\hfill{\footnotesize \path{lean/qtsl-aux/QtslAux/QtslCompose/Flatten.lean}, line 75}\par\nopagebreak
\begin{leansrc}
theorem step_prodKernel (K₁ : α → α → ℝ) (K₂ : β → β → ℝ)
    (μ₁ : α → ℝ) (μ₂ : β → ℝ) :
    step (prodKernel K₁ K₂) (prodDist μ₁ μ₂)
      = prodDist (step K₁ μ₁) (step K₂ μ₂)
\end{leansrc}
\par\medskip\noindent{\small \texttt{Qtsl.QtslCompose.}\texttt{\bfseries pushKernel\_snd}}\hfill{\footnotesize \path{lean/qtsl-aux/QtslAux/QtslCompose/Flatten.lean}, line 122}\par\nopagebreak
\begin{leansrc}
theorem pushKernel_snd {K₁ : α → α → ℝ} {K₂ : β → β → ℝ}
    (h₁ : ∀ a, ∑ a', K₁ a a' = 1) :
    PushKernel Prod.snd (prodKernel K₁ K₂) K₂
\end{leansrc}
\par\medskip\noindent{\small \texttt{Qtsl.QDistTransport.}\texttt{\bfseries strict\_cell\_exact}}\hfill{\footnotesize \path{lean/qtsl-aux/QtslAux/QDistTransport.lean}, line 193}\par\nopagebreak
\begin{leansrc}
theorem strict_cell_exact [DecidableEq β] {π : α → β} {K : α → α → ℝ}
    {Q : β → β → ℝ} (h : PushKernel π K Q) (μ₀ : α → ℝ) (n : ℕ) :
    push π (law K μ₀ n) = law Q (push π μ₀) n
\end{leansrc}
\par\medskip\noindent{\small \texttt{Qtsl.QtslCompose.}\texttt{\bfseries aggReduce}}\hfill{\footnotesize \path{lean/qtsl-aux/QtslAux/QtslCompose/AggReduce.lean}, line 37}\par\nopagebreak
\begin{leansrc}
theorem aggReduce {X A : Type*} (agg : X → A) (scope : Set X) (K : Set A)
    (G : A → ℝ) (hG : ∀ a ∈ K, G a ≤ 0) (hin : ∀ x ∈ scope, agg x ∈ K) :
    ∀ x ∈ scope, G (agg x) ≤ 0
\end{leansrc}
\par\medskip\noindent{\small \texttt{Qtsl.QtslCompose.}\texttt{\bfseries coupling\_row}}\hfill{\footnotesize \path{lean/qtsl-aux/QtslAux/QtslCompose/AggReduce.lean}, line 80}\par\nopagebreak
\begin{leansrc}
theorem coupling_row {m : Multiset C} {f : C → ℝ} {l u : ℝ}
    (hl : ∀ c ∈ m, l ≤ f c) (hu : ∀ c ∈ m, f c ≤ u) :
    l * m.card ≤ query m f ∧ query m f ≤ u * m.card
\end{leansrc}
\par\medskip\noindent{\small \texttt{Qtsl.QtslCompose.}\texttt{\bfseries count\_row\_nonneg}}\hfill{\footnotesize \path{lean/qtsl-aux/QtslAux/QtslCompose/AggReduce.lean}, line 96}\par\nopagebreak
\begin{leansrc}
theorem count_row_nonneg (m : Multiset C) : (0 : ℝ) ≤ m.card
\end{leansrc}
\par\medskip\noindent{\small \texttt{Qtsl.QtslCompose.}\texttt{\bfseries count\_row\_le}}\hfill{\footnotesize \path{lean/qtsl-aux/QtslAux/QtslCompose/AggReduce.lean}, line 99}\par\nopagebreak
\begin{leansrc}
theorem count_row_le {m : Multiset C} {M : ℕ} (h : m.card ≤ M) :
    (m.card : ℝ) ≤ M
\end{leansrc}
\par\medskip\noindent{\small \texttt{Qtsl.QtslCompose.}\texttt{\bfseries moment\_row}}\hfill{\footnotesize \path{lean/qtsl-aux/QtslAux/QtslCompose/AggReduce.lean}, line 105}\par\nopagebreak
\begin{leansrc}
theorem moment_row (m : Multiset C) (f : C → ℝ) :
    (query m f) ^ 2 ≤ m.card * query m (fun c => f c ^ 2)
\end{leansrc}
\par\medskip\noindent{\small \texttt{Qtsl.QtslCompose.}\texttt{\bfseries query\_count}}\hfill{\footnotesize \path{lean/qtsl-aux/QtslAux/QtslCompose/AggReduce.lean}, line 66}\par\nopagebreak
\begin{leansrc}
theorem query_count [DecidableEq C] (m : Multiset C) (f : C → ℝ) :
    query m f = ∑ c ∈ m.toFinset, (m.count c : ℝ) * f c
\end{leansrc}
\par\medskip\noindent{\small \texttt{Qtsl.QtslCompose.}\texttt{\bfseries query\_const\_one}}\hfill{\footnotesize \path{lean/qtsl-aux/QtslAux/QtslCompose/AggReduce.lean}, line 55}\par\nopagebreak
\begin{leansrc}
theorem query_const_one (m : Multiset C) :
    query m (fun _ => (1 : ℝ)) = m.card
\end{leansrc}
\par\medskip\noindent{\small \texttt{Qtsl.Structural.}\texttt{\bfseries countAtomRounding\_sound}}\hfill{\footnotesize \path{lean/qtsl-aux/QtslAux/Structural.lean}, line 236}\par\nopagebreak
\begin{leansrc}
theorem countAtomRounding_sound (k : X → ℝ) (j : ℤ) (Reach : Set X)
    (hbridge : ∀ x ∈ Reach, ∃ n : ℤ, k x = (n : ℝ)) :
    ∀ x ∈ Reach, x ∈ {x | k x ≤ (j : ℝ)}ᶜ → (j : ℝ) + 1 ≤ k x
\end{leansrc}
\par\medskip\noindent{\small \texttt{Qtsl.Structural.}\texttt{\bfseries countAtomRounding\_tighter}}\hfill{\footnotesize \path{lean/qtsl-aux/QtslAux/Structural.lean}, line 250}\par\nopagebreak
\begin{leansrc}
theorem countAtomRounding_tighter (k : X → ℝ) (j : ℤ) :
    {x | (j : ℝ) + 1 ≤ k x} ⊆ {x | (j : ℝ) < k x}
\end{leansrc}
\par\medskip\noindent{\small \texttt{Qtsl.Structural.}\texttt{\bfseries countAtomRounding\_eq}}\hfill{\footnotesize \path{lean/qtsl-aux/QtslAux/Structural.lean}, line 259}\par\nopagebreak
\begin{leansrc}
theorem countAtomRounding_eq (k : X → ℝ) (j : ℤ) (Reach : Set X)
    (hbridge : ∀ x ∈ Reach, ∃ n : ℤ, k x = (n : ℝ)) :
    Reach ∩ {x | k x ≤ (j : ℝ)}ᶜ = Reach ∩ {x | (j : ℝ) + 1 ≤ k x}
\end{leansrc}
\par\medskip\noindent{\small \texttt{Qtsl.QRelChain.}\texttt{\bfseries sound\_QRelTV}}\hfill{\footnotesize \path{lean/qtsl-aux/QtslAux/QRelChain.lean}, line 736}\par\nopagebreak
\begin{leansrc}
theorem sound_QRelTV [Fintype L] [DecidableEq L] {f : σ₁ → L} {g : σ₂ → L}
    {V : σ₁ × σ₂ → ℝ} {b : ℝ}
    (hK : IsMarkov K) (hμ₀ : IsDist μ₀) (hK₁ : IsMarkov K₁) (hK₂ : IsMarkov K₂)
    (hm₁ : PushKernel Prod.fst K K₁) (hm₂ : PushKernel Prod.snd K K₂)
    (hfreeze : ∀ p ∈ A₁ ×ˢ A₂, K p p = 1)
    (hfr₁ : ∀ a ∈ A₁, K₁ a a = 1) (hfr₂ : ∀ a ∈ A₂, K₂ a a = 1)
    (habs₁ : Tendsto (fun n => ∑ a ∈ A₁, law K₁ (push Prod.fst μ₀) n a)
      atTop (𝓝 1))
    (habs₂ : Tendsto (fun n => ∑ a ∈ A₂, law K₂ (push Prod.snd μ₀) n a)
      atTop (𝓝 1))
    (hV0 : ∀ p, 0 ≤ V p) (hdrift : ∀ p, ∑ q, K p q * V q ≤ V p)
    (hterm : ∀ p ∈ A₁ ×ˢ A₂, f p.1 ≠ g p.2 → 1 ≤ V p)
    (hb : ∑ p, μ₀ p * V p ≤ b) :
    tvDist (push f (finalLaw K₁ (push Prod.fst μ₀) A₁))
      (push g (finalLaw K₂ (push Prod.snd μ₀) A₂)) ≤ b
\end{leansrc}
\par\medskip\noindent{\small \texttt{Qtsl.QRelChain.}\texttt{\bfseries sound\_QRelExpGap}}\hfill{\footnotesize \path{lean/qtsl-aux/QtslAux/QRelChain.lean}, line 777}\par\nopagebreak
\begin{leansrc}
theorem sound_QRelExpGap {v₁ : σ₁ → ℝ} {v₂ : σ₂ → ℝ} {V : σ₁ × σ₂ → ℝ} {b : ℝ}
    (hK : IsMarkov K) (hμ₀ : IsDist μ₀) (hK₁ : IsMarkov K₁) (hK₂ : IsMarkov K₂)
    (hm₁ : PushKernel Prod.fst K K₁) (hm₂ : PushKernel Prod.snd K K₂)
    (hfreeze : ∀ p ∈ A₁ ×ˢ A₂, K p p = 1)
    (hfr₁ : ∀ a ∈ A₁, K₁ a a = 1) (hfr₂ : ∀ a ∈ A₂, K₂ a a = 1)
    (habs₁ : Tendsto (fun n => ∑ a ∈ A₁, law K₁ (push Prod.fst μ₀) n a)
      atTop (𝓝 1))
    (habs₂ : Tendsto (fun n => ∑ a ∈ A₂, law K₂ (push Prod.snd μ₀) n a)
      atTop (𝓝 1))
    (hV0 : ∀ p, 0 ≤ V p) (hdrift : ∀ p, ∑ q, K p q * V q ≤ V p)
    (hterm : ∀ p ∈ A₁ ×ˢ A₂, |v₁ p.1 - v₂ p.2| ≤ V p)
    (hb : ∑ p, μ₀ p * V p ≤ b) :
    |(∑ a₁ ∈ A₁, finalMass K₁ (push Prod.fst μ₀) a₁ * v₁ a₁) -
      ∑ a₂ ∈ A₂, finalMass K₂ (push Prod.snd μ₀) a₂ * v₂ a₂| ≤ b
\end{leansrc}
\par\medskip\noindent{\small \texttt{Qtsl.QRelChain.}\texttt{\bfseries tendsto\_absorbedMass\_one}}\hfill{\footnotesize \path{lean/qtsl-aux/QtslAux/QRelChain.lean}, line 392}\par\nopagebreak
\begin{leansrc}
theorem tendsto_absorbedMass_one (hK : IsMarkov K) (hμ₀ : IsDist μ₀)
    (hfr : ∀ a ∈ A, K a a = 1)
    (hreach : ∀ s, ∃ k, 0 < ∑ a ∈ A, law K (pointMass s) k a) :
    Tendsto (fun n => ∑ a ∈ A, law K μ₀ n a) atTop (𝓝 1)
\end{leansrc}
\par\medskip\noindent{\small \texttt{Qtsl.QRelChain.}\texttt{\bfseries tvDist}}\hfill{\footnotesize \path{lean/qtsl-aux/QtslAux/QRelChain.lean}, line 453}\par\nopagebreak
\begin{leansrc}
def tvDist [Fintype L] (p q : L → ℝ) : ℝ
\end{leansrc}
\par\medskip\noindent{\small \texttt{Qtsl.QRelChain.}\texttt{\bfseries push\_step}}\hfill{\footnotesize \path{lean/qtsl-aux/QtslAux/QRelChain.lean}, line 133}\par\nopagebreak
\begin{leansrc}
theorem push_step [Fintype β] {π : α → β} {K : α → α → ℝ} {Q : β → β → ℝ}
    (h : PushKernel π K Q) (μ : α → ℝ) :
    push π (step K μ) = step Q (push π μ)
\end{leansrc}
\par\medskip\noindent{\small \texttt{Qtsl.QRelChain.}\texttt{\bfseries expect\_step\_le}}\hfill{\footnotesize \path{lean/qtsl-aux/QtslAux/QRelChain.lean}, line 164}\par\nopagebreak
\begin{leansrc}
theorem expect_step_le (hμ : ∀ s, 0 ≤ μ s)
    (hdrift : ∀ s, ∑ t, K s t * V t ≤ V s) :
    ∑ t, step K μ t * V t ≤ ∑ s, μ s * V s
\end{leansrc}
\par\medskip\noindent{\small \texttt{Qtsl.QRelChain.}\texttt{\bfseries expect\_law\_le}}\hfill{\footnotesize \path{lean/qtsl-aux/QtslAux/QRelChain.lean}, line 175}\par\nopagebreak
\begin{leansrc}
theorem expect_law_le (hK : ∀ s t, 0 ≤ K s t) (hμ₀ : ∀ s, 0 ≤ μ₀ s)
    (hdrift : ∀ s, ∑ t, K s t * V t ≤ V s) :
    ∀ n, ∑ s, law K μ₀ n s * V s ≤ ∑ s, μ₀ s * V s
  | 0 => le_rfl
  | n + 1 =>
    (expect_step_le (law_nonneg hK hμ₀ n) hdrift).trans
      (expect_law_le hK hμ₀ hdrift n)

end Drift

section Absorb

variable [Fintype α] {K : α → α → ℝ} {μ₀ : α → ℝ}

theorem law_mono_absorbing (hK : ∀ s t, 0 ≤ K s t) (hμ₀ : ∀ s, 0 ≤ μ₀ s)
    {a : α} (haa : K a a = 1) :
    Monotone fun n => law K μ₀ n a
\end{leansrc}
\par\medskip\noindent{\small \texttt{Qtsl.QRelChain.}\texttt{\bfseries law\_mono\_absorbing}}\hfill{\footnotesize \path{lean/qtsl-aux/QtslAux/QRelChain.lean}, line 189}\par\nopagebreak
\begin{leansrc}
theorem law_mono_absorbing (hK : ∀ s t, 0 ≤ K s t) (hμ₀ : ∀ s, 0 ≤ μ₀ s)
    {a : α} (haa : K a a = 1) :
    Monotone fun n => law K μ₀ n a
\end{leansrc}
\par\medskip\noindent{\small \texttt{Qtsl.QRelChain.}\texttt{\bfseries finalMass}}\hfill{\footnotesize \path{lean/qtsl-aux/QtslAux/QRelChain.lean}, line 198}\par\nopagebreak
\begin{leansrc}
def finalMass (K : α → α → ℝ) (μ₀ : α → ℝ) (a : α) : ℝ
\end{leansrc}
\par\medskip\noindent{\small \texttt{Qtsl.QRelChain.}\texttt{\bfseries finalExpect\_le}}\hfill{\footnotesize \path{lean/qtsl-aux/QtslAux/QRelChain.lean}, line 226}\par\nopagebreak
\begin{leansrc}
theorem finalExpect_le {V gv : α → ℝ} {A : Finset α} {b : ℝ}
    (hK : IsMarkov K) (hμ₀ : IsDist μ₀) (hfr : ∀ a ∈ A, K a a = 1)
    (hgV : ∀ a ∈ A, gv a ≤ V a) (hV0 : ∀ s, 0 ≤ V s)
    (hdrift : ∀ s, ∑ t, K s t * V t ≤ V s) (hb : ∑ s, μ₀ s * V s ≤ b) :
    ∑ a ∈ A, finalMass K μ₀ a * gv a ≤ b
\end{leansrc}
\par\medskip\noindent{\small \texttt{Qtsl.QRelChain.}\texttt{\bfseries tvDist}}\hfill{\footnotesize \path{lean/qtsl-aux/QtslAux/QRelChain.lean}, line 453}\par\nopagebreak
\begin{leansrc}
def tvDist [Fintype L] (p q : L → ℝ) : ℝ
\end{leansrc}
\par\medskip\noindent{\small \texttt{Qtsl.QRelChain.}\texttt{\bfseries tvDist\_push\_le\_disagree}}\hfill{\footnotesize \path{lean/qtsl-aux/QtslAux/QRelChain.lean}, line 455}\par\nopagebreak
\begin{leansrc}
theorem tvDist_push_le_disagree {σ₁ σ₂ : Type*} [Fintype σ₁] [Fintype σ₂]
    [Fintype L] [DecidableEq L] (γ : σ₁ × σ₂ → ℝ) (hγ : ∀ p, 0 ≤ γ p)
    (f : σ₁ → L) (g : σ₂ → L) :
    tvDist (push (fun p : σ₁ × σ₂ => f p.1) γ) (push (fun p : σ₁ × σ₂ => g p.2) γ) ≤
      ∑ p ∈ univ.filter fun p : σ₁ × σ₂ => f p.1 ≠ g p.2, γ p
\end{leansrc}
\par\medskip\noindent{\small \texttt{Qtsl.QRelChain.}\texttt{\bfseries sum\_finalMass\_snd\_eq}}\hfill{\footnotesize \path{lean/qtsl-aux/QtslAux/QRelChain.lean}, line 532}\par\nopagebreak
\begin{leansrc}
theorem sum_finalMass_snd_eq (hK : IsMarkov K) (hμ₀ : IsDist μ₀)
    (hK₁ : IsMarkov K₁) (hK₂ : IsMarkov K₂)
    (hm₁ : PushKernel Prod.fst K K₁) (hm₂ : PushKernel Prod.snd K K₂)
    (hfreeze : ∀ p ∈ A₁ ×ˢ A₂, K p p = 1) (hfr₁ : ∀ a ∈ A₁, K₁ a a = 1)
    (habs₂ : Tendsto (fun n => ∑ a ∈ A₂, law K₂ (push Prod.snd μ₀) n a)
      atTop (𝓝 1))
    {a₁ : σ₁} (ha₁ : a₁ ∈ A₁) :
    ∑ a₂ ∈ A₂, finalMass K μ₀ (a₁, a₂) = finalMass K₁ (push Prod.fst μ₀) a₁
\end{leansrc}
\par\medskip\noindent{\small \texttt{Qtsl.QRelChain.}\texttt{\bfseries sum\_finalMass\_fst\_eq}}\hfill{\footnotesize \path{lean/qtsl-aux/QtslAux/QRelChain.lean}, line 585}\par\nopagebreak
\begin{leansrc}
theorem sum_finalMass_fst_eq (hK : IsMarkov K) (hμ₀ : IsDist μ₀)
    (hK₁ : IsMarkov K₁) (hK₂ : IsMarkov K₂)
    (hm₁ : PushKernel Prod.fst K K₁) (hm₂ : PushKernel Prod.snd K K₂)
    (hfreeze : ∀ p ∈ A₁ ×ˢ A₂, K p p = 1) (hfr₂ : ∀ a ∈ A₂, K₂ a a = 1)
    (habs₁ : Tendsto (fun n => ∑ a ∈ A₁, law K₁ (push Prod.fst μ₀) n a)
      atTop (𝓝 1))
    {a₂ : σ₂} (ha₂ : a₂ ∈ A₂) :
    ∑ a₁ ∈ A₁, finalMass K μ₀ (a₁, a₂) = finalMass K₂ (push Prod.snd μ₀) a₂
\end{leansrc}
\par\medskip\noindent{\small \texttt{Qtsl.QRelChain.}\texttt{\bfseries sum\_finalMass\_prod\_eq\_one}}\hfill{\footnotesize \path{lean/qtsl-aux/QtslAux/QRelChain.lean}, line 638}\par\nopagebreak
\begin{leansrc}
theorem sum_finalMass_prod_eq_one (hK : IsMarkov K) (hμ₀ : IsDist μ₀)
    (hK₁ : IsMarkov K₁) (hK₂ : IsMarkov K₂)
    (hm₁ : PushKernel Prod.fst K K₁) (hm₂ : PushKernel Prod.snd K K₂)
    (hfreeze : ∀ p ∈ A₁ ×ˢ A₂, K p p = 1) (hfr₁ : ∀ a ∈ A₁, K₁ a a = 1)
    (habs₁ : Tendsto (fun n => ∑ a ∈ A₁, law K₁ (push Prod.fst μ₀) n a)
      atTop (𝓝 1))
    (habs₂ : Tendsto (fun n => ∑ a ∈ A₂, law K₂ (push Prod.snd μ₀) n a)
      atTop (𝓝 1)) :
    ∑ p ∈ A₁ ×ˢ A₂, finalMass K μ₀ p = 1
\end{leansrc}
\par\medskip\noindent{\small \texttt{Qtsl.QRelChain.}\texttt{\bfseries push\_finalLaw\_fst}}\hfill{\footnotesize \path{lean/qtsl-aux/QtslAux/QRelChain.lean}, line 667}\par\nopagebreak
\begin{leansrc}
theorem push_finalLaw_fst [Fintype L] [DecidableEq L]
    (hK : IsMarkov K) (hμ₀ : IsDist μ₀) (hK₁ : IsMarkov K₁) (hK₂ : IsMarkov K₂)
    (hm₁ : PushKernel Prod.fst K K₁) (hm₂ : PushKernel Prod.snd K K₂)
    (hfreeze : ∀ p ∈ A₁ ×ˢ A₂, K p p = 1) (hfr₁ : ∀ a ∈ A₁, K₁ a a = 1)
    (habs₂ : Tendsto (fun n => ∑ a ∈ A₂, law K₂ (push Prod.snd μ₀) n a)
      atTop (𝓝 1))
    (f : σ₁ → L) :
    push (fun p : σ₁ × σ₂ => f p.1) (finalLaw K μ₀ (A₁ ×ˢ A₂)) =
      push f (finalLaw K₁ (push Prod.fst μ₀) A₁)
\end{leansrc}
\par\medskip\noindent{\small \texttt{Qtsl.QRelChain.}\texttt{\bfseries sound\_QRelTV\_of\_reach}}\hfill{\footnotesize \path{lean/qtsl-aux/QtslAux/QRelChain.lean}, line 815}\par\nopagebreak
\begin{leansrc}
theorem sound_QRelTV_of_reach [Fintype L] [DecidableEq L] {f : σ₁ → L}
    {g : σ₂ → L} {V : σ₁ × σ₂ → ℝ} {b : ℝ}
    (hK : IsMarkov K) (hμ₀ : IsDist μ₀) (hK₁ : IsMarkov K₁) (hK₂ : IsMarkov K₂)
    (hm₁ : PushKernel Prod.fst K K₁) (hm₂ : PushKernel Prod.snd K K₂)
    (hfreeze : ∀ p ∈ A₁ ×ˢ A₂, K p p = 1)
    (hfr₁ : ∀ a ∈ A₁, K₁ a a = 1) (hfr₂ : ∀ a ∈ A₂, K₂ a a = 1)
    (hreach₁ : ∀ s, ∃ k, 0 < ∑ a ∈ A₁, law K₁ (pointMass s) k a)
    (hreach₂ : ∀ s, ∃ k, 0 < ∑ a ∈ A₂, law K₂ (pointMass s) k a)
    (hV0 : ∀ p, 0 ≤ V p) (hdrift : ∀ p, ∑ q, K p q * V q ≤ V p)
    (hterm : ∀ p ∈ A₁ ×ˢ A₂, f p.1 ≠ g p.2 → 1 ≤ V p)
    (hb : ∑ p, μ₀ p * V p ≤ b) :
    tvDist (push f (finalLaw K₁ (push Prod.fst μ₀) A₁))
      (push g (finalLaw K₂ (push Prod.snd μ₀) A₂)) ≤ b
\end{leansrc}
\par\medskip\noindent{\small \texttt{Qtsl.QRelChain.}\texttt{\bfseries sound\_QRelExpGap\_of\_reach}}\hfill{\footnotesize \path{lean/qtsl-aux/QtslAux/QRelChain.lean}, line 833}\par\nopagebreak
\begin{leansrc}
theorem sound_QRelExpGap_of_reach {v₁ : σ₁ → ℝ} {v₂ : σ₂ → ℝ}
    {V : σ₁ × σ₂ → ℝ} {b : ℝ}
    (hK : IsMarkov K) (hμ₀ : IsDist μ₀) (hK₁ : IsMarkov K₁) (hK₂ : IsMarkov K₂)
    (hm₁ : PushKernel Prod.fst K K₁) (hm₂ : PushKernel Prod.snd K K₂)
    (hfreeze : ∀ p ∈ A₁ ×ˢ A₂, K p p = 1)
    (hfr₁ : ∀ a ∈ A₁, K₁ a a = 1) (hfr₂ : ∀ a ∈ A₂, K₂ a a = 1)
    (hreach₁ : ∀ s, ∃ k, 0 < ∑ a ∈ A₁, law K₁ (pointMass s) k a)
    (hreach₂ : ∀ s, ∃ k, 0 < ∑ a ∈ A₂, law K₂ (pointMass s) k a)
    (hV0 : ∀ p, 0 ≤ V p) (hdrift : ∀ p, ∑ q, K p q * V q ≤ V p)
    (hterm : ∀ p ∈ A₁ ×ˢ A₂, |v₁ p.1 - v₂ p.2| ≤ V p)
    (hb : ∑ p, μ₀ p * V p ≤ b) :
    |(∑ a₁ ∈ A₁, finalMass K₁ (push Prod.fst μ₀) a₁ * v₁ a₁) -
      ∑ a₂ ∈ A₂, finalMass K₂ (push Prod.snd μ₀) a₂ * v₂ a₂| ≤ b
\end{leansrc}
\par\medskip\noindent{\small \texttt{Qtsl.QRelChain.}\texttt{\bfseries exists\_uniform\_reach}}\hfill{\footnotesize \path{lean/qtsl-aux/QtslAux/QRelChain.lean}, line 334}\par\nopagebreak
\begin{leansrc}
theorem exists_uniform_reach (hK : IsMarkov K) (hfr : ∀ a ∈ A, K a a = 1)
    (hreach : ∀ s, ∃ k, 0 < ∑ a ∈ A, law K (pointMass s) k a) :
    ∃ N ε, 0 < ε ∧ 1 ≤ N ∧ ∀ s, ε ≤ ∑ a ∈ A, law K (pointMass s) N a
\end{leansrc}
\par\medskip\noindent{\small \texttt{Qtsl.QRelChain.}\texttt{\bfseries absorbed\_decay}}\hfill{\footnotesize \path{lean/qtsl-aux/QtslAux/QRelChain.lean}, line 354}\par\nopagebreak
\begin{leansrc}
theorem absorbed_decay {μ : α → ℝ} (hK : IsMarkov K) (hμ : IsDist μ)
    (hfr : ∀ a ∈ A, K a a = 1) {N : ℕ} {ε : ℝ}
    (hε : ∀ s, ε ≤ ∑ a ∈ A, law K (pointMass s) N a) :
    1 - ∑ a ∈ A, law K μ N a ≤
      (1 - ε) * (1 - ∑ s ∈ A, μ s)
\end{leansrc}
\par\medskip\noindent{\small \texttt{Qtsl.QFixBall.}\texttt{\bfseries sound\_QFixBall}}\hfill{\footnotesize \path{lean/qtsl-aux/QtslAux/QFixBall.lean}, line 88}\par\nopagebreak
\begin{leansrc}
theorem sound_QFixBall [CompleteSpace α] (hK : K < 1)
    (hlip : LipschitzOnWith K f (closedBall x₀ r))
    (hres : dist (f x₀) x₀ ≤ δ) (hrad : δ + K * r ≤ r) :
    ∃! y, y ∈ closedBall x₀ r ∧ IsFixedPt f y
\end{leansrc}
\par\medskip\noindent{\small \texttt{Qtsl.QFixBall.}\texttt{\bfseries dist\_fixedPoint\_le}}\hfill{\footnotesize \path{lean/qtsl-aux/QtslAux/QFixBall.lean}, line 59}\par\nopagebreak
\begin{leansrc}
theorem dist_fixedPoint_le (hK : K < 1)
    (hlip : LipschitzOnWith K f (closedBall x₀ r))
    (hres : dist (f x₀) x₀ ≤ δ) (hrad : δ + K * r ≤ r)
    {y : α} (hy : y ∈ closedBall x₀ r) (hfy : IsFixedPt f y) :
    dist x₀ y ≤ δ / (1 - K)
\end{leansrc}
\par\medskip\noindent{\small \texttt{Qtsl.QFixBall.}\texttt{\bfseries anchored\_le\_radius}}\hfill{\footnotesize \path{lean/qtsl-aux/QtslAux/QFixBall.lean}, line 103}\par\nopagebreak
\begin{leansrc}
theorem anchored_le_radius (hK : K < 1)
    (hrad : δ + K * r ≤ r) : δ / (1 - K) ≤ r
\end{leansrc}
\par\medskip\noindent{\small \texttt{Qtsl.QFixBall.}\texttt{\bfseries sound\_QFixBall\_affine}}\hfill{\footnotesize \path{lean/qtsl-aux/QtslAux/QFixBall.lean}, line 120}\par\nopagebreak
\begin{leansrc}
theorem sound_QFixBall_affine {M : ℕ}
    (A : EuclideanSpace ℝ (Fin M) →L[ℝ] EuclideanSpace ℝ (Fin M))
    (c x₀ : EuclideanSpace ℝ (Fin M)) {K : ℝ≥0} {δ r : ℝ}
    (hK : K < 1) (hA : ‖A‖₊ ≤ K)
    (hres : dist (A x₀ + c) x₀ ≤ δ) (hrad : δ + K * r ≤ r) :
    ∃! y, y ∈ closedBall x₀ r ∧ IsFixedPt (fun x => A x + c) y
\end{leansrc}
\par\medskip\noindent{\small \texttt{Qtsl.Wasserstein.}\texttt{\bfseries sound\_QFixBall\_wasserstein}}\hfill{\footnotesize \path{lean/qtsl-aux/QtslAux/Wasserstein/Complete.lean}, line 449}\par\nopagebreak
\begin{leansrc}
theorem sound_QFixBall_wasserstein {T : P2 X → P2 X} {μ₀ : P2 X} {K : ℝ≥0}
    {δ r : ℝ} (hK : K < 1)
    (hlip : LipschitzOnWith K T (closedBall μ₀ r))
    (hres : dist (T μ₀) μ₀ ≤ δ) (hrad : δ + K * r ≤ r) :
    ∃! ν, ν ∈ closedBall μ₀ r ∧ IsFixedPt T ν
\end{leansrc}
\par\medskip\noindent{\small \texttt{Qtsl.QFixBall.}\texttt{\bfseries delta\_nonneg}}\hfill{\footnotesize \path{lean/qtsl-aux/QtslAux/QFixBall.lean}, line 34}\par\nopagebreak
\begin{leansrc}
theorem delta_nonneg (hres : dist (f x₀) x₀ ≤ δ) : 0 ≤ δ
\end{leansrc}
\par\medskip\noindent{\small \texttt{Qtsl.QFixBall.}\texttt{\bfseries radius\_nonneg}}\hfill{\footnotesize \path{lean/qtsl-aux/QtslAux/QFixBall.lean}, line 37}\par\nopagebreak
\begin{leansrc}
theorem radius_nonneg (hK : K < 1) (hres : dist (f x₀) x₀ ≤ δ)
    (hrad : δ + K * r ≤ r) : 0 ≤ r
\end{leansrc}
\par\medskip\noindent{\small \texttt{Qtsl.QFixBall.}\texttt{\bfseries mapsTo\_closedBall}}\hfill{\footnotesize \path{lean/qtsl-aux/QtslAux/QFixBall.lean}, line 45}\par\nopagebreak
\begin{leansrc}
theorem mapsTo_closedBall (hK : K < 1)
    (hlip : LipschitzOnWith K f (closedBall x₀ r))
    (hres : dist (f x₀) x₀ ≤ δ) (hrad : δ + K * r ≤ r) :
    Set.MapsTo f (closedBall x₀ r) (closedBall x₀ r)
\end{leansrc}
\par\medskip\noindent{\small \texttt{Qtsl.QFixBall.}\texttt{\bfseries fixedPoint\_unique}}\hfill{\footnotesize \path{lean/qtsl-aux/QtslAux/QFixBall.lean}, line 76}\par\nopagebreak
\begin{leansrc}
theorem fixedPoint_unique (hK : K < 1)
    (hlip : LipschitzOnWith K f (closedBall x₀ r))
    {y z : α} (hy : y ∈ closedBall x₀ r) (hz : z ∈ closedBall x₀ r)
    (hfy : IsFixedPt f y) (hfz : IsFixedPt f z) : y = z
\end{leansrc}
\par\medskip\noindent{\small \texttt{Qtsl.QFixBall.}\texttt{\bfseries lipschitzWith\_affine}}\hfill{\footnotesize \path{lean/qtsl-aux/QtslAux/QFixBall.lean}, line 113}\par\nopagebreak
\begin{leansrc}
theorem lipschitzWith_affine {E : Type*} [NormedAddCommGroup E]
    [NormedSpace ℝ E] (A : E →L[ℝ] E) (c : E) (hA : ‖A‖₊ ≤ K) :
    LipschitzWith K fun x => A x + c
\end{leansrc}
\par\medskip\noindent{\small \texttt{Qtsl.QDistTransport.}\texttt{\bfseries pushKernel\_iff\_eps\_zero}}\hfill{\footnotesize \path{lean/qtsl-aux/QtslAux/QDistTransport.lean}, line 183}\par\nopagebreak
\begin{leansrc}
theorem pushKernel_iff_eps_zero [DecidableEq β] {π : α → β}
    {K : α → α → ℝ} {Q : β → β → ℝ} :
    PushKernel π K Q ↔ EpsPushKernel π K Q 0
\end{leansrc}
\par\medskip\noindent{\small \texttt{Qtsl.QDistTransport.}\texttt{\bfseries epsPushKernel\_comp}}\hfill{\footnotesize \path{lean/qtsl-aux/QtslAux/QDistTransport.lean}, line 168}\par\nopagebreak
\begin{leansrc}
theorem epsPushKernel_comp [DecidableEq β] [DecidableEq γ]
    {π₁ : α → β} {π₂ : β → γ} {K₁ : α → α → ℝ} {K₂ : β → β → ℝ}
    {K₃ : γ → γ → ℝ} {ε₁ ε₂ : ℝ}
    (h₁ : EpsPushKernel π₁ K₁ K₂ ε₁) (h₂ : EpsPushKernel π₂ K₂ K₃ ε₂) :
    EpsPushKernel (π₂ ∘ π₁) K₁ K₃ (ε₁ + ε₂)
\end{leansrc}
\par\medskip\noindent{\small \texttt{Qtsl.QDistTransport.}\texttt{\bfseries sound\_QDistSelfTransport}}\hfill{\footnotesize \path{lean/qtsl-aux/QtslAux/QDistTransport.lean}, line 74}\par\nopagebreak
\begin{leansrc}
theorem sound_QDistSelfTransport [DecidableEq β] (π : α → β)
    {μ ν : α → ℝ} {ε : ℝ} (h : tvDist μ ν ≤ ε) :
    tvDist (push π μ) (push π ν) ≤ ε
\end{leansrc}
\par\medskip\noindent{\small \texttt{Qtsl.QRelChain.}\texttt{\bfseries IsMarkov}}\hfill{\footnotesize \path{lean/qtsl-aux/QtslAux/QRelChain.lean}, line 55}\par\nopagebreak
\begin{leansrc}
structure IsMarkov [Fintype α] (K : α → α → ℝ) : Prop where …
\end{leansrc}
\par\medskip\noindent{\small \texttt{Qtsl.QDistTransport.}\texttt{\bfseries sound\_QDistTransport\_eps}}\hfill{\footnotesize \path{lean/qtsl-aux/QtslAux/QDistTransport.lean}, line 152}\par\nopagebreak
\begin{leansrc}
theorem sound_QDistTransport_eps [DecidableEq β] {π : α → β}
    {K : α → α → ℝ} {Q : β → β → ℝ} {ε : ℝ} {μ₀ : α → ℝ}
    (hK : IsMarkov K) (hQ : IsMarkov Q) (hμ₀ : IsDist μ₀)
    (h : EpsPushKernel π K Q ε) (n : ℕ) :
    tvDist (push π (law K μ₀ n)) (law Q (push π μ₀) n) ≤ n * ε
\end{leansrc}
\par\medskip\noindent{\small \texttt{Qtsl.QDistTransport.}\texttt{\bfseries tvDist\_push\_le}}\hfill{\footnotesize \path{lean/qtsl-aux/QtslAux/QDistTransport.lean}, line 60}\par\nopagebreak
\begin{leansrc}
theorem tvDist_push_le [DecidableEq β] (π : α → β) (p q : α → ℝ) :
    tvDist (push π p) (push π q) ≤ tvDist p q
\end{leansrc}
\par\medskip\noindent{\small \texttt{Qtsl.QDistTransport.}\texttt{\bfseries tvDist\_step\_le}}\hfill{\footnotesize \path{lean/qtsl-aux/QtslAux/QDistTransport.lean}, line 86}\par\nopagebreak
\begin{leansrc}
theorem tvDist_step_le {K : α → α → ℝ} (hK : IsMarkov K) (p q : α → ℝ) :
    tvDist (step K p) (step K q) ≤ tvDist p q
\end{leansrc}
\par\medskip\noindent{\small \texttt{Qtsl.QDistTransport.}\texttt{\bfseries EpsPushKernel}}\hfill{\footnotesize \path{lean/qtsl-aux/QtslAux/QDistTransport.lean}, line 103}\par\nopagebreak
\begin{leansrc}
def EpsPushKernel [DecidableEq β] (π : α → β) (K : α → α → ℝ)
    (Q : β → β → ℝ) (ε : ℝ) : Prop
\end{leansrc}
\par\medskip\noindent{\small \texttt{Qtsl.QDistTransport.}\texttt{\bfseries push\_step\_defect}}\hfill{\footnotesize \path{lean/qtsl-aux/QtslAux/QDistTransport.lean}, line 107}\par\nopagebreak
\begin{leansrc}
theorem push_step_defect [DecidableEq β] {π : α → β} {K : α → α → ℝ}
    {Q : β → β → ℝ} {ε : ℝ} {μ : α → ℝ} (hμ : IsDist μ)
    (h : EpsPushKernel π K Q ε) :
    tvDist (push π (step K μ)) (step Q (push π μ)) ≤ ε
\end{leansrc}
\par\medskip\noindent{\small \texttt{Qtsl.QDistTransport.}\texttt{\bfseries tvDist\_triangle}}\hfill{\footnotesize \path{lean/qtsl-aux/QtslAux/QDistTransport.lean}, line 40}\par\nopagebreak
\begin{leansrc}
theorem tvDist_triangle {L : Type*} [Fintype L] (p q r : L → ℝ) :
    tvDist p r ≤ tvDist p q + tvDist q r
\end{leansrc}
\par\medskip\noindent{\small \texttt{SuFunctor.Universe.BelState.}\texttt{\bfseries push\_comp}}\hfill{\footnotesize \path{lean/su-functor/SuFunctor/Universe/TotalCitizenship.lean}, line 89}\par\nopagebreak
\begin{leansrc}
theorem BelState.push_comp (f : V → U) (g : U → W) (x : BelState V) :
    BelState.push (g ∘ f) x = BelState.push g (BelState.push f x)
\end{leansrc}
\par\medskip\noindent{\small \texttt{Qtsl.QDistTransport.}\texttt{\bfseries eq\_of\_tvDist\_le\_zero}}\hfill{\footnotesize \path{lean/qtsl-aux/QtslAux/QDistTransport.lean}, line 48}\par\nopagebreak
\begin{leansrc}
theorem eq_of_tvDist_le_zero {L : Type*} [Fintype L] {p q : L → ℝ}
    (h : tvDist p q ≤ 0) : p = q
\end{leansrc}
\par\medskip\noindent{\small \texttt{Qtsl.Acceptance.}\texttt{\bfseries towerCitation\_sound\_of\_wf}}\hfill{\footnotesize \path{lean/qtsl-aux/QtslAux/Acceptance.lean}, line 56}\par\nopagebreak
\begin{leansrc}
theorem towerCitation_sound_of_wf
    (C : ι → Row σ J ι) (baseTrue : σ → Prop) (sem : J → Prop)
    (SideOK : ι → Prop)
    (hbase : ∀ i s, C i = Row.base s → baseTrue s)
    (hside : ∀ i, SideOK i)
    (hstep : ∀ i prems concl, C i = Row.ext prems concl → SideOK i →
      (∀ j ∈ prems, Row.holds baseTrue sem (C j)) → sem concl)
    (hwf : WellFounded (cites C)) :
    ∀ i, Row.holds baseTrue sem (C i)
\end{leansrc}
\par\medskip\noindent{\small \texttt{Qtsl.Acceptance.}\texttt{\bfseries towerCitation\_sound}}\hfill{\footnotesize \path{lean/qtsl-aux/QtslAux/Acceptance.lean}, line 74}\par\nopagebreak
\begin{leansrc}
theorem towerCitation_sound [Finite ι]
    (C : ι → Row σ J ι) (baseTrue : σ → Prop) (sem : J → Prop)
    (SideOK : ι → Prop)
    (hbase : ∀ i s, C i = Row.base s → baseTrue s)
    (hside : ∀ i, SideOK i)
    (hstep : ∀ i prems concl, C i = Row.ext prems concl → SideOK i →
      (∀ j ∈ prems, Row.holds baseTrue sem (C j)) → sem concl)
    (hacyc : Acyclic C) :
    ∀ i, Row.holds baseTrue sem (C i)
\end{leansrc}
\par\medskip\noindent{\small \texttt{Qtsl.Acceptance.}\texttt{\bfseries towerCitation\_root}}\hfill{\footnotesize \path{lean/qtsl-aux/QtslAux/Acceptance.lean}, line 86}\par\nopagebreak
\begin{leansrc}
theorem towerCitation_root [Finite ι]
    (C : ι → Row σ J ι) (baseTrue : σ → Prop) (sem : J → Prop)
    (SideOK : ι → Prop)
    (hbase : ∀ i s, C i = Row.base s → baseTrue s)
    (hside : ∀ i, SideOK i)
    (hstep : ∀ i prems concl, C i = Row.ext prems concl → SideOK i →
      (∀ j ∈ prems, Row.holds baseTrue sem (C j)) → sem concl)
    (hacyc : Acyclic C)
    {Jroot : J} (hroot : ∃ i prems, C i = Row.ext prems Jroot) :
    sem Jroot
\end{leansrc}
\par\medskip\noindent{\small \texttt{Qtsl.Acceptance.Row.}\texttt{\bfseries cited}}\hfill{\footnotesize \path{lean/qtsl-aux/QtslAux/Acceptance.lean}, line 36}\par\nopagebreak
\begin{leansrc}
def Row.cited : Row σ J ι → List ι
  | .base _ => []
  | .ext prems _ => prems

def Row.holds (baseTrue : σ → Prop) (sem : J → Prop) : Row σ J ι → Prop
  | .base s => baseTrue s
  | .ext _ concl => sem concl

def cites (C : ι → Row σ J ι) (j i : ι) : Prop
\end{leansrc}
\par\medskip\noindent{\small \texttt{Qtsl.Acceptance.}\texttt{\bfseries acyclic\_wellFounded}}\hfill{\footnotesize \path{lean/qtsl-aux/QtslAux/Acceptance.lean}, line 50}\par\nopagebreak
\begin{leansrc}
theorem acyclic_wellFounded [Finite ι] (C : ι → Row σ J ι)
    (hacyc : Acyclic C) : WellFounded (cites C)
\end{leansrc}
\par\medskip\noindent{\small \texttt{Qtsl.Acceptance.}\texttt{\bfseries citePos\_sound}}\hfill{\footnotesize \path{lean/qtsl-aux/QtslAux/Acceptance.lean}, line 101}\par\nopagebreak
\begin{leansrc}
theorem citePos_sound {X : Type*} {R D : Set X} {P : Set X → Prop}
    (hmono : Monotone P) (hsub : R ⊆ D) (hR : P R) : P D
\end{leansrc}
\par\medskip\noindent{\small \texttt{Qtsl.Acceptance.}\texttt{\bfseries citeCmpl\_subset}}\hfill{\footnotesize \path{lean/qtsl-aux/QtslAux/Acceptance.lean}, line 105}\par\nopagebreak
\begin{leansrc}
theorem citeCmpl_subset {X : Type*} {R D : Set X} (hsub : R ⊆ D) :
    Dᶜ ⊆ Rᶜ
\end{leansrc}
\par\medskip\noindent{\small \texttt{Qtsl.Acceptance.}\texttt{\bfseries citeCmpl\_sound}}\hfill{\footnotesize \path{lean/qtsl-aux/QtslAux/Acceptance.lean}, line 109}\par\nopagebreak
\begin{leansrc}
theorem citeCmpl_sound {X : Type*} {R D : Set X} {P : Set X → Prop}
    (hanti : Antitone P) (hsub : R ⊆ D) (hR : P Rᶜ) : P Dᶜ
\end{leansrc}
\par\medskip\noindent Modules cited: \texttt{QtslAux.QGuard}, \texttt{QtslAux.LayerD.MPDischarge}, \texttt{SuFunctor.Site.FellAcceptance}, \texttt{SuFunctor.Site.FellBehaviorSpec}, \texttt{QtslAux.Enclosure}.
\par\medskip\noindent Also cited in this chapter, listed under an earlier one: \texttt{ProofObject.sound}, \texttt{SuFunctor.Plain.CheckerFace}, \texttt{semanticallyValid\_iff\_face}.

\section{Partial Completeness}
Citations from Chapter~\ref{ch:qtsl-partial-completeness}.

\par\medskip\noindent{\small \texttt{PPLC.Qtsl.}\texttt{\bfseries sound\_u10f}}\hfill{\footnotesize \path{lean/qtsl/Qtsl/TimedCore.lean}, line 49}\par\nopagebreak
\begin{leansrc}
theorem sound_u10f (B : Behavior X) (T : ℝ≥0) (ν χ : Set X) (x : X)
    (hbox : B.BoxInv ν x) :
    B.AService T (χ ∩ ν) x ↔ B.AService T χ x
\end{leansrc}
\par\medskip\noindent{\small \texttt{PPLC.Qtsl.}\texttt{\bfseries sound\_firstService}}\hfill{\footnotesize \path{lean/qtsl/Qtsl/TimedCore.lean}, line 138}\par\nopagebreak
\begin{leansrc}
theorem sound_firstService (hB3 : B.FusionClosed) (hn : 0 < n)
    (α : Fin n → Set X) (T : Fin n → ℝ≥0) (x : X)
    (hatt : ∀ i, B.FirstVisits (α i))
    (hx : ∀ i, B.AService (T i) (α i) x) :
    B.AService (Finset.univ.inf' (Finset.univ_nonempty_iff.mpr ⟨⟨0, hn⟩⟩) T)
      (B.fsTarget α T) x
\end{leansrc}
\par\medskip\noindent{\small \texttt{PPLC.Qtsl.}\texttt{\bfseries sound\_fsResidual}}\hfill{\footnotesize \path{lean/qtsl/Qtsl/TimedCore.lean}, line 149}\par\nopagebreak
\begin{leansrc}
theorem sound_fsResidual (hB3 : B.FusionClosed) (hn : 0 < n)
    (α : Fin n → Set X) (T : Fin n → ℝ≥0) (c : ℝ≥0) (x : X)
    (hc : ∀ i, c < T i)
    (hatt : ∀ i, B.FirstVisits (α i))
    (hx : ∀ i, B.AService (T i) (α i) x) :
    B.AService c (B.fsTarget α T ∪ B.residualCell α T c) x
\end{leansrc}
\par\medskip\noindent{\small \texttt{PPLC.Qtsl.TimedCoreCountermodel.}\texttt{\bfseries firstService\_attainment\_necessary}}\hfill{\footnotesize \path{lean/qtsl/Qtsl/TimedCore.lean}, line 496}\par\nopagebreak
\begin{leansrc}
theorem firstService_attainment_necessary :
    ∃ (S : Type) (B : Behavior S) (α : Fin 2 → Set S) (x : S),
      B.FusionClosed ∧
      (∀ i, B.AService 2 (α i) x) ∧
      ¬ B.AService 2 (B.fsTarget α (fun _ => 2)) x
\end{leansrc}

